%% 第5巻　変化をつかむ：微分と積分　本文（第4回編集委員会で完成）
%% 序章 p.8〜11、第1部 p.12〜113、第2部 p.114〜175、第3部 p.176〜253、巻末 p.254〜
%% 5-0 序章　スピードメーターは何を計算している？（第11回：6ページ）
\makeatletter
\gdef\bk@partlabel{第5巻}\gdef\bk@partname{変化をつかむ}
\makeatother
\bkchapter{5}{0}{スピードメーターは\\何を計算している？}{\bklevel{2}}{スピードメーターは何を計算している？}

\noindent 車のスピードメーターの針が「50」を指しています。時速\qty{50}{km}です。でも、この車は1時間も走っていません。走りはじめて、まだ10秒です。

\begin{bkfigure}
\begin{tikzpicture}[x=1mm, y=1mm]
  % スピードメーター（0〜180 km/h）
  \fill[bkInk!88] (0,0) circle (34);
  \fill[bkInk!75] (0,0) circle (31);
  \foreach \v in {0,10,...,180} {
    \pgfmathsetmacro\ang{210-\v*4/3}
    \pgfmathtruncatemacro\isbig{mod(\v,20)==0}
    \ifnum\isbig=1
      \draw[white, line width=1.2pt] (\ang:24.5) -- (\ang:29.5);
      \node[white, font=\sffamily\bfseries\footnotesize] at (\ang:19.5) {\v};
    \else
      \draw[white, line width=0.6pt] (\ang:26.5) -- (\ang:29.5);
    \fi}
  \pgfmathsetmacro\need{210-50*4/3}
  \draw[bkOrange, line width=2.2pt, line cap=round] (0,0) -- (\need:27);
  \fill[bkOrange!70!black] (0,0) circle (3);
  \node[white, font=\sffamily\footnotesize] at (0,-12) {km/h};
  % 説明
  \node[anchor=west, text width=58mm, align=left, font=\small] at (42,11)
    {針は、いまの速さ「時速\qty{50}{km}」を指している。\\[2pt]
     でも、時速\qty{50}{km}は「1時間に\qty{50}{km}進む速さ」。この車は、まだ1時間走っていない。};
  \node[anchor=west, text width=58mm, align=left, font=\sffamily\bfseries\small, text=bkDeep] at (42,-18)
    {走っていない1時間の距離を、\\メーターはどうやって知ったのだろう？};
\end{tikzpicture}
\end{bkfigure}

\bksection{doguu}{「時速\qty{50}{km}」は、何を言っているの？}

時速\qty{50}{km}は「1時間あたり\qty{50}{km}」という意味で、\qty{50}{km/h}と書きます（h は hour の頭の文字で「時間」）。でも、この車は1時間走ったわけではありません。メーターが言っているのは、次のことです。

\begin{taisetsu}
\textbf{もし、いまの速さのまま1時間走りつづけたら、\qty{50}{km}進む。}
\end{taisetsu}

「いまの速さのまま」なら、1時間より短い時間でも、進む距離が計算できます。30分は1時間の半分なので、進む距離も半分の\qty{25}{km}です。1分では、1時間が60分（\qty{60}{min/h}、min は minute の頭の文字で「分」）なので、
\[ 50\,\mathrm{km/h}\div60\,\mathrm{min/h}=0.833\cdots\,\mathrm{km/min} \]
h どうしが約分されて消え、「1分あたり約\qty{0.83}{km}」（約\qty{833}{m}）になります。

\begin{bkfigure}
\begin{tikzpicture}[x=1mm, y=1mm]
  % 道路と、いまの車。点線の車は「いまの速さのまま」走ったときの位置（はばは目安で、長さの比ではない）
  \fill[bkGray!18] (-4,-2.5) rectangle (128,0);
  \draw[bkGray, line width=0.8pt] (-4,0) -- (128,0);
  \bkcar{(0,0)}{0.8}{bkOrange}
  \node[anchor=south, font=\sffamily\scriptsize, text=bkOrange!80!black] at (9.6,14.5) {いま};
  \draw[bk arrow, bkBlue, dash pattern=on 3pt off 2pt] (26,7) -- (124,7);
  \node[anchor=south, font=\sffamily\footnotesize, text=bkBlue] at (75,8.5) {いまの速さ（時速\qty{50}{km}）のまま、走りつづけたら……};
  \foreach \x/\lab in {40/1秒後, 64/10秒後, 88/1分後, 112/1時間後} {
    \draw[bkGray, line width=0.7pt] (\x,0) -- (\x,-4);
    \node[anchor=north, font=\sffamily\scriptsize] at (\x,-4.8) {\lab};}
\end{tikzpicture}
\end{bkfigure}

\begin{bkcenter}
\begin{bktable}{colspec={Q[l,m,wd=28mm] *{5}{Q[c,m,wd=17mm]}}, rowsep=1.5pt, colsep=3pt}
走りつづけた時間 & 1秒 & 10秒 & 1分 & 30分 & 1時間 \\
進む距離 & 約\qty{14}{m} & 約\qty{139}{m} & 約\qty{833}{m} & \qty{25}{km} & \qty{50}{km} \\
\end{bktable}
\bktabcaption{時速\qty{50}{km}のまま走ると、どれだけ進むか}
\end{bkcenter}

1秒に進む距離は、単位つきの式で計算します（2-10）。\qty{1}{km}は\qty{1000}{m}なので「1 km あたり 1000 m」で \qty{1000}{m/km}、1時間は $60\times60=3600$ 秒なので「1時間あたり3600秒」で \qty{3600}{s/h} です（s は second の頭の文字で「秒」）。
\[ 50\,\mathrm{km/h}\times1000\,\mathrm{m/km}\div3600\,\mathrm{s/h}=13.88\cdots\,\mathrm{m/s} \]
km どうし、h どうしが約分されて消え、m/s が残ります。だから1秒では約\qty{14}{m}進みます。くわしい計算のしかたは、5-1でもう一度ていねいに見ます。

%% ---------------- p.2 瞬間の速さ ----------------
\bksection{doguu}{「いまの速さ」は、どうやって計算する？}

ここで、こまったことが起きます。2-3で学んだ速さの式は、
\[ \text{速さ}=\text{進んだ道のり}\div\text{かかった時間} \]
でした。この式は、「10秒のあいだに\qty{100}{m}進んだ」のように、\textbf{ある時間のあいだ}の速さです。このあいだに速くなったり、おそくなったりしても、ならした値になります。だから\textbf{平均の速さ}といいます。\bkidx{はやさ}{速さ}

スピードメーターが表示しているのは、「いま、この瞬間」の速さです。瞬間とは、時間の長さが0の「時刻の1点」のことです。時間の長さが0だと、進んだ道のりも0です。すると、式は
\[ 0\div0 \]
になってしまいます。0でわる計算はできません（1-14）。

\begin{bkfigure}
\begin{tikzpicture}[x=1mm, y=1mm]
  % 時間の線。上：あいだ（平均の速さ）、下：瞬間（時刻の1点）
  \draw[bk arrow, bkInk, line width=0.8pt] (0,14) -- (122,14);
  \node[anchor=west, font=\sffamily\scriptsize] at (122.5,14) {時間};
  \fill[bkBlue!25] (24,11.8) rectangle (84,16.2);
  \draw[bkBlue, line width=0.8pt] (24,11.8) rectangle (84,16.2);
  \node[anchor=south, font=\sffamily\footnotesize, text=bkBlue] at (54,17.5) {あいだ（たとえば10秒）};
  \node[anchor=north, font=\sffamily\footnotesize, text=bkInk, align=center] at (54,10) {進んだ道のり÷かかった時間 で計算できる（平均の速さ）};
  \draw[bk arrow, bkInk, line width=0.8pt] (0,-6) -- (122,-6);
  \node[anchor=west, font=\sffamily\scriptsize] at (122.5,-6) {時間};
  \fill[bkOrange] (54,-6) circle (1.6);
  \node[anchor=south, font=\sffamily\footnotesize, text=bkOrange!80!black] at (54,-3) {瞬間（時刻の1点。長さ0）};
  \node[anchor=north, font=\sffamily\footnotesize, text=bkInk] at (54,-9.5) {道のり0÷時間0 になって、このままでは計算できない};
\end{tikzpicture}
\end{bkfigure}

それでもメーターは、「いまの速さ」を表示しています。どうやっているのでしょう。答えは5-1にあります。ヒントは、「0にする」のではなく、「とても短いあいだ」を考えることです。

この「瞬間の変化の速さ」を計算する道具が\textbf{微分}（びぶん）です。そして反対に、速さから進んだ距離を求める道具が\textbf{積分}（せきぶん）です。この巻では、この2つの道具を、表計算とグラフを使って、少しずつ組み立てていきます。

\begin{taisetsu}
\textbf{微分}：いま、どれだけの速さで変わっているか（瞬間の変化の速さ）\\
\textbf{積分}：少しずつの変化をたし合わせて、全部でどれだけになったか
\end{taisetsu}

\bkfukidashi{mina}{0÷0 はできないのに、メーターは計算している。ふしぎだね。}

%% ---------------- p.3 おふろで見る微分と積分 ----------------
\bksection{doguu}{おふろの水で見る、微分と積分}

微分と積分は、車の速さだけの道具ではありません。もう1つ、おふろに水をためる場面で考えてみましょう。

じゃ口を少しずつ開けながら、おふろに水を入れました。1分ごとに、たまった水の量を調べると、次のようになりました。

\begin{bkfigure}
\begin{tikzpicture}[x=1mm, y=1mm]
  % 左：浴そう。1分ごとに入った水を、色を変えて積み重ねる（6, 9, 12, 12, 12 L。高さは 1 L＝0.6 mm）
  \def\yy{0}
  \foreach \v/\c/\k in {6/bkBlue!18/1, 9/bkBlue!30/2, 12/bkBlue!42/3, 12/bkBlue!54/4, 12/bkBlue!66/5} {
    \pgfmathsetmacro\ya{\yy}\pgfmathsetmacro\yb{\yy+\v*0.6}
    \fill[\c] (2,\ya) rectangle (34,\yb);
    \node[font=\sffamily\scriptsize, text=bkInk] at (18,{(\ya+\yb)/2}) {\k 分めに \v\,L};
    \xdef\yy{\yb}}
  \draw[bkInk, line width=1.1pt, rounded corners=1.5mm] (0,38) -- (0,-1.2) -- (36,-1.2) -- (36,38);
  % じゃ口（右から浴そうの上へのびる管と、下向きの口）
  \fill[bkGray] (24,44) rectangle (40,47);
  \fill[bkGray] (24,40) rectangle (27.5,47);
  \draw[bkBlue!60, line width=1.2pt] (25.75,39.5) -- (25.75,31.5);
  \node[anchor=north, font=\sffamily\footnotesize] at (18,-3) {5分で\qty{51}{L}たまった};
  % 右：表。ふえた量は、となりの列のあいだに書く。2〜3分の1組だけ、ひき算の線をかく
  \begin{scope}[shift={(46,0)}]
    \foreach \t [count=\i from 0] in {0,1,2,3,4,5} {\node[font=\sffamily\small] at (26+\i*11,36) {\t};}
    \foreach \v [count=\i from 0] in {0,6,15,27,39,51} {\node[font=\sffamily\small] (w\i) at (26+\i*11,26) {\v};}
    \foreach \v [count=\i from 0] in {6,9,12,12,12} {\node[font=\sffamily\small, text=bkOrange!80!black] (d\i) at (31.5+\i*11,6) {\v};}
    \node[anchor=east, font=\sffamily\scriptsize, align=right] at (20,36) {時刻（分）};
    \node[anchor=east, font=\sffamily\scriptsize, align=right] at (20,26) {たまった量（L）};
    \node[anchor=east, font=\sffamily\scriptsize, align=right] at (20,6) {1分ごとに\\ふえた量（L）};
    \draw[bkLine, line width=0.5pt] (21,31) -- (86,31) (21,21) -- (86,21) (21,11) -- (86,11);
    \draw[bkOrange!80!black, line width=0.6pt] (w2.south) -- (d2.north) (w3.south) -- (d2.north);
    \node[anchor=west, font=\sffamily\scriptsize, text=bkOrange!80!black, inner sep=0.8pt] at (59.5,15.5) {$27-15=12$};
  \end{scope}
\end{tikzpicture}
\bkfigcaption{左：1分ごとに入った水が、積み重なってたまる。右：たまった量と、1分ごとにふえた量。}
\end{bkfigure}

\textbf{ひき算で、ふえ方がわかる}\par
たまった量の表から、となりどうしのひき算（あと${}-{}$まえ、1-5・2-3）をすると、1分ごとにふえた量がわかります。たとえば、2分から3分までは $27-15=12$ で、\qty{12}{L}ふえました。これは「1分あたり\qty{12}{L}」という、\textbf{水がたまる速さ}です。この「ふえ方（変化の速さ）を求める」考えを、とても細かい時間まで進めたものが\textbf{微分}です。

\textbf{たし算で、全部の量がわかる}\par
反対に、1分ごとにふえた量をはじめから順にたしていくと、たまった量がわかります。
\[ 6+9+12+12+12=51 \]
で、5分で\qty{51}{L}です。この「少しずつの量をたし合わせて全部を求める」考えを、とても細かい時間まで進めたものが\textbf{積分}です。

ここでは1分ごとに区切りました。でも、ほんとうは1分のあいだにも、水の入り方は少しずつ変わっています。区切りを0.1分、0.01分……と、どんどん細かくしていくと、微分と積分になります。微分はひき算のなかま、積分はたし算のなかまで、この2つはおたがいに逆の計算です（5-11）。

%% ---------------- p.4 この巻が通る路線 ----------------
\bksection{doguu}{この巻が通る路線と駅}

この巻の章は、4本の路線の駅です。前の巻で乗った路線が、この巻で「微分」と「積分」という大きな駅に着きます。

\begin{bkfigure}
\begin{tikzpicture}[x=1mm, y=1mm]
  \foreach \r/\y/\nm/\stA/\stB/\stC in {
      1/60/変化/{1-5・2-3・3-5}/{5-1〜5-8・5-17・5-18}/{6-1（微分方程式）},
      2/42/たし合わせ/{2-3・4-5}/{5-9〜5-13・5-16}/{6-3・8-11},
      3/24/倍々/{2-7・3-11}/{5-14}/{6-4・9-5},
      4/6/回転と波/{4-2・4-8}/{5-15}/{6-10・9-1}} {
    \draw[bkRoute\r, line width=3pt, line cap=round] (8,\y) -- (132,\y);
    \node[inner sep=0pt] at (4,\y) {\bkrouteicon[6.5mm]{\r}};
    \node[anchor=south west, font=\sffamily\bfseries\small, text=bkInk] at (10,\y+1.8) {\nm};
    \fill[white] (34,\y) circle (2.2);\draw[bkRoute\r, line width=1.1pt] (34,\y) circle (2.2);
    \node[anchor=north, font=\sffamily\scriptsize, text=bkGray] at (34,\y-3) {\stA};
    \fill[white, rounded corners=1.5mm] (56,\y-4.2) rectangle (104,\y+4.2);
    \draw[bkRoute\r, line width=1.6pt, rounded corners=1.5mm] (56,\y-4.2) rectangle (104,\y+4.2);
    \node[font=\sffamily\bfseries\scriptsize, text=bkInk] at (80,\y) {\stB};
    \fill[white] (124,\y) circle (2.2);\draw[bkRoute\r, line width=1.1pt] (124,\y) circle (2.2);
    \node[anchor=north, font=\sffamily\scriptsize, text=bkGray] at (124,\y-3) {\stC};}
  \node[font=\sffamily\bfseries\small, text=bkGray] at (34,72) {前の巻};
  \node[font=\sffamily\bfseries\small, text=bkDeep] at (80,72) {この巻（第5巻）};
  \node[font=\sffamily\bfseries\small, text=bkGray] at (124,72) {この先};
\end{tikzpicture}
\bkfigcaption{第5巻を通る4本の路線。四角の中がこの巻の駅。}
\end{bkfigure}

\begin{bkcenter}
\begin{bktable}{colspec={Q[c,m,wd=8mm] Q[l,m,wd=30mm] Q[l,m,wd=82mm]}, rowsep=2pt}
 & 路線 & この巻で着く駅 \\
\bkrouteicon{1} & 変化 & ひき算（1-5）→ 傾き（3-5）→ \textbf{微分}。瞬間の変化の速さ \\
\bkrouteicon{2} & たし合わせ & かけ算の面積（2-3）→ \textbf{積分}。細い長方形をたし合わせる \\
\bkrouteicon{3} & 倍々 & 指数関数（3-11）→ \textbf{e}。微分しても形が変わらない関数 \\
\bkrouteicon{4} & 回転と波 & 三角関数（4-8）→ \textbf{sin・cosの微分}。波は微分すると1/4周期ずれる \\
\end{bktable}
\end{bkcenter}

変化の路線（\bkrouteicon[0.9em]{1}）とたし合わせの路線（\bkrouteicon[0.9em]{2}）は、5-11で1つにつながります。おふろの水で見た「ひき算とたし算は逆の計算」が、「積分は微分の逆」という、この巻でいちばん大事な発見になります。

%% ---------------- p.5 この巻でできるようになること ----------------
\bksection{goal}{この巻でできるようになること}

\begin{bkfigure}
\begin{tikzpicture}[x=1mm, y=1mm,
  q/.style={draw=bkVol, line width=1.2pt, fill=white, rounded corners=2.5mm, minimum width=34mm, minimum height=17mm, align=center, font=\sffamily\small}]
  \node[q] (x) at (0,0) {\textbf{位置} $x$\\[1pt]{\footnotesize どこにいる？}};
  \node[q] (v) at (52,0) {\textbf{速度} $v$\\[1pt]{\footnotesize どれだけ速く動く？}};
  \node[q] (a) at (104,0) {\textbf{加速度} $a$\\[1pt]{\footnotesize 速さがどう変わる？}};
  \draw[bk arrow, bkOrange] ([yshift=3mm]x.east) -- node[above, font=\sffamily\bfseries\small] {微分} ([yshift=3mm]v.west);
  \draw[bk arrow, bkOrange] ([yshift=3mm]v.east) -- node[above, font=\sffamily\bfseries\small] {微分} ([yshift=3mm]a.west);
  \draw[bk arrow, bkBlue] ([yshift=-3mm]v.west) -- node[below, font=\sffamily\bfseries\small] {積分} ([yshift=-3mm]x.east);
  \draw[bk arrow, bkBlue] ([yshift=-3mm]a.west) -- node[below, font=\sffamily\bfseries\small] {積分} ([yshift=-3mm]v.east);
  \node[font=\footnotesize, text=bkGray] at (26,-15) {5-9〜5-13};
  \node[font=\footnotesize, text=bkGray] at (26,15) {5-1〜5-5};
  \node[font=\footnotesize, text=bkGray] at (78,15) {5-5};
  \node[font=\footnotesize, text=bkGray] at (78,-15) {5-13};
\end{tikzpicture}
\bkfigcaption{微分は右へ、積分は左へ。2つは逆向きの道具。}
\end{bkfigure}

位置・速度・加速度の3つは、「微分」で右へ、「積分」で左へ移れます。たとえば、位置が時間とともにどう変わるかがわかれば、微分して速度がわかります。反対に、速度がわかれば、積分して位置がわかります。おふろの水なら、「たまった量」が位置、「たまる速さ」が速度にあたります。

\begin{bkcenter}
\begin{bktable}{colspec={Q[l,m,wd=20mm] Q[l,m,wd=58mm] Q[l,m,wd=42mm]}, rowsep=2pt}
部 & できるようになること & 物理・工学の例 \\
第1部 微分 & 瞬間の速さを計算する。グラフの傾きを式で求める。いちばん高い所を見つける & スピードメーター、投げ上げたボールの最高点 \\
第2部 積分 & 速さのグラフの面積から進んだ距離を求める。たまった量を計算する & ばねの仕事、水の量、電気の使用量 \\
第3部 特別な関数 & e・sin・cos を微分・積分する。むずかしい関数を多項式で近似する & 放射線の減り方、ばねの振動、振り子 \\
\end{bktable}
\end{bkcenter}

\begin{column}{なぜ、微分と積分を学ぶの？}
物理の法則の多くは、「いまの状態が、次の瞬間の変化を決める」という形で書かれています。たとえば「落ちる物の速さは、1秒ごとに\qty{9.8}{m/s}ずつ増える」。この「変化の速さ」を書くことばが微分で、変化を積み重ねて全体を計算する道具が積分です。この巻の道具がそろうと、次の第6巻で「未来を計算する」微分方程式に進めます。\bkidx{びぶんほうていしき}{微分方程式}
\end{column}

%% ---------------- p.6 この巻の読み方 ----------------
\bksection{reidai}{この巻の読み方}

\bkfukidashi{sora}{微分って、むずかしそう。でも「平均の速さ」なら2-3でやったよ。}
\bkfukidashi{delta}{この巻では、まず表計算で数を出して、きまりを見つけてから式にするよ。手を動かせば、ちゃんとわかるよ。}

\bksubsection{1}{章のはじめの「思い出そう」}

この巻の章には、「思い出そう」の囲みがあります。その章で使う、前の巻で学んだことをまとめたものです。札の番号の章にもどれば、くわしく書いてあります。
\begin{omoidasou}
\bkrecall{2-3}{速さ＝進んだ道のり÷かかった時間。「変わった量」は \mbox{$\Delta$（デルタ）}で書く。}
\bkrecall{3-5}{グラフの傾き＝$\dfrac{y\text{の変化量}}{x\text{の変化量}}=\dfrac{\Delta y}{\Delta x}$}
\end{omoidasou}

\bksubsection{2}{表計算とノートを用意しよう}

表計算（2-18）をたくさん使います。「やってみよう」の表は、コンピュータで作ってみよう。問題はノートに解きます。電卓も使ってかまいません。

\bksubsection{3}{式の変形は、1行ずつ}

式の変形は1行ずつ書いてあります。どの行も、前の行に「展開する」「両辺を同じ数でわる」などのきまり（2-14・3-8）を1つ使っただけです。わからない行は、前の行とくらべてみよう。

\bksubsection{4}{単位の換算は、単位つきの式で}

時速と秒速、分と秒のように単位をなおすときは、\qty{1000}{m/km}（1 km あたり 1000 m）、\qty{60}{s/min}（1分あたり60秒）、\qty{3600}{s/h}（1時間あたり3600秒）のような「単位つきの数」をかけたりわったりします（2-10）。同じ単位どうしは約分されて消え、ほしい単位だけが残ります。単位が正しく残ったかどうかで、計算のまちがいにも気づけます。

\bksubsection{5}{わからなくなったら}

はじめての考えは、1回読んですぐにわからなくても大丈夫です。例題を、ノートに1行ずつ写しながら、自分の手で計算してみよう。どの行でわからなくなったかが見つかれば、その行の「（2-14）」のような番号の章にもどって確かめればよいのです。


%% ======================== 第1部 微分 ========================
\bkPartOpener{1}{第5巻\hspace{0.8em}第1部}{微分}%
  {172/{スピードメーターの「いまの速さ」は、どう決める？},
   160/{投げ上げたボールがいちばん高くなるのは、いつ？},
   148/{次の瞬間の位置は、どうやって予想する？}}%
  {% 曲線と接線（白い線）
   \begin{scope}[white, line width=1.2pt, line cap=round]
     \draw (118,116) parabola bend (118,116) (168,156);
     \draw[line width=0.9pt] (128,113) -- (172,149);
     \fill (150,131) circle (1.6);
     \draw[line width=0.7pt, dash pattern=on 2pt off 1.5pt] (150,131) -- (162,131) -- (162,140.8);
     \node[white, font=\sffamily\bfseries\small] at (165,128) {傾き};
   \end{scope}}%
  {この部では、時間をどんどん短くして、「いまの瞬間」の変化の速さを計算する道具をつくるよ。それが微分。スピードメーターも、ロケットの制御も、この道具で動いているよ。}
\bkPartContents{1}{%% 自動生成（gen.py）
\bktocchap{5-1}{平均の速さと瞬間の速さ}{2時間で100 km走った車は、ずっと時速50 km？}{1}{2}{2}{18}{C}
\bktocchap{5-2}{時間をどんどん短くしてみよう}{0.1秒、0.01秒…と縮めると、ある値に近づく}{1}{2}{2}{40}{C}
\bktocchap{5-3}{かぎりなく近づく：極限と接線}{曲線にぴったりふれる直線の傾き}{1}{3}{3}{60}{C}
\bktocchap{5-4}{表から導関数を見つけよう}{表をながめると、微分のきまりが見えてくる}{1}{3}{3}{78}{C}
\bktocchap{5-5}{位置→速度→加速度}{2回微分すると、何がわかる？}{1}{3}{3}{104}{C}
\bktocchap{5-6}{微分の道具箱：積と合成関数}{歯車をつなぐと、回る速さはかけ算になる}{1}{4}{4}{126}{C}
\bktocchap{5-7}{いちばん高いのはどこ？　最大と最小}{いちばん高い点では、速度が0になる}{1}{3}{3}{154}{C}
\bktocchap{5-8}{短い時間なら、かけ算で近似できる}{「次の位置」を少しずつ予想する}{1}{2}{2}{174}{C}
}%
{「いまの速さ」を、グラフの傾きとして計算できるようになります。速度・加速度の関係、いちばん高い所の見つけ方、次の瞬間の予想のしかたは、第6巻の微分方程式の土台になります。}
\bkpartbegin{1}{微分}
%% 5-1 平均の速さと瞬間の速さ（第11回：ページ制限なし。偶数ページ）
\bkchapter{5}{1}{平均の速さと瞬間の速さ}{\bklevel{2}}{平均の速さと瞬間の速さ}
\bkmichishirube{2-3,3-4,3-5}{5-2,5-3}

\begin{goal}
\begin{itemize}[label={\color{bkTheme}●}, leftmargin=1.6\zw, itemsep=0pt, topsep=1pt]
\item \bklevel{2}\enspace 平均の速さを、表やグラフから求められる。
\item \bklevel{2}\enspace 平均の速さが、グラフの「2点を結ぶ直線の傾き」になることを説明できる。\bkidx{かたむき}{傾き}
\item \bklevel{2}\enspace 時間を短くして、瞬間の速さを見積もれる。
\end{itemize}
\end{goal}

\begin{hatena}
\qty{100}{km}はなれた町まで、車で2時間。$100\div2=50$ で時速\qty{50}{km}です。でも、ずっと時速\qty{50}{km}で走っていたのでしょうか？\par
\begin{bkfigure}
\begin{tikzpicture}
% 第7回：区間をはっきり分けた折れ線にする（0〜0.5h 時速60、0.5〜0.7h 止まる、0.7〜1h 時速50、
%   1〜1.25h じゅうたい時速20、1.25〜1.5h 時速80、1.5〜2h 時速60。合計100 km）。
%   0.5時間ごとの点は例題1の表と同じ。注記は帯の上と右下の凡例に置き、引き出し線は使わない
\begin{axis}[bkplot, width=100mm, height=50mm, xmin=0, xmax=2.2, ymin=0, ymax=112,
  xtick={0.5,1,1.5,2}, ytick={25,50,75,100}, xlabel={時間（h）}, ylabel={進んだ距離（km）},
  ylabel style={anchor=south east, xshift=3mm}]
  \fill[bkRed!12] (axis cs:0.5,0) rectangle (axis cs:0.7,104);
  \fill[bkOrange!16] (axis cs:1.0,0) rectangle (axis cs:1.25,104);
  \node[font=\sffamily\scriptsize, text=bkRed!80!black, anchor=south, align=center, inner sep=1pt] at (axis cs:0.6,104) {止まった\\（休けい）};
  \node[font=\sffamily\scriptsize, text=bkOrange!80!black, anchor=south, align=center, inner sep=1pt] at (axis cs:1.125,104) {じゅうたい\\（のろのろ）};
  \addplot[bkthin, dash pattern=on 4pt off 2pt, domain=0:2] {50*x};
  \addplot[bkcurve, line join=round] coordinates {(0,0) (0.5,30) (0.7,30) (1.0,45) (1.25,50) (1.5,70) (2,100)};
  \addplot[only marks, mark=*, mark size=1.5pt, bkInk] coordinates {(0.5,30) (1,45) (1.5,70) (2,100)};
  \node[on layer=axis foreground, anchor=south east, fill=white, draw=bkLine, line width=0.4pt, rounded corners=1mm, inner sep=1.2mm,
    font=\sffamily\scriptsize, text=bkInk, align=left] at (axis cs:2.17,4)
    {\tikz[baseline=-0.6ex]\draw[bkBlue, line width=1.7pt] (0,0)--(5mm,0);\enspace じっさいの車\\[1pt]
     \tikz[baseline=-0.6ex]\draw[bkGray, line width=0.9pt, dash pattern=on 4pt off 2pt] (0,0)--(5mm,0);\enspace ずっと時速\qty{50}{km}};
\end{axis}
\end{tikzpicture}
\end{bkfigure}
水平な所は、信号や休けいで止まっていた所（メーターは0）。傾きがゆるい所は、じゅうたいの所（メーターは20くらい）。「時速\qty{50}{km}」と、スピードメーターが指す速さは、何がちがうのでしょう。
\end{hatena}

\clearpage
%% ---------------- やってみよう・思い出そう ----------------
\begin{yattemiyou}[坂をころがるボールの速さを調べよう]
\textbf{用意するもの}\par
長い板（または本を重ねてつくった坂）、ボール、ものさし、スマートフォン（動画をとる）
\begin{enumerate}[label={\bkcircled[bkGreen]{\arabic*}}, leftmargin=2.4\zw, itemsep=1pt, topsep=3pt]
\item 板を少しだけ傾け、ものさしを横に置く。ボールを静かに放して、ころがるようすを動画にとる。
\item 動画を0.5秒ずつ止めて、ボールの位置を読み、表に書く。
\item となりどうしの時刻のあいだで、「進んだ距離÷かかった時間」を計算しよう。
\end{enumerate}
\begin{bkcenter}
\begin{bktable}{colspec={Q[l,m,wd=34mm] *{5}{Q[c,m,wd=13mm]}}, rowsep=1.5pt}
時刻 $t$（s） & 0 & 0.5 & 1.0 & 1.5 & 2.0 \\
位置 $x$（cm） & 0 & 5 & 20 & 45 & 80 \\
\end{bktable}
\bktabcaption{ある班の結果。坂の上のスタート地点を0にした。}
\end{bkcenter}
\end{yattemiyou}

\noindent ③の計算をしてみましょう。0.5秒ごとに進んだ距離は、「あとの位置${}-{}$まえの位置」です。それを、かかった時間0.5秒でわります。

\begin{bkcenter}
\begin{bktable}{colspec={Q[l,m,wd=27mm] *{4}{Q[c,m,wd=23.5mm]}}, rowsep=1.5pt, colsep=2pt, cells={font=\small}}
あいだ（s） & 0〜0.5 & 0.5〜1.0 & 1.0〜1.5 & 1.5〜2.0 \\
進んだ距離（cm） & $5-0=5$ & $20-5=15$ & $45-20=25$ & $80-45=35$ \\
距離÷時間（cm/s） & $5\div0.5=10$ & $15\div0.5=30$ & $25\div0.5=50$ & $35\div0.5=70$ \\
\end{bktable}
\bktabcaption{0.5秒ごとの「距離÷時間」}
\end{bkcenter}

なぜ0.5でわるのでしょう。0.5秒で\qty{5}{cm}進んだボールが、同じ速さで1秒（0.5秒の2倍）走れば、2倍の\qty{10}{cm}進みます。だから「1秒あたり\qty{10}{cm}」で、\qty{10}{cm/s}です。0.5でわることは、2倍することと同じです（$5\div0.5=10$）。

「距離÷時間」は、10、30、50、70と増えています。ボールは、だんだん速くなっています。では、ちょうど1.0秒の「瞬間」の速さは、いくつなのでしょう。この章では、それを考えます。

\begin{omoidasou}
\bkrecall{2-3}{速さ${}={}$進んだ道のり÷かかった時間。変わった量を \mbox{$\Delta$（デルタ）}で書く。$\Delta x$ は「$x$ の変化量」。}
\bkrecall{3-4}{横軸を時間 $t$、たて軸を位置 $x$ にしたグラフを、$x$-$t$ グラフという。}
\bkrecall{3-5}{グラフの傾き $=\dfrac{\text{たての変化量}}{\text{横の変化量}}$。直線が右上がりなら正、右下がりなら負。}
\bkrecall{2-10}{\qty{1000}{m/km}（1 km あたり 1000 m）、\qty{3600}{s/h}（1時間あたり3600秒）。\\換算は、単位つきの数をかけたりわったりして、同じ単位を約分して消す。}
\end{omoidasou}

%% ---------------- 数学の道具 ----------------
\bksection{doguu}{数学の道具}

\bksubsection{1}{平均の速さは $\dfrac{\Delta x}{\Delta t}$}

\label{bk:5-1:heikin}\bkhatenaref の車で考えます。出発して0.5時間（30分）のとき、車は出発点から\qty{30}{km}の所にいました。1時間のときは\qty{45}{km}の所にいました。このあいだの速さを求めましょう。

\begin{bkfigure}
\begin{tikzpicture}[x=1mm, y=1mm]
  % 道路を数直線にする（1 km＝1.1 mm）。0.5時間の車と1時間の車。Δx は 30→45 km
  \draw[bkGray, line width=0.8pt] (0,0) -- (118,0);
  \foreach \k in {0,10,...,100} {\draw[bkGray, line width=0.6pt] (\k*1.1,0) -- (\k*1.1,-1.8);
    \node[anchor=north, font=\sffamily\scriptsize] at (\k*1.1,-2.4) {\k};}
  \node[anchor=west, font=\sffamily\scriptsize] at (113,-4.2) {（km）};
  \bkcar{(30*1.1-11.3,0.3)}{0.42}{bkBlue!70}
  \bkcar{(45*1.1-11.3,0.3)}{0.42}{bkOrange}
  \node[anchor=south, font=\sffamily\scriptsize, text=bkBlue] at (30*1.1-4,8.5) {0.5時間のとき};
  \node[anchor=south, font=\sffamily\scriptsize, text=bkOrange!80!black] at (45*1.1+3,8.5) {1時間のとき};
  \draw[bkInk, line width=0.6pt, {Bar[width=4pt]}-{Bar[width=4pt]}] (30*1.1,-9.5) -- (45*1.1,-9.5);
  \node[anchor=north, font=\sffamily\footnotesize] at (37.5*1.1,-10.8) {$\Delta x=45-30=\qty{15}{km}$};
  \node[anchor=north west, font=\sffamily\footnotesize, align=left] at (66,-8.5) {かかった時間\\$\Delta t=1-0.5=\qty{0.5}{h}$};
\end{tikzpicture}
\bkfigcaption{道路を数直線にして、0.5時間と1時間の車の位置をかいた。}
\end{bkfigure}

進んだ距離もかかった時間も、「あと${}-{}$まえ」の変化量です（2-3）。位置 $x$ の変化量を $\Delta x$、時刻 $t$ の変化量を $\Delta t$ と書きます。\bkidx{へんかりょう}{変化量}
\[ \Delta x=45-30=15\quad\text{（km）},\qquad \Delta t=1-0.5=0.5\quad\text{（h）} \]
速さは、距離を時間でわって
\[ \frac{\Delta x}{\Delta t}=\frac{15}{0.5}=30\quad\text{（km/h）} \]
\textbf{時速\qty{30}{km}}です。0.5でわるのは、やってみようと同じで、「0.5時間で\qty{15}{km}なら、1時間ではその2倍」だからです。

この時速\qty{30}{km}は、「このあいだ、ずっと時速\qty{30}{km}で走ったとすると、ちょうど同じ\qty{15}{km}進む」という速さです。じっさいには、このあいだに車は止まったり（休けい）、時速\qty{50}{km}で走ったりしています。それを\textbf{ならした}値なので、\textbf{平均の速さ}といいます。\bkidx{へいきんのはやさ}{平均の速さ}

いつでも使えるように、文字で書きます。時刻 $t_1$ に位置 $x_1$ にいて、時刻 $t_2$ に位置 $x_2$ にいたとします。$t_1$ の右下の小さな1は「1つめ」という印で、「ティー・いち」と読みます。かけ算ではありません。
\[ \Delta t=t_2-t_1,\qquad \Delta x=x_2-x_1,\qquad \text{平均の速さ}=\frac{\text{進んだ距離}}{\text{かかった時間}}=\frac{\Delta x}{\Delta t} \]

同じ車でも、0時間から2時間までをとると、$\Delta x=100-0=100$、$\Delta t=2-0=2$ なので、平均の速さは $100\div2=\qty{50}{km/h}$ です。0.5時間から1時間までの\qty{30}{km/h}とは、ちがう値です。平均の速さは、\textbf{どのあいだをとるか}で変わります。

\bksubsection{2}{平均の速さは、2点を結ぶ直線の傾き}

\label{bk:5-1:katamuki}やってみようのボールの $x$-$t$ グラフに、時刻1.0秒の点Pと、時刻2.0秒の点Qをとり、PとQを結ぶ直線を引きます。

\begin{bkfigure}
\begin{tikzpicture}
\begin{axis}[bkplot, width=78mm, height=48mm, xmin=0, xmax=2.3, ymin=0, ymax=95,
  xtick={0,0.5,1,1.5,2}, ytick={0,20,40,60,80}, xlabel={$t$（s）}, ylabel={$x$（cm）}]
  \addplot[bkcurve, domain=0:2.1] {20*x^2};
  \addplot[bkcurve2, domain=0.7:2.2, line width=1.2pt] {20+60*(x-1)};
  \addplot[only marks, mark=*, mark size=2pt, bkInk] coordinates {(1,20) (2,80)};
  \draw[bkGray, line width=0.7pt, dash pattern=on 2pt off 1.5pt] (axis cs:1,20) -- (axis cs:2,20) -- (axis cs:2,80);
  \node[anchor=north, font=\sffamily\footnotesize] at (axis cs:1.5,20) {$\Delta t=\qty{1.0}{s}$};
  \node[anchor=west, font=\sffamily\footnotesize] at (axis cs:2,50) {$\Delta x=\qty{60}{cm}$};
  \node[anchor=south east, font=\sffamily\small] at (axis cs:1,20) {P};
  \node[anchor=south east, font=\sffamily\small] at (axis cs:2,80) {Q};
\end{axis}
\end{tikzpicture}
\bkfigcaption{やってみようのボール。PとQを結ぶ直線（オレンジ）の傾きは $\dfrac{60}{1.0}=60$。}
\end{bkfigure}

直線の傾きは、「たての変化量÷横の変化量」でした（3-5）。$x$-$t$ グラフでは、たては位置 $x$、横は時刻 $t$ なので、
\[ \text{傾き}=\frac{\text{たての変化量}}{\text{横の変化量}}=\frac{x_2-x_1}{t_2-t_1}=\frac{\Delta x}{\Delta t} \]
これは、平均の速さとまったく同じ式です。PからQまでは、たてに $80-20=\qty{60}{cm}$、横に $2.0-1.0=\qty{1.0}{s}$ なので、傾きは $60\div1.0=60$。1.0秒から2.0秒までの平均の速さは\qty{60}{cm/s}です。

もう1つの見方もあります。ずっと同じ速さで走る物の $x$-$t$ グラフは、まっすぐな直線になります（1秒ごとに同じ距離ずつ進むから）。直線PQは、「Pの時刻にPの位置を出て、ずっと同じ速さで走り、Qの時刻にQの位置に着く物」のグラフです。その速さが平均の速さです。\bkhatenaref のグラフの点線（ずっと時速\qty{50}{km}）も、出発の点と着いた点を結ぶ直線になっていました。

このように、$x$-$t$ グラフの上の\textbf{2点を結ぶ直線の傾き}が、そのあいだの平均の速さになります。

\bksubsection{3}{時間を短くすると、瞬間の速さに近づく}

\label{bk:5-1:shunkan}ボールの1.0秒の瞬間の速さを知りたいとします。1.0秒から2.0秒までの平均（\qty{60}{cm/s}）では、あいだが長すぎます。そのあいだにも、ボールはどんどん速くなっているからです。

それなら、あいだを0にすればよいのでしょうか。1.0秒から1.0秒までにすると、進んだ距離は0、かかった時間も0なので、$0\div0$ になって計算できません（1-14）。そこで、あいだを\textbf{0にはしないで、どんどん短くしていく}ことにします。

短いあいだの位置を知るには、表にない時刻の位置が必要です。表をよく見ると、時刻が0.5秒から1.0秒へ2倍になると、位置は5から20へ4倍（$2\times2$）、1.5秒（0.5秒の3倍）では45で9倍（$3\times3$）です。位置は時刻の2乗に比例しています（3-9）。比例定数は、1.0秒のときの $20\div1.0^2=20$ なので、
\[ x=20t^2\quad\text{（cm）} \]
という式で表せます。表の数と合うか、確かめましょう。

\begin{bkcenter}
\begin{bktable}{colspec={Q[l,m,wd=34mm] *{5}{Q[c,m,wd=14mm]}}, rowsep=1.5pt}
時刻 $t$（s） & 0 & 0.5 & 1.0 & 1.5 & 2.0 \\
$t^2$ & 0 & 0.25 & 1 & 2.25 & 4 \\
位置 $x=20t^2$（cm） & 0 & 5 & 20 & 45 & 80 \\
\end{bktable}
\bktabcaption{式で計算した位置は、やってみようの表とぴったり合う。}
\end{bkcenter}

式があれば、表にない時刻の位置も計算できます。たとえば $t=1.1$ なら、$1.1^2=1.21$ なので $20\times1.21=24.2$ です。
\begin{bkcenter}
\begin{bktable}{colspec={Q[l,m,wd=34mm] *{3}{Q[c,m,wd=22mm]}}, rowsep=1.5pt}
時刻 $t$（s） & 1.1 & 1.01 & 1.001 \\
$t^2$ & 1.21 & 1.0201 & 1.002001 \\
位置 $x=20t^2$（cm） & 24.2 & 20.402 & 20.04002 \\
\end{bktable}
\bktabcaption{表にない時刻の位置を、式で計算した。}
\end{bkcenter}

ここからは、「時刻 $t=1.1$ のときの位置」を $x(1.1)$ と書きます。「エックス・かっこ1.1」と読みます（3-3の $f(3)$ と同じ書き方）。$x(1.1)=24.2$、$x(1.0)=20$ です。かけ算の $x\times1.1$ ではないので、気をつけよう。

これで、1.0秒からのあいだを短くしていけます。たとえば1.0秒から1.1秒までなら、進んだ距離は $\Delta x=x(1.1)-x(1.0)=24.2-20=4.2$、平均の速さは $4.2\div0.1=42$ です。ほかの行も、同じ「あと${}-{}$まえ」と「÷$\Delta t$」で計算します。

\begin{bkcenter}
\begin{bktable}{colspec={Q[c,m,wd=22mm] Q[c,m,wd=14mm] Q[c,m,wd=24mm] Q[c,m,wd=20mm] Q[c,m,wd=26mm]}, rowsep=1.5pt, colsep=3pt}
あいだ & $\Delta t$（s） & あとの位置（cm） & $\Delta x$（cm） & 平均の速さ（cm/s） \\
1.0〜2.0 s & 1 & 80 & 60 & 60 \\
1.0〜1.5 s & 0.5 & 45 & 25 & 50 \\
1.0〜1.1 s & 0.1 & 24.2 & 4.2 & 42 \\
1.0〜1.01 s & 0.01 & 20.402 & 0.402 & 40.2 \\
1.0〜1.001 s & 0.001 & 20.04002 & 0.04002 & 40.02 \\
\end{bktable}
\bktabcaption{1.0秒からのあいだを短くしていくと、平均の速さは40に近づく。}
\end{bkcenter}

平均の速さは、60、50、42、40.2、40.02と、40に近づいていきます。40より大きい分（ずれ）を見ると、20、10、2、0.2、0.02です。$\Delta t$ を $\dfrac{1}{10}$ にすると、ずれも $\dfrac{1}{10}$ になっています。だから、$\Delta t$ を短くすれば、ずれはいくらでも0に近くできます。ずれが消えた先の\qty{40}{cm/s}を、\textbf{1.0秒の瞬間の速さ}と決めます。\bkidx{しゅんかんのはやさ}{瞬間の速さ}

\begin{bkfigure}
\begin{tikzpicture}
\begin{axis}[bkplot, width=82mm, height=50mm, xmin=0, xmax=2.3, ymin=0, ymax=95,
  xtick={0,0.5,1,1.5,2}, ytick={0,20,40,60,80}, xlabel={$t$（s）}, ylabel={$x$（cm）}]
  \addplot[bkcurve, domain=0:2.05] {20*x^2};
  \addplot[bkthin, domain=0.7:2.1] {20+60*(x-1)};
  \addplot[bkthin, domain=0.65:2.15] {20+50*(x-1)};
  \addplot[bkcurve2, domain=0.55:2.2, line width=1.2pt] {20+40*(x-1)};
  \addplot[only marks, mark=*, mark size=1.8pt, bkInk] coordinates {(1,20) (2,80) (1.5,45)};
  \node[anchor=south east, font=\sffamily\small] at (axis cs:1,20) {P};
  % 線の右はしに、少しあけてラベル（3本の高さをずらして重ならないようにする）
  \node[anchor=west, font=\sffamily\scriptsize, text=bkGray] at (axis cs:2.25,88.5) {傾き60};
  \node[anchor=west, font=\sffamily\scriptsize, text=bkGray] at (axis cs:2.25,77.5) {傾き50};
  \node[anchor=west, font=\sffamily\scriptsize\bfseries, text=bkOrange!80!black] at (axis cs:2.25,66.5) {傾き40};
\end{axis}
\end{tikzpicture}
\bkfigcaption{2つめの点を P に近づけると、直線の傾きは 60 → 50 → … → 40 と変わっていく。}
\end{bkfigure}

グラフで見ると、2つめの点を点Pにどんどん近づけることにあたります。直線は、だんだん「点Pでグラフにふれるだけの直線」（オレンジ）に近づき、その傾きは40になります。この直線のことは、5-3でくわしく調べます。

\bksubsection{4}{単位つきの式で、時速と秒速をなおす}

\label{bk:5-1:kansan}速さの単位には、時速（km/h）と秒速（m/s）があります。くらべたり計算したりする前に、単位をそろえます（2-10）。そのとき使うのが、次の「単位つきの数」です。

\begin{bkcenter}
\begin{bktable}{colspec={Q[c,m,wd=26mm] Q[l,m,wd=40mm] Q[l,m,wd=50mm]}, rowsep=1.5pt}
単位つきの数 & 読み方 & なぜその数か \\
\qty{1000}{m/km} & 1 km あたり 1000 m & \qty{1}{km}＝\qty{1000}{m}だから \\
\qty{60}{s/min} & 1分あたり60秒 & 1分は60秒だから \\
\qty{60}{min/h} & 1時間あたり60分 & 1時間は60分だから \\
\qty{3600}{s/h} & 1時間あたり3600秒 & $60\times60=3600$（秒）だから \\
\end{bktable}
\bktabcaption{換算に使う単位つきの数（2-10）}
\end{bkcenter}

\noindent\textbf{時速\qty{72}{km}を、秒速になおす}\par
時速\qty{72}{km}は「1時間あたり\qty{72}{km}」です。これを「1秒あたり何m」にします。2つの段に分けます。

\begin{bkfigure}
\begin{tikzpicture}[x=1mm, y=1mm]
  % 1時間（3600秒）の帯と、72 km（72000 m）の帯。はじめの3こと最後の1こだけかき、あいだは「……」でつなぐ
  \foreach \y/\c/\lab in {13/bkBlue/1秒, 0/bkGreen/\qty{20}{m}} {
    \foreach \x in {22,36,50,108} {\fill[\c!10] (\x,\y) rectangle ++(14,7);
      \draw[\c!70!black, line width=0.7pt] (\x,\y) rectangle ++(14,7);
      \node[font=\sffamily\scriptsize] at (\x+7,\y+3.5) {\lab};}
    \node[font=\sffamily\small, text=bkGray] at (86,\y+3.5) {…………};}
  \fill[bkOrange!35] (22,13) rectangle ++(14,7);\draw[bkOrange!80!black, line width=0.9pt] (22,13) rectangle ++(14,7);
  \node[font=\sffamily\scriptsize] at (29,16.5) {1秒};
  \fill[bkOrange!35] (22,0) rectangle ++(14,7);\draw[bkOrange!80!black, line width=0.9pt] (22,0) rectangle ++(14,7);
  \node[font=\sffamily\scriptsize] at (29,3.5) {\qty{20}{m}};
  \node[anchor=east, font=\sffamily\footnotesize] at (20,16.5) {時間};
  \node[anchor=east, font=\sffamily\footnotesize] at (20,3.5) {距離};
  \node[anchor=south, font=\sffamily\footnotesize] at (72,21) {1時間＝3600秒（1秒が3600こ）};
  \node[anchor=north, font=\sffamily\footnotesize] at (72,-1.5) {\qty{72}{km}＝\qty{72000}{m}（3600こに等しく分ける → 1こ $72000\div3600=\qty{20}{m}$）};
\end{tikzpicture}
\bkfigcaption{1時間を3600この1秒に分けると、\qty{72000}{m}も3600こに分かれる。1秒あたり\qty{20}{m}。}
\end{bkfigure}

\textbf{①}\enspace km を m にする。「1 km あたり 1000 m」をかけます。分子の km と分母の km が約分されて消えます。
\[ 72\,\frac{\cancel{\mathrm{km}}}{\mathrm{h}}\times1000\,\frac{\mathrm{m}}{\cancel{\mathrm{km}}}=72000\,\frac{\mathrm{m}}{\mathrm{h}} \]
\textbf{②}\enspace「1時間あたり」を「1秒あたり」にする。1時間は3600秒なので、\qty{3600}{s/h}でわります。わり算は、ひっくり返した数（逆数）をかけるのと同じなので（2-4）、$\div3600\,\mathrm{s/h}$ は $\times\dfrac{1}{3600}\,\mathrm{h/s}$ です。h が約分されて消えます。
\[ 72000\,\frac{\mathrm{m}}{\cancel{\mathrm{h}}}\times\frac{1}{3600}\,\frac{\cancel{\mathrm{h}}}{\mathrm{s}}=20\,\frac{\mathrm{m}}{\mathrm{s}} \]
まとめて1行で書くと、次のようになります。
\[ 72\,\mathrm{km/h}\times1000\,\mathrm{m/km}\div3600\,\mathrm{s/h}=20\,\mathrm{m/s} \]

\noindent\textbf{なぜ「÷3.6」でよいのか}\par
上の式の数だけを見ると、$72\times1000\div3600$ です。
\[ 1000\div3600=\frac{1000}{3600}=\frac{1}{3.6} \]
なので、$72\times1000\div3600=72\div3.6$ です。だから、\textbf{時速を秒速にするには÷3.6}、反対に\textbf{秒速を時速にするには×3.6}になります（2-10）。÷3.6 は近道ですが、どちらにかけるか迷ったら、単位つきの式にもどれば必ずわかります。

\noindent\textbf{秒速を時速にする（反対向き）}\par
秒速\qty{15}{m}を時速になおすときは、①s を h にするために\qty{3600}{s/h}をかけ（s が約分）、②m を km にするために\qty{1000}{m/km}でわります（m が約分）。
\[ 15\,\mathrm{m/s}\times3600\,\mathrm{s/h}\div1000\,\mathrm{m/km}=54\,\mathrm{km/h} \]
$15\times3.6=54$ と同じ答えです。

\bksubsection{5}{スピードメーターは「とても短い時間の平均」}

車のスピードメーターは、タイヤ（またはタイヤにつながった軸）が回る回数を数えています。タイヤが1回転すると、車はタイヤのまわりの長さだけ進みます。

たとえば、タイヤのまわりの長さが\qty{2}{m}（1回転あたり\qty{2}{m}）で、0.5秒のあいだに7回転したとします。進んだ距離は $\qty{2}{m}\times7=\qty{14}{m}$ なので、
\[ \text{平均の速さ}=\frac{14\,\mathrm{m}}{0.5\,\mathrm{s}}=28\,\mathrm{m/s} \]
時速になおすと、
\[ 28\,\mathrm{m/s}\times3600\,\mathrm{s/h}\div1000\,\mathrm{m/km}=100.8\,\mathrm{km/h} \]

\begin{bkfigure}
\begin{tikzpicture}[x=1mm, y=1mm]
  % 第7回：右へころがるタイヤは時計回り。左は、1回転で「まわりの長さ」だけ進むようすを、
  %   実際の長さの比でかく（半径7 mm のタイヤ → 1回転で 2π×7＝43.98 mm 進む）。右は7回転ぶんの帯
  \draw[bkGray, line width=0.8pt] (-10,-7) -- (52,-7);
  \fill[bkInk!20] (0,0) circle (7);
  \draw[bkInk, line width=1pt] (0,0) circle (7);
  \fill[white] (0,0) circle (3.6);
  \draw[bkInk, line width=0.8pt] (0,0) circle (3.6);
  \fill[bkOrange] (0,-5.6) circle (1.1);
  \draw[bk arrow, bkOrange, line width=1pt] (130:10) arc[start angle=130, end angle=50, radius=10];
  \node[font=\sffamily\scriptsize, text=bkOrange!80!black, anchor=south] at (0,10.6) {回る向き};
  \draw[bkInk!45, dash pattern=on 2pt off 1.5pt, line width=0.8pt] (43.98,0) circle (7);
  \fill[bkOrange!55] (43.98,-5.6) circle (1.1);
  \draw[bk arrow, bkBlue] (10,0) -- (33,0);
  \node[font=\sffamily\scriptsize, text=bkBlue, anchor=south] at (21.5,0.8) {進む向き};
  \draw[bkGray, line width=0.6pt, {Bar[width=4pt]}-{Bar[width=4pt]}] (0,-10) -- (43.98,-10);
  \node[font=\sffamily\footnotesize, anchor=north] at (22,-11) {1回転で\qty{2}{m}（タイヤのまわりの長さ）};
  \begin{scope}[shift={(62,0)}]
    \foreach \k in {0,...,6} {\draw[bkGray, fill=bkBlue!8, line width=0.6pt] (\k*9,-7) rectangle ++(9,4);
      \node[font=\sffamily\scriptsize] at (\k*9+4.5,-5) {2 m};}
    \draw[bk arrow, bkBlue] (0,1) -- (63,1);
    \node[font=\sffamily\footnotesize, anchor=north] at (31.5,-11) {0.5秒で7回転 → $2\times7=\qty{14}{m}$};
  \end{scope}
\end{tikzpicture}
\end{bkfigure}

0.5秒はとても短いので、そのあいだに速さはほとんど変わりません。だから、この平均の速さを、その瞬間の速さとして表示しているのです。メーターは、この計算を何度もくり返して、表示を新しくしています。スピードメーターが表示する「瞬間の速さ」は、じつは\textbf{とても短い時間の平均の速さ}です。序章（5-0）の「走っていない1時間の距離を、どうやって知ったのか」の答えは、「とても短い時間に進んだ距離から、1時間あたりの量を計算した」です。

\begin{taisetsu}
\textbf{平均の速さ} $=\dfrac{\Delta x}{\Delta t}$\quad（グラフの2点を結ぶ直線の傾き）\\
\textbf{瞬間の速さ}：$\Delta t$ を0にしないで、どんどん短くしたときに、平均の速さが近づく値\\
\textbf{換算}：$\mathrm{km/h}\times1000\,\mathrm{m/km}\div3600\,\mathrm{s/h}=\mathrm{m/s}$（÷3.6 と同じ）
\end{taisetsu}

\begin{kotoba}
\small
\bkword{平均の速さ}{あるあいだに進んだ距離を、かかった時間でわったもの。そのあいだ、ずっと同じ速さで走ったとしたときの速さ。}
\bkword{瞬間の速さ}{あいだをどんどん短くしたときに、平均の速さが近づいていく値。スピードメーターが表示する速さ。}
\bkword{$\Delta t$、$\Delta x$}{時間と位置の変化量。「デルタ・ティー」「デルタ・エックス」と読む。}
\bkword{$x(1.1)$}{時刻 $t=1.1$ のときの位置 $x$。かけ算ではない。}
\end{kotoba}

\begin{bkfigure}
\begin{tikzpicture}[x=1mm, y=1mm]
  % この章の流れ。箱どうしを矢印でつなぐ
  \node[bk node, text width=30mm, align=center, font=\sffamily\small] (a) at (0,0) {あるあいだの\\平均の速さ $\dfrac{\Delta x}{\Delta t}$};
  \node[bk node, text width=30mm, align=center, font=\sffamily\small] (b) at (42,0) {2点を結ぶ\\直線の傾き};
  \node[bk node, text width=30mm, align=center, font=\sffamily\small] (c) at (84,0) {あいだを\\どんどん短く};
  \node[bk node, text width=30mm, align=center, font=\sffamily\small, fill=bkOrange!12] (d) at (84,-24) {瞬間の速さ\\（近づく先）};
  \node[bk node, text width=34mm, align=center, font=\sffamily\small] (e) at (36,-24) {スピードメーター\\（とても短い時間）};
  \draw[bk link] (a) -- (b);
  \draw[bk link] (b) -- (c);
  \draw[bk link] (c) -- (d);
  \draw[bk link] (d) -- (e);
\end{tikzpicture}
\bkfigcaption{この章の流れ。平均の速さから、瞬間の速さへ。}
\end{bkfigure}

%% ---------------- 例題 ----------------
\bksection{reidai}{例題でたしかめよう}

\begin{reidai}[表から平均の速さを求める]
\bkhatenaref の車の位置を、0.5時間ごとに記録しました。1時間から1.5時間までの平均の速さを求めましょう。
\begin{bkcenter}
\begin{bktable}{colspec={Q[l,m,wd=24mm] *{5}{Q[c,m,wd=11mm]}}, rowsep=1pt}
$t$（h） & 0 & 0.5 & 1.0 & 1.5 & 2.0 \\
$x$（km） & 0 & 30 & 45 & 70 & 100 \\
\end{bktable}
\end{bkcenter}
\kangaekata 時間と位置の変化量を「あと${}-{}$まえ」で求めて、わり算します。\par
\shiki $\Delta t=1.5-1.0=0.5$\quad（h）\par
\hspace*{2.6\zw}$\Delta x=70-45=25$\quad（km）\par
\hspace*{2.6\zw}$\dfrac{\Delta x}{\Delta t}=\dfrac{25}{0.5}=50$\par
\kotae \qty{50}{km/h}
\end{reidai}

\begin{jibun}
例題1の表で、次のあいだの平均の速さを求めましょう。\quad（1）\mbox{1.5〜\qty{2.0}{h}}\quad（2）\mbox{0〜\qty{1.0}{h}}\quad（3）\mbox{0.5〜\qty{1.5}{h}}
\bkjans{（1）\nobreak \qty{60}{km/h}\quad（2）\nobreak \qty{45}{km/h}\quad（3）\nobreak \qty{40}{km/h}}
\end{jibun}

\begin{reidai}[グラフから平均の速さを求める]
\noindent\begin{minipage}[c]{0.55\linewidth}
右の $x$-$t$ グラフの点A（\qty{2}{s}、\qty{6}{m}）から点B（\qty{5}{s}、\qty{21}{m}）までの、平均の速さを求めましょう。\par\vspace{2pt}
\kangaekata 平均の速さは、AとBを結ぶ直線の傾きです。たての変化量（位置）を、横の変化量（時間）でわります。\par
\shiki $\Delta t=5-2=3$\quad（s）、\quad$\Delta x=21-6=15$\quad（m）\par
\hspace*{2.6\zw}$\dfrac{\Delta x}{\Delta t}=\dfrac{15}{3}=5$\par
\kotae \qty{5}{m/s}
\end{minipage}\hfill
\begin{minipage}[c]{0.42\linewidth}
\begin{tikzpicture}
\begin{axis}[bkplot, width=49mm, height=40mm, xmin=0, xmax=6.4, ymin=0, ymax=27, xtick={0,1,...,6}, ytick={0,6,12,18,24},
  xlabel={$t$（s）}, ylabel={$x$（m）}]
  \addplot[bkcurve, domain=0:6, line width=1.3pt] {x^2/2+x+2};
  \addplot[only marks, mark=*, mark size=1.8pt, bkInk] coordinates {(2,6) (5,21)};
  \addplot[bkcurve2, line width=1pt, domain=1.4:5.6] {6+5*(x-2)};
  \node[anchor=north west, font=\sffamily\scriptsize] at (axis cs:2,6) {A};
  \node[anchor=south east, font=\sffamily\scriptsize] at (axis cs:5,21) {B};
\end{axis}
\end{tikzpicture}
\end{minipage}
\end{reidai}

\begin{jibun}
\noindent\begin{minipage}[c]{0.55\linewidth}
右の $x$-$t$ グラフの上の3点、C（\qty{1}{s}、\qty{4}{m}）、D（\qty{4}{s}、\qty{16}{m}）、E（\qty{6}{s}、\qty{26}{m}）について、2点を結ぶ直線の傾きを求め、次の平均の速さを答えましょう。\par
（1）CからD\quad（2）DからE
\end{minipage}\hfill
\begin{minipage}[c]{0.42\linewidth}
\begin{tikzpicture}
\begin{axis}[bkplot, width=49mm, height=40mm, xmin=0, xmax=6.8, ymin=0, ymax=30, xtick={0,1,...,6}, ytick={0,8,16,24},
  xlabel={$t$（s）}, ylabel={$x$（m）}]
  \addplot[bkcurve, domain=0:6.5, line width=1.3pt] {0.2*x^2+3*x+0.8};
  \addplot[only marks, mark=*, mark size=1.8pt, bkInk] coordinates {(1,4) (4,16) (6,26)};
  \node[anchor=north west, font=\sffamily\scriptsize] at (axis cs:1,4) {C};
  \node[anchor=north west, font=\sffamily\scriptsize] at (axis cs:4,16) {D};
  \node[anchor=north west, font=\sffamily\scriptsize] at (axis cs:6,26) {E};
\end{axis}
\end{tikzpicture}
\end{minipage}
\bkjans{（1）\nobreak \qty{4}{m/s}（$\dfrac{16-4}{4-1}=4$）\quad（2）\nobreak \qty{5}{m/s}（$\dfrac{26-16}{6-4}=5$）}
\end{jibun}

\begin{reidai}[あいだを短くして、瞬間の速さを見積もる]
$x=20t^2$（cm）で動くボールの、2.0秒の瞬間の速さを見積もりましょう。\par
\kangaekata 2.0秒から、$\Delta t=0.1$ と $0.01$ のあいだの平均の速さを計算し、近づいていく値を読みます。はじめに、2.0秒の位置を求めます。\par
\shiki $x(2.0)=20\times2.0^2=20\times4=80$\par
\hspace*{2.6\zw}$\Delta t=0.1$ のとき：$x(2.1)=20\times2.1^2=20\times4.41=88.2$\par
\hspace*{2.6\zw}\phantom{$\Delta t=0.1$ のとき：}$\dfrac{88.2-80}{0.1}=\dfrac{8.2}{0.1}=82$\par
\hspace*{2.6\zw}$\Delta t=0.01$ のとき：$x(2.01)=20\times2.01^2=20\times4.0401=80.802$\par
\hspace*{2.6\zw}\phantom{$\Delta t=0.01$ のとき：}$\dfrac{80.802-80}{0.01}=\dfrac{0.802}{0.01}=80.2$\par
\begin{bkcenter}
\begin{bktable}{colspec={Q[c,m,wd=26mm] *{3}{Q[c,m,wd=18mm]}}, rowsep=1pt}
$\Delta t$ & 0.1 & 0.01 & → \\
平均の速さ & 82 & 80.2 & 80 \\
80からのずれ & 2 & 0.2 & 0 \\
\end{bktable}
\end{bkcenter}
\kotae 82、80.2と80に近づくので、およそ\qty{80}{cm/s}
\end{reidai}

\begin{jibun}
例題3のボールについて、$\Delta t=0.1$ と $0.01$ のあいだの平均の速さを計算して、次の時刻の瞬間の速さを見積もりましょう。\par
（1）0.5秒\quad（2）1.5秒
\bkjans{（1）\nobreak $x(0.5)=5$。$\dfrac{7.2-5}{0.1}=22$、$\dfrac{5.202-5}{0.01}=20.2$。およそ\qty{20}{cm/s}\quad（2）\nobreak $x(1.5)=45$。$\dfrac{51.2-45}{0.1}=62$、$\dfrac{45.602-45}{0.01}=60.2$。およそ\qty{60}{cm/s}}
\end{jibun}

\begin{reidai}[時速と秒速をなおす]
（1）時速\qty{72}{km}は、秒速\mbox{何 m}？\quad（2）秒速\qty{15}{m}は、時速\mbox{何 km}？\par
\kangaekata 単位つきの数 \qty{1000}{m/km} と \qty{3600}{s/h} をかけたりわったりして、消したい単位を約分します。\par
\shiki（1）km を消すために \qty{1000}{m/km} をかけ、h を消すために \qty{3600}{s/h} でわる。\par
\hspace*{4.3\zw}$72\,\mathrm{km/h}\times1000\,\mathrm{m/km}\div3600\,\mathrm{s/h}$\par
\hspace*{4.3\zw}$=72\times1000\div3600\ \mathrm{m/s}=20\,\mathrm{m/s}$\par
\hspace*{2.6\zw}（2）s を消すために \qty{3600}{s/h} をかけ、m を消すために \qty{1000}{m/km} でわる。\par
\hspace*{4.3\zw}$15\,\mathrm{m/s}\times3600\,\mathrm{s/h}\div1000\,\mathrm{m/km}$\par
\hspace*{4.3\zw}$=15\times3600\div1000\ \mathrm{km/h}=54\,\mathrm{km/h}$\par
\kotae （1）\qty{20}{m/s}\quad（2）\qty{54}{km/h}
\end{reidai}

\begin{jibun}
（1）時速\qty{90}{km}は秒速\mbox{何 m}でしょう。\quad（2）秒速\qty{12}{m}は時速\mbox{何 km}でしょう。単位つきの式で計算し、÷3.6・×3.6 でも確かめよう。
\bkjans{（1）\nobreak \qty{25}{m/s}\quad（2）\nobreak \qty{43.2}{km/h}}
\end{jibun}

\begin{reidai}[タイヤの回転から速さを求める]
まわりの長さ\qty{1.5}{m}のタイヤが、0.25秒に4回転しました。秒速\mbox{何 m}、時速\mbox{何 km}でしょう。\par
\kangaekata 1回転で、タイヤのまわりの長さ\qty{1.5}{m}だけ進みます。進んだ距離を求めてから、時間でわります。\par
\shiki 進んだ距離：$\qty{1.5}{m}\times4=\qty{6}{m}$\par
\hspace*{2.6\zw}秒速：$\dfrac{6\,\mathrm{m}}{0.25\,\mathrm{s}}=24\,\mathrm{m/s}$\par
\hspace*{2.6\zw}時速：$24\,\mathrm{m/s}\times3600\,\mathrm{s/h}\div1000\,\mathrm{m/km}=86.4\,\mathrm{km/h}$\par
\begin{bkcenter}
\begin{tikzpicture}[x=1mm, y=1mm]
  % 1回転ごとに 1.5 m。4回転で 6 m（長さ 1 m＝16 mm）
  \foreach \k in {0,...,3} {\fill[bkBlue!12] (\k*24,0) rectangle ++(24,5); \draw[bkBlue!70!black, line width=0.6pt] (\k*24,0) rectangle ++(24,5);
    \node[font=\sffamily\scriptsize] at (\k*24+12,2.5) {\the\numexpr\k+1\relax 回転め \qty{1.5}{m}};}
  \draw[bkGray, line width=0.6pt, {Bar[width=4pt]}-{Bar[width=4pt]}] (0,-3) -- (96,-3);
  \node[anchor=north, font=\sffamily\footnotesize] at (48,-4) {0.25秒で $\qty{1.5}{m}\times4=\qty{6}{m}$};
\end{tikzpicture}
\end{bkcenter}
\kotae \qty{24}{m/s}、\qty{86.4}{km/h}
\end{reidai}

\begin{jibun}
まわりの長さ\qty{2.0}{m}のタイヤが、次のように回転しました。それぞれ、秒速\mbox{何 m}、時速\mbox{何 km}でしょう。\par
（1）0.4秒に5回転\quad（2）1秒に10回転
\bkjans{（1）\nobreak $\qty{2.0}{m}\times5=\qty{10}{m}$、$10\div0.4=\qty{25}{m/s}$、\qty{90}{km/h}\quad（2）\nobreak \qty{20}{m/s}、\qty{72}{km/h}}
\end{jibun}

\begin{reidai}[分速を、時速と秒速になおす]
電車が2分で\qty{2.4}{km}進みました。（1）1分あたり\mbox{何 km}（km/min）でしょう。（2）時速\mbox{何 km}でしょう。（3）秒速\mbox{何 m}でしょう。\par
\kangaekata 1時間は60分なので \qty{60}{min/h}、1分は60秒なので \qty{60}{s/min} を使います。\par
\shiki（1）$\dfrac{2.4\,\mathrm{km}}{2\,\mathrm{min}}=1.2\,\mathrm{km/min}$\par
\hspace*{2.6\zw}（2）min を消すために \qty{60}{min/h} をかける。\par
\hspace*{4.3\zw}$1.2\,\mathrm{km/min}\times60\,\mathrm{min/h}=72\,\mathrm{km/h}$\par
\hspace*{2.6\zw}（3）km を m に、min を s にする。\par
\hspace*{4.3\zw}$1.2\,\mathrm{km/min}\times1000\,\mathrm{m/km}\div60\,\mathrm{s/min}=20\,\mathrm{m/s}$\par
\begin{bkcenter}
\begin{tikzpicture}[x=1mm, y=1mm]
  % 時間の帯（60分＝100 mm、目もり1つが1分）。はじめの2分に色
  \fill[bkOrange!50] (0,6) rectangle ({2*100/60},11);
  \foreach \k in {1,...,59} {\draw[bkGray!60, line width=0.25pt] ({\k*100/60},6) -- ({\k*100/60},11);}
  \draw[bkInk, line width=0.6pt] (0,6) rectangle (100,11);
  \node[anchor=south west, font=\sffamily\scriptsize] at (0,11.8) {はじめの2分（色の部分）で\qty{2.4}{km}、1分あたり\qty{1.2}{km}};
  \draw[bkGray, line width=0.6pt, {Bar[width=4pt]}-{Bar[width=4pt]}] (0,3) -- (100,3);
  \node[anchor=north, font=\sffamily\scriptsize] at (50,1.8) {60分（1時間）では $\qty{1.2}{km/min}\times\qty{60}{min}=\qty{72}{km}$（目もり1つが1分）};
\end{tikzpicture}
\end{bkcenter}
\kotae （1）\qty{1.2}{km/min}\quad（2）\qty{72}{km/h}\quad（3）\qty{20}{m/s}
\end{reidai}

\begin{jibun}
自転車が4分で\qty{3}{km}進みました。（1）\mbox{何 km/min}でしょう。（2）時速\mbox{何 km}でしょう。（3）秒速\mbox{何 m}でしょう。
\bkjans{（1）\nobreak \qty{0.75}{km/min}\quad（2）\nobreak \qty{45}{km/h}\quad（3）\nobreak \qty{12.5}{m/s}}
\end{jibun}

\begin{reidai}[行きと帰りで速さがちがうときの平均]
片道\qty{30}{km}の道を、行きは時速\qty{30}{km}、帰りは時速\qty{15}{km}で走りました。往復の平均の速さを求めましょう。\par
\kangaekata 平均の速さは「全部の距離÷全部の時間」です。まず、行きと帰りにかかった時間を求めます。\par
\shiki 行きの時間：$30\div30=1$\quad（h）\par
\hspace*{2.6\zw}帰りの時間：$30\div15=2$\quad（h）\par
\hspace*{2.6\zw}全部の距離：$30+30=60$\quad（km）\par
\hspace*{2.6\zw}平均の速さ：$\dfrac{60}{1+2}=\dfrac{60}{3}=20$\par
\begin{bkcenter}
\begin{tikzpicture}[x=1mm, y=1mm]
  \fill[bkGray!20] (10,-1.5) rectangle (100,1.5);
  \node[bk node, font=\sffamily\small] (A) at (0,0) {家};
  \node[bk node, font=\sffamily\small] (B) at (110,0) {町};
  \draw[bk arrow, bkBlue] (14,6) -- (96,6);
  \node[anchor=south, font=\sffamily\scriptsize, text=bkBlue] at (55,6.8) {行き：時速\qty{30}{km}で\qty{30}{km} → 1時間};
  \draw[bk arrow, bkOrange] (96,-6) -- (14,-6);
  \node[anchor=north, font=\sffamily\scriptsize, text=bkOrange!80!black] at (55,-6.8) {帰り：時速\qty{15}{km}で\qty{30}{km} → 2時間};
\end{tikzpicture}
\end{bkcenter}
\kotae \qty{20}{km/h}（$\dfrac{30+15}{2}=22.5$ ではない）
\end{reidai}

\begin{jibun}
片道\qty{12}{km}の道を、行きは時速\qty{12}{km}、帰りは時速\qty{6}{km}で走りました。往復の平均の速さを求めましょう。
\bkjans{行き1時間、帰り2時間。$\dfrac{24}{3}=\qty{8}{km/h}$}
\end{jibun}

\clearpage
%% ---------------- まちがえやすいところ・物理工学 ----------------
\bksection{tsumazuki}{まちがえやすいところをチェック}

\begin{tsumazuki}{平均の速さは「速さの平均」ではない}
\qty{120}{km}はなれた町まで、行きは時速\qty{60}{km}、帰りは時速\qty{40}{km}で走りました。往復の平均の速さは？
\par\vspace{2pt}
\bkng{$\dfrac{60+40}{2}=\qty{50}{km/h}$}\bkok{行き2時間、帰り3時間。$\dfrac{240}{2+3}=\qty{48}{km/h}$}
\par\vspace{2pt}
平均の速さは、いつも「全部の距離÷全部の時間」で求めます。おそい帰りのほうが時間が長いので、平均はおそいほうに引っぱられます。
\end{tsumazuki}

\begin{tsumazuki}{単位をそろえずにわり算する}
$\qty{3}{km}$ を5分で走ったとき、$3\div5=0.6$ は「\qty{0.6}{km/min}」です。「\qty{0.6}{m/s}」ではありません。
\bkng{$3\div5=0.6$ なので \qty{0.6}{m/s}}\bkok{$0.6\,\mathrm{km/min}\times1000\,\mathrm{m/km}\div60\,\mathrm{s/min}=10\,\mathrm{m/s}$}
数だけでなく、単位もいっしょに書いて計算すると、まちがいに気づけます。
\end{tsumazuki}

\begin{tsumazuki}{換算で、かけるとわるを逆にする}
\bkng{$72\,\mathrm{km/h}\times3600\,\mathrm{s/h}=259200$}\bkok{$72\,\mathrm{km/h}\div3600\,\mathrm{s/h}\times1000\,\mathrm{m/km}=20\,\mathrm{m/s}$}
左の式では h が約分されず、$\mathrm{km\cdot s/h^2}$ というおかしな単位が残ります。単位が m/s になるかで確かめよう。
\end{tsumazuki}

\begin{tsumazuki}{$x(1.1)$ をかけ算と読む}
\bkng{$x(1.1)=x\times1.1$}\bkok{$x(1.1)$ は「$t=1.1$ のときの位置」。$x=20t^2$ なら $20\times1.1^2=24.2$}
かっこの前の $x$ は「位置を求めるきまり」の名前で、かけられる数ではありません（3-3の $f(3)$ と同じ）。
\end{tsumazuki}

\begin{tsumazuki}{瞬間の速さを $0\div0$ で求めようとする}
\bkng{あいだを0にして、$\dfrac{0}{0}$ を計算する}\bkok{あいだを0.1、0.01、0.001と短くして、近づく先を読む}
0でわる計算はできません（1-14）。あいだは0にしないで、短くしていきます。
\end{tsumazuki}

\begin{butsuri}
\begin{itemize}[leftmargin=1.2\zw, itemsep=2pt, topsep=0pt, label={\color{bkPurple}●}]
\item \textbf{区間の平均速度の取りしまり}\par
外国の高速道路などには、2つの地点を通った時刻から、そのあいだの平均の速さを計算して、速すぎる車を見つけるしくみがあります。
\item \textbf{スマートフォンの速さ}\par
地図アプリは、短い時間ごとの位置（GPS）の変化などから、速さを計算して表示しています。
\item \textbf{100 m 走の記録}\par
9.58秒の世界記録の平均の速さは、$100\div9.58$ で約\qty{10.4}{m/s}。でも、速さがいちばん大きいところでは約\qty{12}{m/s}あったと計算されています。スタートの直後はおそいので、平均はそれより小さくなるのです。
\end{itemize}
\end{butsuri}

\begin{motto}[2]
\textbf{もどるときの速さはマイナス}\par
ボールが坂をのぼって、また下りてくるとき、下りるあいだは位置 $x$ が小さくなるので、$\Delta x$ はマイナスになります（2-2：マイナスは反対の向き）。たとえば2秒で\qty{6}{m}もどったら、$\dfrac{\Delta x}{\Delta t}=\dfrac{-6}{2}=\qty{-3}{m/s}$。このように向きもふくめた速さを\textbf{速度}といいます。$x$-$t$ グラフでは、右下がりの直線の傾きにあたります。速度は5-5でくわしく使います。\bkidx{そくど}{速度}
\end{motto}

\clearpage
%% ---------------- 章末問題 ----------------
\bkendproblems
\bkproblemgroup{基本}{2}

\bkproblem 電車の位置を2分ごとに記録しました。
\begin{bkcenter}
\begin{bktable}{colspec={Q[l,m,wd=24mm] *{4}{Q[c,m,wd=12mm]}}, rowsep=1pt}
$t$（min） & 0 & 2 & 4 & 6 \\
$x$（km） & 0 & 1.2 & 4.0 & 6.0 \\
\end{bktable}
\end{bkcenter}
（1）0〜6分の平均の速さは、\mbox{何 km/min} でしょう。\quad（2）2〜4分の平均の速さは、\mbox{何 km/min} でしょう。\par
（3）（1）の速さは、時速\mbox{何 km}でしょう。単位つきの式で求めましょう。
\bkpans{（1）\nobreak \qty{1.0}{km/min}\quad（2）\nobreak \qty{1.4}{km/min}\quad（3）\nobreak $1.0\,\mathrm{km/min}\times60\,\mathrm{min/h}=\qty{60}{km/h}$}

\bkproblem $x$-$t$ グラフの上に、3つの点A（\qty{0}{s}、\qty{0}{m}）、B（\qty{4}{s}、\qty{12}{m}）、C（\qty{10}{s}、\qty{24}{m}）があります。AからB、BからC、AからCの平均の速さを、それぞれ求めましょう。
\bkpans{A→B \qty{3}{m/s}、B→C \qty{2}{m/s}、A→C \qty{2.4}{m/s}}

\bkproblem 単位つきの式で求めましょう。（1）時速\qty{36}{km}は秒速\mbox{何 m}？\quad（2）秒速\qty{5}{m}は時速\mbox{何 km}？\quad（3）時速\qty{300}{km}の新幹線は、秒速\mbox{何 m}？（小数第1位まで）
\bkpans{（1）\nobreak \qty{10}{m/s}\quad（2）\nobreak \qty{18}{km/h}\quad（3）\nobreak \qty{83.3}{m/s}}

\bkproblem タイヤのまわりの長さが\qty{1.8}{m}の車で、タイヤが0.2秒に3回転しました。このときの速さは、秒速\mbox{何 m}ですか。時速\mbox{何 km}ですか。
\bkpans{\qty{27}{m/s}、\qty{97.2}{km/h}}

\bkproblem $x=20t^2$（cm）で動くボールの、2.5秒の瞬間の速さを見積もりましょう。$\Delta t=0.1$ と $0.01$ のあいだの平均の速さを計算して、近づく値を答えましょう。
\bkpans{$x(2.5)=125$、$x(2.6)=135.2$ より 102、$x(2.51)=126.002$ より 100.2。およそ\qty{100}{cm/s}}

\bkproblemgroup{チャレンジ}{2}

\bkproblem 位置が $x=5t^2$（m）で表される物体があります。
（1）3.0秒から3.1秒までの平均の速さを求めましょう。\quad（2）3.0秒から3.01秒までの平均の速さを求めましょう。\par
（3）3.0秒の瞬間の速さは、およそいくつと考えられますか。
\bkpans{（1）\nobreak \qty{30.5}{m/s}（$\dfrac{48.05-45}{0.1}$）\quad（2）\nobreak \qty{30.05}{m/s}\quad（3）\nobreak およそ\qty{30}{m/s}}

\bkproblem 片道\qty{6}{km}の山道を、上りは時速\qty{4}{km}、下りは時速\qty{6}{km}で歩きました。往復の平均の速さを求めましょう。
\bkpans{上り1.5時間、下り1時間。$\dfrac{12}{2.5}=\qty{4.8}{km/h}$}

\bkproblem 秒速\qty{340}{m}で進む音は、時速\mbox{何 km}でしょう。単位つきの式で求めましょう。また、時速\qty{1000}{km}で飛ぶ飛行機と、どちらが速いでしょう。
\bkpans{$340\,\mathrm{m/s}\times3600\,\mathrm{s/h}\div1000\,\mathrm{m/km}=\qty{1224}{km/h}$。音のほうが速い}

\begin{furikaeri}
\begin{checklist}
\item 平均の速さを $\dfrac{\Delta x}{\Delta t}$ で求められる。\hfill{\footnotesize\sffamily\color{bkTheme}→\pageref{bk:5-1:heikin}ページ}
\item 平均の速さが、2点を結ぶ直線の傾きだと説明できる。\hfill{\footnotesize\sffamily\color{bkTheme}→\pageref{bk:5-1:katamuki}ページ}
\item あいだを短くして、瞬間の速さを見積もれる。\hfill{\footnotesize\sffamily\color{bkTheme}→\pageref{bk:5-1:shunkan}ページ}
\end{checklist}
\end{furikaeri}

\bkfukidashi{mina}{瞬間の速さは「近づいていく値」なんだね。次の章では、近づいていくようすを式で追いかけてみよう。}

%% 5-2 時間をどんどん短くしてみよう（第11回：ページ制限なし。偶数ページ）
\bkchapter{5}{2}{時間をどんどん短くしてみよう}{\bklevel{2}}{時間をどんどん短くしてみよう}
\bkmichishirube{3-8,3-9,5-1}{5-3,5-4}

\begin{goal}
\begin{itemize}[label={\color{bkTheme}●}, leftmargin=1.6\zw, itemsep=0pt, topsep=1pt]
\item \bklevel{2}\enspace $\Delta t$ を 0.1、0.01 と小さくして、平均の速さを表で計算できる。
\item \bklevel{2}\enspace 平均の速さを $h$ の式で表し、$h$ が小さいときの値を読める。
\item \bklevel{2}\enspace $x=at^2$ の形の運動で、瞬間の速さを求められる。
\end{itemize}
\end{goal}

\begin{hatena}
\noindent\begin{minipage}[c]{0.71\linewidth}
手からボールを落とすと、だんだん速くなりながら落ちていきます。落ちはじめてから $t$ 秒のあいだに落ちる距離は、
\[ x=5t^2\quad\text{（m）} \]
と表せます（ほんとうは $4.9t^2$ ですが、この章では計算しやすいように5にします）。1秒たった瞬間、ボールの速さはいくつでしょう。1秒で\qty{5}{m}落ちたから、秒速\qty{5}{m}？
\end{minipage}\hfill
\begin{minipage}[c]{0.26\linewidth}\raggedleft
\begin{tikzpicture}[x=1mm, y=0.9mm]
  \draw[bkGray, line width=0.8pt] (0,0) -- (0,-47);
  \foreach \t/\x in {0/0,1/5,2/20,3/45} {
    \draw[bkGray, line width=0.6pt] (-1.5,-\x) -- (1.5,-\x);
    \node[anchor=east, font=\sffamily\scriptsize] at (-2,-\x) {\x\ m};
    \shade[ball color=bkOrange] (8,-\x) circle (2.2);
    \node[anchor=west, font=\sffamily\scriptsize] at (11,-\x) {\t 秒};}
\end{tikzpicture}
\end{minipage}
\end{hatena}

図の数を、式で確かめてみましょう。2秒なら $5\times2^2=5\times4=20$、3秒なら $5\times3^2=5\times9=45$ です。1秒ごとに落ちる距離は、5、15、25 m と、どんどん大きくなっています。

5-1で調べたように、0秒から1秒までの平均の速さは $\dfrac{5-0}{1-0}=\qty{5}{m/s}$ です。でも、ボールはだんだん速くなっているので、1秒の瞬間の速さは、それより大きいはずです。この章では、あいだを短くして、瞬間の速さを追いかけます。

%% ---------------- やってみよう・思い出そう ----------------
\begin{yattemiyou}[表計算であいだを短くしよう]
表計算のシートに、次のように入力します。A列にあいだ $\Delta t$、B列に $1+\Delta t$ 秒の位置、C列に平均の速さを計算させます。
\begin{bkcenter}
\bksheet{3}{28}{%
1 & $\Delta t$ & $1+\Delta t$ 秒の位置 & 平均の速さ \\
2 & 1 & \texttt{=5*(1+A2)\^{}2} & \texttt{=(B2-5)/A2} \\
3 & 0.1 & （上をコピー） & （上をコピー） \\}
\end{bkcenter}
\begin{enumerate}[label={\bkcircled[bkGreen]{\arabic*}}, leftmargin=2.4\zw, itemsep=1pt, topsep=3pt]
\item A列に 1、0.1、0.01、0.001、0.0001 と入れ、B列とC列の式を下へコピーしよう。
\item C列の数は、何に近づいていくかな？
\end{enumerate}
\end{yattemiyou}

\noindent 式の意味を確かめておきます。\texttt{\^{}2} は2乗（2-18）です。B列の \texttt{=5*(1+A2)\^{}2} は、$x=5t^2$ の $t$ に $1+\Delta t$ を入れた「$1+\Delta t$ 秒の位置」です。C列の \texttt{=(B2-5)/A2} は、「あとの位置${}-{}$1秒の位置（5）」で進んだ距離を出し、それをあいだ $\Delta t$ でわった、平均の速さです。

\begin{bkcenter}
\bksheet{3}{28}{%
1 & $\Delta t$ & $1+\Delta t$ 秒の位置 & 平均の速さ \\
2 & 1 & 20 & 15 \\
3 & 0.1 & 6.05 & 10.5 \\
4 & 0.01 & 5.1005 & 10.05 \\
5 & 0.001 & 5.010005 & 10.005 \\
6 & 0.0001 & 5.00100005 & 10.0005 \\}
\bktabcaption{計算の結果。平均の速さは 10 に近づいていく。}
\end{bkcenter}

たとえば3行めは、$\Delta t=0.1$ なので1.1秒の位置 $5\times1.1^2=5\times1.21=6.05$ を求め、$(6.05-5)\div0.1=1.05\div0.1=10.5$ です。電卓でも確かめてみよう。

\begin{omoidasou}
\bkrecall{5-1}{平均の速さ $=\dfrac{\Delta x}{\Delta t}$。あいだを短くすると、瞬間の速さに近づく。$x(1.1)$ は「$t=1.1$ のときの $x$」。}
\bkrecall{3-9}{$y=ax^2$ のグラフは放物線。$x$ が2倍になると $y$ は4倍になる。}
\bkrecall{3-8}{$(a+b)^2=a^2+2ab+b^2$、$(a-b)^2=a^2-2ab+b^2$\\たとえば $(1+h)^2=1+2h+h^2$}
\bkrecall{2-12}{まとめてわるときは、1つずつわってよい：$(A+B)\div h=A\div h+B\div h$}
\end{omoidasou}

%% ---------------- 数学の道具 ----------------
\bksection{doguu}{数学の道具}

\bksubsection{1}{前からも後ろからも、10に近づく}

\leavevmode\label{bk:5-2:hyou}やってみようでは、1秒より後ろの時刻（$1+\Delta t$ 秒）をとりました。1秒より前の時刻（$1-\Delta t$ 秒）から1秒までの平均の速さも、計算してみましょう。たとえば $\Delta t=0.1$ なら、0.9秒の位置は $5\times0.9^2=5\times0.81=4.05$ です。0.9秒から1秒までに $5-4.05=0.95$（m）進んだので、
\[ \frac{5-4.05}{0.1}=\frac{0.95}{0.1}=9.5 \]

\begin{bkcenter}
\begin{bktable}{colspec={Q[l,m,wd=44mm] *{4}{Q[c,m,wd=16mm]}}, rowsep=1.5pt}
$\Delta t$ & 1 & 0.1 & 0.01 & 0.001 \\
1秒から $1+\Delta t$ 秒まで & 15 & 10.5 & 10.05 & 10.005 \\
$1-\Delta t$ 秒から1秒まで & 5 & 9.5 & 9.95 & 9.995 \\
\end{bktable}
\bktabcaption{1秒の前と後ろの、平均の速さ（m/s）}
\end{bkcenter}

\begin{bkfigure}
\begin{tikzpicture}[x=1mm, y=1mm]
  % 数直線（9.4〜10.6、1＝100 mm）。後ろからの値は上、前からの値は下に書き、高さをずらす
  \draw[bk arrow, bkInk, line width=0.8pt] (-2,0) -- (124,0);
  \foreach \v in {9.5,10,10.5} {\draw[bkInk, line width=0.6pt] ({(\v-9.4)*100},-1.2) -- ({(\v-9.4)*100},1.2);}
  \fill[bkOrange] (60,0) circle (1.3);
  \node[anchor=south, font=\sffamily\bfseries\small, text=bkOrange!80!black] at (59,2.2) {10};
  \foreach \v/\hh in {10.5/4, 10.05/10} {\fill[bkBlue] ({(\v-9.4)*100},0) circle (0.9);
    \draw[bkBlue, line width=0.5pt] ({(\v-9.4)*100},0) -- ({(\v-9.4)*100},\hh-0.6);
    \node[anchor=south, font=\sffamily\scriptsize, text=bkBlue] at ({(\v-9.4)*100},\hh) {\v};}
  \foreach \v/\hh in {9.5/4, 9.95/10} {\fill[bkGreen!70!black] ({(\v-9.4)*100},0) circle (0.9);
    \draw[bkGreen!70!black, line width=0.5pt] ({(\v-9.4)*100},0) -- ({(\v-9.4)*100},-\hh+0.6);
    \node[anchor=north, font=\sffamily\scriptsize, text=bkGreen!70!black] at ({(\v-9.4)*100},-\hh) {\v};}
  \draw[bk arrow, bkBlue, line width=0.9pt] (118,16) -- (70,16);
  \node[anchor=south, font=\sffamily\footnotesize, text=bkBlue] at (94,17) {後ろから（10より大きい側）};
  \draw[bk arrow, bkGreen!70!black, line width=0.9pt] (2,-16) -- (50,-16);
  \node[anchor=north, font=\sffamily\footnotesize, text=bkGreen!70!black] at (26,-17) {前から（10より小さい側）};
\end{tikzpicture}
\bkfigcaption{数直線で見ると、両側から10をはさみながら近づいていく。}
\end{bkfigure}

後ろからは10より少し大きい値で、前からは10より少し小さい値で、どちらも10に近づいていきます。10を両側からはさみうちにしているので、近づく先は10しかありません。だから、\textbf{1秒の瞬間の速さは\qty{10}{m/s}}だと考えられます。

\bkhatenaref の「秒速\qty{5}{m}？」は、0秒から1秒までの平均の速さでした。瞬間の速さ10はその2倍です。なぜ2倍なのでしょう。あとでわかるように、このボールの速さは0秒で0、1秒で10と、時間に比例して同じ割合でふえていきます。同じ割合でふえるとき、0から10までをならした平均は、まん中の $(0+10)\div2=5$ です。だから、平均の5と瞬間の10は、ちょうど2倍のちがいになります。

\bksubsection{2}{文字 $h$ を使って、近づき方を式にする}

\leavevmode\label{bk:5-2:h}表では、$\Delta t$ を変えるたびに計算し直しました。0.1でも、0.01でも、同じ計算のくり返しです。そこで、あいだを文字 $h$ で表して、1回の計算ですませましょう。$h$ は、0.1 や 0.01 など、どんな数でも入れられる箱です（2-11）。ただし、あとで $h$ でわるので、$h$ は0でない数とします。

\noindent\textbf{①\enspace $1+h$ 秒の位置}\par
$x=5t^2$ の $t$ に、かっこごと $1+h$ を入れます（代入、2-13）。
\[ x(1+h)=5(1+h)^2 \]

\noindent\textbf{②\enspace かっこをはずす（展開）}\par
$(1+h)^2$ は、1辺が $1+h$ の正方形の面積です。図のように4つに分けると、
\[ (1+h)^2=1+h+h+h^2=1+2h+h^2 \]

\noindent\begin{minipage}[c]{0.52\linewidth}
\begin{tikzpicture}[x=1mm, y=1mm]
  % 1辺 1+h の正方形（1→34 mm、h→9 mm。h は見やすいように大きめにかいた）
  \fill[bkBlue!14] (0,9) rectangle (34,43);
  \fill[bkOrange!25] (34,9) rectangle (43,43);
  \fill[bkOrange!25] (0,0) rectangle (34,9);
  \fill[bkRed!35] (34,0) rectangle (43,9);
  \draw[bkInk, line width=0.8pt] (0,0) rectangle (43,43);
  \draw[bkInk, line width=0.6pt] (34,0) -- (34,43) (0,9) -- (43,9);
  \node[font=\sffamily\small] at (17,26) {$1\times1=1$};
  \node[font=\sffamily\small] at (38.5,26) {$h$};
  \node[font=\sffamily\small] at (17,4.5) {$h$};
  \node[font=\sffamily\scriptsize] at (38.5,4.5) {$h^2$};
  \node[anchor=south, font=\sffamily\footnotesize] at (17,44.5) {1};
  \node[anchor=south, font=\sffamily\footnotesize] at (38.5,44.5) {$h$};
  \node[anchor=east, font=\sffamily\footnotesize] at (-1.5,26) {1};
  \node[anchor=east, font=\sffamily\footnotesize] at (-1.5,4.5) {$h$};
\end{tikzpicture}
\end{minipage}\hfill
\begin{minipage}[c]{0.45\linewidth}\raggedright
大きな正方形 1 と、細長い長方形 $1\times h=h$ が2まい、小さな正方形 $h\times h=h^2$ が1まいです。

$h=0.1$ で確かめると、$1.1^2=1.21$、\\$1+2\times0.1+0.01=1.21$ で合います。
\end{minipage}

\medskip
だから、$x(1+h)=5(1+2h+h^2)$。かっこの中の1つずつに5をかけて（分配法則、2-12）、
\[ x(1+h)=5+10h+5h^2 \]

\noindent\textbf{③\enspace 進んだ距離 $\Delta x$（あと${}-{}$まえ）}\par
1秒の位置は $x(1)=5$ なので、
\[ \Delta x=(5+10h+5h^2)-5=10h+5h^2 \]

\noindent\textbf{④\enspace $h$ でわる}\par
まとめてわるときは、1つずつわってよいので（2-12）、
\[ \frac{\Delta x}{h}=\frac{10h+5h^2}{h}=\frac{10h}{h}+\frac{5h^2}{h}=10+5h \]
$\dfrac{10h}{h}$ は、分子と分母の $h$ を約分して10です。$\dfrac{5h^2}{h}=\dfrac{5\times h\times h}{h}$ は、$h$ を1つ約分して $5h$ です。

\medskip
こうして、1秒から $1+h$ 秒までの平均の速さは、
\[ \frac{\Delta x}{\Delta t}=10+5h \]
という式になりました。表と合っているか、確かめましょう。
\begin{align*}
h=0.1\ \text{なら}&\quad 10+5\times0.1=10+0.5=10.5\\
h=0.01\ \text{なら}&\quad 10+5\times0.01=10+0.05=10.05
\end{align*}
表の10.5、10.05とぴったり合います。

\begin{bkfigure}
\begin{tikzpicture}[x=1mm, y=1mm]
  % 第6回：左に行の見出し、箱の中は数だけ。矢印は箱の名前どうしでつなぎ、両端を離す
  \node[anchor=east, font=\sffamily\small, text=bkInk, inner sep=0pt] at (14,3.2) {$h$};
  \node[anchor=east, font=\sffamily\small, text=bkInk, inner sep=0pt] at (14,-3.2) {$10+5h$};
  \foreach \h/\v/\xx/\nm in {1/15/27/b1, 0.1/10.5/61/b2, 0.01/10.05/95/b3, 0.001/10.005/129/b4} {
    \node[bk node, minimum width=22mm, inner xsep=2mm, font=\sffamily\small] (\nm) at (\xx,0) {$\h$\\[1pt]$\v$};}
  \foreach \a/\b in {b1/b2,b2/b3,b3/b4} \draw[bk link] (\a) -- (\b);
  \draw[bk link, bkOrange, shorten <=1.2mm] (b4.south) -- ++(0,-6.5) coordinate (ar);
  \node[anchor=north, font=\sffamily\bfseries\small, text=bkOrange!80!black, inner sep=0pt] at ([yshift=-1.5mm]ar) {10 に近づく};
  \node[anchor=north west, font=\sffamily\small, text=bkInk, inner sep=0pt] at (16,-9.5) {$h$ が小さくなると、$5h$ の部分が、どんどん小さくなる};
\end{tikzpicture}
\end{bkfigure}

$h$ を小さくすると、$10+5h$ の $5h$ の部分がどんどん0に近づき、全体は10に近づきます。表では「10に近づいていそう」としか言えませんでしたが、式を見れば、近づく先が10だとはっきりわかります。

\noindent\textbf{10と $5h$ は、何を表している？}\par
$\Delta x=10h+5h^2$ の2つの部分にも、意味があります。$10h$ は、「1秒のときの速さ\qty{10}{m/s}のまま、$h$ 秒走ったときの距離」です（距離＝速さ×時間）。$5h^2$ は、そのあいだにボールが速くなった分だけ、よけいに進んだ距離です。$h$ が小さいと、$h^2$ は $h$ よりずっと小さくなります（$h=0.01$ なら $h^2=0.0001$）。だから、短いあいだでは「よけいな分」はほとんど消えて、平均の速さは10にほぼ等しくなるのです。

\bksubsection{3}{どの時刻でも計算できる}

\leavevmode\label{bk:5-2:t}1秒のかわりに、どの時刻 $t$ でも同じ計算ができます。$t$ 秒から $t+h$ 秒までを考えます。ここでは $t$ は「いま考えている時刻」で、計算のあいだは決まった数として扱います。$h$ は、そこからのあいだです。

②の正方形の図の「1」を「$t$」にかえると、$(t+h)^2=t^2+2th+h^2$ です（3-8）。①〜④と同じ段で、1行ずつ計算します（1・2行め：代入して展開、3行め：あと${}-{}$まえ、4行め：1つずつ $h$ でわる）。
\begin{align*}
x(t+h)&=5(t+h)^2=5(t^2+2th+h^2)\\
&=5t^2+10th+5h^2\\
\Delta x&=(5t^2+10th+5h^2)-5t^2=10th+5h^2\\
\frac{\Delta x}{h}&=\frac{10th+5h^2}{h}=10t+5h
\end{align*}
$t=1$ を入れると $10+5h$ で、さっきの式と同じになります。$h$ を0に近づけると $5h$ が消えて、$t$ 秒の瞬間の速さ $v$ は
\[ v=10t\quad\text{（m/s）} \]
です（$v$ は速さを表す文字。velocity の頭の文字、2-11）。

\begin{bkcenter}
\begin{bktable}{colspec={Q[l,m,wd=52mm] *{5}{Q[c,m,wd=9mm]}}, rowsep=1.5pt}
時刻 $t$（s） & 0 & 1 & 2 & 3 & 4 \\
位置 $x=5t^2$（m） & 0 & 5 & 20 & 45 & 80 \\
瞬間の速さ $v=10t$（m/s） & 0 & 10 & 20 & 30 & 40 \\
\end{bktable}
\bktabcaption{落ちるボールの位置と、瞬間の速さ}
\end{bkcenter}

瞬間の速さは、1秒ごとに\qty{10}{m/s}ずつ増えています。2-0で「落ちるボールの速さは、1秒ごとに約\qty{10}{m/s}ずつふえる」と書いたのは、このことです。式から、それが確かめられました。

\bksubsection{4}{$x=at^2$ なら $v=2at$}

\leavevmode\label{bk:5-2:a}5のかわりに、どんな数 $a$ でも同じ計算ができます。$x=at^2$ とします（$a$ は決まった数）。
\begin{align*}
x(t+h)-x(t)&=a(t+h)^2-at^2\\
&=at^2+2ath+ah^2-at^2=2ath+ah^2\\
\frac{\Delta x}{h}&=\frac{2ath+ah^2}{h}=2at+ah
\end{align*}
$h$ を0に近づけると $ah$ が消えて、
\[ v=2at \]
です。2という数は、展開の $(t+h)^2=t^2+\underline{2th}+h^2$ の「2」から来ています。正方形の図で、細長い長方形が\textbf{2まい}あったからです。

\begin{bkcenter}
\begin{bktable}{colspec={Q[c,m,wd=26mm] Q[c,m,wd=20mm] Q[c,m,wd=40mm]}, rowsep=1.5pt}
位置の式 & $a$ & 瞬間の速さ $v=2at$ \\
$x=5t^2$ & 5 & $v=2\times5\times t=10t$ \\
$x=3t^2$ & 3 & $v=2\times3\times t=6t$ \\
$x=20t^2$ & 20 & $v=2\times20\times t=40t$ \\
$x=4.9t^2$ & 4.9 & $v=2\times4.9\times t=9.8t$ \\
\end{bktable}
\bktabcaption{$x=at^2$ の瞬間の速さ}
\end{bkcenter}

3行めの $x=20t^2$ は、5-1のボールです。$v=40t$ に $t=1$ を入れると40で、5-1で表から見つけた\qty{40}{cm/s}と合います。

4行めの $x=4.9t^2$ が、ほんとうの落ちる物の式です。$v=9.8t$ なので、落ちる物は、1秒ごとに\qty{9.8}{m/s}ずつ速くなります（5-5）。\bkhatenaref で「ほんとうは4.9」と書いた4.9は、この9.8の半分です。

\noindent\textbf{同じ速さで動く物なら}\par
位置が $x=20t$（m）の物は、どうでしょう。
\[ \frac{x(t+h)-x(t)}{h}=\frac{20(t+h)-20t}{h}=\frac{20t+20h-20t}{h}=\frac{20h}{h}=20 \]
$h$ がどんな数でも、平均の速さは20のままです。グラフがまっすぐな直線なので、どのあいだをとっても傾きは20です。だから、瞬間の速さもいつも\qty{20}{m/s}です。

\bksubsection{5}{グラフで見ると}

\noindent\begin{minipage}[c]{0.5\linewidth}
\begin{tikzpicture}
\begin{axis}[bkplot, width=66mm, height=48mm, xmin=0, xmax=3.4, ymin=0, ymax=56, xtick={0,1,2,3}, ytick={0,10,...,50},
  xlabel={$t$（s）}, ylabel={$x$（m）}]
  % 第6回：接線は細くして曲線の上に重ね、曲線が見えるようにする。接線は枠（ymin〜ymax）の中で止める
  \addplot[bkcurve, domain=0:3.1] {5*x^2};
  \foreach \t/\a/\b in {1/0.56/2.0, 2/1.25/2.9, 3/2.1/3.36} {\edef\tmp{\noexpand\addplot[bkcurve2, line width=1pt, domain=\a:\b] {5*\t*\t+10*\t*(x-\t)};}\tmp}
  \addplot[only marks, mark=*, mark size=1.6pt, bkInk] coordinates {(1,5) (2,20) (3,45)};
\end{axis}
\end{tikzpicture}
\end{minipage}\hfill
\begin{minipage}[c]{0.46\linewidth}\raggedright
$x$-$t$ グラフの各点に、瞬間の速さを傾きにもつ直線（オレンジ）を引きました。傾きは、1秒で10、2秒で20、3秒で30。グラフがだんだん急になっていくようすが、傾きの数で表せました。
\end{minipage}

\bksubsection{6}{前からの近づき方も、式で確かめる}

$1-h$ 秒から1秒までの平均の速さも、$(a-b)^2=a^2-2ab+b^2$（3-8）を使って、式にしてみましょう。こんどは1秒が「あと」、$1-h$ 秒が「まえ」です。
\begin{align*}
x(1-h)&=5(1-h)^2=5(1-2h+h^2)=5-10h+5h^2 \\
\Delta x&=5-(5-10h+5h^2)=5-5+10h-5h^2=10h-5h^2 \\
\frac{\Delta x}{h}&=\frac{10h-5h^2}{h}=10-5h
\end{align*}
2行めでは、マイナスのかっこをはずすと、中の符号が全部かわることを使いました（2-12）。$h=0.1$ なら $10-0.5=9.5$ で、表と合います。

後ろからは $10+5h$、前からは $10-5h$。$h$ を0に近づけると、どちらも10になります。近づく先が同じなので、「1秒の瞬間の速さは10」と決められます。

\begin{taisetsu}
平均の速さを $h$ の式で表し、$h$ を0に近づける。残った値が瞬間の速さ。\\
$x=5t^2$ のとき、平均の速さは $10t+5h$、瞬間の速さは $v=10t$\\
$x=at^2$ のとき、瞬間の速さは $v=2at$
\end{taisetsu}

\begin{kotoba}
\small
\bkword{$h$}{あいだの長さを表す文字。0ではない、小さい数。}
\bkword{近づく}{$h$ を0に近づけたとき、式の値がある数にどんどん近くなること。5-3で「極限」という名前をつける。}
\bkword{$v$}{速さを表す文字（velocity の頭の文字）。}
\end{kotoba}

%% ---------------- 例題 ----------------
\bksection{reidai}{例題でたしかめよう}

\begin{reidai}[表で近づき方を見る]
$x=5t^2$（m）のボールについて、2秒から $2+\Delta t$ 秒までの平均の速さを、$\Delta t=0.1$ と $0.01$ で求め、2秒の瞬間の速さを答えましょう。\par
\kangaekata はじめに2秒の位置を求めます。次に、後ろの時刻の位置を計算して、（あと${}-{}$まえ）÷あいだ を求めます。\par
\shiki $x(2)=5\times2^2=5\times4=20$\par
\hspace*{2.6\zw}$x(2.1)=5\times2.1^2=5\times4.41=22.05$\par
\hspace*{2.6\zw}$x(2.01)=5\times2.01^2=5\times4.0401=20.2005$\par
\hspace*{2.6\zw}$(22.05-20)\div0.1=2.05\div0.1=20.5$\par
\hspace*{2.6\zw}$(20.2005-20)\div0.01=0.2005\div0.01=20.05$\par
\kotae 20.5、20.05と20に近づくので、\qty{20}{m/s}
\end{reidai}

\begin{jibun}
\bkreidairef{1}のボールの3秒の瞬間の速さを、同じ方法で求めましょう。
\bkjans{$x(3)=45$。$x(3.1)=48.05$ より 30.5、$x(3.01)=45.3005$ より 30.05。\qty{30}{m/s}}
\end{jibun}

\begin{reidai}[$h$ の式で表す]
$x=5t^2$ のボールについて、2秒から $2+h$ 秒までの平均の速さを $h$ の式で表し、2秒の瞬間の速さを求めましょう。\par
\kangaekata 代入して展開してから、「あと${}-{}$まえ」をとり、$h$ でわります。$(2+h)^2=4+4h+h^2$（正方形の図の1を2にかえる）。\par
\shiki $x(2+h)=5(2+h)^2=5(4+4h+h^2)=20+20h+5h^2$\par
\hspace*{2.6\zw}$\Delta x=(20+20h+5h^2)-20=20h+5h^2$\par
\hspace*{2.6\zw}$\dfrac{\Delta x}{h}=\dfrac{20h+5h^2}{h}=20+5h$\par
\kotae 平均の速さは $20+5h$。$h$ を0に近づけて、瞬間の速さは\qty{20}{m/s}
\end{reidai}

\begin{jibun}
$x=5t^2$ について、3秒から $3+h$ 秒までの平均の速さを $h$ の式で表し、3秒の瞬間の速さを求めましょう。
\bkjans{$5(3+h)^2=45+30h+5h^2$ より $30+5h$。\qty{30}{m/s}}
\end{jibun}

\begin{reidai}[ちがう式の運動]
位置が $x=3t^2$（m）で表される物体の、1秒の瞬間の速さを求めましょう。\par
\kangaekata 5のかわりに3になっただけです。同じ4つの段（代入・展開、あと${}-{}$まえ、$h$ でわる、$h$ を0に近づける）で計算します。\par
\shiki $x(1)=3\times1^2=3$\par
\hspace*{2.6\zw}$x(1+h)=3(1+h)^2=3(1+2h+h^2)=3+6h+3h^2$\par
\hspace*{2.6\zw}$\Delta x=(3+6h+3h^2)-3=6h+3h^2$\par
\hspace*{2.6\zw}$\dfrac{\Delta x}{h}=\dfrac{6h+3h^2}{h}=6+3h$\par
\kotae $h$ を0に近づけて、\qty{6}{m/s}（$v=2at=2\times3\times1=6$ とも合う）
\end{reidai}

\begin{jibun}
位置が $x=2t^2$（m）の物体の、3秒の瞬間の速さを求めましょう。
\bkjans{$2(3+h)^2=18+12h+2h^2$ より $12+2h$。\qty{12}{m/s}}
\end{jibun}

\begin{reidai}[項が2つある式]
位置が $x=2t^2+t$（m）の物体の、1秒の瞬間の速さを求めましょう。\par
\kangaekata $t$ が2か所にあります。どちらの $t$ にも、かっこごと $1+h$ を入れることを忘れずに。\par
\shiki $x(1)=2\times1^2+1=3$\par
\hspace*{2.6\zw}$x(1+h)=2(1+h)^2+(1+h)$\par
\hspace*{2.6\zw}$\phantom{x(1+h)}=2+4h+2h^2+1+h=3+5h+2h^2$\par
\hspace*{2.6\zw}$\Delta x=(3+5h+2h^2)-3=5h+2h^2$\par
\hspace*{2.6\zw}$\dfrac{\Delta x}{h}=\dfrac{5h+2h^2}{h}=5+2h$\par
\kotae $h$ を0に近づけて、\qty{5}{m/s}
\end{reidai}

\begin{jibun}
位置が $x=t^2+3t$（m）の物体の、2秒の瞬間の速さを求めましょう。
\bkjans{$x(2)=10$、$x(2+h)=10+7h+h^2$ より $7+h$。\qty{7}{m/s}}
\end{jibun}

\begin{reidai}[$t$ 秒の速さの式を使う]
$x=5t^2$（m）のボールの $t$ 秒の瞬間の速さを $t$ の式で表し、速さが\qty{30}{m/s}になるのは何秒後か求めましょう。そのとき、何m落ちていますか。\par
\kangaekata $t$ 秒から $t+h$ 秒までの平均の速さは $10t+5h$ でした（\pageref{bk:5-2:t}ページ）。$h$ を0に近づけて速さの式を作り、それが30になる $t$ を求めます。\par
\shiki $h$ を0に近づけて $v=10t$\par
\hspace*{2.6\zw}$10t=30$。両辺を10でわって（2-14）$t=3$\par
\hspace*{2.6\zw}$x=5\times3^2=5\times9=45$\par
\noindent\begin{minipage}[c]{0.5\linewidth}
\begin{tikzpicture}
\begin{axis}[bkplot, width=56mm, height=38mm, xmin=0, xmax=4.3, ymin=0, ymax=44, xtick={0,1,2,3,4}, ytick={0,10,20,30,40},
  xlabel={$t$（s）}, ylabel={$v$（m/s）}]
  \addplot[bkcurve2, domain=0:4] {10*x};
  \draw[bkGray, line width=0.6pt, dash pattern=on 2pt off 1.5pt] (axis cs:0,30) -- (axis cs:3,30) -- (axis cs:3,0);
  \addplot[only marks, mark=*, mark size=1.7pt, bkInk] coordinates {(3,30)};
\end{axis}
\end{tikzpicture}
\end{minipage}\hfill
\begin{minipage}[c]{0.46\linewidth}\raggedright
$v=10t$ のグラフ（速さと時間のグラフ）。たて軸の30から横に進み、グラフに当たった所の $t$ を読むと3秒です。
\end{minipage}\par
\kotae $v=10t$、3秒後、\qty{45}{m}
\end{reidai}

\begin{jibun}
$x=2t^2$（m）の物体の $t$ 秒の瞬間の速さを $t$ の式で表し、速さが\qty{12}{m/s}になる時刻と、そのときの位置を求めましょう。
\bkjans{$2(t+h)^2-2t^2=4th+2h^2$ より $v=4t$。$4t=12$ より3秒後、$x=\qty{18}{m}$}
\end{jibun}

\begin{reidai}[上に投げたボール（速さがマイナスになる）]
ボールを真上に投げました。手の位置を0、上向きを正として、$t$ 秒後の高さは $x=10t-5t^2$（m）です。1.5秒の瞬間の速さを求めましょう。\par
\kangaekata $-5t^2$ の項にも、かっこごと入れます。マイナスのかっこをはずすときは、中の符号が全部かわります（2-12）。\par
\shiki $x(1.5)=10\times1.5-5\times2.25=15-11.25=3.75$\par
\hspace*{2.6\zw}$x(1.5+h)=10(1.5+h)-5(1.5+h)^2$\par
\hspace*{2.6\zw}$\phantom{x(1.5+h)}=15+10h-5(2.25+3h+h^2)$\par
\hspace*{2.6\zw}$\phantom{x(1.5+h)}=15+10h-11.25-15h-5h^2=3.75-5h-5h^2$\par
\hspace*{2.6\zw}$\Delta x=(3.75-5h-5h^2)-3.75=-5h-5h^2$\par
\hspace*{2.6\zw}$\dfrac{\Delta x}{h}=\dfrac{-5h-5h^2}{h}=-5-5h$\par
\kotae $h$ を0に近づけて、$\qty{-5}{m/s}$。マイナスは下向き。1.5秒には、ボールは\qty{5}{m/s}で落ちてきている。
\end{reidai}

\begin{jibun}
\bkreidairef{6}のボールの、0.5秒の瞬間の速さを求めましょう。上向きと下向きのどちらに動いていますか。
\bkjans{$x(0.5)=3.75$、$x(0.5+h)=3.75+5h-5h^2$ より $5-5h$。\qty{5}{m/s}で上向き}
\end{jibun}

\begin{reidai}[$v=2at$ を使う]
5-1の坂をころがるボールの位置は $x=20t^2$（cm）でした。$v=2at$ を使って $t$ 秒の瞬間の速さを式で表し、速さが\qty{100}{cm/s}になる時刻と、そのときの位置を求めましょう。\par
\kangaekata $x=at^2$ の $a$ が20です。$v=2at$ の $a$ に20を入れます。\par
\shiki $v=2\times20\times t=40t$\par
\hspace*{2.6\zw}$40t=100$。両辺を40でわって $t=2.5$\par
\hspace*{2.6\zw}$x=20\times2.5^2=20\times6.25=125$\par
\kotae $v=40t$（cm/s）、2.5秒、\qty{125}{cm}
\end{reidai}

\begin{jibun}
ほんとうの落ちる物の式は $x=4.9t^2$（m）です。$v=2at$ を使って $t$ 秒の瞬間の速さを式で表し、速さが\qty{19.6}{m/s}になる時刻と、それまでに落ちた距離を求めましょう。
\bkjans{$v=9.8t$。$9.8t=19.6$ より2秒、$4.9\times2^2=\qty{19.6}{m}$}
\end{jibun}

\begin{column}{ガリレオは、どうやって時間をはかった？}
ガリレオのころ（400年ほど前）には、ストップウォッチはありませんでした。ガリレオは、大きな水おけの底に細い管をつけ、球が坂をころがっているあいだだけ水を流し、出てきた水の重さをはかりました。大きな水おけなので、水面の高さはほとんど変わらず、水はほぼ同じ速さで流れ出ます。だから、水の重さは時間に比例します。水の重さが2倍なら、時間も2倍です。こうして、「進む距離は時間の2乗に比例する」ことを確かめたのです。「はかれない量を、はかれる量に置きかえる」という、計測の工夫です。
\end{column}

\clearpage
%% ---------------- まちがえやすいところ・物理工学 ----------------
\bksection{tsumazuki}{まちがえやすいところをチェック}

\begin{tsumazuki}{いきなり $h=0$ を入れてしまう}
\vspace{7pt}% 1行目が分数で高いので、題とのあいだをほかの箱とそろえる
$\dfrac{10h+5h^2}{h}$ に $h=0$ を入れると $\dfrac{0}{0}$ になり、計算できません。0でわることはできないからです（1-14）。先に $h$ でわって $10+5h$ にしてから、$h$ を0に近づけます。$h$ そのものは、最後まで0ではありません。
\end{tsumazuki}

\begin{tsumazuki}{展開で $2h$ をわすれる}
\bkng{$(1+h)^2=1+h^2$}\bkok{$(1+h)^2=(1+h)(1+h)=1+h+h+h^2=1+2h+h^2$}
1辺 $1+h$ の正方形を4つに分けた図を思い出すと、まちがえません（3-8）。$h=0.1$ で確かめると、$1.1^2=1.21$ ですが、$1+0.1^2=1.01$ で合いません。
\end{tsumazuki}

\begin{tsumazuki}{「あと${}-{}$まえ」の「まえ」をひき忘れる}
\bkng{$\dfrac{x(1+h)}{h}=\dfrac{5+10h+5h^2}{h}$}\bkok{$\dfrac{x(1+h)-x(1)}{h}=\dfrac{10h+5h^2}{h}$}
平均の速さは、進んだ距離（変化量）÷時間です。$x(1+h)$ は「位置」で、進んだ距離ではありません。
\end{tsumazuki}

\begin{tsumazuki}{2か所の $t$ の片方にだけ入れる}
\bkng{$x=2t^2+t$ で、$x(1+h)=2(1+h)^2+1$}\bkok{$x(1+h)=2(1+h)^2+(1+h)$}
$t$ がある所は、全部かっこごと入れかえます（2-13）。
\end{tsumazuki}

\begin{butsuri}
\begin{itemize}[leftmargin=1.2\zw, itemsep=2pt, topsep=0pt, label={\color{bkPurple}●}]
\item \textbf{ガリレオの斜面の実験}\par
400年ほど前、ガリレオは坂をころがる球を調べ、進む距離が時間の2乗に比例する（$x=at^2$）ことを見つけました。この章の計算から、速さは $v=2at$ で、時間に比例して増えることがわかります。
\item \textbf{車の加速性能}\par
「時速\qty{100}{km}まで○秒」という性能は、瞬間の速さがどれだけ早く増えるかを表しています。
\item \textbf{落下の衝撃}\par
$v=9.8t$ なので、落ちる時間が長いほど速くなります。2秒落ちると約\qty{20}{m/s}（時速約\qty{70}{km}）にもなります。高い所から物を落とさないのは、このためです。
\end{itemize}
\end{butsuri}

\begin{column}{表計算で $h$ を小さくしすぎると}
やってみようの表計算で、$h$ をもっと小さくしてみましょう。B列は、やってみようと同じ「（あと${}-{}$まえ）÷あいだ」で計算した平均の速さ、C列は、$h$ の式 $10+5h$ で計算した値です。
\begin{bkcenter}
\bksheet{3}{30}{%
1 & $h$ & 平均の速さ & $10+5h$ \\
2 & 0.1 & 10.5 & 10.5 \\
3 & 0.001 & 10.005 & 10.005 \\
4 & 0.00001 & 10.00005 & 10.00005 \\
5 & $10^{-8}$ & {\color{bkRed}\bfseries 9.9999999392…} & 10.00000005 \\
6 & $10^{-10}$ & {\color{bkRed}\bfseries 10.0000008274…} & 10.0000000005 \\
7 & $10^{-12}$ & {\color{bkRed}\bfseries 10.0008890058…} & 10.000000000005 \\}
\bktabcaption{B列は、$h$ が小さすぎると10からはなれる（赤）。C列は10に近づいていく。}
\end{bkcenter}
$h=10^{-8}$ より小さくすると、B列の値は10からはなれてしまいます。コンピュータは数を16けたほどしか覚えていないので、とても近い2つの数をひき算すると、下のけたのずれが目立つのです（2-19）。$h$ の式 $10+5h$ を使えば、このずれは起きません。式で考えることの強みです。\par\vspace{4pt}
\textbf{ためしてみよう}\par
自分の表計算でも、$h$ を $\dfrac{1}{10}$ ずつ小さくして、B列が何回めから10からはなれはじめるか調べてみよう。
\end{column}

\clearpage
%% ---------------- 章末問題 ----------------
\bkendproblems
\bkproblemgroup{基本}{2}

\bkproblem 位置が $x=5t^2$（m）のボールについて、4秒から $4+\Delta t$ 秒までの平均の速さを、$\Delta t=0.1$ と $0.01$ で求めましょう。4秒の瞬間の速さは、いくつと考えられますか。
\bkpans{40.5（$x(4.1)=84.05$）、40.05（$x(4.01)=80.4005$）。\qty{40}{m/s}}

\bkproblem 次の式を展開しましょう。\quad（1）$(1+h)^2$\quad（2）$(3+h)^2$\quad（3）$(t+h)^2$
\bkpans{（1）\nobreak $1+2h+h^2$\quad（2）\nobreak $9+6h+h^2$\quad（3）\nobreak $t^2+2th+h^2$}

\bkproblem 位置が $x=4t^2$（m）の物体について、2秒から $2+h$ 秒までの平均の速さを $h$ の式で表し、2秒の瞬間の速さを求めましょう。
\bkpans{$4(2+h)^2=16+16h+4h^2$ より $16+4h$。\qty{16}{m/s}}

\bkproblem 位置が $x=t^2$（m）の物体の、5秒の瞬間の速さを求めましょう。
\bkpans{$(5+h)^2=25+10h+h^2$ より $10+h$。\qty{10}{m/s}}

\bkproblem 位置が $x=3t^2$（m）の物体の、$t$ 秒の瞬間の速さを $t$ の式で表しましょう。また、4秒の瞬間の速さを求めましょう。
\bkpans{$3(t+h)^2-3t^2=6th+3h^2$ より $v=6t$。4秒で\qty{24}{m/s}}

\bkproblem 位置が $x=2t^2+3t$（m）の物体の、1秒の瞬間の速さを求めましょう。
\bkpans{$x(1)=5$、$x(1+h)=5+7h+2h^2$ より $7+2h$。\qty{7}{m/s}}

\bkproblem 次の物体の、$t$ 秒の瞬間の速さを $v=2at$ を使って求めましょう。\quad（1）$x=6t^2$\quad（2）$x=0.5t^2$
\bkpans{（1）\nobreak $v=12t$\quad（2）\nobreak $v=t$}

\bkproblemgroup{チャレンジ}{2}

\bkproblem 位置が $x=3t^2+2t$（m）の物体の、2秒の瞬間の速さを求めましょう。
\bkpans{$x(2)=16$、$x(2+h)=16+14h+3h^2$ より $14+3h$。\qty{14}{m/s}}

\bkproblem $x=5t^2$（m）で落ちるボールの瞬間の速さは $v=10t$（m/s）でした。速さが\qty{25}{m/s}になるのは何秒後ですか。また、それまでに何m落ちていますか。
\bkpans{2.5秒後。$5\times2.5^2=\qty{31.25}{m}$}

\bkproblem 位置が $x=t^2-4t$（m）の物体があります。1秒の瞬間の速さを求めましょう。答えがマイナスになったら、それは何を表していますか。
\bkpans{$x(1)=-3$、$x(1+h)=-3-2h+h^2$ より $-2+h$。$\qty{-2}{m/s}$。$x$ が小さくなる向き（もどる向き）に、\qty{2}{m/s}で動いている}

\bkproblem 同じ所を同時に出発した2つの物体があります。Aの位置は $x=5t^2$、Bの位置は $x=20t$（m）です。2つの速さが同じになるのは何秒後ですか。そのとき、前にいるのはどちらで、何mはなれていますか。
\bkpans{Aの速さ $10t$、Bは\qty{20}{m/s}。$10t=20$ より2秒後。A $20$\,m、B $40$\,m で、Bが\qty{20}{m}前}

\begin{furikaeri}
\begin{checklist}
\item $\Delta t$ を小さくして、平均の速さを表で計算できる。\hfill{\footnotesize\sffamily\color{bkTheme}→\pageref{bk:5-2:hyou}ページ}
\item 平均の速さを $h$ の式で表せる。\hfill{\footnotesize\sffamily\color{bkTheme}→\pageref{bk:5-2:h}ページ}
\item $x=at^2$ の瞬間の速さを求められる。\hfill{\footnotesize\sffamily\color{bkTheme}→\pageref{bk:5-2:a}ページ}
\end{checklist}
\end{furikaeri}

\bkfukidashi{sora}{$h$ を0に「近づける」のと、0を「入れる」のはちがうんだね。次の章で、この考えに名前をつけるよ。}

%% 5-3 かぎりなく近づく：極限と接線（第11回：ページ制限なし。偶数ページ）
\bkchapter{5}{3}{かぎりなく近づく：\\極限と接線}{\bklevel{3}}{かぎりなく近づく：極限と接線}
\bkmichishirube{3-5,4-2,5-2}{5-4,5-8}

\begin{goal}
\begin{itemize}[label={\color{bkTheme}●}, leftmargin=1.6\zw, itemsep=0pt, topsep=1pt]
\item \bklevel{3}\enspace $\lim$ の記号を読み、極限の値を求められる。
\item \bklevel{3}\enspace 微分係数を、接線の傾きとして求められる。
\item \bklevel{3}\enspace 導関数を求め、接線の式を書ける。\bkidx{せっせんのしき}{接線の式}
\end{itemize}
\end{goal}

\begin{hatena}
\noindent\begin{minipage}[c]{0.55\linewidth}
スケートボードのランプ（すべり台のような曲がった坂）を考えます。坂の急さは、場所ごとにちがいます。下のほうはゆるやかで、上へ行くほど急です。

では、点Pのところの「坂の急さ」は、どんな数で表せばよいでしょう。まっすぐな坂なら傾きで表せました（3-5）。でも、曲がった坂には、決まった傾きがありません。
\end{minipage}\hfill
\begin{minipage}[c]{0.42\linewidth}
\begin{tikzpicture}
\begin{axis}[width=58mm, height=46mm, axis lines=none, xmin=-0.2, xmax=2.3, ymin=-0.3, ymax=4.6]
  \fill[bkYellow!25] (axis cs:0,0) -- plot[domain=0:2, samples=40] ({\x},{\x^2}) -- (axis cs:2,0) -- cycle;
  \addplot[bkcurve, domain=0:2] {x^2};
  \draw[bkGray, line width=1pt] (axis cs:-0.2,0) -- (axis cs:2.3,0);
  \addplot[only marks, mark=*, mark size=2pt, bkInk] coordinates {(1.2,1.44)};
  \node[anchor=south east, font=\sffamily\small] at (axis cs:1.2,1.44) {P};
  \addplot[bkcurve2, line width=1pt, dashed, domain=0.65:1.72] {1.44+2.4*(x-1.2)};
  \node[anchor=west, font=\sffamily\scriptsize, text=bkGray] at (axis cs:1.78,2.9) {?};
\end{axis}
\end{tikzpicture}
\end{minipage}
\end{hatena}

5-2では、時間のあいだ $h$ を0に近づけて、瞬間の速さを求めました。同じ考え方で、曲がったグラフの1点での「傾き」が決められます。この章では、その考え方に名前と記号をつけます。

\clearpage
%% ---------------- やってみよう・思い出そう ----------------
\begin{yattemiyou}[グラフにぴったりふれる直線を引こう]
\textbf{用意するもの}\par
方眼紙、定規
\begin{enumerate}[label={\bkcircled[bkGreen]{\arabic*}}, leftmargin=2.4\zw, itemsep=1pt, topsep=3pt]
\item 下の表の点を方眼紙にとり、なめらかな線でつないで、$y=x^2$ のグラフをかこう。1目もりは0.5にする。
\item 点P（1，1）と点Q（2，4）を結ぶ直線を引き、傾きを求めよう。
\item Qを（1.5，2.25）にかえて、同じように傾きを求めよう。
\item 定規を点Pに当て、グラフにぴったりふれるように直線を引こう。その傾きは、いくつになったかな？
\end{enumerate}
\begin{bkcenter}
\begin{bktable}{colspec={Q[l,m,wd=14mm] *{9}{Q[c,m,wd=9mm]}}, rowsep=1pt, colsep=3pt}
$x$ & $-2$ & $-1.5$ & $-1$ & $-0.5$ & 0 & 0.5 & 1 & 1.5 & 2 \\
$y=x^2$ & 4 & 2.25 & 1 & 0.25 & 0 & 0.25 & 1 & 2.25 & 4 \\
\end{bktable}
\end{bkcenter}
\end{yattemiyou}

\noindent 傾きは「たての変化量÷横の変化量」です（3-5）。②の傾きは $\dfrac{4-1}{2-1}=\dfrac{3}{1}=3$、③は $\dfrac{2.25-1}{1.5-1}=\dfrac{1.25}{0.5}=2.5$ です。QをPに近づけると、傾きは小さくなっていきます。④で引いた直線の傾きは、2に近くなったはずです。この章で、ちょうど2になることを確かめます。

\begin{omoidasou}
\bkrecall{5-2}{$x=5t^2$ の、1秒から $1+h$ 秒までの平均の速さは $10+5h$。$h$ を0に近づけると10。}
\bkrecall{3-3}{$f(x)$ は「関数 $f$ に $x$ を入れた値」。\\$f(x)=x^2$ なら $f(3)=3^2=9$、$f(a+h)=(a+h)^2$。}
\bkrecall{3-5}{傾きが $m$ で、点 $(x_1,\ y_1)$ を通る直線の式は $y-y_1=m(x-x_1)$。}
\bkrecall{4-2}{円の接線は、円に1点だけでふれる直線。接線は、ふれた点で半径に垂直。}
\end{omoidasou}

\noindent 3-3の $f(x)$ は、「$x$ を入れると、きまりにしたがって数が出てくる箱」のようなものです。$f(x)=x^2$ という箱に3を入れると9が、$a+h$ を入れると $(a+h)^2$ が出てきます。5-1・5-2で使った $x(1.1)$ も、同じ書き方でした。

\begin{bkfigure}
\begin{tikzpicture}[x=1mm, y=1mm]
  \node[bk node, minimum width=30mm, minimum height=12mm, font=\sffamily\small, fill=bkTheme!10] (f) at (60,0) {$f$\enspace（2乗する）};
  \node[font=\sffamily\small] (in1) at (14,5) {3};
  \node[font=\sffamily\small] (in2) at (14,-5) {$a+h$};
  \node[font=\sffamily\small] (o1) at (108,5) {$f(3)=9$};
  \node[font=\sffamily\small] (o2) at (108,-5) {$f(a+h)=(a+h)^2$};
  \draw[bk link] (in1) -- ([yshift=5mm]f.west);
  \draw[bk link] (in2) -- ([yshift=-5mm]f.west);
  \draw[bk link] ([yshift=5mm]f.east) -- (o1);
  \draw[bk link] ([yshift=-5mm]f.east) -- (o2);
\end{tikzpicture}
\bkfigcaption{関数 $f(x)=x^2$ は、入れた数を2乗して出す箱。}
\end{bkfigure}

%% ---------------- 数学の道具 ----------------
\bksection{doguu}{数学の道具}

\bksubsection{1}{かぎりなく近づく：極限と $\lim$}

\leavevmode\label{bk:5-3:lim}5-2では、「$h$ を0に近づけると、$10+5h$ は10に近づく」と考えました。表にすると、次のようになります。

\begin{bkcenter}
\begin{bktable}{colspec={Q[l,m,wd=22mm] *{5}{Q[c,m,wd=17mm]}}, rowsep=1.5pt}
$h$ & 0.1 & 0.01 & 0.001 & 0.0001 & …… \\
$5h$ & 0.5 & 0.05 & 0.005 & 0.0005 & …… \\
$10+5h$ & 10.5 & 10.05 & 10.005 & 10.0005 & → 10 \\
\end{bktable}
\bktabcaption{$h$ を0に近づけると、$10+5h$ は10に近づく}
\end{bkcenter}

$h$ は0にはなりません。0でない数のまま、0にいくらでも近くなります。このようすを「\textbf{$h$ を0にかぎりなく近づける}」といい、そのとき式の値が近づいていく先の10を、\textbf{極限}（きょくげん）といいます。記号で、次のように書きます。\bkidx{きょくげん}{極限}
\[ \lim_{h\to0}(10+5h)=10 \]

\begin{bkfigure}
\begin{tikzpicture}[x=1mm, y=1mm]
  % lim の書き方。式を3つの部分に分けて並べ、説明は各部分の真上・真下に置く（線は交差しない）
  \node[font=\Large, inner sep=1pt] (A) at (0,0) {$\displaystyle\lim_{h\to0}$};
  \node[font=\Large, anchor=base west, inner sep=1pt] (B) at ([xshift=2mm]A.base east) {$(10+5h)$};
  \node[font=\Large, anchor=base west, inner sep=1pt] (C) at ([xshift=2mm]B.base east) {$=\ 10$};
  \node[anchor=south east, font=\sffamily\footnotesize, text=bkBlue] (a1) at ([yshift=7mm, xshift=3mm]A.north) {リミット（かぎりなく近づけると）};
  \node[anchor=north, font=\sffamily\footnotesize, text=bkBlue] (a2) at ([yshift=-7mm]A.south) {$h$ を0に近づけるとき};
  \node[anchor=south west, font=\sffamily\footnotesize, text=bkGreen!70!black] (b1) at ([yshift=7mm, xshift=-3mm]B.north) {この式の値が近づく先は};
  \node[anchor=north, font=\sffamily\footnotesize, text=bkOrange!80!black] (c1) at ([yshift=-7mm, xshift=6mm]C.south) {10（極限）};
  \draw[bkBlue, line width=0.5pt] ([yshift=5.5mm]A.north) -- ([yshift=1.5mm]A.north);
  \draw[bkBlue, line width=0.5pt] ([yshift=1.5mm]a2.north) -- ([yshift=-1.5mm]A.south);
  \draw[bkGreen!70!black, line width=0.5pt] ([yshift=5.5mm]B.north) -- ([yshift=1.5mm]B.north);
  \draw[bkOrange!80!black, line width=0.5pt] ([yshift=1.5mm, xshift=-6mm]c1.north) -- ([yshift=-1.5mm]C.south);
\end{tikzpicture}
\bkfigcaption{$\lim$ の式の読み方}
\end{bkfigure}

「リミット、$h$ が0に近づくとき、$10+5h$ は10」と読みます。$\lim$ は、英語で「限り・限界」という意味の limit の頭の3文字です。$\lim$ の下の $h\to0$ は「$h$ を0に近づける」という意味で、矢印「→」は「近づく」と読みます。\bkidx{りみっと}{$\lim$（リミット）}

\noindent\textbf{「$=10$」と書いてよいの？}\par
$10+5h$ は、$h$ が0でない限り、ちょうど10にはなりません。でも、$\lim$ の式は「$10+5h$ の値」ではなく、「近づく先」を表しています。近づく先は、ちょうど10です。だから「$=10$」と書いてよいのです。

\begin{seikaku}
\textbf{「かぎりなく近づく」とは}\par
$h$ を0でない数のまま小さくしていくと、$10+5h$ と10の差 $5h$ は、0.001より小さく、0.000001より小さく……と、どんなに小さい数を決めても、それより小さくできます。このことを「$10+5h$ は10にかぎりなく近づく」といいます。大学の数学では、この言い方をもっと正確な式で書きます。
\end{seikaku}

\noindent\textbf{分数の形の極限}\par
$\displaystyle\lim_{h\to0}\frac{6h+h^2}{h}$ のように分母に $h$ があるときは、$h$ に0を入れると $\dfrac{0}{0}$ になって計算できません。そこで、$h$ を0に近づける前に、分子の1つずつを $h$ でわります（$h$ は0ではないので、わってよい。2-12）。\bkidx{ぶんぼ}{分母}
\[ \lim_{h\to0}\frac{6h+h^2}{h}=\lim_{h\to0}\Bigl(\frac{6h}{h}+\frac{h^2}{h}\Bigr)=\lim_{h\to0}(6+h)=6 \]
数で確かめると、$h=0.1$ なら $\dfrac{0.6+0.01}{0.1}=\dfrac{0.61}{0.1}=6.1$、$h=0.01$ なら $\dfrac{0.06+0.0001}{0.01}=6.01$ で、6に近づいています。

\bkkigocards{\bkkigocard{$\lim$}{リミット}{かぎりなく近づけたときの、近づく先の値（極限）}{limit}{5-3}}
{\bkkigocard{$\dfrac{\dd y}{\dd x}$}{ディーワイ・ディーエックス}{$y$ を $x$ で微分したもの（導関数）}{derivative}{5-3}}

\bksubsection{2}{曲線の接線}

\leavevmode\label{bk:5-3:sessen}円の接線は、円に1点だけでふれる直線でした（4-2）。曲がったグラフにも、1点でぴったりふれる直線を考えます。

グラフの上に点Pをとり、もう1つの点QをPの近くにとって、直線PQを引きます。Qをグラフにそって、Pにかぎりなく近づけていくと、直線PQはある1本の直線に近づいていきます。この直線を、点Pでの\textbf{接線}（せっせん）といい、Pを\textbf{接点}といいます。\bkidx{せっせん}{接線}

\begin{bkfigure}
\begin{tikzpicture}
\begin{axis}[bkplot, width=82mm, height=62mm, xmin=-0.3, xmax=2.4, ymin=-0.3, ymax=4.6, xtick={1,2}, ytick={1,2,3,4}, xlabel={$x$}, ylabel={$y$}]
  % 第6回：直線はx軸より上・枠の中で止め、名前は各直線の右の端に置く
  \addplot[bkcurve, domain=-0.3:2.1] {x^2};
  \addplot[bkthin, domain=0.72:2.2] {1+3*(x-1)};
  \addplot[bkthin, domain=0.64:2.4] {1+2.5*(x-1)};
  \addplot[bkthin, domain=0.6:2.35] {1+2.2*(x-1)};
  \addplot[bkcurve2, line width=1.4pt, domain=0.55:2.3] {1+2*(x-1)};
  \addplot[only marks, mark=*, mark size=1.8pt, bkInk] coordinates {(1,1) (2,4) (1.5,2.25) (1.2,1.44)};
  \node[anchor=north west, font=\sffamily\small] at (axis cs:1,1) {P};
  \node[anchor=south east, font=\sffamily\scriptsize, inner sep=1.5pt] at (axis cs:2,4) {Q};
  \node[anchor=south, font=\sffamily\scriptsize, text=bkGray, inner sep=1.5pt] at (axis cs:2.2,4.6) {傾き3};
  \node[anchor=west, font=\sffamily\scriptsize, text=bkGray, xshift=1mm] at (axis cs:2.4,4.5) {傾き2.5};
  \node[anchor=west, font=\sffamily\scriptsize, text=bkGray, xshift=1mm] at (axis cs:2.35,3.97) {傾き2.2};
  \node[anchor=west, font=\sffamily\scriptsize\bfseries, text=bkOrange!80!black, xshift=1mm] at (axis cs:2.3,3.6) {接線（傾き2）};
\end{axis}
\end{tikzpicture}
\bkfigcaption{$y=x^2$ の点P（1，1）。Qを（2，4）→（1.5，2.25）→（1.2，1.44）と近づけると、直線PQは接線（オレンジ）に近づく。}
\end{bkfigure}

\noindent\textbf{もう1つの見方：拡大すると、まっすぐに見える}\par
曲線を、点Pのまわりで虫めがねのように拡大していきます。下の3つの図は、どれも同じ大きさのわくに、Pのまわりをかいたものです。左から、横はば2、0.4、0.08の部分です。

\begin{bkfigure}
\begin{tikzpicture}
  % 3つのわくは同じ大きさ（38 mm 四方）。中心は P(1,1)。青い曲線 y=x^2、オレンジの点線は接線 y=2x-1
  \foreach \w/\xs/\lab in {1/0/横はば2, 0.2/46/横はば0.4, 0.04/92/横はば0.08} {
  \pgfmathsetmacro\xa{1-\w}\pgfmathsetmacro\xb{1+\w}\pgfmathsetmacro\ya{1-2*\w}\pgfmathsetmacro\yb{1+2*\w}
  \begin{axis}[at={(\xs mm,0)}, scale only axis, width=38mm, height=38mm, xmin=\xa, xmax=\xb, ymin=\ya, ymax=\yb,
    xtick=\empty, ytick=\empty, axis line style={bkGray, line width=0.6pt}, clip=true]
    \addplot[bkcurve, domain=\xa:\xb, samples=60] {x^2};
    \addplot[bkOrange, line width=0.9pt, dash pattern=on 3pt off 2pt, domain=\xa:\xb] {2*x-1};
    \addplot[only marks, mark=*, mark size=1.5pt, bkInk] coordinates {(1,1)};
  \end{axis}
  \node[font=\sffamily\scriptsize, anchor=north] at ({\xs mm+19mm},-2mm) {\lab};}
\end{tikzpicture}
\bkfigcaption{青が $y=x^2$、オレンジの点線が接線。拡大するほど、曲線は接線とほとんど重なる。}
\end{bkfigure}

拡大するほど、曲線はまっすぐに見え、接線とほとんど重なっていきます。つまり接線は、「その点のまわりをとても大きく拡大したとき、曲線が見せるまっすぐな線」です。だから、接線の傾きを、その点での「坂の急さ」と決めればよいのです。

\bksubsection{3}{接線の傾き：微分係数}

\leavevmode\label{bk:5-3:bibunkeisu}直線PQの傾きを、式で書きます。関数を $y=f(x)$ とし、Pの $x$ 座標を $a$、Qの $x$ 座標を $a+h$ とします。$a$ は「Pの場所」を表す決まった数で、$h$ は「PからQまでの横のはば」です。

\noindent\begin{minipage}[c]{0.5\linewidth}
\begin{tikzpicture}
\begin{axis}[bkplot, width=62mm, height=54mm, xmin=0, xmax=2.2, ymin=0, ymax=3.2, xtick={0.8,1.6}, xticklabels={$a$, $a+h$},
  ytick={0.64,2.56}, yticklabels={$f(a)$, $f(a+h)$}, xlabel={$x$}, ylabel={$y$}]
  \addplot[bkcurve, domain=0:1.75] {x^2};
  \draw[bkGray, line width=0.6pt, dash pattern=on 2pt off 1.5pt] (axis cs:0.8,0) -- (axis cs:0.8,0.64) -- (axis cs:0,0.64);
  \draw[bkGray, line width=0.6pt, dash pattern=on 2pt off 1.5pt] (axis cs:1.6,0) -- (axis cs:1.6,2.56) -- (axis cs:0,2.56);
  \addplot[bkcurve2, line width=1pt, domain=0.6:1.85] {0.64+2.4*(x-0.8)};
  \draw[bkInk, line width=0.7pt] (axis cs:0.8,0.64) -- (axis cs:1.6,0.64) -- (axis cs:1.6,2.56);
  \addplot[only marks, mark=*, mark size=1.7pt, bkInk] coordinates {(0.8,0.64) (1.6,2.56)};
  \node[anchor=south east, font=\sffamily\small] at (axis cs:0.8,0.64) {P};
  \node[anchor=east, font=\sffamily\small, xshift=-1.5mm] at (axis cs:1.6,2.56) {Q};
  \node[anchor=north, font=\sffamily\footnotesize] at (axis cs:1.2,0.6) {$h$};
\end{axis}
\end{tikzpicture}
\end{minipage}\hfill
\begin{minipage}[c]{0.47\linewidth}\raggedright
Pの高さは $f(a)$、Qの高さは $f(a+h)$。\par
横の変化量は $(a+h)-a=h$。\par
たての変化量は $f(a+h)-f(a)$（Qの右の、たての黒い線）。
\end{minipage}

\[ \text{PQの傾き}=\frac{\text{たての変化量}}{\text{横の変化量}}=\frac{f(a+h)-f(a)}{(a+h)-a}=\frac{f(a+h)-f(a)}{h} \]
これを\textbf{平均変化率}といいます。「$x$ が $a$ から $a+h$ まで変わるあいだに、$y$ が平均して $x$ の1あたりどれだけ変わったか」を表すからです。5-1の平均の速さは、$f$ が位置、$x$ が時間のときの平均変化率です。

Qを P に近づけることは、横のはば $h$ を0に近づけることです。だから、接線の傾きは
\[ f'(a)=\lim_{h\to0}\frac{f(a+h)-f(a)}{h} \]
です。これを、$x=a$ での\textbf{微分係数}（びぶんけいすう）といい、$f'(a)$（エフ・ダッシュ・エー）と書きます。$f$ の右上の「$'$」（ダッシュ）が「傾き」の印です。\bkidx{びぶんけいすう}{微分係数}\bkidx{へいきんへんかりつ}{平均変化率}

\noindent\textbf{やってみようの答え}\par
$f(x)=x^2$、$a=1$ とします。1行ずつ計算します。
\begin{align*}
f(1)&=1^2=1\\
f(1+h)&=(1+h)^2=1+2h+h^2 &&\text{（5-2 の正方形の図）}\\
f(1+h)-f(1)&=(1+2h+h^2)-1=2h+h^2\\
f'(1)&=\lim_{h\to0}\frac{2h+h^2}{h}=\lim_{h\to0}(2+h)=2
\end{align*}
定規で引いた直線の傾きは、ちょうど2です。図の傾き3（$h=1$）、2.5（$h=0.5$）、2.2（$h=0.2$）は、$2+h$ に $h$ を入れた値になっています。

\bksubsection{4}{どの点でも計算する：導関数}

\leavevmode\label{bk:5-3:doukansuu}$f(x)=x^2$ の微分係数を、いろいろな点で求めてみましょう。$a=1$ では2でした。同じ計算をすると、次のようになります（$a=3$ は例題2、$a=-1$ はじぶんでとこう）。

\begin{bkcenter}
\begin{bktable}{colspec={Q[l,m,wd=30mm] *{5}{Q[c,m,wd=13mm]}}, rowsep=1.5pt}
点の $x$ 座標 $a$ & $-1$ & 0 & 1 & 2 & 3 \\
微分係数 $f'(a)$ & $-2$ & 0 & 2 & 4 & 6 \\
\end{bktable}
\bktabcaption{$f(x)=x^2$ の、いろいろな点での接線の傾き}
\end{bkcenter}

どの点でも、$f'(a)$ は $a$ の2倍になっていそうです。1つずつ計算するかわりに、$a$ のところを $x$ のままにして計算すれば、どの点の接線の傾きも、一度に求められます。
\begin{align*}
f'(x)&=\lim_{h\to0}\frac{(x+h)^2-x^2}{h}=\lim_{h\to0}\frac{x^2+2xh+h^2-x^2}{h}\\
&=\lim_{h\to0}\frac{2xh+h^2}{h}=\lim_{h\to0}(2x+h)=2x
\end{align*}
表で見た「2倍」が、式で確かめられました。この $f'(x)$ を、$f(x)$ の\textbf{導関数}（どうかんすう）といいます。「$x$ を入れると、その点の接線の傾きが出てくる関数」です。導関数を求めることを、「$f(x)$ を\textbf{微分する}」といいます。$f(x)=x^2$ を微分すると $2x$ です。$x=1$ なら傾き2、$x=3$ なら傾き6、$x=-1$ なら傾き $-2$（右下がり）。\bkidx{どうかんすう}{導関数}\bkidx{びぶん}{微分}

\begin{bkfigure}
\begin{tikzpicture}
\begin{axis}[bkplot, name=top, width=62mm, height=46mm, xmin=-2.2, xmax=2.4, ymin=-0.5, ymax=4.6, xtick={-2,-1,1,2}, ytick={1,2,3,4}, xlabel={$x$}, ylabel={$y=x^2$}]
  \addplot[bkcurve, domain=-2.1:2.1] {x^2};
  % 第6回：接線は枠（ymax 4.6）の中で止める
  \foreach \a/\l/\r in {-1.5/-2/-1, 0/-0.5/0.5, 1/0.5/1.5, 2/1.6/2.15} {\edef\tmp{\noexpand\addplot[bkcurve2, line width=0.9pt, domain=\l:\r] {\a*\a+2*\a*(x-\a)};}\tmp}
\end{axis}
\begin{axis}[bkplot, at={($(top.east)+(12mm,0)$)}, anchor=west, width=62mm, height=46mm, xmin=-2.2, xmax=2.4, ymin=-4.4, ymax=4.6, xtick={-2,-1,1,2}, ytick={-4,-2,2,4}, xlabel={$x$}, ylabel={$y'=2x$}]
  \addplot[bkcurve2, domain=-2.1:2.1] {2*x};
  \addplot[only marks, mark=*, mark size=1.6pt, bkInk] coordinates {(-1.5,-3) (0,0) (1,2) (2,4)};
\end{axis}
\end{tikzpicture}
\bkfigcaption{左：$y=x^2$ と、いろいろな点の接線。右：傾きを点で打つと、直線 $y=2x$ にならぶ。}
\end{bkfigure}

導関数には、いくつかの書き方があります。

\begin{bkcenter}
\begin{bktable}{colspec={Q[c,m,wd=14mm] Q[l,m,wd=56mm] Q[l,m,wd=52mm]}, rowsep=1.5pt, cells={font=\small}, row{1}={font=\sffamily\bfseries\small}}
書き方 & 読み方 & よく使う場面 \\
$f'(x)$ & エフ・ダッシュ・エックス & 関数 $f$ を微分したとき \\
$y'$ & ワイ・ダッシュ & $y=\cdots$ の式を微分したとき \\
$\dfrac{\dd y}{\dd x}$ & ディーワイ・ディーエックス & $y$ を $x$ で微分したとき \\
$\dfrac{\dd x}{\dd t}$ & ディーエックス・ディーティー & 位置 $x$ を時間 $t$ で微分（速度） \\
\end{bktable}
\bktabcaption{導関数の書き方。どれも同じものを表す。}
\end{bkcenter}

$\dfrac{\dd y}{\dd x}$ は、平均変化率 $\dfrac{\Delta y}{\Delta x}$ の $\Delta$ を $\dd$ にかえた形です。「$\Delta x$ をかぎりなく小さくしたときの $\dfrac{\Delta y}{\Delta x}$」という意味を、形で表しています。分数のように見えますが、$\dd y$ を $\dd x$ でわった数ではなく、ひとまとまりの記号です。落ちるボール $x=5t^2$ なら、$\dfrac{\dd x}{\dd t}=10t$ が瞬間の速さです（5-2 の $v=10t$）。

\bksubsection{5}{接線の式}

\leavevmode\label{bk:5-3:shiki}接線は、点 $(a,\ f(a))$ を通り、傾きが $f'(a)$ の直線です。直線の式は、通る1点と傾きがわかれば作れます（3-5）。なぜそうなるのか、確かめておきます。

\noindent\begin{minipage}[c]{0.5\linewidth}
\begin{tikzpicture}
\begin{axis}[bkplot, width=62mm, height=52mm, xmin=0, xmax=4, ymin=-1.5, ymax=6.5, xtick={1,3}, xticklabels={$x_1$, $x$},
  ytick={1,5}, yticklabels={$y_1$, $y$}, xlabel={}, ylabel={}]
  \addplot[bkcurve2, line width=1.1pt, domain=0.3:3.6] {1+2*(x-1)};
  \draw[bkInk, line width=0.7pt] (axis cs:1,1) -- (axis cs:3,1) -- (axis cs:3,5);
  \addplot[only marks, mark=*, mark size=1.7pt, bkInk] coordinates {(1,1) (3,5)};
  \node[anchor=south east, font=\sffamily\footnotesize] at (axis cs:1,1) {A};
  \node[anchor=south east, font=\sffamily\footnotesize] at (axis cs:3,5) {B};
  \node[anchor=north, font=\sffamily\footnotesize] at (axis cs:2,0.9) {$x-x_1$};
  \node[anchor=west, font=\sffamily\footnotesize] at (axis cs:3.05,3) {$y-y_1$};
\end{axis}
\end{tikzpicture}
\end{minipage}\hfill
\begin{minipage}[c]{0.47\linewidth}\raggedright
傾き $m$ の直線が点A $(x_1,\ y_1)$ を通るとします。直線の上のほかの点を B $(x,\ y)$ とすると、AからBまでの傾きも $m$ です。
\[ \frac{y-y_1}{x-x_1}=m \]
両辺に $x-x_1$ をかけて（2-14）、
\[ y-y_1=m(x-x_1) \]
\end{minipage}

\medskip
数で確かめましょう。点 $(1,\ 1)$ を通り、傾き2の直線の上には、点 $(3,\ 5)$ があります（右へ2進むと、上へ $2\times2=4$ 上がる）。$\dfrac{5-1}{3-1}=\dfrac{4}{2}=2$ で、傾きは2です。

接線の場合は、$x_1$ が $a$、$y_1$ が $f(a)$、$m$ が $f'(a)$ なので、
\[ y-f(a)=f'(a)(x-a) \]
$y=x^2$ の点（1，1）での接線は、$f'(1)=2$ なので $y-1=2(x-1)$。右辺を展開して $y-1=2x-2$、両辺に1をたして $y=2x-1$ です。$x=1$ を入れると $y=1$ で、接点を通っています。

\begin{taisetsu}
微分係数（接線の傾き）\quad $f'(a)=\displaystyle\lim_{h\to0}\frac{f(a+h)-f(a)}{h}$\\
導関数 $f'(x)$ を求めることを「微分する」という。$(x^2)'=2x$\\
接線の式\quad $y-f(a)=f'(a)(x-a)$
\end{taisetsu}

\begin{kotoba}
\small
\bkword{極限}{かぎりなく近づけたときに、式の値が近づいていく先の値。}
\bkword{接線}{曲線の1点に、ぴったりふれる直線。その点のまわりを拡大すると、曲線と重なって見える。}
\bkword{微分係数}{1つの点での接線の傾き。$f'(a)$。}
\bkword{導関数}{$x$ を入れると、その点の接線の傾きが出てくる関数。$f'(x)$。}
\end{kotoba}

%% ---------------- 例題 ----------------
\bksection{reidai}{例題でたしかめよう}

\begin{reidai}[極限を求める]
$\displaystyle\lim_{h\to0}\frac{8h+2h^2}{h}$ を求めましょう。\par
\kangaekata $h=0$ を入れると $\dfrac{0}{0}$ になるので、$h$ を0に近づける前に、分子の1つずつを $h$ でわります。\par
\shiki $\dfrac{8h+2h^2}{h}=\dfrac{8h}{h}+\dfrac{2h^2}{h}=8+2h$\par
\hspace*{2.6\zw}$h$ を0に近づけると、$2h$ は0に近づくので、$8+2h$ は8に近づく。\par
\kotae 8
\end{reidai}

\begin{jibun}
$\displaystyle\lim_{h\to0}\frac{12h-3h^2}{h}$ を求めましょう。
\bkjans{$\lim(12-3h)=12$}
\end{jibun}

\begin{reidai}[微分係数を定義から求める]
$f(x)=x^2$ の、$x=3$ での微分係数 $f'(3)$ を求めましょう。\par
\kangaekata $f'(a)$ の式に $a=3$ を入れます。$f(3+h)$ は、$x$ に $3+h$ をかっこごと入れた値です。\par
\shiki $f(3)=3^2=9$\par
\hspace*{2.6\zw}$f(3+h)=(3+h)^2=9+6h+h^2$\par
\hspace*{2.6\zw}$f'(3)=\displaystyle\lim_{h\to0}\frac{(9+6h+h^2)-9}{h}=\lim_{h\to0}\frac{6h+h^2}{h}$\par
\hspace*{2.6\zw}$\phantom{f'(3)}=\displaystyle\lim_{h\to0}(6+h)=6$\par
\hspace*{2.6\zw}確かめ：$h=0.1$ なら $6+0.1=6.1$、$h=0.01$ なら $6.01$ で、6に近づく。\par
\kotae $f'(3)=6$
\end{reidai}

\begin{jibun}
$f(x)=x^2$ の、次の点での微分係数を、定義から求めましょう。\quad（1）$x=-1$\quad（2）$x=2$
\bkjans{（1）\nobreak $(-1+h)^2=1-2h+h^2$ より $\lim(-2+h)=-2$\quad（2）\nobreak $(2+h)^2=4+4h+h^2$ より $\lim(4+h)=4$}
\end{jibun}

\begin{reidai}[導関数を求める]
$f(x)=3x^2$ の導関数を、定義から求めましょう。\par
\kangaekata $f'(x)=\displaystyle\lim_{h\to0}\frac{f(x+h)-f(x)}{h}$ の分子を、先に計算します。\par
\shiki $f(x+h)=3(x+h)^2=3(x^2+2xh+h^2)=3x^2+6xh+3h^2$\par
\hspace*{2.6\zw}$f(x+h)-f(x)=3x^2+6xh+3h^2-3x^2=6xh+3h^2$\par
\hspace*{2.6\zw}$f'(x)=\displaystyle\lim_{h\to0}\frac{6xh+3h^2}{h}=\lim_{h\to0}(6x+3h)=6x$\par
\kotae $f'(x)=6x$
\end{reidai}

\begin{jibun}
$f(x)=x^2+x$ の導関数を、定義から求めましょう。
\bkjans{分子は $2xh+h^2+h$。$\lim(2x+h+1)=2x+1$}
\end{jibun}

\begin{reidai}[接線の式]
$y=x^2$ のグラフについて、（1）$x=2$ の点での接線の式を求めましょう。（2）傾きが6になる接線の、接点と式を求めましょう。\par
\kangaekata $(x^2)'=2x$ を使います。\par
\hspace*{4.3\zw}（1）接点は $(2,\ 2^2)=(2,\ 4)$、傾きは $f'(2)=2\times2=4$。\par
\hspace*{4.3\zw}（2）接点を $x=a$ とすると、傾きは $f'(a)=2a$。これが6になる $a$ を求める。\par
\shiki（1）$y-4=4(x-2)$\par
\hspace*{4.3\zw}$y-4=4x-8$\par
\hspace*{4.3\zw}$y=4x-8+4=4x-4$\par
\hspace*{2.6\zw}（2）$2a=6$ より $a=3$。接点は $(3,\ 3^2)=(3,\ 9)$\par
\hspace*{4.3\zw}$y-9=6(x-3)$ より $y=6x-18+9=6x-9$\par
\kotae （1）$y=4x-4$\quad（2）接点 $(3,\ 9)$、$y=6x-9$
\end{reidai}

\begin{jibun}
$y=x^2$ のグラフについて、（1）$x=-1$ の点での接線の式と、（2）傾きが2になる接線の接点と式を求めましょう。
\bkjans{（1）\nobreak 接点 $(-1,\ 1)$、傾き $-2$。$y=-2x-1$\quad（2）\nobreak $2a=2$ より接点 $(1,\ 1)$。$y=2x-1$}
\end{jibun}

\begin{reidai}[導関数から接線を引く]
$f(x)=x^2-4x$ の導関数を定義から求め、$x=1$ の点での接線の式を求めましょう。\par
\kangaekata $-4x$ の $x$ にも、かっこごと $x+h$ を入れます。\par
\shiki $f(x+h)-f(x)=(x+h)^2-4(x+h)-(x^2-4x)$\par
\hspace*{2.6\zw}$\phantom{f(x+h)-f(x)}=x^2+2xh+h^2-4x-4h-x^2+4x$\par
\hspace*{2.6\zw}$\phantom{f(x+h)-f(x)}=2xh+h^2-4h$\par
\hspace*{2.6\zw}$f'(x)=\displaystyle\lim_{h\to0}\frac{2xh+h^2-4h}{h}=\lim_{h\to0}(2x+h-4)=2x-4$\par
\hspace*{2.6\zw}接点 $(1,\ f(1))=(1,\ 1-4)=(1,\ -3)$、傾き $f'(1)=2\times1-4=-2$\par
\hspace*{2.6\zw}$y-(-3)=-2(x-1)$。$y+3=-2x+2$ より $y=-2x-1$\par
\kotae $f'(x)=2x-4$、$y=-2x-1$
\end{reidai}

\begin{jibun}
$f(x)=x^2+2x$ の導関数を定義から求め、次の点での接線の式を求めましょう。\quad（1）$x=1$\quad（2）$x=-2$
\bkjans{$f'(x)=2x+2$。（1）\nobreak 接点 $(1,\ 3)$、傾き4。$y=4x-1$\quad（2）\nobreak 接点 $(-2,\ 0)$、傾き $-2$。$y=-2x-4$}
\end{jibun}

\begin{reidai}[$x$-$t$ グラフの接線]
落ちるボールの位置は $x=5t^2$（m）で、$\dfrac{\dd x}{\dd t}=10t$ です。$x$-$t$ グラフの $t=1$ の点での接線の式を求めましょう。その傾きは何を表していますか。\par
\kangaekata 横軸が $t$、たて軸が $x$ なので、接線の式も $t$ と $x$ で書きます。「$y-f(a)=f'(a)(x-a)$」の $y$ が $x$ に、$x$ が $t$ にかわります。\par
\shiki 接点は $(1,\ 5\times1^2)=(1,\ 5)$、傾きは $10\times1=10$\par
\hspace*{2.6\zw}$x-5=10(t-1)$。$x-5=10t-10$ より $x=10t-5$\par
\noindent\begin{minipage}[c]{0.5\linewidth}
\begin{tikzpicture}
\begin{axis}[bkplot, width=56mm, height=40mm, xmin=0, xmax=2.2, ymin=-6, ymax=22, xtick={0,1,2}, ytick={-5,0,5,10,15,20},
  xlabel={$t$（s）}, ylabel={$x$（m）}]
  \addplot[bkcurve, domain=0:2] {5*x^2};
  \addplot[bkcurve2, line width=1pt, domain=0.1:2.2] {10*x-5};
  \addplot[only marks, mark=*, mark size=1.7pt, bkInk] coordinates {(1,5)};
\end{axis}
\end{tikzpicture}
\end{minipage}\hfill
\begin{minipage}[c]{0.46\linewidth}\raggedright
青が $x=5t^2$、オレンジが接線 $x=10t-5$。接線は点 $(1,\ 5)$ でグラフにふれています。$t=0$ を入れると $x=-5$ なので、たて軸と $-5$ で交わります。
\end{minipage}\par
\kotae $x=10t-5$。傾き10は、1秒の瞬間の速さ\qty{10}{m/s}を表す。
\end{reidai}

\begin{jibun}
$x=5t^2$（m）の $x$-$t$ グラフの、次の点での接線の式を求めましょう。\quad（1）$t=3$\quad（2）$t=0.5$
\bkjans{（1）\nobreak 接点 $(3,\ 45)$、傾き30。$x=30t-45$\quad（2）\nobreak 接点 $(0.5,\ 1.25)$、傾き5。$x=5t-1.25$}
\end{jibun}

%% ---------------- まちがえやすいところ・物理工学 ----------------
\bksection{tsumazuki}{まちがえやすいところをチェック}

\begin{tsumazuki}{$f(a+h)$ を $f(a)+h$ としてしまう}
$f(x)=x^2$ のとき、$f(3+h)$ は「$x$ に $3+h$ を入れた値」です。
\bkng{$f(3+h)=f(3)+h=9+h$}\bkok{$f(3+h)=(3+h)^2=9+6h+h^2$}
かっこごと入れるのが代入です（2-13）。$h=1$ で確かめると、$f(4)=16$ ですが、$9+1=10$ で合いません。
\end{tsumazuki}

\begin{tsumazuki}{接線の式で、傾きと接点を取りちがえる}
$y=x^2$ の $x=3$ の点の接線は、接点 $(3,\ 9)$、傾き $f'(3)=6$ です。
\bkng{$y-6=9(x-3)$}\bkok{$y-9=6(x-3)$}
かっこの前の数が傾き、左の $y$ からひく数が接点の $y$ 座標です。
\end{tsumazuki}

\begin{tsumazuki}{接線は、グラフと1点だけで交わる？}
\vspace{7pt}% 1行目が横に図のある段なので、題とのあいだをほかの箱とそろえる
\noindent\begin{minipage}[c]{20\zw}
円の接線は1点だけで円にふれます。でも、曲線の接線は、接点からはなれた所で、もう一度グラフと交わることがあります（右図）。接線かどうかは、「接点の近くで、グラフにぴったりふれているか」で決まります。
\end{minipage}\hfill
\begin{minipage}[c]{0.36\linewidth}
\begin{tikzpicture}
\begin{axis}[width=44mm, height=30mm, axis lines=none, xmin=0, xmax=6.5, ymin=-1.4, ymax=1.5, clip=false]
  % 第6回：x=4 の接線（枠の中で止める）。x=1.35 で、もう一度グラフと交わる
  \addplot[bkcurve, domain=0:6.3, samples=80, line width=1.2pt] {sin(deg(x))};
  \addplot[bkcurve2, domain=0.8:4.9, line width=1pt] {-0.7568-0.6536*(x-4)};
  \addplot[only marks, mark=*, mark size=1.4pt, bkInk] coordinates {(4,-0.7568)};
  \addplot[only marks, mark=o, mark size=1.6pt, bkInk, line width=0.8pt] coordinates {(1.35,0.976)};
  \node[anchor=south west, font=\sffamily\scriptsize, inner sep=1pt] at (axis cs:4.1,-0.48) {接点};
  \node[anchor=south west, font=\sffamily\scriptsize, inner sep=1pt] at (axis cs:1.6,1.3) {交わる};
\end{axis}
\end{tikzpicture}
\end{minipage}
\end{tsumazuki}

\begin{butsuri}
\begin{itemize}[leftmargin=1.2\zw, itemsep=2pt, topsep=0pt, label={\color{bkPurple}●}]
\item \textbf{ハンマー投げ}\par
回っているハンマーから手をはなすと、ハンマーはその瞬間の円の接線の向きに飛んでいきます。
\item \textbf{カーブの設計}\par
道路や線路は、まっすぐな所から曲がった所へ、向きが急に変わらないようにつないであります。つなぎめで、まっすぐな部分と曲がった部分の接線の向きが同じになるように設計するのです。
\item \textbf{速度計}\par
$x$-$t$ グラフの接線の傾きが、その瞬間の速さです。スピードメーターは、接線の傾きを（とても短い時間の平均で）計算する機械だといえます。
\end{itemize}
\end{butsuri}

\begin{motto}[3]
\textbf{とがった所には、接線が引けない}\par
$y=|x|$（$x$ の大きさ。$|{-3}|=3$）のグラフは、原点でとがっています。原点の右では傾き1、左では傾き $-1$ なので、右から近づけた傾きと左から近づけた傾きが一致しません。拡大しても、とがったままで、まっすぐには見えません。このような点では接線が決まらず、「微分できない」といいます。
\end{motto}

\begin{column}{ニュートンとライプニッツ}
微分の考え方は、350年ほど前に、イギリスのニュートンとドイツのライプニッツが、それぞれ別々に作りあげました。ニュートンは、惑星や落ちる物の運動を調べるために、時間で微分したものを $\dot{x}$（エックス・ドット）と、文字の上に点を打って書きました。ライプニッツは、$\dfrac{\dd y}{\dd x}$ という書き方を考えました。いまでも、物理では時間の微分に $\dot{x}$ を、数学や工学では $\dfrac{\dd y}{\dd x}$ をよく使います（5-5）。\bkidx{にゅーとん}{ニュートン}\bkidx{らいぷにっつ}{ライプニッツ}
\end{column}

%% ---------------- 章末問題 ----------------
\bkendproblems
\bkproblemgroup{基本}{3}

\bkproblem 次の極限を求めましょう。\quad（1）$\displaystyle\lim_{h\to0}\frac{5h+h^2}{h}$\quad（2）$\displaystyle\lim_{h\to0}\frac{4h-2h^2}{h}$
\bkpans{（1）\nobreak 5\quad（2）\nobreak 4}

\bkproblem $f(x)=x^2$ の、次の点での微分係数を定義から求めましょう。\\（1）$x=4$\quad（2）$x=-3$
\bkpans{（1）\nobreak $\lim(8+h)=8$\quad（2）\nobreak $\lim(-6+h)=-6$}

\bkproblem 次の関数の導関数を、定義から求めましょう。\\（1）$f(x)=2x^2$\quad（2）$f(x)=x^2-3x$
\bkpans{（1）\nobreak $4x$\quad（2）\nobreak $2x-3$}

\bkproblem $y=x^2$ のグラフの、$x=3$ の点での接線の式を求めましょう。
\bkpans{$y=6x-9$}

\bkproblem $f(x)=x^2+4x$ の導関数を定義から求め、$x=-1$ の点での接線の式を求めましょう。
\bkpans{$f'(x)=2x+4$。接点 $(-1,\ -3)$、傾き2。$y=2x-1$}

\bkproblem 次の極限を求めましょう。\quad（1）$\displaystyle\lim_{h\to0}\frac{3h+h^2}{h}$\quad（2）$\displaystyle\lim_{h\to0}\frac{-7h+4h^2}{h}$
\bkpans{（1）\nobreak 3\quad（2）\nobreak $-7$}

\bkproblemgroup{チャレンジ}{3}

\bkproblem 落ちるボールの位置は $x=5t^2$（m）で、$\dfrac{\dd x}{\dd t}=10t$ です。$x$-$t$ グラフの $t=2$ の点での接線の式を求めましょう。また、その傾きは何を表していますか。
\bkpans{$x=20t-20$。傾き20は、2秒の瞬間の速さ\qty{20}{m/s}}

\bkproblem $y=x^2$ のグラフの接線のうち、傾きが $-4$ になるものがあります。接点の座標と、接線の式を求めましょう。
\bkpans{$2a=-4$ より $a=-2$。接点 $(-2,\ 4)$、$y=-4x-4$}

\bkproblem $f(x)=x^2-2x$ の導関数を定義から求め、$x=3$ の点での接線の式を求めましょう。
\bkpans{$f'(x)=2x-2$。接点 $(3,\ 3)$、傾き4。$y=4x-9$}

\bkproblem 位置が $x=t^2-2t$（m）の物体があります。$x$-$t$ グラフの $t=3$ の点での接線の式を求め、その傾きが何を表すか答えましょう。
\bkpans{$\dfrac{\dd x}{\dd t}=2t-2$。接点 $(3,\ 3)$、傾き4。$x=4t-9$。傾きは3秒の瞬間の速さ\qty{4}{m/s}}

\bkproblem $f(x)=2x^2$ の導関数を定義から求め、$x=1$ の点での接線の式を求めましょう。
\bkpans{$f'(x)=4x$。接点 $(1,\ 2)$、傾き4。$y=4x-2$}

\begin{furikaeri}
\begin{checklist}
\item $\lim$ の記号を読み、極限の値を求められる。\hfill{\footnotesize\sffamily\color{bkTheme}→\pageref{bk:5-3:lim}ページ}
\item 微分係数を、接線の傾きとして求められる。\hfill{\footnotesize\sffamily\color{bkTheme}→\pageref{bk:5-3:bibunkeisu}ページ}
\item 導関数を求め、接線の式を書ける。\hfill{\footnotesize\sffamily\color{bkTheme}→\pageref{bk:5-3:shiki}ページ}
\end{checklist}
\end{furikaeri}

\bkfukidashi{delta}{$x^2$ を微分すると $2x$。では、$x^3$ や $x^4$ はどうなるかな？　次の章で、表をながめてきまりを見つけよう。}

%% 5-4 表から導関数を見つけよう（第11回：ページ制限なし。説明・図・例題を書き足した）
% 第11回（X13）：立方体の図を描き直す。斜めから見た箱 \bkvebox{x0}{y0}{z0}{dx}{dy}{dz}{色}
%   （x は右、y は奥、z は上。見える3つの面＝前・上・右をぬる。奥の物から順に描く）
\newcommand\bkvebox[7]{%
  \pgfmathsetmacro\bkveX{#1+#4}\pgfmathsetmacro\bkveY{#2+#5}\pgfmathsetmacro\bkveZ{#3+#6}%
  \filldraw[fill=#7!42, draw=#7!60!black, line width=0.5pt, line join=round] (#1,#2,#3) -- (\bkveX,#2,#3) -- (\bkveX,#2,\bkveZ) -- (#1,#2,\bkveZ) -- cycle;
  \filldraw[fill=#7!18, draw=#7!60!black, line width=0.5pt, line join=round] (#1,#2,\bkveZ) -- (\bkveX,#2,\bkveZ) -- (\bkveX,\bkveY,\bkveZ) -- (#1,\bkveY,\bkveZ) -- cycle;
  \filldraw[fill=#7!68, draw=#7!60!black, line width=0.5pt, line join=round] (\bkveX,#2,#3) -- (\bkveX,\bkveY,#3) -- (\bkveX,\bkveY,\bkveZ) -- (\bkveX,#2,\bkveZ) -- cycle;}
% 1辺 a の立方体と、それを h だけ大きくしたときにふえる7つの部分 \bkvecube{a}{h}{右のすき間}{前のすき間}{上のすき間}（すき間0なら組み立てた形）
\newcommand\bkvecube[5]{%
  \bkvebox{0}{0}{0}{#1}{#1}{#1}{bkBlue}%
  \bkvebox{#1+#3}{0}{0}{#2}{#1}{#1}{bkOrange}%
  \bkvebox{0}{0}{#1+#5}{#1}{#1}{#2}{bkOrange}%
  \bkvebox{0}{-#4-#2}{0}{#1}{#2}{#1}{bkOrange}%
  \bkvebox{#1+#3}{-#4-#2}{0}{#2}{#2}{#1}{bkGreen}%
  \bkvebox{0}{-#4-#2}{#1+#5}{#1}{#2}{#2}{bkGreen}%
  \bkvebox{#1+#3}{0}{#1+#5}{#2}{#1}{#2}{bkGreen}%
  \bkvebox{#1+#3}{-#4-#2}{#1+#5}{#2}{#2}{#2}{bkPurple}}
\bkchapter{5}{4}{表から導関数を見つけよう}{\bklevel{3}}{表から導関数を見つけよう}
\bkmichishirube{2-7,3-8,5-3}{5-5,5-6,5-11}

\begin{goal}
\begin{itemize}[label={\color{bkTheme}●}, leftmargin=1.6\zw, itemsep=0pt, topsep=1pt]
\item \bklevel{3}\enspace 表計算の結果から、微分のきまりを見つけられる。
\item \bklevel{3}\enspace $(x^n)'=nx^{n-1}$ を使って微分できる。
\item \bklevel{3}\enspace 定数倍と和のきまりを使って、多項式を微分できる。
\end{itemize}
\end{goal}

\begin{hatena}
\noindent\begin{minipage}[c]{0.6\linewidth}
5-3で、$x^2$ を微分すると $2x$ になることを、極限の計算で確かめました。では、$x^3$ や $x^4$ は、微分するとどうなるでしょう。

1つずつ $\lim$ を計算するのは大変です。何かきまりがあれば、すぐに微分できそうです。
\end{minipage}\hfill
\begin{minipage}[c]{0.36\linewidth}
\begin{tikzpicture}[x=1mm, y=1mm]
  \node[draw=bkTheme, line width=1pt, rounded corners=2mm, minimum width=38mm, minimum height=8mm, font=\sffamily\small] at (0,12) {$x^2$ \ →\ $2x$};
  \node[draw=bkTheme, line width=1pt, rounded corners=2mm, minimum width=38mm, minimum height=8mm, font=\sffamily\small] at (0,2) {$x^3$ \ →\ ？};
  \node[draw=bkTheme, line width=1pt, rounded corners=2mm, minimum width=38mm, minimum height=8mm, font=\sffamily\small] at (0,-8) {$x^4$ \ →\ ？};
\end{tikzpicture}
\end{minipage}
\end{hatena}

この章では、表計算で数をたくさん出して、そのならび方からきまりを見つけます。見つけたきまりは、式で確かめます。さらに、立方体の図でも確かめます。「数で見つけて、式と図で確かめる」のは、科学者がいつもしていることです。

\clearpage
%% ---------------- やってみよう・思い出そう ----------------
\begin{yattemiyou}[表計算で接線の傾きをならべよう]
5-3で、接線の傾き（微分係数）は、平均変化率 $\dfrac{f(x+h)-f(x)}{h}$ の $h$ を0に近づけた値でした。表計算では $h$ を0にはできないので、$h=0.001$ のような小さな数にします。$h$ が小さければ、平均変化率は接線の傾きにとても近くなります（5-2の表で、$\Delta t=0.001$ のとき、平均の速さは10.005で、瞬間の速さ10にとても近かったね）。

\begin{enumerate}[label={\bkcircled[bkGreen]{\arabic*}}, leftmargin=2.4\zw, itemsep=1pt, topsep=3pt]
\item A列に $x$ の値（0、1、2、3、4）を入れる。
\item B2に \texttt{=((A2+0.001)\^{}2-A2\^{}2)/0.001} を入れる（式の読み方は下の表）。C2には、2つの \texttt{\^{}2} を \texttt{\^{}3} にかえた式、D2には \texttt{\^{}4} にかえた式を入れる。
\item B2〜D2の式を、6行めまで下へコピーする（2-18）。
\end{enumerate}
\end{yattemiyou}

B2の式は、次のように読みます。
\begin{bkcenter}
\begin{bktable}{colspec={Q[l,m,wd=40mm] Q[l,m,wd=78mm]}, rowsep=1.5pt, cells={font=\small}, row{1}={font=\sffamily\bfseries\small}}
式の部分 & 意味 \\
\texttt{(A2+0.001)\^{}2} & $f(x+h)$。$x$ に $h=0.001$ をたしてから2乗する \\
\texttt{A2\^{}2} & $f(x)$。$x$ をそのまま2乗する \\
\texttt{( … - … )/0.001} & ひき算した差を $h=0.001$ でわる（平均変化率） \\
\end{bktable}
\end{bkcenter}
計算の結果は、次のようになります。
\begin{bkcenter}
\bksheet{4}{22}{%
1 & $x$ & $x^2$ の傾き & $x^3$ の傾き & $x^4$ の傾き \\
2 & 0 & 0.001 & 0.000001 & 0.000000 \\
3 & 1 & 2.001 & 3.003001 & 4.006004 \\
4 & 2 & 4.001 & 12.006001 & 32.024008 \\
5 & 3 & 6.001 & 27.009001 & 108.054012 \\
6 & 4 & 8.001 & 48.012001 & 256.096016 \\}
\bktabcaption{計算の結果。小数のはしたを取ると、何に見えるかな？}
\end{bkcenter}
たとえば、C4の 12.006001 は、「$x^3$ のグラフの、$x=2$ の所の接線の傾きは、およそ12」という意味です。

\begin{omoidasou}
\bkrecall{5-3}{導関数 $f'(x)=\displaystyle\lim_{h\to0}\frac{f(x+h)-f(x)}{h}$。$(x^2)'=2x$。}
\bkrecall{5-2}{$x=5t^2$ の瞬間の速さは $v=10t$。}
\bkrecall{3-8}{$(x+h)^2=x^2+2xh+h^2$。かっこをかけるときは、全部の組み合わせをかけてたす。}
\bkrecall{2-7}{$x^0=1$、$x^{-1}=\dfrac{1}{x}$、$x^{-2}=\dfrac{1}{x^2}$。$x^m\times x^n=x^{m+n}$。}
\bkrecall{2-12}{まとめてわるときは、1つずつわってよい：$(A+B)\div h=A\div h+B\div h$}
\bkrecall{2-18}{表計算では、式を下へコピーすると、番地がずれて同じ計算をくり返せる。}
\end{omoidasou}

%% ---------------- 数学の道具 ----------------
\bksection{doguu}{数学の道具}

\bksubsection{1}{小数のはしたは、どこから来たの？}

\leavevmode\label{bk:5-4:hyou}やってみようの表には、2.001 や 12.006001 のように、小数のはしたがついています。このはしたは、$h$ を0にできず、0.001 にしたために出たものです。

たしかめてみましょう。$x^3$ の、$x=2$ の傾きを、$h$ をかえて計算します。

\begin{bkcenter}
\begin{bktable}{colspec={Q[c,m,wd=26mm] Q[c,m,wd=34mm] Q[c,m,wd=30mm]}, rowsep=1.5pt}
$h$ & 平均変化率 & はした \\
0.01 & 12.0601 & 0.0601 \\
0.001 & 12.006001 & 0.006001 \\
0.0001 & 12.00060001 & 0.00060001 \\
\end{bktable}
\bktabcaption{$h$ を10分の1にすると、はしたもおよそ10分の1になる。}
\end{bkcenter}

$h$ を小さくするほど、はしたは0に近づき、平均変化率は12に近づきます。だから、ほんとうの接線の傾き（$h$ を0に近づけた極限）は12だと考えられます。同じように、表の数のはしたを取ると、次のようになります。

\begin{bkcenter}
\begin{bktable}{colspec={Q[c,m,wd=22mm] *{5}{Q[c,m,wd=12mm]}}, rowsep=2pt}
$x$ & 0 & 1 & 2 & 3 & 4 \\
$x^2$ の傾き & 0 & 2 & 4 & 6 & 8 \\
$x^3$ の傾き & 0 & 3 & 12 & 27 & 48 \\
$x^4$ の傾き & 0 & 4 & 32 & 108 & 256 \\
\end{bktable}
\bktabcaption{小数のはしたを取った接線の傾き}
\end{bkcenter}

\bksubsection{2}{ならび方から、きまりを見つける}

\leavevmode\label{bk:5-4:kimari0}この表から、きまりを探しましょう。いきなり全部を見ずに、1か所ずつ見ていきます。

\noindent\textbf{①\enspace $x=1$ の列を見る}\par
\noindent $x=1$ の列は、上から 2、3、4 です。これは、$x^2$、$x^3$、$x^4$ の右肩の数（\textbf{指数}）とまったく同じです。1は何回かけても1なので、$x=1$ では、$x$ の部分がじゃまをせず、前にかかる数だけが見えているのかもしれません。

\noindent\textbf{②\enspace $x^2$ の行を見る}\par
\noindent 0、2、4、6、8 は、$x$ の値 0、1、2、3、4 の2倍です。つまり $2\times x$。5-3で求めた $(x^2)'=2x$ と合っています。

\noindent\textbf{③\enspace $x^3$ の行を見る}\par
\noindent ①から、前にかかる数は3かもしれません。そこで、$x^3$ の行を3でわってみます。

\begin{bkcenter}
\begin{bktable}{colspec={Q[c,m,wd=26mm] *{5}{Q[c,m,wd=15mm]}}, rowsep=1.5pt}
$x$ & 0 & 1 & 2 & 3 & 4 \\
$x^3$ の傾き & 0 & 3 & 12 & 27 & 48 \\
3でわると & 0 & 1 & 4 & 9 & 16 \\
$x^2$ の値 & $0^2=0$ & $1^2=1$ & $2^2=4$ & $3^2=9$ & $4^2=16$ \\
\end{bktable}
\bktabcaption{3でわった行は、$x^2$ の値とぴったり同じ。}
\end{bkcenter}

3でわると $x^2$ になったので、$x^3$ の傾きは $3\times x^2$ です。

\noindent\textbf{④\enspace $x^4$ の行を見る}\par
\noindent 同じように4でわると、0、1、8、27、64。これは $0^3$、$1^3$、$2^3$、$3^3$、$4^3$ です。だから $x^4$ の傾きは $4\times x^3$ です。

まとめると、次のようになっていそうです。
\[ (x^2)'=2x^1,\qquad (x^3)'=3x^2,\qquad (x^4)'=4x^3 \]

\begin{bkfigure}
\begin{tikzpicture}[x=1mm, y=1mm]
  \node[font=\Large] (a) at (0,0) {$x^{\textcolor{bkOrange}{\mathbf{3}}}$};
  \node[font=\Large, anchor=west] (b) at (40,0) {$\textcolor{bkOrange}{\mathbf{3}}\,x^{\textcolor{bkBlue}{3-1}}$};
  \node[font=\Large, anchor=west] (c) at (b.east) {$=3x^2$};
  \draw[bk link] (a) -- node[bk link label, above] {微分する} (b);
  \draw[bkOrange!80!black, line width=0.5pt] ([xshift=2.2mm]b.south west) -- ++(0,-3.5) -- ++(-6,0);
  \node[font=\sffamily\footnotesize, text=bkOrange!80!black, anchor=east] at ([xshift=-3.8mm, yshift=-3.5mm]b.south west) {① 指数を前に出す};
  \draw[bkBlue!80!black, line width=0.5pt] ([xshift=-1.5mm, yshift=-1mm]b.north east) -- ++(0,3) -- ++(6,0);
  \node[font=\sffamily\footnotesize, text=bkBlue!80!black, anchor=west] at ([xshift=4.5mm, yshift=2mm]b.north east) {② 指数を1へらす};
\end{tikzpicture}
\bkfigcaption{見えてきたきまり。「前に出す」と「1へらす」の2つを、いつもセットで行う。}
\end{bkfigure}

まだ「表で見えた」だけです。たまたま0〜4で合っただけかもしれません。次に、式で確かめましょう。

\bksubsection{3}{$x^3$ のきまりを、式で確かめる}

\leavevmode\label{bk:5-4:x3}5-3で $x^2$ を微分したのと同じ手順で、$x^3$ の導関数を極限で求めます。そのために、まず $(x+h)^3$ を展開します。$(x+h)^3=(x+h)^2\times(x+h)$ と考え、$(x+h)^2=x^2+2xh+h^2$（3-8）を使います。

左のかっこの3つの項（$x^2$、$2xh$、$h^2$）の1つずつに、右のかっこの2つの項（$x$、$h$）をかけます。全部で $3\times2=6$ 個の組み合わせができます。

\begin{bkcenter}
\begin{bktable}{colspec={Q[c,m,wd=26mm] Q[c,m,wd=26mm] Q[c,m,wd=26mm]}, rowsep=2pt, cell{1}{1}={bg=white}}
かける数 & $\times\,x$ & $\times\,h$ \\
$x^2$ & $x^3$ & $x^2h$ \\
$2xh$ & $2x^2h$ & $2xh^2$ \\
$h^2$ & $xh^2$ & $h^3$ \\
\end{bktable}
\bktabcaption{6つの組み合わせ。同じ形の項（$x^2h$ どうし、$xh^2$ どうし）をまとめる。}
\end{bkcenter}

\begin{align*}
(x+h)^3&=x^3+x^2h+2x^2h+2xh^2+xh^2+h^3\\
&=x^3+3x^2h+3xh^2+h^3
\end{align*}
$x^2h$ が $1+2=3$ 個、$xh^2$ が $2+1=3$ 個あるので、どちらも3がつきます。

これで、平均変化率が計算できます。$x^3$ を引くと、$x^3$ の項が消えます。
\[ \frac{(x+h)^3-x^3}{h}=\frac{3x^2h+3xh^2+h^3}{h} \]
分子の3つの項を、1つずつ $h$ でわります（2-12。$h$ は0ではないので、わってよい）。
\[ \frac{3x^2h}{h}+\frac{3xh^2}{h}+\frac{h^3}{h}=3x^2+3xh+h^2 \]
最後に、$h$ を0に近づけます。$3xh$ と $h^2$ は、$h$ がかかっているので0に近づきます。たとえば $x=2$ なら、

\begin{bkcenter}
\begin{bktable}{colspec={Q[c,m,wd=16mm] Q[c,m,wd=16mm] Q[c,m,wd=18mm] Q[c,m,wd=22mm] Q[c,m,wd=26mm]}, rowsep=1.5pt}
$h$ & $3x^2$ & $3xh$ & $h^2$ & 合計 \\
0.1 & 12 & 0.6 & 0.01 & 12.61 \\
0.01 & 12 & 0.06 & 0.0001 & 12.0601 \\
0.001 & 12 & 0.006 & 0.000001 & 12.006001 \\
\end{bktable}
\bktabcaption{$x=2$ のとき。$3x^2=12$ は $h$ によらず同じで、残りの2つだけが0に近づく。}
\end{bkcenter}

合計の列は、1つ前の表の「平均変化率」と同じ数です。表計算のはしたの正体は、$3xh+h^2$ だったのです。$h$ を0に近づけると、残るのは $3x^2$ だけです。
\[ (x^3)'=\lim_{h\to0}(3x^2+3xh+h^2)=3x^2 \]
表で見つけたきまりのとおりです。

\bksubsection{4}{立方体の図で確かめる}

\leavevmode\label{bk:5-4:cube}$x^3$ は、1辺 $x$ の立方体の体積です。式の計算を、図でも確かめてみましょう。

1辺 $x$ の立方体の、たて・横・高さを、どれも $h$ だけのばします。すると、1辺 $x+h$ の立方体になります（下の図）。ふえた体積は、大きな立方体の体積から、もとの立方体の体積を引いた $(x+h)^3-x^3$ です。

\begin{bkfigure}
\begin{tikzpicture}[x={(1mm,0mm)}, y={(0.48mm,0.32mm)}, z={(0mm,1mm)}]
  % 左：もとの立方体（1辺 x）
  \bkvebox{0}{0}{0}{18}{18}{18}{bkBlue}
  \node[font=\sffamily\small, anchor=north] at (9,0,-1.6) {$x$};
  \node[font=\sffamily\small, anchor=east] at (-1.6,0,9) {$x$};
  \draw[bk arrow, bkGray] (30,9,9) -- (48,9,9);
  \node[font=\sffamily\scriptsize, text=bkGray, align=center, anchor=south] at (39,9,10.5) {たて・横・高さを\\$h$ ずつのばす};
  % 右：1辺 x+h の立方体（ふえた部分を色分け）。前の面は y=-6
  \begin{scope}[shift={(72,0,0)}]
    \bkvecube{18}{6}{0}{0}{0}
    \draw[bkInk, line width=0.5pt] (0,-6,-1) -- (0,-6,-2.8) (18,-6,-1) -- (18,-6,-2.8) (24,-6,-1) -- (24,-6,-2.8);
    \node[font=\sffamily\small, anchor=north] at (9,-6,-2.2) {$x$};
    \node[font=\sffamily\small, anchor=north] at (21,-6,-2.2) {$h$};
    \draw[bkInk, line width=0.5pt] (-1,-6,0) -- (-2.8,-6,0) (-1,-6,18) -- (-2.8,-6,18) (-1,-6,24) -- (-2.8,-6,24);
    \node[font=\sffamily\small, anchor=east] at (-2.6,-6,9) {$x$};
    \node[font=\sffamily\small, anchor=east] at (-2.6,-6,21) {$h$};
  \end{scope}
\end{tikzpicture}
\bkfigcaption{左：1辺 $x$ の立方体（青）。右：1辺を $x+h$ にした立方体。ふえた部分を色で分けた。}
\end{bkfigure}

右の立方体では、青い立方体はかくれて見えません。ふえた部分がまわりをおおっているからです。そこで、ふえた部分を、少しずつ引きはなしてみます。

\begin{bkfigure}
\begin{tikzpicture}[x={(1mm,0mm)}, y={(0.48mm,0.32mm)}, z={(0mm,1mm)}]
  \bkvecube{18}{6}{10}{22}{10}
  \begin{scope}[x=1mm, y=1mm, lg/.style={font=\sffamily\footnotesize, anchor=west}]
    \fill[bkBlue!42] (60,30) rectangle (64,34);   \node[lg] at (65.5,32) {もとの立方体（1こ）};
    \fill[bkOrange!42] (60,22) rectangle (64,26); \node[lg] at (65.5,24) {うすい板（3まい）};
    \fill[bkGreen!42] (60,14) rectangle (64,18);  \node[lg] at (65.5,16) {細い角柱（3本）};
    \fill[bkPurple!42] (60,6) rectangle (64,10);  \node[lg] at (65.5,8) {小さな立方体（1こ）};
  \end{scope}
\end{tikzpicture}
\bkfigcaption{1辺 $x+h$ の立方体を、色ごとに引きはなした形。青い立方体のまわりに、3種類の部分がついていた。}
\end{bkfigure}

青い立方体（もとの立方体）のまわりについていた部分は、3種類あります。同じ種類を集めて、ならべてみましょう。

\begin{bkfigure}
\begin{tikzpicture}[x={(1mm,0mm)}, y={(0.48mm,0.32mm)}, z={(0mm,1mm)}, lb/.style={font=\sffamily\footnotesize, anchor=north, align=center}]
  % うすい板 3まい（向きはそろえて、平らに置く）
  \foreach \k in {2,1,0} {\bkvebox{0}{0}{\k*12}{16}{16}{5}{bkOrange}}
  % 細い角柱 3本
  \foreach \k in {2,1,0} {\bkvebox{44}{0}{\k*10}{16}{5}{5}{bkGreen}}
  % 小さな立方体 1こ
  \bkvebox{86}{0}{0}{5}{5}{5}{bkPurple}
  \node[lb] at (12,0,-2.5) {\textbf{うすい板 3まい}\\1まいは $x\times x\times h=x^2h$\\合計 $\textcolor{bkOrange!80!black}{3x^2h}$};
  \node[lb] at (53,0,-2.5) {\textbf{細い角柱 3本}\\1本は $x\times h\times h=xh^2$\\合計 $\textcolor{bkGreen!60!black}{3xh^2}$};
  \node[lb] at (90,0,-2.5) {\textbf{小さな立方体 1こ}\\$h\times h\times h=h^3$\\\phantom{合計}};
\end{tikzpicture}
\bkfigcaption{ふえた部分を種類ごとにならべた。板3まいは、向きはちがっても同じ形（$x\times x\times h$）。}
\end{bkfigure}

ふえた体積を全部たすと、
\[ (x+h)^3-x^3=3x^2h+3xh^2+h^3 \]
これは、前の小見出しで展開して求めた式と、ぴったり同じです。展開の「3」は、板が3まい、角柱が3本あることだったのです。

\noindent\textbf{$h$ を小さくすると、どうなる？}\par
\noindent $x=1$ として、$h=0.1$ と $h=0.01$ で、それぞれの体積を計算します。

\begin{bkcenter}
\begin{bktable}{colspec={Q[c,m,wd=14mm] Q[c,m,wd=24mm] Q[c,m,wd=24mm] Q[c,m,wd=24mm] Q[c,m,wd=26mm]}, rowsep=1.5pt, cells={font=\small}}
$h$ & 板 $3x^2h$ & 角柱 $3xh^2$ & 小さな立方体 $h^3$ & ふえた体積の合計 \\
0.1 & 0.3 & 0.03 & 0.001 & 0.331 \\
0.01 & 0.03 & 0.0003 & 0.000001 & 0.030301 \\
\end{bktable}
\bktabcaption{$x=1$ のとき。$h$ を10分の1にすると、板は10分の1、角柱は100分の1、小さな立方体は1000分の1になる。}
\end{bkcenter}

$h=0.1$ では、板の分が合計のおよそ91\%です。$h=0.01$ では、およそ99\%になります。$h$ を小さくするほど、ふえた体積は、ほとんど板3まいの分だけになります。

なぜ、こんなにちがうのでしょう。板は、厚さだけが $h$ です。$h$ を $\dfrac{1}{10}$ にすると、厚さが $\dfrac{1}{10}$ になるだけなので、体積も $\dfrac{1}{10}$ です。角柱は、たてと横の2つが $h$ なので、$\dfrac{1}{10}\times\dfrac{1}{10}=\dfrac{1}{100}$ になります。小さな立方体は3つとも $h$ なので、$\dfrac{1}{10}\times\dfrac{1}{10}\times\dfrac{1}{10}=\dfrac{1}{1000}$ です。

平均変化率は、ふえた体積を $h$ でわったものです。表の合計を $h$ でわると、$0.331\div0.1=3.31$、$0.030301\div0.01=3.0301$ で、3に近づきます。板の分を $h$ でわった $3x^2h\div h=3x^2$ は、$x=1$ で3のまま変わりません。角柱と小さな立方体の分は、$h$ でわっても0に近づいていきます。

だから、体積のふえる速さ（$x^3$ の導関数）は、板3まい分の $3x^2$ です。式でも、表でも、図でも、同じ答え $(x^3)'=3x^2$ になりました。

\begin{bkcenter}
\begin{bktable}{colspec={Q[l,m,wd=22mm] Q[l,m,wd=50mm] Q[l,m,wd=46mm]}, rowsep=1.5pt, cells={font=\small}, row{1}={font=\sffamily\bfseries\small}}
方法 & したこと & $3x^2$ の「3」の正体 \\
表 & $h=0.001$ で傾きを計算し、ならべた & 3でわると $x^2$ になった \\
式 & $(x+h)^3$ を展開し、$h\to0$ にした & $x^2h$ の項が3個あった \\
図 & 立方体をのばし、ふえた部分を分けた & うすい板が3まいあった \\
\end{bktable}
\bktabcaption{3つの方法で、同じきまりを確かめた。}
\end{bkcenter}

\bksubsection{5}{$x^n$ を微分するきまり}

\leavevmode\label{bk:5-4:kimari}$x^2$、$x^3$、$x^4$ で見つけたきまりは、指数がもっと大きくても成り立ちます（理由は、この章のもっと知りたい）。$n$ を自然数（1、2、3、……）とすると、
\[ (x^n)'=nx^{n-1} \]

\begin{bkcenter}
\begin{bktable}{colspec={Q[c,m,wd=26mm] *{6}{Q[c,m,wd=13mm]}}, rowsep=2pt}
関数 & $x$ & $x^2$ & $x^3$ & $x^4$ & $x^5$ & $x^n$ \\
導関数 & $1$ & $2x$ & $3x^2$ & $4x^3$ & $5x^4$ & $nx^{n-1}$ \\
\end{bktable}
\bktabcaption{指数を前に出して、指数を1へらす。}
\end{bkcenter}

\noindent\textbf{$x$ を微分すると1}\par
\noindent $x=x^1$ と考えて、きまりを使うと $1\times x^{1-1}=1\times x^0=1\times1=1$ です（$x^0=1$ は2-7）。グラフでも確かめましょう。$y=x$ のグラフは、右へ1進むと上へ1上がる直線なので、どこでも傾きは1です（下の図の左）。

\noindent\textbf{定数を微分すると0}\par
\noindent $y=5$ のように、$x$ が変わっても変わらない数を\textbf{定数}といいます。$y=5$ のグラフは横にまっすぐな直線で、右へ進んでも上がりも下がりもしないので、傾きは0です（下の図の右）。平均変化率で計算しても、$\dfrac{5-5}{h}=0$ です。だから、\textbf{定数を微分すると0}です。\bkidx{ていすう}{定数}

\begin{bkfigure}
\begin{tikzpicture}
\begin{axis}[bkplot, name=g1, width=56mm, height=40mm, xmin=-0.3, xmax=4.4, ymin=-0.3, ymax=4.6, xtick={1,2,3,4}, ytick={1,2,3,4}, xlabel={$x$}, ylabel={$y=x$}]
  \addplot[bkcurve, domain=0:4.2] {x};
  \draw[bkOrange!80!black, line width=0.9pt] (axis cs:1.5,1.5) -- (axis cs:2.5,1.5) -- (axis cs:2.5,2.5);
  \node[font=\sffamily\scriptsize, text=bkOrange!80!black, anchor=north] at (axis cs:2,1.45) {右へ1};
  \node[font=\sffamily\scriptsize, text=bkOrange!80!black, anchor=west] at (axis cs:2.6,2) {上へ1};
\end{axis}
\begin{axis}[bkplot, at={($(g1.east)+(16mm,0)$)}, anchor=west, width=56mm, height=40mm, xmin=-0.3, xmax=4.4, ymin=-0.3, ymax=6.6, xtick={1,2,3,4}, ytick={5}, xlabel={$x$}, ylabel={$y=5$}]
  \addplot[bkcurve, domain=0:4.2] {5};
  \draw[bkOrange!80!black, line width=0.9pt, -{Stealth[length=2mm]}] (axis cs:1.5,4.1) -- (axis cs:2.5,4.1);
  \node[font=\sffamily\scriptsize, text=bkOrange!80!black, anchor=north, align=center] at (axis cs:2,3.6) {右へ1進んでも、\\上へ0};
\end{axis}
\end{tikzpicture}
\bkfigcaption{左：$y=x$ はどこでも傾き1。右：$y=5$ はどこでも傾き0。}
\end{bkfigure}

\noindent\textbf{マイナスの指数でも使える}\par
\noindent $\dfrac{1}{x}$ は、2-7のきまりで $x^{-1}$ と書けます。この指数 $-1$ にも、きまりを使ってみましょう。前に $-1$ を出して、指数を1へらすと $-1-1=-2$ です。
\[ \Bigl(\frac{1}{x}\Bigr)'=(x^{-1})'=-1\times x^{-2}=-\frac{1}{x^2} \]
ほんとうに合っているか、表計算で確かめます。B列の式を、次のようにします。
\begin{center}\texttt{=(1/(A2+0.001)-1/A2)/0.001}\end{center}
$x=0$ では $\dfrac{1}{0}$ が計算できない（0ではわれない）ので、$x=1$、2、4 で調べます。

\noindent\begin{minipage}[c]{0.5\linewidth}
\begin{bktable}{colspec={Q[c,m,wd=8mm] Q[c,m,wd=24mm] Q[c,m,wd=20mm]}, rowsep=1.5pt, cells={font=\small}}
$x$ & 表計算 & $-\dfrac{1}{x^2}$ \\
1 & $-0.999000$ & $-1$ \\
2 & $-0.249875$ & $-0.25$ \\
4 & $-0.062484$ & $-0.0625$ \\
\end{bktable}
\end{minipage}\hfill
\begin{minipage}[c]{0.46\linewidth}
表計算の値は、$-\dfrac{1}{x^2}$ とよく合っています。答えがマイナスになるのは、$y=\dfrac{1}{x}$ のグラフが右下がりだからです（下の図）。$x$ が大きくなると、$\dfrac{1}{x}$ は小さくなります。
\end{minipage}

\begin{bkfigure}
\begin{tikzpicture}
\begin{axis}[bkplot, width=70mm, height=44mm, xmin=0, xmax=4.4, ymin=0, ymax=2.6, xtick={1,2,3,4}, ytick={1,2}, xlabel={$x$}, ylabel={$y=\dfrac{1}{x}$}]
  \addplot[bkcurve, domain=0.4:4.3, samples=60] {1/x};
  \addplot[bkcurve2, line width=1pt, domain=0.25:1.9] {1-(x-1)};
  \addplot[only marks, mark=*, mark size=1.6pt, bkInk] coordinates {(1,1)};
  \node[font=\sffamily\scriptsize, text=bkOrange!80!black, anchor=west, inner sep=1.5pt] at (axis cs:2.2,1.9) {オレンジ：$x=1$ での接線（傾き $-1$）};
\end{axis}
\end{tikzpicture}
\bkfigcaption{$y=\frac{1}{x}$ のグラフと、$x=1$ での接線。右下がりなので、傾きはマイナス。}
\end{bkfigure}

このように、$n$ がマイナスの整数でも、きまり $(x^n)'=nx^{n-1}$ は使えます。指数を1へらすときは、$-1$ が $-2$ に、$-2$ が $-3$ になります（0の向きではなく、数直線の左へ1つ進む）。

\bksubsection{6}{定数倍と和の微分}

\leavevmode\label{bk:5-4:wa}\textbf{定数倍}\par
\noindent 5-2では、$x=5t^2$ の速さが $v=10t$ でした。これは、$(t^2)'=2t$ のちょうど5倍です。表でくらべてみましょう。

\begin{bkcenter}
\begin{bktable}{colspec={Q[l,m,wd=34mm] *{3}{Q[c,m,wd=14mm]}}, rowsep=1.5pt}
$t$ & 1 & 2 & 3 \\
$t^2$ の傾き $2t$ & 2 & 4 & 6 \\
$5t^2$ の傾き $10t$ & 10 & 20 & 30 \\
\end{bktable}
\bktabcaption{どの $t$ でも、$5t^2$ の傾きは $t^2$ の傾きの5倍。}
\end{bkcenter}

なぜ5倍になるのでしょう。$5t^2$ のグラフは、$t^2$ のグラフの高さを全部5倍にしたものです。高さが全部5倍なので、同じだけ右へ進んだときに上がる量も5倍です。下の図は、$t=1$ の所の接線です。$t^2$ の接線は、右へ1進むと上へ2上がります。$5t^2$ の接線は、右へ1進むと上へ10上がります。

\begin{bkfigure}
\begin{tikzpicture}
\begin{axis}[bkplot, name=c1, width=56mm, height=46mm, xmin=0, xmax=2.2, ymin=0, ymax=5.5, xtick={1,2}, ytick={1,2,3,4,5}, xlabel={$t$}, ylabel={$t^2$}]
  \addplot[bkcurve, domain=0:2.2] {x^2};
  \addplot[bkcurve2, line width=1pt, domain=0.5:2.2] {1+2*(x-1)};
  \addplot[only marks, mark=*, mark size=1.5pt, bkInk] coordinates {(1,1)};
  \draw[bkOrange!80!black, line width=0.8pt, dashed] (axis cs:1,1) -- (axis cs:2,1) -- (axis cs:2,3);
  \node[font=\sffamily\scriptsize, text=bkOrange!80!black, anchor=north] at (axis cs:1.5,0.9) {右へ1};
  \node[font=\sffamily\scriptsize, text=bkOrange!80!black, anchor=west] at (axis cs:2.03,2) {上へ2};
\end{axis}
\begin{axis}[bkplot, at={($(c1.east)+(16mm,0)$)}, anchor=west, width=56mm, height=46mm, xmin=0, xmax=2.2, ymin=0, ymax=22, xtick={1,2}, ytick={5,10,15,20}, xlabel={$t$}, ylabel={$5t^2$}]
  \addplot[bkcurve, domain=0:2.1] {5*x^2};
  \addplot[bkcurve2, line width=1pt, domain=0.6:2.2] {5+10*(x-1)};
  \addplot[only marks, mark=*, mark size=1.5pt, bkInk] coordinates {(1,5)};
  \draw[bkOrange!80!black, line width=0.8pt, dashed] (axis cs:1,5) -- (axis cs:2,5) -- (axis cs:2,15);
  \node[font=\sffamily\scriptsize, text=bkOrange!80!black, anchor=north] at (axis cs:1.5,4.6) {右へ1};
  \node[font=\sffamily\scriptsize, text=bkOrange!80!black, anchor=west] at (axis cs:2.03,10) {上へ10};
\end{axis}
\end{tikzpicture}
\bkfigcaption{$t=1$ での接線（オレンジ）。左：$t^2$ は傾き2。右：$5t^2$ は傾き10で、5倍。}
\end{bkfigure}
式でも、平均変化率の5を外に出せます。
\[ \frac{5(t+h)^2-5t^2}{h}=5\times\frac{(t+h)^2-t^2}{h}\ \longrightarrow\ 5\times2t=10t \]
このように、\textbf{数をかけた関数を微分するときは、数はそのままにして、残りを微分すればよい}のです。

\noindent\textbf{和}\par
\noindent 水そうに、2つの蛇口から水を入れることを考えます。Aの蛇口で水位が毎分\qty{2}{cm}、Bの蛇口で毎分\qty{3}{cm}上がるなら、両方を開けると、毎分 $2+3=\qty{5}{cm}$ 上がります。変わる速さは、たし算できます。

式でも確かめましょう。$f(x)=x^2+x$ の平均変化率は、2-12のきまりで2つに分けられます。
\begin{align*}
\frac{\{(x+h)^2+(x+h)\}-(x^2+x)}{h}&=\frac{(x+h)^2-x^2}{h}+\frac{(x+h)-x}{h}\\
&\longrightarrow\ 2x+1
\end{align*}
1つめは $x^2$ の平均変化率で $2x$ に、2つめは $x$ の平均変化率で1に近づきます。だから、\textbf{和を微分するときは、1つずつ微分してたせばよい}のです。

\Needspace{4\baselineskip}\noindent\textbf{差}\par
\noindent ひき算は、「$-1$ をかけてたす」と考えられます。$x^2-x=x^2+(-1)\times x$ なので、定数倍と和のきまりから、$(x^2-x)'=2x+(-1)\times1=2x-1$。ひき算も、1つずつ微分してひけばよいのです。

\begin{taisetsu}
$(x^n)'=nx^{n-1}$\qquad 定数を微分すると0\\
定数倍：$\{cf(x)\}'=cf'(x)$\qquad 和：$\{f(x)+g(x)\}'=f'(x)+g'(x)$
\end{taisetsu}

$3x^3-2x^2+5x-7$ のように、$x^n$ に数をかけたものをたし合わせた式を、\textbf{多項式}といいます。多項式は、1つの項ずつ微分してたせば微分できます。\bkidx{たこうしき}{多項式}

\begin{bkcenter}
\begin{bktable}{colspec={Q[c,m,wd=16mm] Q[l,m,wd=58mm] Q[c,m,wd=30mm]}, rowsep=1.5pt, cells={font=\small}, row{1}={font=\sffamily\bfseries\small}}
項 & 考え方 & 微分すると \\
$3x^3$ & 3はそのまま、$x^3$ を微分して $3x^2$ & $3\times3x^2=9x^2$ \\
$-2x^2$ & $-2$ はそのまま、$x^2$ を微分して $2x$ & $-2\times2x=-4x$ \\
$5x$ & 5はそのまま、$x$ を微分して1 & $5\times1=5$ \\
$-7$ & 定数なので0 & $0$ \\
\end{bktable}
\bktabcaption{$3x^3-2x^2+5x-7$ を、1つの項ずつ微分する。}
\end{bkcenter}
\[ (3x^3-2x^2+5x-7)'=9x^2-4x+5 \]

\begin{motto}[3]
\textbf{くり返し微分する}\par
微分した式も、また微分できます。
\[ x^4\ \xrightarrow{\text{微分}}\ 4x^3\ \xrightarrow{\text{微分}}\ 12x^2\ \xrightarrow{\text{微分}}\ 24x\ \xrightarrow{\text{微分}}\ 24\ \xrightarrow{\text{微分}}\ 0 \]
$4x^3$ の微分は、4はそのままで $x^3$ を微分して $4\times3x^2=12x^2$ です。多項式は、微分するたびに指数が1つずつへるので、いつか定数になり、その次は0になります。2回微分することには、物理で大事な意味があります（5-5）。\bkidx{にかいびぶん}{2回微分}
\end{motto}

\bksubsection{7}{グラフで見くらべる：$f(x)$ と $f'(x)$}

\leavevmode\label{bk:5-4:graph}\noindent\begin{minipage}[c]{0.5\linewidth}
\begin{tikzpicture}
% 第6回：上下の間を広げ、目もり数字に線や曲線がかからないようにする（±1 の数字は点線の横へ、−2・−3 は曲線のない側へ）
\begin{axis}[bkplot, name=top, width=62mm, height=38mm, xmin=-2.3, xmax=2.5, ymin=-2.6, ymax=2.8, xtick={-2,-1,1,2}, xticklabels={,,,$2$}, ytick={-2,2}, xlabel={$x$}, ylabel={$f(x)=x^3-3x$}]
  \addplot[bkcurve, domain=-2.06:2.06] {x^3-3*x};
  \draw[bkGray, dashed, line width=0.5pt] (axis cs:-1,-2.6) -- (axis cs:-1,2.6) (axis cs:1,-2.6) -- (axis cs:1,2.6);
  \node[anchor=north east, font=\sffamily\footnotesize, text=bkInk, inner sep=1.5pt, xshift=-0.5mm] at (axis cs:-2,0) {$-2$};
  \node[anchor=north east, font=\sffamily\footnotesize, text=bkInk, inner sep=1.5pt, xshift=-1mm] at (axis cs:-1,0) {$-1$};
  \node[anchor=north west, font=\sffamily\footnotesize, text=bkInk, inner sep=1.5pt, xshift=1mm] at (axis cs:1,0) {$1$};
\end{axis}
\begin{axis}[bkplot, at={($(top.south)+(0,-12mm)$)}, anchor=north, width=62mm, height=38mm, xmin=-2.3, xmax=2.5, ymin=-4, ymax=9.5, xtick={-2,-1,1,2}, xticklabels={$-2$,,,$2$}, ytick={-3,3,6,9}, yticklabels={,$3$,$6$,$9$}, xlabel={$x$}, ylabel={$f'(x)=3x^2-3$}]
  \addplot[bkcurve2, domain=-2.04:2.04] {3*x^2-3};
  \draw[bkGray, dashed, line width=0.5pt] (axis cs:-1,-4) -- (axis cs:-1,9.5) (axis cs:1,-4) -- (axis cs:1,9.5);
  \node[anchor=north east, font=\sffamily\footnotesize, text=bkInk, inner sep=1.5pt, xshift=-1mm] at (axis cs:-1,0) {$-1$};
  \node[anchor=north west, font=\sffamily\footnotesize, text=bkInk, inner sep=1.5pt, xshift=1mm] at (axis cs:1,0) {$1$};
  \node[anchor=north east, font=\sffamily\footnotesize, text=bkInk, inner sep=1.5pt] at (axis cs:0,-3) {$-3$};
\end{axis}
\end{tikzpicture}
\end{minipage}\hfill
\begin{minipage}[c]{0.47\linewidth}\raggedright
$f(x)=x^3-3x$（青）を、ここまでのきまりで微分すると、$f'(x)=3x^2-3$ です。上下にならべると、次のことがわかります。
\begin{itemize}[leftmargin=1.2\zw, itemsep=1pt, topsep=2pt, label={\color{bkTheme}●}]
\item $f'(x)$ が正の所（$x<-1$、$1<x$）では、$f(x)$ は右上がり。
\item $f'(x)$ が負の所（$-1<x<1$）では、$f(x)$ は右下がり。
\item $f'(x)=0$ の所（$x=\pm1$）で、$f(x)$ は山のてっぺんや谷の底になる。
\end{itemize}
\end{minipage}

\vspace{2pt}
たとえば $x=2$ では、$f'(2)=3\times4-3=9$ で、$f(x)$ は急な右上がりです。$x=0$ では $f'(0)=-3$ で、右下がりです。$x=1$ では $f'(1)=3-3=0$ で、グラフは一瞬だけ平らになり、谷の底になります。この見方は、5-7で「いちばん高い所」や「いちばん低い所」を見つけるのに使います。



\begin{motto}[3]
\textbf{どんな $n$ でも成り立つわけ}\par\vspace{4pt}
\noindent\begin{minipage}[c]{0.3\linewidth}\centering\small
$1$\\ $1\quad1$\\ $1\quad2\quad1$\\ $1\quad{\color{bkOrange}\mathbf{3}}\quad3\quad1$\\ $1\quad{\color{bkOrange}\mathbf{4}}\quad6\quad4\quad1$\\[2pt]
{\footnotesize $(x+h)^n$ の係数}
\end{minipage}\hfill
\begin{minipage}[c]{22\zw}
$(x+h)^n$ を展開したときの係数は、左の三角形のように、上の2つの数をたしてつくれます。たとえば $n=3$ なら、
\[ (x+h)^3=x^3+3x^2h+3xh^2+h^3 \]
左から2番めの数（オレンジ）は、いつも $n$ です。
\end{minipage}\par\vspace{3pt}
$(x+h)^n$ は、$(x+h)$ を $n$ 個かけたものです。展開すると、$h$ を1つもふくまない項は $x^n$ です。$h$ を1つだけふくむ項は、$n$ 個のかっこのどれか1つから $h$ を、残りから $x$ をえらんだもので、$x^{n-1}h$ が $n$ 個できます（立方体なら、板が3まい）。残りの項は $h^2$ 以上をふくむので、$h$ でわっても $h$ が残り、$h\to0$ で0になります（角柱や小さな立方体）。だから $(x^n)'=nx^{n-1}$ です。
\end{motto}

\begin{kotoba}
\small
\bkword{指数}{$x^3$ の右肩の3のように、何回かけるかを表す数。}
\bkword{定数}{5や $-7$ のように、$x$ が変わっても変わらない数。微分すると0。}
\bkword{多項式}{$3x^3-2x^2+5x-7$ のように、$x^n$ に数をかけたものをたし合わせた式。}
\bkword{$f'(x)$}{「エフ・ダッシュ・エックス」。$f(x)$ の導関数。$\dfrac{\dd f}{\dd x}$ とも書く。}
\end{kotoba}

%% ---------------- 例題 ----------------
\bksection{reidai}{例題でたしかめよう}

\begin{reidai}[$x^n$ と定数の微分]
次の関数を微分しましょう。\quad（1）$x^5$\quad（2）$x$\quad（3）$7$\par
\kangaekata（1）（2）は $(x^n)'=nx^{n-1}$ を使います。指数を前に出して、指数を1へらします。（3）は定数です。\par
\shiki（1）指数5を前に出し、指数を $5-1=4$ にする。$5x^4$\par
\hspace*{2.6\zw}（2）$x=x^1$ なので、$1\times x^{1-1}=1\times x^0=1$\par
\hspace*{2.6\zw}（3）7は定数なので、微分すると0\par
\kotae （1）$5x^4$\quad（2）$1$\quad（3）$0$
\end{reidai}

\begin{jibun}
次の関数を微分しましょう。\quad（1）$x^6$\quad（2）$x^{10}$\quad（3）$-3$
\bkjans{（1）\nobreak $6x^5$\quad（2）\nobreak $10x^9$\quad（3）\nobreak $0$}
\end{jibun}

\begin{reidai}[多項式の微分と微分係数]
$f(x)=3x^3-2x^2+5x-7$ を微分し、$f'(1)$ と $f'(-1)$ を求めましょう。\par
\kangaekata 1つの項ずつ微分します。数はそのまま、$x^n$ を微分します。$f'(-1)$ は、$x$ に $-1$ をかっこごと入れます（2-13）。\par
\shiki $f'(x)=3\times3x^2-2\times2x+5\times1-0=9x^2-4x+5$\par
\hspace*{2.6\zw}$f'(1)=9\times1^2-4\times1+5=9-4+5=10$\par
\hspace*{2.6\zw}$f'(-1)=9\times(-1)^2-4\times(-1)+5=9+4+5=18$\par
\kotae $f'(x)=9x^2-4x+5$、$f'(1)=10$、$f'(-1)=18$
\end{reidai}

\begin{jibun}
$f(x)=2x^3+x^2-4x+1$ を微分し、$f'(2)$ と $f'(-1)$ を求めましょう。
\bkjans{$f'(x)=6x^2+2x-4$、$f'(2)=24$、$f'(-1)=0$}
\end{jibun}

\begin{reidai}[分数の係数がある多項式]
$f(x)=\dfrac{1}{3}x^3-x^2+2$ を微分しましょう。\par
\kangaekata 分数がかかっていても、定数倍のきまりは同じです。分数はそのままにして、$x^n$ を微分し、最後に約分します。\par
\shiki $\dfrac{1}{3}x^3$ の微分は $\dfrac{1}{3}\times3x^2=x^2$（$\dfrac{1}{3}\times3=1$）\par
\hspace*{2.6\zw}$-x^2$ の微分は $-2x$、2の微分は0\par
\kotae $f'(x)=x^2-2x$
\end{reidai}

\begin{jibun}
$f(x)=\dfrac{1}{2}x^4-\dfrac{2}{3}x^3+5$ を微分しましょう。
\bkjans{$\dfrac{1}{2}\times4x^3-\dfrac{2}{3}\times3x^2=2x^3-2x^2$}
\end{jibun}

\begin{reidai}[位置から速さを求める]
投げ下ろしたボールの位置が $x=4.9t^2+2t+1$（m）のとき、$t$ 秒の瞬間の速さ $v$ と、3秒の速さを求めましょう。\par
\kangaekata 瞬間の速さは、位置を時間で微分したもの（$v=\dfrac{\dd x}{\dd t}$、5-3）です。ここでは $x$ が関数の名前で、$t$ が変わる数です。$t$ について、1つの項ずつ微分します。\par
\shiki $4.9t^2$ の微分は $4.9\times2t=9.8t$、$2t$ の微分は2、1の微分は0\par
\hspace*{2.6\zw}$v=9.8t+2$\par
\hspace*{2.6\zw}$t=3$ のとき $v=9.8\times3+2=29.4+2=31.4$\par
\kotae $v=9.8t+2$、\qty{31.4}{m/s}
\end{reidai}

\begin{jibun}
真上に投げたボールの高さが $x=-5t^2+20t$（m）のとき、$v$ を求めましょう。1秒と3秒の速さはいくつですか。
\bkjans{$v=-10t+20$。1秒：\qty{10}{m/s}、3秒：$\qty{-10}{m/s}$（下向きに\qty{10}{m/s}）}
\end{jibun}

\begin{reidai}[マイナスの指数]
$f(x)=\dfrac{1}{x^2}$ を微分し、$f'(1)$ を求めましょう。\par
\kangaekata $\dfrac{1}{x^2}=x^{-2}$（2-7）と書きなおして、きまりを使います。指数を1へらすと、$-2-1=-3$ です。\par
\shiki $(x^{-2})'=-2\times x^{-2-1}=-2x^{-3}=-\dfrac{2}{x^3}$\par
\hspace*{2.6\zw}$f'(1)=-\dfrac{2}{1^3}=-2$\par
\kotae $f'(x)=-\dfrac{2}{x^3}$、$f'(1)=-2$
\end{reidai}

\begin{jibun}
$f(x)=\dfrac{1}{x}$ を微分し、$f'(2)$ を求めましょう。
\bkjans{$f'(x)=-\dfrac{1}{x^2}$、$f'(2)=-\dfrac{1}{4}$}
\end{jibun}

\begin{reidai}[接線の式]
$y=x^3-2x$ のグラフの、$x=1$ の点での接線の式を求めましょう。\par
\kangaekata 接線の式は $y-f(a)=f'(a)(x-a)$（5-3）。接点の $y$ 座標 $f(1)$ と、傾き $f'(1)$ を順に求めます。\par
\shiki 接点：$x=1$ のとき $y=1^3-2\times1=1-2=-1$。接点は $(1,\ -1)$\par
\hspace*{2.6\zw}傾き：$y'=3x^2-2$ なので、$3\times1^2-2=1$\par
\hspace*{2.6\zw}$y-(-1)=1\times(x-1)$ より $y+1=x-1$。$y=x-2$\par
\kotae $y=x-2$
\end{reidai}

\begin{jibun}
$y=x^3-2x$ のグラフの、$x=2$ の点での接線の式を求めましょう。
\bkjans{接点 $(2,\ 4)$、傾き $3\times4-2=10$。$y=10x-16$}
\end{jibun}

\begin{reidai}[展開してから微分する]
次の関数を微分しましょう。\quad（1）$(x+2)(x-3)$\quad（2）$(2x+1)^2$\par
\kangaekata かっこのままでは、まだきまりが使えません。展開して $x^n$ の和にしてから、1つの項ずつ微分します（3-8）。\par
\shiki（1）$(x+2)(x-3)=x^2-3x+2x-6=x^2-x-6$\par
\hspace*{4.3\zw}微分して $2x-1$\par
\hspace*{2.6\zw}（2）$(2x+1)^2=(2x)^2+2\times2x\times1+1^2=4x^2+4x+1$\par
\hspace*{4.3\zw}微分して $4\times2x+4=8x+4$\par
\kotae （1）$2x-1$\quad（2）$8x+4$
\end{reidai}

\begin{jibun}
次の関数を微分しましょう。\quad（1）$(x-1)(x+4)$\quad（2）$(x+3)^2$
\bkjans{（1）\nobreak $x^2+3x-4$ より $2x+3$\quad（2）\nobreak $x^2+6x+9$ より $2x+6$}
\end{jibun}

\begin{reidai}[速さが0になるとき]
真上に投げたボールの高さが $x=-5t^2+30t$（m）です。速さが0になる時刻と、そのときの高さを求めましょう。\par
\kangaekata ボールは上がりながらおそくなり、いちばん高い所で一瞬止まって、向きが変わります。その瞬間は、速さ $v$ が0です。\par
\shiki $v=-5\times2t+30=-10t+30$\par
\hspace*{2.6\zw}$-10t+30=0$ より $10t=30$、$t=3$\par
\hspace*{2.6\zw}$x=-5\times3^2+30\times3=-45+90=45$\par
\kotae 3秒後、\qty{45}{m}（いちばん高い所）
\end{reidai}

\begin{jibun}
真上に投げたボールの高さが $x=-5t^2+20t$（m）です。速さが0になる時刻と、そのときの高さを求めましょう。
\bkjans{$v=-10t+20$。2秒後、\qty{20}{m}}
\end{jibun}

\begin{reidai}[山のてっぺんと谷の底]
$f(x)=x^3-12x$ について、$f'(x)=0$ となる $x$ を求め、そのときの $f(x)$ の値を答えましょう。それぞれ、グラフの山のてっぺんですか、谷の底ですか。\par
\kangaekata 数学の道具7のように、$f'(x)=0$ の所で、グラフは一瞬平らになります。その前後で $f'(x)$ の符号（右上がりか右下がりか）を調べます。\par
\shiki $f'(x)=3x^2-12$。$3x^2-12=0$ より $x^2=4$\par
\hspace*{2.6\zw}2乗して4になる数は、2と $-2$ の2つ（2-8）\par
\hspace*{2.6\zw}$f(-2)=(-2)^3-12\times(-2)=-8+24=16$\par
\hspace*{2.6\zw}$f(2)=2^3-12\times2=8-24=-16$\par
\hspace*{2.6\zw}$x=0$ では $f'(0)=-12$（右下がり）、$x=3$ では $f'(3)=15$（右上がり）\par
\hspace*{2.6\zw}右上がり→平ら→右下がりの $x=-2$ は山、右下がり→平ら→右上がりの $x=2$ は谷\par
\kotae $x=-2$ で $f(x)=16$（山のてっぺん）、$x=2$ で $f(x)=-16$（谷の底）
\end{reidai}

\begin{jibun}
$f(x)=x^3-27x$ について、$f'(x)=0$ となる $x$ を求め、そのときの $f(x)$ の値を答えましょう。それぞれ、山のてっぺんですか、谷の底ですか。
\bkjans{$f'(x)=3x^2-27=0$ より $x=-3,\ 3$。$f(-3)=54$（山のてっぺん）、$f(3)=-54$（谷の底）}
\end{jibun}

\begin{reidai}[傾きから接点を探す]
$y=x^3$ のグラフで、接線の傾きが12になる点を、すべて求めましょう。\par
\kangaekata 接線の傾きは $y'$ の値です（5-3）。$y'=12$ となる $x$ を求めてから、$y$ 座標を計算します。\par
\shiki $y'=3x^2$。$3x^2=12$ より $x^2=4$、$x=2,\ -2$\par
\hspace*{2.6\zw}$x=2$ のとき $y=2^3=8$、$x=-2$ のとき $y=(-2)^3=-8$\par
\kotae $(2,\ 8)$ と $(-2,\ -8)$ の2つ
\end{reidai}

\begin{jibun}
$y=x^3$ のグラフで、接線の傾きが3になる点を、すべて求めましょう。
\bkjans{$3x^2=3$ より $x=1,\ -1$。$(1,\ 1)$ と $(-1,\ -1)$}
\end{jibun}

\begin{reidai}[文字の定数をふくむ式]
円の面積 $S=\pi r^2$ を、半径 $r$ で微分しましょう。\par
\kangaekata $\pi$ は 3.14159… という決まった数（定数）です。文字で書いてあっても、定数倍のきまりで、そのままにしておきます。\par
\shiki $\dfrac{\dd S}{\dd r}=\pi\times(r^2)'=\pi\times2r=2\pi r$\par
\kotae $\dfrac{\dd S}{\dd r}=2\pi r$
\end{reidai}

\begin{jibun}
1辺 $a$ の正方形の面積 $S=a^2$ と、1辺 $a$ の立方体の体積 $V=a^3$ を、それぞれ $a$ で微分しましょう。
\bkjans{$\dfrac{\dd S}{\dd a}=2a$、$\dfrac{\dd V}{\dd a}=3a^2$}
\end{jibun}

%% ---------------- まちがえやすいところ・物理工学 ----------------
\bksection{tsumazuki}{まちがえやすいところをチェック}

\begin{tsumazuki}{指数を前に出し忘れる・1へらし忘れる}
\bkng{$(x^3)'=x^2$}\bkng{$(x^3)'=3x^3$}\bkok{$(x^3)'=3x^2$}
「前に出す」と「1へらす」はセット。表の値（$x=2$ で12）で確かめよう。$x^2$ なら $2^2=4$、$3x^3$ なら $3\times8=24$ で、どちらも12になりません。
\end{tsumazuki}

\begin{tsumazuki}{定数を残してしまう}
\bkng{$(x^2+5)'=2x+5$}\bkok{$(x^2+5)'=2x$}
グラフを上に5ずらしても、坂の急さ（傾き）は同じです。たした定数は、微分すると消えます。
\end{tsumazuki}

\begin{tsumazuki}{マイナスの指数を、0の向きにへらす}
\bkng{$(x^{-2})'=-2x^{-1}$}\bkok{$(x^{-2})'=-2x^{-3}$}
「1へらす」は、1をひくことです。$-2-1=-3$ です。$-2$ から1ひくと、数直線の左へ進むので、$-1$ ではなく $-3$ になります。
\end{tsumazuki}

\begin{tsumazuki}{前にかかっている数を忘れる}
\bkng{$(3x^2)'=2x$}\bkok{$(3x^2)'=3\times2x=6x$}
前の3はそのまま残して、$x^2$ の微分 $2x$ にかけます。$3x^2$ のグラフは $x^2$ の高さを3倍にしたものなので、傾きも3倍です。
\end{tsumazuki}

\begin{tsumazuki}{$x$ の微分を0にしてしまう}
\bkng{$(5x)'=0$}\bkok{$(5x)'=5$}
0になるのは、$x$ をふくまない定数だけです。$5x$ は $x$ が1ふえると5ふえるので、傾きは5です。
\end{tsumazuki}

\begin{column}{ケプラーも、表からきまりを見つけた}
400年ほど前、ドイツのケプラーは、惑星が太陽を1周する時間（公転周期 $T$ 年）と、太陽からのきょり（$a$。地球を1とする）の表をながめて、きまりを探しました。何年も計算を続けて、$T^2$ と $a^3$ がいつも同じになることに気づきました（1619年）。火星は $a=1.524$、$T=1.881$ で、$T^2=3.54$、$a^3=3.54$。木星は $a=5.203$、$T=11.86$ で、$T^2=140.7$、$a^3=140.9$ です。その理由は、のちにニュートンが式で説明しました。「数で見つけて、式で確かめる」は、この章でしたことと同じ流れです。
\end{column}

\begin{butsuri}
\noindent\begin{minipage}[c]{0.74\linewidth}
\textbf{体積のふえ方は、表面積}\par
半径 $r$ の球の体積 $V=\dfrac{4}{3}\pi r^3$（4-5）を $r$ で微分すると、
\[ \frac{\dd V}{\dd r}=\frac{4}{3}\pi\times3r^2=4\pi r^2 \]
で、球の表面積（4-5）と同じです。球に、厚さ $h$ のうすいペンキをぬると、ふえる体積は、ほぼ「表面積×厚さ」です。立方体で、ふえた体積がほぼ板3まい（表面の分）だったのと同じわけです。円の面積 $\pi r^2$ の微分が円周 $2\pi r$ になるのも、同じわけです。\bkidx{えんのめんせき}{円の面積}
\end{minipage}\hfill
\begin{minipage}[c]{0.22\linewidth}\centering
\begin{tikzpicture}[x=1mm, y=1mm]
  \shade[ball color=bkBlue!40] (0,0) circle (10);
  \draw[bkOrange, line width=2.2pt] (0,0) circle (11.2);
  \draw[bkInk, line width=0.6pt] (0,0) -- (10,0);
  \node[font=\sffamily\scriptsize, anchor=south, inner sep=1.5pt] at (5,0) {$r$};
  \node[font=\sffamily\scriptsize, text=bkOrange!80!black, anchor=north] at (0,-13.5) {うすいペンキ};
\end{tikzpicture}
\end{minipage}
\end{butsuri}

\begin{column}{コンピュータに微分させる}
式を文字のまま計算するソフト（数式処理ソフト）を使うと、コンピュータが微分してくれます。たとえば \mbox{\texttt{diff(3*x\^{}3-2*x\^{}2+5*x-7, x)}} と入力すると、\mbox{\texttt{9*x\^{}2-4*x+5}} と答えが出ます。計算のまちがいを確かめるのに便利です。でも、どんな意味の計算なのかがわかっていないと、答えが正しいか判断できません。この巻では、自分の手で計算します。
\end{column}

%% ---------------- 章末問題 ----------------
\bkendproblems
\bkproblemgroup{基本}{3}

\bkproblem 次の関数を微分しましょう。\quad（1）$x^8$\quad（2）$4x^3$\quad（3）$x^2+3x-2$
\bkpans{（1）\nobreak $8x^7$\quad（2）\nobreak $12x^2$\quad（3）\nobreak $2x+3$}

\bkproblem 次の関数を微分しましょう。\quad（1）$\dfrac{1}{4}x^4$\quad（2）$\dfrac{2}{3}x^3-\dfrac{1}{2}x^2+x$
\bkpans{（1）\nobreak $x^3$\quad（2）\nobreak $2x^2-x+1$}

\bkproblem 次の関数の、$x=2$ での微分係数を求めましょう。\quad（1）$5x^3$\quad（2）$x^3-4x^2+3x$
\bkpans{（1）\nobreak $15x^2$ より60\quad（2）\nobreak $3x^2-8x+3$ より $-1$}

\bkproblem $f(x)=x^3-6x$ について、$f'(x)$ と $f'(2)$ を求めましょう。
\bkpans{$f'(x)=3x^2-6$、$f'(2)=6$}

\bkproblem 位置が $x=3t^2-12t+5$（m）の物体について、次を求めましょう。\par
（1）速さ $v$（$t$ の式）\quad（2）1秒の速さ\quad（3）速さが0になる時刻
\bkpans{（1）\nobreak $v=6t-12$\quad（2）\nobreak $\qty{-6}{m/s}$\quad（3）\nobreak 2秒}

\bkproblem $f(x)=\dfrac{1}{x^3}$ を微分しましょう。
\bkpans{$-\dfrac{3}{x^4}$}

\bkproblem 展開してから微分しましょう。\quad（1）$(x+1)(x+5)$\quad（2）$(3x-1)^2$
\bkpans{（1）\nobreak $2x+6$\quad（2）\nobreak $9x^2-6x+1$ より $18x-6$}

\bkproblem $f(x)=x^4-2x^2$ の $f'(x)$ と、$f'(-1)$、$f'(0)$、$f'(2)$ を求めましょう。
\bkpans{$f'(x)=4x^3-4x$。$f'(-1)=0$、$f'(0)=0$、$f'(2)=24$}

\bkproblem $y=\dfrac{1}{x}$ のグラフの、$x=1$ の点での接線の式を求めましょう（例題5・6）。
\bkpans{接点 $(1,\ 1)$、傾き $-1$。$y=-x+2$}

\bkproblem 接線の傾きが、$y=x^2$ のグラフで $-6$、$y=x^3$ のグラフで27になる点を、すべて求めましょう。
\bkpans{$y=x^2$：$2x=-6$ より $(-3,\ 9)$。$y=x^3$：$3x^2=27$ より $(3,\ 27)$ と $(-3,\ -27)$}

\bkproblemgroup{チャレンジ}{3}

\bkproblem 半径 $r$ の円の面積 $S=\pi r^2$ を $r$ で微分しましょう。その結果は、円の何を表していますか。
\bkpans{$\dfrac{\dd S}{\dd r}=2\pi r$。円周の長さ}

\bkproblem $y=x^3$ のグラフの、$x=1$ の点での接線の式を求めましょう（例題6と同じ手順です）。
\bkpans{接点 $(1,\ 1)$、傾き $3$。$y=3x-2$}

\bkproblem 真上に投げたボールの高さが $x=-4.9t^2+19.6t$（m）です。いちばん高くなる時刻とその高さ、投げた瞬間の速さを求めましょう。
\bkpans{$v=-9.8t+19.6$。2秒後に\qty{19.6}{m}。投げた瞬間（$t=0$）は\qty{19.6}{m/s}}

\bkproblem $f(x)=x^3-3x$ について、$f'(x)=0$ となる $x$ を2つ求め、そのときの $f(x)$ の値を答えましょう。それぞれ、グラフの山のてっぺんと谷の底のどちらですか（例題9）。
\bkpans{$f'(x)=3x^2-3=0$ より $x^2=1$、$x=1,\ -1$。$f(-1)=2$（山のてっぺん）、$f(1)=-2$（谷の底）}

\begin{furikaeri}
\begin{checklist}
\item 表の結果から、微分のきまりを見つけられる。\hfill{\footnotesize\sffamily\color{bkTheme}→\pageref{bk:5-4:kimari0}ページ}
\item $(x^n)'=nx^{n-1}$ を使って微分できる。\hfill{\footnotesize\sffamily\color{bkTheme}→\pageref{bk:5-4:kimari}ページ}
\item 定数倍と和のきまりで、多項式を微分できる。\hfill{\footnotesize\sffamily\color{bkTheme}→\pageref{bk:5-4:wa}ページ}
\end{checklist}
\end{furikaeri}

\bkfukidashi{mina}{位置を微分すると速さ。では、速さをもう一度微分すると何になるのかな？}

%% 5-5 位置→速度→加速度（第11回：ページ制限なし。説明・図・例題を書き足した）
\bkchapter{5}{5}{位置→速度→加速度}{\bklevel{3}}{位置→速度→加速度}
\bkmichishirube{2-10,5-1,5-4}{5-13,6-1,6-15}

\begin{goal}
\begin{itemize}[label={\color{bkTheme}●}, leftmargin=1.6\zw, itemsep=0pt, topsep=1pt]
\item \bklevel{3}\enspace 速度と速さのちがい（向き）を説明できる。
\item \bklevel{3}\enspace 位置を2回微分して、加速度を求められる。
\item \bklevel{3}\enspace $\dot{x}$、$\ddot{x}$、$\dfrac{\dd^2x}{\dd t^2}$ の書き方を読める。
\end{itemize}
\end{goal}

\begin{hatena}
\noindent\begin{minipage}[c]{0.49\linewidth}
エレベーターに乗ると、上へ動き出すとき、体がぐっと重く感じます。止まるときは、ふわっと軽く感じます。でも、同じ速さで動いているあいだは、止まっているときと同じで、何も感じません。

体は、速さそのものではなく、何を感じているのでしょう。
\end{minipage}\hfill
\begin{minipage}[c]{0.48\linewidth}\centering
\begin{tikzpicture}[x=1mm, y=1mm]
  % 第11回（X14）：カゴは上面のまん中のロープで、機械室の滑車（巻き上げ機）から吊るす。滑車の反対側につり合いおもり
  % 昇降路のかべ
  \draw[bkGray, line width=0.9pt] (0,0) -- (0,47) (32,0) -- (32,47);
  \draw[bkGray, line width=0.6pt] (0,0) -- (32,0);
  % 機械室
  \draw[bkGray, line width=0.9pt, fill=bkGray!10] (0,47) rectangle (32,64);
  \fill[white] (10,46.4) rectangle (12,47.6) (24,46.4) rectangle (26,47.6);
  % 巻き上げ機（モーター）と滑車
  \draw[bkInk!70, line width=0.6pt, fill=bkGray!45, rounded corners=0.8] (1.6,51.5) rectangle (7,58.5);
  \draw[bkInk!70, line width=1pt] (7,55) -- (18,55);
  \draw[bkInk!80, line width=0.8pt, fill=bkGray!25] (18,55) circle (7);
  \fill[bkInk!80] (18,55) circle (1);
  % ロープ：カゴの上面のまん中 → 滑車の左 → 上を回って → 滑車の右 → おもり
  \draw[bkInk, line width=0.9pt] (11,29) -- (11,55) arc[start angle=180, end angle=0, radius=7] -- (25,35);
  % カゴ
  \draw[bkInk!80, line width=0.8pt, fill=bkPaper] (2,4) rectangle (20,29);
  \fill[bkInk!80] (9.8,29) rectangle (12.2,30.4);
  \bkperson{(9,6)}{14}{bkRoute6}
  \node[font=\sffamily\scriptsize, anchor=north east, text=bkGray, inner sep=1pt] at (19.2,28.2) {カゴ};
  % つり合いおもり
  \draw[bkInk!80, line width=0.7pt, fill=bkGray!55] (22.8,24) rectangle (27.2,35);
  % 名前
  \node[font=\sffamily\scriptsize, anchor=west, text=bkGray] at (0.8,61.6) {機械室};
  \node[font=\sffamily\scriptsize, anchor=north, text=bkGray, inner sep=1pt] at (25.4,22.6) {おもり};
  % 動く向きとラベル（かべの外）
  \draw[bk arrow, bkOrange] (36,6) -- (36,40);
  \node[anchor=west, font=\sffamily\scriptsize, align=left] at (38,40) {止まる\\軽い};
  \node[anchor=west, font=\sffamily\scriptsize, align=left] at (38,23) {一定の速さ\\ふつう};
  \node[anchor=west, font=\sffamily\scriptsize, align=left] at (38,6) {動き出す\\重い};
\end{tikzpicture}
\end{minipage}
\end{hatena}

体が感じているのは、速さそのものではなく、「速さが変わること」です。この章では、速さが変わるようすを数で表す方法を学びます。位置を微分すると速度、速度をもう一度微分すると加速度になります。

\clearpage
%% ---------------- やってみよう・思い出そう ----------------
\begin{yattemiyou}[車の速さの変わり方を調べよう]
止まっていた車が走り出し、しばらく同じ速さで走ってから、ブレーキをかけました。スピードメーターの値（km/h）を1秒ごとに記録しました。
\begin{enumerate}[label={\bkcircled[bkGreen]{\arabic*}}, leftmargin=2.4\zw, itemsep=1pt, topsep=3pt]
\item スピードメーターの値を、m/s になおそう（2-10）。たとえば\qty{36}{km/h}は、
\[ \qty{36}{km/h}\times\qty{1000}{m/km}\div\qty{3600}{s/h}=\qty{10}{m/s} \]
\qty{1000}{m/km} は「1 km あたり 1000 m」、\qty{3600}{s/h} は「1 時間あたり 3600 秒」という意味です。km どうし、h どうしが約分されて消え、m/s が残ります。
\item 1秒ごとに、速さがどれだけ変わったか（あとの速さ${}-{}$まえの速さ）を計算しよう。
\end{enumerate}
\begin{bkcenter}
\begin{bktable}{colspec={Q[l,m,wd=21mm] *{7}{Q[c,m,wd=7mm]} *{2}{Q[c,m,wd=8.6mm]}}, rowsep=1.5pt, colsep=1.1mm}
時刻 $t$（s） & 0 & 1 & 2 & 3 & 4 & 5 & 6 & 7 & 8 \\
メーター（km/h） & 0 & 9 & 18 & 27 & 36 & 36 & 36 & 25.2 & 14.4 \\
速さ $v$（m/s） & 0 & 2.5 & 5 & 7.5 & 10 & 10 & 10 & 7 & 4 \\
1秒ごとの変化 & & 2.5 & 2.5 & 2.5 & 2.5 & 0 & 0 & $-3$ & $-3$ \\
\end{bktable}
\bktabcaption{4行めは「あとの速さ${}-{}$まえの速さ」（単位は m/s）。}
\end{bkcenter}
\end{yattemiyou}

% 第7回：左の段の幅を字の幅の整数倍にして、字間がのびないようにする
\noindent\begin{minipage}[c]{14\zw}
走り出したときは、1秒ごとに速さが\qty{2.5}{m/s}ずつふえています。同じ速さのときは0、ブレーキをかけたときは、1秒ごとに\qty{3}{m/s}ずつへっています。速さと時間のグラフ（$v$-$t$ グラフ、2-3）にすると、この「1秒ごとの変化」は、グラフの傾きです。
\end{minipage}\hfill
\begin{minipage}[c]{0.57\linewidth}\centering
\begin{tikzpicture}
\begin{axis}[bkplot, width=70mm, height=40mm, xmin=0, xmax=8.6, ymin=0, ymax=12.5, xtick={1,...,8}, ytick={2.5,5,7.5,10}, xlabel={$t$（s）}, ylabel={$v$（m/s）}]
  \addplot[bkcurve, mark=*, mark size=1.4pt] coordinates {(0,0) (1,2.5) (2,5) (3,7.5) (4,10) (5,10) (6,10) (7,7) (8,4)};
  \node[font=\sffamily\scriptsize, text=bkOrange!80!black, anchor=north west] at (axis cs:2.3,4.6) {傾き $2.5$};
  \node[font=\sffamily\scriptsize, text=bkOrange!80!black, anchor=south] at (axis cs:5,10.7) {傾き $0$};
  \node[font=\sffamily\scriptsize, text=bkOrange!80!black, anchor=south west] at (axis cs:7.2,8) {傾き $-3$};
\end{axis}
\end{tikzpicture}
\bkfigcaption{やってみようの車の $v$-$t$ グラフ}
\end{minipage}

\begin{omoidasou}
\bkrecall{5-4}{$(x^n)'=nx^{n-1}$、定数を微分すると0。定数倍・和は1つずつ微分する。}
\bkrecall{5-1}{位置が小さくなる向きに動くとき、$\dfrac{\Delta x}{\Delta t}$ はマイナス。向きをふくめた速さが速度。}
\bkrecall{2-2}{マイナスは「反対の向き」を表す。$+$ の向きは、はじめに決めて変えない。}
\bkrecall{2-10}{換算の数も単位つきの量（\qty{1000}{m/km}、\qty{3600}{s/h}）。\\単位も計算できる。$(\mathrm{m/s})\div\mathrm{s}=\mathrm{m/s^2}$}
\bkrecall{2-3}{速さが一定なら、$v$-$t$ グラフの長方形の面積が道のり。}
\end{omoidasou}

%% ---------------- 数学の道具 ----------------
\bksection{doguu}{数学の道具}

\bksubsection{1}{速度は、位置を微分したもの}

\leavevmode\label{bk:5-5:sokudo}まっすぐな道を走る車の位置を、数直線の上の $x$（m）で表します。はじめに、右を $+$ の向きと決めます（2-2）。

\begin{bkfigure}
\begin{tikzpicture}[x=1mm, y=1mm]
  % 上：右へ動く車
  \draw[bkInk, line width=0.7pt, -{Stealth[length=3mm]}] (0,22) -- (122,22);
  \foreach \k/\l in {10/0,40/10,70/20,100/30} {\draw[bkInk, line width=0.6pt] (\k,20.8) -- (\k,23.2); \node[font=\sffamily\scriptsize, anchor=north] at (\k,20.3) {\l};}
  \node[font=\sffamily\scriptsize, anchor=north] at (119,20.3) {$x$（m）};
  \bkcar{(34,24.5)}{0.55}{bkBlue}
  \draw[bk arrow, bkOrange] (52,30) -- (68,30);
  \node[font=\sffamily\footnotesize, anchor=west, align=left] at (72,31) {右へ動く：$\Delta x>0$、速度は $+$};
  % 下：左へ動く車
  \draw[bkInk, line width=0.7pt, -{Stealth[length=3mm]}] (0,0) -- (122,0);
  \foreach \k/\l in {10/0,40/10,70/20,100/30} {\draw[bkInk, line width=0.6pt] (\k,-1.2) -- (\k,1.2); \node[font=\sffamily\scriptsize, anchor=north] at (\k,-1.7) {\l};}
  \node[font=\sffamily\scriptsize, anchor=north] at (119,-1.7) {$x$（m）};
  \begin{scope}[xscale=-1]\bkcar{(-88,2.5)}{0.55}{bkGreen}\end{scope}
  \draw[bk arrow, bkOrange] (70,8) -- (54,8);
  \node[font=\sffamily\footnotesize, anchor=east, align=right] at (50,9) {左へ動く：$\Delta x<0$、速度は $-$};
\end{tikzpicture}
\bkfigcaption{$+$ の向き（右）に動くと $\Delta x$ は正、反対の向きに動くと $\Delta x$ は負。}
\end{bkfigure}

$x$ がふえる向きに動けば、$\Delta x$ も $\dfrac{\Delta x}{\Delta t}$ も正です。$x$ がへる向きに動けば、$\Delta x$ も $\dfrac{\Delta x}{\Delta t}$ も負になります。この向きもふくめた瞬間の速さを\textbf{速度}といい、$v$ で表します。5-3で学んだように、瞬間の値は、位置 $x$ を時間 $t$ で微分すると求まります。\bkidx{そくど}{速度}
\[ v=\frac{\dd x}{\dd t} \]
速度から向き（$+$ か $-$ か）を取った大きさが、ふだん「速さ」とよんでいるものです。

\begin{bkcenter}
\begin{bktable}{colspec={Q[c,m,wd=26mm] Q[c,m,wd=26mm] Q[c,m,wd=26mm]}, rowsep=1.5pt}
速度 & 向き & 速さ \\
\qty{10}{m/s} & $+$ の向き（右） & \qty{10}{m/s} \\
$\qty{-10}{m/s}$ & $-$ の向き（左） & \qty{10}{m/s} \\
\qty{0}{m/s} & 止まっている & \qty{0}{m/s} \\
\end{bktable}
\bktabcaption{速度が $\qty{10}{m/s}$ と $\qty{-10}{m/s}$ は、速さは同じで向きが反対。}
\end{bkcenter}

\noindent\textbf{例：真上に投げたボール}\par
\noindent ボールを真上に、秒速\qty{20}{m}で投げ上げます。上向きを $+$ にすると、$t$ 秒後の高さは
\[ x=20t-5t^2\quad\text{（m）} \]
です。なぜこの式になるのでしょう。もし重力がなければ、ボールは秒速\qty{20}{m}のまま上がり続けて、$t$ 秒で $20t$（m）上がります。でも、実際には重力で引っぱられ、手をはなした物が落ちるのと同じだけ下がります。5-2のように、$t$ 秒で落ちる距離は $5t^2$（m）でした（ほんとうは $4.9t^2$ ですが、ここでも計算しやすいように5にします）。だから、高さは $20t$ から $5t^2$ を引いた $20t-5t^2$ になります。

微分すると（5-4）、
\[ v=\frac{\dd x}{\dd t}=20-5\times2t=20-10t \]

\noindent\begin{minipage}[c]{0.58\linewidth}
\begin{bktable}{colspec={Q[c,m,wd=17mm] *{5}{Q[c,m,wd=8.6mm]}}, rowsep=1.5pt, colsep=1.2mm, cells={font=\small}}
$t$（s） & 0 & 1 & 2 & 3 & 4 \\
$x$（m） & 0 & 15 & 20 & 15 & 0 \\
$v$（m/s） & 20 & 10 & 0 & $-10$ & $-20$ \\
\end{bktable}

\vspace{4pt}
1秒後の速度は $20-10=10$ で、上向きに\qty{10}{m/s}。3秒後は $20-30=-10$ で、下向きに\qty{10}{m/s}。2秒後は $v=0$ で、ボールは一瞬止まり、いちばん高い所（\qty{20}{m}）にいます。1秒後と3秒後は、高さも速さも同じですが、向きが反対です。速度を使うと、この向きのちがいが $+$ と $-$ で表せます。
\end{minipage}\hfill
\begin{minipage}[c]{0.38\linewidth}\centering
\begin{tikzpicture}[x=1mm, y=1.6mm]
  \draw[bkInk, line width=0.6pt] (0,0) -- (0,22);
  \foreach \h in {0,5,10,15,20} {\draw[bkInk, line width=0.5pt] (-1,\h) -- (1,\h); \node[font=\sffamily\scriptsize, anchor=east] at (-1.5,\h) {\h};}
  \node[font=\sffamily\scriptsize, anchor=south] at (0,22.5) {$x$（m）};
  % 上がるとき（左の列）と、下りるとき（右の列）
  \foreach \t/\h/\v in {0/0/20,1/15/10} {\shade[ball color=bkOrange] (8,\h) circle[radius=1.6mm]; \draw[bk arrow, bkBlue] (8,{\h+1.2}) -- (8,{\h+1.2+\v/3});}
  \shade[ball color=bkOrange] (14,20) circle[radius=1.6mm];
  \foreach \t/\h/\v in {3/15/10,4/0/20} {\shade[ball color=bkOrange] (20,\h) circle[radius=1.6mm]; \draw[bk arrow, bkRed] (20,{\h-1.2}) -- (20,{\h-1.2-\v/3});}
  \node[font=\sffamily\scriptsize, anchor=west] at (9.8,0) {0秒};
  \node[font=\sffamily\scriptsize, anchor=west] at (9.8,15) {1秒};
  \node[font=\sffamily\scriptsize, anchor=south] at (14,21.2) {2秒};
  \node[font=\sffamily\scriptsize, anchor=west] at (21.8,15) {3秒};
  \node[font=\sffamily\scriptsize, anchor=west] at (21.8,1.2) {4秒};
\end{tikzpicture}
\end{minipage}

\vspace{2pt}
\noindent{\small 右の図：青の矢印は上向きの速度、赤の矢印は下向きの速度。矢印が長いほど速い。}

\bksubsection{2}{加速度は、速度を微分したもの}

\leavevmode\label{bk:5-5:kasokudo}やってみようの車では、走り出したときに、1秒ごとに速度が\qty{2.5}{m/s}ずつふえていました。このように、速度が変わる速さを\textbf{加速度}（かそくど）といい、$a$ で表します。\bkidx{かそくど}{加速度}

\noindent\textbf{平均の加速度}\par
\noindent 時間 $\Delta t$ のあいだに速度が $\Delta v$ 変わったとき、$\dfrac{\Delta v}{\Delta t}$ を\textbf{平均の加速度}といいます。平均の速さ $\dfrac{\Delta x}{\Delta t}$（5-1）の $x$ を $v$ にかえただけです。やってみようの車の0秒から4秒までなら、
\[ \frac{\Delta v}{\Delta t}=\frac{\qty{10}{m/s}-\qty{0}{m/s}}{\qty{4}{s}}=\frac{\qty{10}{m/s}}{\qty{4}{s}}=\qty{2.5}{m/s^2} \]

\noindent\textbf{単位 \unit{m/s^2} の意味}\par
\noindent 速度の単位 m/s を、時間の単位 s でわるので、$(\mathrm{m/s})\div\mathrm{s}=\mathrm{m/s^2}$ です（2-10）。「メートル毎秒毎秒」と読みます。\qty{2.5}{m/s^2} は、「\textbf{1秒あたり、速度が\qty{2.5}{m/s}ずつ変わる}」という意味です。s が2回出てくるのは、1回めが速度の「毎秒」、2回めが「1秒ごとに変わる」の「毎秒」だからです。

\noindent\textbf{瞬間の加速度}\par
\noindent 平均の速さで $\Delta t$ を0に近づけると、瞬間の速度になりました（5-2・5-3）。同じように、平均の加速度で $\Delta t$ を0に近づけると、その瞬間の加速度になります。つまり、加速度は速度を時間で微分したものです。
\[ a=\frac{\dd v}{\dd t} \]
グラフで見ると、$x$-$t$ グラフの接線の傾きが速度 $v$、$v$-$t$ グラフの接線の傾きが加速度 $a$ です。やってみようの車の加速度は、走り出したときが \qty{2.5}{m/s^2}、一定の速さのときが \qty{0}{m/s^2}、ブレーキのときが $\qty{-3}{m/s^2}$ です。

\noindent\textbf{加速度のマイナス}\par
\noindent ブレーキのときの $\qty{-3}{m/s^2}$ は、速度が1秒ごとに\qty{3}{m/s}ずつへることを表しています。$+$ の向きに走っている車の加速度が $-$ なら、車はおそくなります（まちがえやすいところも見よう）。

\noindent\textbf{身近な乗り物の加速度}\par
\noindent 0から時速\qty{100}{km}まで10秒で加速する車を考えます。まず、速度を m/s になおします。
\[ \qty{100}{km/h}\times\qty{1000}{m/km}\div\qty{3600}{s/h}\approx\qty{27.8}{m/s} \]
（$100000\div3600=27.77\cdots$ を、小数第1位までのがい数にした。$\approx$ は「ほぼ等しい」を表す記号で、5-8 でくわしく学びます）
平均の加速度は、$\dfrac{\qty{27.8}{m/s}}{\qty{10}{s}}=\qty{2.78}{m/s^2}$ です。

\bksubsection{3}{はてな？の答え：体は加速度を感じている}

\leavevmode\label{bk:5-5:hatena}\bkhatenaref のエレベーターの速度を、グラフにしてみます。上向きを $+$ とします。このエレベーターは、2秒かけて\qty{2}{m/s}まで速くなり、6秒間その速さで上がり、2秒かけて止まります。

\begin{bkfigure}
\begin{tikzpicture}
\begin{axis}[bkplot, name=e1, width=96mm, height=34mm, xmin=0, xmax=10.8, ymin=0, ymax=2.9, xtick={2,4,6,8,10}, ytick={1,2}, xlabel={$t$（s）}, ylabel={$v$（m/s）}]
  \addplot[bkcurve] coordinates {(0,0) (2,2) (8,2) (10,0)};
  \node[font=\sffamily\scriptsize, text=bkOrange!80!black, anchor=north west] at (axis cs:1.15,0.95) {傾き1};
  \node[font=\sffamily\scriptsize, text=bkOrange!80!black, anchor=south] at (axis cs:5,2.1) {傾き0};
  \node[font=\sffamily\scriptsize, text=bkOrange!80!black, anchor=north east] at (axis cs:8.85,0.95) {傾き$-1$};
\end{axis}
\begin{axis}[bkplot, at={($(e1.south)+(0,-11mm)$)}, anchor=north, width=96mm, height=34mm, xmin=0, xmax=10.8, ymin=-1.7, ymax=1.7, xtick={2,4,6,8,10}, xticklabels={}, ytick={-1,1}, xlabel={$t$（s）}, ylabel={$a$（\unit{m/s^2}）}]
  \addplot[bkcurve3, line width=1.4pt] coordinates {(0,1) (2,1)};
  \addplot[bkcurve3, line width=1.4pt] coordinates {(2,0) (8,0)};
  \addplot[bkcurve3, line width=1.4pt] coordinates {(8,-1) (10,-1)};
  \node[font=\sffamily\scriptsize, anchor=west, inner sep=1pt] at (axis cs:2.3,1) {重く感じる};
  \node[font=\sffamily\scriptsize, anchor=south, inner sep=3pt] at (axis cs:5,0) {ふつう};
  \node[font=\sffamily\scriptsize, anchor=east, inner sep=1pt] at (axis cs:7.7,-1) {軽く感じる};
\end{axis}
\end{tikzpicture}
\bkfigcaption{上：エレベーターの $v$-$t$ グラフ。下：その傾き（加速度）のグラフ。}
\end{bkfigure}

速度のグラフ（上）は、0秒から8秒まで、ずっと $+$（上向き）です。でも、体の感じ方は3つに分かれます。感じ方が分かれる所は、下の加速度のグラフが分かれる所とぴったり同じです。
\begin{itemize}[leftmargin=1.2\zw, itemsep=1pt, topsep=2pt, label={\color{bkTheme}●}]
\item 動き出すとき（0〜2秒）：加速度は $+\qty{1}{m/s^2}$（上向き）。重く感じる。
\item 一定の速さのとき（2〜8秒）：加速度は0。止まっているときと同じ。
\item 止まるとき（8〜10秒）：加速度は $\qty{-1}{m/s^2}$（下向き）。軽く感じる。
\end{itemize}
体が感じていたのは、速度ではなく\textbf{加速度}だったのです。

\begin{motto}[3]
\textbf{体重計に乗って、エレベーターに乗ると}\par
エレベーターの中で体重計に乗ると、上へ動き出すときは目もりが大きく、止まるときは小さくなります。上向きに加速するには、床が足を、ふだんより強くおし上げなければならないからです。床から強くおされるのを、体は「重い」と感じます。くわしくは、力と運動を学ぶ第6巻で考えます。
\end{motto}

\bksubsection{4}{2回微分する}

\leavevmode\label{bk:5-5:nikai}加速度は、位置 $x$ を微分して速度 $v$ にし、その $v$ をもう一度微分したものです。これを、位置を\textbf{2回微分する}といいます。

\begin{bkfigure}
\begin{tikzpicture}[x=1mm, y=1mm, q/.style={bk node, minimum width=30mm, minimum height=12mm}]
  \node[q] (x) at (0,0) {位置 $x$\\（m）};
  \node[q] (v) at (50,0) {速度 $v$\\（m/s）};
  \node[q] (a) at (100,0) {加速度 $a$\\（\unit{m/s^2}）};
  \draw[bk link] (x) -- node[bk link label, above] {時間で微分} (v);
  \draw[bk link] (v) -- node[bk link label, above] {時間で微分} (a);
\end{tikzpicture}
\bkfigcaption{位置を1回微分すると速度、2回微分すると加速度。単位は、1回微分するごとに s でわった形になる。}
\end{bkfigure}

式で書くと、次のようになります。
\[ a=\frac{\dd v}{\dd t}=\frac{\dd}{\dd t}\Bigl(\frac{\dd x}{\dd t}\Bigr)=\frac{\dd^2x}{\dd t^2} \]
$\dfrac{\dd}{\dd t}$ は「時間で微分する」という命令の記号です。命令を2回つづけたので、上には $\dd$ が2つ、下には $\dd t$ が2つあります。これを短く書いたのが $\dfrac{\dd^2x}{\dd t^2}$ です。上の $\dd^2$ は「$\dd$ を2回」、下の $\dd t^2$ は「$\dd t$ を2回」という意味で、$t^2$ を微分するという意味ではありません。

\begin{bkcenter}
\begin{bktable}{colspec={Q[l,m,wd=26mm] Q[c,m,wd=18mm] Q[c,m,wd=18mm] Q[l,m,wd=58mm]}, rowsep=1.5pt, cells={font=\small}, row{1}={font=\sffamily\bfseries\small}}
書き方 & 1回微分 & 2回微分 & 読み方・使う分野 \\
数学 & $x'$ & $x''$ & ダッシュ、ツーダッシュ \\
ライプニッツ & $\dfrac{\dd x}{\dd t}$ & $\dfrac{\dd^2x}{\dd t^2}$ & ディーツー・エックス・ディーティーにじょう \\
ニュートン & $\dot{x}$ & $\ddot{x}$ & エックス・ドット、ツードット（物理で、時間の微分に使う） \\
\end{bktable}
\bktabcaption{1回微分は速度、2回微分は加速度。専門書では場面で使い分ける。}
\end{bkcenter}

ニュートンの書き方では、文字の上の点の数が、時間で微分した回数です。$\dot{x}$ は1回、$\ddot{x}$ は2回です。短く書けるので、物理の本でよく使われます。

\bkkigocards{\bkkigocard{$\dot{x}$、$\ddot{x}$}{エックス・ドット、エックス・ツードット}{時間で1回微分（速度）・2回微分（加速度）}{x-dot}{5-5}}
{\bkkigocard{$\dfrac{\dd^2x}{\dd t^2}$}{ディーツー・エックス・ディーティーにじょう}{$x$ を $t$ で2回微分したもの}{second derivative}{5-5}}

\noindent\textbf{例}\par
\noindent 位置が $x=t^3$（m）のとき、1回微分して $v=\dot{x}=3t^2$、もう1回微分して $a=\ddot{x}=3\times2t=6t$ です。2秒では、$v=3\times2^2=\qty{12}{m/s}$、$a=6\times2=\qty{12}{m/s^2}$ です。

\bksubsection{5}{落ちる物の加速度}

\leavevmode\label{bk:5-5:g}手をはなしたボールの、落ちた距離はほんとうは $x=4.9t^2$（m）です（下向きを $+$）。2回微分すると、
\[ v=\dot{x}=4.9\times2t=9.8t,\qquad a=\ddot{x}=9.8 \]
表で見ると、速度は1秒ごとに\qty{9.8}{m/s}ずつふえています。

\begin{bkcenter}
\begin{bktable}{colspec={Q[l,m,wd=30mm] *{4}{Q[c,m,wd=12mm]}}, rowsep=1.5pt}
$t$（s） & 0 & 1 & 2 & 3 \\
$x=4.9t^2$（m） & 0 & 4.9 & 19.6 & 44.1 \\
$v=9.8t$（m/s） & 0 & 9.8 & 19.6 & 29.4 \\
1秒ごとの $v$ の変化 & & 9.8 & 9.8 & 9.8 \\
\end{bktable}
\bktabcaption{落ちる物の速度は、1秒ごとに\qty{9.8}{m/s}ずつふえる。}
\end{bkcenter}

加速度は時間によらず、いつも \qty{9.8}{m/s^2} です。これを\textbf{重力加速度}といい、$g$ で表します。9.8という数は、計算で決めたものではなく、地球の上で、物を落としてはかった値です。月の上では、重力が弱いので、約\qty{1.6}{m/s^2}です。\bkidx{じゅうりょくかそくど}{重力加速度 $g$}

空気の抵抗がなければ、重い鉄の球も軽いプラスチックの球も、同じ加速度 $g$ で落ちます。空気のない月の上で、宇宙飛行士がハンマーと鳥の羽を同時に落とすと、2つは同時に地面に着きました（1971年、アポロ15号）。

真上に投げ上げたボール $x=v_0t-4.9t^2$（上向きを $+$。$v_0$ は投げたときの速度）でも、2回微分すると $a=-9.8$ です。上がっているときも、いちばん高い所でも、落ちているときも、加速度はいつも下向きに $g$ です。

\bksubsection{6}{グラフで見る：$x$・$v$・$a$}

\leavevmode\label{bk:5-5:graph}\noindent\begin{minipage}[c]{0.56\linewidth}
\begin{tikzpicture}
% 第6回：段の間を9mmにして、v・a の名前が上のグラフにかからないようにする。
% v のグラフの目もり数字は、線のない側に置く（1は下、2は左下、3・4は上）
\begin{axis}[bkplot, name=g1, width=66mm, height=28mm, xmin=0, xmax=4.4, ymin=0, ymax=24, xtick={1,2,3,4}, ytick={10,20}, xlabel={$t$}, ylabel={$x$}]
  \addplot[bkcurve, domain=0:4] {20*x-5*x^2};
\end{axis}
\begin{axis}[bkplot, at={($(g1.south)+(0,-9mm)$)}, anchor=north, name=g2, width=66mm, height=28mm, xmin=0, xmax=4.4, ymin=-22, ymax=24, xtick={1,2,3,4}, xticklabels={}, ytick={-20,20}, xlabel={$t$}, ylabel={$v$}]
  \addplot[bkcurve2, domain=0:4] {20-10*x};
  \node[anchor=north, font=\sffamily\footnotesize, text=bkInk, inner sep=2pt] at (axis cs:1,0) {$1$};
  \node[anchor=north east, font=\sffamily\footnotesize, text=bkInk, inner sep=1.5pt] at (axis cs:2,0) {$2$};
  \node[anchor=south, font=\sffamily\footnotesize, text=bkInk, inner sep=2.5pt] at (axis cs:3,0) {$3$};
  \node[anchor=south, font=\sffamily\footnotesize, text=bkInk, inner sep=2.5pt] at (axis cs:4,0) {$4$};
\end{axis}
\begin{axis}[bkplot, at={($(g2.south)+(0,-9mm)$)}, anchor=north, width=66mm, height=28mm, xmin=0, xmax=4.4, ymin=-14, ymax=6, xtick={1,2,3,4}, ytick={-10}, xlabel={$t$}, ylabel={$a$}]
  \addplot[bkcurve3, domain=0:4] {-10};
\end{axis}
\end{tikzpicture}
\end{minipage}\hfill
\begin{minipage}[c]{0.4\linewidth}
1で使った、投げ上げたボールの位置・速度・加速度のグラフです。
\[ x=20t-5t^2 \]
\begin{itemize}[leftmargin=1.2\zw, itemsep=1pt, topsep=2pt, label={\color{bkTheme}●}]
\item $x$-$t$ の接線の傾きが $v$。
\item $v$-$t$ の傾きが $a$。
\item 最も高い所（$t=2$）で $v=0$。でも $a$ は0ではなく、$-10$。
\end{itemize}
\end{minipage}

\vspace{2pt}
上のグラフ（$x$）は、山の形です。山をのぼっている0〜2秒は傾きが正なので、まん中のグラフ（$v$）は $+$ です。山のてっぺんの2秒で傾きは0、山を下りる2〜4秒は傾きが負なので、$v$ は $-$ です。まん中の $v$ のグラフは、どこでも同じ傾き $-10$ の直線です。だから、下の $a$ のグラフは、いつも $-10$ の横の直線になります。

\bksubsection{7}{加速度が一定の運動の式}

\leavevmode\label{bk:5-5:toukasoku}落ちるボールの式 $x=4.9t^2$ の4.9は、どこから来たのでしょう。また、加速度がいつも同じ値 $a$ の運動は、どんな式で表せるでしょう。ここでは、「2回微分すると $a$ になる式」を、微分を逆にたどって作ってみます。

\noindent\textbf{①\enspace 2回微分すると $a$ になる式}\par
\noindent $t^2$ を2回微分すると、$t^2\to2t\to2$ で、2になります。2ではなく $a$ にしたいので、$t^2$ に $\dfrac{a}{2}$ をかけます。
\[ \frac{1}{2}at^2\ \xrightarrow{\text{微分}}\ \frac{1}{2}a\times2t=at\ \xrightarrow{\text{微分}}\ a \]
$\dfrac{1}{2}$ は、$t^2$ を微分したときに前に出てくる2を、打ち消すためについています。

\noindent\textbf{②\enspace はじめの速度 $v_0$ を入れる}\par
\noindent ①の式の速度は $at$ なので、$t=0$ で速度が0です。はじめに速度 $v_0$ で動いていたときは、$v_0t$ をたします。$v_0t$ は、1回微分すると $v_0$、2回微分すると0なので、加速度は $a$ のまま変わりません。

\noindent\textbf{③\enspace はじめの位置 $x_0$ を入れる}\par
\noindent 定数 $x_0$ をたしても、微分すると0なので、速度も加速度も変わりません。$t=0$ での位置が $x_0$ になります。

まとめると、加速度が一定の値 $a$ の運動は、
\[ x=x_0+v_0t+\frac{1}{2}at^2,\qquad v=v_0+at \]
と表せます（この式の別の求め方は5-13）。$t=0$ を入れると、$x=x_0$、$v=v_0$ で、たしかに「はじめの位置」と「はじめの速度」になっています。

落ちるボールは、$x_0=0$、$v_0=0$、$a=9.8$ の場合です。$\dfrac{1}{2}\times9.8=4.9$ なので、$x=4.9t^2$ になります。\textbf{4.9は、重力加速度9.8の半分}だったのです。

\noindent\textbf{例：ブレーキで止まる車}\par
\noindent 秒速\qty{20}{m}で走る車が、加速度 $\qty{-4}{m/s^2}$ でブレーキをかけます。$x_0=0$、$v_0=20$、$a=-4$ なので、
\[ v=20-4t,\qquad x=20t+\frac{1}{2}\times(-4)\times t^2=20t-2t^2 \]

\begin{bkcenter}
\begin{bktable}{colspec={Q[c,m,wd=18mm] *{6}{Q[c,m,wd=10mm]}}, rowsep=1.5pt, cells={font=\small}}
$t$（s） & 0 & 1 & 2 & 3 & 4 & 5 \\
$v$（m/s） & 20 & 16 & 12 & 8 & 4 & 0 \\
$x$（m） & 0 & 18 & 32 & 42 & 48 & 50 \\
\end{bktable}
\bktabcaption{ブレーキをかけた車の速度と位置。速度は1秒ごとに\qty{4}{m/s}ずつへる。}
\end{bkcenter}

$v=0$ になるのは、$20-4t=0$ より $4t=20$、$t=5$ 秒後です。そのあいだに進む距離は、$x=20\times5-2\times5^2=100-50=\qty{50}{m}$ です。

\vspace{4pt}
\noindent\begin{minipage}[c]{0.5\linewidth}
\textbf{グラフでも確かめる}\par
速さが一定なら、$v$-$t$ グラフの長方形の面積が道のりでした（2-3）。速度がへっていくときも、グラフの下の面積が道のりになります（くわしくは5-9・5-13）。右の三角形の面積は、
\[ \frac{1}{2}\times\qty{5}{s}\times\qty{20}{m/s}=\qty{50}{m} \]
で、式の答えと同じです。
\end{minipage}\hfill
\begin{minipage}[c]{0.46\linewidth}\centering
\begin{tikzpicture}
\begin{axis}[bkplot, width=58mm, height=38mm, xmin=0, xmax=5.8, ymin=0, ymax=24, xtick={1,...,5}, ytick={10,20}, xlabel={$t$（s）}, ylabel={$v$（m/s）}]
  \fill[bkYellow!45] (axis cs:0,0) -- (axis cs:0,20) -- (axis cs:5,0) -- cycle;
  \addplot[bkcurve, domain=0:5] {20-4*x};
  \node[font=\sffamily\scriptsize, anchor=center] at (axis cs:1.4,5.5) {面積50};
\end{axis}
\end{tikzpicture}
\end{minipage}

\begin{motto}[3]
\textbf{3回微分すると}\par
加速度をもう一度微分した $\dddot{x}$ を、加加速度（ジャーク）といいます。加速度がどれだけ急に変わるかを表し、乗り物の乗り心地を設計するときに使います。はてな？のエレベーターのグラフでは、加速度が2秒と8秒でいきなり変わっています。実際のエレベーターは、加速度を少しずつ変えて、がくんとしないようにしています。
\end{motto}

\begin{taisetsu}
速度 $v=\dfrac{\dd x}{\dd t}=\dot{x}$（m/s）\qquad 加速度 $a=\dfrac{\dd v}{\dd t}=\dfrac{\dd^2x}{\dd t^2}=\ddot{x}$（\unit{m/s^2}）\\
位置を1回微分すると速度、2回微分すると加速度\\
加速度が一定なら\quad $x=x_0+v_0t+\dfrac{1}{2}at^2$、\enspace$v=v_0+at$
\end{taisetsu}

\begin{kotoba}
\small
\bkword{速度}{向きもふくめた速さ。位置を時間で微分したもの。}
\bkword{加速度}{速度の変わる速さ。速度を時間で微分したもの。単位は \unit{m/s^2}。}
\bkword{重力加速度}{落ちる物の加速度。地球の上で $g=\qty{9.8}{m/s^2}$。}
\end{kotoba}

%% ---------------- 例題 ----------------
\bksection{reidai}{例題でたしかめよう}

\begin{reidai}[速度と加速度を求める]
位置が $x=3t^2+2t$（m）の物体の、速度 $v$ と加速度 $a$ を求めましょう。\par
\kangaekata $x$ を $t$ で1回微分すると $v$、その $v$ をもう1回微分すると $a$ です。\par
\shiki $v=\dfrac{\dd x}{\dd t}=3\times2t+2\times1=6t+2$\par
\hspace*{2.6\zw}$a=\dfrac{\dd v}{\dd t}=6\times1+0=6$\par
\kotae $v=6t+2$（m/s）、$a=6$（\unit{m/s^2}）
\end{reidai}

\begin{jibun}
位置が $x=2t^3$（m）の物体の、$v$ と $a$ を求めましょう。また、1秒での $v$ と $a$ の値を答えましょう。
\bkjans{$v=6t^2$、$a=12t$。1秒で $v=\qty{6}{m/s}$、$a=\qty{12}{m/s^2}$}
\end{jibun}

\begin{reidai}[平均の加速度]
止まっていた車が、8秒で時速\qty{72}{km}になりました。平均の加速度を求めましょう。\par
\kangaekata 加速度の単位は \unit{m/s^2} なので、速度を m/s になおしてから、$\dfrac{\Delta v}{\Delta t}$ を計算します。\par
\shiki $\qty{72}{km/h}\times\qty{1000}{m/km}\div\qty{3600}{s/h}=\qty{20}{m/s}$\par
\hspace*{2.6\zw}$\Delta v=\qty{20}{m/s}-\qty{0}{m/s}=\qty{20}{m/s}$、$\Delta t=\qty{8}{s}$\par
\hspace*{2.6\zw}$\dfrac{\Delta v}{\Delta t}=\dfrac{\qty{20}{m/s}}{\qty{8}{s}}=\qty{2.5}{m/s^2}$\par
\kotae \qty{2.5}{m/s^2}
\end{reidai}

\begin{jibun}
止まっていた電車が、30秒で時速\qty{54}{km}になりました。平均の加速度を求めましょう。
\bkjans{$\qty{54}{km/h}\times\qty{1000}{m/km}\div\qty{3600}{s/h}=\qty{15}{m/s}$。$15\div30=\qty{0.5}{m/s^2}$}
\end{jibun}

\begin{reidai}[ある時刻の値]
位置が $x=t^3-6t^2+9t$（m）の物体の、2秒での速度と加速度を求めましょう。\par
\kangaekata まず $v$ と $a$ を $t$ の式で求めてから、$t=2$ を入れます。\par
\shiki $v=3t^2-6\times2t+9=3t^2-12t+9$\par
\hspace*{2.6\zw}$a=3\times2t-12=6t-12$\par
\hspace*{2.6\zw}$t=2$ で $v=3\times2^2-12\times2+9=12-24+9=-3$\par
\hspace*{2.6\zw}$a=6\times2-12=0$\par
\kotae 速度 $\qty{-3}{m/s}$、加速度 \qty{0}{m/s^2}
\end{reidai}

\begin{jibun}
位置が $x=t^3-3t$（m）の物体の、2秒での速度と加速度を求めましょう。
\bkjans{$v=3t^2-3$、$a=6t$。速度 \qty{9}{m/s}、加速度 \qty{12}{m/s^2}}
\end{jibun}

\begin{reidai}[投げ上げ]
真上に投げたボールの高さが $x=15t-5t^2$（m、上向きが $+$）のとき、（1）1秒と2秒の速度、（2）速度が0になる時刻、（3）加速度を求めましょう。\par
\kangaekata 速度の符号で向きがわかります。$+$ なら上向き、$-$ なら下向きです。\par
\shiki $v=15-5\times2t=15-10t$、\enspace$a=-10$\par
\hspace*{2.6\zw}（1）$t=1$ で $v=15-10=5$、$t=2$ で $v=15-20=-5$\par
\hspace*{2.6\zw}（2）$15-10t=0$ より $10t=15$、$t=1.5$\par
\hspace*{2.6\zw}（3）$a=-10$（時刻によらない）\par
\kotae （1）上向きに\qty{5}{m/s}、下向きに\qty{5}{m/s}\quad（2）1.5秒（最も高い所）\quad（3）$\qty{-10}{m/s^2}$
\end{reidai}

\begin{jibun}
高さが $x=25t-5t^2$（m）のボールについて、1秒と4秒の速度、速度が0になる時刻、加速度を求めましょう。
\bkjans{\qty{15}{m/s}、$\qty{-15}{m/s}$。2.5秒。$\qty{-10}{m/s^2}$}
\end{jibun}

\begin{reidai}[加速度が一定の運動の式をつくる]
はじめの位置が\qty{2}{m}、はじめの速度が\qty{3}{m/s}、加速度が \qty{4}{m/s^2} で一定の物体があります。位置 $x$ と速度 $v$ を $t$ の式で表し、2秒での位置と速度を求めましょう。\par
\kangaekata $x=x_0+v_0t+\dfrac{1}{2}at^2$、$v=v_0+at$ に、$x_0=2$、$v_0=3$、$a=4$ を入れます。\par
\shiki $x=2+3t+\dfrac{1}{2}\times4\times t^2=2+3t+2t^2$\par
\hspace*{2.6\zw}$v=3+4t$\par
\hspace*{2.6\zw}$t=2$ で $x=2+6+8=16$、$v=3+8=11$\par
\hspace*{2.6\zw}確かめ：$x$ を微分すると $3+4t$ で、$v$ と同じ\par
\kotae $x=2t^2+3t+2$、$v=4t+3$。2秒で\qty{16}{m}、\qty{11}{m/s}
\end{reidai}

\begin{jibun}
止まっていた物体が、はじめの位置0から、加速度 \qty{6}{m/s^2} で一定のまま動きはじめました。$x$ と $v$ を $t$ の式で表し、3秒での位置と速度を求めましょう。
\bkjans{$x=3t^2$、$v=6t$。3秒で\qty{27}{m}、\qty{18}{m/s}}
\end{jibun}

\begin{reidai}[ブレーキで止まるまで]
秒速\qty{15}{m}で走る自転車が、加速度 $\qty{-3}{m/s^2}$ でブレーキをかけました。止まるまでの時間と、そのあいだに進む距離を求めましょう。\par
\kangaekata 加速度が一定なので、$x=v_0t+\dfrac{1}{2}at^2$、$v=v_0+at$ を使います（$x_0=0$）。止まるのは $v=0$ のときです。\par
\shiki $v=15+(-3)\times t=15-3t$。$15-3t=0$ より $t=5$\par
\hspace*{2.6\zw}$x=15t+\dfrac{1}{2}\times(-3)\times t^2=15t-1.5t^2$\par
\hspace*{2.6\zw}$t=5$ を入れて $x=15\times5-1.5\times25=75-37.5=37.5$\par
\kotae 5秒、\qty{37.5}{m}
\end{reidai}

\begin{jibun}
秒速\qty{20}{m}で走る車が、加速度 $\qty{-5}{m/s^2}$ でブレーキをかけました。止まるまでの時間と距離を求めましょう。
\bkjans{$20-5t=0$ より4秒。$20\times4-2.5\times16=\qty{40}{m}$}
\end{jibun}

\noindent\textbf{速さが2倍になると、止まるまでの距離は？}\par
\noindent 加速度 $\qty{-5}{m/s^2}$ でブレーキをかけたとき、はじめの速さと、止まるまでの距離をくらべます（じぶんでとこう6と同じ計算）。
\begin{bkcenter}
\begin{bktable}{colspec={Q[l,m,wd=40mm] *{3}{Q[c,m,wd=16mm]}}, rowsep=1.5pt}
はじめの速さ（m/s） & 10 & 20 & 30 \\
止まるまでの時間（s） & 2 & 4 & 6 \\
止まるまでの距離（m） & 10 & 40 & 90 \\
\end{bktable}
\bktabcaption{速さが2倍になると距離は4倍、3倍になると距離は9倍。}
\end{bkcenter}
速さが2倍になると、止まるまでの時間が2倍になり、そのあいだの平均の速さも2倍になります。だから、距離は $2\times2=4$ 倍です。車のスピードを出しすぎると危ないのは、このためです。

\begin{reidai}[速度が0になる時刻が2つあるとき]
位置が $x=t^3-6t^2+9t$（m）の物体の速度が0になる時刻を求め、それぞれの時刻の加速度を答えましょう。\par
\kangaekata $v=0$ は $t$ の2次方程式になります。因数分解で解きます（3-8・3-10）。\par
\shiki $v=3t^2-12t+9=3(t^2-4t+3)=3(t-1)(t-3)$\par
\hspace*{2.6\zw}$v=0$ より $t-1=0$ または $t-3=0$。$t=1,\ 3$\par
\hspace*{2.6\zw}$a=6t-12$ より、$t=1$ で $a=6-12=-6$、$t=3$ で $a=18-12=6$\par
\kotae 1秒（加速度 $\qty{-6}{m/s^2}$）と3秒（加速度 \qty{6}{m/s^2}）
\end{reidai}

\begin{jibun}
位置が $x=t^3-3t^2$（m）の物体の速度が0になる時刻と、そのときの加速度を求めましょう。
\bkjans{$v=3t^2-6t=3t(t-2)$ より0秒と2秒。$a=6t-6$ より $\qty{-6}{m/s^2}$ と \qty{6}{m/s^2}}
\end{jibun}

\begin{reidai}[落ちる物の速さ]
落ちる物の位置は $x=4.9t^2$（m、下向きが $+$）です。速さが\qty{19.6}{m/s}になるのは何秒後ですか。そのとき、何m落ちていますか。\par
\kangaekata 速さ $v$ を $t$ の式で表し、$v=19.6$ となる $t$ を求めます。\par
\shiki $v=4.9\times2t=9.8t$。$9.8t=19.6$ より $t=19.6\div9.8=2$\par
\hspace*{2.6\zw}$x=4.9\times2^2=4.9\times4=19.6$\par
\kotae 2秒後、\qty{19.6}{m}（加速度はいつも $g=\qty{9.8}{m/s^2}$）
\end{reidai}

\begin{jibun}
落ちる物（$x=4.9t^2$）の速さが\qty{29.4}{m/s}になるのは何秒後ですか。そのとき、何m落ちていますか。
\bkjans{$9.8t=29.4$ より3秒後。$4.9\times9=\qty{44.1}{m}$}
\end{jibun}

%% ---------------- まちがえやすいところ・物理工学 ----------------
\bksection{tsumazuki}{まちがえやすいところをチェック}

\begin{tsumazuki}{速度が0なら、加速度も0？}
投げ上げたボールは、最も高い所で一瞬止まります（$v=0$）。でも、加速度は $-10$ のままです。もし加速度も0なら、速度は0のまま変わらないので、ボールは空中に止まったままになるはずです。速度が0でも、速度が変わり続けていれば、加速度は0ではありません。
\end{tsumazuki}

\begin{tsumazuki}{加速度がマイナスなら、おそくなる？}
加速度の符号は、速度が変わる「向き」です。速度も加速度もマイナスなら、マイナスの向きにどんどん速くなります。投げ上げたボールが落ちてくるとき（$v<0$、$a<0$）は、速くなっていきます。速くなるかおそくなるかは、$v$ と $a$ の符号が同じ（速くなる）か、ちがう（おそくなる）かで決まります。
\end{tsumazuki}

\begin{tsumazuki}{時速を秒速になおすとき、かける・わるを逆にする}
\bkng{$72\times3.6=259.2$ で、\qty{259.2}{m/s}}\bkok{$\qty{72}{km/h}\times\qty{1000}{m/km}\div\qty{3600}{s/h}=\qty{20}{m/s}$}
単位をつけて計算すると、km と h が約分されて m/s が残るかどうかで、まちがいに気づけます。秒速\qty{259}{m}は、新幹線よりずっと速く、車としておかしいことにも気づけます。
\end{tsumazuki}

\begin{butsuri}
\begin{itemize}[leftmargin=1.2\zw, itemsep=2pt, topsep=0pt, label={\color{bkPurple}●}]
\item \textbf{加速度センサー}\\
スマートフォンには、加速度をはかる小さな部品が入っています。画面の向きが変わるのは、重力が下向きに引っぱる向きを、この部品が感じとっているからです。車のエアバッグは、ぶつかったときの大きな加速度を感じて開きます。
\item \textbf{力と加速度}\\
物にはたらく力の合計は、質量×加速度に等しい（$F=ma$）。これがニュートンの運動の法則です。加速度は位置の2回微分なので、運動の法則は「2回微分をふくむ式」になります。これが第6巻の微分方程式です（6-15）。
\item \textbf{乗り心地}\\
加速度がいきなり変わると、乗っている人はがくんと感じます。エレベーターや電車は、加速度をゆっくり変えるように動かしています。
\end{itemize}
\end{butsuri}

%% ---------------- 章末問題 ----------------
\bkendproblems
\bkproblemgroup{基本}{3}

\bkproblem 位置が $x=4t^2-3t+2$（m）の物体の、$v$ と $a$ を求めましょう。
\bkpans{$v=8t-3$（m/s）、$a=8$（\unit{m/s^2}）}

\bkproblem 位置が $x=2t^3-3t^2$（m）の物体の、2秒での速度と加速度を求めましょう。
\bkpans{$v=6t^2-6t$ より \qty{12}{m/s}、$a=12t-6$ より \qty{18}{m/s^2}}

\bkproblem 真上に投げたボールの高さが $x=30t-5t^2$（m）のとき、（1）1秒と4秒の速度、（2）最も高くなる時刻、（3）加速度を求めましょう。
\bkpans{（1）\nobreak \qty{20}{m/s}、$\qty{-10}{m/s}$\quad（2）\nobreak 3秒\quad（3）\nobreak $\qty{-10}{m/s^2}$}

\bkproblem 止まっていた自転車が、5秒で時速\qty{36}{km}になりました。平均の加速度を求めましょう。
\bkpans{$\qty{36}{km/h}\times\qty{1000}{m/km}\div\qty{3600}{s/h}=\qty{10}{m/s}$、$10\div5=\qty{2}{m/s^2}$}

\bkproblem 時速\qty{72}{km}で走る車が、加速度 $\qty{-5}{m/s^2}$ でブレーキをかけました。止まるまでの時間と距離を求めましょう。
\bkpans{\qty{20}{m/s} なので、$20-5t=0$ より4秒、$20\times4-2.5\times16=\qty{40}{m}$}

\bkproblem 位置が $x=t^3-12t$（m）の物体について、速度が0になる時刻（$t>0$）と、そのときの加速度を求めましょう。
\bkpans{$v=3t^2-12=3(t-2)(t+2)$ より2秒。$a=6t$ より \qty{12}{m/s^2}}

\bkproblem 落ちる物（$x=4.9t^2$）の、3秒後の速さと加速度を求めましょう。
\bkpans{$v=9.8t$ より \qty{29.4}{m/s}。加速度 \qty{9.8}{m/s^2}}

\bkproblemgroup{チャレンジ}{3}

\bkproblem 位置が $x=t^3-9t^2+24t$（m）の物体について、速度が0になる時刻を2つ求め、それぞれの時刻の加速度を答えましょう（2次方程式は3-10）。
\bkpans{$v=3t^2-18t+24=3(t-2)(t-4)$ より2秒と4秒。$a=6t-18$ より、2秒で $\qty{-6}{m/s^2}$、4秒で \qty{6}{m/s^2}}

\bkproblem 空気の抵抗がないとき、落ちる物は $x=4.9t^2$（m）です。速さが時速\qty{100}{km}になるのは、落ちはじめてから何秒後ですか。小数第1位まで求めましょう。
\bkpans{時速\qty{100}{km}は約\qty{27.8}{m/s}。$9.8t=27.8$ より、約2.8秒後}

\bkproblem 止まっていた電車が、加速度 \qty{0.8}{m/s^2} で動きはじめました。時速\qty{72}{km}になるまでの時間と、そのあいだに進む距離を求めましょう。
\bkpans{\qty{20}{m/s}。$v=0.8t=20$ より25秒。$x=\dfrac{1}{2}\times0.8\times25^2=\qty{250}{m}$}

\bkproblem エレベーターが、止まった状態から3秒かけて上向きに\qty{1.5}{m/s}になり、7秒間その速さで上がり、2秒かけて止まりました。（1）それぞれのあいだの加速度を求めましょう。（2）体が重く感じるのは、どのあいだですか（はてな？の答え）。
\bkpans{（1）\nobreak 0〜3秒：\qty{0.5}{m/s^2}、3〜10秒：\qty{0}{m/s^2}、10〜12秒：$\qty{-0.75}{m/s^2}$\quad（2）\nobreak 0〜3秒（動き出すとき）}

\begin{furikaeri}
\begin{checklist}
\item 速度と速さのちがいを説明できる。\hfill{\footnotesize\sffamily\color{bkTheme}→\pageref{bk:5-5:sokudo}ページ}
\item 位置を2回微分して、加速度を求められる。\hfill{\footnotesize\sffamily\color{bkTheme}→\pageref{bk:5-5:nikai}ページ}
\item $\dot{x}$、$\ddot{x}$、$\dfrac{\dd^2x}{\dd t^2}$ を読める。\hfill{\footnotesize\sffamily\color{bkTheme}→\pageref{bk:5-5:nikai}ページ}
\end{checklist}
\end{furikaeri}

\bkfukidashi{sora}{$x^3$ や $x^4$ は微分できるようになったけど、$(2x+1)^5$ みたいな式は、展開しないとだめなのかな？}

%% 5-6 微分の道具箱：積と合成関数（第11回：ページ制限なし。説明・図・例題を書き足した）
% 第7回：歯車  \bkgear{(中心)}{歯数 z}{モジュール m（mm）}{1つの歯の向き（度）}{ぬり}{線}
%   ピッチ円の半径 mz/2、歯先 +m、歯元 -1.25m。歯は台形で、外形を1本の線でかく。
%   かみ合う2つは、中心どうしを (m×(z1+z2)/2) はなし、中心線の上で「片方は歯、もう片方は歯みぞ」にする
\newcommand\bkgear[6]{\begin{scope}[shift={#1}]
  \pgfmathsetmacro\bkgR{#3*#2/2}\pgfmathsetmacro\bkgA{\bkgR+#3}%
  \pgfmathsetmacro\bkgF{\bkgR-1.25*#3}\pgfmathsetmacro\bkgP{360/#2}%
  \pgfmathsetmacro\bkgT{27*#3/\bkgA}\pgfmathsetmacro\bkgB{54*#3/\bkgF}%
  \filldraw[fill=#5, draw=#6, line width=0.5pt, line join=round]
    (#4-\bkgB:\bkgF)
    \foreach \k [evaluate=\k as \c using #4+\k*\bkgP] in {0,...,\numexpr#2-1\relax}{
      -- (\c-\bkgB:\bkgF) -- (\c-\bkgT:\bkgA) -- (\c+\bkgT:\bkgA) -- (\c+\bkgB:\bkgF)
      arc[start angle=\c+\bkgB, end angle=\c+\bkgP-\bkgB, radius=\bkgF]} -- cycle;
  \filldraw[fill=white, draw=#6, line width=0.5pt] (0,0) circle ({max(0.9,0.22*\bkgR)});
  \fill[#6] (0,0) circle (0.55);
\end{scope}}
\bkchapter{5}{6}{微分の道具箱：\\積と合成関数}{\bklevel{4}}{微分の道具箱：積と合成関数}
\bkmichishirube{2-5,3-3,5-4}{5-7,5-14,5-15,5-16}

\begin{goal}
\begin{itemize}[label={\color{bkTheme}●}, leftmargin=1.6\zw, itemsep=0pt, topsep=1pt]
\item \bklevel{3}\enspace 積の微分を、長方形の面積で説明し、計算できる。
\item \bklevel{4}\enspace 式の中から「関数の中の関数」（合成関数）を見つけられる。
\item \bklevel{4}\enspace 合成関数の微分を使って、$(2x+1)^5$ や $\sqrt{x^2+9}$ を微分できる。
\end{itemize}
\end{goal}

\begin{hatena}
\noindent\begin{minipage}[c]{0.49\linewidth}
3つの歯車A・B・Cが、かみ合って回っています。Aが1回転すると、Bは3回転します。Bが1回転すると、Cは2回転します。

Aを1秒に1回転させると、Cは1秒に何回転するでしょう。

回る速さの比が、つながっていくと……？
\end{minipage}\hfill
\begin{minipage}[c]{0.49\linewidth}\centering
\begin{tikzpicture}[x=1mm, y=1mm]
  % 第7回：歯数 A 48・B 16・C 8（A→B 3倍、B→C 2倍）。モジュール 0.6 mm、中心距離 19.2 mm・7.2 mm
  \bkgear{(0,0)}{48}{0.6}{0}{bkBlue!22}{bkBlue!75!black}
  \bkgear{(19.2,0)}{16}{0.6}{191.25}{bkOrange!28}{bkOrange!75!black}
  \bkgear{(26.4,0)}{8}{0.6}{180}{bkGreen!25}{bkGreen!75!black}
  % 回る向き（かみ合うとなりどうしは逆向き）
  \draw[-{Stealth[length=4pt]}, bkInk!75, line width=0.9pt] (130:10) arc[start angle=130, end angle=50, radius=10];
  \draw[-{Stealth[length=3.5pt]}, bkInk!75, line width=0.8pt] ([shift={(19.2,0)}]50:7.4) arc[start angle=50, end angle=130, radius=7.4];
  \draw[-{Stealth[length=3.5pt]}, bkInk!75, line width=0.8pt] ([shift={(26.4,0)}]70:4.4) arc[start angle=70, end angle=-70, radius=4.4];
  % 文字は歯先の円から 1.5 mm 以上はなす
  \node[font=\sffamily\small, anchor=north] at (0,-16.2) {\textbf{A}\enspace{\scriptsize 歯48まい}};
  \node[font=\sffamily\bfseries\small, anchor=north] at (19.2,-6.8) {B};
  \node[font=\sffamily\scriptsize, anchor=north] at (19.2,-11.2) {歯16まい};
  \node[font=\sffamily\bfseries\small, anchor=south west] at (32.3,0.2) {C};
  \node[font=\sffamily\scriptsize, anchor=north west] at (32.3,-0.2) {歯8まい};
  \node[font=\sffamily\scriptsize, text=bkGray, align=left, anchor=south west] at (13.2,8.6) {A：1回転 → B：3回転\\B：1回転 → C：2回転};
\end{tikzpicture}
\end{minipage}
\end{hatena}

Aが1回転するあいだに、Bは3回転します。そのBが3回転するあいだに、Cは1回転ごとに2回転するので、$3\times2=6$ 回転です。回る速さの比は、\textbf{かけ算でつながります}。この章の後半で学ぶ「合成関数の微分」は、これとまったく同じしくみです。前半では、かけ算の形をした式の微分（積の微分）を学びます。

\clearpage
%% ---------------- やってみよう・思い出そう ----------------
\begin{yattemiyou}[自転車のギアで、回る速さの比を調べよう]
\par\addvspace{3pt}\textbf{用意するもの}\par\noindent 変速ギアつきの自転車（大人といっしょに）、メジャー
\begin{enumerate}[label={\bkcircled[bkGreen]{\arabic*}}, leftmargin=2.4\zw, itemsep=1pt, topsep=3pt]
\item 自転車をスタンドで立て、ペダルを手でゆっくり1回転させる。後ろのタイヤは何回転するかな？
\item 前のギアの歯の数と、後ろのギアの歯の数を数え、①の結果とくらべよう。
\item タイヤのまわりの長さをはかり、ペダル1回転で何m進むかを計算しよう。
\item 後ろのギアを変えて、①〜③をくり返そう。
\end{enumerate}
\end{yattemiyou}

\noindent たとえば、前のギアの歯が42枚、後ろが14枚なら、ペダルが1回転するあいだに、チェーンは歯42枚分進みます。後ろのギアは、歯14枚分で1回転なので、$42\div14=3$ 回転します。後ろのギアとタイヤは同じ軸についているので、タイヤも3回転です。タイヤが1回転で\qty{2}{m}進むなら、ペダル1回転で進む距離は、次のように求められます。

\begin{center}
（ペダル1回転でタイヤが回る回数）×（タイヤ1回転で進む距離）\\
＝ 3 × \qty{2}{m} ＝ \qty{6}{m}
\end{center}

「ペダル→タイヤ」と「タイヤ→道」の2つの比を、かけ算でつなぎました。

\begin{omoidasou}
\bkrecall{5-4}{$(x^n)'=nx^{n-1}$。定数倍・和は1つずつ微分する。}
\bkrecall{3-3}{関数の出力を、別の関数に入れることができる。$u=3x+1$、$y=u^2$ なら $y=(3x+1)^2$。}
\bkrecall{2-11}{計算の順番：かっこの中を先に計算する。}
\bkrecall{2-5}{比はかけ算でつなげる。$\mathrm{A}:\mathrm{B}=1:3$、$\mathrm{B}:\mathrm{C}=1:2$ なら $\mathrm{A}:\mathrm{C}=1:6$。}
\bkrecall{3-11}{分数の指数：$x^{\frac{1}{2}}=\sqrt{x}$、$x^{-\frac{1}{2}}=\dfrac{1}{\sqrt{x}}$}
\bkrecall{4-5}{球の体積 $\dfrac{4}{3}\pi r^3$、円の面積 $\pi r^2$。円すいの体積は、底面積×高さ×$\dfrac{1}{3}$。}
\end{omoidasou}

%% ---------------- 数学の道具 ----------------
\bksection{doguu}{数学の道具}

\bksubsection{1}{お店の売り上げで考える}

\leavevmode\label{bk:5-6:mise}かけ算の形をした量が、どうふえるかを考えます。まずは、身近な例からです。

パン屋さんで、1個100円のパンが、1日に50個売れています。売り上げは、$100\times50=5000$ 円です。次の日、値段を2円上げて102円にしたら、51個売れました。売り上げは $102\times51=5202$ 円で、202円ふえました。この202円は、どこからきたのでしょう。

\noindent\begin{minipage}[c]{0.59\linewidth}\centering
\begin{tikzpicture}[x=1mm, y=1mm]
  \fill[bkBlue!15] (0,0) rectangle (52,30);
  \fill[bkOrange!35] (52,0) rectangle (68,30);
  \fill[bkGreen!30] (0,30) rectangle (52,40);
  \fill[bkGray!40] (52,30) rectangle (68,40);
  \draw[bkInk, line width=0.8pt] (0,0) rectangle (68,40);
  \draw[bkInk, line width=0.5pt] (52,0) -- (52,40) (0,30) -- (68,30);
  \node[font=\sffamily\small, align=center] at (26,15) {100円×50個\\＝5000円};
  \node[font=\sffamily\scriptsize, align=center] at (60,15) {2円\\×50個\\＝100円};
  \node[font=\sffamily\scriptsize] at (26,35) {100円×1個＝100円};
  \node[font=\sffamily\scriptsize, align=center] at (60,35) {2円×1個\\＝2円};
  \node[font=\sffamily\scriptsize, anchor=north] at (26,-1.5) {100円};
  \node[font=\sffamily\scriptsize, anchor=north] at (60,-1.5) {2円};
  \node[font=\sffamily\scriptsize, anchor=east] at (-1.5,15) {50個};
  \node[font=\sffamily\scriptsize, anchor=east] at (-1.5,35) {1個};
  \node[font=\sffamily\footnotesize, anchor=north] at (34,-6.5) {1個の値段（横）と、売れた個数（たて）};
\end{tikzpicture}
\end{minipage}\hfill
\begin{minipage}[c]{0.38\linewidth}
売り上げ＝値段×個数 を、長方形の面積と考えます。横が値段、たてが個数です（図は、ふえた分を大きくかいている）。

横が2円、たてが1個のびると、ふえた面積は3つの部分に分かれます。
\end{minipage}

\begin{bkcenter}
\begin{bktable}{colspec={Q[l,m,wd=50mm] Q[c,m,wd=34mm] Q[c,m,wd=20mm]}, rowsep=1.5pt, cells={font=\small}, row{1}={font=\sffamily\bfseries\small}}
部分 & 計算 & ふえた額 \\
値上げの分（右の帯） & 2円×50個 & 100円 \\
ふえた1個の分（上の帯） & 100円×1個 & 100円 \\
ふえた1個の値上げの分（右上） & 2円×1個 & 2円 \\
合計 & & 202円 \\
\end{bktable}
\bktabcaption{ふえた202円の内わけ。右上の小さな四角だけ、ずっと小さい。}
\end{bkcenter}

右上の小さな四角は、「値段のふえた分」と「個数のふえた分」をかけたものです。ふえた分どうしをかけるので、とても小さくなります。この「ふえた分×ふえた分」が、積の微分のかぎになります。

\bksubsection{2}{かけ算の微分：長方形の面積で考える}

\leavevmode\label{bk:5-6:seki}$x$ によって変わる2つの量 $f(x)$ と $g(x)$ の積 $f(x)g(x)$ を微分します。お店の例と同じように、横 $f$、たて $g$ の長方形の面積 $S=fg$ だと考えましょう。

\noindent\begin{minipage}[c]{0.44\linewidth}
\begin{tikzpicture}[x=1mm, y=1mm]
  % 第7回：のびた分の帯を広げ（横 16 mm・たて 8 mm）、帯の中の文字にゆとりをもたせる
  \fill[bkBlue!15] (0,0) rectangle (36,22);
  \fill[bkOrange!35] (36,0) rectangle (52,22);
  \fill[bkGreen!30] (0,22) rectangle (36,30);
  \fill[bkGray!40] (36,22) rectangle (52,30);
  \draw[bkInk, line width=0.8pt] (0,0) rectangle (52,30);
  \draw[bkInk, line width=0.5pt] (36,0) -- (36,30) (0,22) -- (52,22);
  \node[font=\sffamily\small] at (18,11) {$fg$};
  \node[font=\sffamily\scriptsize] at (44,11) {$g\,\Delta f$};
  \node[font=\sffamily\scriptsize] at (18,26) {$f\,\Delta g$};
  \node[font=\sffamily\scriptsize] at (44,26) {$\Delta f\Delta g$};
  \node[font=\sffamily\small, anchor=north] at (18,-1) {$f$};
  \node[font=\sffamily\small, anchor=north] at (44,-1) {$\Delta f$};
  \node[font=\sffamily\small, anchor=east] at (-1,11) {$g$};
  \node[font=\sffamily\small, anchor=east] at (-1,26) {$\Delta g$};
\end{tikzpicture}
\end{minipage}\hfill
\begin{minipage}[c]{0.53\linewidth}
$x$ が $\Delta x$ だけ変わると、横が $\Delta f$、たてが $\Delta g$ だけのびます。ふえる面積は、お店の例と同じ3つの部分です。
\[ \Delta S=g\,\Delta f+f\,\Delta g+\Delta f\,\Delta g \]
\end{minipage}

\vspace{4pt}
平均変化率 $\dfrac{\Delta S}{\Delta x}$ にするために、両辺を $\Delta x$ でわります。右辺は3つの項を1つずつわります（2-12）。3つめの項は、$\Delta f$ のほうだけを $\Delta x$ でわった形に書きます（$\dfrac{\Delta f\,\Delta g}{\Delta x}=\dfrac{\Delta f}{\Delta x}\times\Delta g$）。
\[ \frac{\Delta S}{\Delta x}=g\,\frac{\Delta f}{\Delta x}+f\,\frac{\Delta g}{\Delta x}+\frac{\Delta f}{\Delta x}\times\Delta g \]
ここで $\Delta x$ を0に近づけます。1つめと2つめの項は、$\dfrac{\Delta f}{\Delta x}\to f'$、$\dfrac{\Delta g}{\Delta x}\to g'$ なので、$gf'$ と $fg'$ に近づきます（5-3）。では、3つめの項 $\dfrac{\Delta f}{\Delta x}\times\Delta g$ は、何に近づくでしょう。

\bksubsection{3}{3つめの項は、なぜ0に近づくの？}

\leavevmode\label{bk:5-6:mittsume}3つめの項を、前と後ろに分けて考えます。

\noindent\textbf{前の部分\enspace$\dfrac{\Delta f}{\Delta x}$}\par
\noindent これは $f$ の平均変化率です。$\Delta x\to0$ で、$f'$ という決まった数に近づきます。たとえば $f=x^2$、$x=1$ なら、$f'=2x=2$ に近づきます。

\noindent\textbf{後ろの部分\enspace$\Delta g$}\par
\noindent これは「$g$ の変化そのもの」で、$\Delta x$ でわっていません。$x$ の変化 $\Delta x$ が0に近づけば、$g$ の変化 $\Delta g$ も0に近づきます。$x$ がほとんど変わらなければ、$g$ もほとんど変わらないからです。

だから、3つめの項は「2 × ほとんど0」で、0に近づきます。

\noindent\textbf{$\dfrac{\Delta g}{\Delta x}$ とのちがい}\par
\noindent ここで、「$\Delta g$ が0に近づくなら、$\dfrac{\Delta g}{\Delta x}$ も0に近づくのでは？」と思うかもしれません。でも、そうはなりません。$\dfrac{\Delta g}{\Delta x}$ は、分子の $\Delta g$ も、分母の $\Delta x$ も、どちらも小さくなります。小さい数を小さい数でわるので、わった値は0にならず、$g'$ に近づきます。$f=x^2$、$g=x^3$、$x=1$ で計算してみましょう。

\begin{bkcenter}
\begin{bktable}{colspec={Q[c,m,wd=16mm] Q[c,m,wd=24mm] Q[c,m,wd=24mm] Q[c,m,wd=20mm] Q[c,m,wd=30mm]}, rowsep=1.5pt, cells={font=\small}}
$\Delta x$ & $\Delta g$ & $\dfrac{\Delta g}{\Delta x}$ & $\dfrac{\Delta f}{\Delta x}$ & $\dfrac{\Delta f}{\Delta x}\times\Delta g$ \\
0.1 & 0.331 & 3.31 & 2.1 & 0.6951 \\
0.01 & 0.030301 & 3.0301 & 2.01 & 0.06090501 \\
0.001 & 0.003003001 & 3.003001 & 2.001 & 0.006009005 \\
近づく先 & 0 & 3（$g'=3x^2$） & 2（$f'=2x$） & 0 \\
\end{bktable}
\bktabcaption{$f=x^2$、$g=x^3$、$x=1$ のとき。$\Delta g$ は0に、$\dfrac{\Delta g}{\Delta x}$ は3に近づく。}
\end{bkcenter}

$\Delta g$ の列は、$\Delta x$ を $\dfrac{1}{10}$ にするたびに、およそ $\dfrac{1}{10}$ になり、0に近づきます。$\dfrac{\Delta g}{\Delta x}$ の列は、3に近づいて、0にはなりません。いちばん右の列は「2に近い数」×「0に近い数」なので、0に近づきます。

5-2の落ちるボールを思い出しましょう。1秒から1.001秒までの0.001秒のあいだに、ボールは $5\times1.001^2-5=0.010005$（m）しか落ちません。落ちた距離は0に近いのです。でも、落ちた距離÷時間は $0.010005\div0.001=10.005$ で、速さ\qty{10}{m/s}に近づきました。「変化そのもの」と「変化÷時間」は、ちがうものなのです。

\noindent\textbf{図で見ると}\par
\noindent 長方形の図で、$\Delta x$ を半分、また半分にしていきます。$\Delta x$ が半分になると、横ののび $\Delta f$ も、たてののび $\Delta g$ も、およそ半分になります。下の図は、1ますを1として、のびを4ます、2ます、1ますにしたものです。

\begin{bkfigure}
\begin{tikzpicture}[x=1mm, y=1mm]
\foreach \ox/\d/\lab in {0/4/{はじめ（のび4ます）}, 48.4/2/{$\Delta x$ を半分に（2ます）}, 92/1/{さらに半分に（1ます）}} {
  \begin{scope}[shift={(\ox,0)}]
    \fill[bkBlue!15] (0,0) rectangle (28.8,19.2);
    \fill[bkOrange!40] (28.8,0) rectangle ({28.8+2.4*\d},19.2);
    \fill[bkGreen!35] (0,19.2) rectangle (28.8,{19.2+2.4*\d});
    \fill[bkInk!55] (28.8,19.2) rectangle ({28.8+2.4*\d},{19.2+2.4*\d});
    \draw[bkInk!30, line width=0.2pt, step=2.4] (0,0) grid ({28.8+2.4*\d},{19.2+2.4*\d});
    \draw[bkInk, line width=0.7pt] (0,0) rectangle ({28.8+2.4*\d},{19.2+2.4*\d});
    \node[font=\sffamily\scriptsize, anchor=north] at ({(28.8+2.4*\d)/2},-1.5) {\lab};
  \end{scope}}
\end{tikzpicture}
\bkfigcaption{横12ます・たて8ますの長方形を、のばしていく。オレンジ・緑は帯、こい灰色は右上の小さな四角。}
\end{bkfigure}

\begin{bkcenter}
\begin{bktable}{colspec={Q[l,m,wd=44mm] Q[c,m,wd=22mm] Q[c,m,wd=22mm] Q[c,m,wd=22mm]}, rowsep=1.5pt, cells={font=\small}}
 & はじめ & 半分に & さらに半分に \\
のび（$\Delta f$ と $\Delta g$） & 4 & 2 & 1 \\
右の帯（$8\times$ のび） & 32 & 16 & 8 \\
上の帯（$12\times$ のび） & 48 & 24 & 12 \\
右上の小さな四角（のび×のび） & 16 & 4 & 1 \\
\end{bktable}
\bktabcaption{のびが半分になると、帯は半分に、右上の小さな四角は4分の1になる。}
\end{bkcenter}

帯の面積は、のびを半分にすると半分になります。右上の小さな四角は、たても横も半分になるので、$\dfrac{1}{2}\times\dfrac{1}{2}=\dfrac{1}{4}$ になります。のびでわってみると、右の帯は $32\div4=8$、$16\div2=8$、$8\div1=8$ でいつも8、上の帯はいつも12です。右上の四角は $16\div4=4$、$4\div2=2$、$1\div1=1$ と、半分ずつになって0に近づきます。右上の小さな四角は、まわりの帯よりずっと速く小さくなり、わり算しても消えていくのです。

\noindent\textbf{積の微分の式}\par
\noindent これで、3つの項の行き先がわかりました。
\[ \frac{\Delta S}{\Delta x}=g\,\frac{\Delta f}{\Delta x}+f\,\frac{\Delta g}{\Delta x}+\frac{\Delta f}{\Delta x}\times\Delta g\ \longrightarrow\ gf'+fg'+0 \]
だから、
\[ \{f(x)g(x)\}'=f'(x)g(x)+f(x)g'(x) \]
これを\textbf{積の微分}といいます。「片方ずつ微分して、たす」と覚えましょう。1つめは「$f$ だけが変わった分（右の帯）」、2つめは「$g$ だけが変わった分（上の帯）」です。\bkidx{せきのびぶん}{積の微分}

\noindent\textbf{確かめ}\par
\noindent $x^2\times x^3=x^5$ なので、5-4のきまりで、答えは $5x^4$ のはずです。積の微分を使うと、
\[ (x^2)'\times x^3+x^2\times(x^3)'=2x\times x^3+x^2\times3x^2=2x^4+3x^4=5x^4 \]
合っています。$x=1$ では $2+3=5$ です。上の表の3つの項をたすと、$\Delta x=0.001$ で $2.001+3.003001+0.006009\cdots=5.010010\cdots$ となり、5に近づいています。

\bksubsection{4}{関数の中の関数：合成関数}

\leavevmode\label{bk:5-6:gousei}$y=(3x+1)^2$ を、$x$ に数を入れて計算してみましょう。$x=1$ なら、かっこの中を先に計算して（2-11）、$3\times1+1=4$。それを2乗して $4^2=16$ です。

\begin{bkcenter}
\begin{bktable}{colspec={Q[l,m,wd=40mm] *{3}{Q[c,m,wd=14mm]}}, rowsep=1.5pt}
$x$ & 0 & 1 & 2 \\
① かっこの中 $u=3x+1$ & 1 & 4 & 7 \\
② それを2乗 $y=u^2$ & 1 & 16 & 49 \\
\end{bktable}
\bktabcaption{$y=(3x+1)^2$ は、①と②の2つのはたらきを、順に行っている。}
\end{bkcenter}

このように、$y=(3x+1)^2$ は、2つのはたらきを順に行います。まず $x$ を $u=3x+1$ にし、次にその $u$ を2乗して $y=u^2$ にします。

\begin{bkfigure}
\begin{tikzpicture}[x=1mm, y=1mm, m/.style={draw=bkTheme, line width=1pt, fill=white, rounded corners=2mm, minimum width=30mm, minimum height=11mm, align=center, font=\sffamily\small}]
  \node[font=\sffamily\bfseries] (x) at (0,0) {$x$};
  \node[m] (g) at (30,0) {内側の関数\\$u=3x+1$};
  \node[m] (f) at (76,0) {外側の関数\\$y=u^2$};
  \node[font=\sffamily\bfseries] (y) at (108,0) {$y$};
  \draw[bk arrow, bkGray] (x) -- (g);
  \draw[bk arrow, bkGray] (g) -- node[above, font=\sffamily\small] {$u$} (f);
  \draw[bk arrow, bkGray] (f) -- (y);
  \node[font=\sffamily\scriptsize, text=bkGray, anchor=north] at (0,-4) {1};
  \node[font=\sffamily\scriptsize, text=bkGray, anchor=north] at (53,-4) {4};
  \node[font=\sffamily\scriptsize, text=bkGray, anchor=north] at (108,-4) {16};
\end{tikzpicture}
\bkfigcaption{2つの関数を順につないだものが、合成関数。灰色の数は $x=1$ のとき。}
\end{bkfigure}

このように、関数の出力を別の関数に入れてつくった関数を、\textbf{合成関数}（ごうせいかんすう）といいます（3-3）。先にはたらくほうを\textbf{内側の関数}、後にはたらくほうを\textbf{外側の関数}とよびます。かっこの中は先に計算するので、多くの場合、かっこの中が内側です。\bkidx{ごうせいかんすう}{合成関数}

\begin{bkcenter}
\begin{bktable}{colspec={Q[c,m,wd=32mm] Q[c,m,wd=32mm] Q[c,m,wd=32mm]}, rowsep=2pt}
式 & 内側 $u=$ & 外側 $y=$ \\
$(2x-5)^4$ & $2x-5$ & $u^4$ \\
$\sqrt{x^2+9}$ & $x^2+9$ & $\sqrt{u}=u^{\frac{1}{2}}$ \\
$\dfrac{1}{2x+1}$ & $2x+1$ & $\dfrac{1}{u}=u^{-1}$ \\
\end{bktable}
\bktabcaption{合成関数を、内側と外側に分ける。「かっこの中」や「ルートの中」や「分母」が内側になることが多い。}
\end{bkcenter}

\bksubsection{5}{合成関数の微分：変わる速さの比をかけ算でつなぐ}

\leavevmode\label{bk:5-6:chain}$y=(3x+1)^2$ で、$x$ を少しだけ変えてみます。$x$ が変わると $u$ が変わり、$u$ が変わると $y$ が変わります。$x=1$ から $\Delta x=0.01$ だけ変えると、

\begin{bkfigure}
\begin{tikzpicture}[x=1mm, y=1mm]
  \foreach \yy/\nm/\len/\col in {30/x/4/bkBlue, 15/u/12/bkGreen, 0/y/96/bkOrange} {
    \draw[bkInk, line width=0.6pt, -{Stealth[length=2.5mm]}] (10,\yy) -- (124,\yy);
    \node[font=\sffamily\small, anchor=east] at (8,\yy) {$\nm$};
    \draw[\col, line width=2.4pt] (14,\yy) -- ({14+\len},\yy);
    \draw[bkInk, line width=0.6pt] (14,{\yy-1.5}) -- (14,{\yy+1.5}) ({14+\len},{\yy-1.5}) -- ({14+\len},{\yy+1.5});}
  \node[font=\sffamily\footnotesize, anchor=south west, inner sep=1pt] at (13,32) {1 から 1.01 まで：$\Delta x=0.01$};
  \node[font=\sffamily\footnotesize, anchor=south west, inner sep=1pt] at (13,17) {4 から 4.03 まで：$\Delta u=0.03$（$\Delta x$ の3倍）};
  \node[font=\sffamily\footnotesize, anchor=south west, inner sep=1pt] at (13,2) {16 から 16.2409 まで：$\Delta y=0.2409$（$\Delta u$ の約8倍）};
\end{tikzpicture}
\bkfigcaption{$x$・$u$・$y$ の変化を、太い線の長さで表した（3本とも、同じ拡大のしかた）。}
\end{bkfigure}

$u=3x+1$ は、$x$ が変わった分の3倍だけ変わります。$y=u^2$ は、$u=4$ のあたりでは、$u$ が変わった分のおよそ8倍だけ変わります（$u^2$ の微分は $2u$ で、$u=4$ なら $2\times4=8$。5-4のきまりを、文字 $u$ で使った）。だから、$y$ は $x$ が変わった分の、およそ $8\times3=24$ 倍変わります。実際、$\dfrac{\Delta y}{\Delta x}=\dfrac{0.2409}{0.01}=24.09$ で、24に近い数です。歯車の回る速さの比をかけ算でつないだのと同じです。

式で書くと、変化の比は、次のようにかけ算でつながります。
\[ \frac{\Delta y}{\Delta x}=\frac{\Delta y}{\Delta u}\times\frac{\Delta u}{\Delta x} \]
右辺は、分数のかけ算で $\Delta u$ を約分すると、左辺になります。$\Delta x$ を0に近づけると、$\Delta u$ も0に近づくので、3つの比はそれぞれ導関数になります。
\[ \frac{\dd y}{\dd x}=\frac{\dd y}{\dd u}\times\frac{\dd u}{\dd x} \]
これを\textbf{合成関数の微分}（英語で chain rule：くさりのきまり）といいます。「外側を微分して、内側の微分をかける」と覚えましょう。\bkidx{ごうせいかんすうのびぶん}{合成関数の微分}

\noindent\textbf{手順}\par
\noindent $y=(3x+1)^2$ を、合成関数の微分で微分します。
\begin{enumerate}[label={\bkcircled{\arabic*}}, leftmargin=2.4\zw, itemsep=1pt, topsep=2pt]
\item 内側と外側に分ける：$u=3x+1$、$y=u^2$
\item 外側を $u$ で微分する：$\dfrac{\dd y}{\dd u}=2u$
\item 内側を $x$ で微分する：$\dfrac{\dd u}{\dd x}=3$
\item かけ算する：$\dfrac{\dd y}{\dd x}=2u\times3=6u$
\item $u$ を $x$ の式にもどす：$6u=6(3x+1)$
\end{enumerate}
⑤を忘れると、答えに $u$ が残ってしまいます。問題は $x$ の式なので、答えも $x$ の式にします。

\noindent\textbf{確かめ}\par
\noindent 展開すると $(3x+1)^2=9x^2+6x+1$。微分して $18x+6$。合成関数の微分の答え $6(3x+1)=18x+6$ と同じです。$x=1$ なら $6\times4=24$ で、上の図の24.09とも合います。展開できない式や、展開が大変な式でも、合成関数の微分なら同じように計算できます。

\noindent\textbf{表計算でも確かめよう}\par
\noindent $y=(3x+1)^2$ について、$h=0.001$ の平均変化率と、合成関数の微分の答え $6(3x+1)$ をならべます。
\begin{bkcenter}
\bksheet{3}{30}{%
1 & $x$ & 平均変化率（$h=0.001$） & $6(3x+1)$ \\
2 & 0 & 6.009 & 6 \\
3 & 1 & 24.009 & 24 \\
4 & 2 & 42.009 & 42 \\}
\end{bkcenter}
はしたの 0.009 は $h$ が0でない分で、$h$ を小さくすると消えていきます。

\begin{seikaku}
上の説明では、$\Delta u$ でわっています。$\Delta x$ が0でなくても、$\Delta u$ が0になることがあります（たとえば内側の関数が一定のとき）。大学の数学では、その場合も合わせて、きちんと証明します。結果の式は、どの場合でも正しく成り立ちます。
\end{seikaku}

\begin{column}{ライプニッツの記号が便利なわけ}
合成関数の微分 $\dfrac{\dd y}{\dd x}=\dfrac{\dd y}{\dd u}\times\dfrac{\dd u}{\dd x}$ は、分数の約分のように見えます。$\dfrac{\dd y}{\dd x}$ はほんとうの分数ではありませんが、もとは $\dfrac{\Delta y}{\Delta x}$ という分数の極限なので、このように計算してよいのです。ライプニッツは、計算のきまりが自然に見えるように記号を工夫しました。300年以上たったいまでも、工学の専門書ではこの記号が使われています。
\end{column}

\bksubsection{6}{よく使う形}

\leavevmode\label{bk:5-6:katachi}\textbf{1次式の $n$ 乗}\par
\noindent 内側が $u=ax+b$ なら $\dfrac{\dd u}{\dd x}=a$ です。外側 $u^n$ の微分は $nu^{n-1}$ なので、
\[ \{(ax+b)^n\}'=n(ax+b)^{n-1}\times a \]
たとえば $\{(2x+1)^5\}'=5(2x+1)^4\times2=10(2x+1)^4$。展開すると6つの項が出てくる式でも、1行で微分できます。

\noindent\textbf{ルートの微分}\par
\noindent $(x^n)'=nx^{n-1}$ は、$n$ が分数のときも成り立ちます。$n$ が $\dfrac{1}{2}$ のときの理由は、下のもっと知りたいで確かめます（ほかの分数は5-14）。$\sqrt{x}=x^{\frac{1}{2}}$ なので、指数の $\dfrac{1}{2}$ を前に出します。そして、指数を1へらして、$\dfrac{1}{2}-1=-\dfrac{1}{2}$ にします。
\[ (\sqrt{x})'=\frac{1}{2}x^{\frac{1}{2}-1}=\frac{1}{2}x^{-\frac{1}{2}}=\frac{1}{2\sqrt{x}} \]
表計算で確かめると、$x=4$ では、
\[ \frac{\sqrt{4.001}-2}{0.001}=0.24998\cdots \]
で、$\dfrac{1}{2\sqrt{4}}=\dfrac{1}{4}=0.25$ とよく合います。

\begin{motto}[4]
\textbf{ルートの微分を、極限で確かめる}\par
\[ \frac{\sqrt{x+h}-\sqrt{x}}{h} \]
\noindent の分子と分母に $\sqrt{x+h}+\sqrt{x}$ をかけます。分子は $(a-b)(a+b)=a^2-b^2$（3-8）を使うと、$(\sqrt{x+h})^2-(\sqrt{x})^2=(x+h)-x=h$ になります。
\[ \frac{\sqrt{x+h}-\sqrt{x}}{h}=\frac{h}{h(\sqrt{x+h}+\sqrt{x})}=\frac{1}{\sqrt{x+h}+\sqrt{x}} \]
$h\to0$ で、分母は $\sqrt{x}+\sqrt{x}=2\sqrt{x}$ に近づくので、$(\sqrt{x})'=\dfrac{1}{2\sqrt{x}}$ です。
\end{motto}

\noindent\textbf{分数の形}\par
\noindent $\dfrac{1}{f(x)}=\{f(x)\}^{-1}$ と考えると、外側は $u^{-1}$ で、その微分は $-1\times u^{-2}=-u^{-2}$ です（5-4）。内側の微分 $f'(x)$ をかけて、
\[ \Bigl\{\frac{1}{f(x)}\Bigr\}'=-\{f(x)\}^{-2}\times f'(x)=-\frac{f'(x)}{\{f(x)\}^2} \]

\bksubsection{7}{つながった変化の速さ}

\leavevmode\label{bk:5-6:rensa}合成関数の微分は、「ある量が時間とともに変わると、それにつながった量がどんな速さで変わるか」を計算するのに使えます。

\noindent\begin{minipage}[c]{0.66\linewidth}
\textbf{例：風船をふくらませる}\par
半径が1秒に\qty{2}{cm}ずつ、$r=2t$（cm）でふえていくとき、体積 $V=\dfrac{4}{3}\pi r^3$ のふえる速さを求めます。$V$ は $r$ の関数で、$r$ は $t$ の関数です。$t\to r\to V$ の順につながった合成関数で、内側が $r=2t$、外側が $V=\dfrac{4}{3}\pi r^3$ です。
\[ \frac{\dd V}{\dd t}=\frac{\dd V}{\dd r}\times\frac{\dd r}{\dd t}=4\pi r^2\times2 \]
$\dfrac{\dd V}{\dd r}=\dfrac{4}{3}\pi\times3r^2=4\pi r^2$（5-4）、$\dfrac{\dd r}{\dd t}=2$ です。
\end{minipage}\hfill
\begin{minipage}[c]{0.3\linewidth}\centering
\begin{tikzpicture}[x=1mm, y=1mm]
  \shade[ball color=bkRed!60] (0,0) circle (8);
  \draw[bkRed!50!black, line width=0.6pt, dashed] (0,0) circle (12);
  \draw[bkInk!70, line width=0.8pt] (0,-8) -- (0,-15);
  \draw[bk arrow, bkGray, line width=0.8pt] (8.5,0) -- (11.5,0);
\end{tikzpicture}
\end{minipage}

\vspace{2pt}
$t=1$ 秒（$r=\qty{2}{cm}$）なら $4\pi\times2^2\times2=32\pi$ で、約\qty{100.5}{cm^3/s}。$t=2$ 秒（$r=\qty{4}{cm}$）なら $4\pi\times4^2\times2=128\pi$ で、約\qty{402}{cm^3/s} です。半径は同じ速さでのびているのに、体積のふえる速さは4倍になりました。$r^2$ が $2^2=4$ から $4^2=16$ へ、4倍になったからです。風船が大きくなるほど、体積はどんどん速くふえます。

\bksubsection{8}{円すい形のタンクの水位}

\leavevmode\label{bk:5-6:tank}\noindent\begin{minipage}[c]{0.64\linewidth}
先がとがった円すい形のタンクに、毎秒\qty{50}{cm^3}の水を入れます。このタンクは、深さ\qty{40}{cm}の所で水面の半径が\qty{20}{cm}、深さ\qty{20}{cm}の所で\qty{10}{cm}というように、水面の半径がいつも深さの半分になる形です。だから、水の深さが $h$（cm）のとき、水面の半径は $\dfrac{h}{2}$ です。

水の体積は、円すいの体積の公式（4-5）から
\[ V=\frac{1}{3}\times\pi\Bigl(\frac{h}{2}\Bigr)^2\times h=\frac{1}{3}\times\frac{\pi h^2}{4}\times h=\frac{\pi}{12}h^3 \]
\end{minipage}\hfill
\begin{minipage}[c]{0.32\linewidth}\centering
\begin{tikzpicture}[x=1mm, y=1mm]
  \fill[bkBlue!25] (0,0) -- (-7,14) -- (7,14) -- cycle;
  \draw[bkInk, line width=0.9pt] (-14,28) -- (0,0) -- (14,28);
  \draw[bkInk, line width=0.6pt] (0,28) ellipse (14 and 2.5);
  \draw[bkBlue!60!black, line width=0.6pt] (0,14) ellipse (7 and 1.3);
  \draw[bkGray, line width=0.6pt, <->] (10,0) -- (10,14);
  \node[font=\sffamily\scriptsize, anchor=west] at (11,7) {$h$};
  \draw[bkGray, line width=0.6pt, <->] (0,15.8) -- (7,15.8);
  \node[font=\sffamily\scriptsize, anchor=south] at (3.5,16) {$\frac{h}{2}$};
  \draw[bk arrow, bkBlue, line width=1pt] (-6,36) -- (-2,30);
  \node[font=\sffamily\scriptsize, anchor=south] at (-6,36) {毎秒\qty{50}{cm^3}};
\end{tikzpicture}
\end{minipage}

\vspace{2pt}
$V$ も $h$ も、時間とともに変わります。$V$ は $h$ の関数で、$h$ は時間 $t$ の関数なので、合成関数の微分が使えます。$\dfrac{\dd V}{\dd h}=\dfrac{\pi}{12}\times3h^2=\dfrac{\pi}{4}h^2$ なので、
\[ \frac{\dd V}{\dd t}=\frac{\dd V}{\dd h}\times\frac{\dd h}{\dd t}=\frac{\pi}{4}h^2\times\frac{\dd h}{\dd t} \]
水を毎秒\qty{50}{cm^3}入れるので、$\dfrac{\dd V}{\dd t}=50$ です。両辺を $\dfrac{\pi}{4}h^2$ でわると（2-14）、水位の上がる速さは
\[ \frac{\dd h}{\dd t}=50\div\frac{\pi}{4}h^2=50\times\frac{4}{\pi h^2}=\frac{200}{\pi h^2} \]
深さ\qty{10}{cm}では $\dfrac{200}{100\pi}=\dfrac{2}{\pi}$ で約\qty{0.64}{cm/s}、深さ\qty{20}{cm}では $\dfrac{200}{400\pi}=\dfrac{1}{2\pi}$ で約\qty{0.16}{cm/s}です。同じ量の水を入れても、上のほうほど水面が広いので、水位はゆっくり上がります。

\bksubsection{9}{振り子の長さと周期}

\leavevmode\label{bk:5-6:furiko}\noindent\begin{minipage}[c]{0.7\linewidth}
長さ $L$（m）の振り子が1往復する時間（周期）$T$（s）は、
\[ T=2\pi\sqrt{\frac{L}{g}}\qquad(g=9.8) \]
です（2-8で、周期が長さの $\sqrt{\phantom{L}}$ に比例することを表で調べた。この式は6-6で求めます）。振り子時計は、おもりの位置を調整ねじで動かして $L$ を変え、時計の進みおくれを直します。$L$ を変えると、$T$ はどんな速さで変わるでしょう。\bkidx{しゅうき}{周期}
\end{minipage}\hfill
\begin{minipage}[c]{0.26\linewidth}\centering
\begin{tikzpicture}[x=1mm, y=1mm]
  \fill[bkGray!50] (-8,0) rectangle (8,2);
  \draw[bkInk, line width=0.7pt] (0,0) -- (-9,-28);
  \draw[bkInk!40, line width=0.5pt, dashed] (0,0) -- (0,-30);
  \shade[ball color=bkOrange] (-9,-28) circle (3);
  \draw[bkGray, line width=0.5pt, <->] (3,0) -- (3,-28);
  \node[font=\sffamily\scriptsize, anchor=west] at (4,-14) {$L$};
\end{tikzpicture}
\end{minipage}

\vspace{2pt}
内側を $u=\dfrac{L}{g}$、外側を $T=2\pi\sqrt{u}$ と考えます。
\begin{enumerate}[label={\bkcircled{\arabic*}}, leftmargin=2.4\zw, itemsep=1pt, topsep=2pt]
\item 内側の微分：$u=\dfrac{1}{g}\times L$ なので、$\dfrac{\dd u}{\dd L}=\dfrac{1}{g}$
\item 外側の微分：$2\pi\times\dfrac{1}{2\sqrt{u}}=\dfrac{\pi}{\sqrt{u}}$（ルートの微分）
\item かけ算して、$u$ をもどす：$\dfrac{\dd T}{\dd L}=\dfrac{\pi}{\sqrt{L/g}}\times\dfrac{1}{g}=\dfrac{\pi}{g\sqrt{L/g}}$
\end{enumerate}
分母の $g\sqrt{L/g}$ は、$g=\sqrt{g^2}$ と考えると $\sqrt{g^2\times\dfrac{L}{g}}=\sqrt{gL}$ です。だから、
\[ \frac{\dd T}{\dd L}=\frac{\pi}{\sqrt{gL}} \]
$L=\qty{1}{m}$ なら $\dfrac{\pi}{\sqrt{9.8}}=\dfrac{3.14}{3.13}\approx1.0$（s/m）です（$\approx$ は「ほぼ等しい」。5-5・5-8）。長さを\qty{1}{mm}（\qty{0.001}{m}）のばすと、周期は約0.001秒長くなります（この見積もり方は5-8）。長さ\qty{1}{m}の振り子の周期は約2秒なので、1日（86400秒）に約43000回往復します。1回ごとに0.001秒おくれると、1日に約43秒おくれる計算です。調整ねじが細かく作られているのは、このためです。

\bksubsection{10}{どの道具を使う？}

\begin{bkfigure}
\begin{tikzpicture}[x=1mm, y=1mm, q/.style={bk node, minimum width=47mm, minimum height=15mm},
  r/.style={bk node, draw=bkOrange, fill=bkOrange!8, minimum width=51mm, minimum height=15mm}]
  % 第6回：たてに組み直す（はい は右、いいえ は下）。箱の中の余白を広げ、字を大きくした
  \node[q] (q1) at (23.5,0) {2つの式の\\かけ算になっている？};
  \node[r] (a1) at (95,0) {積の微分\\$f'g+fg'$};
  \node[q] (q2) at (23.5,-30) {かっこやルートの中に\\$x$ の式が入っている？};
  \node[r] (a2) at (95,-30) {合成関数の微分\\外側の微分×内側の微分};
  \node[r, minimum width=47mm] (a3) at (23.5,-60) {$(x^n)'=nx^{n-1}$\\定数倍・和のきまり};
  \draw[bk link] (q1) -- node[bk link label, above] {はい} (a1);
  \draw[bk link] (q1) -- node[bk link label, right] {いいえ} (q2);
  \draw[bk link] (q2) -- node[bk link label, above] {はい} (a2);
  \draw[bk link] (q2) -- node[bk link label, right] {いいえ} (a3);
  \draw[bk link, dashed] (a1) -- node[bk link label, right, align=left] {かけ算の中に\\かっこがあれば} (a2);
\end{tikzpicture}
\bkfigcaption{微分の道具のえらび方。組み合わせて使うことも多い（例題5）。}
\end{bkfigure}

\begin{bkcenter}
\begin{bktable}{colspec={Q[c,m,wd=24mm] Q[c,m,wd=16mm] Q[c,m,wd=14mm] Q[c,m,wd=18mm] Q[c,m,wd=14mm] Q[c,m,wd=25mm]}, rowsep=2pt}
式 & 内側 $u$ & 外側 & 外側の微分 & 内側の微分 & 答え \\
$(5x+2)^3$ & $5x+2$ & $u^3$ & $3u^2$ & $5$ & $15(5x+2)^2$ \\
$(x^2+1)^4$ & $x^2+1$ & $u^4$ & $4u^3$ & $2x$ & $8x(x^2+1)^3$ \\
$\sqrt{3x+1}$ & $3x+1$ & $\sqrt{u}$ & $\dfrac{1}{2\sqrt{u}}$ & $3$ & $\dfrac{3}{2\sqrt{3x+1}}$ \\
$\dfrac{1}{x^2+4}$ & $x^2+4$ & $u^{-1}$ & $-u^{-2}$ & $2x$ & $-\dfrac{2x}{(x^2+4)^2}$ \\
\end{bktable}
\bktabcaption{合成関数の微分の手順。左から順にうめていけば、まちがえない。}
\end{bkcenter}

\bkkigocards{\bkkigocard{$f\circ g$}{エフ・まる・ジー}{合成関数。先に $g$、次に $f$。$(f\circ g)(x)=f(g(x))$}{composite}{5-6}}
{\bkkigocard{$x^{\frac{1}{2}}$}{エックスの2分の1乗}{$\sqrt{x}$ のこと。分数の指数}{x to the one-half}{3-11}}

\begin{senmonsho}[大学の教科書ふうの書き方]
\begin{bkhonbun}
\par\addvspace{3pt}\textbf{定理 3.4（合成関数の微分法）}\par\noindent 関数 $g$ が $x$ で微分可能で、関数 $f$ が $g(x)$ で微分可能ならば、合成関数 $f\circ g$ も $x$ で微分可能であり、
\[ (f\circ g)'(x)=f'\bigl(g(x)\bigr)\,g'(x) \tag{3.7} \]
が成り立つ。これを連鎖律（chain rule）という。
\end{bkhonbun}
\bkyaku{$f\circ g$（エフ・まる・ジー）は、「先に $g$、次に $f$」の順につないだ合成関数のことで、$(f\circ g)(x)=f(g(x))$ と書く。内側の $g$ と外側の $f$ がどちらも微分できれば、合成関数も微分できる。その微分は「外側の微分（内側の値のところで）×内側の微分」になる。「微分可能」は「接線が引ける（とがっていない）」こと（5-3）。}
\end{senmonsho}

\begin{taisetsu}
積の微分\quad $\{f(x)g(x)\}'=f'(x)g(x)+f(x)g'(x)$\\[2pt]
合成関数の微分\quad $\dfrac{\dd y}{\dd x}=\dfrac{\dd y}{\dd u}\times\dfrac{\dd u}{\dd x}$\quad（外側の微分×内側の微分）
\end{taisetsu}

\begin{motto}[4]
\textbf{商の微分}\par
$\dfrac{f}{g}=f\times\dfrac{1}{g}$ と考えて、積の微分と分数の形の微分を組み合わせると、
\[ \Bigl(\frac{f}{g}\Bigr)'=f'\times\frac{1}{g}+f\times\Bigl(-\frac{g'}{g^2}\Bigr)=\frac{f'g}{g^2}-\frac{fg'}{g^2}=\frac{f'g-fg'}{g^2} \]
となります（3つめの式は、1つめの分数の分子と分母に $g$ をかけて、分母を $g^2$ にそろえた）。新しく覚えなくても、積と合成関数の微分から作れます。
\end{motto}

\begin{kotoba}
\small
\bkword{積の微分}{かけ算の形の関数の微分。片方ずつ微分してたす。}
\bkword{合成関数}{関数の出力を、別の関数に入れてつくった関数。内側と外側がある。}
\bkword{合成関数の微分}{外側を微分して、内側の微分をかける。chain rule（チェーン・ルール）ともいう。}
\end{kotoba}

%% ---------------- 例題 ----------------
\bksection{reidai}{例題でたしかめよう}

\begin{reidai}[積の微分]
$y=(2x+1)(x^2-3)$ を微分しましょう。\par
\kangaekata $f=2x+1$、$g=x^2-3$ の積と見ます。$f'=2$、$g'=2x$ です。$\{fg\}'=f'g+fg'$ に入れます。\par
\shiki $y'=(2x+1)'(x^2-3)+(2x+1)(x^2-3)'$\par
\hspace*{2.6\zw}$\phantom{y'}=2(x^2-3)+(2x+1)\times2x$\par
\hspace*{2.6\zw}$\phantom{y'}=2x^2-6+4x^2+2x=6x^2+2x-6$\par
\hspace*{2.6\zw}確かめ：展開すると $2x^3+x^2-6x-3$。微分して $6x^2+2x-6$\par
\kotae $y'=6x^2+2x-6$
\end{reidai}

\begin{jibun}
$y=(x+2)(x^2+1)$ を、積の微分で微分しましょう。
\bkjans{$(x^2+1)+(x+2)\times2x=3x^2+4x+1$}
\end{jibun}

\begin{reidai}[かっこの $n$ 乗]
次の関数を微分しましょう。\par
（1）$(2x-5)^4$\quad（2）$(4-x)^3$\quad（3）$(x^2+1)^4$\par
\kangaekata かっこの中が内側です。外側 $u^n$ の微分は $nu^{n-1}$。それに内側の微分をかけます。\par
\shiki（1）内側 $u=2x-5$ の微分は2。$4(2x-5)^3\times2=8(2x-5)^3$\par
\hspace*{2.6\zw}（2）内側 $u=4-x$ の微分は $-1$。$3(4-x)^2\times(-1)=-3(4-x)^2$\par
\hspace*{2.6\zw}（3）内側 $u=x^2+1$ の微分は $2x$。$4(x^2+1)^3\times2x=8x(x^2+1)^3$\par
\kotae （1）$8(2x-5)^3$\quad（2）$-3(4-x)^2$\quad（3）$8x(x^2+1)^3$
\end{reidai}

\begin{jibun}
次の関数を微分しましょう。\par
（1）$(3x+2)^5$\quad（2）$(2-3x)^4$\quad（3）$(x^2+3)^3$
\bkjans{（1）\nobreak $15(3x+2)^4$\quad（2）\nobreak $-12(2-3x)^3$\quad（3）\nobreak $6x(x^2+3)^2$}
\end{jibun}

\begin{reidai}[ルートをふくむ式]
$y=\sqrt{x^2+9}$ を微分し、$x=4$ での値を求めましょう。\par
\kangaekata 内側 $u=x^2+9$（微分すると $2x$）、外側 $\sqrt{u}$（微分すると $\dfrac{1}{2\sqrt{u}}$）。\par
\shiki $y'=\dfrac{1}{2\sqrt{x^2+9}}\times2x=\dfrac{2x}{2\sqrt{x^2+9}}=\dfrac{x}{\sqrt{x^2+9}}$\par
\hspace*{2.6\zw}$x=4$ で $\dfrac{4}{\sqrt{16+9}}=\dfrac{4}{\sqrt{25}}=\dfrac{4}{5}$\par
\kotae $y'=\dfrac{x}{\sqrt{x^2+9}}$、$x=4$ で $\dfrac{4}{5}$
\end{reidai}

\begin{jibun}
$y=\sqrt{2x+1}$ を微分し、$x=4$ での値を求めましょう。
\bkjans{$\dfrac{1}{2\sqrt{2x+1}}\times2=\dfrac{1}{\sqrt{2x+1}}$。$x=4$ で $\dfrac{1}{3}$}
\end{jibun}

\begin{reidai}[分数の形]
次の関数を微分しましょう。\quad（1）$\dfrac{1}{2x+1}$\quad（2）$\dfrac{1}{(x+3)^2}$\par
\kangaekata 分母を、マイナスの指数で書きなおします（2-7）。指数を1へらすと、$-1$ は $-2$ に、$-2$ は $-3$ になります。\par
\shiki（1）$y=(2x+1)^{-1}$。内側 $2x+1$ の微分は2\par
\hspace*{4.3\zw}$y'=-1\times(2x+1)^{-2}\times2=-\dfrac{2}{(2x+1)^2}$\par
\hspace*{2.6\zw}（2）$y=(x+3)^{-2}$。内側 $x+3$ の微分は1\par
\hspace*{4.3\zw}$y'=-2\times(x+3)^{-3}\times1=-\dfrac{2}{(x+3)^3}$\par
\kotae （1）$-\dfrac{2}{(2x+1)^2}$\quad（2）$-\dfrac{2}{(x+3)^3}$
\end{reidai}

\begin{jibun}
次の関数を微分しましょう。\quad（1）$\dfrac{1}{x^2+1}$\quad（2）$\dfrac{1}{(2x-1)^2}$
\bkjans{（1）\nobreak $-(x^2+1)^{-2}\times2x=-\dfrac{2x}{(x^2+1)^2}$\quad（2）\nobreak $-2(2x-1)^{-3}\times2=-\dfrac{4}{(2x-1)^3}$}
\end{jibun}

\begin{reidai}[積と合成関数の組み合わせ]
$y=x(x+1)^3$ を微分しましょう。\par
\kangaekata $x$ と $(x+1)^3$ の積です。$(x+1)^3$ の微分には、合成関数の微分を使います（内側 $x+1$ の微分は1）。\par
\shiki $\{(x+1)^3\}'=3(x+1)^2\times1=3(x+1)^2$\par
\hspace*{2.6\zw}$y'=(x)'\times(x+1)^3+x\times3(x+1)^2=(x+1)^3+3x(x+1)^2$\par
\hspace*{2.6\zw}まとめる：どちらの項にも $(x+1)^2$ があるので、くくり出す\par
\hspace*{2.6\zw}$(x+1)^2\{(x+1)+3x\}=(x+1)^2(4x+1)$\par
\kotae $y'=(x+1)^3+3x(x+1)^2$\quad（まとめると $(x+1)^2(4x+1)$）
\end{reidai}

\begin{jibun}
$y=x^2(2x-1)^3$ を微分しましょう。
\bkjans{$2x(2x-1)^3+x^2\times3(2x-1)^2\times2=2x(2x-1)^3+6x^2(2x-1)^2$}
\end{jibun}

\begin{reidai}[つながった変化の速さ]
正方形の1辺が $a=3t$（cm）でのびていきます。2秒のとき、面積 $S=a^2$ がふえる速さを求めましょう。\par
\kangaekata $S$ は $a$ の関数、$a$ は $t$ の関数です。$\dfrac{\dd S}{\dd t}=\dfrac{\dd S}{\dd a}\times\dfrac{\dd a}{\dd t}$ を使います。\par
\shiki $\dfrac{\dd S}{\dd a}=2a$、$\dfrac{\dd a}{\dd t}=3$ なので、$\dfrac{\dd S}{\dd t}=2a\times3=6a$\par
\hspace*{2.6\zw}$t=2$ で $a=3\times2=6$ なので、$6\times6=36$\par
\kotae \qty{36}{cm^2/s}
\end{reidai}

\begin{jibun}
円の半径が $r=3t$（cm）でのびていきます。2秒のとき、面積 $S=\pi r^2$ がふえる速さを求めましょう。
\bkjans{$2\pi r\times3$、$r=6$ で $36\pi\approx\qty{113}{cm^2/s}$}
\end{jibun}

\begin{reidai}[水位の上がる速さ]
上の円すい形のタンク（$V=\dfrac{\pi}{12}h^3$）に、毎秒\qty{50}{cm^3}の水を入れます。深さ\qty{10}{cm}のとき、水位の上がる速さを求めましょう。\par
\kangaekata $\dfrac{\dd V}{\dd t}=\dfrac{\dd V}{\dd h}\times\dfrac{\dd h}{\dd t}$ に、わかっている数を入れて、$\dfrac{\dd h}{\dd t}$ について解きます。\par
\shiki $\dfrac{\dd V}{\dd h}=\dfrac{\pi}{4}h^2$。$h=10$ で $\dfrac{\pi}{4}\times100=25\pi$\par
\hspace*{2.6\zw}$50=25\pi\times\dfrac{\dd h}{\dd t}$。両辺を $25\pi$ でわって、$\dfrac{\dd h}{\dd t}=\dfrac{50}{25\pi}=\dfrac{2}{\pi}$\par
\kotae $\dfrac{2}{\pi}\approx\qty{0.64}{cm/s}$
\end{reidai}

\begin{jibun}
同じタンクに毎秒\qty{100}{cm^3}の水を入れます。深さ\qty{20}{cm}のとき、水位の上がる速さを求めましょう。
\bkjans{$100=\dfrac{\pi}{4}\times400\times\dfrac{\dd h}{\dd t}$ より $\dfrac{1}{\pi}\approx\qty{0.32}{cm/s}$}
\end{jibun}

\begin{reidai}[振り子の長さと周期]
振り子の周期 $T=2\pi\sqrt{\dfrac{L}{g}}$ について、$L=\qty{1}{m}$ のときの $\dfrac{\dd T}{\dd L}$ を求めましょう（$g=9.8$）。糸を\qty{1}{cm}のばすと、周期はおよそ何秒のびますか。\par
\kangaekata 数学の道具9で求めた $\dfrac{\dd T}{\dd L}=\dfrac{\pi}{\sqrt{gL}}$ に、$L=1$ を入れます。のびる時間は、およそ「変わる速さ×のばした長さ」です（5-8）。\par
\shiki $\dfrac{\dd T}{\dd L}=\dfrac{\pi}{\sqrt{9.8\times1}}=\dfrac{3.14}{3.13}\approx1.0$\par
\hspace*{2.6\zw}\qty{1}{cm}＝\qty{0.01}{m} なので、$\Delta T\approx1.0\times0.01=0.01$\par
\kotae 約1.0（s/m）。約0.01秒のびる
\end{reidai}

\begin{jibun}
$L=\qty{0.64}{m}$ のときの $\dfrac{\dd T}{\dd L}$ を求めましょう（$\sqrt{9.8\times0.64}=2.50$）。
\bkjans{$\dfrac{3.14}{2.50}\approx1.3$（s/m）}
\end{jibun}

\begin{reidai}[分数の形を、積として微分する]
$y=\dfrac{x}{x+1}$ を微分しましょう。\par
\kangaekata $x\times(x+1)^{-1}$ と見て、積の微分と合成関数の微分を使います。$\{(x+1)^{-1}\}'=-(x+1)^{-2}$ です。\par
\shiki $y'=1\times(x+1)^{-1}+x\times\{-(x+1)^{-2}\}=\dfrac{1}{x+1}-\dfrac{x}{(x+1)^2}$\par
\hspace*{2.6\zw}分母を $(x+1)^2$ にそろえる：$\dfrac{x+1}{(x+1)^2}-\dfrac{x}{(x+1)^2}=\dfrac{(x+1)-x}{(x+1)^2}$\par
\kotae $y'=\dfrac{1}{(x+1)^2}$
\end{reidai}

\begin{jibun}
$y=\dfrac{x}{2x+1}$ を微分しましょう。
\bkjans{$(2x+1)^{-1}-2x(2x+1)^{-2}=\dfrac{1}{(2x+1)^2}$}
\end{jibun}

%% ---------------- まちがえやすいところ・物理工学 ----------------
\bksection{tsumazuki}{まちがえやすいところをチェック}

\begin{tsumazuki}{積の微分を「微分どうしの積」にしてしまう}
\bkng{$\{x^2\cdot x^3\}'=(x^2)'(x^3)'=2x\cdot3x^2=6x^3$}\bkok{$2x\cdot x^3+x^2\cdot3x^2=5x^4$}
長方形の図では、ふえる面積は「横がのびた分」と「たてがのびた分」の和でした。微分どうしをかけると、右上の小さな四角だけを見ていることになります。
\end{tsumazuki}

\begin{tsumazuki}{内側の微分をかけ忘れる}
\bkng{$\{(2x-5)^4\}'=4(2x-5)^3$}\bkok{$4(2x-5)^3\times2$}
歯車の比を1つかけ忘れたのと同じです。必ず内側の微分をかけます。
\end{tsumazuki}

\begin{tsumazuki}{答えに $u$ を残してしまう}
\bkng{$\{(3x+1)^2\}'=2u\times3=6u$}\bkok{$6u=6(3x+1)$}
$u$ は、計算のために自分でつけた名前です。最後に $u=3x+1$ をもどして、$x$ の式で答えます。
\end{tsumazuki}

\begin{butsuri}
\begin{itemize}[leftmargin=1.2\zw, itemsep=2pt, topsep=0pt, label={\color{bkPurple}●}]
\item \textbf{変速機と減速機}\\
車や自転車の変速機、ロボットのモーターにつける減速機は、歯車で回る速さの比を変えています。いくつもの歯車をつないだときの比は、合成関数の微分と同じように、かけ算で計算します。
\item \textbf{ロボットの腕}\\
関節のモーターの回る速さから、腕の先が動く速さを求めるのに、合成関数の微分を使います（第7巻）。
\item \textbf{電力}\\
電力は電圧と電流の積 $P=VI$ です。電圧も電流も時間とともに変わるとき、電力の変わる速さは積の微分で計算します。\bkidx{でんりょく}{電力}
\end{itemize}
\end{butsuri}

%% ---------------- 章末問題 ----------------
\bkendproblems
\bkproblemgroup{基本}{4}

\bkproblem $y=(x-1)(x^2+x+1)$ を積の微分で微分し、展開してから微分した結果と同じになることを確かめましょう。
\bkpans{$(x^2+x+1)+(x-1)(2x+1)=3x^2$。展開すると $x^3-1$ なので、微分して $3x^2$}

\bkproblem 次の関数を微分しましょう。\quad（1）$(4x+1)^3$\quad（2）$(1-x)^5$
\bkpans{（1）\nobreak $12(4x+1)^2$\quad（2）\nobreak $-5(1-x)^4$}

\bkproblem $y=\sqrt{x^2+16}$ を微分し、$x=3$ での値を求めましょう。
\bkpans{$y'=\dfrac{x}{\sqrt{x^2+16}}$、$x=3$ で $\dfrac{3}{5}$}

\bkproblem 球の半径が $r=t$（cm）でのびていきます。球の体積は $\mbox{$V=\dfrac{4}{3}\pi r^3$}$ です。3秒のとき、体積がふえる速さを求めましょう。
\bkpans{$4\pi r^2\times1$、$r=3$ で $36\pi\approx\qty{113}{cm^3/s}$}

\bkproblem $y=\dfrac{1}{4x-3}$ を微分しましょう。
\bkpans{$-\dfrac{4}{(4x-3)^2}$}

\bkproblem $y=(x^2+1)^3$ を微分しましょう。
\bkpans{$3(x^2+1)^2\times2x=6x(x^2+1)^2$}

\bkproblem $y=\dfrac{x}{x+2}$ を微分しましょう。
\bkpans{$(x+2)^{-1}-x(x+2)^{-2}=\dfrac{2}{(x+2)^2}$}

\bkproblemgroup{チャレンジ}{4}

\bkproblem $y=x(2x+1)^4$ を微分しましょう。
\bkpans{$(2x+1)^4+8x(2x+1)^3$（まとめると $(2x+1)^3(10x+1)$）}

\bkproblem $y=\dfrac{1}{(3x-1)^2}$ を微分し、$x=1$ での値を求めましょう。
\bkpans{$y=(3x-1)^{-2}$ より $y'=-6(3x-1)^{-3}=-\dfrac{6}{(3x-1)^3}$。$x=1$ で $-\dfrac{3}{4}$}

\bkproblem 長さ $L$ の振り子の周期は $T=2\pi\sqrt{\dfrac{L}{g}}$ です。$g=9.8$ として、$L=\qty{0.25}{m}$ のときの $\dfrac{\dd T}{\dd L}$ を求めましょう（小数第1位まで）。
\bkpans{$\dfrac{\pi}{\sqrt{9.8\times0.25}}=\dfrac{\pi}{1.565}\approx2.0$（s/m）}

\bkproblem 立方体の1辺が $a=2t+1$（cm）でのびていきます。1秒のとき、体積 $V=a^3$ がふえる速さを求めましょう。
\bkpans{$\dfrac{\dd V}{\dd t}=3a^2\times2$。$t=1$ で $a=3$ なので\qty{54}{cm^3/s}}

\bkproblem $y=(x^2+1)(x-2)^2$ を微分しましょう（例題5）。
\bkpans{$2x(x-2)^2+(x^2+1)\times2(x-2)$（まとめると $2(x-2)(2x^2-2x+1)$）}

\bkproblem 数学の道具8の円すい形のタンク（$V=\dfrac{\pi}{12}h^3$）に、毎秒\qty{30}{cm^3}の水を入れます。深さ\qty{6}{cm}のとき、水位の上がる速さを求めましょう。
\bkpans{$30=\dfrac{\pi}{4}\times36\times\dfrac{\dd h}{\dd t}$ より $\dfrac{10}{3\pi}\approx\qty{1.06}{cm/s}$}

\bkproblem $y=(2x-1)^3$ のグラフの、$x=1$ の点での接線の式を求めましょう（5-4 例題6と同じ手順）。
\bkpans{接点 $(1,\ 1)$。$y'=6(2x-1)^2$ より傾き6。$y=6x-5$}

\bkproblem 数学の道具7の風船（半径 $r=2t$ cm）で、3秒のときの体積のふえる速さを求めましょう。1秒のときの何倍ですか。
\bkpans{$4\pi\times6^2\times2=288\pi\approx\qty{905}{cm^3/s}$。1秒（$32\pi$）の9倍}

\begin{furikaeri}
\begin{checklist}
\item 積の微分を、長方形の面積で説明できる。\hfill{\footnotesize\sffamily\color{bkTheme}→\pageref{bk:5-6:seki}ページ}
\item 合成関数の内側と外側を見分けられる。\hfill{\footnotesize\sffamily\color{bkTheme}→\pageref{bk:5-6:gousei}ページ}
\item 合成関数の微分を使って計算できる。\hfill{\footnotesize\sffamily\color{bkTheme}→\pageref{bk:5-6:chain}ページ}
\end{checklist}
\end{furikaeri}

\bkfukidashi{delta}{微分の道具がそろったね。次の章では、微分を使って「いちばん高い所」や「いちばん大きくなる形」を見つけるよ。}

%% 5-7 いちばん高いのはどこ？：最大と最小（第11回：書き足し。ページ数は自由・偶数）
\bkchapter{5}{7}{いちばん高いのはどこ？\\最大と最小}{\bklevel{3}}{いちばん高いのはどこ？　最大と最小}
\bkmichishirube{2-2,3-10,5-5,5-6}{5-8,6-13}

\begin{goal}
\begin{itemize}[label={\color{bkTheme}●}, leftmargin=1.6\zw, itemsep=0pt, topsep=1pt]
\item \bklevel{3}\enspace 導関数の符号から、グラフが上がるか下がるかを判断できる。
\item \bklevel{3}\enspace 増減表で、山（極大）と谷（極小）を見つけられる。
\item \bklevel{3}\enspace 箱の体積や面積がいちばん大きくなる形を、微分で求められる。
\end{itemize}
\end{goal}

\begin{hatena}
\noindent\begin{minipage}[c]{0.6\linewidth}\raggedright
1辺\qty{12}{cm}の正方形の紙の4すみから、同じ大きさの正方形を切りとって折り曲げ、ふたのない箱をつくります。

切りとる正方形を小さくすると、箱は浅くなります。大きくすると、底がせまくなります。体積がいちばん大きくなるのは、何 cm の正方形を切りとったときでしょう。
\end{minipage}\hfill
\begin{minipage}[c]{0.37\linewidth}\centering
\begin{tikzpicture}[x=1mm, y=1mm, scale=0.62]
  \fill[bkYellow!20] (0,0) rectangle (36,36);
  \foreach \cx/\cy in {0/0,27/0,0/27,27/27} \fill[bkGray!40] (\cx,\cy) rectangle ++(9,9);
  \draw[bkInk, line width=0.8pt] (0,0) rectangle (36,36);
  \draw[bkInk, line width=0.5pt, dashed] (9,0) -- (9,36) (27,0) -- (27,36) (0,9) -- (36,9) (0,27) -- (36,27);
  \node[font=\sffamily\scriptsize, anchor=north] at (18,-0.5) {\qty{12}{cm}};
  \node[font=\sffamily\scriptsize] at (4.5,4.5) {$x$};
  \draw[bk arrow, bkGray, line width=0.9pt] (39,18) -- (45,18);
  \begin{scope}[shift={(48,6)}, x={(1mm,0mm)}, y={(0.5mm,0.35mm)}, z={(0mm,1mm)}]
    \draw[bkInk, line width=0.7pt, fill=bkYellow!35] (0,0,0) -- (18,0,0) -- (18,0,9) -- (0,0,9) -- cycle;
    \draw[bkInk, line width=0.7pt, fill=bkYellow!20] (18,0,0) -- (18,18,0) -- (18,18,9) -- (18,0,9) -- cycle;
    \draw[bkInk, line width=0.7pt] (0,0,9) -- (0,18,9) -- (18,18,9);
  \end{scope}
\end{tikzpicture}
\end{minipage}
\end{hatena}

いちばん大きい所や、いちばん小さい所を見つけることは、物をつくるときにいつも必要になります。材料を少なくしたい、たくさん入れたい、いちばん速く着きたい……。この章では、微分を使って、こうした「いちばん」を見つけます。

\clearpage
%% ---------------- やってみよう・思い出そう ----------------
\begin{yattemiyou}[箱をつくって体積をくらべよう]
\par\addvspace{3pt}\textbf{用意するもの}\par\noindent 方眼紙（1辺\qty{12}{cm}の正方形を5まい）、はさみ、テープ
\begin{enumerate}[label={\bkcircled[bkGreen]{\arabic*}}, leftmargin=2.4\zw, itemsep=1pt, topsep=3pt]
\item 切りとる正方形の1辺 $x$ を1、2、3、4、\qty{5}{cm}にして、5つの箱をつくろう。
\item 底の1辺・深さをはかり、体積（底の面積×深さ）を計算して、表にまとめよう。
\item どの箱の体積がいちばん大きいかな？
\end{enumerate}
\end{yattemiyou}

\noindent\begin{minipage}[c]{0.5\linewidth}\raggedright
\textbf{底の1辺は $12-2x$}\quad 右の図のように、紙の1辺\qty{12}{cm}から、左はしで $x$、右はしでも $x$ を切りとります。残りが箱の底の1辺なので、$12-x-x=12-2x$ です。深さは、折り曲げた部分の $x$ です。
\end{minipage}\hfill
\begin{minipage}[c]{0.46\linewidth}\centering
\begin{tikzpicture}[x=1mm, y=1mm]
  \fill[bkYellow!20] (0,0) rectangle (48,48);
  \fill[bkBlue!18] (12,12) rectangle (36,36);
  \foreach \cx/\cy in {0/0,36/0,0/36,36/36} \fill[bkGray!40] (\cx,\cy) rectangle ++(12,12);
  \draw[bkInk, line width=0.8pt] (0,0) rectangle (48,48);
  \draw[bkInk, line width=0.5pt, dashed] (12,0) -- (12,48) (36,0) -- (36,48) (0,12) -- (48,12) (0,36) -- (48,36);
  \draw[bkInk, line width=0.5pt, |<->|] (0,-3) -- (12,-3);
  \draw[bkInk, line width=0.5pt, |<->|] (12,-3) -- (36,-3);
  \draw[bkInk, line width=0.5pt, |<->|] (36,-3) -- (48,-3);
  \node[font=\sffamily\footnotesize, anchor=north] at (6,-4.5) {$x$};
  \node[font=\sffamily\footnotesize, anchor=north] at (24,-4.5) {$12-2x$};
  \node[font=\sffamily\footnotesize, anchor=north] at (42,-4.5) {$x$};
  \node[font=\sffamily\footnotesize, align=center] at (24,24) {箱の底};
  \node[font=\sffamily\scriptsize, align=center] at (24,42) {折って\\かべに};
\end{tikzpicture}
\end{minipage}

\medskip
体積は「底の面積×深さ」なので、
\[ V=(12-2x)\times(12-2x)\times x=x(12-2x)^2 \]
です。$x=2$ なら、底の1辺は $12-4=8$、体積は $8\times8\times2=128$ です。

\begin{bkcenter}
\begin{bktable}{colspec={Q[l,m,wd=34mm] *{5}{Q[c,m,wd=12mm]}}, rowsep=1.5pt}
$x$（cm） & 1 & 2 & 3 & 4 & 5 \\
底の1辺 $12-2x$（cm） & 10 & 8 & 6 & 4 & 2 \\
体積 $V$（\unit{cm^3}） & 100 & 128 & 108 & 64 & 20 \\
\end{bktable}
\bktabcaption{切りとる正方形の1辺と、箱の体積}
\end{bkcenter}

5つの中では、$x=2$ のときがいちばん大きくなりました。でも、$x=1.9$ や $x=2.1$ のほうが大きいかもしれません。電卓で計算すると、$x=1.9$ で $127.756$、$x=2.1$ で $127.764$ で、どちらも128より少しだけ小さくなりました。では、$x=2.01$ では？　表をいくら細かくしても、「ほんとうにいちばん」とは言いきれません。微分を使うと、これをはっきり決められます。

\noindent\textbf{$x$ が使える範囲}\quad 切りとる正方形の1辺 $x$ は0より大きい数です。また、底の1辺 $12-2x$ も0より大きくないと、箱になりません。$12-2x>0$ より $x<6$。だから、$x$ は $0<x<6$ の範囲の数です。

\begin{omoidasou}
\bkrecall{5-3}{$f'(a)$ は、$x=a$ での接線の傾き。$f'(a)=\displaystyle\lim_{h\to0}\frac{f(a+h)-f(a)}{h}$。}
\bkrecall{5-4}{$(x^n)'=nx^{n-1}$。定数倍・和は1つずつ微分する。}
\bkrecall{5-5}{$f'(x)$ をもう一度微分したものが $f''(x)$（2回微分）。位置を2回微分すると加速度。}
\bkrecall{5-6}{積の微分 $(fg)'=f'g+fg'$。$\{(ax+b)^n\}'=n(ax+b)^{n-1}\times a$。}
\bkrecall{2-2}{$a>0$ は正、$a<0$ は負。$(+)\times(-)=(-)$、$(-)\times(-)=(+)$。}
\bkrecall{3-8}{$(a-b)^2=a^2-2ab+b^2$。たとえば $(12-2x)^2=144-48x+4x^2$。}
\bkrecall{3-10}{和が $-8$、積が12になる2つの数は $-2$ と $-6$ なので、\\$x^2-8x+12=(x-2)(x-6)$。$=0$ になる $x$ は2と6。また、$x^2-1=(x+1)(x-1)$。}
\end{omoidasou}

%% ---------------- 数学の道具 ----------------
\Needspace{8\baselineskip}
\bksection{doguu}{数学の道具}

\bksubsection{1}{導関数の符号で、上がる・下がるがわかる}

\label{bk:5-7:fugou}$f(x)=x^3-3x$ のグラフで考えましょう。グラフの上のいろいろな点に、接線を引いてみます。

\noindent\begin{minipage}[c]{0.5\linewidth}\centering
\begin{tikzpicture}
\begin{axis}[bkplot, width=62mm, height=54mm, xmin=-2.3, xmax=2.5, ymin=-3, ymax=3.6, xtick={-1,1}, ytick={-2,2}, xlabel={$x$}, ylabel={$y$}]
  \addplot[bkcurve, domain=-2.1:2.1] {x^3-3*x};
  \foreach \a in {-1.7,1.7} {\edef\tmp{\noexpand\addplot[bkGreen!80!black, line width=1.1pt, domain=\a-0.18:\a+0.18] {\a^3-3*\a+(3*\a*\a-3)*(x-\a)};}\tmp}
  \foreach \a in {-0.5,0.5} {\edef\tmp{\noexpand\addplot[bkRed, line width=1.1pt, domain=\a-0.4:\a+0.4] {\a^3-3*\a+(3*\a*\a-3)*(x-\a)};}\tmp}
  \node[font=\sffamily\scriptsize, text=bkGreen!60!black, anchor=east] at (axis cs:-1.95,-1.3) {傾き＋};
  \node[font=\sffamily\scriptsize, text=bkRed!80!black, anchor=south west] at (axis cs:0.3,0.6) {傾き－};
  \node[font=\sffamily\scriptsize, text=bkGreen!60!black, anchor=west] at (axis cs:1.95,-1.3) {傾き＋};
\end{axis}
\end{tikzpicture}
\end{minipage}\hfill
\begin{minipage}[c]{0.46\linewidth}\raggedright
グラフが右上がりの所（緑）では、接線も右上がりで、傾きは正です。右下がりの所（赤）では、接線も右下がりで、傾きは負です。接線の傾きは $f'(x)$ でした（5-3）。
\end{minipage}

\medskip
だから、次のことが言えます。
\begin{itemize}[leftmargin=1.6\zw, itemsep=1pt, topsep=2pt, label={\color{bkTheme}●}]
\item $f'(x)>0$ の所では、$x$ がふえると $f(x)$ もふえる（\textbf{増加}）。
\item $f'(x)<0$ の所では、$x$ がふえると $f(x)$ はへる（\textbf{減少}）。
\end{itemize}
\bkidx{ぞうか・げんしょう}{増加・減少}

\noindent\textbf{数で確かめる}\quad $f'(x)=3x^2-3$ です。
\begin{itemize}[leftmargin=1.6\zw, itemsep=1pt, topsep=2pt, label={\color{bkTheme}●}]
\item $x=1.7$ では $f'(1.7)=3\times2.89-3=5.67>0$。\\$x$ を少しふやすと、$f(1.7)=-0.187$ → $f(1.8)=0.432$ で、ふえている。
\item $x=0.5$ では $f'(0.5)=0.75-3=-2.25<0$。\\$x$ を少しふやすと、$f(0.5)=-1.375$ → $f(0.6)=-1.584$ で、へっている。
\end{itemize}

\noindent\textbf{なぜそう言えるのか}\quad $f'(a)$ は、$\dfrac{f(a+h)-f(a)}{h}$ の、$h$ を0に近づけたときの近づく先でした（5-3）。$f'(a)$ が正なら、$h$ が小さい正の数のとき、この分数も正です。分母の $h$ が正なので、分子 $f(a+h)-f(a)$ も正、つまり $f(a+h)>f(a)$ です。少し右へ行くと、値が大きくなるのです。

\noindent\textbf{例}\quad 投げ上げたボールの高さ $x=20t-5t^2$ なら、$v=\dot{x}=20-10t$（5-5）。$t<2$ では $v>0$ で上がり、$t>2$ では $v<0$ で下がります。上がるのから下がるのに変わる $t=2$（$v=0$）で、いちばん高くなります。

\bksubsection{2}{増減表}

\label{bk:5-7:zougen}$f(x)=x^3-3x$ が、どこで上がり、どこで下がるかを、表にまとめましょう。次の順に調べます。\bkidx{ぞうげんひょう}{増減表}

\noindent\textbf{手順1　微分して、$f'(x)=0$ となる $x$ を求める}
\[ f'(x)=3x^2-3=3(x^2-1)=3(x+1)(x-1) \]
（3でくくり、$x^2-1=(x+1)(x-1)$ と因数分解した。3-10）。$f'(x)=0$ になるのは、$x+1=0$ か $x-1=0$ のときなので、$x=-1$ と $x=1$ です。

\noindent\textbf{手順2　$f'(x)$ の符号を調べる}\quad $f'(x)$ は、3と $(x+1)$ と $(x-1)$ のかけ算です。それぞれの符号を表にすると、かけ算の符号がわかります（2-2）。

\begin{bkcenter}
\begin{tblr}{colspec={Q[c,m,wd=15mm] Q[c,m,wd=18mm] Q[c,m,wd=12mm] Q[c,m,wd=24mm] Q[c,m,wd=12mm] Q[c,m,wd=18mm]}, colsep=4pt, hlines={0.6pt,bkLine}, vline{2}={0.6pt,bkLine}, rowsep=2pt, row{1}={bg=bkTheme!10, font=\small}}
$x$ & $x<-1$ & $-1$ & $-1<x<1$ & $1$ & $x>1$ \\
$x+1$ & $-$ & $0$ & $+$ & $+$ & $+$ \\
$x-1$ & $-$ & $-$ & $-$ & $0$ & $+$ \\
$f'(x)$ & $+$ & $0$ & $-$ & $0$ & $+$ \\
\end{tblr}
\end{bkcenter}

\noindent たとえば $x<-1$ では、$(x+1)$ も $(x-1)$ も負なので、$3\times(-)\times(-)=(+)$ です。数を入れて確かめると、$f'(-2)=3\times(-1)\times(-3)=9>0$、$f'(0)=3\times1\times(-1)=-3<0$、$f'(2)=3\times3\times1=9>0$ で、表と合っています。

\noindent\textbf{手順3　表にまとめ、$f(x)$ の値を入れる}\quad 上がる所は $\nearrow$、下がる所は $\searrow$ と書きます。$f(-1)=(-1)^3-3\times(-1)=-1+3=2$、$f(1)=1-3=-2$ です。このような表を\textbf{増減表}といいます。

\begin{bkcenter}
\begin{tblr}{colspec={Q[c,m,wd=14mm] *{5}{Q[c,m,wd=16mm]}}, hlines={0.6pt,bkLine}, vline{2}={0.6pt,bkLine}, rowsep=2pt, row{1}={bg=bkTheme!10}}
$x$ & $\cdots$ & $-1$ & $\cdots$ & $1$ & $\cdots$ \\
$f'(x)$ & $+$ & $0$ & $-$ & $0$ & $+$ \\
$f(x)$ & $\nearrow$ & {{\bfseries 2}\\{\scriptsize 極大}} & $\searrow$ & {{\bfseries $-2$}\\{\scriptsize 極小}} & $\nearrow$ \\
\end{tblr}
\end{bkcenter}

\noindent\begin{minipage}[c]{0.42\linewidth}\centering
\begin{tikzpicture}
\begin{axis}[bkplot, width=54mm, height=44mm, xmin=-2.3, xmax=2.5, ymin=-3, ymax=3.4, xtick={-1,1}, ytick={-2,2}, xlabel={$x$}, ylabel={$y$}]
  \addplot[bkcurve, domain=-2.1:2.1] {x^3-3*x};
  \addplot[bkcurve2, line width=1pt, domain=-1.5:-0.6] {2};
  \addplot[bkcurve2, line width=1pt, domain=0.6:1.5] {-2};
  \addplot[only marks, mark=*, mark size=1.8pt, bkInk] coordinates {(-1,2) (1,-2)};
\end{axis}
\end{tikzpicture}
\end{minipage}\hfill
\begin{minipage}[c]{0.55\linewidth}\raggedright
$x=-1$ で、上がるのから下がるのに変わるので、そこが山のてっぺんです。これを\textbf{極大}（きょくだい）といい、そのときの値 $f(-1)=2$ を\textbf{極大値}といいます。$x=1$ は谷の底で\textbf{極小}（きょくしょう）、$f(1)=-2$ が\textbf{極小値}です。極大・極小では、接線が横向き（傾き0、オレンジ）になっています。\bkidx{きょくだい・きょくしょう}{極大・極小}
\end{minipage}

\medskip
\noindent\textbf{「極大」は「いちばん大きい」ではない}\quad 極大値2は、$x=-1$ の\textbf{近くの中で}いちばん大きい値です。グラフ全体では、右へ行くと $f(3)=27-9=18$ のように、2よりずっと大きくなります。山のてっぺんが、全体でいちばん高いとはかぎらないのです。

\bksubsection{3}{範囲が決まっているときの最大と最小}

\label{bk:5-7:saidai}箱の問題では、$x$ は $0<x<6$ の数しか使えませんでした。このように範囲が決まっているときは、極大・極小のほかに、\textbf{範囲のはしの値}もくらべます。山のてっぺんより、はしのほうが高いことがあるからです。

\noindent\begin{minipage}[c]{0.52\linewidth}\raggedright
\textbf{例}\quad $f(x)=x^3-3x$ を、$0\leqq x\leqq3$ の範囲で考えます（$\leqq$ は「小さいか等しい」。2-2）。この範囲にある $f'(x)=0$ の点は $x=1$ だけです（$x=-1$ は範囲の外）。はしの $x=0$、$x=3$ と、$x=1$ の値をくらべます。
\[ f(0)=0,\quad f(1)=-2,\quad f(3)=18 \]
最大値は、はしの $x=3$ のときの18、最小値は $x=1$ のときの $-2$ です。\bkidx{さいだいち}{最大値}\bkidx{さいしょうち}{最小値}
\end{minipage}\hfill
\begin{minipage}[c]{0.44\linewidth}\centering
\begin{tikzpicture}
\begin{axis}[bkplot, width=54mm, height=56mm, xmin=-0.4, xmax=3.6, ymin=-3.5, ymax=21, xtick={1,3}, ytick={-2,18}, xlabel={$x$}, ylabel={$y$}]
  \addplot[bkGray, line width=0.8pt, dashed, domain=-0.3:0] {x^3-3*x};
  \addplot[bkcurve, domain=0:3] {x^3-3*x};
  \addplot[only marks, mark=*, mark size=1.8pt, bkRed] coordinates {(3,18)};
  \addplot[only marks, mark=*, mark size=1.8pt, bkBlue] coordinates {(1,-2) (0,0)};
  \draw[bkGray, dashed, line width=0.5pt] (axis cs:0,18) -- (axis cs:3,18);
  \node[font=\sffamily\scriptsize, anchor=east, text=bkRed!80!black] at (axis cs:2.85,18) {最大（はし）};
  \node[font=\sffamily\scriptsize, anchor=north, text=bkBlue] at (axis cs:1,-2.4) {最小};
\end{axis}
\end{tikzpicture}
\end{minipage}

\bksubsection{4}{はてな？の箱を解く}

\label{bk:5-7:hako}\bkhatenaref の箱の体積は $V=x(12-2x)^2$（$0<x<6$）でした。展開してから微分します。

\noindent\textbf{手順1　展開する}\quad $(12-2x)^2=144-48x+4x^2$（3-8）。これに $x$ をかけます。
\[ V=x(144-48x+4x^2)=144x-48x^2+4x^3 \]
\textbf{手順2　微分する}\quad $(x^n)'=nx^{n-1}$ を1つずつ使います（5-4）。
\[ V'=144-96x+12x^2=12x^2-96x+144 \]
\textbf{手順3　$V'=0$ を解く}\quad どの項も12でわれるので、12でくくってから因数分解します（3-10）。
\[ V'=12(x^2-8x+12)=12(x-2)(x-6) \]
$V'=0$ になるのは $x=2$ と $x=6$。でも、$x=6$ は範囲 $0<x<6$ に入りません（底の1辺が0で、箱にならない）。だから、調べるのは $x=2$ だけです。

\noindent\textbf{手順4　増減表}\quad $x=1$ では $V'=12\times(-1)\times(-5)=60>0$、$x=3$ では $V'=12\times1\times(-3)=-36<0$。2の前で $+$、後で $-$ です。

\begin{bkcenter}
\begin{tblr}{colspec={Q[c,m,wd=14mm] *{5}{Q[c,m,wd=14mm]}}, hlines={0.6pt,bkLine}, vline{2}={0.6pt,bkLine}, rowsep=2pt, row{1}={bg=bkTheme!10}}
$x$ & $0$ & $\cdots$ & $2$ & $\cdots$ & $6$ \\
$V'$ & & $+$ & $0$ & $-$ & \\
$V$ & & $\nearrow$ & {{\bfseries 128}\\{\scriptsize 最大}} & $\searrow$ & \\
\end{tblr}
\end{bkcenter}
\noindent\begin{minipage}[c]{0.4\linewidth}\raggedright
$x=2$ までは体積がふえ、2をこえるとへります。だから、$x=\qty{2}{cm}$ のとき体積は最大で、$2\times8^2=\qty{128}{cm^3}$ です。やってみようの表で見た2が、ちょうどいちばんでした。これが、\bkhatenaref の答えです。
\end{minipage}\hfill
\begin{minipage}[c]{0.57\linewidth}\centering
\begin{tikzpicture}
\begin{axis}[bkplot, width=62mm, height=42mm, xmin=0, xmax=6.4, ymin=0, ymax=150, xtick={1,...,6}, ytick={50,100,128}, xlabel={$x$（cm）}, ylabel={$V$（\unit{cm^3}）}]
  \addplot[bkcurve, domain=0:6] {x*(12-2*x)^2};
  \addplot[only marks, mark=*, mark size=1.6pt, bkInk] coordinates {(1,100) (2,128) (3,108) (4,64) (5,20)};
  \addplot[bkcurve2, line width=1pt, domain=1.2:2.8] {128};
\end{axis}
\end{tikzpicture}
\bkfigcaption{箱の体積。$x=2$ で接線が横向きになる。}
\end{minipage}

\bksubsection{5}{2回微分で、山と谷を見分ける}

\label{bk:5-7:nikai}\noindent\begin{minipage}[c]{0.68\linewidth}\raggedright
山のてっぺんの近くでは、接線の傾きは「正 → 0 → 負」と、だんだん小さくなっていきます。傾き $f'(x)$ がへっていくので、$f'(x)$ を微分した $f''(x)$（2回微分、5-5）は負です。谷の底では反対に、傾きは「負 → 0 → 正」とふえていくので、$f''(x)$ は正です。
\end{minipage}\hfill
\begin{minipage}[c]{0.29\linewidth}\centering
\begin{tikzpicture}[x=0.66mm, y=1mm]
  \draw[bkBlue, line width=1.5pt] (0,0) parabola bend (12,10) (24,0);
  \foreach \xx/\sl in {4/1.1,12/0,20/-1.1} \draw[bkOrange, line width=0.8pt] ({\xx-3},{10-(\xx-12)^2*10/144-3*\sl}) -- ({\xx+3},{10-(\xx-12)^2*10/144+3*\sl});
  \node[font=\sffamily\scriptsize, align=center] at (12,-6) {傾き\\正→0→負\\$f''<0$};
  \draw[bkBlue, line width=1.5pt] (30,10) parabola bend (42,0) (54,10);
  \node[font=\sffamily\scriptsize, align=center] at (42,-6) {傾き\\負→0→正\\$f''>0$};
\end{tikzpicture}
\end{minipage}

\medskip
\noindent\textbf{数で確かめる}\quad $f(x)=x^3-3x$ の山 $x=-1$ の近くで、傾き $f'(x)=3x^2-3$ を計算すると、$f'(-1.5)=3.75$、$f'(-1)=0$、$f'(-0.5)=-2.25$ と、へっていきます。

\begin{itemize}[leftmargin=1.6\zw, itemsep=1pt, topsep=2pt, label={\color{bkTheme}●}]
\item $f'(a)=0$、$f''(a)<0$ なら、$x=a$ は極大（山）
\item $f'(a)=0$、$f''(a)>0$ なら、$x=a$ は極小（谷）
\end{itemize}
$f(x)=x^3-3x$ なら $f''(x)=6x$。$f''(-1)=-6<0$ なので $x=-1$ は極大、$f''(1)=6>0$ なので $x=1$ は極小。増減表と同じ結果です。投げ上げたボールでは、加速度 $\ddot{x}=-10<0$ なので、$\dot{x}=0$ の点はいつも山（最高点）になります。

\bksubsection{6}{「いちばん」を求める手順}

\label{bk:5-7:tejun}文章の問題で最大・最小を求めるときは、次の順に進めます。

\begin{bkfigure}
\begin{tikzpicture}[x=1mm, y=1mm]
  \node[bk node, text width=34mm] (t1) at (0,12) {\textbf{① 式にする}\\大きく（小さく）したい\\量を、$x$ の式で表す};
  \node[bk node, text width=34mm] (t2) at (46,12) {\textbf{② 範囲を決める}\\$x$ が使える範囲\\（長さは正、など）};
  \node[bk node, text width=34mm] (t3) at (92,12) {\textbf{③ 微分して $=0$}\\$f'(x)=0$ を解き、\\範囲の中の解をえらぶ};
  \node[bk node, text width=34mm] (t4) at (23,-14) {\textbf{④ 増減表}\\$f'$ の符号を調べ、\\上がる・下がるを書く};
  \node[bk node, text width=34mm] (t5) at (69,-14) {\textbf{⑤ くらべる}\\極大・極小と、はしの\\値をくらべて決める};
  \draw[bk link] (t1) -- (t2);
  \draw[bk link] (t2) -- (t3);
  \draw[bk link] (t3.south) -- (t4.north east);
  \draw[bk link] (t4) -- (t5);
\end{tikzpicture}
\bkfigcaption{最大・最小の問題の手順}
\end{bkfigure}

\begin{taisetsu}
$f'(x)>0$ なら増加、$f'(x)<0$ なら減少。\\
$f'(x)=0$ の前後で符号が変われば、そこが極大・極小。\\
範囲が決まっているときは、はしの値もくらべて最大・最小を決める。
\end{taisetsu}

\begin{kotoba}
\small
\bkword{増減表}{$f'(x)$ の符号と、$f(x)$ が上がるか下がるかをまとめた表。}
\bkword{極大・極小}{増加から減少に変わる所（山）と、減少から増加に変わる所（谷）。その近くでいちばん大きい・小さい所。}
\bkword{最大値・最小値}{決められた範囲の中で、いちばん大きい値・いちばん小さい値。}
\end{kotoba}

\begin{column}{表計算で「いちばん」をさがす}
式が複雑で $f'(x)=0$ が解けないときは、表計算で $x$ を細かくふって値をくらべます。箱の体積なら、$x$ を0.1きざみにしても0.01きざみにしても、$x=2$ で最大だとわかります。工学では、このように計算機で最大・最小をさがす方法（最適化）がよく使われます。そのもとにあるのも、「傾きが0の所をさがす」考え方です。
\end{column}

%% ---------------- 例題 ----------------
\bksection{reidai}{例題でたしかめよう}

\begin{reidai}[いちばん高い所]
真上に投げたボールの高さが $x=20t-5t^2$（m）のとき、最も高くなる時刻と高さを求めましょう。\par
\kangaekata いちばん高い所では、上がるのから下がるのに変わるので、速度 $\dot{x}$ が0になります（①）。\par
\shiki $\dot{x}=20-5\times2t=20-10t$\par
\hspace*{2.6\zw}$20-10t=0$ より $10t=20$、$t=2$\par
\hspace*{2.6\zw}$t<2$ で $\dot{x}>0$（上がる）、$t>2$ で $\dot{x}<0$（下がる）\par
\hspace*{2.6\zw}$x=20\times2-5\times2^2=40-20=20$\par
\kotae 2秒後に\qty{20}{m}
\end{reidai}

\begin{jibun}
高さが $x=30t-5t^2$（m）のボールの、最も高くなる時刻と高さを求めましょう。
\bkjans{$\dot{x}=30-10t=0$ より3秒後、$90-45=\qty{45}{m}$}
\end{jibun}

\begin{reidai}[増減表と極値]
$f(x)=x^3-6x^2+9x$ の増減表をかき、極大値と極小値を求めましょう。\par
\kangaekata ②の手順1〜3で進めます。\par
\shiki $f'(x)=3x^2-12x+9=3(x^2-4x+3)=3(x-1)(x-3)$\par
\hspace*{2.6\zw}$f'=0$ は $x=1,\,3$\par
\hspace*{2.6\zw}符号：$f'(0)=9>0$、$f'(2)=-3<0$、$f'(4)=9>0$\par
\hspace*{2.6\zw}$f(1)=1-6+9=4$、$f(3)=27-54+27=0$\par
\begin{bkcenter}
\begin{tblr}{colspec={Q[c,m,wd=14mm] *{5}{Q[c,m,wd=13mm]}}, hlines={0.6pt,bkLine}, vline{2}={0.6pt,bkLine}, rowsep=1.5pt, row{1}={bg=bkTheme!10}}
$x$ & $\cdots$ & $1$ & $\cdots$ & $3$ & $\cdots$ \\
$f'(x)$ & $+$ & $0$ & $-$ & $0$ & $+$ \\
$f(x)$ & $\nearrow$ & {{\bfseries 4}\\{\scriptsize 極大}} & $\searrow$ & {{\bfseries 0}\\{\scriptsize 極小}} & $\nearrow$ \\
\end{tblr}
\end{bkcenter}
\kotae $x=1$ で極大値4、$x=3$ で極小値0
\end{reidai}

\begin{jibun}
$f(x)=x^3-12x$ の増減表をかき、極大値と極小値を求めましょう。
\bkjans{$f'=3x^2-12=3(x+2)(x-2)$。$x=-2$ で極大値16、$x=2$ で極小値 $-16$}
\end{jibun}

\begin{reidai}[範囲のある最大・最小（2次の式）]
$f(x)=x^2-4x+5$ の、$0\leqq x\leqq5$ での最大値と最小値を求めましょう。\par
\kangaekata 範囲の中の $f'(x)=0$ の点と、はしの $x=0$、$x=5$ の値をくらべます（③）。\par
\shiki $f'(x)=2x-4=0$ より $x=2$（範囲の中）\par
\hspace*{2.6\zw}$x<2$ で $f'<0$、$x>2$ で $f'>0$ なので、$x=2$ は谷\par
\hspace*{2.6\zw}$f(0)=5$、$f(2)=4-8+5=1$、$f(5)=25-20+5=10$\par
\begin{bkfigure}
\begin{tikzpicture}
\begin{axis}[bkplot, width=62mm, height=44mm, xmin=-0.3, xmax=5.8, ymin=0, ymax=12, xtick={2,5}, ytick={1,5,10}, xlabel={$x$}, ylabel={$y$}]
  \addplot[bkcurve, domain=0:5] {x^2-4*x+5};
  \addplot[only marks, mark=*, mark size=1.8pt, bkRed] coordinates {(5,10)};
  \addplot[only marks, mark=*, mark size=1.8pt, bkBlue] coordinates {(2,1)};
  \addplot[only marks, mark=*, mark size=1.5pt, bkInk] coordinates {(0,5)};
  \node[font=\sffamily\scriptsize, anchor=east, text=bkRed!80!black] at (axis cs:4.8,10.6) {最大（はし）};
  \node[font=\sffamily\scriptsize, anchor=north, text=bkBlue] at (axis cs:2,0.8) {最小（谷）};
\end{axis}
\end{tikzpicture}
\end{bkfigure}
\kotae 最大値10（$x=5$）、最小値1（$x=2$）
\end{reidai}

\begin{jibun}
$f(x)=x^2-6x+10$ の、$0\leqq x\leqq5$ での最大値と最小値を求めましょう。
\bkjans{$f'=2x-6=0$ より $x=3$。$f(0)=10$、$f(3)=1$、$f(5)=5$。最大値10（$x=0$）、最小値1（$x=3$）}
\end{jibun}

\begin{reidai}[範囲のある最大・最小（3次の式）]
$f(x)=x^3-3x^2-9x$ の、$-2\leqq x\leqq4$ での最大値と最小値を求めましょう。\par
\kangaekata 極大・極小が2つとも範囲に入るときは、それと、はしの2つの、合わせて4つの値をくらべます。\par
\shiki $f'(x)=3x^2-6x-9=3(x^2-2x-3)=3(x+1)(x-3)$\par
\hspace*{2.6\zw}$f'=0$ は $x=-1,\,3$（どちらも範囲の中）\par
\hspace*{2.6\zw}$f(-2)=-8-12+18=-2$\par
\hspace*{2.6\zw}$f(-1)=-1-3+9=5$（極大）\par
\hspace*{2.6\zw}$f(3)=27-27-27=-27$（極小）\par
\hspace*{2.6\zw}$f(4)=64-48-36=-20$\par
\begin{bkfigure}
\begin{tikzpicture}
\begin{axis}[bkplot, width=70mm, height=48mm, xmin=-2.5, xmax=4.8, ymin=-31, ymax=9, xtick={-1,3,4}, ytick={-27,5}, xlabel={$x$}, ylabel={$y$}]
  \addplot[bkcurve, domain=-2:4] {x^3-3*x^2-9*x};
  \addplot[only marks, mark=*, mark size=1.8pt, bkRed] coordinates {(-1,5)};
  \addplot[only marks, mark=*, mark size=1.8pt, bkBlue] coordinates {(3,-27)};
  \addplot[only marks, mark=*, mark size=1.5pt, bkInk] coordinates {(-2,-2) (4,-20)};
  \node[font=\sffamily\scriptsize, anchor=south, text=bkRed!80!black] at (axis cs:-1,5.6) {最大};
  \node[font=\sffamily\scriptsize, anchor=west, text=bkBlue] at (axis cs:3.2,-28) {最小};
\end{axis}
\end{tikzpicture}
\end{bkfigure}
\kotae 最大値5（$x=-1$）、最小値 $-27$（$x=3$）
\end{reidai}

\begin{jibun}
$f(x)=x^3-12x$ の、$-3\leqq x\leqq3$ での最大値と最小値を求めましょう。
\bkjans{$f(-3)=9$、$f(-2)=16$、$f(2)=-16$、$f(3)=-9$。最大値16（$x=-2$）、最小値 $-16$（$x=2$）}
\end{jibun}

\begin{reidai}[囲む面積を最大にする]
（1）長さ\qty{40}{m}のロープで長方形を囲みます。面積がいちばん大きくなるのは、たてと横が何mのときでしょう。\par
（2）まっすぐなかべを1辺にして、長さ\qty{40}{m}のロープで残りの3辺を囲みます。かべと直角な辺を $x$ m として、面積が最大になる $x$ と面積を求めましょう。\par
\begin{bkfigure}
\begin{tikzpicture}[x=1mm, y=1mm]
  \fill[bkGreen!15] (0,0) rectangle (30,20);
  \draw[bkOrange!80!black, line width=1.2pt] (0,0) rectangle (30,20);
  \node[font=\sffamily\footnotesize, anchor=east] at (-1,10) {$x$};
  \node[font=\sffamily\footnotesize, anchor=north] at (15,-1) {$20-x$};
  \node[font=\sffamily\footnotesize, anchor=south] at (15,21) {（1）};
  \begin{scope}[shift={(55,0)}]
    \fill[bkGreen!15] (0,0) rectangle (40,20);
    \fill[pattern=north east lines, pattern color=bkGray] (-1,20) rectangle (41,23);
    \draw[bkInk, line width=1.2pt] (-1,20) -- (41,20);
    \draw[bkOrange!80!black, line width=1.2pt] (0,20) -- (0,0) -- (40,0) -- (40,20);
    \node[font=\sffamily\footnotesize, anchor=east] at (-1,10) {$x$};
    \node[font=\sffamily\footnotesize, anchor=north] at (20,-1) {$40-2x$};
    \node[font=\sffamily\footnotesize, anchor=south] at (20,24) {（2）かべ};
  \end{scope}
\end{tikzpicture}
\end{bkfigure}
\kangaekata（1）たて＋横は、ロープの半分の\qty{20}{m}。横は $20-x$。（2）ロープは $x$ が2本と横が1本なので、横は $40-2x$。\par
\shiki（1）$S=x(20-x)=20x-x^2$、$S'=20-2x=0$ より $x=10$。前後で $+$ から $-$\par
\hspace*{2.6\zw}（2）$S=x(40-2x)=40x-2x^2$、$S'=40-4x=0$ より $x=10$。前後で $+$ から $-$\par
\kotae（1）たて・横とも\qty{10}{m}（正方形）、\qty{100}{m^2}\par
\hspace*{2.6\zw}（2）$x=\qty{10}{m}$（横\qty{20}{m}）、\qty{200}{m^2}
\end{reidai}

\begin{jibun}
（1）長さ\qty{24}{m}のロープで長方形を囲みます。面積が最大になるたて・横と面積を求めましょう。\par
（2）かべを1辺にして、長さ\qty{24}{m}のロープで残りの3辺を囲みます。かべと直角な辺を $x$ として、面積が最大になる $x$ と面積を求めましょう。
\bkjans{（1）\nobreak $S=x(12-x)$、$x=6$。たて・横とも\qty{6}{m}、\qty{36}{m^2}\quad（2）\nobreak $S=x(24-2x)$、$S'=24-4x=0$ より $x=\qty{6}{m}$、$6\times12=\qty{72}{m^2}$}
\end{jibun}

\begin{reidai}[箱の体積の最大]
1辺\qty{18}{cm}の正方形の紙の4すみから1辺 $x$ の正方形を切りとって、ふたのない箱をつくります。体積が最大になる $x$ と体積を求めましょう。\par
\kangaekata $V=x(18-2x)^2$（$0<x<9$）。展開してもできますが、積の微分と合成関数の微分（5-6）を使うと、展開しなくてすみます。$x$ と $(18-2x)^2$ のかけ算と見ます。\par
\shiki $(x)'=1$、$\{(18-2x)^2\}'=2(18-2x)\times(-2)$\par
\hspace*{2.6\zw}$V'=1\times(18-2x)^2+x\times2(18-2x)\times(-2)$\par
\hspace*{2.6\zw}$\phantom{V'}=(18-2x)^2-4x(18-2x)$\par
\hspace*{2.6\zw}$\phantom{V'}=(18-2x)\{(18-2x)-4x\}=(18-2x)(18-6x)$\par
\hspace*{2.6\zw}$0<x<9$ で $V'=0$ は $x=3$（$18-2x=0$ の $x=9$ は範囲外）\par
\hspace*{2.6\zw}$x=2$ で $V'=14\times6>0$、$x=4$ で $V'=10\times(-6)<0$\par
\hspace*{2.6\zw}$V=3\times(18-6)^2=3\times144=432$\par
\kotae $x=\qty{3}{cm}$、\qty{432}{cm^3}
\end{reidai}

\begin{jibun}
1辺\qty{24}{cm}の正方形の紙で、同じように箱をつくります。体積が最大になる $x$ と体積を求めましょう。
\bkjans{$V'=(24-2x)^2-4x(24-2x)=(24-2x)(24-6x)$ より $x=\qty{4}{cm}$、$4\times16^2=\qty{1024}{cm^3}$}
\end{jibun}

\begin{reidai}[材料がいちばん少ない缶]
体積が $54\pi$ \unit{cm^3} の、ふたのある円柱の缶をつくります。表面積が最小になる半径 $r$ と高さを求めましょう。\par
\noindent\begin{minipage}[c]{0.5\linewidth}\raggedright
\kangaekata 缶を切り開くと、右の図のように、円が2まい（ふたと底）と、長方形が1まい（側面）になります。長方形の横は、円のまわりの長さ $2\pi r$ です。
\end{minipage}\hfill
\begin{minipage}[c]{0.46\linewidth}\centering
\begin{tikzpicture}[x=1mm, y=1mm]
  \fill[bkBlue!20, draw=bkInk, line width=0.5pt] (0,0) rectangle (44,16);
  \fill[bkYellow!40, draw=bkInk, line width=0.5pt] (22,23) circle (7);
  \fill[bkYellow!40, draw=bkInk, line width=0.5pt] (22,-7) circle (7);
  \draw[bkInk, line width=0.5pt] (22,23) -- (29,23);
  \node[font=\sffamily\scriptsize, anchor=south] at (25.5,23) {$r$};
  \node[font=\sffamily\footnotesize, align=center] at (22,8) {側面（横 $2\pi r$）};
  \node[font=\sffamily\footnotesize, anchor=west] at (44.5,8) {$h$};
\end{tikzpicture}
\end{minipage}\par
\hspace*{2.6\zw}体積 $\pi r^2\times h=54\pi$ より、高さは $h=\dfrac{54}{r^2}$。表面積は\par
\hspace*{2.6\zw}$S=2\times\pi r^2+2\pi r\times h=2\pi r^2+2\pi r\times\dfrac{54}{r^2}=2\pi r^2+\dfrac{108\pi}{r}$\par
\shiki $\dfrac{108\pi}{r}=108\pi r^{-1}$ の微分は $-108\pi r^{-2}$（5-4）なので\par
\hspace*{2.6\zw}$S'=4\pi r-\dfrac{108\pi}{r^2}=0$\par
\hspace*{2.6\zw}両辺に $r^2$ をかけて $4\pi r^3=108\pi$、$r^3=27$、$r=3$\par
\hspace*{2.6\zw}$r=2$ で $S'<0$、$r=4$ で $S'>0$ なので最小。$h=\dfrac{54}{9}=6$\par
\kotae 半径\qty{3}{cm}、高さ\qty{6}{cm}（高さが直径と等しい）
\end{reidai}

\begin{jibun}
体積が $128\pi$ \unit{cm^3} の、ふたのある円柱の缶で、表面積が最小になる半径と高さを求めましょう。
\bkjans{$h=\dfrac{128}{r^2}$、$S=2\pi r^2+\dfrac{256\pi}{r}$。$S'=4\pi r-\dfrac{256\pi}{r^2}=0$ より $r^3=64$、$r=\qty{4}{cm}$、高さ\qty{8}{cm}}
\end{jibun}

\begin{reidai}[電池から取り出す電力を最大にする]
電圧\qty{12}{V}の電池の中には、\qty{2}{\ohm}の抵抗（内部抵抗）があります。この電池に抵抗 $R$ の器具をつなぎます。器具で使われる電力 $P$ が最大になる $R$ と、そのときの電力を求めましょう。\par
\noindent\begin{minipage}[c]{0.56\linewidth}\raggedright
\kangaekata 電流 $I$ は、内部抵抗と $R$ の両方を通ります。直列につないだ抵抗はたし算でまとめられるので、全体の抵抗は $R+2$。$I=\dfrac{V}{R}$（2-15）より、$I=\dfrac{12}{R+2}$。電力は「電圧×電流」で、器具にかかる電圧は $RI$（2-15）なので、
\end{minipage}\hfill
\begin{minipage}[c]{0.4\linewidth}\centering
\begin{tikzpicture}[x=1mm, y=1mm]
  \draw[bkGray, dashed, line width=0.5pt, rounded corners=1.5mm] (-6,-6) rectangle (6,20);
  \node[font=\sffamily\scriptsize, text=bkGray, anchor=north] at (0,-12) {電池};
  \draw[bkInk, line width=0.8pt] (0,-6) -- (0,1) (0,4) -- (0,8) (0,15) -- (0,24) -- (36,24) -- (36,15) (36,3) -- (36,-10) -- (0,-10) -- (0,-6);
  \draw[bkInk, line width=1.2pt] (-3.5,1) -- (3.5,1);
  \draw[bkInk, line width=1.2pt] (-1.8,4) -- (1.8,4);
  \fill[white, draw=bkInk, line width=0.8pt] (-2,8) rectangle (2,15);
  \fill[bkOrange!35, draw=bkInk, line width=0.8pt] (33.5,3) rectangle (38.5,15);
  \node[font=\sffamily\scriptsize, anchor=west] at (7,11.5) {\qty{2}{\ohm}};
  \node[font=\sffamily\scriptsize, anchor=west] at (7,2.5) {\qty{12}{V}};
  \node[font=\sffamily\footnotesize, anchor=west] at (39.5,9) {$R$};
  \draw[bk arrow, bkRed, line width=0.9pt] (14,26.5) -- (24,26.5);
  \node[font=\sffamily\scriptsize, text=bkRed, anchor=south] at (19,27.5) {$I$};
\end{tikzpicture}
\end{minipage}\par
\hspace*{2.6\zw}$P=RI\times I=R\times\dfrac{144}{(R+2)^2}=144R(R+2)^{-2}$\par
\shiki 積の微分と合成関数の微分（5-6）で\par
\hspace*{2.6\zw}$P'=144\{1\times(R+2)^{-2}+R\times(-2)(R+2)^{-3}\}$\par
\hspace*{2.6\zw}$(R+2)^{-2}=(R+2)\times(R+2)^{-3}$ なので、$(R+2)^{-3}$ でくくると\par
\hspace*{2.6\zw}$P'=144(R+2)^{-3}\{(R+2)-2R\}=\dfrac{144(2-R)}{(R+2)^3}$\par
\hspace*{2.6\zw}$R<2$ で $P'>0$、$R>2$ で $P'<0$ なので、$R=2$ で最大\par
\kotae $R=\qty{2}{\ohm}$（内部抵抗と同じ）、$P=\dfrac{144\times2}{4^2}=\dfrac{288}{16}=\qty{18}{W}$
\end{reidai}

\noindent\textbf{表で確かめる}\quad $R$ を変えて $P=\dfrac{144R}{(R+2)^2}$ を計算すると、下の表のように、$R=2$ のときがいちばん大きくなっています。抵抗が小さすぎると、電流は大きくても器具の電圧が小さく、抵抗が大きすぎると、電流が小さくなってしまうのです。
\begin{bkcenter}
\begin{bktable}{colspec={Q[c,m,wd=16mm] *{5}{Q[c,m,wd=12mm]}}, rowsep=1.5pt}
$R$（\unit{\ohm}） & 1 & 2 & 3 & 4 & 8 \\
$P$（W） & 16 & 18 & 17.28 & 16 & 11.52 \\
\end{bktable}
\end{bkcenter}

\begin{jibun}
電圧\qty{6}{V}、内部抵抗\qty{3}{\ohm}の電池に抵抗 $R$ をつなぎます。$R$ で使われる電力 $P$ を $R$ の式で表し、$P$ が最大になる $R$ と電力を求めましょう。
\bkjans{$I=\dfrac{6}{R+3}$、$P=RI^2=\dfrac{36R}{(R+3)^2}$。$P'=\dfrac{36(3-R)}{(R+3)^3}$ より $R=\qty{3}{\ohm}$、$P=\dfrac{108}{36}=\qty{3}{W}$}
\end{jibun}

%% ---------------- まちがえやすいところ・物理工学 ----------------
\bksection{tsumazuki}{まちがえやすいところをチェック}

\begin{tsumazuki}{$f'(x)=0$ なら、いつも極大か極小？}
\bkng{$f(x)=x^3$ で $f'(0)=0$ だから、$x=0$ は極大か極小}
\bkok{$x=0$ の前後で $f'>0$ のまま。極大でも極小でもない}
\vspace{2pt}\noindent\begin{minipage}[c]{0.76\linewidth}\raggedright
\strut$f(x)=x^3$ では $f'(x)=3x^2$ なので、$f'(0)=0$ です。でも、$x=0$ の前でも後でも $f'>0$ で、ずっと上がり続けています（とちゅうで一瞬だけ平らになる）。$f'(x)=0$ でも、前後で符号が変わらなければ、極大でも極小でもありません。増減表をかいて、符号が変わるかを確かめましょう。
\end{minipage}\hfill
\begin{minipage}[c]{0.21\linewidth}\centering
\begin{tikzpicture}
\begin{axis}[bkplot, width=28mm, height=30mm, xmin=-1.5, xmax=1.6, ymin=-2, ymax=2.2, xtick=\empty, ytick=\empty, xlabel={$x$}, ylabel={$y$}]
  \addplot[bkcurve, domain=-1.25:1.25] {x^3};
  \addplot[bkcurve2, line width=1pt, domain=-0.6:0.6] {0};
\end{axis}
\end{tikzpicture}
\end{minipage}
\end{tsumazuki}

\begin{tsumazuki}{範囲のはしを調べ忘れる}
\bkng{$f(x)=x^3-3x$（$-2\leqq x\leqq3$）の最大値は、極大値の2}
\bkok{はしの $f(3)=27-9=18$ もくらべて、最大値は18}
\noindent\begin{minipage}[c]{0.62\linewidth}\raggedright
範囲が決まっているとき、極大値がいちばん大きいとはかぎりません（右の図。\bkreidairef{3}でも、はしの $x=5$ が最大）。極大・極小の値と、はしの値を全部くらべます。
\end{minipage}\hfill
\begin{minipage}[c]{0.34\linewidth}\centering
\begin{tikzpicture}
\begin{axis}[bkplot, width=40mm, height=36mm, xmin=-2.4, xmax=3.6, ymin=-4, ymax=21, xtick={-1,3}, ytick={2,18}, xlabel={$x$}, ylabel={$y$}]
  \addplot[bkcurve, line width=1.3pt, domain=-2:3] {x^3-3*x};
  \addplot[only marks, mark=*, mark size=1.6pt, bkRed] coordinates {(3,18)};
  \addplot[only marks, mark=*, mark size=1.4pt, bkInk] coordinates {(-1,2)};
\end{axis}
\end{tikzpicture}
\end{minipage}
\end{tsumazuki}

\begin{tsumazuki}{範囲の外の解も答えにしてしまう}
\bkng{はてな？の箱で、$V'=0$ だから $x=2$ と $x=6$ が答え}
\bkok{$x=6$ は $0<x<6$ に入らない（底の1辺が0）。答えは $x=2$ だけ}
$f'(x)=0$ の解が出たら、問題の範囲に入っているかを、いつも確かめます。とくに長さ・時間・体積は、0やマイナスになると意味がなくなります。
\end{tsumazuki}

\begin{butsuri}
\begin{itemize}[leftmargin=1.2\zw, itemsep=2pt, topsep=0pt, label={\color{bkPurple}●}]
\item \textbf{缶の形}\\
同じ量が入る円柱の缶のうち、材料（表面積）がいちばん少ないのは、高さが直径と等しい形です（\bkreidairef{7}）。ほんとうの缶は、ふたの厚みや、持ちやすさ・積みやすさも考えて形を決めています。
\item \textbf{いちばん多く電気をとり出す}\\
器具の抵抗を電池の内部抵抗と同じにすると、器具で使える電力がいちばん大きくなります（\bkreidairef{8}）。電気工学で「整合」（マッチング）とよばれる考え方で、アンテナやスピーカーをつなぐときに使われます。
\item \textbf{ボールをいちばん遠くへ}\\
空気の抵抗がなければ、地面から45°の向きに投げたとき、いちばん遠くまで届きます。三角関数の微分（5-15）を学ぶと、この章の方法で確かめられます。
\end{itemize}
\end{butsuri}

%% ---------------- 章末問題 ----------------
\bkendproblems
\bkproblemgroup{基本}{3}

\bkproblem $f(x)=x^3-3x^2$ の増減表をかき、極大値と極小値を求めましょう。
\bkpans{$f'=3x^2-6x=3x(x-2)$。$x=0$ で極大値0、$x=2$ で極小値 $-4$}

\bkproblem 高さ $x=40t-5t^2$（m）のボールの、最高点の時刻と高さを求めましょう。
\bkpans{$\dot{x}=40-10t=0$ より4秒後、$160-80=\qty{80}{m}$}

\bkproblem $f(x)=-x^2+6x$ の、$0\leqq x\leqq4$ での最大値と最小値を求めましょう。
\bkpans{$f'=-2x+6$ より $x=3$。$f(0)=0$、$f(3)=9$、$f(4)=8$。最大値9（$x=3$）、最小値0（$x=0$）}

\bkproblem 長さ\qty{60}{cm}の針金で長方形をつくります。面積が最大になるたて・横と、その面積を求めましょう。
\bkpans{$S=x(30-x)$、$x=15$。たて・横とも\qty{15}{cm}、\qty{225}{cm^2}}

\bkproblem $f(x)=2x^3-3x^2-12x$ の極大値と極小値を求めましょう。
\bkpans{$f'=6x^2-6x-12=6(x-2)(x+1)$。$x=-1$ で極大値7、$x=2$ で極小値 $-20$}

\bkproblem $x>0$ で、$f(x)=x+\dfrac{9}{x}$ の最小値を求めましょう。
\bkpans{$f'=1-\dfrac{9}{x^2}=0$ より $x^2=9$、$x>0$ なので $x=3$。最小値 $3+3=6$}

\bkproblemgroup{チャレンジ}{4}

\bkproblem 1辺\qty{30}{cm}の正方形の紙の4すみから1辺 $x$（cm）の正方形を切りとって、ふたのない箱をつくります。体積 $V=x(30-2x)^2$ が最大になる $x$ と、そのときの体積を求めましょう（積の微分と合成関数の微分を使ってもよい）。
\bkpans{$V'=(30-2x)^2-4x(30-2x)=(30-2x)(30-6x)$。$x=\qty{5}{cm}$、$5\times20^2=\qty{2000}{cm^3}$}

\bkproblem 体積が $16\pi$ \unit{cm^3} の円柱の缶をつくります。半径を $r$ とすると、表面積は $S=2\pi r^2+\dfrac{32\pi}{r}$ です。$S$ が最小になる $r$ と、そのときの高さを求めましょう。
\bkpans{$S'=4\pi r-\dfrac{32\pi}{r^2}=0$ より $r^3=8$、$r=\qty{2}{cm}$。高さ $\dfrac{16}{2^2}=\qty{4}{cm}$（直径と等しい）}

\bkproblem 電圧\qty{9}{V}、内部抵抗\qty{1.5}{\ohm}の電池に抵抗 $R$ をつなぐと、$R$ の電力は $P=\dfrac{81R}{(R+1.5)^2}$（W）です。$P$ が最大になる $R$ と電力を求めましょう。
\bkpans{$P'=\dfrac{81(1.5-R)}{(R+1.5)^3}$ より $R=\qty{1.5}{\ohm}$、$P=\dfrac{121.5}{9}=\qty{13.5}{W}$}

\begin{furikaeri}
\begin{checklist}
\item 導関数の符号から、上がる・下がるを判断できる。\hfill{\footnotesize\sffamily\color{bkTheme}→\pageref{bk:5-7:fugou}ページ}
\item 増減表をかいて、極大・極小を求められる。\hfill{\footnotesize\sffamily\color{bkTheme}→\pageref{bk:5-7:zougen}ページ}
\item 範囲のはしもくらべて、最大・最小を求められる。\hfill{\footnotesize\sffamily\color{bkTheme}→\pageref{bk:5-7:saidai}ページ}
\end{checklist}
\end{furikaeri}

\bkfukidashi{mina}{接線は、グラフの「近くのようす」をよく表しているんだね。次の章では、接線を使って、少し先の値を予想するよ。}

%% 5-8 短い時間なら、かけ算で近似できる（第11回：書き足し。ページ数は自由・偶数）
\makeatletter
% 立方体の図の直方体（前・右・上の3面）  \bkvhbox{色}{x0}{y0}{z0}{x1}{y1}{z1}（x・y・z の向きは scope で決める）
\newcommand\bkvhbox[7]{%
  \fill[#1!45, draw=bkInk, line width=0.45pt] (#2,#3,#4) -- (#5,#3,#4) -- (#5,#3,#7) -- (#2,#3,#7) -- cycle;
  \fill[#1!25, draw=bkInk, line width=0.45pt] (#5,#3,#4) -- (#5,#6,#4) -- (#5,#6,#7) -- (#5,#3,#7) -- cycle;
  \fill[#1!65, draw=bkInk, line width=0.45pt] (#2,#3,#7) -- (#5,#3,#7) -- (#5,#6,#7) -- (#2,#6,#7) -- cycle;}
\makeatother
\bkchapter{5}{8}{短い時間なら、\\かけ算で近似できる}{\bklevel{2}}{短い時間なら、かけ算で近似できる}
\bkmichishirube{2-6,5-3,5-5}{5-9,5-17,6-2}

\begin{goal}
\begin{itemize}[label={\color{bkTheme}●}, leftmargin=1.6\zw, itemsep=0pt, topsep=1pt]
\item \bklevel{2}\enspace 「次の位置$\approx$今の位置＋速度×$\Delta t$」で、少し先を予想できる。
\item \bklevel{2}\enspace 接線を使って、$f(a+h)\approx f(a)+f'(a)h$ と近似できる。
\item \bklevel{3}\enspace はかった値のずれが、計算にひびく大きさを見積もれる。
\end{itemize}
\end{goal}

\begin{hatena}
\noindent\begin{minipage}[c]{0.55\linewidth}\raggedright
車のカーナビは、人工衛星の電波（GPS）で、いまの位置を知っています。でも、長いトンネルの中では電波がとどきません。

それでも、カーナビの地図の上の車は、トンネルの中をちゃんと進んでいきます。電波がないのに、どうやって位置を決めているのでしょう。
\end{minipage}\hfill
\begin{minipage}[c]{0.41\linewidth}
\begin{tikzpicture}[x=1mm, y=1mm]
  \fill[bkGray!35] (-2,0) rectangle (50,16);
  \fill[bkGray!60] (-2,16) arc (180:0:26 and 10) -- (50,16) -- cycle;
  \fill[bkInk!80] (8,2) rectangle (42,14);
  \draw[white, line width=0.8pt, dash pattern=on 3pt off 2pt] (8,8) -- (42,8);
  \fill[bkRoute2] (20,5) rectangle ++(9,5);
  \fill[bkInk] (22,4.6) circle (1) (27,4.6) circle (1);
  \draw[bk arrow, bkOrange, line width=1pt] (31,7.5) -- (39,7.5);
  \node[font=\sffamily\scriptsize] at (24,-3) {トンネルの中};
\end{tikzpicture}
\end{minipage}
\end{hatena}

カーナビは、トンネルに入る直前の位置と、タイヤの回転からわかる速さを使って、「少し時間がたったら、速さ×時間だけ進んでいるはず」と計算しています。曲がった向きは、回転をはかるセンサー（ジャイロセンサー）でわかります。この章では、「短い時間なら、速さ×時間で進む距離がほぼわかる」という考え方を、微分を使って説明します。

\clearpage
%% ---------------- やってみよう・思い出そう ----------------
\begin{yattemiyou}[速度から、少しずつ位置を予想しよう]
落ちるボールの速度は $v=10t$（m/s）です（5-2）。位置0から、0.5秒ごとに「次の位置＝今の位置＋今の速度×0.5」と計算します。表計算では、D3に \texttt{=D2+C2*0.5} と入れ、下へコピーします（2-18）。
\begin{bkcenter}
\bksheet{4}{22}{%
1 & $t$（s） & $v=10t$ & 予想した位置 & ほんとう $5t^2$ \\
2 & 0 & 0 & 0 & 0 \\
3 & 0.5 & 5 & 0 & 1.25 \\
4 & 1.0 & 10 & 2.5 & 5 \\
5 & 1.5 & 15 & 7.5 & 11.25 \\
6 & 2.0 & 20 & 15 & 20 \\}
\end{bkcenter}
\begin{enumerate}[label={\bkcircled[bkGreen]{\arabic*}}, leftmargin=2.4\zw, itemsep=1pt, topsep=3pt]
\item 0.5秒ごとの表の数を、電卓で1行ずつ確かめよう。
\item 0.1秒ごとにして、2秒の位置を予想しよう。
\item 0.01秒ごとにすると、どうなるかな？
\end{enumerate}
\end{yattemiyou}

\noindent\textbf{表の読み方}\quad 予想した位置の列は、1つ上の行の「位置＋速度×0.5」です。
\begin{itemize}[leftmargin=1.6\zw, itemsep=1pt, topsep=2pt, label={\color{bkTheme}●}]
\item 0.5秒：$0+0\times0.5=0$。0秒の速度が0なので、予想では1歩も進みません。
\item 1.0秒：$0+5\times0.5=2.5$。0.5秒の速度5を使います。
\item 1.5秒：$2.5+10\times0.5=7.5$。\quad 2.0秒：$7.5+15\times0.5=15$。
\end{itemize}
0.5秒ごとでは、2秒の予想は\qty{15}{m}で、ほんとうの\qty{20}{m}より小さくなりました。ボールは、0.5秒のあいだにもどんどん速くなっているのに、予想では区間の\textbf{はじめの速度}のまま進むとしたからです。

0.1秒ごとにすると\qty{19.0}{m}、0.01秒ごとにすると\qty{19.9}{m}になり、ほんとうの値に近づきます。区間を短くするほど、そのあいだに速度が変わる量も小さくなるので、予想は正確になります。

\begin{omoidasou}
\bkrecall{5-3}{$f'(a)=\displaystyle\lim_{h\to0}\frac{f(a+h)-f(a)}{h}$。接線の式は $y=f(a)+f'(a)(x-a)$。接点の近くでは、接線はグラフにぴったりよりそう。}
\bkrecall{5-5}{速度は位置を時間で微分したもの。$v=\dfrac{\dd x}{\dd t}$}
\bkrecall{2-14}{等式の両辺に同じ数をかけても、等式はなりたつ。}
\bkrecall{2-6}{はかった値にはずれ（誤差）がある。$\qty{5.0}{cm}$ と書けば、およそ $\pm\qty{0.05}{cm}$ のずれをふくむ。}
\bkrecall{2-18}{表計算では、上の行の結果を使って下の行を計算し、式をコピーしてくり返せる。}
\end{omoidasou}

%% ---------------- 数学の道具 ----------------
\bksection{doguu}{数学の道具}

\bksubsection{1}{次の位置$\approx$今の位置＋速度×$\Delta t$}

\label{bk:5-8:tsugi}速度は $v=\dfrac{\dd x}{\dd t}$ でした（5-5）。これは、「時間 $\Delta t$ のあいだに位置が $\Delta x$ 変わったとき、$\dfrac{\Delta x}{\Delta t}$ は、$\Delta t$ を0に近づけると $v$ に近づく」という意味でした（5-3）。だから、$\Delta t$ がとても小さいときは、
\[ \frac{\Delta x}{\Delta t}\approx v \]
です。「$\approx$」は「\textbf{ほぼ等しい}」という記号で、「ニアリーイコール」と読みます。両辺に $\Delta t$ をかけると（2-14）、
\[ \Delta x\approx v\,\Delta t \]
となります。短い時間に進む距離（$\Delta x$）は、速度×時間でほぼ求められる、ということです。\bkidx{きんじ}{近似（$\approx$）}

\bkkigocards{\bkkigocard{$\approx$}{ニアリーイコール}{ほぼ等しい（近似）}{nearly equal}{5-8}}

\noindent\textbf{例}\quad 位置\qty{100}{m}の所を\qty{15}{m/s}で走る車は、0.2秒のあいだに $15\times0.2=\qty{3}{m}$ ほど進みます。0.2秒後の位置は、$100+3=\qty{103}{m}$ ほどです。

\begin{bkfigure}
\begin{tikzpicture}[x=1mm, y=1mm]
  \draw[bk axis] (-4,0) -- (100,0);
  \node[font=\sffamily\small, anchor=west] at (101,0) {m};
  \foreach \p in {98,...,104} \draw[bkInk, line width=0.5pt] ({14*(\p-98)+5},-1) -- ({14*(\p-98)+5},1);
  \foreach \p in {98,100,102,104} \node[font=\sffamily\footnotesize, anchor=north] at ({14*(\p-98)+5},-2) {\p};
  \bkcar{(28.4,2.5)}{0.33}{bkRoute2}
  \fill[bkBlue] (33,0) circle (1.2);
  \fill[bkOrange] (75,0) circle (1.2);
  \draw[bk arrow, bkOrange] (33,-9) -- (75,-9);
  \node[font=\sffamily\footnotesize, anchor=north] at (54,-10.5) {速度×$\Delta t=15\times0.2=3$};
  \node[font=\sffamily\footnotesize, anchor=south, text=bkBlue] at (33,9.5) {今の位置 100};
  \node[font=\sffamily\footnotesize, anchor=south, text=bkOrange!80!black] at (75,3) {次の位置 103};
\end{tikzpicture}
\bkfigcaption{今の位置に「速度×$\Delta t$」をたすと、次の位置が予想できる。}
\end{bkfigure}

だから、少し先の位置は
\[ \text{次の位置}\approx\text{今の位置}+\text{今の速度}\times\Delta t \]
と予想できます。$\Delta t$ のあいだにも速度は変わるので、ぴったりではありません。

\noindent\textbf{$\Delta t$ が小さいほどよいわけ}\quad 落ちるボール（$v=10t$）は、1秒ごとに速度が\qty{10}{m/s}ふえます。$\Delta t=0.5$ 秒なら、そのあいだに速度は\qty{5}{m/s}も変わってしまいます。$\Delta t=0.1$ 秒なら\qty{1}{m/s}、$\Delta t=0.01$ 秒なら\qty{0.1}{m/s}しか変わりません。$\Delta t$ が小さいほど、「速度がずっと同じ」と考えても、ほんとうとのちがいが小さくなるのです。

この計算をくり返して少しずつ先を予想するのが、やってみようの計算です。この方法を\textbf{オイラー法}といいます。第6巻では、この方法でいろいろな運動を計算します（6-2）。\bkidx{おいらーほう}{オイラー法}

\noindent\textbf{$\Delta t$ と予想のずれ}\quad やってみようの結果をまとめると、次のようになります。

\noindent\begin{minipage}[c]{0.47\linewidth}\centering
\begin{bktable}{colspec={Q[c,m,wd=12mm] Q[c,m,wd=19mm] Q[c,m,wd=12mm]}, rowsep=1.5pt}
$\Delta t$（s） & {2秒の予想\\（m）} & {ずれ\\（m）} \\
0.5 & 15 & 5 \\
0.1 & 19.0 & 1.0 \\
0.01 & 19.9 & 0.1 \\
\end{bktable}
\bktabcaption{$\Delta t$ を $\frac{1}{10}$ にすると、\\ずれもほぼ $\frac{1}{10}$ になる。}
\end{minipage}\hfill
\begin{minipage}[c]{0.51\linewidth}\centering
\begin{tikzpicture}
\begin{axis}[bkplot, width=52mm, height=42mm, xmin=0, xmax=2.3, ymin=0, ymax=22, xtick={0.5,1,1.5,2}, ytick={5,10,15,20}, xlabel={$t$（s）}, ylabel={$x$（m）}]
  \addplot[bkcurve, domain=0:2] {5*x^2};
  \addplot[bkcurve2, line width=1.2pt, mark=*, mark size=1.4pt] coordinates {(0,0) (0.5,0) (1,2.5) (1.5,7.5) (2,15)};
  \addplot[bkcurve3, line width=1pt] coordinates {(0,0) (0.1,0) (0.2,0.1) (0.3,0.3) (0.4,0.6) (0.5,1) (0.6,1.5) (0.7,2.1) (0.8,2.8) (0.9,3.6) (1,4.5) (1.1,5.5) (1.2,6.6) (1.3,7.8) (1.4,9.1) (1.5,10.5) (1.6,12) (1.7,13.6) (1.8,15.3) (1.9,17.1) (2,19)};
  \node[font=\sffamily\scriptsize, text=bkOrange!80!black, anchor=west] at (axis cs:2.1,13.5) {$0.5$};
  \node[font=\sffamily\scriptsize, text=bkGreen!70!black, anchor=west] at (axis cs:2.1,18.2) {$0.1$};
  \node[font=\sffamily\scriptsize, text=bkBlue, anchor=east] at (axis cs:1.7,18) {ほんとう};
\end{axis}
\end{tikzpicture}
\bkfigcaption{予想（オレンジ・緑、数は $\Delta t$）と\\ほんとうの位置（青）。}
\end{minipage}

\medskip
ずれは、どれも $\Delta t$ のちょうど10倍（$0.5\times10=5$、$0.1\times10=1$、$0.01\times10=0.1$）で、$\Delta t$ に比例して小さくなります。10倍正確にしたければ、10倍細かく計算すればよいのです。計算の回数は10倍になりますが、コンピュータなら一瞬です。

\begin{motto}[3]
\textbf{ずれが $\Delta t$ の10倍になるわけ}

\noindent\begin{minipage}[c]{0.4\linewidth}\centering
\begin{tikzpicture}
\begin{axis}[bkplot, width=50mm, height=44mm, xmin=0, xmax=2.4, ymin=0, ymax=23, xtick={0.5,1,1.5,2}, ytick={5,10,15,20}, xlabel={$t$}, ylabel={$v$}, x tick label style={font=\sffamily\scriptsize}]
  \foreach \k in {1,2,3} {\edef\tmp{\noexpand\fill[bkBlue!25, draw=bkBlue!70, line width=0.4pt] (axis cs:{0.5*\k},0) rectangle (axis cs:{0.5*\k+0.5},{5*\k});}\tmp}
  \foreach \k in {0,1,2,3} {\edef\tmp{\noexpand\fill[bkOrange!55] (axis cs:{0.5*\k},{5*\k}) -- (axis cs:{0.5*\k+0.5},{5*\k}) -- (axis cs:{0.5*\k+0.5},{5*\k+5}) -- cycle;}\tmp}
  \addplot[bkcurve, domain=0:2.1] {10*x};
\end{axis}
\end{tikzpicture}
\end{minipage}\hfill
\begin{minipage}[c]{0.56\linewidth}\raggedright
$\Delta t=0.5$ のときの $v$-$t$ グラフ（$v=10t$）。青い長方形は「はじめの速度×0.5」で、面積をたすと予想の15になります。オレンジのすき間が、予想で足りなかった分です。
\end{minipage}

\medskip
速度のグラフ（$v$-$t$ グラフ）で考えます。「速度×$\Delta t$」は、たて（速度）×横（$\Delta t$）の細い長方形の面積です（2-3）。予想は、区間のはじめの速度を使うので、長方形の階段がグラフの直線より少し下になり、すき間（小さな三角形）の分だけ少なくなります。1つのすき間は、底辺 $\Delta t$、高さ $10\times\Delta t$（$\Delta t$ のあいだにふえた速度）の三角形で、面積は $\Delta t\times10\Delta t\div2=5\times\Delta t\times\Delta t$。2秒までに、すき間は $2\div\Delta t$ こあるので、合計は $(2\div\Delta t)\times5\times\Delta t\times\Delta t=10\times\Delta t$ です。グラフの面積と距離の関係は、次の5-9でくわしく学びます。
\end{motto}

\bksubsection{2}{接線で近似する}

\label{bk:5-8:sessen}同じ考え方を、ふつうの関数 $y=f(x)$ にも使ってみましょう。①では、「位置の変わる速さ（速度）×時間の変化」で、位置の変化を求めました。関数でも、「$y$ の変わる速さ $f'(a)$ × $x$ の変化 $h$」で、$y$ の変化がほぼわかるはずです。
\[ f(a+h)\approx f(a)+f'(a)\,h \]
これをグラフで確かめます。$y=f(x)$ のグラフの点 $(a,\ f(a))$ での接線は、$y=f(a)+f'(a)(x-a)$ でした（5-3）。$x=a+h$ とすると $x-a=h$ なので、接線の上では $y=f(a)+f'(a)h$ です。$h$ が小さければ、グラフと接線はほとんど重なっているので、グラフの高さ $f(a+h)$ は、接線の高さ $f(a)+f'(a)h$ にほぼ等しくなります。

\noindent\begin{minipage}[c]{0.5\linewidth}\centering
\begin{tikzpicture}
\begin{axis}[bkplot, width=60mm, height=46mm, xmin=0, xmax=4.1, ymin=0, ymax=2.3, xtick={1,3}, xticklabels={$a$,$a+h$}, ytick=\empty, xlabel={$x$}, ylabel={$y$}]
  \addplot[bkcurve, domain=0.02:3.05] {sqrt(x)};
  \addplot[bkcurve2, line width=1pt, domain=0.1:3.05] {1+0.5*(x-1)};
  \draw[bkGray, dashed, line width=0.5pt] (axis cs:1,0) -- (axis cs:1,1) -- (axis cs:3,1) (axis cs:3,0) -- (axis cs:3,1);
  \draw[bkInk, line width=0.6pt, <->] (axis cs:1,0.25) -- (axis cs:3,0.25);
  \node[font=\sffamily\scriptsize, anchor=south] at (axis cs:2,0.27) {$h$};
  \draw[bkInk, line width=0.6pt, <->] (axis cs:3,1) -- (axis cs:3,2);
  \node[font=\sffamily\scriptsize, anchor=west] at (axis cs:3.12,1.35) {$f'(a)h$};
  \draw[bkRed, line width=0.7pt, decorate, decoration={brace, amplitude=2.5pt}] (axis cs:3.16,2) -- (axis cs:3.16,1.732);
  \node[font=\sffamily\scriptsize, anchor=west, text=bkRed] at (axis cs:3.28,1.866) {ずれ};
  \addplot[only marks, mark=*, mark size=1.6pt, bkInk] coordinates {(1,1)};
\end{axis}
\end{tikzpicture}
\end{minipage}\hfill
\begin{minipage}[c]{0.46\linewidth}
図で、青がグラフ、オレンジが接線です。$a$ から $h$ 進むと、接線は $f'(a)\times h$ だけ上がります（傾き×横の長さ）。グラフは、それより少しだけ低い所にあります。このちがいが近似の\textbf{ずれ}です。$h$ が小さくなるほど、ずれは小さくなります。
\end{minipage}

\medskip
このように、接線を使って関数の値を求めることを\textbf{1次近似}（いちじきんじ）といいます。「1次」は、$h$ の1乗までしか使わないという意味です。\bkidx{いちじきんじ}{1次近似}

\noindent\textbf{例}\quad $\sqrt{4.1}$ を近似します。4.1は、答えがすぐわかる4（$\sqrt{4}=2$）に近いので、$a=4$、$h=0.1$ とします。
\begin{enumerate}[leftmargin=2.6\zw, itemsep=1pt, topsep=2pt]
\item $f(x)=\sqrt{x}$ なので、$f(4)=\sqrt{4}=2$。
\item $f'(x)=\dfrac{1}{2\sqrt{x}}$（5-6）なので、$f'(4)=\dfrac{1}{2\times2}=\dfrac{1}{4}$。
\item $\sqrt{4.1}\approx2+\dfrac{1}{4}\times0.1=2+0.025=2.025$
\end{enumerate}
電卓では $\sqrt{4.1}=2.02485\cdots$。小数第3位まで合っています。

\begin{bkcenter}
\begin{bktable}{colspec={Q[c,m,wd=14mm] Q[c,m,wd=34mm] Q[c,m,wd=30mm] Q[c,m,wd=22mm]}, rowsep=1.5pt}
$h$ & 1次近似 $2+\frac{1}{4}h$ & ほんとう $\sqrt{4+h}$ & ずれ \\
0.1 & 2.025 & 2.02485… & 0.00015 \\
1 & 2.25 & 2.23607… & 0.014 \\
5 & 3.25 & 3 & 0.25 \\
\end{bktable}
\bktabcaption{$h$ が小さいとよく合い、$h$ が大きいとずれる。}
\end{bkcenter}

\bksubsection{3}{近似のずれは、どれくらい？}

\label{bk:5-8:zure}$f(x)=x^2$ で、近似のずれをくわしく調べましょう。$f'(x)=2x$ なので、ほんとうの値と1次近似は、次のようになります。
\begin{align*}
f(a+h)&=(a+h)^2=a^2+2ah+h^2 && \text{（ほんとう。3-8の展開）}\\
f(a)+f'(a)h&=a^2+2a\times h=a^2+2ah && \text{（1次近似）}
\end{align*}
ひき算すると、ずれはちょうど $h^2$ です。

\noindent\begin{minipage}[c]{0.4\linewidth}\centering
\begin{tikzpicture}[x=1mm, y=1mm]
  \fill[bkBlue!22] (0,0) rectangle (30,30);
  \fill[bkOrange!40] (30,0) rectangle (38,30);
  \fill[bkOrange!40] (0,30) rectangle (30,38);
  \fill[bkRed!55] (30,30) rectangle (38,38);
  \draw[bkInk, line width=0.6pt] (0,0) rectangle (38,38) (30,0) -- (30,38) (0,30) -- (38,30);
  \node[font=\sffamily\small] at (15,15) {$a^2$};
  \node[font=\sffamily\footnotesize] at (34,15) {$ah$};
  \node[font=\sffamily\footnotesize] at (15,34) {$ah$};
  \node[font=\sffamily\footnotesize, anchor=north] at (15,-1.5) {$a$};
  \node[font=\sffamily\footnotesize, anchor=north] at (34,-1.5) {$h$};
  \node[font=\sffamily\footnotesize, anchor=west] at (39.5,34) {$h^2$};
\end{tikzpicture}
\end{minipage}\hfill
\begin{minipage}[c]{0.56\linewidth}
図で見ると、1辺 $a+h$ の正方形の面積は、もとの正方形 $a^2$（青）、細い長方形 $ah$ が2つ（オレンジ）、すみの小さな正方形 $h^2$（赤）に分かれます。1次近似は、青とオレンジだけを数えて、赤のすみ $h^2$ を数えていません。$h$ が小さいと、すみの正方形は細い長方形よりずっと小さくなります。
\end{minipage}

\begin{bkcenter}
\begin{bktable}{colspec={Q[c,m,wd=16mm] Q[c,m,wd=30mm] Q[c,m,wd=30mm] Q[c,m,wd=26mm]}, rowsep=1.5pt}
$h$ & {$3^2+2\times3h$\\（近似）} & {$(3+h)^2$\\（ほんとう）} & ずれ $h^2$ \\
1 & 15 & 16 & 1 \\
0.1 & 9.6 & 9.61 & 0.01 \\
0.01 & 9.06 & 9.0601 & 0.0001 \\
\end{bktable}
\bktabcaption{$a=3$ のとき。$h$ を $\frac{1}{10}$ にすると、ずれは $\frac{1}{100}$ になる。}
\end{bkcenter}

$h$ を10分の1にすると、ずれは $\dfrac{1}{10}\times\dfrac{1}{10}=\dfrac{1}{100}$ になります。ずれは $h$ よりずっと速く小さくなるので、$h$ が小さいときの1次近似はとても正確です。

\bksubsection{4}{よく使う近似}

\label{bk:5-8:yoku}$h$ が1よりずっと小さいとき、次の近似がよく使われます。
\[ (1+h)^n\approx1+nh,\qquad \sqrt{1+h}\approx1+\frac{h}{2},\qquad \frac{1}{1+h}\approx1-h \]
どれも、②の1次近似から作れます。

\noindent\textbf{1つめ}\quad $f(x)=x^n$、$a=1$ とします。$f(1)=1^n=1$。$f'(x)=nx^{n-1}$（5-4）なので、$f'(1)=n\times1^{n-1}=n$。だから
\[ (1+h)^n=f(1+h)\approx f(1)+f'(1)h=1+nh \]
$n=2$ で確かめると、ほんとうは $(1+h)^2=1+2h+h^2$ で、近似は $1+2h$。すみの $h^2$（③）だけがちがいます。$h=0.01$ なら、$1.01^2=1.0201$ と $1.02$ で、ほぼ同じです。

\noindent\textbf{2つめ・3つめ}\quad $\sqrt{1+h}=(1+h)^{\frac{1}{2}}$（3-11）、$\dfrac{1}{1+h}=(1+h)^{-1}$（2-7・2-12）なので、1つめの $n$ に $\dfrac{1}{2}$ と $-1$ を入れたものです。$(x^n)'=nx^{n-1}$ は、$n$ が分数やマイナスでも成り立つので（5-4・5-6）、そのまま使えます。
\[ \sqrt{1+h}\approx1+\frac{1}{2}h,\qquad \frac{1}{1+h}\approx1+(-1)\times h=1-h \]

\noindent\textbf{例}\quad 利率1\,\%で10年預けたお金は、何倍になるでしょう。1\,\%の利息がつくと、お金は1年で $1+0.01=1.01$ 倍になります（2-5）。10年なら、1.01を10回かけて $1.01^{10}$ 倍です。$h=0.01$、$n=10$ として近似すると
\[ 1.01^{10}=(1+0.01)^{10}\approx1+10\times0.01=1.1 \]
およそ1.1倍です（ほんとうは $1.1046\cdots$ 倍）。1.01を10回かけるより、ずっと楽に見当がつきます。

\bksubsection{5}{はかった値のずれが、計算にひびく大きさ}

\label{bk:5-8:gosa}はかった値には、ずれ（誤差）があります（2-6）。$x$ のずれが $\Delta x$ のとき、それを使って計算した $y=f(x)$ のずれ $\Delta y$ は、1次近似から見積もれます。$f(x+\Delta x)\approx f(x)+f'(x)\Delta x$ なので、$y$ の変わる量は\bkidx{ごさ}{誤差}
\[ \Delta y=f(x+\Delta x)-f(x)\approx f'(x)\,\Delta x \]

\noindent\textbf{例}\quad 立方体の1辺をはかったら \qty{10.0}{cm} で、$\pm\qty{0.1}{cm}$ のずれがありそうです。体積 $V=x^3$ は、どれくらいずれるでしょう。
\begin{enumerate}[leftmargin=2.6\zw, itemsep=1pt, topsep=2pt]
\item $V'=3x^2$。$x=10$ で $3\times10^2=300$。
\item $\Delta V\approx300\times0.1=\qty{30}{cm^3}$
\end{enumerate}

\noindent\begin{minipage}[c]{0.5\linewidth}\raggedright
図で見ると、1辺が\qty{0.1}{cm}のびたとき、ふえる部分は、ほぼ $10\times10\times0.1=10$（\unit{cm^3}）の板が3まいで、$10\times3=30$ です（5-4の立方体の図）。電卓で確かめると $10.1^3=1030.301$ で、ふえたのは $30.301$。ほとんど同じです。わずかなちがいは、板と板のあいだのすみの細い部分です。
\end{minipage}\hfill
\begin{minipage}[c]{0.46\linewidth}\centering
\begin{tikzpicture}[x=1mm, y=1mm]
  \begin{scope}[x={(1mm,0mm)}, y={(0.42mm,0.3mm)}, z={(0mm,1mm)}]
    \bkvhbox{bkGreen}{0}{30}{0}{22}{33}{22}
    \bkvhbox{bkBlue}{0}{0}{0}{22}{22}{22}
    \bkvhbox{bkOrange}{30}{0}{0}{33}{22}{22}
    \bkvhbox{bkRed}{0}{0}{29}{22}{22}{32}
  \end{scope}
  \node[font=\sffamily\footnotesize, anchor=north] at (11,-1.5) {10 cm};
  \node[font=\sffamily\footnotesize, anchor=north] at (31.5,-1.5) {0.1 cm};
\end{tikzpicture}
\bkfigcaption{もとの立方体（青）と、ふえる3まいの板。\\板のあつさは大きくかいている。}
\end{minipage}

\medskip
\noindent\textbf{ずれを割合で見る}\quad 体積は\qty{1000}{cm^3}なので、ずれの30は $30\div1000=0.03$、つまり3\,\%です。1辺のずれは $0.1\div10=0.01$ で1\,\%なのに、体積では3倍の3\,\%になりました。

これは、たまたまではありません。$V=x^3$ なら、$\Delta V\approx3x^2\Delta x$ なので、
\[ \frac{\Delta V}{V}\approx\frac{3x^2\,\Delta x}{x^3}=3\times\frac{\Delta x}{x} \]
（$x^2$ と $x^3$ を約分すると、分母に $x$ が1つ残る）。体積のずれの割合は、いつも1辺のずれの割合の3倍です。同じように、面積 $S=x^2$ なら2倍になります。\textbf{$n$ 乗の量は、ずれの割合が $n$ 倍にふくらむ}のです。

\begin{taisetsu}
$\Delta t$ が小さいとき\quad 次の位置 $\approx$ 今の位置＋速度×$\Delta t$\\
1次近似\quad $f(a+h)\approx f(a)+f'(a)\,h$\quad（接線で近似する）\\
ずれの見積もり\quad $\Delta y\approx f'(x)\,\Delta x$。$x^n$ の量では、ずれの割合が $n$ 倍。
\end{taisetsu}

\begin{kotoba}
\small
\bkword{$\approx$}{「ほぼ等しい」という記号。ニアリーイコール。}
\bkword{1次近似}{接線を使って、関数の値を近くで近似すること。}
\bkword{オイラー法}{「次の値＝今の値＋変わる速さ×$\Delta t$」をくり返して、少しずつ先を予想する方法。}
\bkword{誤差}{はかった値や近似した値と、ほんとうの値とのずれ。}
\end{kotoba}

\begin{column}{丸い地球も、近くで見れば平ら}
\noindent\begin{minipage}[c]{0.64\linewidth}\raggedright
地球は丸いのに、校庭は平らに見えます。地球の半径（約\qty{6400}{km}）にくらべて、校庭の\qty{100}{m}はとても短いからです。100 m先で、まっすぐな線（接線）から地面が下がるのは、わずか約\qty{0.8}{mm}です。曲がったものも、ごく近くだけを見ればまっすぐとみなせる――これが1次近似の考え方です。工学では、本当は曲がった関係をもつ部品（センサーやばね）も、使う範囲がせまいときは直線（比例）とみなして計算します。
\end{minipage}\hfill
\begin{minipage}[c]{0.32\linewidth}\centering
\begin{tikzpicture}[x=0.72mm, y=0.72mm]
  \draw[bkBlue, line width=1.4pt] (0,0) arc (90:55:40);
  \draw[bkBlue, line width=1.4pt] (0,0) arc (90:125:40);
  \draw[bkOrange, line width=0.9pt] (-24,0) -- (24,0);
  \fill[bkInk] (0,0) circle (0.9);
  \node[font=\sffamily\scriptsize, text=bkOrange!80!black, anchor=south] at (14,0.8) {接線};
  \node[font=\sffamily\scriptsize, text=bkBlue, anchor=north] at (0,-6) {地面（円）};
\end{tikzpicture}
\end{minipage}
\end{column}

%% ---------------- 例題 ----------------
\bksection{reidai}{例題でたしかめよう}

\begin{reidai}[次の位置を予想する]
（1）位置\qty{100}{m}の所を、速度\qty{15}{m/s}で走っている車の、0.2秒後の位置を予想しましょう。\par
（2）位置\qty{100}{m}の所を、速度\qty{-5}{m/s}で動いている物体の、0.4秒後の位置を予想しましょう。\par
\kangaekata 0.2秒や0.4秒は短いので、そのあいだの速度はほぼ一定と見ます。①の「次の位置$\approx$今の位置＋速度×$\Delta t$」に入れます。速度がマイナスのときは、位置が小さくなる向きに動くので、マイナスのままかけます。\par
\shiki（1）$100+15\times0.2=100+3=103$\par
\hspace*{2.6\zw}（2）$100+(-5)\times0.4=100+(-2)=98$\par
\begin{bkfigure}
\begin{tikzpicture}[x=1mm, y=1mm]
  \draw[bk axis] (-4,0) -- (104,0);
  \node[font=\sffamily\small, anchor=west] at (105,0) {m};
  \foreach \p in {96,...,104} \draw[bkInk, line width=0.5pt] ({12*(\p-96)},-1) -- ({12*(\p-96)},1);
  \foreach \p in {96,98,100,102,104} \node[font=\sffamily\footnotesize, anchor=north] at ({12*(\p-96)},-2) {\p};
  \fill[bkInk] (48,0) circle (1.2);
  \draw[bk arrow, bkBlue] (48,4) -- (84,4);
  \node[font=\sffamily\footnotesize, anchor=south, text=bkBlue] at (66,5) {（1）$+15\times0.2=+3$};
  \draw[bk arrow, bkRed] (48,-9) -- (24,-9);
  \node[font=\sffamily\footnotesize, anchor=north, text=bkRed!80!black] at (36,-10.5) {（2）$(-5)\times0.4=-2$};
\end{tikzpicture}
\end{bkfigure}
\kotae（1）およそ\qty{103}{m}\quad（2）およそ\qty{98}{m}
\end{reidai}

\begin{jibun}
（1）位置\qty{50}{m}の所を、速度\qty{8}{m/s}で走っている自転車の、0.5秒後の位置を予想しましょう。\par
（2）位置\qty{20}{m}の所を、速度\qty{-6}{m/s}で動いている物体の、0.5秒後の位置を予想しましょう。
\bkjans{（1）\nobreak $50+8\times0.5=\qty{54}{m}$\quad（2）\nobreak $20+(-6)\times0.5=\qty{17}{m}$}
\end{jibun}

\begin{reidai}[くり返して予想する]
速度が $v=2t$（m/s）の物体が、$t=0$ で位置0にいます。1秒ごとに「次の位置＝今の位置＋今の速度×1」と計算して、3秒の位置を予想しましょう。ほんとうの位置 $t^2$ とくらべましょう。\par
\kangaekata ①の式を、$\Delta t=1$ で3回くり返します。使う速度は、各1秒の\textbf{はじめ}の値です（0秒で0、1秒で2、2秒で4）。\par
\shiki 1秒の位置：$0+0\times1=0$\par
\hspace*{2.6\zw}2秒の位置：$0+2\times1=2$\par
\hspace*{2.6\zw}3秒の位置：$2+4\times1=6$\par
\begin{bkcenter}
\begin{bktable}{colspec={Q[c,m,wd=10mm] Q[c,m,wd=24mm] Q[c,m,wd=20mm] Q[c,m,wd=16mm] Q[c,m,wd=18mm]}, rowsep=1pt, cells={font=\small}, row{1}={font=\sffamily\bfseries\footnotesize}}
$t$（s） & はじめの速度 $2t$ & 速度×1（m） & 予想（m） & ほんとう $t^2$ \\
0 & 0 & 0 & 0 & 0 \\
1 & 2 & 2 & 0 & 1 \\
2 & 4 & 4 & 2 & 4 \\
3 & & & 6 & 9 \\
\end{bktable}
\end{bkcenter}
\kotae 予想\qty{6}{m}（ほんとうは $3^2=\qty{9}{m}$。$\Delta t$ を小さくすると近づく）
\end{reidai}

\begin{jibun}
\bkreidairef{2}で、0.5秒ごとに計算して、2秒の位置を予想しましょう（ほんとうは\qty{4}{m}）。
\bkjans{速度は0、1、2、3。0 \mbox{→ 0} \mbox{→ 0.5} \mbox{→ 1.5} \mbox{→ 3}。予想\qty{3}{m}}
\end{jibun}

\begin{reidai}[速度がへっていくとき]
速度 $v=20-4t$（m/s）でブレーキをかけている車が、$t=0$ で位置0にいます。1秒ごとに予想して、3秒の位置を求めましょう。ほんとうの位置 $x=20t-2t^2$ とくらべましょう。\par
\kangaekata 例題2と同じく、①の式を1秒ずつくり返します。各1秒のはじめの速度は、$t=0$ で20、$t=1$ で16、$t=2$ で12です。\par
\shiki 1秒：$0+20\times1=20$\par
\hspace*{2.6\zw}2秒：$20+16\times1=36$\par
\hspace*{2.6\zw}3秒：$36+12\times1=48$\par
\hspace*{2.6\zw}ほんとう：$20\times3-2\times3^2=60-18=42$\par
\begin{bkcenter}
\begin{bktable}{colspec={Q[c,m,wd=10mm] Q[c,m,wd=26mm] Q[c,m,wd=20mm] Q[c,m,wd=16mm] Q[c,m,wd=22mm]}, rowsep=1pt, cells={font=\small}, row{1}={font=\sffamily\bfseries\footnotesize}}
$t$（s） & はじめの速度 $20-4t$ & 速度×1（m） & 予想（m） & ほんとう $20t-2t^2$ \\
0 & 20 & 20 & 0 & 0 \\
1 & 16 & 16 & 20 & 18 \\
2 & 12 & 12 & 36 & 32 \\
3 & & & 48 & 42 \\
\end{bktable}
\end{bkcenter}
\kotae 予想\qty{48}{m}（ほんとうは\qty{42}{m}。速度がへっていくのに、はじめの大きい速度を使ったので、大きめになる）
\end{reidai}

\begin{jibun}
速度 $v=10-2t$（m/s）の物体が、$t=0$ で位置0にいます。1秒ごとに予想して3秒の位置を求め、ほんとうの位置 $x=10t-t^2$ とくらべましょう。
\bkjans{速度は10、8、6。0 \mbox{→ 10} \mbox{→ 18} \mbox{→ 24}。予想\qty{24}{m}、ほんとうは $30-9=\qty{21}{m}$}
\end{jibun}

\begin{reidai}[接線で近似する]
$\sqrt{9.2}$ を1次近似で求めましょう。\par
\kangaekata 9.2に近くて、$\sqrt{\phantom{9}}$ がすぐわかる数は9です。②の1次近似で、$f(x)=\sqrt{x}$、$a=9$、$h=0.2$ とします。\par
\shiki $f(9)=\sqrt{9}=3$\par
\hspace*{2.6\zw}$f'(9)=\dfrac{1}{2\sqrt{9}}=\dfrac{1}{2\times3}=\dfrac{1}{6}$\quad（5-6）\par
\hspace*{2.6\zw}$\sqrt{9.2}\approx3+\dfrac{1}{6}\times0.2=3+0.0333\cdots=3.033$\par
\kotae およそ3.033（電卓では3.0331…）
\end{reidai}

\begin{jibun}
$\sqrt{16.3}$ を1次近似で求めましょう。
\bkjans{$a=16$、$f'(16)=\dfrac{1}{8}$。$4+\dfrac{1}{8}\times0.3=4.0375$（電卓では4.0373…）}
\end{jibun}

\begin{reidai}[$(1+h)^n\approx1+nh$]
$1.02^5$ を近似しましょう。\par
\kangaekata $1.02=1+0.02$ と見ると、④の $(1+h)^n\approx1+nh$ が使えます。$h=0.02$、$n=5$ を入れます。\par
\shiki $(1+0.02)^5\approx1+5\times0.02=1+0.1=1.1$\par
\kotae およそ1.1（ほんとうは1.10408…）
\end{reidai}

\begin{jibun}
$0.99^{10}$ を近似しましょう。$0.99=1+(-0.01)$ なので、$h=-0.01$ です。
\bkjans{$1+10\times(-0.01)=0.9$（ほんとうは0.9043…）}
\end{jibun}

\begin{reidai}[マイナスや分数の指数でも使える]
（1）$\dfrac{1}{1.02}$\quad（2）$\sqrt{1.04}$\quad（3）$\dfrac{1}{0.96}$ を、$(1+h)^n\approx1+nh$ で近似しましょう。\par
\kangaekata $\dfrac{1}{1+h}=(1+h)^{-1}$、$\sqrt{1+h}=(1+h)^{\frac{1}{2}}$ と書きかえます（④）。（3）は $0.96=1+(-0.04)$ なので、$h=-0.04$、$n=-1$ です。まず、$(1+h)^n$ の形に書きかえて、$n$ と $h$ を読みとります。\par
\begin{bkcenter}
\begin{bktable}{colspec={Q[c,m,wd=12mm] Q[c,m,wd=34mm] Q[c,m,wd=12mm] Q[c,m,wd=14mm]}, rowsep=1pt, cells={font=\small}, row{1}={font=\sffamily\bfseries\footnotesize}}
 & 書きかえ & $n$ & $h$ \\
（1） & $\dfrac{1}{1.02}=(1+0.02)^{-1}$ & $-1$ & $0.02$ \\
（2） & $\sqrt{1.04}=(1+0.04)^{\frac{1}{2}}$ & $\dfrac{1}{2}$ & $0.04$ \\
（3） & $\dfrac{1}{0.96}=(1-0.04)^{-1}$ & $-1$ & $-0.04$ \\
\end{bktable}
\end{bkcenter}
\shiki（1）$n=-1$、$h=0.02$：$1+(-1)\times0.02=0.98$\par
\hspace*{2.6\zw}（2）$n=\dfrac{1}{2}$、$h=0.04$：$1+\dfrac{1}{2}\times0.04=1.02$\par
\hspace*{2.6\zw}（3）$n=-1$、$h=-0.04$：$1+(-1)\times(-0.04)=1.04$\par
\kotae（1）およそ0.98（ほんとうは0.98039…）\par
\hspace*{2.6\zw}（2）およそ1.02（ほんとうは1.01980…）\par
\hspace*{2.6\zw}（3）およそ1.04（ほんとうは1.04166…）
\end{reidai}

\begin{jibun}
（1）$\dfrac{1}{1.05}$\quad（2）$\sqrt{1.06}$\quad（3）$\sqrt{0.98}$ を近似しましょう。
\bkjans{（1）\nobreak $1-0.05=0.95$（ほんとうは0.95238…）\quad（2）\nobreak $1+0.03=1.03$（ほんとうは1.02956…）\quad（3）\nobreak $1+\dfrac{1}{2}\times(-0.02)=0.99$（ほんとうは0.98994…）}
\end{jibun}

\begin{reidai}[ずれの見積もり]
球の半径をはかったら\qty{5.0}{cm}で、$\pm\qty{0.1}{cm}$ のずれがあります。このとき、球の体積 $V=\dfrac{4}{3}\pi r^3$ のずれを見積もりましょう。\par
\kangaekata ⑤の $\Delta V\approx V'\,\Delta r$ を使います。まず $V$ を $r$ で微分し（5-4）、$r=5$、$\Delta r=0.1$ を入れます。$\dfrac{4}{3}\times3=4$ なので、$V'=4\pi r^2$ です。\par
\shiki $V'=4\pi r^2=4\pi\times25=100\pi$\par
\hspace*{2.6\zw}$\Delta V\approx100\pi\times0.1=10\pi\approx31.4$\par
\hspace*{2.6\zw}体積：$\dfrac{4}{3}\pi\times125\approx524$。割合：$31.4\div524\approx0.06$\par
\kotae およそ $\pm\qty{31}{cm^3}$（約6\,\%。半径のずれ2\,\%の3倍）
\end{reidai}

\begin{jibun}
正方形の1辺をはかったら\qty{20.0}{cm}で、$\pm\qty{0.2}{cm}$ のずれがあります。面積のずれを見積もりましょう。また、それは面積の何\,\%ですか。
\bkjans{$S'=2x=40$、$\Delta S\approx40\times0.2=\qty{8}{cm^2}$。$8\div400=0.02$ で2\,\%（1辺のずれ1\,\%の2倍）}
\end{jibun}

\begin{reidai}[ずれを何\,\%以内にすればよいか]
1辺がおよそ\qty{20}{cm}の立方体の体積を、ずれ1.5\,\%以内で求めたいと思います。1辺をはかるときのずれは、何cm以内にすればよいでしょう。\par
\kangaekata ⑤で調べたように、体積（3乗の量）のずれの割合は、1辺のずれの割合の3倍です。逆に考えて、1辺のずれの割合は、体積のずれの割合の $\dfrac{1}{3}$ にすればよいのです。\par
\shiki $\dfrac{\Delta V}{V}\approx3\times\dfrac{\Delta x}{x}\leqq0.015$ より $\dfrac{\Delta x}{x}\leqq0.005$\par
\hspace*{2.6\zw}$\Delta x\leqq0.005\times20=0.1$\par
\kotae \qty{0.1}{cm}以内（\qty{1}{mm}以内）
\end{reidai}

\begin{jibun}
半径がおよそ\qty{10}{cm}の円の面積を、ずれ1\,\%以内で求めたいと思います。半径をはかるときのずれは、何cm以内にすればよいでしょう。
\bkjans{面積（2乗の量）のずれの割合は2倍。$2\times\dfrac{\Delta r}{r}\leqq0.01$ より $\dfrac{\Delta r}{r}\leqq0.005$、$\Delta r\leqq\qty{0.05}{cm}$}
\end{jibun}

%% ---------------- まちがえやすいところ・物理工学 ----------------
\bksection{tsumazuki}{まちがえやすいところをチェック}

\begin{tsumazuki}{$h$ が大きくても近似を使ってしまう}
\bkng{$\sqrt{9}\approx\sqrt{4}+\dfrac{1}{4}\times5=3.25$}
\bkok{$\sqrt{9}=3$。$h=5$ は大きすぎて、近似が使えない}
接線がグラフによりそっているのは、接点の近くだけです（②の表）。1次近似は、$h$ が小さいときにだけ使います。目安は、$h$ が $a$ よりずっと小さいときです。
\end{tsumazuki}

\begin{tsumazuki}{「$\approx$」を「$=$」と書いてしまう}
\bkng{$\sqrt{4.1}=2.025$}
\bkok{$\sqrt{4.1}\approx2.025$}
近似した値は、ほんとうの値とは少しちがうので、「$\approx$」で書きます。どのくらいずれているかを、いつも意識しましょう。
\end{tsumazuki}

\begin{tsumazuki}{くり返しの計算で、おわりの速度を使ってしまう}
\bkng{$v=2t$、1秒ごと：1秒の位置を $0+2\times1=2$}
\bkok{1秒の位置は $0+0\times1=0$（0秒の速度を使う）}
「次の位置＝今の位置＋\textbf{今の}速度×$\Delta t$」です。区間のはじめ（今）の速度を使います。表計算なら、1つ上の行の速度を使います。
\end{tsumazuki}

\begin{butsuri}
\begin{itemize}[leftmargin=1.2\zw, itemsep=2pt, topsep=0pt, label={\color{bkPurple}●}]
\item \textbf{カーナビの推測航法}\\
トンネルの中では、最後にわかった位置に「速度×時間」を少しずつたして、位置を予想しています。向きの変化は、回転の速さをはかる半導体のセンサー（ジャイロセンサー）ではかります。船や飛行機でも、昔から使われてきた方法です。
\item \textbf{ゲームの物理}\\
ゲームの中のボールやキャラクターは、1秒に60回ほど、「次の位置＝今の位置＋速度×$\dfrac{1}{60}$秒」と計算して動かしています。オイラー法そのものです。
\item \textbf{天気予報}\\
空気の動きや温度を、数十秒〜数分ごとに「次の状態＝今の状態＋変化の速さ×時間」と計算して、明日の天気を予想しています（第6巻）。
\end{itemize}
\end{butsuri}

%% ---------------- 章末問題 ----------------
\bkendproblems
\bkproblemgroup{基本}{2}

\bkproblem 位置\qty{30}{m}の所を、速度\qty{-4}{m/s}で動いている物体の、0.5秒後の位置を予想しましょう。
\bkpans{$30+(-4)\times0.5=\qty{28}{m}$}

\bkproblem $\sqrt{25.5}$ を1次近似で求めましょう。
\bkpans{$f'(25)=\dfrac{1}{10}$。$5+\dfrac{1}{10}\times0.5=5.05$}

\bkproblem $1.03^4$ を近似しましょう。
\bkpans{$1+4\times0.03=1.12$}

\bkproblem 円の半径をはかったら\qty{10.0}{cm}で、$\pm\qty{0.1}{cm}$ のずれがあります。このとき、円の面積 $S=\pi r^2$ のずれを見積もりましょう。
\bkpans{$S'=2\pi r=20\pi$、$\Delta S\approx20\pi\times0.1=2\pi\approx\qty{6.3}{cm^2}$}

\bkproblem $\dfrac{1}{0.98}$ と $\sqrt{0.96}$ を、$(1+h)^n\approx1+nh$ で近似しましょう。
\bkpans{$h=-0.02$、$n=-1$ で1.02。$h=-0.04$、$n=\frac{1}{2}$ で0.98}

\bkproblem 速度 $v=6-2t$（m/s）の物体が、$t=0$ で位置\qty{10}{m}にいます。1秒ごとに予想して2秒の位置を求め、ほんとうの位置 $x=10+6t-t^2$ とくらべましょう。
\bkpans{速度は6、4。10 \mbox{→ 16} \mbox{→ 20}。予想\qty{20}{m}、ほんとうは $10+12-4=\qty{18}{m}$}

\bkproblemgroup{チャレンジ}{3}

\bkproblem 速度が $v=3t^2$（m/s）の物体が、$t=0$ で位置0にいます。0.5秒ごとに予想して、2秒の位置を求めましょう。ほんとうの位置は $t^3$ です。
\bkpans{速度は0、0.75、3、6.75。0 \mbox{→ 0} \mbox{→ 0.375} \mbox{→ 1.875} \mbox{→ 5.25}。予想\qty{5.25}{m}（ほんとうは\qty{8}{m}）}

\bkproblem $f(x)=x^3$ で、$f(2.01)$ を1次近似で求めましょう。ほんとうの値とのずれも計算しましょう。
\bkpans{$f'(2)=12$。$8+12\times0.01=8.12$。ほんとうは8.120601なので、ずれは0.000601}

\bkproblem 半径およそ\qty{5.0}{cm}の球の体積を、ずれ3\,\%以内で求めたいと思います。半径をはかるときのずれは、何cm以内にすればよいでしょう。
\bkpans{$\dfrac{\Delta V}{V}\approx\dfrac{3\Delta r}{r}\leqq0.03$ より $\dfrac{\Delta r}{r}\leqq0.01$、$\Delta r\leqq\qty{0.05}{cm}$}

\bkproblem 落ちるボールの速度は $v=10t$（m/s）です。位置0から、0.25秒ごとに「次の位置＝今の位置＋今の速度×0.25」と計算して、1秒の位置を予想しましょう。ほんとうの位置 $5t^2$ とくらべましょう。
\bkpans{速度は0、2.5、5、7.5。0 \mbox{→ 0} \mbox{→ 0.625} \mbox{→ 1.875} \mbox{→ 3.75}。予想\qty{3.75}{m}（ほんとうは\qty{5}{m}）}

\begin{furikaeri}
\begin{checklist}
\item 次の位置を、今の位置＋速度×$\Delta t$ で予想できる。\hfill{\footnotesize\sffamily\color{bkTheme}→\pageref{bk:5-8:tsugi}ページ}
\item 接線を使って、1次近似ができる。\hfill{\footnotesize\sffamily\color{bkTheme}→\pageref{bk:5-8:sessen}ページ}
\item 測定のずれが、計算にひびく大きさを見積もれる。\hfill{\footnotesize\sffamily\color{bkTheme}→\pageref{bk:5-8:gosa}ページ}
\end{checklist}
\end{furikaeri}

\bkfukidashi{sora}{速度×時間をどんどんたしていくと、進んだ距離がわかる。これって、次の章の「グラフの面積」につながるのかな？}


%% ======================== 第2部 積分 ========================
\bkPartOpener{2}{第5巻\hspace{0.8em}第2部}{積分}%
  {172/{速さのグラフから、進んだ距離はわかる？},
   160/{曲がった形の面積は、どうやって求める？},
   148/{電気のメーターは、何をたし合わせている？}}%
  {% 細い長方形（白い線）
   \begin{scope}[white, line width=1pt]
     \foreach \i in {0,...,7} {\pgfmathsetmacro\hh{4+0.45*(\i+1)*(\i+1)}
       \draw (120+\i*6,112) rectangle ++(6,\hh);}
     \draw[line width=1.4pt] (120,112) parabola (168,112+32.8);
   \end{scope}}%
  {この部では、細い長方形をたくさんたし合わせて、たまった量を計算する道具をつくるよ。それが積分。最後に、積分が微分の逆になっていることを発見するよ。}
\bkPartContents{2}{%% 自動生成（gen.py）
\bktocchap{5-9}{グラフの面積は進んだ距離}{速さのグラフの面積が、進んだ道のりになる}{2}{2}{2}{196}{C}
\bktocchap{5-10}{細い長方形に分けてたそう}{曲がった形の面積を、細い長方形で求める}{2}{2}{4}{228}{C}
\bktocchap{5-11}{積分は微分の逆}{面積がふえる速さは、その場所の高さ}{2}{3}{3}{262}{C}
\bktocchap{5-12}{物理での積分：仕事と流量}{ばねをのばす仕事は、三角形の面積}{2}{3}{3}{284}{C}
\bktocchap{5-13}{積分定数は「はじめの状態」}{スタート地点がわからないと、位置は決まらない}{2}{3}{3}{310}{C}
}%
{速さから進んだ距離を、流量からたまった水の量を、電力から電力量を計算できるようになります。「はじめの状態」と組み合わせれば、加速度から位置まで求められます。}
\bkpartbegin{2}{積分}
%% 5-9 グラフの面積は進んだ距離（第11回：積分の最初の章として書き足した。ページ数は自由・偶数）
\bkchapter{5}{9}{グラフの面積は進んだ距離}{\bklevel{2}}{グラフの面積は進んだ距離}
\bkmichishirube{2-3,4-5,5-5,5-8}{5-10,5-11,5-13}

\begin{goal}
\begin{itemize}[label={\color{bkTheme}●}, leftmargin=1.6\zw, itemsep=0pt, topsep=1pt]
\item \bklevel{2}\enspace 速さが変わるときも、「細かく分けて、かけて、たす」で進んだ距離を求められる。
\item \bklevel{2}\enspace 速度のグラフ（$v$-$t$ グラフ）を長方形・三角形・台形に分けて、進んだ距離を計算できる。
\item \bklevel{2}\enspace 速度がマイナスのとき、面積の符号の意味（変位と道のり）を説明できる。
\end{itemize}
\end{goal}

\begin{hatena}
\noindent\begin{minipage}[c]{0.58\linewidth}\raggedright
トラックには、走るあいだの速さを記録する機械（タコグラフ）がついています。記録されるのは速さだけです。

右は、ある日の記録を簡単にしたグラフです（横は時間 h、たては速さ km/h）。この記録だけから、\mbox{何km}走ったかがわかります。どうやって求めるのでしょう。
\end{minipage}\hfill
\begin{minipage}[c]{0.4\linewidth}\centering
\begin{tikzpicture}
\begin{axis}[bkplot, width=50mm, height=38mm, xmin=0, xmax=6.8, ymin=0, ymax=74, xtick={1,2,3,4,5,6}, ytick={20,40,60}, xlabel={$t$}, ylabel={$v$}]
  \addplot[bkcurve, line width=1.3pt] coordinates {(0,0) (0.5,40) (2,40) (2.5,0) (3,0) (3.5,60) (5.5,60) (6,0)};
\end{axis}
\end{tikzpicture}
\end{minipage}
\end{hatena}

2-3では、速さが一定のとき、「速さ×時間＝道のり」が、速さと時間のグラフの長方形の面積になることを学びました。でも、ほんとうの乗り物の速さは、いつも変わっています。

この章では、速さが変わるときにも、グラフの面積が進んだ距離になることを、1段ずつ確かめます。カギは「\textbf{細かく分けて、かけて、たす}」の3つの段です。これが、\textbf{積分}（せきぶん）の考え方の始まりです。

%% ---------------- やってみよう・思い出そう ----------------
\begin{yattemiyou}[方眼のマスを数えて距離を求めよう]
自転車の速さを、2秒ごとに記録しました。記録と記録のあいだは、速さがまっすぐ（一定の割合で）変わったとします。
\begin{bkcenter}
\begin{bktable}{colspec={Q[l,m,wd=26mm] *{6}{Q[c,m,wd=10mm]}}, rowsep=1.5pt}
時刻 $t$（s） & 0 & 2 & 4 & 6 & 8 & 10 \\
速さ $v$（m/s） & 0 & 2 & 4 & 4 & 4 & 2 \\
\end{bktable}
\end{bkcenter}
\begin{enumerate}[label={\bkcircled[bkGreen]{\arabic*}}, leftmargin=2.4\zw, itemsep=1pt, topsep=3pt]
\item 方眼紙に、横を時刻、たてを速さにして点をとり、となりの点を直線で結ぼう（$v$-$t$ グラフ）。1マスは、たて\qty{1}{m/s}、横\qty{1}{s}にしよう。
\item グラフと横軸のあいだにある方眼のマスを数えよう。ななめの線で半分に切られたマスは、2つで1マスと数えよう。
\item 1マスは、何mの距離にあたるかな？　10秒で何m進んだかな？
\end{enumerate}
\end{yattemiyou}

\noindent やってみようのグラフをかくと、下のようになります。

\begin{bkfigure}
\begin{tikzpicture}
\begin{axis}[bkplot, width=112mm, height=52mm, xmin=0, xmax=11.2, ymin=0, ymax=5.6, xtick={1,...,10}, ytick={1,...,5}, xlabel={$t$（s）}, ylabel={$v$（m/s）}]
  \fill[bkYellow!40] (axis cs:0,0) -- (axis cs:4,4) -- (axis cs:4,0) -- cycle;
  \fill[bkBlue!20] (axis cs:4,0) rectangle (axis cs:8,4);
  \fill[bkGreen!25] (axis cs:8,0) -- (axis cs:8,4) -- (axis cs:10,2) -- (axis cs:10,0) -- cycle;
  \addplot[bkcurve, line width=1.4pt, mark=*, mark size=1.4pt] coordinates {(0,0) (2,2) (4,4) (6,4) (8,4) (10,2)};
  \node[font=\sffamily\footnotesize, fill=bkYellow!40, inner sep=1pt] at (axis cs:3.05,0.5) {8マス};
  \node[font=\sffamily\footnotesize, fill=bkBlue!20, inner sep=1pt] at (axis cs:6,2) {16マス};
  \node[font=\sffamily\footnotesize, fill=bkGreen!25, inner sep=1pt] at (axis cs:9,1) {6マス};
\end{axis}
\end{tikzpicture}
\bkfigcaption{0〜4秒は三角形、4〜8秒は長方形、8〜10秒は台形。}
\end{bkfigure}

\noindent\begin{minipage}[c]{0.66\linewidth}
\textbf{1マスは何m？}\quad 1マスは、たてが\qty{1}{m/s}、横が\qty{1}{s}です。「1秒に\qty{1}{m}進む速さで、1秒走った」ということなので、進んだ距離は\qty{1}{m}です。だから、\textbf{1マスは\qty{1}{m}}にあたります。
\end{minipage}\hfill
\begin{minipage}[c]{0.3\linewidth}\centering
\begin{tikzpicture}[x=1mm, y=1mm]
  \fill[bkBlue!20] (0,0) rectangle (14,14);
  \draw[bkBlue!80!black, line width=0.8pt] (0,0) rectangle (14,14);
  \node[font=\sffamily\small] at (7,7) {\qty{1}{m}};
  \node[font=\sffamily\footnotesize, anchor=north] at (7,-1.5) {横 \qty{1}{s}};
  \node[font=\sffamily\footnotesize, anchor=east] at (-1.5,7) {たて \qty{1}{m/s}};
\end{tikzpicture}
\end{minipage}

\noindent\begin{minipage}[c]{0.62\linewidth}\raggedright
\textbf{マスを数える}\quad 0〜4秒の三角形を大きくしたのが右の図です。まるごとのマス（黄色）が $0+1+2+3=6$ こ、ななめの線で半分に切られたマス（オレンジ）が4こあります。半分のマスは2つで1マスなので、4こで2マスです。合わせて $6+2=8$ マスです。
\end{minipage}\hfill
\begin{minipage}[c]{0.34\linewidth}\centering
\begin{tikzpicture}[x=7mm, y=7mm]
  \foreach \k in {1,2,3} {\foreach \j in {1,...,\k} {\fill[bkYellow!55] (\k,\j-1) rectangle (\k+1,\j);}}
  \foreach \k in {0,1,2,3} {\fill[bkOrange!40] (\k,\k) -- (\k+1,\k) -- (\k+1,\k+1) -- cycle;}
  \draw[bkLine, line width=0.4pt] (0,0) grid (4,4);
  \draw[bkBlue!80!black, line width=1.2pt] (0,0) -- (4,4);
  \draw[bkInk, line width=0.6pt] (0,0) -- (4.3,0) (0,0) -- (0,4.3);
  \foreach \k in {1,2,3,4} \node[font=\sffamily\scriptsize, anchor=north] at (\k,-0.1) {\k};
\end{tikzpicture}
\end{minipage}

\medskip
4〜8秒の長方形は、たて4マス、横4マスで $4\times4=16$ マスです。8〜10秒の台形は、まるごとのマスが $3+2=5$ こと、半分のマスが2こ（＝1マス）で6マスです。

全部で $8+16+6=30$ マスです。1マスが\qty{1}{m}なので、自転車は10秒で\qty{30}{m}進んだことになります。グラフの面積（マスの数）が、進んだ距離になっています。なぜそうなるのかを、この章で確かめていきます。

\begin{omoidasou}
\bkrecall{2-3}{道のり＝速さ×時間。速さが一定なら、$v$-$t$ グラフ（速さと時間のグラフ）の長方形の面積が道のり。}
\bkrecall{4-5}{三角形の面積＝底辺×高さ÷2（同じ三角形を2つ合わせると長方形になるから）。台形の面積＝（上底＋下底）×高さ÷2（同じ台形を2つ、さかさに合わせると平行四辺形になるから）。}
\bkrecall{5-5}{加速度は「1秒ごとに速度が何 m/s 変わるか」。加速度が一定の $a$ なら $v=v_0+at$、$x=x_0+v_0t+\dfrac{1}{2}at^2$。}
\bkrecall{5-8}{短い時間 $\Delta t$ のあいだは、速度はほとんど変わらない。だから $\Delta x\approx v\,\Delta t$（進む距離$\approx$速度×時間）。}
\bkrecall{2-10}{換算は「1をかける」。\qty{1000}{m/km}（1 kmあたり1000 m）、\qty{3600}{s/h}（1時間あたり3600秒）、\qty{60}{s/min}（1分あたり60秒）。}
\end{omoidasou}

%% ---------------- 数学の道具 ----------------
\bksection{doguu}{数学の道具}

\bksubsection{1}{速さが一定なら、長方形の面積}

\label{bk:5-9:chouhou}車が、\qty{6}{m/s}（1秒に\qty{6}{m}進む速さ）で2秒走ります。1秒めに\qty{6}{m}、2秒めにも\qty{6}{m}進むので、合わせて $6+6=6\times2=\qty{12}{m}$ です。

\begin{bkfigure}
\begin{tikzpicture}[x=1mm, y=1mm]
  % 数直線（1 m を 7 mm）
  \draw[bk axis] (-3,0) -- (92,0);
  \node[font=\sffamily\small, anchor=west] at (93,0) {m};
  \foreach \p in {0,6,12} {\draw[bkInk, line width=0.8pt] ({7*\p},-1.5) -- ({7*\p},1.5);
    \node[font=\sffamily\small, anchor=north] at ({7*\p},-2.5) {\p};}
  \foreach \p in {1,...,11} \draw[bkInk, line width=0.4pt] ({7*\p},-0.8) -- ({7*\p},0.8);
  \foreach \p/\t in {0/0,6/1,12/2} {\bkcar{({7*\p-3.6},3.5)}{0.3}{bkRoute2}
    \node[font=\sffamily\footnotesize, anchor=south] at ({7*\p+0.9},10) {\t 秒};}
  \draw[bk arrow, bkOrange] (3,-8) -- (39,-8);
  \draw[bk arrow, bkOrange] (45,-8) -- (81,-8);
  \node[font=\sffamily\footnotesize, anchor=north] at (21,-9.5) {1秒で\qty{6}{m}};
  \node[font=\sffamily\footnotesize, anchor=north] at (63,-9.5) {1秒で\qty{6}{m}};
\end{tikzpicture}
\bkfigcaption{\qty{6}{m/s}で2秒。1秒ごとに\qty{6}{m}ずつ進む。}
\end{bkfigure}

\noindent\begin{minipage}[c]{0.46\linewidth}\centering
\begin{tikzpicture}
\begin{axis}[bkplot, width=58mm, height=42mm, xmin=0, xmax=7.6, ymin=0, ymax=7.4, xtick={1,...,7}, ytick={3,6}, xlabel={$t$（s）}, ylabel={$v$（m/s）}]
  \fill[bkBlue!18] (axis cs:0,0) rectangle (axis cs:2,6);
  \fill[bkOrange!25] (axis cs:2,0) rectangle (axis cs:7,3);
  \addplot[bkcurve, line width=1.4pt] coordinates {(0,6) (2,6)};
  \addplot[bkcurve, line width=1.4pt] coordinates {(2,3) (7,3)};
  \node[font=\sffamily\small, fill=bkBlue!18, inner sep=1pt] at (axis cs:1,3.6) {\qty{12}{m}};
  \node[font=\sffamily\small, fill=bkOrange!25, inner sep=1pt] at (axis cs:4.5,1.5) {\qty{15}{m}};
\end{axis}
\end{tikzpicture}
\end{minipage}\hfill
\begin{minipage}[c]{0.5\linewidth}\raggedright
$v$-$t$ グラフでは、たて6（速さ）、横2（時間）の長方形の面積 $6\times2=12$ が、進んだ距離です。長方形を1秒ずつのたての列に切ると、1列が「1秒で進む距離」になっています。

単位も計算してみましょう。秒（s）どうしが約分されて、距離の単位 m が残ります。
\[ \qty{6}{m/s}\times\qty{2}{s}=6\,\frac{\mathrm{m}}{\cancel{\mathrm{s}}}\times2\,\cancel{\mathrm{s}}=\qty{12}{m} \]
\end{minipage}

\medskip
左のグラフのように、はじめの2秒は\qty{6}{m/s}、次の5秒は\qty{3}{m/s}で進んだとき、距離は
\[ 6\times2+3\times5=12+15=\qty{27}{m} \]
です。速さがとちゅうで変わっても、速さが一定の部分ごとに長方形に分けて、たし合わせればよいのです。

\bksubsection{2}{速さが変わり続けるとき：細かく分けて考える}

\label{bk:5-9:wakeru}止まっていた車が走り出して、1秒ごとに\qty{2}{m/s}ずつ速くなります（加速度\qty{2}{m/s^2}、5-5）。$t$ 秒後の速さは $v=2t$（m/s）です。1秒後は2、2秒後は4、4秒後は\qty{8}{m/s}です。

この車は、4秒で何m進むでしょう。

「速さ×時間」を使いたくても、速さがずっと変わっているので、どの速さをかければよいかわかりません。0をかければ0、8をかければ32。どちらもおかしいですね。

そこで、5-8の考え方を使います。\textbf{短い時間のあいだなら、速さはほとんど変わらない}ので、そのあいだは「速さ×時間」で距離を求めてよいのです。次の3つの段で計算します。

\begin{bkfigure}
\begin{tikzpicture}[x=1mm, y=1mm]
  \node[bk node, text width=30mm] (s1) at (0,0) {\textbf{① 分ける}\\4秒を1秒ずつ\\4つの区間に分ける};
  \node[bk node, text width=30mm] (s2) at (45,0) {\textbf{② かける}\\区間ごとに\\速さ×1秒};
  \node[bk node, text width=30mm] (s3) at (90,0) {\textbf{③ たす}\\4つの距離を\\全部たす};
  \draw[bk link] (s1) -- (s2);
  \draw[bk link] (s2) -- (s3);
\end{tikzpicture}
\bkfigcaption{速さが変わるときの距離の求め方。この3つの段を「積分」の考え方という。}
\end{bkfigure}

②で、区間のあいだの速さには、ひとまず\textbf{区間のはじめの速さ}を使います。たとえば1〜2秒の区間は、はじめ（1秒）の速さ\qty{2}{m/s}で1秒進んだと考えて、$2\times1=\qty{2}{m}$ です。

\begin{bkcenter}
\begin{bktable}{colspec={Q[l,m,wd=34mm] *{4}{Q[c,m,wd=14mm]} Q[c,m,wd=14mm]}, rowsep=1.5pt}
区間（s） & 0〜1 & 1〜2 & 2〜3 & 3〜4 & 合計 \\
はじめの速さ（m/s） & 0 & 2 & 4 & 6 & \\
速さ×1秒（m） & 0 & 2 & 4 & 6 & {\bfseries 12} \\
\end{bktable}
\bktabcaption{区間のはじめの速さを使うと、合計は\qty{12}{m}。}
\end{bkcenter}

\noindent 1つ1つの「速さ×1秒」は、$v$-$t$ グラフでは、はば1、高さが速さの\textbf{細い長方形}の面積です。それをならべると、階段の形になります（下の図の左）。0〜1秒の長方形は、高さが0なので見えません。

\begin{bkfigure}
\begin{tikzpicture}
\begin{axis}[bkplot, name=w1, width=60mm, height=54mm, xmin=0, xmax=4.7, ymin=0, ymax=9.4, xtick={1,2,3,4}, ytick={2,4,6,8}, xlabel={$t$}, ylabel={$v$}, title={\sffamily\small はじめの速さ（小さめ）}]
  \foreach \k in {1,2,3} {\edef\tmp{\noexpand\fill[bkBlue!25, draw=bkBlue!70, line width=0.5pt] (axis cs:\k,0) rectangle (axis cs:\k+1,{2*\k});}\tmp}
  \foreach \k in {0,1,2,3} {\edef\tmp{\noexpand\fill[bkOrange!50] (axis cs:\k,{2*\k}) -- (axis cs:\k+1,{2*\k}) -- (axis cs:\k+1,{2*\k+2}) -- cycle;}\tmp}
  \addplot[bkcurve, domain=0:4.2] {2*x};
  \foreach \k in {1,2,3} {\edef\tmp{\noexpand\node[font=\noexpand\sffamily\noexpand\footnotesize] at (axis cs:\k+0.5,\k) {\the\numexpr2*\k\relax};}\tmp}
\end{axis}
\begin{axis}[bkplot, at={(w1.east)}, anchor=west, xshift=12mm, width=60mm, height=54mm, xmin=0, xmax=4.7, ymin=0, ymax=9.4, xtick={1,2,3,4}, ytick={2,4,6,8}, xlabel={$t$}, ylabel={$v$}, title={\sffamily\small おわりの速さ（大きめ）}]
  \foreach \k in {0,1,2,3} {\edef\tmp{\noexpand\fill[bkBlue!25, draw=bkBlue!70, line width=0.5pt] (axis cs:\k,0) rectangle (axis cs:\k+1,{2*\k+2});}\tmp}
  \foreach \k in {0,1,2,3} {\edef\tmp{\noexpand\fill[bkRed!35] (axis cs:\k,{2*\k}) -- (axis cs:\k,{2*\k+2}) -- (axis cs:\k+1,{2*\k+2}) -- cycle;}\tmp}
  \addplot[bkcurve, domain=0:4.2] {2*x};
  \node[font=\sffamily\footnotesize] at (axis cs:0.72,0.5) {2};
  \foreach \k in {1,2,3} {\edef\tmp{\noexpand\node[font=\noexpand\sffamily\noexpand\footnotesize] at (axis cs:\k+0.5,\k) {\the\numexpr2*\k+2\relax};}\tmp}
\end{axis}
\end{tikzpicture}
\bkfigcaption{左：はじめの速さの階段（合計12）。オレンジは、数えもれたすき間。\\右：おわりの速さの階段（合計20）。赤は、はみ出した分。}
\end{bkfigure}

\noindent\textbf{12 mは少なすぎる}\quad 車は、どの区間でも、とちゅうで速くなっています。たとえば1〜2秒のあいだ、速さは2より大きく、4より小さい値です。それなのに、ずっと2のままとして計算したので、ほんとうより少なくなります。グラフでは、階段と直線のあいだのオレンジのすき間が、数えもれた分です。

\noindent\textbf{おわりの速さなら多すぎる}\quad 反対に、区間の\textbf{おわりの速さ}（1秒で2、2秒で4、3秒で6、4秒で8）を使うと、合計は
\[ 2+4+6+8=\qty{20}{m} \]
です。こんどはずっと速いとして計算したので、ほんとうより多くなります（右の図の赤い部分がはみ出した分）。

だから、ほんとうの距離は、12 mより多く、20 mより少ないことがわかります。
\[ 12<\text{ほんとうの距離}<20 \]

\bksubsection{3}{区間を短くすると、階段は三角形に近づく}

\label{bk:5-9:chikazuku}区間をもっと短くしてみましょう。区間のはばを $\Delta t$（デルタ・ティー）と書きます。$\Delta t=0.5$ なら、4秒は8つの区間に分かれます。1つめの区間のはじめの速さは0、2つめは $2\times0.5=1$、3つめは $2\times1=2$……です。

\begin{align*}
\text{はじめの速さ：}\quad&(0+1+2+3+4+5+6+7)\times0.5=28\times0.5=14\\
\text{おわりの速さ：}\quad&(1+2+3+4+5+6+7+8)\times0.5=36\times0.5=18
\end{align*}
0.5は、どの区間にも共通の「×0.5秒」なので、先に速さをたしてから、まとめて0.5をかけました（分配法則、2-11）。同じように $\Delta t$ を小さくして計算すると、次の表になります。

\begin{bkcenter}
\begin{bktable}{colspec={Q[c,m,wd=17mm] Q[c,m,wd=14mm] Q[c,m,wd=30mm] Q[c,m,wd=30mm] Q[c,m,wd=16mm]}, rowsep=1.5pt}
$\Delta t$（s） & 区間の数 & はじめの速さの和（m） & おわりの速さの和（m） & 2つの差 \\
1 & 4 & 12 & 20 & 8 \\
0.5 & 8 & 14 & 18 & 4 \\
0.25 & 16 & 15 & 17 & 2 \\
0.1 & 40 & 15.6 & 16.4 & 0.8 \\
0.01 & 400 & 15.96 & 16.04 & 0.08 \\
\end{bktable}
\bktabcaption{区間を短くすると、2つの和はどちらも16に近づく。}
\end{bkcenter}

$\Delta t$ を小さくするほど、小さめの和は大きくなり、大きめの和は小さくなって、2つの差は0に近づきます。2つの和は、どちらも16に近づいていきます。ほんとうの距離は、いつもこの2つのあいだにあるので、\textbf{ほんとうの距離は\qty{16}{m}}です。

図でも見てみましょう。はじめの速さの階段で、区間を短くしていきます。

\begin{bkfigure}
\begin{tikzpicture}[x=1mm, y=1mm]
  \foreach \n/\px/\dtv/\sm/\gp in {4/0/1/12/4, 8/46/0.5/14/2, 16/92/0.25/15/1} {
    \begin{scope}[shift={(\px,0)}]
      \foreach \k in {1,...,\n} {
        \fill[bkBlue!25] ({32*(\k-1)/\n},0) rectangle ({32*\k/\n},{24*(\k-1)/\n});
        \fill[bkOrange!55] ({32*(\k-1)/\n},{24*(\k-1)/\n}) -- ({32*\k/\n},{24*(\k-1)/\n}) -- ({32*\k/\n},{24*\k/\n}) -- cycle;}
      \foreach \k in {2,...,\n} \draw[bkBlue!70, line width=0.3pt] ({32*(\k-1)/\n},0) rectangle ({32*\k/\n},{24*(\k-1)/\n});
      \draw[bkBlue!80!black, line width=1pt] (0,0) -- (33,24.75);
      \draw[bkInk, line width=0.6pt, -{Stealth[length=3pt]}] (0,0) -- (37,0);
      \draw[bkInk, line width=0.6pt, -{Stealth[length=3pt]}] (0,0) -- (0,28);
      \node[font=\sffamily\scriptsize, anchor=north, inner sep=1pt] at (32,-0.5) {4};
      \node[font=\sffamily\bfseries\footnotesize, anchor=north west, inner sep=0pt] at (0,-5) {$\Delta t=\dtv$};
      \node[font=\sffamily\footnotesize, anchor=north west, inner sep=0pt, align=left] at (0,-10.5) {階段 \sm\ m\\{\color{bkOrange!80!black}すき間 \gp\ m}};
    \end{scope}}
\end{tikzpicture}
\bkfigcaption{区間を短くすると、オレンジのすき間はどんどん小さくなり、階段は三角形に近づく。}
\end{bkfigure}

\noindent\textbf{すき間が0に近づくわけ}\quad オレンジのすき間は、小さな三角形の集まりです。1つの三角形は、底辺が $\Delta t$、高さが「$\Delta t$ のあいだに速さがふえた分」の $2\times\Delta t$ です。
\begin{itemize}[leftmargin=1.6\zw, itemsep=1pt, topsep=2pt, label={\color{bkTheme}●}]
\item 小さな三角形1つの面積は $\Delta t\times2\Delta t\div2=\Delta t\times\Delta t$
\item 三角形の数は、区間の数と同じ $4\div\Delta t$ こ
\item すき間の合計は $(4\div\Delta t)\times\Delta t\times\Delta t=4\times\Delta t$
\end{itemize}
$\Delta t=1$ なら4、$\Delta t=0.5$ なら2、$\Delta t=0.25$ なら1で、図と合っています。$\Delta t$ を半分にすると、すき間の合計も半分になります。$\Delta t$ を0に近づけると、すき間は0に近づき、階段の面積は、直線の下の\textbf{三角形の面積}にかぎりなく近づきます。

\begin{taisetsu}
① 分ける\quad ② かける（速さ×$\Delta t$）\quad ③ たす\\
$\Delta t$ を小さくしていくと、たした値は、$v$-$t$ グラフの下の面積に近づく。\\
その面積が、ほんとうに進んだ距離。
\end{taisetsu}

\begin{column}{表計算で階段の和を計算しよう}
②・③の車（$v=2t$）を、表計算で $\Delta t=0.1$ にして計算します（2-18）。A列に区間のはじめの時刻、B列にその速さ、C列に「速さ×0.1」、D列にC列の合計をつみ上げます。
\begin{bkcenter}
\bksheet{4}{22}{%
1 & はじめの時刻 & 速さ $2t$ & 速さ×0.1 & 合計 \\
2 & 0 & 0 & 0 & 0 \\
3 & 0.1 & 0.2 & 0.02 & 0.02 \\
4 & 0.2 & 0.4 & 0.04 & 0.06 \\
5 & （下へコピー） & & & \\
41 & 3.9 & 7.8 & 0.78 & 15.6 \\}
\end{bkcenter}
B2に \texttt{=2*A2}、C2に \texttt{=B2*0.1}、D2に \texttt{=C2}、D3に \texttt{=D2+C3} と入れて下へコピーすると、41行めの合計は15.6になり、③の表と合います。0.1を0.01にして、行を400行にすると15.96。ほんとうの16に近づいていきます。
\end{column}

\Needspace{6\baselineskip}
\bksubsection{4}{一定の割合で速くなるなら、三角形の面積}

\label{bk:5-9:sankaku}②・③の車の三角形は、底辺4（時間）、高さ8（4秒後の速さ）なので、
\[ 4\times8\div2=\qty{16}{m} \]
表で近づいていった16と、ぴったり同じです。

\noindent\begin{minipage}[c]{0.52\linewidth}\raggedright
\textbf{なぜ÷2なの？}\quad 右の図のように、同じ三角形を2つ合わせると、たて8、横4の長方形になります。長方形の面積 $4\times8=32$ は「ずっと\qty{8}{m/s}で走った距離」で、三角形はその半分です。
\end{minipage}\hfill
\begin{minipage}[c]{0.44\linewidth}\centering
\begin{tikzpicture}
\begin{axis}[bkplot, width=56mm, height=50mm, xmin=0, xmax=4.9, ymin=0, ymax=9.6, xtick={2,4}, ytick={4,8}, xlabel={$t$}, ylabel={$v$}]
  \fill[bkGray!20] (axis cs:0,0) -- (axis cs:0,8) -- (axis cs:4,8) -- cycle;
  \fill[bkBlue!28] (axis cs:0,0) -- (axis cs:4,8) -- (axis cs:4,0) -- cycle;
  \draw[bkGray, line width=0.6pt, dashed] (axis cs:0,8) -- (axis cs:4,8) -- (axis cs:4,0);
  \addplot[bkcurve, domain=0:4.3] {2*x};
  \node[font=\sffamily\footnotesize, fill=bkBlue!28, inner sep=1pt] at (axis cs:2.85,1.6) {\qty{16}{m}};
  \node[font=\sffamily\footnotesize, text=bkGray!80!black] at (axis cs:1.15,6.4) {同じ形};
\end{axis}
\end{tikzpicture}
\end{minipage}

\medskip
\noindent\begin{minipage}[c]{0.44\linewidth}\centering
\begin{tikzpicture}
\begin{axis}[bkplot, width=56mm, height=50mm, xmin=0, xmax=4.9, ymin=0, ymax=9.6, xtick={2,4}, ytick={4,8}, xlabel={$t$}, ylabel={$v$}]
  \fill[bkBlue!28] (axis cs:0,0) rectangle (axis cs:4,4);
  \fill[bkOrange!55] (axis cs:2,4) -- (axis cs:4,4) -- (axis cs:4,8) -- cycle;
  \fill[white] (axis cs:0,0) -- (axis cs:2,4) -- (axis cs:0,4) -- cycle;
  \fill[bkOrange!25] (axis cs:0,0) -- (axis cs:2,4) -- (axis cs:0,4) -- cycle;
  \draw[bkGray, line width=0.6pt, dashed] (axis cs:0,4) -- (axis cs:4,4);
  \addplot[bkcurve, domain=0:4.3] {2*x};
  \node[font=\sffamily\scriptsize, text=bkOrange!80!black] at (axis cs:3.35,5.5) {A};
  \node[font=\sffamily\scriptsize, text=bkOrange!80!black] at (axis cs:0.6,2.85) {B};
\end{axis}
\end{tikzpicture}
\end{minipage}\hfill
\begin{minipage}[c]{0.52\linewidth}\raggedright
\textbf{平均の速さで考える}\quad 速さは0から8まで、まっすぐふえていきます。2秒の速さ4をまん中にして、1秒の速さ2と3秒の速さ6、0.5秒の速さ1と3.5秒の速さ7……と、前と後ろで組にすると、どの組も平均は4です。だから、4秒間の\textbf{平均の速さは $(0+8)\div2=\qty{4}{m/s}$}。
\[ 4\times4=\qty{16}{m} \]
図では、三角形の上の部分Aを切りとって、左上のBにうめると、たて4、横4の長方形になります。2通りの考え方で、同じ16 mになりました。
\end{minipage}

\medskip
\noindent\textbf{文字の式にする}\quad 止まっていた物が、一定の加速度 $a$ で $t$ 秒間速くなるとします。
\begin{itemize}[leftmargin=1.6\zw, itemsep=1pt, topsep=2pt, label={\color{bkTheme}●}]
\item 三角形の底辺は、時間の $t$。
\item 高さは、$t$ 秒後の速さ。1秒に $a$ ずつ、$t$ 秒ふえるので $at$。
\item ÷2 は、三角形が長方形の半分だから。$\div2$ は $\times\dfrac{1}{2}$ と同じ。
\end{itemize}
\[ x=t\times at\times\frac{1}{2}=\frac{1}{2}at^2 \]
これは、5-5で微分して確かめた式 $x=\dfrac{1}{2}at^2$ と同じです。$a=2$、$t=4$ を入れると $\dfrac{1}{2}\times2\times16=16$ で、②・③の車とも合います。

\Needspace{12\baselineskip}
\bksubsection{5}{はじめから動いていれば、台形の面積}

\label{bk:5-9:daikei}\qty{5}{m/s}で走っていた自転車が、一定の割合で速くなり、4秒後に\qty{15}{m/s}になりました。4秒で何m進むでしょう。$v$-$t$ グラフは、たての辺が5と15、はばが4の\textbf{台形}になります。

\noindent\begin{minipage}[c]{0.46\linewidth}\centering
\begin{tikzpicture}
\begin{axis}[bkplot, width=58mm, height=54mm, xmin=0, xmax=5.2, ymin=0, ymax=18, xtick={4}, ytick={5,15}, xlabel={$t$}, ylabel={$v$}]
  \fill[bkBlue!25] (axis cs:0,0) rectangle (axis cs:4,5);
  \fill[bkOrange!40] (axis cs:0,5) -- (axis cs:4,5) -- (axis cs:4,15) -- cycle;
  \draw[bkGray, line width=0.6pt, dashed] (axis cs:4,0) -- (axis cs:4,15);
  \addplot[bkcurve, domain=0:4.4] {5+2.5*x};
  \node[font=\sffamily\footnotesize] at (axis cs:2,2.5) {\qty{20}{m}};
  \node[font=\sffamily\footnotesize] at (axis cs:3.05,8) {\qty{20}{m}};
\end{axis}
\end{tikzpicture}
\end{minipage}\hfill
\begin{minipage}[c]{0.5\linewidth}\raggedright
\textbf{分けて考える}\quad 台形を、下の長方形（青）と上の三角形（オレンジ）に分けます。
\begin{itemize}[leftmargin=1.2\zw, itemsep=1pt, topsep=2pt, label={\color{bkTheme}●}]
\item 長方形：はじめの速さ5のまま4秒走った分。$5\times4=20$
\item 三角形：速くなったことで、よけいに進んだ分。高さは $15-5=10$（ふえた速さ）。$4\times10\div2=20$
\end{itemize}
合わせて $20+20=\qty{40}{m}$ です。
\end{minipage}

\medskip
\noindent\textbf{台形の公式で}\quad このグラフの台形は、横にねています。たての2つの辺（5と15）が上底と下底、横のはば（4秒）が高さです。
\[ (5+15)\times4\div2=20\times4\div2=\qty{40}{m} \]
\textbf{平均の速さで}\quad 速さは5から15までまっすぐふえるので、平均の速さは $(5+15)\div2=\qty{10}{m/s}$。$10\times4=\qty{40}{m}$。3通りとも同じです。

\noindent\textbf{文字の式にする}\quad はじめの速度 $v_0$、加速度 $a$ で $t$ 秒走ると、$t$ 秒後の速度は $v_0+at$ です。長方形は $v_0\times t$、三角形は高さ $at$ なので $\dfrac{1}{2}at^2$（④と同じ）。
\[ x=v_0t+\frac{1}{2}at^2 \]
5-5の $x=x_0+v_0t+\dfrac{1}{2}at^2$ の、$x_0$（はじめの位置）以外の部分です。面積が表すのは、はじめの位置から\textbf{どれだけ変わったか}なので、$x_0$ は出てきません。

\bksubsection{6}{おそくなるときも、同じ}

\label{bk:5-9:osoku}ブレーキをかけて速さがへるとき、$v$-$t$ グラフは右下がりの直線になります。形は、向きがちがうだけの三角形や台形です。

\begin{bkfigure}
\begin{tikzpicture}
\begin{axis}[bkplot, name=o1, width=58mm, height=46mm, xmin=0, xmax=6.2, ymin=0, ymax=24, xtick={5}, ytick={20}, xlabel={$t$}, ylabel={$v$}, title={\sffamily\small 止まるまで（三角形）}]
  \fill[bkBlue!25] (axis cs:0,0) -- (axis cs:0,20) -- (axis cs:5,0) -- cycle;
  \addplot[bkcurve, domain=0:5] {20-4*x};
  \node[font=\sffamily\footnotesize] at (axis cs:1.5,5) {\qty{50}{m}};
\end{axis}
\begin{axis}[bkplot, at={(o1.east)}, anchor=west, xshift=12mm, width=58mm, height=46mm, xmin=0, xmax=5.2, ymin=0, ymax=18, xtick={4}, ytick={5,15}, xlabel={$t$}, ylabel={$v$}, title={\sffamily\small とちゅうまで（台形）}]
  \fill[bkBlue!25] (axis cs:0,0) -- (axis cs:0,15) -- (axis cs:4,5) -- (axis cs:4,0) -- cycle;
  \draw[bkGray, line width=0.6pt, dashed] (axis cs:4,0) -- (axis cs:4,5);
  \addplot[bkcurve, domain=0:4] {15-2.5*x};
  \node[font=\sffamily\footnotesize] at (axis cs:2,3.6) {\qty{40}{m}};
\end{axis}
\end{tikzpicture}
\bkfigcaption{左：\qty{20}{m/s}から5秒で止まる。右：\qty{15}{m/s}から4秒で\qty{5}{m/s}になる。}
\end{bkfigure}

\begin{itemize}[leftmargin=1.6\zw, itemsep=1pt, topsep=2pt, label={\color{bkTheme}●}]
\item 左：底辺5、高さ20の三角形なので、$5\times20\div2=\qty{50}{m}$。\\平均の速さで考えると、$(20+0)\div2=10$、$10\times5=50$。\\5-5の例（\qty{20}{m/s}の車が加速度 $\qty{-4}{m/s^2}$ で止まるまで）と同じ結果です。
\item 右：上底15、下底5、高さ4の台形。$(15+5)\times4\div2=\qty{40}{m}$。⑤の速くなる自転車（5から15）と、ちょうど左右を入れかえた形なので、面積も同じ40です。
\end{itemize}

\Needspace{6\baselineskip}
\bksubsection{7}{速度がマイナスの部分は、マイナスの面積}

\label{bk:5-9:minus}ろう下で、右向きを＋とします。\qty{4}{m/s}で3秒進み、そのあと\qty{-2}{m/s}で2秒もどりました。\qty{-2}{m/s}は「－の向き（左）に、1秒に\qty{2}{m}」という意味です（5-5）。

\noindent\begin{minipage}[c]{0.5\linewidth}\centering
\begin{tikzpicture}
% 第7回：目もりの数字 3・5 は、赤い四角の中ではなく下に出し、点線でつなぐ
\begin{axis}[bkplot, width=62mm, height=46mm, xmin=0, xmax=6, ymin=-3.6, ymax=5, xtick={3,5}, xticklabels={}, ytick={-2,4}, xlabel={$t$（s）}, ylabel={$v$（m/s）}]
  \fill[bkBlue!20] (axis cs:0,0) rectangle (axis cs:3,4);
  \fill[bkRed!25] (axis cs:3,0) rectangle (axis cs:5,-2);
  \addplot[bkcurve, line width=1.4pt] coordinates {(0,4) (3,4)};
  \addplot[bkcurve, line width=1.4pt] coordinates {(3,-2) (5,-2)};
  \draw[bkGray, line width=0.5pt, dash pattern=on 1.5pt off 1pt] (axis cs:3,-2) -- (axis cs:3,-2.55) (axis cs:5,-2) -- (axis cs:5,-2.55);
  \node[font=\footnotesize, text=bkInk, anchor=north] at (axis cs:3,-2.6) {$3$};
  \node[font=\footnotesize, text=bkInk, anchor=north] at (axis cs:5,-2.6) {$5$};
  \node[font=\sffamily\small] at (axis cs:1.5,2) {$+12$};
  \node[font=\sffamily\small] at (axis cs:4,-1) {$-4$};
\end{axis}
\end{tikzpicture}
\end{minipage}\hfill
\begin{minipage}[c]{0.46\linewidth}\raggedright
横軸より上の長方形は、たて4、横3なので、面積は $4\times3=12$。右へ\qty{12}{m}進んだ、という意味です。

横軸より下の長方形は、たて $-2$、横2なので、面積を $(-2)\times2=-4$ と考えます。左へ\qty{4}{m}もどった、という意味です。
\end{minipage}

\begin{bkfigure}
\begin{tikzpicture}[x=1mm, y=1mm]
  \draw[bk axis] (-4,0) -- (104,0);
  \node[font=\sffamily\small, anchor=west] at (105,0) {m};
  \foreach \p in {0,...,13} \draw[bkInk, line width=0.4pt] ({7*\p},-0.9) -- ({7*\p},0.9);
  \foreach \p in {0,4,8,12} \node[font=\sffamily\small, anchor=north] at ({7*\p},-2) {\p};
  \fill[bkInk] (0,0) circle (1.3);
  \fill[bkRed] (56,0) circle (1.3);
  \draw[bk arrow, bkBlue] (0,5) -- (84,5);
  \node[font=\sffamily\footnotesize, anchor=south] at (36,6.5) {3秒で $+\qty{12}{m}$};
  \draw[bk arrow, bkRed] (84,13) -- (56,13);
  \node[font=\sffamily\footnotesize, anchor=south] at (70,14.5) {2秒で $-\qty{4}{m}$};
  \node[font=\sffamily\footnotesize, anchor=north] at (0,-7) {はじめ};
  \node[font=\sffamily\footnotesize, anchor=north, text=bkRed!80!black] at (56,-7) {最後の位置};
\end{tikzpicture}
\bkfigcaption{はじめの位置0から、12まで進み、8までもどった。}
\end{bkfigure}

最後の位置は、はじめの位置から $12+(-4)=\qty{8}{m}$ 右です。このように、はじめの位置から最後の位置までの変化を\textbf{変位}（へんい）といいます。変位は、向きをふくめた量で、上の面積の合計（マイナスの面積もそのままたす）が変位です。\bkidx{へんい}{変位}

一方、実際に動いた長さは、もどった分もたして $12+4=\qty{16}{m}$ です。これを\textbf{道のり}といいます。道のりは、マイナスの面積も、プラスの大きさにしてたします。\bkidx{みちのり}{道のり}

\noindent\textbf{線が横軸を横切るとき}\quad ボールを\qty{20}{m/s}で真上に投げ上げます（上向きを＋）。速度は $v=20-10t$（5-5）。2秒までは＋（上がる）、2秒で0（いちばん高い）、2秒からは－（落ちる）です。

\noindent\begin{minipage}[c]{0.42\linewidth}\centering
\begin{tikzpicture}
\begin{axis}[bkplot, width=52mm, height=50mm, xmin=0, xmax=4.8, ymin=-24, ymax=26, xtick={1}, ytick={-20,20}, xlabel={$t$}, ylabel={$v$}]
  \fill[bkBlue!25] (axis cs:0,0) -- (axis cs:0,20) -- (axis cs:2,0) -- cycle;
  \fill[bkRed!25] (axis cs:2,0) -- (axis cs:4,-20) -- (axis cs:4,0) -- cycle;
  \draw[bkGray, line width=0.4pt] (axis cs:3,0) -- (axis cs:3,1.2) (axis cs:4,0) -- (axis cs:4,1.2);
  \addplot[bkcurve, domain=0:4.2] {20-10*x};
  \node[font=\sffamily\footnotesize] at (axis cs:0.6,5) {$+20$};
  \node[font=\sffamily\footnotesize] at (axis cs:3.55,-6.5) {$-20$};
  \node[font=\sffamily\footnotesize, anchor=south] at (axis cs:2,1.2) {2};
  \node[font=\sffamily\footnotesize, anchor=south] at (axis cs:3,1.2) {3};
  \node[font=\sffamily\footnotesize, anchor=south] at (axis cs:4,1.2) {4};
\end{axis}
\end{tikzpicture}
\end{minipage}\hfill
\begin{minipage}[c]{0.54\linewidth}\raggedright
\begin{itemize}[leftmargin=1.2\zw, itemsep=1pt, topsep=2pt, label={\color{bkTheme}●}]
\item 0〜2秒は、横軸より上の三角形。\\$2\times20\div2=20$ で、\qty{20}{m}上がる。
\item 2〜4秒は、横軸より下の三角形。\\$-(2\times20\div2)=-20$ で、\qty{20}{m}落ちる。
\item 4秒までの変位は $20+(-20)=0$。\\投げた所にもどってきた。
\item 道のりは $20+20=\qty{40}{m}$。
\end{itemize}
\end{minipage}

\Needspace{12\baselineskip}
\bksubsection{8}{面積をつみ上げて、位置のグラフをつくる}

\label{bk:5-9:tsumiage}面積は、とちゅうまでの分も計算できます。やってみようの自転車で、2秒ごとの面積を求め、0秒からの合計（出発点からの位置）を表にします。

\begin{bkcenter}
\begin{bktable}{colspec={Q[l,m,wd=30mm] *{6}{Q[c,m,wd=12mm]}}, rowsep=1.5pt, colsep=4pt, row{3}={font=\small}}
時刻 $t$（s） & 0 & 2 & 4 & 6 & 8 & 10 \\
速さ $v$（m/s） & 0 & 2 & 4 & 4 & 4 & 2 \\
2秒ごとの形 & & 三角形 & 台形 & 長方形 & 長方形 & 台形 \\
2秒ごとの面積（m） & & 2 & 6 & 8 & 8 & 6 \\
0秒からの合計＝位置 $x$（m） & 0 & 2 & 8 & 16 & 24 & 30 \\
\end{bktable}
\bktabcaption{2秒ごとの面積をつみ上げると、位置になる。}
\end{bkcenter}

\noindent 2秒ごとの面積の求め方：0〜2秒は三角形 $2\times2\div2=2$。2〜4秒は台形 $(2+4)\times2\div2=6$。4〜6秒と6〜8秒は長方形 $4\times2=8$。8〜10秒は台形 $(4+2)\times2\div2=6$。合計の行は、前の合計に次の面積をたしていきます（$0+2=2$、$2+6=8$、$8+8=16$……）。

\noindent\begin{minipage}[c]{0.52\linewidth}\centering
\begin{tikzpicture}
\begin{axis}[bkplot, width=64mm, height=46mm, xmin=0, xmax=11, ymin=0, ymax=34, xtick={2,4,6,8,10}, ytick={10,20,30}, xlabel={$t$（s）}, ylabel={$x$（m）}]
  \addplot[bkcurve, domain=0:4] {x^2/2};
  \addplot[bkcurve, domain=4:8] {8+4*(x-4)};
  \addplot[bkcurve, domain=8:10] {24+4*(x-8)-0.5*(x-8)^2};
  \addplot[only marks, mark=*, mark size=1.5pt, bkInk] coordinates {(0,0) (2,2) (4,8) (6,16) (8,24) (10,30)};
\end{axis}
\end{tikzpicture}
\end{minipage}\hfill
\begin{minipage}[c]{0.44\linewidth}\raggedright
表の点を $x$-$t$ グラフにすると、左のようになります。速さがふえている0〜4秒はだんだん急に、速さ一定の4〜8秒はまっすぐ、おそくなる8〜10秒はだんだんゆるやかになります。

4〜8秒の傾きは\\$(24-8)\div(8-4)=4$ で、このときの速さ\qty{4}{m/s}と同じです。
\end{minipage}

\medskip
$x$-$t$ グラフの傾きは、もとの速さになっています（5-3・5-5で学んだ「位置を微分すると速度」）。面積をつみ上げる計算と、傾きを求める計算は、ちょうど反対向きの関係になっています。

\begin{bkcenter}
\begin{bktable}{colspec={Q[l,m,wd=22mm] Q[l,m,wd=50mm] Q[l,m,wd=50mm]}, rowsep=2pt, cells={font=\small}, row{1}={font=\sffamily\bfseries\small}}
 & 微分（5-1〜5-8） & 面積をつみ上げる（5-9〜） \\
向き & 位置 $x$ → 速度 $v$ & 速度 $v$ → 位置の変化 \\
短い区間で & わり算 $\Delta x\div\Delta t$ & かけ算 $v\times\Delta t$ \\
そのあと & $\Delta t$ を0に近づける & たして、$\Delta t$ を0に近づける \\
グラフでは & 傾き & 面積 \\
\end{bktable}
\bktabcaption{微分と、面積をつみ上げる計算のくらべ。5-11で「積分は微分の逆」として整理する。}
\end{bkcenter}

\bksubsection{9}{はてな？の答え：タコグラフ}

\label{bk:5-9:taco}\bkhatenaref のグラフは、時間の単位が h（時間）、速さの単位が km/h です。km/h×h＝km なので、面積がそのまま km になります。単位の換算はいりません。グラフを、まっすぐな線ごとの区間に分けます。

\begin{bkcenter}
\begin{bktable}{colspec={Q[c,m,wd=20mm] Q[c,m,wd=16mm] Q[c,m,wd=48mm] Q[c,m,wd=16mm]}, rowsep=1.5pt}
区間（h） & 形 & 面積の計算 & 距離（km） \\
0〜0.5 & 三角形 & $0.5\times40\div2$ & 10 \\
0.5〜2 & 長方形 & $40\times1.5$ & 60 \\
2〜2.5 & 三角形 & $0.5\times40\div2$ & 10 \\
2.5〜3 & （休けい） & 速さ0 & 0 \\
3〜3.5 & 三角形 & $0.5\times60\div2$ & 15 \\
3.5〜5.5 & 長方形 & $60\times2$ & 120 \\
5.5〜6 & 三角形 & $0.5\times60\div2$ & 15 \\
\end{bktable}
\bktabcaption{全部で $10+60+10+0+15+120+15=230$。}
\end{bkcenter}

\begin{bkfigure}
\begin{tikzpicture}
\begin{axis}[bkplot, width=110mm, height=44mm, xmin=0, xmax=6.6, ymin=0, ymax=72, xtick={0.5,1,...,6}, xticklabels={0.5,1,1.5,2,2.5,3,3.5,4,4.5,5,5.5,6}, ytick={20,40,60}, xlabel={$t$（h）}, ylabel={$v$（km/h）}, x tick label style={font=\sffamily\scriptsize}]
  \fill[bkYellow!45] (axis cs:0,0) -- (axis cs:0.5,40) -- (axis cs:0.5,0) -- cycle;
  \fill[bkBlue!22] (axis cs:0.5,0) rectangle (axis cs:2,40);
  \fill[bkYellow!45] (axis cs:2,0) -- (axis cs:2,40) -- (axis cs:2.5,0) -- cycle;
  \fill[bkYellow!45] (axis cs:3,0) -- (axis cs:3.5,60) -- (axis cs:3.5,0) -- cycle;
  \fill[bkBlue!22] (axis cs:3.5,0) rectangle (axis cs:5.5,60);
  \fill[bkYellow!45] (axis cs:5.5,0) -- (axis cs:5.5,60) -- (axis cs:6,0) -- cycle;
  \addplot[bkcurve, line width=1.3pt] coordinates {(0,0) (0.5,40) (2,40) (2.5,0) (3,0) (3.5,60) (5.5,60) (6,0)};
  \node[font=\sffamily\footnotesize] at (axis cs:1.25,20) {60};
  \node[font=\sffamily\footnotesize] at (axis cs:4.5,30) {120};
  \node[font=\sffamily\scriptsize] at (axis cs:0.36,8) {10};
  \node[font=\sffamily\scriptsize] at (axis cs:2.14,8) {10};
  \node[font=\sffamily\scriptsize] at (axis cs:3.36,10) {15};
  \node[font=\sffamily\scriptsize] at (axis cs:5.64,10) {15};
\end{axis}
\end{tikzpicture}
\bkfigcaption{タコグラフの面積を、三角形（黄）と長方形（青）に分けた。数は面積（km）。}
\end{bkfigure}

\noindent このトラックは、6時間で\qty{230}{km}走りました。速さだけの記録から、走った距離がわかりました。ほんとうのタコグラフの線はもっとでこぼこですが、細かく分けて（5-10）たせば、同じように距離が求められます。

\begin{motto}[2]
\textbf{グラフが曲がっていても、マスを数えればおよそわかる}

\begin{bkfigure}
\begin{tikzpicture}[x=6.5mm, y=4.4mm]
  \foreach \k/\lo/\hi in {0/0/3,1/2/5,2/4/7,3/6/8,4/7/8,5/7/8,6/6/8,7/4/7,8/2/5,9/0/3} {
    \fill[bkYellow!55] (\k,0) rectangle (\k+1,\lo);
    \fill[bkOrange!35] (\k,\lo) rectangle (\k+1,\hi);}
  \draw[bkLine, line width=0.4pt] (0,0) grid (10,8);
  \draw[bkBlue!80!black, line width=1.3pt, domain=0:10, samples=60] plot (\x,{3*\x-0.3*\x*\x});
  \draw[bkInk, line width=0.6pt, -{Stealth[length=3pt]}] (0,0) -- (10.7,0);
  \draw[bkInk, line width=0.6pt, -{Stealth[length=3pt]}] (0,0) -- (0,8.7);
  \foreach \k in {2,4,6,8,10} \node[font=\sffamily\scriptsize, anchor=north] at (\k,-0.15) {\k};
  \foreach \k in {2,4,6,8} \node[font=\sffamily\scriptsize, anchor=east] at (-0.15,\k) {\k};
  \node[font=\sffamily\scriptsize, anchor=west] at (10.7,0) {$t$（s）};
  \node[font=\sffamily\scriptsize, anchor=south] at (0,8.7) {$v$（m/s）};
\end{tikzpicture}
\bkfigcaption{黄：まるごと入っているマス（38こ）。オレンジ：線がとおり、一部だけ入っているマス（24こ）。}
\end{bkfigure}

ほんとうの乗り物の $v$-$t$ グラフは、上の図のように曲がっています。それでも、マスを数えれば、およその距離がわかります。一部だけ入っているマスを、ならして半分と見ると、$38+24\div2=50$ で、約\qty{50}{m}です。

\medskip
この図の速さは $v=3t-0.3t^2$ で、5-10・5-11の方法で計算すると、ちょうど\qty{50}{m}になります。マスを細かくするほど、数えた値はほんとうの値に近づきます。「細かく分けて、たす」は、曲がったグラフでもそのまま使えるのです。
\end{motto}

\begin{column}{ガリレオより前の発見}
一定の割合で速くなるときに進む距離は、「はじめの速さとおわりの速さの平均で、ずっと走ったときの距離」に等しくなります。④・⑤の「平均の速さで考える」と同じことです。このきまりは、700年ほど前に、イギリスのオックスフォードの学者たちが見つけ、フランスのオレームが図（いまの $v$-$t$ グラフのようなもの）を使って説明しました。ガリレオは、これを実験で確かめ、落ちる物の運動にあてはめました。
\end{column}

\begin{taisetsu}
$v$-$t$ グラフと横軸のあいだの面積 ＝ そのあいだの位置の変化（変位）\\
速さが変わるときは、細かく分けて、かけて（速さ×$\Delta t$）、たす。\\
横軸より下の部分は、マイナスの面積として数える。
\end{taisetsu}

\begin{motto}[2]
\textbf{$a$-$t$ グラフの面積は、速度の変化}

加速度は「1秒あたりの速度の変化」でした（5-5）。だから、加速度のグラフ（$a$-$t$ グラフ）の面積は、そのあいだの速度の変化になります。単位も、$\unit{m/s^2}\times\unit{s}=\unit{m/s}$ と、速度の単位になります。加速度が\qty{2}{m/s^2}で3秒なら、$2\times3=6$ で、速度は\qty{6}{m/s}ふえます。スマートフォンの加速度センサーの値から速度を計算するのは、この考え方です。
\end{motto}

\begin{kotoba}
\small
\bkword{$v$-$t$ グラフ}{横軸に時間、たて軸に速度をとったグラフ。面積が位置の変化になる。}
\bkword{変位}{はじめの位置から、最後の位置までの変化。向きをふくむ（マイナスにもなる）。}
\bkword{道のり}{実際に動いた長さの合計。向きは考えない（いつもプラス）。}
\end{kotoba}

%% ---------------- 例題 ----------------
\bksection{reidai}{例題でたしかめよう}

\begin{reidai}[長方形に分ける]
電車が、\qty{20}{m/s}で30秒走ったあと、\qty{10}{m/s}で20秒走りました。進んだ距離を求めましょう。\par
\kangaekata 速さが一定の区間ごとに、①の長方形の面積（速さ×時間）を求めて、たし合わせます。\par
\begin{bkfigure}
\begin{tikzpicture}
\begin{axis}[bkplot, width=66mm, height=36mm, xmin=0, xmax=58, ymin=0, ymax=25, xtick={30,50}, ytick={10,20}, xlabel={$t$}, ylabel={$v$}]
  \fill[bkBlue!25] (axis cs:0,0) rectangle (axis cs:30,20);
  \fill[bkOrange!35] (axis cs:30,0) rectangle (axis cs:50,10);
  \addplot[bkcurve, line width=1.3pt] coordinates {(0,20) (30,20)};
  \addplot[bkcurve, line width=1.3pt] coordinates {(30,10) (50,10)};
  \node[font=\sffamily\footnotesize] at (axis cs:15,10) {600};
  \node[font=\sffamily\footnotesize] at (axis cs:40,5) {200};
\end{axis}
\end{tikzpicture}
\end{bkfigure}
\shiki はじめの30秒：$20\times30=600$\par
\hspace*{2.6\zw}次の20秒：$10\times20=200$\par
\hspace*{2.6\zw}合計：$600+200=800$\par
\kotae \qty{800}{m}
\end{reidai}

\begin{jibun}
自転車が、\qty{5}{m/s}で40秒走ったあと、\qty{3}{m/s}で60秒走りました。進んだ距離を求めましょう。
\bkjans{$5\times40+3\times60=200+180=\qty{380}{m}$}
\end{jibun}

\begin{reidai}[単位をそろえてから]
車が、時速\qty{90}{km}で2分走りました。進んだ距離は何mですか。\par
\kangaekata 速さは「1時間あたり」、時間は「分」で、単位がそろっていません。速さを m/s に、時間を s にそろえてから、かけます（2-10）。1 kmは1000 mなので \qty{1000}{m/km}、1時間は $60\times60=3600$ 秒なので \qty{3600}{s/h} を使います。\par
\shiki 速さ：$\qty{90}{km/h}\times\qty{1000}{m/km}=\qty{90000}{m/h}$\par
\hspace*{2.6\zw}$\qty{90000}{m/h}\div\qty{3600}{s/h}=\qty{25}{m/s}$\par
\hspace*{2.6\zw}時間：$\qty{2}{min}\times\qty{60}{s/min}=\qty{120}{s}$\par
\hspace*{2.6\zw}距離：$\qty{25}{m/s}\times\qty{120}{s}=\qty{3000}{m}$\par
\hspace*{2.6\zw}単位の約分：$90\,\dfrac{\cancel{\mathrm{km}}}{\mathrm{h}}\times1000\,\dfrac{\mathrm{m}}{\cancel{\mathrm{km}}}\div3600\,\dfrac{\mathrm{s}}{\mathrm{h}}=25\,\dfrac{\mathrm{m}}{\mathrm{s}}$\par
\hspace*{2.6\zw}確かめ（時間を h にそろえる）：$\qty{2}{min}\div\qty{60}{min/h}=\dfrac{1}{30}$\,h\par
\hspace*{2.6\zw}$\qty{90}{km/h}\times\dfrac{1}{30}\,\mathrm{h}=\qty{3}{km}$（＝\qty{3000}{m}）\par
\kotae \qty{3000}{m}（\qty{3}{km}）
\end{reidai}

\begin{jibun}
自転車が、時速\qty{18}{km}で3分走りました。進んだ距離は何mですか。
\bkjans{$\qty{18}{km/h}\times\qty{1000}{m/km}\div\qty{3600}{s/h}=\qty{5}{m/s}$、$\qty{3}{min}\times\qty{60}{s/min}=\qty{180}{s}$。$5\times180=\qty{900}{m}$}
\end{jibun}

\begin{reidai}[階段の和ではさむ]
止まっていた車が走り出し、速さが $v=3t$（m/s）でふえていきます。2秒間を0.5秒ずつの区間に分け、（1）はじめの速さを使った和、（2）おわりの速さを使った和を求め、（3）三角形の面積（ほんとうの距離）とくらべましょう。\par
\kangaekata ②・③と同じく、分けて（0.5秒ずつ4区間）、かけて（速さ×0.5）、たします。区間のはしの時刻は0、0.5、1、1.5、2秒で、そのときの速さは $3t$ に入れて0、1.5、3、4.5、6です。\par
\shiki（1）はじめの速さ：$(0+1.5+3+4.5)\times0.5=9\times0.5=4.5$\par
\hspace*{2.6\zw}（2）おわりの速さ：$(1.5+3+4.5+6)\times0.5=15\times0.5=7.5$\par
\hspace*{2.6\zw}（3）2秒後の速さは6。三角形：$2\times6\div2=6$\par
\begin{bkfigure}
\begin{tikzpicture}
\begin{axis}[bkplot, name=s1, width=48mm, height=38mm, xmin=0, xmax=2.4, ymin=0, ymax=7.2, xtick={0.5,1,1.5,2}, ytick={3,6}, xlabel={$t$}, ylabel={$v$}, x tick label style={font=\sffamily\scriptsize}, title={\sffamily\footnotesize（1）はじめの速さ}]
  \foreach \k in {1,2,3} {\edef\tmp{\noexpand\fill[bkBlue!25, draw=bkBlue!70, line width=0.4pt] (axis cs:{0.5*\k},0) rectangle (axis cs:{0.5*\k+0.5},{1.5*\k});}\tmp}
  \addplot[bkcurve, line width=1.2pt, domain=0:2.1] {3*x};
\end{axis}
\begin{axis}[bkplot, at={(s1.east)}, anchor=west, xshift=12mm, width=48mm, height=38mm, xmin=0, xmax=2.4, ymin=0, ymax=7.2, xtick={0.5,1,1.5,2}, ytick={3,6}, xlabel={$t$}, ylabel={$v$}, x tick label style={font=\sffamily\scriptsize}, title={\sffamily\footnotesize（2）おわりの速さ}]
  \foreach \k in {0,1,2,3} {\edef\tmp{\noexpand\fill[bkBlue!25, draw=bkBlue!70, line width=0.4pt] (axis cs:{0.5*\k},0) rectangle (axis cs:{0.5*\k+0.5},{1.5*\k+1.5});}\tmp}
  \addplot[bkcurve, line width=1.2pt, domain=0:2.1] {3*x};
\end{axis}
\end{tikzpicture}
\end{bkfigure}
\kotae（1）\qty{4.5}{m}\quad（2）\qty{7.5}{m}\quad（3）\qty{6}{m}。$4.5<6<7.5$ で、ほんとうの距離は2つの和のあいだにある
\end{reidai}

\begin{jibun}
止まっていた車の速さが $v=4t$（m/s）でふえていきます。4秒間を1秒ずつの区間に分け、（1）はじめの速さを使った和、（2）おわりの速さを使った和を求め、（3）三角形の面積とくらべましょう。
\bkjans{（1）\nobreak $(0+4+8+12)\times1=\qty{24}{m}$\quad（2）\nobreak $(4+8+12+16)\times1=\qty{40}{m}$\quad（3）\nobreak $4\times16\div2=\qty{32}{m}$。$24<32<40$}
\end{jibun}

\begin{reidai}[止まっていた物が速くなる（三角形）]
止まっていた車が、4秒で\qty{12}{m/s}まで一定の割合で速くなりました。この4秒で進んだ距離を求めましょう。\par
\kangaekata $v$-$t$ グラフは、原点（0秒で速さ0）から（4，12）までの直線。下は、底辺4（時間）、高さ12（4秒後の速さ）の三角形です（④）。\par
\shiki 底辺4、高さ12の三角形：$4\times12\div2=24$\par
\hspace*{2.6\zw}確かめ（平均の速さ）：$(0+12)\div2=6$、$6\times4=24$\par
\begin{bkfigure}
\begin{tikzpicture}
\begin{axis}[bkplot, width=50mm, height=36mm, xmin=0, xmax=5, ymin=0, ymax=14.5, xtick={4}, ytick={12}, xlabel={$t$}, ylabel={$v$}]
  \fill[bkBlue!25] (axis cs:0,0) -- (axis cs:4,12) -- (axis cs:4,0) -- cycle;
  \draw[bkGray, dashed, line width=0.5pt] (axis cs:0,12) -- (axis cs:4,12);
  \addplot[bkcurve, line width=1.3pt, domain=0:4] {3*x};
  \node[font=\sffamily\footnotesize] at (axis cs:2.9,3) {24};
\end{axis}
\end{tikzpicture}
\end{bkfigure}
\kotae \qty{24}{m}
\end{reidai}

\begin{jibun}
止まっていた電車が、20秒で\qty{15}{m/s}まで一定の割合で速くなりました。この20秒で進んだ距離を求めましょう。
\bkjans{$20\times15\div2=\qty{150}{m}$}
\end{jibun}

\begin{reidai}[加速度がわかっているとき]
止まっていた物体が、加速度\qty{2}{m/s^2}で5秒間速くなりました。進んだ距離を求めましょう。\par
\noindent\begin{minipage}[c]{0.6\linewidth}\raggedright
\kangaekata まず、5秒後の速さ（三角形の高さ）を求めます。加速度\qty{2}{m/s^2}は「1秒ごとに\qty{2}{m/s}ずつ速くなる」なので、5秒で $2\times5$ ふえます。\par
\shiki 5秒後の速さ：$2\times5=\qty{10}{m/s}$\par
\hspace*{2.6\zw}三角形の面積：$5\times10\div2=25$\par
\hspace*{2.6\zw}確かめ（$\frac{1}{2}at^2$）：$\dfrac{1}{2}\times2\times5^2=25$\par
\end{minipage}\hfill
\begin{minipage}[c]{0.37\linewidth}\centering
\begin{tikzpicture}
\begin{axis}[bkplot, width=41mm, height=38mm, xmin=0, xmax=6.2, ymin=0, ymax=12.5, xtick={1,...,5}, ytick={2,4,6,8,10}, xlabel={$t$}, ylabel={$v$}, y tick label style={font=\sffamily\scriptsize}]
  \fill[bkBlue!25] (axis cs:0,0) -- (axis cs:5,10) -- (axis cs:5,0) -- cycle;
  \addplot[bkcurve, line width=1.3pt, domain=0:5] {2*x};
  \node[font=\sffamily\footnotesize] at (axis cs:3.7,2.4) {25};
\end{axis}
\end{tikzpicture}
\end{minipage}\par
\kotae \qty{25}{m}
\end{reidai}

\begin{jibun}
止まっていた物体が、加速度\qty{1.5}{m/s^2}で4秒間速くなりました。進んだ距離を求めましょう。
\bkjans{4秒後の速さ $1.5\times4=\qty{6}{m/s}$。$4\times6\div2=\qty{12}{m}$}
\end{jibun}

\begin{reidai}[はじめから動いている（台形）]
\qty{5}{m/s}で走っていた自転車が、一定の割合で速くなり、4秒後に\qty{15}{m/s}になりました。この4秒で進んだ距離を求めましょう。\par
\noindent\begin{minipage}[c]{0.6\linewidth}\raggedright
\kangaekata ⑤の台形です。たての辺（上底・下底）が5と15、横のはば（高さ）が4です。\par
\shiki $(5+15)\times4\div2=20\times4\div2=40$\par
\hspace*{2.6\zw}確かめ（長方形＋三角形）：$5\times4+4\times(15-5)\div2=20+20=40$\par
\end{minipage}\hfill
\begin{minipage}[c]{0.37\linewidth}\centering
\begin{tikzpicture}
\begin{axis}[bkplot, width=41mm, height=38mm, xmin=0, xmax=5, ymin=0, ymax=17.5, xtick={4}, ytick={5,15}, xlabel={$t$}, ylabel={$v$}]
  \fill[bkBlue!25] (axis cs:0,0) -- (axis cs:0,5) -- (axis cs:4,15) -- (axis cs:4,0) -- cycle;
  \draw[bkGray, dashed, line width=0.5pt] (axis cs:4,0) -- (axis cs:4,15) (axis cs:0,15) -- (axis cs:4,15);
  \addplot[bkcurve, line width=1.3pt, domain=0:4] {5+2.5*x};
  \node[font=\sffamily\footnotesize] at (axis cs:2,3.5) {40};
\end{axis}
\end{tikzpicture}
\end{minipage}\par
\kotae \qty{40}{m}
\end{reidai}

\begin{jibun}
\qty{4}{m/s}で走っていた自転車が、一定の割合で速くなり、3秒後に\qty{10}{m/s}になりました。この3秒で進んだ距離を求めましょう。
\bkjans{$(4+10)\times3\div2=\qty{21}{m}$}
\end{jibun}

\begin{reidai}[ブレーキでおそくなる]
（1）\qty{15}{m/s}で走っていた車が、一定の割合でおそくなり、4秒後に\qty{5}{m/s}になりました。この4秒で進んだ距離は？\par
（2）時速\qty{72}{km}で走っていた車が、加速度 $\qty{-4}{m/s^2}$ でブレーキをかけました。止まるまでに進む距離は？\par
\kangaekata（1）は右下がりの台形（⑥）。（2）は、速さを m/s にしてから、止まるまでの時間（三角形の底辺）を求めます。加速度 $-4$ は「1秒ごとに\qty{4}{m/s}ずつおそくなる」です。\par
\begin{bkfigure}
\begin{tikzpicture}
\begin{axis}[bkplot, name=e1, width=46mm, height=36mm, xmin=0, xmax=5, ymin=0, ymax=17.5, xtick={4}, ytick={5,15}, xlabel={$t$}, ylabel={$v$}, title={\sffamily\footnotesize（1）}]
  \fill[bkBlue!25] (axis cs:0,0) -- (axis cs:0,15) -- (axis cs:4,5) -- (axis cs:4,0) -- cycle;
  \draw[bkGray, dashed, line width=0.5pt] (axis cs:4,0) -- (axis cs:4,5);
  \addplot[bkcurve, line width=1.3pt, domain=0:4] {15-2.5*x};
  \node[font=\sffamily\footnotesize] at (axis cs:2,3.6) {40};
\end{axis}
\begin{axis}[bkplot, at={(e1.east)}, anchor=west, xshift=12mm, width=46mm, height=36mm, xmin=0, xmax=6.2, ymin=0, ymax=23, xtick={5}, ytick={20}, xlabel={$t$}, ylabel={$v$}, title={\sffamily\footnotesize（2）}]
  \fill[bkBlue!25] (axis cs:0,0) -- (axis cs:0,20) -- (axis cs:5,0) -- cycle;
  \addplot[bkcurve, line width=1.3pt, domain=0:5] {20-4*x};
  \node[font=\sffamily\footnotesize] at (axis cs:1.3,5) {50};
\end{axis}
\end{tikzpicture}
\end{bkfigure}
\shiki（1）$(15+5)\times4\div2=40$\par
\hspace*{2.6\zw}（2）$\qty{72}{km/h}\times\qty{1000}{m/km}\div\qty{3600}{s/h}=\qty{20}{m/s}$\par
\hspace*{4.6\zw}止まるまでの時間：$20\div4=\qty{5}{s}$\par
\hspace*{4.6\zw}三角形の面積：$5\times20\div2=50$\par
\kotae（1）\qty{40}{m}\quad（2）\qty{50}{m}
\end{reidai}

\begin{jibun}
（1）\qty{12}{m/s}で走っていた車が、一定の割合でおそくなり、2秒後に\qty{6}{m/s}になりました。この2秒で進んだ距離は？\par
（2）時速\qty{54}{km}で走っていた車が、加速度 $\qty{-3}{m/s^2}$ でブレーキをかけました。止まるまでに進む距離は？
\bkjans{（1）\nobreak $(12+6)\times2\div2=\qty{18}{m}$\quad（2）\nobreak \qty{15}{m/s}、$15\div3=\qty{5}{s}$、$5\times15\div2=\qty{37.5}{m}$}
\end{jibun}

\begin{reidai}[3つの区間：駅から駅まで]
電車が駅を出て、20秒で\qty{15}{m/s}まで一定の割合で速くなり、\qty{15}{m/s}で60秒走り、10秒で一定の割合でおそくなって次の駅に止まりました。駅と駅のあいだの距離を求めましょう。\par
\kangaekata $v$-$t$ グラフは、三角形・長方形・三角形をつなげた台形です。区間ごとに求めてたします。\par
\begin{bkfigure}
\begin{tikzpicture}
\begin{axis}[bkplot, width=84mm, height=36mm, xmin=0, xmax=98, ymin=0, ymax=18.5, xtick={20,80,90}, ytick={15}, xlabel={$t$}, ylabel={$v$}]
  \fill[bkYellow!45] (axis cs:0,0) -- (axis cs:20,15) -- (axis cs:20,0) -- cycle;
  \fill[bkBlue!22] (axis cs:20,0) rectangle (axis cs:80,15);
  \fill[bkYellow!45] (axis cs:80,0) -- (axis cs:80,15) -- (axis cs:90,0) -- cycle;
  \addplot[bkcurve, line width=1.3pt] coordinates {(0,0) (20,15) (80,15) (90,0)};
  \node[font=\sffamily\footnotesize] at (axis cs:14.5,4) {150};
  \node[font=\sffamily\footnotesize] at (axis cs:50,7.5) {900};
  \node[font=\sffamily\footnotesize] at (axis cs:83.4,4) {75};
\end{axis}
\end{tikzpicture}
\end{bkfigure}
\shiki 速くなる20秒：$20\times15\div2=150$\par
\hspace*{2.6\zw}一定の60秒：$15\times60=900$\par
\hspace*{2.6\zw}おそくなる10秒：$10\times15\div2=75$\par
\hspace*{2.6\zw}合計：$150+900+75=1125$\par
\hspace*{2.6\zw}確かめ（大きな台形：上底60、下底90、高さ15）：\par
\hspace*{2.6\zw}$(60+90)\times15\div2=1125$\par
\kotae \qty{1125}{m}
\end{reidai}

\begin{jibun}
電車が、10秒で\qty{10}{m/s}まで一定の割合で速くなり、\qty{10}{m/s}で30秒走り、10秒で一定の割合でおそくなって止まりました。進んだ距離を求めましょう。
\bkjans{$50+300+50=\qty{400}{m}$}
\end{jibun}

\begin{reidai}[マイナスの面積]
エレベーターが、\qty{2}{m/s}で6秒上がり、そのあと\qty{-3}{m/s}で2秒下がりました（上向きが＋）。変位と、動いた道のりを求めましょう。\par
\noindent\begin{minipage}[c]{0.6\linewidth}\raggedright
\kangaekata 下がるときの速度はマイナスなので、⑦のマイナスの面積です。変位は符号のままたし、道のりは大きさ（プラス）をたします。\par
\shiki 上がる：$2\times6=12$、下がる：$(-3)\times2=-6$\par
\hspace*{2.6\zw}変位：$12+(-6)=6$\quad 道のり：$12+6=18$\par
\end{minipage}\hfill
\begin{minipage}[c]{0.37\linewidth}\centering
\begin{tikzpicture}
\begin{axis}[bkplot, width=41mm, height=40mm, xmin=0, xmax=9.6, ymin=-4.6, ymax=3.4, xtick={6}, ytick={-3,2}, xlabel={$t$}, ylabel={$v$}]
  \fill[bkBlue!22] (axis cs:0,0) rectangle (axis cs:6,2);
  \fill[bkRed!25] (axis cs:6,0) rectangle (axis cs:8,-3);
  \addplot[bkcurve, line width=1.3pt] coordinates {(0,2) (6,2)};
  \addplot[bkcurve, line width=1.3pt] coordinates {(6,-3) (8,-3)};
  \draw[bkGray, line width=0.4pt] (axis cs:8,0) -- (axis cs:8,0.5);
  \node[font=\sffamily\footnotesize, anchor=south] at (axis cs:8,0.5) {8};
  \node[font=\sffamily\footnotesize] at (axis cs:3,1) {$+12$};
  \node[font=\sffamily\footnotesize] at (axis cs:7,-1.5) {$-6$};
\end{axis}
\end{tikzpicture}
\end{minipage}\par
\kotae 変位\qty{6}{m}（はじめより上）、道のり\qty{18}{m}
\end{reidai}

\begin{jibun}
ボールが\qty{3}{m/s}で2秒ころがり、かべではね返って\qty{-2}{m/s}で4秒ころがりました。変位と道のりを求めましょう。
\bkjans{変位 $6+(-8)=\qty{-2}{m}$（はじめの位置より\qty{2}{m}手前）、道のり $6+8=\qty{14}{m}$}
\end{jibun}

\begin{reidai}[線が横軸を横切る：投げ上げ]
ビルの屋上から、ボールを\qty{10}{m/s}で真上に投げ上げました（上向きが＋、加速度 $\qty{-10}{m/s^2}$）。速度は $v=10-10t$ です。\par
（1）いちばん高くなるまでに、何m上がりますか。（2）3秒後の変位はいくつですか。\par
\noindent\begin{minipage}[c]{0.6\linewidth}\raggedright
\kangaekata $v=0$ になる時刻で、上の三角形と下の三角形に分けます。$10-10t=0$ より $t=1$。3秒後の速度は $10-30=-20$。\par
\shiki（1）0〜1秒の三角形：$1\times10\div2=5$\par
\hspace*{2.6\zw}（2）1〜3秒の三角形：$-(2\times20\div2)=-20$\par
\hspace*{4.6\zw}変位：$5+(-20)=-15$\par
\end{minipage}\hfill
\begin{minipage}[c]{0.37\linewidth}\centering
\begin{tikzpicture}
\begin{axis}[bkplot, width=41mm, height=44mm, xmin=0, xmax=3.7, ymin=-23, ymax=13, xtick={2}, ytick={-20,10}, xlabel={$t$}, ylabel={$v$}]
  \fill[bkBlue!25] (axis cs:0,0) -- (axis cs:0,10) -- (axis cs:1,0) -- cycle;
  \fill[bkRed!25] (axis cs:1,0) -- (axis cs:3,-20) -- (axis cs:3,0) -- cycle;
  \draw[bkGray, line width=0.4pt] (axis cs:1,0) -- (axis cs:1,1) (axis cs:3,0) -- (axis cs:3,1);
  \addplot[bkcurve, line width=1.3pt, domain=0:3.2] {10-10*x};
  \node[font=\sffamily\footnotesize, anchor=south] at (axis cs:1,1) {1};
  \node[font=\sffamily\footnotesize, anchor=south] at (axis cs:3,1) {3};
  \node[font=\sffamily\scriptsize] at (axis cs:0.3,2.6) {$+5$};
  \node[font=\sffamily\footnotesize] at (axis cs:2.55,-6) {$-20$};
\end{axis}
\end{tikzpicture}
\end{minipage}\par
\kotae（1）\qty{5}{m}\quad（2）\qty{-15}{m}（投げた所より\qty{15}{m}下）
\end{reidai}

\begin{jibun}
ボールを\qty{15}{m/s}で真上に投げ上げました（上向きが＋、加速度 $\qty{-10}{m/s^2}$）。速度は $v=15-10t$ です。（1）いちばん高くなるまでに何m上がりますか。（2）3秒後の変位はいくつですか。
\bkjans{（1）\nobreak $t=1.5$、$1.5\times15\div2=\qty{11.25}{m}$\quad（2）\nobreak 3秒後の速度 $-15$。$11.25+(-11.25)=\qty{0}{m}$（投げた所にもどる）}
\end{jibun}

\begin{reidai}[面積をつみ上げる]
物体の速さを、1秒ごとに記録しました（それぞれのあいだは一定の割合で変わる）。0秒の位置を0として、1秒ごとの位置を求めましょう。\par
\begin{bkcenter}
\begin{bktable}{colspec={Q[l,m,wd=22mm] *{4}{Q[c,m,wd=10mm]}}, rowsep=1pt}
$t$（s） & 0 & 1 & 2 & 3 \\
$v$（m/s） & 2 & 4 & 4 & 0 \\
\end{bktable}
\end{bkcenter}
\noindent\begin{minipage}[c]{0.6\linewidth}\raggedright
\kangaekata ⑧のように、1秒ごとの面積（台形・長方形・三角形）を求め、前の位置につみ上げます。\par
\shiki 0〜1秒（台形）：$(2+4)\times1\div2=3$\par
\hspace*{2.6\zw}1〜2秒（長方形）：$4\times1=4$\par
\hspace*{2.6\zw}2〜3秒（三角形）：$1\times4\div2=2$\par
\hspace*{2.6\zw}位置：$0+3=3$、$3+4=7$、$7+2=9$\par
\end{minipage}\hfill
\begin{minipage}[c]{0.37\linewidth}\centering
\begin{tikzpicture}
\begin{axis}[bkplot, width=41mm, height=38mm, xmin=0, xmax=3.7, ymin=0, ymax=5.3, xtick={1,2,3}, ytick={2,4}, xlabel={$t$}, ylabel={$v$}]
  \fill[bkYellow!45] (axis cs:0,0) -- (axis cs:0,2) -- (axis cs:1,4) -- (axis cs:1,0) -- cycle;
  \fill[bkBlue!22] (axis cs:1,0) rectangle (axis cs:2,4);
  \fill[bkGreen!25] (axis cs:2,0) -- (axis cs:2,4) -- (axis cs:3,0) -- cycle;
  \addplot[bkcurve, line width=1.3pt] coordinates {(0,2) (1,4) (2,4) (3,0)};
  \node[font=\sffamily\footnotesize] at (axis cs:0.5,1.2) {3};
  \node[font=\sffamily\footnotesize] at (axis cs:1.5,2) {4};
  \node[font=\sffamily\footnotesize] at (axis cs:2.3,1) {2};
\end{axis}
\end{tikzpicture}
\end{minipage}\par
\kotae 1秒で\qty{3}{m}、2秒で\qty{7}{m}、3秒で\qty{9}{m}
\end{reidai}

\begin{jibun}
速さが、0秒で0、2秒で\qty{6}{m/s}、4秒で\qty{6}{m/s}、6秒で\qty{2}{m/s}（あいだは一定の割合で変わる）の物体の、2秒ごとの位置を求めましょう。
\bkjans{\mbox{2秒}で\qty{6}{m}、\mbox{4秒}で\qty{18}{m}、\mbox{6秒}で\qty{26}{m}}
\end{jibun}

%% ---------------- まちがえやすいところ・物理工学 ----------------
\bksection{tsumazuki}{まちがえやすいところをチェック}

\begin{tsumazuki}{単位をそろえずに面積を計算する}
\bkng{時速\qty{72}{km}で10秒：$72\times10=720$ で\qty{720}{m}}
\bkok{$\qty{72}{km/h}\times\qty{1000}{m/km}\div\qty{3600}{s/h}=\qty{20}{m/s}$。$20\times10=\qty{200}{m}$}
「1時間あたりの速さ」に「秒」をかけても、意味のある数になりません。かける前に、単位を m/s と s（または km/h と h）にそろえます。面積の単位が、距離の単位になっているかを確かめましょう。
\end{tsumazuki}

\begin{tsumazuki}{三角形の「÷2」をわすれる}
\bkng{止まっていた車が4秒で\qty{12}{m/s}：$12\times4=\qty{48}{m}$}
\bkok{$4\times12\div2=\qty{24}{m}$}
\noindent\begin{minipage}[c]{0.62\linewidth}\raggedright
$12\times4$ は「はじめから\qty{12}{m/s}で4秒走った」ときの距離で、右の図の長方形全体です。ほんとうは0からだんだん速くなるので、その半分の三角形（青）です。平均の速さ（$12\div2=6$）で確かめるとまちがえません。
\end{minipage}\hfill
\begin{minipage}[c]{0.34\linewidth}\centering
\begin{tikzpicture}
\begin{axis}[bkplot, width=40mm, height=32mm, xmin=0, xmax=5, ymin=0, ymax=14.5, xtick={4}, ytick={12}, xlabel={$t$}, ylabel={$v$}]
  \fill[bkGray!20] (axis cs:0,0) rectangle (axis cs:4,12);
  \fill[bkBlue!30] (axis cs:0,0) -- (axis cs:4,12) -- (axis cs:4,0) -- cycle;
  \addplot[bkcurve, line width=1.2pt, domain=0:4] {3*x};
\end{axis}
\end{tikzpicture}
\end{minipage}
\end{tsumazuki}

\begin{tsumazuki}{変位と道のりをまちがえる}
\bkng{\qty{4}{m/s}で3秒、\qty{-2}{m/s}で2秒。道のりは $12+(-4)=\qty{8}{m}$}
\bkok{道のりは $12+4=\qty{16}{m}$。変位が\qty{8}{m}}
もどった分も、実際に動いた長さです。道のりはマイナスの面積もプラスにしてたし、変位は符号のままたします。
\end{tsumazuki}

\begin{tsumazuki}{$x$-$t$ グラフの面積を求めてしまう}
距離になるのは、$v$-$t$ グラフ（速度のグラフ）の面積です。$x$-$t$ グラフ（位置のグラフ）では、ある時刻の位置は、グラフの高さそのものです。たて軸が何かを、いつも確かめましょう。
\end{tsumazuki}

\begin{column}{コンピュータは階段で計算する}
車の走行距離計や、スマートフォンの歩きの記録アプリは、0.1秒ごとなどの短い時間ごとに速さをはかり、「速さ×時間」をつぎつぎにたして、進んだ距離を出しています。この章の「細かく分けて、かけて、たす」そのものです。区間が短いので、階段とほんとうのグラフのすき間はとても小さく、正確な距離が求められます。5-8のオイラー法も、同じ計算でした。
\end{column}

\begin{butsuri}
\begin{itemize}[leftmargin=1.2\zw, itemsep=2pt, topsep=0pt, label={\color{bkPurple}●}]
\item \textbf{タコグラフ}\\
記録された速さのグラフの面積から、走った距離を計算します（はてな？）。運転時間や休けい時間も、グラフからわかります。
\item \textbf{電車の運転}\\
電車は、加速・一定の速さ・ブレーキを組み合わせて走ります。$v$-$t$ グラフの面積が駅と駅のあいだの距離になるように、運転のしかたを決めています（下の図）。
\item \textbf{ブレーキの距離}\\
同じ加速度でブレーキをかけると、速さが2倍なら止まるまでの時間も2倍です。三角形の底辺も高さも2倍になるので、面積（止まるまでの距離）は $2\times2=4$ 倍になります（下の右の図）。
\end{itemize}
\begin{bkfigure}
\begin{tikzpicture}
\begin{axis}[bkplot, name=r1, width=70mm, height=40mm, xmin=0, xmax=128, ymin=0, ymax=25, xtick={30,90,110}, ytick={20}, xlabel={$t$}, ylabel={$v$}]
  \fill[bkPurple!15] (axis cs:0,0) -- (axis cs:30,20) -- (axis cs:90,20) -- (axis cs:110,0) -- cycle;
  \addplot[bkcurve, line width=1.5pt] coordinates {(0,0) (30,20) (90,20) (110,0)};
  \node[font=\sffamily\scriptsize] at (axis cs:12,22) {加速};
  \node[font=\sffamily\scriptsize] at (axis cs:60,22.5) {一定の速さ};
  \node[font=\sffamily\scriptsize] at (axis cs:106,22) {ブレーキ};
  \node[font=\sffamily\scriptsize] at (axis cs:60,9) {面積＝駅と駅の距離};
\end{axis}
\begin{axis}[bkplot, at={(r1.east)}, anchor=west, xshift=12mm, width=40mm, height=40mm, xmin=0, xmax=4.8, ymin=0, ymax=24, xtick={2,4}, ytick={10,20}, xlabel={$t$}, ylabel={$v$}]
  \fill[bkPurple!15] (axis cs:0,0) -- (axis cs:0,20) -- (axis cs:4,0) -- cycle;
  \fill[bkOrange!45] (axis cs:0,0) -- (axis cs:0,10) -- (axis cs:2,0) -- cycle;
  \addplot[bkcurve, line width=1.3pt, domain=0:4] {20-5*x};
  \addplot[bkcurve2, line width=1.3pt, domain=0:2] {10-5*x};
\end{axis}
\end{tikzpicture}
\bkfigcaption{左：電車の運転曲線（ランカーブ）。右：速さ2倍（むらさき）は、距離4倍。}
\end{bkfigure}
\end{butsuri}

%% ---------------- 章末問題 ----------------
\bkendproblems
\bkproblemgroup{基本}{2}

\bkproblem 時速\qty{54}{km}で2分走りました。進んだ距離は何mですか。
\bkpans{$\qty{54}{km/h}\times\qty{1000}{m/km}\div\qty{3600}{s/h}=\qty{15}{m/s}$、$\qty{2}{min}\times\qty{60}{s/min}=\qty{120}{s}$。$15\times120=\qty{1800}{m}$}

\bkproblem 止まっていた物体が、一定の加速度\qty{3}{m/s^2}で6秒動きました。$v$-$t$ グラフの面積から、進んだ距離を求めましょう。
\bkpans{\mbox{6秒後}の速さ $3\times6=\qty{18}{m/s}$。$6\times18\div2=\qty{54}{m}$}

\bkproblem \qty{10}{m/s}の自転車が、一定の割合でおそくなり、5秒後に\qty{4}{m/s}になりました。この5秒で進んだ距離を求めましょう。
\bkpans{$(10+4)\times5\div2=\qty{35}{m}$}

\bkproblem 電車が、30秒で\qty{20}{m/s}まで一定の割合で速くなり、\qty{20}{m/s}で60秒走り、20秒で一定の割合でおそくなって止まりました。駅と駅のあいだの距離を求めましょう。
\bkpans{$300+1200+200=\qty{1700}{m}$}

\bkproblem \qty{4}{m/s}で5秒進み、そのあと\qty{-3}{m/s}で4秒もどりました。変位と道のりを求めましょう。
\bkpans{変位 $20+(-12)=\qty{8}{m}$、道のり $20+12=\qty{32}{m}$}

\bkproblem 速さが、0秒で0、1秒で\qty{2}{m/s}、2秒で\qty{2}{m/s}、3秒で\qty{4}{m/s}、4秒で0（あいだは一定の割合で変わる）の物体の、1秒ごとの位置を求めましょう（0秒の位置は0）。
\bkpans{1秒で\qty{1}{m}、2秒で\qty{3}{m}、3秒で\qty{6}{m}、4秒で\qty{8}{m}}

\bkproblemgroup{チャレンジ}{2}

\bkproblem ボールを\qty{20}{m/s}で真上に投げ上げます（上向きが＋、加速度 \qty{-10}{m/s^2}）。\par
（1）$v$-$t$ グラフをかきましょう。\quad（2）いちばん高くなるまでの距離を、面積で求めましょう。\quad（3）4秒後の変位はいくつですか。
\bkpans{（1）\nobreak $v=20-10t$ の直線（2秒で0）\quad（2）\nobreak $2\times20\div2=\qty{20}{m}$\quad（3）\nobreak $20+(-20)=\qty{0}{m}$（もとの高さにもどる）}

\bkproblem 時速\qty{36}{km}と時速\qty{72}{km}の車が、どちらも加速度 \qty{-5}{m/s^2} でブレーキをかけて止まりました。止まるまでの距離をそれぞれ求め、くらべましょう。
\bkpans{\qty{10}{m/s}：\mbox{2秒}で $2\times10\div2=\qty{10}{m}$。\qty{20}{m/s}：\mbox{4秒}で $4\times20\div2=\qty{40}{m}$。速さ\mbox{2倍}で距離は\mbox{4倍}}

\bkproblem 止まっていた物体の速さが $v=6t$（m/s）でふえていきます。2秒間を0.5秒ずつに分け、はじめの速さを使った和と、おわりの速さを使った和を求めましょう。ほんとうの距離はいくつで、2つの和とどんな関係にありますか。
\bkpans{はじめ $(0+3+6+9)\times0.5=\qty{9}{m}$、おわり $(3+6+9+12)\times0.5=\qty{15}{m}$。ほんとうは $2\times12\div2=\qty{12}{m}$ で、2つの和のあいだ}

\bkproblem バスが、0.25時間で時速\qty{40}{km}まで一定の割合で速くなり、時速\qty{40}{km}で1.5時間走り、0.25時間で一定の割合でおそくなって止まりました。進んだ距離は何kmですか。
\bkpans{加速と減速は三角形で、どちらも $0.25\times40\div2=5$。一定の速さの所は $40\times1.5=60$。合計 $5+60+5=\qty{70}{km}$}

\begin{furikaeri}
\begin{checklist}
\item 速さが変わるときも、細かく分けて、かけて、たして距離を求められる。\hfill{\footnotesize\sffamily\color{bkTheme}→\pageref{bk:5-9:wakeru}ページ}
\item 長方形・三角形・台形の面積で、進んだ距離がわかる。\hfill{\footnotesize\sffamily\color{bkTheme}\mbox{→\pageref{bk:5-9:sankaku}ページ}}
\item マイナスの面積の意味（変位と道のり）を説明できる。\hfill{\footnotesize\sffamily\color{bkTheme}→\pageref{bk:5-9:minus}ページ}
\end{checklist}
\end{furikaeri}

\bkfukidashi{mina}{グラフがまっすぐなら、三角形や台形で計算できたね。じゃあ、グラフが曲がっていたら？　階段の考え方が、そのまま使えるのかな？}

%% 5-10 細い長方形に分けてたそう（第11回：ページ制限なし。X16・X18・X19 に対応して書き足した）
\bkchapter{5}{10}{細い長方形に分けてたそう}{\bklevelrange{2}{4}}{細い長方形に分けてたそう}
\bkmichishirube{4-5,5-3,5-9}{5-11,5-12}

\begin{goal}
\begin{itemize}[label={\color{bkTheme}●}, leftmargin=1.6\zw, itemsep=0pt, topsep=1pt]
\item \bklevel{2}\enspace 曲線の下の面積を、細い長方形の和で近似できる。
\item \bklevel{3}\enspace $\Sigma$ の記号を読み書きできる。
\item \bklevel{4}\enspace 定積分 $\displaystyle\int_a^b f(x)\,\dd x$ の意味を説明し、$x^2$ の定積分を求められる。
\end{itemize}
\end{goal}

\begin{hatena}
\noindent\begin{minipage}[c]{0.56\linewidth}
落ちるボールの速さは $v=10t$ で、グラフは直線でした。だから、進んだ距離は三角形の面積で求められました（5-9）。

では、速さが $v=t^2$ のように、グラフが曲がっているときはどうでしょう。右の色をぬった部分（0秒から1秒までに進んだ距離）は、三角形でも台形でもありません。どうやって面積を求めればよいでしょう。
\end{minipage}\hfill
\begin{minipage}[c]{0.4\linewidth}
\begin{tikzpicture}
\begin{axis}[bkplot, width=54mm, height=42mm, xmin=0, xmax=1.25, ymin=0, ymax=1.2, xtick={1}, ytick={1}, xlabel={$t$}, ylabel={$v$}]
  \addplot[fill=bkYellow!40, draw=none, domain=0:1] {x^2} \closedcycle;
  \addplot[bkcurve, domain=0:1.1] {x^2};
  \node[font=\sffamily\bfseries\small] at (axis cs:0.72,0.2) {?};
\end{axis}
\end{tikzpicture}
\end{minipage}
\end{hatena}

5-9では、三角形の面積を「細い長方形の階段」で考えました。曲がったグラフでも、同じ方法が使えます。細い長方形に分けて、たし合わせるのです。この「細かく分けてたす」計算が、\textbf{積分}です。

この章は、微分の章（5-1〜5-3）と同じように、1段ずつ進みます。はじめに方眼のマスを数え、次に長方形を4本だけならべて手で計算し、それから本数をふやします。最後に、本数をかぎりなくふやした先に名前をつけます。

%% ---------------- やってみよう・思い出そう ----------------
\begin{yattemiyou}[方眼のマスを数えて、面積を見積もろう]
はてな？の $v=t^2$ を、数学でよく使う文字で $y=x^2$ と書きます（横軸が $x$、たて軸が $y$。グラフの形は同じです）。右の方眼は、たても横も1目もりが0.1です。$x=0$ から1までの、グラフの下の部分のマスを数えます。

\noindent\begin{minipage}[c]{0.44\linewidth}\raggedright
\begin{enumerate}[label={\bkcircled[bkGreen]{\arabic*}}, leftmargin=2.4\zw, itemsep=2pt, topsep=3pt]
\item 全部がグラフの下に入っているマス（青）は、何こかな？
\item グラフが通っていて、一部だけ入っているマス（黄）は、何こかな？
\item 黄のマスは、2つで1マスと数えよう（5-9と同じ）。全部で何マスかな？
\item 1マスの面積は $0.1\times0.1=0.01$。グラフの下の面積は、およそいくつかな？
\end{enumerate}
\end{minipage}\hfill
\begin{minipage}[c]{0.53\linewidth}\centering
\begin{tikzpicture}[x=48mm, y=48mm]
  % 第11回：列ごとに、全部入るマスの数（青）と、一部だけ入るマスの数（黄）。Python で数えた
  \foreach \i/\f/\p in {1/0/1,2/0/1,3/0/1,4/0/2,5/1/2,6/2/2,7/3/2,8/4/3,9/6/3,10/8/2} {
    {\ifnum\f>0 \fill[bkBlue!35] ({(\i-1)/10},0) rectangle ({\i/10},{\f/10});\fi
     \fill[bkYellow!60] ({(\i-1)/10},{\f/10}) rectangle ({\i/10},{(\f+\p)/10});}}
  \draw[bkGray!70, line width=0.3pt] (0,0) grid[step=0.1] (1,1);
  \draw[bkBlue!80!black, line width=1.3pt] plot[domain=0:1, samples=60] (\x,{\x*\x});
  \draw[bkInk, line width=0.7pt, -{Stealth[length=4pt]}] (0,0) -- (1.09,0) node[right, font=\sffamily\small, inner sep=1pt] {$x$};
  \draw[bkInk, line width=0.7pt, -{Stealth[length=4pt]}] (0,0) -- (0,1.09) node[above, font=\sffamily\small, inner sep=1pt] {$y$};
  \foreach \v in {0.5,1} {
    \node[font=\sffamily\footnotesize, anchor=north, inner sep=1pt] at (\v,-0.015) {\v};
    \node[font=\sffamily\footnotesize, anchor=east, inner sep=1pt] at (-0.015,\v) {\v};}
  \node[font=\sffamily\footnotesize, anchor=north east, inner sep=1pt] at (-0.01,-0.01) {0};
  \node[font=\sffamily\footnotesize, text=bkBlue!80!black, anchor=west, inner sep=1pt] at (1.02,0.93) {$y=x^2$};
\end{tikzpicture}
\end{minipage}
\end{yattemiyou}

\noindent 青のマスは24こ、黄のマスは19こです。黄は2つで1マスと数えるので、$19\div2=9.5$ マス。全部で $24+9.5=33.5$ マスです。1マスが0.01なので、面積はおよそ
\[ 33.5\times0.01=0.335 \]
です。でも、黄のマスを「2つで1マス」と数えたのは、だいたいの決め方でした。マスの数え方で、答えが少し変わってしまいます。もっとはっきり決まる方法で、ほんとうの面積を追いつめていきましょう。

\begin{omoidasou}
\bkrecall{5-9}{$v$-$t$ グラフの面積は進んだ距離。方眼のマスや細い長方形の階段で考えた。}
\bkrecall{2-4}{分数のかけ算は、分子どうし・分母どうしをかける。$\left(\dfrac{2}{4}\right)^2=\dfrac{2}{4}\times\dfrac{2}{4}=\dfrac{4}{16}$}
\bkrecall{2-12}{$n^2\times n=n^3$（かける回数をたす）。まとめてわるときは、1つずつわってよい：$(A+B)\div h=A\div h+B\div h$}
\bkrecall{5-3}{$\displaystyle\lim_{n\to\infty}$ は「$n$ をかぎりなく大きくしたときの、近づく先」。$\infty$ は「無限大」。}
\bkrecall{4-5}{円の面積は、細いおうぎ形に切って並べかえ、長方形に近づけて求めた。}
\end{omoidasou}

%% ---------------- 数学の道具 ----------------
\bksection{doguu}{数学の道具}

\bksubsection{1}{長方形を4本ならべて、手で計算する}

\label{bk:5-10:hasami}マスのかわりに、たてに細長い長方形をならべます。まず、長方形を4本だけにして、1段ずつ考えましょう。

\noindent{\sffamily\bfseries\color{bkTheme}段1　はばを決める}\par
0から1までを、4つの等しいはばに分けます。長さ1を4つに等しく分けるので、1本のはばは
\[ 1\div4=\frac{1}{4} \]
です（わり算は分数で書ける、2-4）。

\noindent{\sffamily\bfseries\color{bkTheme}段2　$k$ 本めの右はしの場所を決める}\par
左から数えて1本め、2本め、3本め、4本めの長方形があります。この「何本め」を、文字 $k$ で表します。$k$ 本めの長方形の右はしは、0から、はば $\dfrac{1}{4}$ を $k$ こ進んだ所です。

\begin{bkfigure}
\begin{tikzpicture}[x=1mm, y=1mm]
  \draw[bkInk, line width=0.8pt] (0,0) -- (106,0);
  \foreach \k/\lab in {0/0, 1/\frac{1}{4}, 2/\frac{2}{4}, 3/\frac{3}{4}, 4/\frac{4}{4}} {
    \draw[bkInk, line width=0.8pt] (25*\k,-1.6) -- (25*\k,1.6);
    \node[font=\small, anchor=north, inner sep=1pt] at (25*\k,-2.6) {$\lab$};}
  \node[font=\sffamily\footnotesize, anchor=north, inner sep=1pt] at (100,-9.6) {（＝1）};
  \foreach \j in {1,2,3} {
    \draw[bkOrange!90!black, line width=0.9pt, -{Stealth[length=4pt]}] (25*\j-24,2.4) to[bend left=55] (25*\j-1,2.4);
    \node[font=\footnotesize, text=bkOrange!80!black, anchor=south, inner sep=1pt] at (25*\j-12.5,10.6) {$+\frac{1}{4}$};}
  \node[font=\sffamily\footnotesize, anchor=west, align=left, inner sep=0pt] at (78,9) {3本めの右はしは、\\$\frac{1}{4}$ を3こ進んだ $\frac{3}{4}$};
\end{tikzpicture}
\bkfigcaption{0から1までを4等分した目もり。はばを $k$ こ進むと、$k$ 本めの右はしに着く。}
\end{bkfigure}

だから、$k$ 本めの右はしの場所は
\[ x=\frac{1}{4}\times k=\frac{k}{4} \]
です。$k=3$ なら $\dfrac{3}{4}$、$k=4$ なら $\dfrac{4}{4}=1$ で、ちょうど右の終わりに着きます。

\noindent{\sffamily\bfseries\color{bkTheme}段3　長方形の高さを決める}\par
長方形の上の辺を、どの高さにするかを決めます。ここでは、長方形の\textbf{右上の角が、ちょうどグラフにのる}ようにします。これを\textbf{右はしの長方形}とよびます。右上の角は、右はし $x=\dfrac{k}{4}$ の真上にあるグラフの点です。グラフは $y=x^2$ なので、その点の高さは、$x^2$ の $x$ に $\dfrac{k}{4}$ を入れた
\[ y=\left(\frac{k}{4}\right)^2 \]
です。これが $k$ 本めの長方形の高さです。2乗がつくのは、グラフの式が $y=x^2$ だからです。

\begin{bkfigure}
\begin{tikzpicture}[x=84mm, y=42mm]
  \foreach \k in {1,...,4} {
    \fill[bkOrange!25] ({(\k-1)/4},0) rectangle ({\k/4},{(\k/4)^2});
    \draw[bkOrange!80!black, line width=0.5pt] ({(\k-1)/4},0) rectangle ({\k/4},{(\k/4)^2});
    \node[font=\small, anchor=south, inner sep=1pt] at ({(\k-0.5)/4},{(\k/4)^2+0.025}) {$\left(\frac{\k}{4}\right)^2$};
    \node[font=\sffamily\footnotesize, anchor=north, inner sep=1pt] at ({(\k-0.5)/4},-0.15) {\k 本め};}
  \draw[bkBlue!80!black, line width=1.3pt] plot[domain=0:1.04, samples=50] (\x,{\x*\x});
  \foreach \k in {1,...,4} {\fill[bkInk] ({\k/4},{(\k/4)^2}) circle (1.7pt);}
  \draw[bkInk, line width=0.7pt, -{Stealth[length=4pt]}] (0,0) -- (1.1,0) node[right, font=\sffamily\small, inner sep=1pt] {$x$};
  \draw[bkInk, line width=0.7pt, -{Stealth[length=4pt]}] (0,0) -- (0,1.16) node[above, font=\sffamily\small, inner sep=1pt] {$y$};
  \foreach \k in {1,...,4} {\node[font=\footnotesize, anchor=north, inner sep=1pt] at ({\k/4},-0.02) {$\frac{\k}{4}$};}
  \node[font=\sffamily\footnotesize, anchor=north east, inner sep=1pt] at (-0.005,-0.01) {0};
  \node[font=\sffamily\footnotesize, anchor=east, inner sep=1pt] at (-0.01,1) {1};
  \node[font=\sffamily\small, text=bkBlue!80!black, anchor=west, inner sep=1pt] at (1.05,1.05) {$y=x^2$};
\end{tikzpicture}
\bkfigcaption{右はしの長方形。黒い点（右上の角）がグラフにのる。長方形の上の数が高さ。}
\end{bkfigure}

\noindent{\sffamily\bfseries\color{bkTheme}段4　面積を計算して、たす}\par
長方形の面積は、高さ×はばです。分数のかけ算は、分子どうし・分母どうしをかけます（2-4）。表にまとめましょう。

\begin{bkcenter}
\begin{bktable}{colspec={Q[l,m,wd=34mm] *{4}{Q[c,m,wd=17mm]}}, rowsep=2pt}
$k$（何本め） & 1 & 2 & 3 & 4 \\
右はし $x=\dfrac{k}{4}$ & $\dfrac{1}{4}$ & $\dfrac{2}{4}$ & $\dfrac{3}{4}$ & $\dfrac{4}{4}$ \\
高さ $\left(\dfrac{k}{4}\right)^2$ & $\dfrac{1}{16}$ & $\dfrac{4}{16}$ & $\dfrac{9}{16}$ & $\dfrac{16}{16}$ \\
はば & $\dfrac{1}{4}$ & $\dfrac{1}{4}$ & $\dfrac{1}{4}$ & $\dfrac{1}{4}$ \\
面積＝高さ×はば & $\dfrac{1}{64}$ & $\dfrac{4}{64}$ & $\dfrac{9}{64}$ & $\dfrac{16}{64}$ \\
\end{bktable}
\bktabcaption{右はしの長方形4本。たとえば3本めは、$\frac{9}{16}\times\frac{1}{4}=\frac{9\times1}{16\times4}=\frac{9}{64}$}
\end{bkcenter}

4本の面積をたすと、分母がそろっているので、分子だけをたせます。
\[ \frac{1}{64}+\frac{4}{64}+\frac{9}{64}+\frac{16}{64}=\frac{30}{64}=0.46875 \]

\noindent{\sffamily\bfseries\color{bkTheme}左はしの長方形でも計算する}\par
こんどは、長方形の\textbf{左上の角}がグラフにのるようにします（\textbf{左はしの長方形}）。$k$ 本めの左はしは、右はしより1つ前の目もり $\dfrac{k-1}{4}$ なので、高さは $\left(\dfrac{0}{4}\right)^2$、$\left(\dfrac{1}{4}\right)^2$、$\left(\dfrac{2}{4}\right)^2$、$\left(\dfrac{3}{4}\right)^2$ です。1本めは高さが0なので、見えません。
\[ \frac{0}{64}+\frac{1}{64}+\frac{4}{64}+\frac{9}{64}=\frac{14}{64}=0.21875 \]

\begin{bkfigure}
\begin{tikzpicture}
\begin{axis}[bkplot, name=b1, width=60mm, height=48mm, xmin=0, xmax=1.15, ymin=0, ymax=1.15, xtick={0.25,0.5,0.75,1}, ytick={0.5,1}, xlabel={$x$}, ylabel={$y$}, title={\sffamily\small 右はしの長方形（大きめ）}]
  \foreach \k in {1,...,4} {\edef\tmp{\noexpand\fill[bkOrange!30, draw=bkOrange!70, line width=0.5pt] (axis cs:{(\k-1)/4},0) rectangle (axis cs:{\k/4},{(\k/4)^2});}\tmp}
  \addplot[bkcurve, domain=0:1.05] {x^2};
\end{axis}
\begin{axis}[bkplot, at={(b1.east)}, anchor=west, xshift=14mm, width=60mm, height=48mm, xmin=0, xmax=1.15, ymin=0, ymax=1.15, xtick={0.25,0.5,0.75,1}, ytick={0.5,1}, xlabel={$x$}, ylabel={$y$}, title={\sffamily\small 左はしの長方形（小さめ）}]
  \foreach \k in {1,...,4} {\edef\tmp{\noexpand\fill[bkBlue!25, draw=bkBlue!70, line width=0.5pt] (axis cs:{(\k-1)/4},0) rectangle (axis cs:{\k/4},{((\k-1)/4)^2});}\tmp}
  \addplot[bkcurve, domain=0:1.05] {x^2};
\end{axis}
\end{tikzpicture}
\bkfigcaption{$n=4$ の場合。右はしの長方形はグラフの上にはみ出し、左はしの長方形はグラフの下にすき間がある。}
\end{bkfigure}

右はしの長方形は、グラフの上に少しはみ出しています。だから、ほんとうの面積より大きい。左はしの長方形は、グラフの下にすき間があります。だから、ほんとうの面積より小さい。ほんとうの面積を $S$ とすると、
\[ 0.21875<S<0.46875 \]
です。$S$ は、2つの数のあいだに\textbf{はさまれて}います。やってみようの0.335も、このあいだにあります。

\bksubsection{2}{長方形をふやして、はさみうちにする}

4本では、はさんだはばが $0.46875-0.21875=0.25$ もあって、ゆるすぎます。長方形を細くして、本数をふやしましょう。本数を $n$ 本にしても、段1〜段4は同じです。

\begin{bkcenter}
\begin{bktable}{colspec={Q[l,m,wd=30mm] Q[c,m,wd=26mm] Q[c,m,wd=26mm] Q[c,m,wd=34mm]}, rowsep=2pt, row{1}={font=\sffamily\bfseries\small}}
 & 4本 & 10本 & $n$ 本 \\
はば & $\dfrac{1}{4}$ & $\dfrac{1}{10}$ & $\dfrac{1}{n}$\quad（$1\div n$） \\
$k$ 本めの右はし & $\dfrac{k}{4}$ & $\dfrac{k}{10}$ & $\dfrac{k}{n}$ \\
$k$ 本めの高さ & $\left(\dfrac{k}{4}\right)^2$ & $\left(\dfrac{k}{10}\right)^2$ & $\left(\dfrac{k}{n}\right)^2$ \\
$k$ 本めの面積 & $\left(\dfrac{k}{4}\right)^2\times\dfrac{1}{4}$ & $\left(\dfrac{k}{10}\right)^2\times\dfrac{1}{10}$ & $\left(\dfrac{k}{n}\right)^2\times\dfrac{1}{n}$ \\
\end{bktable}
\bktabcaption{本数が4でも10でも $n$ でも、考え方は同じ。4 のところが $n$ に変わるだけ。}
\end{bkcenter}

10本なら、表計算（2-18）で計算できます。A列に $k$、B列に右はし、C列に高さ、D列に面積を入れます。

\begin{bkcenter}
\bksheet{4}{24}{%
1 & $k$ & 右はし $x$ & 高さ $x^2$ & 面積 \\
2 & 1 & 0.1 & 0.01 & 0.001 \\
3 & 2 & 0.2 & 0.04 & 0.004 \\
4 & 3 & 0.3 & 0.09 & 0.009 \\
5 & 4 & 0.4 & 0.16 & 0.016 \\
6 & 5 & 0.5 & 0.25 & 0.025 \\
7 & 6 & 0.6 & 0.36 & 0.036 \\
8 & 7 & 0.7 & 0.49 & 0.049 \\
9 & 8 & 0.8 & 0.64 & 0.064 \\
10 & 9 & 0.9 & 0.81 & 0.081 \\
11 & 10 & 1 & 1 & 0.1 \\
12 & & & 合計 & 0.385 \\}
\bktabcaption{B2 に \texttt{=A2/10}（右はし $\frac{k}{10}$）、C2 に \texttt{=B2\^{}2}（2乗）、D2 に \texttt{=C2*0.1}（はば0.1をかける）と入れて、下へコピーする。}
\end{bkcenter}

D12 には \texttt{=SUM(D2:D11)} と入れます。SUM は英語で「和（たし算の答え）」という意味で、「D2 から D11 までを全部たす」という命令です。右はしの和は0.385になりました。左はしの和は、同じように計算して0.285です。

$n$ を100、1000とふやすと、次の表のようになります（「2つの平均」は、長方形のかわりに台形をならべた和です。あとの例題8で使います）。

\begin{bkcenter}
\bksheet{4}{24}{%
1 & $n$ & 左はしの和 & 右はしの和 & 2つの平均 \\
2 & 4 & 0.21875 & 0.46875 & 0.34375 \\
3 & 10 & 0.285 & 0.385 & 0.335 \\
4 & 100 & 0.32835 & 0.33835 & 0.33335 \\
5 & 1000 & 0.3328335 & 0.3338335 & 0.3333335 \\}
\bktabcaption{$n$ を大きくすると、左はしの和と右はしの和が、両がわから同じ数に近づいていく。}
\end{bkcenter}

\noindent{\sffamily\bfseries\color{bkTheme}2つの和の差は、いつも $\dfrac{1}{n}$}\par
表を見ると、右はしの和と左はしの和の差は、$n=4$ で0.25、$n=10$ で0.1、$n=100$ で0.01です。いつも $\dfrac{1}{n}$ になっています。そのわけを、図で見ましょう。

\begin{bkfigure}
\begin{tikzpicture}[x=40mm, y=34mm]
  \begin{scope}
    \fill[bkGreen!45] (0.25,0) rectangle (0.5,0.0625);
    \fill[bkBlue!35] (0.5,0) rectangle (0.75,0.25);
    \fill[bkPurple!35] (0.75,0) rectangle (1,0.5625);
    \draw[bkInk!70, line width=0.4pt] (0.25,0) rectangle (0.5,0.0625) (0.5,0) rectangle (0.75,0.25) (0.75,0) rectangle (1,0.5625);
    \draw[bkBlue!80!black, line width=1.1pt] plot[domain=0:1.03, samples=40] (\x,{\x*\x});
    \draw[bkInk, line width=0.6pt] (0,0) -- (1.06,0) (0,0) -- (0,1.08);
    \node[font=\sffamily\small, anchor=south, inner sep=1pt] at (0.5,1.1) {左はしの長方形};
    \node[font=\sffamily\footnotesize, anchor=north, inner sep=1pt] at (0,-0.02) {0};
    \node[font=\sffamily\footnotesize, anchor=north, inner sep=1pt] at (1,-0.02) {1};
  \end{scope}
  \draw[bk arrow, bkGray, line width=0.9pt] (1.12,0.4) -- (1.36,0.4);
  \node[font=\sffamily\footnotesize, anchor=south, align=center, inner sep=1pt] at (1.24,0.47) {1つ左へ\\ずらす};
  \begin{scope}[shift={(1.5,0)}]
    \fill[bkGreen!45] (0,0) rectangle (0.25,0.0625);
    \fill[bkBlue!35] (0.25,0) rectangle (0.5,0.25);
    \fill[bkPurple!35] (0.5,0) rectangle (0.75,0.5625);
    \fill[bkRed!50] (0.75,0) rectangle (1,1);
    \draw[bkInk!70, line width=0.4pt] (0,0) rectangle (0.25,0.0625) (0.25,0) rectangle (0.5,0.25) (0.5,0) rectangle (0.75,0.5625) (0.75,0) rectangle (1,1);
    \draw[bkBlue!80!black, line width=1.1pt] plot[domain=0:1.03, samples=40] (\x,{\x*\x});
    \draw[bkInk, line width=0.6pt] (0,0) -- (1.06,0) (0,0) -- (0,1.08);
    \node[font=\sffamily\small, anchor=south, inner sep=1pt] at (0.5,1.1) {右はしの長方形};
    \node[font=\sffamily\footnotesize, anchor=north, inner sep=1pt] at (0,-0.02) {0};
    \node[font=\sffamily\footnotesize, anchor=north, inner sep=1pt] at (1,-0.02) {1};
    \node[font=\sffamily\footnotesize, anchor=west, align=left, inner sep=0pt] at (1.1,0.55) {赤の1本だけ\\多い。\\はば $\frac{1}{4}$、\\高さ1};
  \end{scope}
\end{tikzpicture}
\bkfigcaption{同じ色の長方形は、同じ形。右はしの長方形は、左はしの長方形を1つ左へずらして、赤の1本を足したもの。}
\end{bkfigure}

左はしの長方形の2本め（緑）は、右はしの長方形の1本めと同じ高さ $\left(\frac{1}{4}\right)^2$ です。3本め（青）と2本め、4本め（むらさき）と3本めも同じです。だから、ちがうのは右はしの4本め（赤）だけです。赤の長方形は、高さ $1^2=1$、はば $\frac{1}{n}$ なので、面積は $1\times\dfrac{1}{n}=\dfrac{1}{n}$ です。

$n=1000$ なら、差はたったの0.001です。ほんとうの面積 $S$ は
\[ 0.3328335<S<0.3338335 \]
のあいだにはさまれています。$n$ をもっと大きくすると、差 $\dfrac{1}{n}$ はいくらでも小さくなり、2つの和は同じ数に近づきます。表の数から、近づく先は $\dfrac{1}{3}=0.3333\cdots$ のようです。ほんとうに $\dfrac{1}{3}$ なのかを、④で式を使って確かめます。

\bksubsection{3}{たし算をまとめて書く：$\Sigma$}

\label{bk:5-10:sigma}長方形が1000本あると、たし算を全部書くのは大変です。そこで、「番号を変えながら、全部たす」ことを短く書く記号を使います。

\begin{bkfigure}
\begin{tikzpicture}[x=1mm, y=1mm]
  \node[inner sep=0pt, font=\LARGE] (S) at (0,0) {$\displaystyle\sum$};
  \node[inner sep=1pt, anchor=south, font=\normalsize] (up) at ([yshift=0.8mm]S.north) {$4$};
  \node[inner sep=1pt, anchor=north, font=\normalsize] (lo) at ([yshift=-0.8mm]S.south) {$k=1$};
  \node[inner sep=1pt, anchor=west, font=\LARGE] (b) at ([xshift=1mm]S.east) {$k^2$};
  \tikzset{lab/.style={font=\sffamily\footnotesize, text=bkInk, inner sep=1pt, align=left}, ar/.style={bkTheme, line width=0.7pt, -{Stealth[length=4pt]}, shorten >=1.2mm}}
  \node[lab, anchor=east] (l1) at (-16,9) {おわりの番号（4まで）};
  \node[lab, anchor=east] (l2) at (-16,-9) {はじめの番号（$k=1$ から）};
  \node[lab, anchor=west] (l3) at (20,0) {たすもの\\（$k$ に番号を入れる）};
  \draw[ar] (l1.east) -- (up.west);
  \draw[ar] (l2.east) -- (lo.west);
  \draw[ar] (l3.west) -- (b.east);
\end{tikzpicture}
\bkfigcaption{$\Sigma$ の読み方。「$k$ を1から4まで1ずつ変えて、$k^2$ を全部たす」}
\end{bkfigure}

この記号は、次のように計算します。

\begin{bkcenter}
\begin{bktable}{colspec={Q[c,m,wd=22mm] Q[c,m,wd=26mm] Q[c,m,wd=20mm]}, rowsep=1.5pt, row{1}={font=\sffamily\bfseries\small}}
$k$ に入れる数 & $k^2$ & 値 \\
1 & $1^2$ & 1 \\
2 & $2^2$ & 4 \\
3 & $3^2$ & 9 \\
4 & $4^2$ & 16 \\
\SetCell[c=2]{r} 全部たす & & 30 \\
\end{bktable}
\end{bkcenter}
\[ \sum_{k=1}^{4}k^2=1^2+2^2+3^2+4^2=1+4+9+16=30 \]

$\Sigma$ は、ギリシャ文字の「シグマ」で、英語の sum（和）の頭文字 S にあたります。表計算の SUM と同じ意味です。上の式は「シグマ、$k$ イコール1から4まで、$k$ の2乗」と読みます。\bkidx{しぐま}{$\Sigma$（シグマ）}

\noindent\textbf{例1}\quad $\displaystyle\sum_{k=1}^{3}k=1+2+3=6$\quad（$k$ そのものをたす）\par
\noindent\textbf{例2}\quad $\displaystyle\sum_{k=1}^{4}(2k+1)=3+5+7+9=24$\quad（$k=1,2,3,4$ を順に入れる。$k=1$ なら $2\times1+1=3$）

\noindent{\sffamily\bfseries\color{bkTheme}長方形の和を $\Sigma$ で書く}\par
4本の右はしの長方形の和は、$k$ 本めの面積 $\left(\dfrac{k}{4}\right)^2\times\dfrac{1}{4}$ を、$k=1$ から4までたしたものです。
\[ \sum_{k=1}^{4}\left(\frac{k}{4}\right)^2\times\frac{1}{4}=\left(\frac{1}{4}\right)^2\times\frac{1}{4}+\left(\frac{2}{4}\right)^2\times\frac{1}{4}+\left(\frac{3}{4}\right)^2\times\frac{1}{4}+\left(\frac{4}{4}\right)^2\times\frac{1}{4} \]
$n$ 本なら、4を $n$ にかえて、
\[ S_n=\sum_{k=1}^{n}\left(\frac{k}{n}\right)^2\times\frac{1}{n} \]
と書けます。$S_n$ は「$n$ 本のときの、右はしの長方形の和」という名前です。

\noindent{\sffamily\bfseries\color{bkTheme}1から $n$ までの和}\par
あとで使うので、$\displaystyle\sum_{k=1}^{n}k=1+2+\cdots+n$ を計算しておきます。$n=4$ で、ブロックの階段を考えます。

\begin{bkfigure}
\begin{tikzpicture}[x=6mm, y=6mm]
  \foreach \j in {1,...,4} {\foreach \i in {1,...,\j} {{\draw[bkBlue!80!black, fill=bkBlue!30, line width=0.5pt] (\j-1,\i-1) rectangle (\j,\i);}}}
  \foreach \j in {1,...,4} {\node[font=\sffamily\footnotesize, anchor=north, inner sep=1pt] at (\j-0.5,-0.15) {\j};}
  \node[font=\sffamily\small, anchor=north, inner sep=1pt] at (2,-1.1) {$1+2+3+4$};
  \draw[bk arrow, bkGray, line width=0.9pt] (5,2.5) -- (7,2.5);
  \node[font=\sffamily\footnotesize, anchor=south, align=center, inner sep=1pt] at (6,2.8) {同じ階段を\\さかさにのせる};
  \begin{scope}[shift={(8,0)}]
    \foreach \j in {1,...,4} {\foreach \i in {1,...,\j} {{\draw[bkBlue!80!black, fill=bkBlue!30, line width=0.5pt] (\j-1,\i-1) rectangle (\j,\i);}}}
    \foreach \j in {1,...,4} {\foreach \i in {\j,...,4} {{\draw[bkOrange!80!black, fill=bkOrange!35, line width=0.5pt] (\j-1,\i) rectangle (\j,\i+1);}}}
    \draw[bkGray, line width=0.5pt, <->] (0,-0.5) -- (4,-0.5);
    \node[font=\sffamily\footnotesize, anchor=north, inner sep=1pt] at (2,-0.7) {よこ4};
    \draw[bkGray, line width=0.5pt, <->] (4.5,0) -- (4.5,5);
    \node[font=\sffamily\footnotesize, anchor=west, inner sep=1pt] at (4.7,2.5) {たて5（$4+1$）};
  \end{scope}
\end{tikzpicture}
\bkfigcaption{階段を2つ組み合わせると、よこ4・たて5の長方形になる。}
\end{bkfigure}

長方形のブロックは $4\times5=20$ こ。階段はその半分なので、
\[ 1+2+3+4=20\div2=10 \]
です。$n$ 段の階段でも、2つ組み合わせると、よこ $n$、たて $n+1$ の長方形になります。だから
\[ \sum_{k=1}^{n}k=1+2+\cdots+n=\frac{n(n+1)}{2} \]
です。$n=3$ なら $\dfrac{3\times4}{2}=6$ で、例1と合います。

\bksubsection{4}{$n$ を大きくした先を、式で求める}

\label{bk:5-10:kyokugen}②の表では、右はしの和 $S_n$ は $\dfrac{1}{3}$ に近づくようでした。$S_n$ を $n$ の式で表して、ほんとうに $\dfrac{1}{3}$ になるかを確かめます。

\noindent{\sffamily\bfseries\color{bkTheme}段1　1本の面積を、1つの分数にする}\par
分数の2乗は、分子と分母をそれぞれ2乗します（2-4）。そのあと $\dfrac{1}{n}$ をかけると、分母は $n^2\times n=n^3$ になります（2-12）。
\[ \left(\frac{k}{n}\right)^2\times\frac{1}{n}=\frac{k^2}{n^2}\times\frac{1}{n}=\frac{k^2}{n^3} \]
$n=4$、$k=2$ で確かめると、$\dfrac{2^2}{4^3}=\dfrac{4}{64}$。①の表の2本めと合います。

\noindent{\sffamily\bfseries\color{bkTheme}段2　分母が同じなので、分子だけをたす}\par
どの長方形も分母が $n^3$ でそろっているので、分子だけをたせます。①で $\dfrac{1+4+9+16}{64}$ としたのと同じです。
\[ S_n=\frac{1^2}{n^3}+\frac{2^2}{n^3}+\cdots+\frac{n^2}{n^3}=\frac{1^2+2^2+\cdots+n^2}{n^3} \]

\noindent{\sffamily\bfseries\color{bkTheme}段3　2乗の和を、$n$ の式にする}\par
分子の $1^2+2^2+\cdots+n^2$ には、次の式があります。まず、表で合っているかを確かめましょう。
\[ 1^2+2^2+\cdots+n^2=\frac{n(n+1)(2n+1)}{6} \]

\begin{bkcenter}
\begin{bktable}{colspec={Q[l,m,wd=44mm] *{6}{Q[c,m,wd=10mm]}}, rowsep=1.5pt}
$n$ & 1 & 2 & 3 & 4 & 5 & 6 \\
$1^2+2^2+\cdots+n^2$ & 1 & 5 & 14 & 30 & 55 & 91 \\
$\dfrac{n(n+1)(2n+1)}{6}$ & 1 & 5 & 14 & 30 & 55 & 91 \\
\end{bktable}
\bktabcaption{たとえば $n=4$ なら $\frac{4\times5\times9}{6}=\frac{180}{6}=30$。どの $n$ でも合う。}
\end{bkcenter}

なぜこの式になるのかは、もっと知りたいで説明します。

\begin{motto}[4]
\noindent\textbf{2乗の和の式のわけ}\par
5-4の立方体の図で、1辺 $x$ の立方体を $h$ だけ大きくすると、$(x+h)^3=x^3+3x^2h+3xh^2+h^3$ でした。$x=k$、$h=1$ とすると、
\[ (k+1)^3-k^3=3k^2+3k+1 \]
です。$k=1,\,2,\,3,\,4$ を入れて、たてにたしてみます。
\begin{align*}
2^3-1^3&=3\times1^2+3\times1+1\\
3^3-2^3&=3\times2^2+3\times2+1\\
4^3-3^3&=3\times3^2+3\times3+1\\
5^3-4^3&=3\times4^2+3\times4+1
\end{align*}
左がわをたすと、$2^3$ と $-2^3$、$3^3$ と $-3^3$、$4^3$ と $-4^3$ が消えて、$5^3-1^3=124$ だけが残ります。右がわは、次の式です。
\[ 3\times(1^2+2^2+3^2+4^2)+3\times(1+2+3+4)+4 \]
\noindent $1+2+3+4=10$（③）なので、
\[ 124=3\times(1^2+2^2+3^2+4^2)+30+4 \]
これから $1^2+2^2+3^2+4^2=(124-34)\div3=30$。4を $n$ にかえて同じ計算をすると、$(n+1)^3-1=3\times(1^2+\cdots+n^2)+3\times\dfrac{n(n+1)}{2}+n$ となり、これを整理すると上の式になります。
\end{motto}

\noindent{\sffamily\bfseries\color{bkTheme}段4　$S_n$ を、$\dfrac{1}{n}$ の式に直す}\par
段2と段3から、
\[ S_n=\frac{n(n+1)(2n+1)}{6}\div n^3=\frac{n(n+1)(2n+1)}{6\times n\times n\times n} \]
分母の $n\times n\times n$ の $n$ を1つずつ、分子のかっこに配ります。
\[ S_n=\frac{1}{6}\times\frac{n}{n}\times\frac{n+1}{n}\times\frac{2n+1}{n} \]
ここで、$\dfrac{n}{n}=1$。そして、まとめてわるときは1つずつわってよい（2-12）ので、
\[ \frac{n+1}{n}=\frac{n}{n}+\frac{1}{n}=1+\frac{1}{n},\qquad \frac{2n+1}{n}=\frac{2n}{n}+\frac{1}{n}=2+\frac{1}{n} \]
です。だから
\[ S_n=\frac{1}{6}\Bigl(1+\frac{1}{n}\Bigr)\Bigl(2+\frac{1}{n}\Bigr) \]
となります。$n=4$ を入れると、$\dfrac{1}{6}\times1.25\times2.25=0.46875$。$n=10$ なら $\dfrac{1}{6}\times1.1\times2.1=0.385$。②の表とぴったり合います。

\noindent{\sffamily\bfseries\color{bkTheme}段5　$n$ をかぎりなく大きくする}\par
$n$ を大きくすると、$\dfrac{1}{n}$ は0.25、0.1、0.01、0.001、……と、0にかぎりなく近づきます（5-3）。だから、$1+\dfrac{1}{n}$ は1に、$2+\dfrac{1}{n}$ は2に近づきます。
\[ \lim_{n\to\infty}S_n=\frac{1}{6}\times1\times2=\frac{2}{6}=\frac{1}{3} \]
表計算で見つけた $\dfrac{1}{3}$ が、式でも確かめられました。左はしの和は、右はしの和より $\dfrac{1}{n}$ 小さいだけ（②）なので、$\dfrac{1}{3}-0=\dfrac{1}{3}$ に近づきます。左右からはさまれたほんとうの面積 $S$ は、ちょうど $\dfrac{1}{3}$ です。

\bkhatenaref の答え：$v=t^2$ のグラフの、0秒から1秒までの面積は $\dfrac{1}{3}$ なので、ボールは0秒から1秒までに $\dfrac{1}{3}\,\text{m}$（約\qty{0.33}{m}）進みます。やってみようの0.335とも、よく合っています。

\bksubsection{5}{はばを0に近づけた先が、定積分}

\label{bk:5-10:teisekibun}はばをせまくしていくようすを、図で見てみましょう。

\begin{bkfigure}
\begin{tikzpicture}[x=1mm, y=1mm]
  \foreach \n/\px/\w/\sm/\ov in {4/0/{$\frac{1}{4}$}/0.469/0.135, 8/36/{$\frac{1}{8}$}/0.398/0.065, 16/72/{$\frac{1}{16}$}/0.365/0.032} {
    \begin{scope}[shift={(\px,8)}]
      % 長方形（右はし）。はみ出しは こい色＋斜線
      \foreach \k in {1,...,\n} {
        \fill[bkOrange!70] ({26*(\k-1)/\n},0) rectangle ({26*\k/\n},{22*(\k/\n)^2});
        \fill[pattern=north east lines, pattern color=white] ({26*(\k-1)/\n},0) rectangle ({26*\k/\n},{22*(\k/\n)^2});}
      \fill[bkOrange!25] (0,0) -- plot[domain=0:1, samples=40] ({26*\x},{22*\x*\x}) -- (26,0) -- cycle;
      \foreach \k in {1,...,\n} {\draw[bkOrange!80!black, line width=0.3pt] ({26*(\k-1)/\n},0) rectangle ({26*\k/\n},{22*(\k/\n)^2});}
      \draw[bkBlue!80!black, line width=1pt] plot[domain=0:1.04, samples=40] ({26*\x},{22*\x*\x});
      \draw[bkInk, line width=0.6pt] (0,0) -- (29,0) (0,0) -- (0,25);
      \node[font=\sffamily\footnotesize, anchor=north, inner sep=1pt] at (26,-0.3) {1};
      \node[font=\sffamily\bfseries\footnotesize, anchor=north west, inner sep=0pt, align=left] at (0,-4.2) {はば \w};
      \node[font=\sffamily\footnotesize, anchor=north west, inner sep=0pt, align=left] at (0,-10) {和 \sm\\{\color{bkOrange!80!black}はみ出し \ov}};
    \end{scope}
    \draw[bk arrow, bkGray, line width=0.8pt] (\px+28.5,19) -- (\px+34.5,19);}
  % はばを0に近づけた先
  \begin{scope}[shift={(108,8)}]
    \fill[bkYellow!40] (0,0) -- plot[domain=0:1, samples=40] ({26*\x},{22*\x*\x}) -- (26,0) -- cycle;
    \draw[bkBlue!80!black, line width=1pt] plot[domain=0:1.04, samples=40] ({26*\x},{22*\x*\x});
    \draw[bkInk, line width=0.6pt] (0,0) -- (29,0) (0,0) -- (0,25);
    \node[font=\sffamily\footnotesize, anchor=north, inner sep=1pt] at (26,-0.3) {1};
    \node[font=\sffamily\bfseries\small] at (17.5,4.5) {$\frac{1}{3}$};
    \node[font=\sffamily\bfseries\footnotesize, anchor=north west, inner sep=0pt, align=left] at (0,-4.2) {はば → 0};
    \node[font=\sffamily\footnotesize, anchor=north west, inner sep=0pt, align=left] at (0,-10) {和 → $\frac{1}{3}$\\{\color{bkOrange!80!black}はみ出し → 0}};
  \end{scope}
\end{tikzpicture}
\bkfigcaption{はばを半分にするたびに、はみ出し（こい色の斜線）は、ほぼ半分になる。はばを0に近づけた先は、曲線の下の面積そのもの。}
\end{bkfigure}

はばを0に近づけていくと、階段の面積は、はてな？の色をぬった部分の面積そのものになります。ここまでにしたことは、いつも次の3段です。

\begin{bkfigure}
\begin{tikzpicture}[x=1mm, y=1mm]
  \tikzset{cv/.style={bkBlue!80!black, line width=1pt}, ax/.style={bkInk, line width=0.6pt}}
  \def\fx(#1){(6+14*(#1)*(#1))}
  \foreach \px/\ttl in {0/{① 細かく分ける}, 46/{② 1本の面積をかける}, 92/{③ 全部たす}} {
    \node[font=\sffamily\bfseries\small, anchor=south, inner sep=1pt] at (\px+16,25) {\ttl};}
  % ① 分ける
  \begin{scope}[shift={(0,0)}]
    \fill[bkYellow!35] (2,0) -- plot[domain=0.0667:1, samples=30] ({30*\x},{6+14*\x*\x}) -- (30,0) -- cycle;
    \foreach \k in {1,...,6} {\draw[bkGray, line width=0.4pt, dash pattern=on 1.2pt off 1pt] ({2+28*\k/6},0) -- ({2+28*\k/6},{6+14*((2+28*\k/6)/30)^2});}
    \draw[cv] plot[domain=0.0667:1.03, samples=30] ({30*\x},{6+14*\x*\x});
    \draw[ax] (0,0) -- (33,0) (0,0) -- (0,23);
    \draw[bkGray, line width=0.5pt, <->] (11.33,-1.6) -- (16,-1.6);
    \node[font=\footnotesize, anchor=north, inner sep=1pt] at (13.67,-2.2) {$\Delta x$};
    \node[font=\footnotesize, anchor=north, inner sep=1pt] at (2,-2.2) {$a$};
    \node[font=\footnotesize, anchor=north, inner sep=1pt] at (30,-2.2) {$b$};
  \end{scope}
  \draw[bk arrow, bkGray, line width=0.8pt] (35,11) -- (42,11);
  % ② かける
  \begin{scope}[shift={(46,0)}]
    \fill[bkYellow!35] (2,0) -- plot[domain=0.0667:1, samples=30] ({30*\x},{6+14*\x*\x}) -- (30,0) -- cycle;
    \fill[bkOrange!45] (16,0) rectangle (20.67,{6+14*(20.67/30)^2});
    \draw[bkOrange!80!black, line width=0.5pt] (16,0) rectangle (20.67,{6+14*(20.67/30)^2});
    \draw[cv] plot[domain=0.0667:1.03, samples=30] ({30*\x},{6+14*\x*\x});
    \draw[ax] (0,0) -- (33,0) (0,0) -- (0,23);
    \draw[bkGray, line width=0.5pt, <->] (16,-1.6) -- (20.67,-1.6);
    \node[font=\footnotesize, anchor=north, inner sep=1pt] at (18.33,-2.2) {$\Delta x$};
    \draw[bkGray, line width=0.5pt, <->] (14.2,0) -- (14.2,{6+14*(20.67/30)^2});
    \node[font=\footnotesize, anchor=east, inner sep=1pt, fill=bkYellow!35] at (13.6,3.2) {$f(x)$};
    \node[font=\sffamily\footnotesize, anchor=north, inner sep=1pt] at (16,-7.5) {高さ×はば＝$f(x)\,\Delta x$};
  \end{scope}
  \draw[bk arrow, bkGray, line width=0.8pt] (81,11) -- (88,11);
  % ③ たす
  \begin{scope}[shift={(92,0)}]
    \foreach \k in {1,...,6} {\fill[bkOrange!35] ({2+28*(\k-1)/6},0) rectangle ({2+28*\k/6},{6+14*((2+28*\k/6)/30)^2});
      \draw[bkOrange!80!black, line width=0.4pt] ({2+28*(\k-1)/6},0) rectangle ({2+28*\k/6},{6+14*((2+28*\k/6)/30)^2});}
    \draw[cv] plot[domain=0.0667:1.03, samples=30] ({30*\x},{6+14*\x*\x});
    \draw[ax] (0,0) -- (33,0) (0,0) -- (0,23);
    \node[font=\sffamily\footnotesize, anchor=north, align=center, inner sep=1pt] at (16,-2.2) {$\sum f(x)\,\Delta x$\\はばを0に近づけて\\$\displaystyle\int_a^b f(x)\,\dd x$};
  \end{scope}
\end{tikzpicture}
\bkfigcaption{積分の3段。① はば $\Delta x$ に細かく分ける。② 1本ずつ「高さ×はば」をかける。③ 全部たして、はばを0に近づける。}
\end{bkfigure}

この近づく先を、次のように書きます。

\begin{bkfigure}
\begin{tikzpicture}[x=1mm, y=1mm]
  \node[inner sep=0pt, font=\Large] (i) at (0,0) {$\displaystyle\int_0^1$};
  \node[inner sep=0pt, font=\Large, anchor=west] (f) at ([xshift=0.6mm]i.east) {$x^2$};
  \node[inner sep=0pt, font=\Large, anchor=west] (d) at ([xshift=1.4mm]f.east) {$\dd x$};
  \node[inner sep=0pt, font=\Large, anchor=west] (e) at ([xshift=1.4mm]d.east) {$=\dfrac{1}{3}$};
  \tikzset{lab/.style={font=\sffamily\footnotesize, text=bkInk, inner sep=1pt, align=left}, ar/.style={bkTheme, line width=0.7pt, -{Stealth[length=4pt]}, shorten >=1.2mm}}
  % 左がわ：∫ と 0・1
  \node[lab, anchor=east] (l1) at ([xshift=-12mm, yshift=6mm]i.center) {$\int$ ＝ たし合わせる\\（sum の S をのばした形）};
  \node[lab, anchor=east] (l2) at ([xshift=-12mm, yshift=-7mm]i.center) {下の 0、上の 1 ＝\\$x=0$ から 1 まで};
  \draw[ar] (l1.east) -- ([xshift=-0.8mm, yshift=2mm]i.center);
  \draw[ar] (l2.east) -- ([xshift=-2mm, yshift=-4.5mm]i.center);
  % 右がわ：式の右に列をつくる。上の2つは式の上を通って下ろす（式に線を重ねない）
  \coordinate (R) at ([xshift=9mm]e.east);
  \node[lab, anchor=west] (l3) at ([yshift=14mm]R |- i.center) {$x^2$ ＝ 長方形の高さ};
  \node[lab, anchor=west] (l4) at ([yshift=9mm]R |- i.center) {$\dd x$ ＝ かぎりなく細いはば};
  \node[lab, anchor=west] (l5) at ([yshift=-1mm]R |- i.center) {$\frac{1}{3}$ ＝ ④で求めた、近づく先};
  \draw[ar] (l3.west) -| (f.north);
  \draw[ar] (l4.west) -| (d.north);
  \draw[ar] (l5.west) -- (e.east |- l5.west);
\end{tikzpicture}
\bkfigcaption{定積分の記号の読み方}
\end{bkfigure}

$\int$ は「インテグラル」と読みます。上の式は「インテグラル、0から1まで、$x$ の2乗、ディーエックス」と読みます。

一般に、グラフ $y=f(x)$ の $x=a$ から $b$ までを細い長方形に分けます。$k$ 本めの長方形の右はしを $x_k$、はばを $\Delta x$ と書きます（④では $x_k=\dfrac{k}{n}$、$\Delta x=\dfrac{1}{n}$。たとえば $n=4$ なら $x_1=\dfrac{1}{4}$、$x_2=\dfrac{2}{4}$、$x_3=\dfrac{3}{4}$、$x_4=\dfrac{4}{4}=1$ で、はばはどれも $\Delta x=\dfrac{1}{4}$）。$\Delta x$ は「$x$ の変化量」（2-3）で、長方形の左はしから右はしまで、$x$ がどれだけ変わったかです。長方形の和は $\sum f(x_k)\,\Delta x$ です。このはばを0に近づけた先を
\[ \sum f(x_k)\,\Delta x\quad\longrightarrow\quad\int_a^b f(x)\,\dd x \]
と書き、$f(x)$ の $a$ から $b$ までの\textbf{定積分}（ていせきぶん）といいます。$\Sigma$ が $\int$ に、はば $\Delta x$ が $\dd x$ に変わったと考えましょう。$a$ を\textbf{下端}、$b$ を\textbf{上端}といいます。\bkidx{ていせきぶん}{定積分}\bkidx{いんてぐらる}{$\int$（インテグラル）}

\noindent{\sffamily\bfseries\color{bkTheme}なぜ $\dd x$ を書くのか}\par
\begin{enumerate}[label={\bkcircled[bkTheme]{\arabic*}}, leftmargin=2.4\zw, itemsep=1pt, topsep=2pt]
\item $f(x)\,\dd x$ は「高さ×はば」、つまり細い長方形1本の面積です。$\dd x$ がないと、たすものが面積になりません。
\item 何を細かく分けたのかを表します。$\displaystyle\int v\,\dd t$ なら、時間 $t$ を細かく分けています。
\item 単位がわかります。$v$（m/s）に $\dd t$（s）をかけるので、$(\text{m/s})\times\text{s}=\text{m}$。答えは距離になります（5-9）。
\end{enumerate}

\noindent{\sffamily\bfseries\color{bkTheme}マイナスになるとき}\par
$f(x)$ がマイナスの所では、長方形の高さがマイナスなので、その面積もマイナスとして数えます（5-9のマイナスの面積と同じ）。

また、下端と上端を入れかえて $\displaystyle\int_1^0 x^2\,\dd x$ と書くと、1から0へ、左向きに分けていくことになります。はば $\Delta x$ は「あと${}-{}$まえ」（2-3）なので、4本に分けるなら $\dfrac{0-1}{4}=-\dfrac{1}{4}$ で、マイナスです。高さ×はばがどれもマイナスになるので、
\[ \int_1^0 x^2\,\dd x=-\frac{1}{3} \]
です。下端と上端を入れかえると、符号が反対になります（5-11でもう一度確かめます）。

\begin{seikaku}
\noindent\textbf{はばと高さの決め方}\par
ここでは、はばをすべて同じにし、高さに右はし（または左はし）の値を使いました。せいかくには、はばはそろっていなくてもよく、高さには小さな区間の中のどこの値を使ってもかまいません。グラフがつながっている（とぎれていない）関数なら、いちばん広いはばを0に近づけたとき、どの決め方でも和は同じ値に近づきます。それが定積分です。
\end{seikaku}

\bkkigocards{\bkkigocard{$\displaystyle\sum_{k=1}^{n}$}{シグマ}{$k$ を1から $n$ まで変えて、全部たす}{sum}{5-10}}
{\bkkigocard{$\displaystyle\int_a^b$}{インテグラル}{$a$ から $b$ まで、細かく分けてたす（定積分）}{integral}{5-10}}

\bksubsection{6}{0から $b$ までの面積}

\label{bk:5-10:b}④では、0から1までを求めました。0から2までなら、どうなるでしょう。同じ4本で、2つの図をならべます。

\begin{bkfigure}
\begin{tikzpicture}
\begin{axis}[bkplot, name=s1, width=56mm, height=46mm, xmin=0, xmax=1.15, ymin=0, ymax=1.15, xtick={0.25,0.5,0.75,1}, xticklabels={$\frac14$,$\frac24$,$\frac34$,1}, ytick={0.5,1}, xlabel={$x$}, ylabel={$y$}, title={\sffamily\small 0から1まで}]
  \foreach \k in {1,2,4} {\edef\tmp{\noexpand\fill[bkOrange!25, draw=bkOrange!70, line width=0.4pt] (axis cs:{(\k-1)/4},0) rectangle (axis cs:{\k/4},{(\k/4)^2});}\tmp}
  \fill[bkRed!40, draw=bkRed!80!black, line width=0.6pt] (axis cs:0.5,0) rectangle (axis cs:0.75,0.5625);
  \addplot[bkcurve, domain=0:1.05] {x^2};
\end{axis}
\begin{axis}[bkplot, at={(s1.east)}, anchor=west, xshift=16mm, width=56mm, height=46mm, xmin=0, xmax=2.3, ymin=0, ymax=4.6, xtick={0.5,1,1.5,2}, ytick={2,4}, xlabel={$x$}, ylabel={$y$}, title={\sffamily\small 0から2まで}]
  \foreach \k in {1,2,4} {\edef\tmp{\noexpand\fill[bkOrange!25, draw=bkOrange!70, line width=0.4pt] (axis cs:{(\k-1)/2},0) rectangle (axis cs:{\k/2},{(\k/2)^2});}\tmp}
  \fill[bkRed!40, draw=bkRed!80!black, line width=0.6pt] (axis cs:1,0) rectangle (axis cs:1.5,2.25);
  \addplot[bkcurve, domain=0:2.1] {x^2};
\end{axis}
\end{tikzpicture}
\bkfigcaption{2つの図は、形がまったく同じ。目もりの数字だけが、横は2倍、たては4倍になっている。}
\end{bkfigure}

赤い3本めの長方形をくらべます。

\begin{bkcenter}
\begin{bktable}{colspec={Q[l,m,wd=28mm] Q[c,m,wd=34mm] Q[c,m,wd=40mm]}, rowsep=2pt, row{1}={font=\sffamily\bfseries\small}}
赤い長方形 & 0から1まで & 0から2まで \\
はば & $1\div4=\dfrac{1}{4}$ & $2\div4=\dfrac{2}{4}$\quad（2倍） \\
右はし & $\dfrac{3}{4}$ & $\dfrac{2}{4}\times3=\dfrac{6}{4}$\quad（2倍） \\
高さ & $\left(\dfrac{3}{4}\right)^2=\dfrac{9}{16}$ & $\left(\dfrac{6}{4}\right)^2=\dfrac{36}{16}$\quad（4倍） \\
面積 & $\dfrac{9}{16}\times\dfrac{1}{4}=\dfrac{9}{64}$ & $\dfrac{36}{16}\times\dfrac{2}{4}=\dfrac{72}{64}$\quad（8倍） \\
\end{bktable}
\end{bkcenter}

はばが2倍になると、右はしの場所も2倍になります。高さは右はしの2乗なので、$2\times2=4$ 倍。面積は、はば2倍×高さ4倍で8倍です。ほかの長方形も同じなので、和も8倍、近づく先も8倍です。
\[ \int_0^2x^2\,\dd x=8\times\frac{1}{3}=\frac{8}{3} \]
8は $2^3$ です。0から $b$ までなら、はばが $b$ 倍、高さが $b^2$ 倍なので、面積は $b\times b^2=b^3$ 倍になります。

\begin{bkcenter}
\begin{bktable}{colspec={Q[l,m,wd=30mm] Q[c,m,wd=36mm] Q[c,m,wd=52mm]}, rowsep=2pt}
 & 0から1まで & 0から $b$ まで \\
はば & $\dfrac{1}{n}$ & $\dfrac{b}{n}$\quad（$b$ 倍） \\
$k$ 本めの右はし & $\dfrac{k}{n}$ & $\dfrac{b}{n}\times k=\dfrac{kb}{n}$\quad（$b$ 倍） \\
$k$ 本めの高さ & $\left(\dfrac{k}{n}\right)^2$ & $\left(\dfrac{kb}{n}\right)^2=b^2\left(\dfrac{k}{n}\right)^2$\quad（$b^2$ 倍） \\
長方形の面積 & $\left(\dfrac{k}{n}\right)^2\times\dfrac{1}{n}$ & $b\times b^2=b^3$ 倍 \\
\end{bktable}
\bktabcaption{はばも高さものびるので、どの長方形の面積も $b^3$ 倍になる。}
\end{bkcenter}
\[ \sum_{k=1}^{n}\left(\frac{kb}{n}\right)^2\times\frac{b}{n}=b^3\times\sum_{k=1}^{n}\left(\frac{k}{n}\right)^2\times\frac{1}{n}\ \longrightarrow\ b^3\times\frac{1}{3} \]
\noindent{\sffamily\bfseries\color{bkTheme}計算の結果}\quad $\displaystyle\int_0^b x^2\,\dd x=\frac{b^3}{3}$\par
\noindent $b=1$ なら $\dfrac{1}{3}$ で、④と合っています。$b=2$ なら $\dfrac{8}{3}$ で、上の計算と合っています。覚えなくても、「0から1までの $\dfrac{1}{3}$ を $b^3$ 倍」と考えれば、いつでも作れます。

\noindent{\sffamily\bfseries\color{bkTheme}0からでないとき}\par
0からでないときは、大きい面積から小さい面積をひきます。

\begin{bkfigure}
\begin{tikzpicture}
\pgfplotsset{mini/.style={bkplot, width=40mm, height=36mm, xmin=0, xmax=2.4, ymin=0, ymax=4.6, xtick={1,2}, xticklabels={$a$,$b$}, ytick=\empty}}
\begin{axis}[mini, name=m1]
  \addplot[fill=bkBlue!30, draw=none, domain=0:2] {x^2} \closedcycle;
  \addplot[bkcurve, line width=1.2pt, domain=0:2.1] {x^2};
\end{axis}
\node[font=\Large, anchor=west] (mi) at ([xshift=2mm]m1.east) {$-$};
\begin{axis}[mini, name=m2, at={(mi.east)}, anchor=west, xshift=2mm]
  \addplot[fill=bkGreen!35, draw=none, domain=0:1] {x^2} \closedcycle;
  \addplot[bkcurve, line width=1.2pt, domain=0:2.1] {x^2};
\end{axis}
\node[font=\Large, anchor=west] (eq) at ([xshift=2mm]m2.east) {$=$};
\begin{axis}[mini, name=m3, at={(eq.east)}, anchor=west, xshift=2mm]
  \addplot[fill=bkOrange!40, draw=none, domain=1:2] {x^2} \closedcycle;
  \addplot[bkcurve, line width=1.2pt, domain=0:2.1] {x^2};
\end{axis}
\end{tikzpicture}
\bkfigcaption{「0から $b$ まで」から「0から $a$ まで」をひくと、「$a$ から $b$ まで」が残る。}
\end{bkfigure}

\begin{align*}
\int_a^b x^2\,\dd x&=\int_0^b x^2\,\dd x-\int_0^a x^2\,\dd x\\
&=\frac{b^3}{3}-\frac{a^3}{3}
\end{align*}

\noindent{\sffamily\bfseries\color{bkTheme}グラフがまっすぐなとき}\par
グラフがまっすぐなら、5-9のように図形で求められます。$y=x$ の下は三角形、$y=1$ の下は長方形です。
\begin{align*}
\int_0^b x\,\dd x&=\frac{1}{2}\times b\times b=\frac{b^2}{2}&&\text{（三角形）}\\
\int_0^b 1\,\dd x&=1\times b=b&&\text{（長方形）}
\end{align*}
$y=x$ のときも、長方形の和で同じ答えになることを、あとの例題5で確かめます。

\begin{bkfigure}
\begin{tikzpicture}
\pgfplotsset{mini2/.style={bkplot, width=50mm, height=38mm, xmin=0, xmax=2.5, ymin=0, ymax=2.5, xtick={2}, xticklabels={$b$}, ytick=\empty, xlabel={$x$}, title style={xshift=3mm}}}
\begin{axis}[mini2, name=t1, title={\sffamily\small $y=x$ の下：三角形}]
  \addplot[fill=bkBlue!25, draw=none, domain=0:2] {x} \closedcycle;
  \addplot[bkcurve, line width=1.2pt, domain=0:2.3] {x};
  \node[font=\sffamily\footnotesize, anchor=west] at (axis cs:2.05,1.0) {高さ $b$};
\end{axis}
\begin{axis}[mini2, name=t2, at={(t1.east)}, anchor=west, xshift=22mm, title={\sffamily\small $y=1$ の下：長方形}]
  \addplot[fill=bkGreen!30, draw=none, domain=0:2] {1} \closedcycle;
  \addplot[bkcurve, line width=1.2pt, domain=0:2.3] {1};
  \node[font=\sffamily\footnotesize, anchor=south] at (axis cs:1,1.05) {高さ1};
\end{axis}
\end{tikzpicture}
\bkfigcaption{グラフがまっすぐなら、三角形（底辺 $b$、高さ $b$）や長方形（よこ $b$、たて1）の面積になる。}
\end{bkfigure}


\begin{senmonsho}[大学の教科書ふうの書き方]
\begin{bkhonbun}
区間 $[a,\,b]$ を $a=x_0<x_1<\cdots<x_n=b$ と分け、各区間 $[x_{k-1},\,x_k]$ から点 $\xi_k$ をとる。$f$ が連続ならば、和 $\displaystyle\sum_{k=1}^{n}f(\xi_k)\,(x_k-x_{k-1})$ は、分け方を細かくすると一定の値に収束する。これを $\displaystyle\int_a^b f(x)\,\dd x$ と書く。
\end{bkhonbun}
\bkyaku{$[a,\,b]$ は「$a$ から $b$ まで（両はしもふくむ）」。$\xi$（クシー）は、小さな区間の中から選んだ1点で、長方形の高さを決める場所。$x_k-x_{k-1}$ は $k$ 番めの長方形のはば $\Delta x$。この和を「リーマン和」とよぶ。}
\end{senmonsho}

\begin{taisetsu}
定積分\quad $\displaystyle\int_a^b f(x)\,\dd x$\\
＝ ① はば $\Delta x$ に分け、② 高さ $f(x)$×はばを1本ずつ計算し、③ $a$ から $b$ までたして、はばを0に近づけた先\\
＝ グラフと横軸のあいだの面積（横軸より下の部分はマイナス）\\
$\displaystyle\int_0^b x^2\,\dd x=\frac{b^3}{3}$\qquad $\displaystyle\int_0^b x\,\dd x=\frac{b^2}{2}$
\end{taisetsu}

\begin{kotoba}
\small
\bkword{$\Sigma$（シグマ）}{番号 $k$ を変えながら、全部たす記号。}
\bkword{定積分}{細かく分けてたした和の、はばを0に近づけた先。グラフと横軸のあいだの面積（下はマイナス）。}
\bkword{$\int$（インテグラル）}{定積分の記号。「たし合わせる」の S をのばした形。}
\bkword{下端・上端}{$\int_a^b$ の $a$（はじめ）と $b$（おわり）。}
\end{kotoba}

%% ---------------- 例題 ----------------
\bksection{reidai}{例題でたしかめよう}

\begin{reidai}[$\Sigma$ を計算する]
$\displaystyle\sum_{k=1}^{4}(2k^2-1)$ を計算しましょう。\par
\kangaekata $k$ に1、2、3、4を順に入れて、できた4つの数をたします。$2k^2$ は「$k$ を2乗してから2倍」です。\par
\shiki $k=1$：$2\times1^2-1=1$\qquad $k=2$：$2\times2^2-1=7$\par
\hspace*{2.6\zw}$k=3$：$2\times3^2-1=17$\qquad $k=4$：$2\times4^2-1=31$\par
\hspace*{2.6\zw}$1+7+17+31=56$\par
\kotae 56
\end{reidai}

\begin{reidai}[長方形の和を $\Sigma$ で書く]
$y=x^2$ のグラフの、$x=0$ から1までの下の面積を、右はしの長方形5本で近似します。和を $\Sigma$ で書き、計算しましょう。\par
\kangaekata ①の段1〜段4と同じです。はばは $1\div5=\dfrac{1}{5}$。$k$ 本めの右はしは、$\dfrac{1}{5}$ を $k$ こ進んだ $\dfrac{k}{5}$。高さは、そこでの $y=x^2$ の値 $\left(\dfrac{k}{5}\right)^2$ です。\par
\shiki $\displaystyle\sum_{k=1}^{5}\left(\frac{k}{5}\right)^2\times\frac{1}{5}=\sum_{k=1}^{5}\frac{k^2}{125}$\quad（$5^2\times5=125$）\par
\hspace*{2.6\zw}$\displaystyle=\frac{1+4+9+16+25}{125}=\frac{55}{125}=0.44$\par
\kotae 0.44（右はしなので、ほんとうの $\frac{1}{3}$ より大きい）
\end{reidai}

\noindent\begin{minipage}[t]{0.485\linewidth}\vspace{0pt}\raggedright
\begin{jibun}
$\displaystyle\sum_{k=1}^{5}k^2$ と $\displaystyle\sum_{k=1}^{3}(k+2)$ を計算しましょう。
\bkjans{55、12}
\end{jibun}
\end{minipage}\hfill
\begin{minipage}[t]{0.485\linewidth}\vspace{0pt}\raggedright
\begin{jibun}
同じ面積を、右はしの長方形3本で近似します。和を $\Sigma$ で書き、計算しましょう。
\bkjans{$\displaystyle\sum_{k=1}^{3}\left(\frac{k}{3}\right)^2\times\frac{1}{3}=\frac{1+4+9}{27}=\frac{14}{27}\approx0.52$}
\end{jibun}
\end{minipage}

\begin{reidai}[長方形の和で近似する]
$y=x$ のグラフの、$x=0$ から2までの下の面積を、はば0.5の右はしの長方形4本で近似しましょう。ほんとうの面積とくらべましょう。\par
\kangaekata 0から2までを、はば0.5で分けると $2\div0.5=4$ 本。右はしは0.5、1、1.5、2です。グラフは $y=x$ なので、高さは右はしの $x$ そのもの（2乗しない）です。\par
\shiki 高さは0.5、1、1.5、2。はばはどれも0.5なので、高さを先にたしてから0.5をかけます。\par
\hspace*{2.6\zw}$(0.5+1+1.5+2)\times0.5=5\times0.5=2.5$\par
\kotae 近似2.5（ほんとうは三角形で $\frac{1}{2}\times2\times2=2$。はみ出した分だけ大きい）
\end{reidai}

\begin{reidai}[$x^2$ の定積分]
$\displaystyle\int_0^3x^2\,\dd x$ を求めましょう。\par
\kangaekata $y=x^2$ のグラフの、0から3までの下の面積です。⑥で、0から $b$ までの面積は $\dfrac{b^3}{3}$ でした（0から1までの $\dfrac{1}{3}$ の $b^3$ 倍）。いまは上端が3なので、$b=3$ を入れます。\par
\shiki $\displaystyle\int_0^3x^2\,\dd x=\frac{3^3}{3}=\frac{27}{3}=9$\par
\bkstep[bkGreen]{\R{確}{たし}かめ}はば1の長方形なら、左はし $0+1+4=5$、右はし $1+4+9=14$。9は、そのあいだにあります。\par
\kotae 9
\end{reidai}

\begin{column}{1から100までの和}
240年ほど前のドイツで、小学生だったガウスは、先生に「1から100までを全部たしなさい」と言われました。ガウスはすぐに、5050と答えたそうです。1と100、2と99、3と98、……のように、はしから組にすると、どの組も和は101。組は $100\div2=50$ 組あるので、$101\times50=5050$。③の「階段を2つ組み合わせる」と同じ考え方で、式にすると $\dfrac{100\times101}{2}=5050$ です。ガウスは大人になって、数学だけでなく、電気や磁気の研究でも大きな仕事をしました。磁石のまわりの磁場の強さ（磁束密度）を表す単位「ガウス」は、彼の名前からとったものです（いまの国際単位ではテスラ（T）を使い、1\,T＝10000ガウスです）。
\end{column}

\noindent\begin{minipage}[t]{0.485\linewidth}\vspace{0pt}\raggedright
\begin{jibun}
$y=x^2$ のグラフの、$x=0$ から2までの下の面積を、はば0.5の右はしの長方形4本で近似しましょう。
\bkjans{$(0.25+1+2.25+4)\times0.5=3.75$（ほんとうは $\frac{8}{3}=2.67$）}
\end{jibun}
\end{minipage}\hfill
\begin{minipage}[t]{0.485\linewidth}\vspace{0pt}\raggedright
\begin{jibun}
$\displaystyle\int_0^2x^2\,\dd x$ と $\displaystyle\int_0^3x^2\,\dd x$ を求めましょう。
\bkjans{$\dfrac{2^3}{3}=\dfrac{8}{3}$、$\dfrac{3^3}{3}=9$}
\end{jibun}
\end{minipage}

\begin{reidai}[$y=x$ の下の面積を、長方形の和で求める]
$y=x$ のグラフの、$x=0$ から2までの下の面積を、右はしの長方形 $n$ 本の和で表し、$n$ をかぎりなく大きくして求めましょう。三角形の面積とくらべましょう。\par
\noindent\begin{minipage}[t]{0.66\linewidth}\raggedright
\kangaekata 段1〜段4で考えます。\par
\hspace*{2.6\zw}はば：長さ2を $n$ 等分して $\dfrac{2}{n}$\par
\hspace*{2.6\zw}右はし：$\dfrac{2}{n}$ を $k$ こ進んで $\dfrac{2k}{n}$\par
\hspace*{2.6\zw}高さ：$y=x$ なので、右はしと同じ $\dfrac{2k}{n}$\par
\hspace*{2.6\zw}面積：$\dfrac{2k}{n}\times\dfrac{2}{n}=\dfrac{4k}{n^2}$
\end{minipage}\hfill
\begin{minipage}[t]{0.32\linewidth}\centering\vspace{0pt}
\begin{tikzpicture}[x=15mm, y=15mm]
  \foreach \k in {1,...,4} {{\fill[bkOrange!30] ({(\k-1)/2},0) rectangle ({\k/2},{\k/2});
    \draw[bkOrange!80!black, line width=0.4pt] ({(\k-1)/2},0) rectangle ({\k/2},{\k/2});}}
  \draw[bkBlue!80!black, line width=1.1pt] (0,0) -- (2.1,2.1);
  \draw[bkInk, line width=0.6pt, -{Stealth[length=3pt]}] (0,0) -- (2.3,0) node[right, font=\footnotesize, inner sep=1pt] {$x$};
  \draw[bkInk, line width=0.6pt, -{Stealth[length=3pt]}] (0,0) -- (0,2.3) node[above, font=\footnotesize, inner sep=1pt] {$y$};
  \node[font=\footnotesize, anchor=north, inner sep=1pt] at (2,-0.05) {2};
  \node[font=\footnotesize, anchor=north, inner sep=1pt] at (1,-0.05) {1};
  \node[font=\footnotesize, text=bkBlue!80!black, anchor=south east, inner sep=1pt] at (1.3,1.55) {$y=x$};
\end{tikzpicture}\par
{\footnotesize\sffamily $n=4$ のとき}
\end{minipage}\par
\shiki $\displaystyle S_n=\sum_{k=1}^{n}\frac{4k}{n^2}=\frac{4}{n^2}\times(1+2+\cdots+n)$\par
\hspace*{2.6\zw}$\displaystyle\phantom{S_n}=\frac{4}{n^2}\times\frac{n(n+1)}{2}=\frac{2(n+1)}{n}=2\Bigl(1+\frac{1}{n}\Bigr)$\quad（③の和）\par
\hspace*{2.6\zw}$n=4$ なら $2\times1.25=2.5$（例題3と同じ）\par
\hspace*{2.6\zw}$n\to\infty$ で $\dfrac{1}{n}\to0$ なので、$2\times1=2$\par
\bkstep[bkGreen]{\R{確}{たし}かめ}三角形の面積は $\dfrac{1}{2}\times2\times2=2$。一致します。\par
\kotae 2
\end{reidai}

\begin{reidai}[曲がったグラフの距離]
速度が $v=3t^2$（m/s）の物体が、0秒から2秒までに進む距離を求めましょう。\par
\kangaekata 距離は $v$-$t$ グラフの面積なので、$\displaystyle\int_0^2 3t^2\,\dd t$ です。$v=3t^2$ のグラフは、$v=t^2$ をたてに3倍にのばした形です。長方形のはばはそのままで、高さがどれも3倍なので、面積の和も3倍です（下の図）。横軸の文字が $x$ から $t$ に変わっても形は同じなので、⑥の $\dfrac{b^3}{3}$ が使えます（$b=2$）。\par
\begin{center}
\begin{tikzpicture}[x=1mm, y=1mm]
  % t=1.25〜1.5 の細い長方形：灰（v=t^2 の高さ）と青（v=3t^2 の高さ）
  \fill[bkBlue!25] (25,0) rectangle (30,{3*1.5*1.5*2.6});
  \fill[bkGray!45] (25,0) rectangle (30,{1.5*1.5*2.6});
  \draw[bkGray!80!black, line width=0.4pt] (25,0) rectangle (30,{1.5*1.5*2.6});
  \draw[bkBlue!80!black, line width=0.4pt] (25,0) rectangle (30,{3*1.5*1.5*2.6});
  \draw[bkGray!80!black, line width=1.1pt] plot[domain=0:2, samples=30] ({20*\x},{2.6*\x*\x});
  \draw[bkBlue!80!black, line width=1.3pt] plot[domain=0:1.62, samples=30] ({20*\x},{7.8*\x*\x});
  \draw[bkInk, line width=0.6pt, -{Stealth[length=4pt]}] (0,0) -- (46,0) node[right, font=\sffamily\small, inner sep=1pt] {$t$};
  \draw[bkInk, line width=0.6pt, -{Stealth[length=4pt]}] (0,0) -- (0,23) node[above, font=\sffamily\small, inner sep=1pt] {$v$};
  \node[font=\sffamily\small, anchor=north] at (40,-0.5) {2};
  \node[font=\sffamily\small, anchor=north] at (30,-0.5) {1.5};
  \node[font=\sffamily\small, text=bkBlue!80!black, anchor=east] at (24,20) {$v=3t^2$};
  \node[font=\sffamily\small, text=bkGray!80!black, anchor=west] at (42,10.4) {$v=t^2$};
  % 高さの比べ（右に寸法線）
  \draw[bkInk, line width=0.5pt, <->] (32,0) -- (32,{1.5*1.5*2.6});
  \draw[bkBlue!80!black, line width=0.5pt, <->] (34.5,0) -- (34.5,{3*1.5*1.5*2.6});
  \node[font=\sffamily\footnotesize, anchor=west, inner sep=0pt] at (36,2.9) {1つ分};
  \node[font=\sffamily\bfseries\footnotesize, text=bkBlue!80!black, anchor=west, inner sep=0pt] at (36,17.5) {高さ3倍};
\end{tikzpicture}
\end{center}
\shiki $\displaystyle\begin{aligned}[t]\int_0^2 3t^2\,\dd t&=3\times\int_0^2 t^2\,\dd t\\&=3\times\frac{2^3}{3}=3\times\frac{8}{3}=8\end{aligned}$\par
\bkstep[bkGreen]{\R{確}{たし}かめ}はば1の長方形なら、左はし $0+3=3$、右はし $3+12=15$。8は、そのあいだにあります。単位は $(\text{m/s})\times\text{s}=\text{m}$ です。\par
\kotae \qty{8}{m}
\end{reidai}

\noindent\begin{minipage}[t]{0.485\linewidth}\vspace{0pt}\raggedright
\begin{jibun}
$y=x$ のグラフの、$x=0$ から3までの下の面積を、右はしの長方形 $n$ 本の和で表し、$n$ をかぎりなく大きくして求めましょう。三角形の面積とくらべましょう。
\bkjans{はば $\dfrac{3}{n}$、高さ $\dfrac{3k}{n}$。$\displaystyle\sum_{k=1}^{n}\frac{9k}{n^2}=\frac{9}{n^2}\times\frac{n(n+1)}{2}=\frac{9}{2}\Bigl(1+\frac{1}{n}\Bigr)\to\frac{9}{2}$。三角形も $\frac{1}{2}\times3\times3=\frac{9}{2}$}
\end{jibun}
\end{minipage}\hfill
\begin{minipage}[t]{0.485\linewidth}\vspace{0pt}\raggedright
\begin{jibun}
速度が $v=6t^2$（m/s）の物体が、0秒から3秒までに進む距離を求めましょう。
\bkjans{$6\times\dfrac{27}{3}=\qty{54}{m}$}
\end{jibun}
\end{minipage}

\begin{reidai}[とちゅうからの面積は、差で求める]
速度が $v=3t^2$（m/s）の物体が、1秒から3秒までに進む距離を求めましょう。\par
\kangaekata ⑥の「0からでないとき」と同じです。0秒から3秒までの面積から、0秒から1秒までの面積をひきます。高さが3倍なので、どちらも $t^2$ の面積の3倍です。\par
\shiki 0〜3秒：$\displaystyle3\times\frac{3^3}{3}=27$\qquad 0〜1秒：$\displaystyle3\times\frac{1^3}{3}=1$\par
\hspace*{2.6\zw}$27-1=26$\par
\kotae \qty{26}{m}
\end{reidai}

\begin{reidai}[台形で近似する]
$y=x^2$ の、$x=0$ から2までの下の面積を、はば0.5の台形4つで近似しましょう。\par
\kangaekata グラフ上の2点を結ぶ台形の面積は、（左の高さ＋右の高さ）÷2×はば（4-5）。左の高さだけで計算したのが左はしの長方形、右の高さだけで計算したのが右はしの長方形なので、台形の和は、2つの長方形の和の平均です（②の表の「2つの平均」）。\par
\shiki 左はし $(0+0.25+1+2.25)\times0.5=1.75$\par
\hspace*{2.6\zw}右はし $(0.25+1+2.25+4)\times0.5=3.75$\par
\hspace*{2.6\zw}平均 $\dfrac{1.75+3.75}{2}=2.75$\par
\kotae 2.75（ほんとうは $\frac{8}{3}=2.666\cdots$。長方形の3.75より、ずっと近い）
\end{reidai}

\noindent\begin{minipage}[t]{0.485\linewidth}\vspace{0pt}\raggedright
\begin{jibun}
速度が $v=t^2$（m/s）の物体が、2秒から4秒までに進む距離を求めましょう。
\bkjans{$\dfrac{64}{3}-\dfrac{8}{3}=\dfrac{56}{3}\approx\qty{18.7}{m}$}
\end{jibun}
\end{minipage}\hfill
\begin{minipage}[t]{0.485\linewidth}\vspace{0pt}\raggedright
\begin{jibun}
$y=x^2$ の、$x=0$ から1までの下の面積を、はば0.5の台形2つで近似しましょう。
\bkjans{左はし0.125、右はし0.625の平均で0.375（ほんとうは $\frac{1}{3}$）}
\end{jibun}
\end{minipage}

\begin{reidai}[うすい板に分けて体積を求める]
底面の半径 $r$、高さ $h$ の円すいの体積を求めましょう。\par
\kangaekata 頂点から $y$ の所で、底面と平行にうすく切ります（いちばん上の頂点で $y=0$、底面で $y=h$）。断面は円です。断面の半径は、頂点からの長さ $y$ に比例します（相似、4-3）。$y=h$ で半径 $r$ なので、$y$ の所では $r\times\dfrac{y}{h}=\dfrac{r}{h}y$。厚さ $\Delta y$ の円板（うすい円柱）を、上から下までたし合わせます。\par
\noindent\begin{minipage}[c]{0.62\linewidth}\raggedright
\shiki 円板の体積＝底面積×厚さ（4-5）\par
\hspace*{2.6\zw}$\pi\Bigl(\dfrac{r}{h}y\Bigr)^2\Delta y=\dfrac{\pi r^2}{h^2}\,y^2\,\Delta y$\par
\hspace*{2.6\zw}$\dfrac{\pi r^2}{h^2}$ は $y$ によらない数なので、外に出す\par
\hspace*{2.6\zw}$\displaystyle\int_0^h\frac{\pi r^2}{h^2}y^2\,\dd y=\frac{\pi r^2}{h^2}\times\frac{h^3}{3}=\frac{1}{3}\pi r^2h$\par
\kotae $\dfrac{1}{3}\pi r^2h$（4-5で学んだ公式と同じ）
\end{minipage}\hfill
\begin{minipage}[c]{0.36\linewidth}\centering
\begin{tikzpicture}[x=1mm, y=1mm]
  % 第6回：頂点 (0,32)、底面の中心 (0,2)、h=30、r=11。y は頂点から下向き。円板は y=18、厚さ Δy=2.4
  \def\yy{18}\def\dy{2.4}
  \pgfmathsetmacro\tp{32-\yy}\pgfmathsetmacro\bt{32-\yy-\dy}\pgfmathsetmacro\rr{11*\yy/30}
  \draw[bkInk, line width=0.5pt, dash pattern=on 1.5pt off 1pt] (11,2) arc[start angle=0, end angle=180, x radius=11, y radius=2.75];
  \fill[bkBlue!40] (-\rr,\tp) -- (-\rr,\bt) arc[start angle=180, end angle=360, x radius=\rr, y radius=\rr/4] -- (\rr,\tp) -- cycle;
  \fill[bkBlue!18] (0,\tp) ellipse [x radius=\rr, y radius=\rr/4];
  \draw[bkBlue!70!black, line width=0.6pt] (0,\tp) ellipse [x radius=\rr, y radius=\rr/4];
  \draw[bkBlue!70!black, line width=0.6pt] (-\rr,\tp) -- (-\rr,\bt) arc[start angle=180, end angle=360, x radius=\rr, y radius=\rr/4] -- (\rr,\tp);
  \draw[bkBlue!70!black, line width=0.6pt] (0,\tp) -- (\rr,\tp);
  \draw[bkInk, line width=0.7pt] (-11,2) arc[start angle=180, end angle=360, x radius=11, y radius=2.75];
  \draw[bkInk, line width=0.7pt] (-11,2) -- (0,32) -- (11,2);
  \draw[bkInk, line width=0.6pt] (0,2) -- (9.2,0.7);
  \fill[bkInk] (0,2) circle (0.5);
  \node[font=\sffamily\footnotesize, anchor=south, inner sep=1pt] at (4.6,1.4) {$r$};
  \draw[bkGray, line width=0.4pt] (-1,32) -- (-18.5,32) (-\rr-0.8,\tp) -- (-14.5,\tp) (-11.8,2) -- (-18.5,2);
  \draw[bkGray, line width=0.5pt, <->] (-13.5,32) -- (-13.5,\tp);
  \node[font=\sffamily\footnotesize, anchor=west] at (-13.1,{(32+\tp)/2}) {$y$};
  \draw[bkGray, line width=0.5pt, <->] (-17,32) -- (-17,2);
  \node[font=\sffamily\footnotesize, anchor=east] at (-17.3,17) {$h$};
  \draw[bkGray, line width=0.4pt] (\rr+0.8,\tp) -- (\rr+4.8,\tp) (\rr+0.8,\bt) -- (\rr+4.8,\bt);
  \draw[bkGray, line width=0.5pt, -{Stealth[length=3pt]}] ({\rr+3.8},{\tp+3}) -- ({\rr+3.8},\tp);
  \draw[bkGray, line width=0.5pt, -{Stealth[length=3pt]}] ({\rr+3.8},{\bt-3}) -- ({\rr+3.8},\bt);
  \node[font=\sffamily\footnotesize, anchor=west] at ({\rr+4.3},{(\tp+\bt)/2}) {$\Delta y$};
  \draw[bkBlue!70!black, line width=0.4pt] ({\rr/2},{\tp+0.2}) -- (6.6,23.4);
  \node[font=\sffamily\footnotesize, text=bkBlue!70!black, anchor=south west, align=left, inner sep=1pt] at (6.2,23.4) {半径\\$\dfrac{r}{h}y$};
\end{tikzpicture}
\end{minipage}
\end{reidai}

\begin{jibun}
底面の半径\qty{3}{cm}、高さ\qty{6}{cm}の円すいの体積を、\bkreidairef{9}と同じように積分で求めましょう。
\bkjans{$\displaystyle\int_0^6\frac{\pi\times9}{36}y^2\,\dd y=\frac{\pi}{4}\times72=18\pi\ \text{cm}^3$}
\end{jibun}

\begin{reidai}[左はしの長方形の和で、極限を求める]
$y=x^2$ のグラフの、$x=0$ から1までの下の面積を、左はしの長方形 $n$ 本の和で表し、$n$ をかぎりなく大きくして求めましょう。\par
\kangaekata 左はしの長方形では、$k$ 本めの高さは、左はし $\dfrac{k-1}{n}$ での値です。番号を1つずらして、0から $n-1$ までの $k$ で書くと、高さは $\left(\dfrac{k}{n}\right)^2$ です（まちがえやすいところ「$\Sigma$ の範囲」）。\par
\begin{center}
\begin{tikzpicture}[x=64mm, y=34mm]
  \foreach \k in {1,2,3} {
    \fill[bkBlue!25] ({\k/4},0) rectangle ({(\k+1)/4},{(\k/4)^2});
    \draw[bkBlue!70!black, line width=0.5pt] ({\k/4},0) rectangle ({(\k+1)/4},{(\k/4)^2});
}
  \draw[bkBlue!80!black, line width=1.2pt] plot[domain=0:1.04, samples=40] (\x,{\x*\x});
  \draw[bkInk, line width=0.6pt, -{Stealth[length=3pt]}] (0,0) -- (1.1,0) node[right, font=\footnotesize, inner sep=1pt] {$x$};
  \draw[bkInk, line width=0.6pt, -{Stealth[length=3pt]}] (0,0) -- (0,1.12) node[above, font=\footnotesize, inner sep=1pt] {$y$};
  \foreach \k in {0,1,2,3} {\node[font=\sffamily\footnotesize, anchor=north, inner sep=1pt] at ({(\k+0.5)/4},-0.04) {$k=\k$};
    \node[font=\footnotesize, anchor=north, inner sep=1pt] at ({(\k+0.5)/4},-0.2) {$\left(\frac{\k}{4}\right)^2$};}
  \node[font=\sffamily\footnotesize, anchor=east, inner sep=1pt] at (-0.02,-0.1) {番号};
  \node[font=\sffamily\footnotesize, anchor=east, inner sep=1pt] at (-0.02,-0.3) {高さ};
\end{tikzpicture}\par
{\footnotesize\sffamily $n=4$ の左はしの長方形。番号 $k$ を0から3にすると、高さは $\left(\frac{k}{4}\right)^2$}
\end{center}
\shiki $\displaystyle\sum_{k=0}^{n-1}\left(\frac{k}{n}\right)^2\times\frac{1}{n}=\sum_{k=1}^{n}\left(\frac{k}{n}\right)^2\times\frac{1}{n}-\left(\frac{n}{n}\right)^2\times\frac{1}{n}$\par
\hspace*{2.6\zw}（$k=0$ の項は0。右はしの和から、いちばん右の1本をひいた形）\par
\hspace*{2.6\zw}$\displaystyle=\frac{1}{6}\Bigl(1+\frac{1}{n}\Bigr)\Bigl(2+\frac{1}{n}\Bigr)-\frac{1}{n}$\quad（④の $S_n$）\par
\hspace*{2.6\zw}$n=4$ なら $0.46875-0.25=0.21875$（①と同じ）\par
\hspace*{2.6\zw}$n\to\infty$ で $\dfrac{1}{3}-0=\dfrac{1}{3}$\par
\kotae $\dfrac{1}{3}$（右はしの和と同じ値に近づく）
\end{reidai}

\begin{jibun}
$y=x$ のグラフの、$x=0$ から2までの下の面積を、左はしの長方形 $n$ 本の和で表し、$n$ をかぎりなく大きくして求めましょう。
\bkjans{右はしの和 $2\Bigl(1+\dfrac{1}{n}\Bigr)$（例題5）から、いちばん右の1本 $2\times\dfrac{2}{n}$ をひいて $2\Bigl(1-\dfrac{1}{n}\Bigr)\to2$}
\end{jibun}

\begin{column}{まん中の高さを使うと}
長方形の高さは、はばのまん中の値で決めてもかまいません（⑤の「せいかくに言うと」）。$y=x^2$ の0から1までを4本に分けると、まん中は $\frac{1}{8}$、$\frac{3}{8}$、$\frac{5}{8}$、$\frac{7}{8}$ なので、和は $\dfrac{1^2+3^2+5^2+7^2}{8^2}\times\dfrac{1}{4}=\dfrac{84}{256}=0.328125$。4本だけで、$\frac{1}{3}$ にかなり近くなります。はみ出す部分とすき間の部分が、ほぼ打ち消しあうからです。\par\vspace{2pt}

\noindent\begin{minipage}[c]{0.42\linewidth}\centering
\begin{tikzpicture}[x=1mm, y=1mm]
  % 1本の長方形を大きくした図。まん中の高さ h_m。左は すき間、右は はみ出し
  \fill[bkGreen!25] (0,0) rectangle (30,16);
  \fill[bkBlue!35] (0,16) -- (0,9.5) .. controls (8,11.6) and (12,13.6) .. (15,16) -- cycle;
  \fill[bkOrange!50] (15,16) .. controls (19,19) and (24,21.5) .. (30,24.5) -- (30,16) -- cycle;
  \draw[bkGreen!60!black, line width=0.5pt] (0,0) rectangle (30,16);
  \draw[bkBlue!80!black, line width=1.2pt] (-3,8.6) .. controls (-1,9) .. (0,9.5) .. controls (8,11.6) and (12,13.6) .. (15,16) .. controls (19,19) and (24,21.5) .. (30,24.5) -- (32,25.8);
  \draw[bkInk, line width=0.6pt] (-4,0) -- (34,0);
  \draw[bkGray, line width=0.4pt, dash pattern=on 1.2pt off 1pt] (15,0) -- (15,16);
  \node[font=\footnotesize, anchor=north, inner sep=1pt] at (15,-0.8) {まん中};
\end{tikzpicture}
\end{minipage}\hfill
\begin{minipage}[c]{0.55\linewidth}\raggedright
左の図は、1本の長方形を大きくしたものです。まん中より右では、グラフが長方形の上にはみ出しています（だいだい色）。左では、長方形の中にグラフより上の部分があり、余分にふくんでいます（青）。だいだい色と青の大きさがほぼ同じなので、打ち消しあって、和がほんとうの面積に近くなります。
\end{minipage}\par\vspace{2pt}
\noindent\begin{minipage}[c]{0.3\linewidth}\centering
\begin{tikzpicture}
\begin{axis}[bkplot, width=36mm, height=34mm, xmin=0, xmax=1.15, ymin=0, ymax=1.15, xtick={0.5,1}, ytick={0.5,1}, xlabel={$x$}, ylabel={$y$}]
  \foreach \k in {1,...,4} {\edef\tmp{\noexpand\fill[bkGreen!30, draw=bkGreen!70!black, line width=0.4pt] (axis cs:{(\k-1)/4},0) rectangle (axis cs:{\k/4},{((2*\k-1)/8)^2});}\tmp}
  \addplot[bkcurve, line width=1.2pt, domain=0:1.05] {x^2};
\end{axis}
\end{tikzpicture}
\end{minipage}\hfill
\begin{minipage}[c]{0.68\linewidth}\centering
\begin{bktable}{colspec={Q[c,m,wd=7mm] *{4}{Q[c,m,wd=15mm]}}, colsep=2pt, rowsep=1pt, cells={font=\footnotesize}, row{1}={font=\sffamily\bfseries\small}}
$n$ & 左はし & 右はし & 台形 & まん中 \\
4 & 0.21875 & 0.46875 & 0.34375 & 0.328125 \\
10 & 0.285 & 0.385 & 0.335 & 0.3325 \\
100 & 0.32835 & 0.33835 & 0.33335 & 0.333325 \\
\end{bktable}
\end{minipage}\par\vspace{2pt}
どの決め方でも、$n$ を大きくすると $\frac{1}{3}$ に近づきます。5-16では、まん中の高さを使って電気の量を計算します。
\end{column}

%% ---------------- まちがえやすいところ・物理工学 ----------------
\bksection{tsumazuki}{まちがえやすいところをチェック}

\begin{tsumazuki}{長方形の高さをまちがえる}
$y=x^2$ の0から1までを $n=4$ 本に分けたときの、$k$ 本めの高さ\par
\bkng{$k^2$（1、4、9、16）}
\bkng{$\dfrac{k}{4}$（2乗をわすれる）}
\bkok{$\left(\dfrac{k}{4}\right)^2$（$\frac{1}{16}$、$\frac{4}{16}$、$\frac{9}{16}$、$\frac{16}{16}$）}
$k$ は「何本め」という番号で、場所ではありません。右はしの場所は $\dfrac{k}{4}$ です。$k^2$ にすると、高さ16の長方形ができます。グラフは高さ1までしかないので、おかしいとわかります。高さは、グラフの式 $y=x^2$ に場所 $\dfrac{k}{4}$ を入れた値です。
\end{tsumazuki}

\begin{tsumazuki}{$b$ 倍にするのを、はばだけにする}
\bkng{$\displaystyle\int_0^2x^2\,\dd x=2\times\frac{1}{3}=\frac{2}{3}$}
\bkok{$\displaystyle\int_0^2x^2\,\dd x=2^3\times\frac{1}{3}=\frac{8}{3}$}
0から2までにすると、はばが2倍になるだけでなく、右はしが2倍になるので、高さも $2^2=4$ 倍になります。
\end{tsumazuki}

\begin{tsumazuki}{$\Sigma$ で、$k$ の範囲をまちがえる}
$\displaystyle\sum_{k=1}^{n}$ は $k=1$ から $n$ までの $n$ 個、$\displaystyle\sum_{k=0}^{n-1}$ は $k=0$ から $n-1$ までの $n$ 個です。左はしの長方形の和は $\displaystyle\sum_{k=0}^{n-1}\left(\frac{k}{n}\right)^2\frac{1}{n}$ で、右はしとは範囲が1つずれています。
\end{tsumazuki}

\begin{tsumazuki}{$\dd x$ をわすれる}
$\displaystyle\int_0^1x^2$ だけでは、何で細かく分けるのかがわかりません。$\dd x$ は「$x$ の方向に細く分けたはば」を表す大事な部分です。$\displaystyle\int v\,\dd t$ なら、時間で分けています。
\end{tsumazuki}

\begin{tsumazuki}{下端と上端を入れかえても同じだと思う}
\bkng{$\displaystyle\int_1^0x^2\,\dd x=\frac{1}{3}$}
\bkok{$\displaystyle\int_1^0x^2\,\dd x=-\frac{1}{3}$}
1から0へ分けていくと、はば $\Delta x$ が「あと${}-{}$まえ」でマイナスになります（⑤）。だから、符号が反対になります。
\end{tsumazuki}

\begin{butsuri}
\begin{itemize}[leftmargin=1.2\zw, itemsep=2pt, topsep=0pt, label={\color{bkPurple}●}]
\item \textbf{コンピュータの積分}\par
コンピュータは、式で積分できないときも、細い長方形や台形に分けてたし合わせて計算します。車の燃費計は、短い時間ごとの「燃料の流れる速さ×時間」をたして、使った燃料を求めています。
\item \textbf{CTスキャン}\par
CTスキャンでは、体をうすく輪切りにした断面の画像が得られます。断面の面積×厚さをたし合わせると、臓器などの体積がわかります（\bkreidairef{9}の円すいと同じ考え方）。
\item \textbf{面積計}\par
地図の上の形のまわりをなぞるだけで面積がわかる「プラニメーター」という道具は、積分の考え方で動いています。
\end{itemize}
\end{butsuri}


%% ---------------- 章末問題 ----------------
\bkendproblems
\bkproblemgroup{基本}{3}

\bkproblem 次を計算しましょう。\quad（1）$\displaystyle\sum_{k=1}^{6}k$\quad（2）$\displaystyle\sum_{k=1}^{4}3k$\quad（3）$\displaystyle\sum_{k=1}^{3}k^3$
\bkpans{（1）\nobreak 21\quad（2）\nobreak 30\quad（3）\nobreak 36}

\bkproblem $y=x^2$ のグラフの、$x=0$ から1までの下の面積を、右はしの長方形6本で近似します。和を $\Sigma$ で書き、計算しましょう（小数第3位まで）。
\bkpans{$\displaystyle\sum_{k=1}^{6}\left(\frac{k}{6}\right)^2\times\frac{1}{6}=\frac{91}{216}\approx0.421$}

\bkproblem $y=x^2$ のグラフの、$x=0$ から1までの下の面積を、はば0.5の左はしの長方形2つと、右はしの長方形2つで近似しましょう。
\bkpans{左はし：$(0+0.25)\times0.5=0.125$、右はし：$(0.25+1)\times0.5=0.625$}

\bkproblem 次の定積分を求めましょう。\quad（1）$\displaystyle\int_0^4x^2\,\dd x$\quad（2）$\displaystyle\int_0^5 x\,\dd x$
\bkpans{（1）\nobreak $\dfrac{64}{3}$\quad（2）\nobreak $\dfrac{25}{2}$}

\bkproblem 速度が $v=t^2$（m/s）の物体が、0秒から6秒までに進む距離を求めましょう。
\bkpans{$\dfrac{216}{3}=\qty{72}{m}$}

\bkproblem $\displaystyle\int_0^3x^2\,\dd x-\int_0^2x^2\,\dd x$ を計算しましょう。これは、どこの面積ですか。
\bkpans{$9-\dfrac{8}{3}=\dfrac{19}{3}$。$y=x^2$ の下の、$x=2$ から3までの面積}

\bkproblem 速度が $v=3t^2$（m/s）の物体が、1秒から2秒までに進む距離を求めましょう。
\bkpans{$8-1=\qty{7}{m}$}

\bkproblemgroup{チャレンジ}{3}

\bkproblem $y=x$ のグラフの下の面積を、長方形 $n$ 本の和で表し、$n$ をかぎりなく大きくして求めましょう。\quad（1）\nobreak $x=0$ から4まで、右はしの長方形\quad（2）\nobreak $x=0$ から3まで、左はしの長方形
\bkpans{（1）\nobreak $\displaystyle\sum_{k=1}^{n}\frac{4k}{n}\times\frac{4}{n}=\frac{16}{n^2}\times\frac{n(n+1)}{2}=8\Bigl(1+\frac{1}{n}\Bigr)\to8$\quad（2）\nobreak 右はしの和 $\dfrac{9}{2}\Bigl(1+\dfrac{1}{n}\Bigr)$ から、いちばん右の1本 $3\times\dfrac{3}{n}$ をひいて $\dfrac{9}{2}\Bigl(1-\dfrac{1}{n}\Bigr)\to\dfrac{9}{2}$}

\bkproblem 高さ $h$、底面が1辺 $a$ の正方形の四角すいの体積が $\dfrac{1}{3}a^2h$ になることを、\bkreidairef{9}と同じように積分で示しましょう。頂点から $y$ の所の断面は、1辺 $\dfrac{a}{h}y$ の正方形です（4-3）。
\bkpans{体積 $=\displaystyle\int_0^h\frac{a^2}{h^2}y^2\,\dd y=\frac{a^2}{h^2}\times\frac{h^3}{3}=\frac{1}{3}a^2h$}

\bkproblem $y=x^2$ の、$x=0$ から3までの下の面積を、はば1の台形3つで近似しましょう。ほんとうの値9とのずれは何\,\%ですか。
\bkpans{左はし $0+1+4=5$、右はし $1+4+9=14$、平均9.5。ずれ $\dfrac{0.5}{9}\approx5.6\,\%$}

\begin{furikaeri}
\begin{checklist}
\item 細い長方形の和で、面積を近似できる。\hfill{\footnotesize\sffamily\color{bkTheme}→\pageref{bk:5-10:hasami}ページ}
\item $\Sigma$ の記号を読み書きできる。\hfill{\footnotesize\sffamily\color{bkTheme}→\pageref{bk:5-10:sigma}ページ}
\item 定積分の意味を説明し、$x^2$ の定積分を求められる。\hfill{\footnotesize\sffamily\color{bkTheme}→\pageref{bk:5-10:teisekibun}ページ}
\end{checklist}
\end{furikaeri}

\bkfukidashi{delta}{0から $x$ までの面積は $\dfrac{x^3}{3}$。これを微分すると $x^2$ にもどるね。偶然かな？}

%% 5-11 積分は微分の逆（第11回：ページ制限なし。X16 の6・X19 に対応して書き足した）
\bkchapter{5}{11}{積分は微分の逆}{\bklevel{3}}{積分は微分の逆}
\bkmichishirube{5-4,5-9,5-10}{5-12,5-13,6-3}

\begin{goal}
\begin{itemize}[label={\color{bkTheme}●}, leftmargin=1.6\zw, itemsep=0pt, topsep=1pt]
\item \bklevel{3}\enspace 「面積がふえる速さは、その場所の高さ」であることを説明できる。
\item \bklevel{3}\enspace 原始関数を求め、不定積分を $+C$ をつけて書ける。
\item \bklevel{3}\enspace 原始関数を使って、定積分を $F(b)-F(a)$ で計算できる。
\end{itemize}
\end{goal}

\begin{hatena}
\noindent\begin{minipage}[c]{0.56\linewidth}\raggedright
5-10で、$y=x^2$ の下の面積を、細い長方形をたくさんたして求めました。0から $b$ までの面積は $\dfrac{b^3}{3}$ でした。

ところで、$(x^3)'=3x^2$（5-4）なので、$\dfrac{x^3}{3}$ を微分すると $\dfrac{3x^2}{3}=x^2$ です。面積の式を微分したら、もとのグラフの式 $x^2$ にもどりました。これは偶然でしょうか。
\end{minipage}\hfill
\begin{minipage}[c]{0.4\linewidth}
\begin{tikzpicture}[x=1mm, y=1mm, m/.style={draw=bkTheme, line width=1pt, fill=white, rounded corners=2mm, minimum width=24mm, minimum height=10mm, font=\sffamily\small}]
  \node[m] (a) at (0,12) {$x^2$};
  \node[m] (b) at (0,-12) {$\dfrac{x^3}{3}$};
  \draw[bk arrow, bkBlue] ([xshift=-4mm]a.south) -- node[left, font=\sffamily\footnotesize] {面積（積分）} ([xshift=-4mm]b.north);
  \draw[bk arrow, bkOrange] ([xshift=4mm]b.north) -- node[right, font=\sffamily\footnotesize] {微分} ([xshift=4mm]a.south);
\end{tikzpicture}
\end{minipage}
\end{hatena}

5-9の自転車でも、面積をつみ上げて位置のグラフをつくると、その傾きがもとの速さになりました。偶然ではなさそうです。この章では、「積分（面積）と微分（傾き）は逆のはたらき」であることを確かめます。これがわかると、細い長方形をたさなくても、微分の表を逆に読むだけで積分できるようになります。

\clearpage
%% ---------------- やってみよう・思い出そう ----------------
\begin{yattemiyou}[面積のふえ方を調べよう]
$y=2x$ のグラフの、0から $x$ までの面積を $S(x)$ とします。グラフがまっすぐなので、面積は三角形で求められます。
\begin{enumerate}[label={\bkcircled[bkGreen]{\arabic*}}, leftmargin=2.4\zw, itemsep=1pt, topsep=3pt]
\item $x=1$、2、3、4 のときの $S(x)$ を求め、表にしよう。
\item $x$ が1ふえるごとに、$S(x)$ がどれだけふえるかを調べよう。下の図の色の帯の面積だよ。
\item $S(x)$ を式で表して微分しよう。何になるかな？
\end{enumerate}
\begin{bkcenter}
\begin{bktable}{colspec={Q[l,m,wd=46mm] *{5}{Q[c,m,wd=9mm]}}, rowsep=1.5pt}
$x$ & 0 & 1 & 2 & 3 & 4 \\
高さ $2x$ & 0 & 2 & 4 & 6 & 8 \\
面積 $S(x)=\frac{1}{2}\times x\times2x$ & 0 & 1 & 4 & 9 & 16 \\
1ふえたときの面積のふえ方 & & 1 & 3 & 5 & 7 \\
\end{bktable}
\end{bkcenter}
\begin{bkfigure}
\begin{tikzpicture}
\begin{axis}[bkplot, width=72mm, height=50mm, xmin=0, xmax=4.6, ymin=0, ymax=9.2, xtick={1,2,3,4}, ytick={2,4,6,8}, xlabel={$x$}, ylabel={$y$}]
  \fill[bkGreen!35] (axis cs:0,0) -- (axis cs:1,0) -- (axis cs:1,2) -- cycle;
  \fill[bkBlue!30] (axis cs:1,0) -- (axis cs:2,0) -- (axis cs:2,4) -- (axis cs:1,2) -- cycle;
  \fill[bkPurple!30] (axis cs:2,0) -- (axis cs:3,0) -- (axis cs:3,6) -- (axis cs:2,4) -- cycle;
  \fill[bkOrange!35] (axis cs:3,0) -- (axis cs:4,0) -- (axis cs:4,8) -- (axis cs:3,6) -- cycle;
  \addplot[bkcurve, domain=0:4.3] {2*x};
  \node[font=\sffamily\small] at (axis cs:0.72,0.5) {1};
  \node[font=\sffamily\small] at (axis cs:1.5,1.5) {3};
  \node[font=\sffamily\small] at (axis cs:2.5,2.5) {5};
  \node[font=\sffamily\small] at (axis cs:3.5,3.5) {7};
  \node[font=\sffamily\small, text=bkBlue!80!black, anchor=west] at (axis cs:4.3,8.6) {$y=2x$};
\end{axis}
\end{tikzpicture}
\bkfigcaption{1ずつの帯。帯は台形（はじめは三角形）で、面積は1、3、5、7。}
\end{bkfigure}
\end{yattemiyou}

\noindent 三角形の底辺は $x$、高さは $2x$ なので、$S(x)=\dfrac{1}{2}\times x\times2x=x\times x=x^2$ です（$\frac{1}{2}\times2=1$）。微分すると $S'(x)=2x$（5-4）で、\textbf{グラフの高さ $2x$ と同じ}になりました。

帯で見ると、たとえば $x=1$ から2までの帯は、左の高さ2・右の高さ4の台形なので、面積は $(2+4)\div2\times1=3$ です。面積のふえ方は、その帯の高さ（の平均）と同じです。高い所ほど、面積は速くふえています。

\begin{omoidasou}
\bkrecall{5-4}{$(x^n)'=nx^{n-1}$。定数を微分すると0。定数倍・和は1つずつ微分できる。}
\bkrecall{5-10}{$\displaystyle\int_a^b f(x)\,\dd x$ は、$a$ から $b$ までの、グラフと横軸のあいだの面積（下はマイナス）。$\displaystyle\int_0^b x^2\,\dd x=\frac{b^3}{3}$}
\bkrecall{5-9}{$v$-$t$ グラフの面積をつみ上げると位置のグラフになり、その傾きはもとの速さ。}
\bkrecall{5-3}{$f'(x)=\displaystyle\lim_{h\to0}\frac{f(x+h)-f(x)}{h}$}
\bkrecall{3-10}{かけて0なら、どちらかが0。$x(x-2)=0$ なら $x=0$ か $x=2$。}
\bkrecall{2-8}{2乗して4になる数は、2と $-2$ の2つ。}
\end{omoidasou}

%% ---------------- 数学の道具 ----------------
\bksection{doguu}{数学の道具}

\bksubsection{1}{右はしを動かせる面積 $S(x)$}

5-10では、「0から1まで」「0から2まで」のように、右はしを決めて面積を求めました。こんどは、右はしを動かして、面積がどう変わるかを見ます。カーテンを右へ引いていくと、かくれる部分がだんだん広がるようすを思いうかべましょう。

\begin{bkfigure}
\begin{tikzpicture}
\pgfplotsset{mini/.style={bkplot, width=44mm, height=40mm, xmin=0, xmax=2.4, ymin=0, ymax=4.6, xtick={1,2}, ytick={2,4}, xlabel={$t$}, ylabel={$y$}}}
\begin{axis}[mini, name=m1, title={\sffamily\small 右はし $x=1$}]
  \addplot[fill=bkBlue!30, draw=none, domain=0:1] {x^2} \closedcycle;
  \addplot[bkcurve, line width=1.2pt, domain=0:2.1] {x^2};
\end{axis}
\begin{axis}[mini, name=m2, at={(m1.east)}, anchor=west, xshift=10mm, title={\sffamily\small 右はし $x=1.5$}]
  \addplot[fill=bkBlue!30, draw=none, domain=0:1.5] {x^2} \closedcycle;
  \addplot[bkcurve, line width=1.2pt, domain=0:2.1] {x^2};
\end{axis}
\begin{axis}[mini, name=m3, at={(m2.east)}, anchor=west, xshift=10mm, title={\sffamily\small 右はし $x=2$}]
  \addplot[fill=bkBlue!30, draw=none, domain=0:2] {x^2} \closedcycle;
  \addplot[bkcurve, line width=1.2pt, domain=0:2.1] {x^2};
\end{axis}
\end{tikzpicture}
\bkfigcaption{グラフ $y=t^2$ の、0から右はし $x$ までの面積。右はしを動かすと、面積が変わる。}
\end{bkfigure}

5-10の $\dfrac{b^3}{3}$ で、面積を計算できます。

\begin{bkcenter}
\begin{bktable}{colspec={Q[l,m,wd=40mm] *{5}{Q[c,m,wd=13mm]}}, rowsep=2pt}
右はし $x$ & 0 & 0.5 & 1 & 1.5 & 2 \\
面積 $S(x)=\dfrac{x^3}{3}$ & 0 & 0.0417 & 0.3333 & 1.125 & 2.6667 \\
\end{bktable}
\bktabcaption{右はし $x$ を決めると、面積が1つに決まる。だから、面積は $x$ の関数。}
\end{bkcenter}

右はし $x$ を1つ決めると、面積が1つに決まります。だから、面積は $x$ の関数です。これを $S(x)$ と書きます。

\noindent{\sffamily\bfseries\color{bkTheme}グラフの横軸を $t$ と書くわけ}\par
文字 $x$ は、もう「右はしの場所」に使いました。グラフの横軸にも $x$ を使うと、どちらの $x$ かわからなくなります。そこで、グラフの横軸は $t$ と書き、グラフを $y=t^2$ と書きます。文字がちがうだけで、5-10の $y=x^2$ と同じグラフです。定積分の記号で書くと、
\[ S(x)=\int_0^x t^2\,\dd t \]
です。「$t$ を0から $x$ まで細かく分けてたした面積」と読みます。

\bksubsection{2}{面積がふえる速さは、その場所の高さ}

\label{bk:5-11:kihon}右はしを $x$ から少しだけ動かして $x+h$ にすると、面積は少しふえます。1あたりのふえ方 $\dfrac{S(x+h)-S(x)}{h}$ は、5-2・5-3で速さを求めたときの「平均の速さ」と同じ形です。$h$ を0に近づけた先が、面積のふえる速さ $S'(x)$ です。

\noindent{\sffamily\bfseries\color{bkTheme}数で確かめる}\par
$y=t^2$、右はし $x=1$ で計算します。$S(x)=\dfrac{x^3}{3}$ を使います。たとえば $h=0.1$ なら、$S(1.1)=\dfrac{1.1^3}{3}=\dfrac{1.331}{3}$、$S(1)=\dfrac{1}{3}$ なので、ふえた面積は $\dfrac{1.331-1}{3}=0.110333\cdots$ です。これを $h=0.1$ でわると 1.10333 です。

\begin{bkcenter}
\begin{bktable}{colspec={Q[c,m,wd=12mm] Q[c,m,wd=36mm] Q[c,m,wd=24mm] Q[c,m,wd=30mm]}, rowsep=2pt, row{1}={font=\sffamily\bfseries\small}}
$h$ & ふえた面積 $S(1+h)-S(1)$ & 1あたり（$\div h$） & 右はしの高さ $(1+h)^2$ \\
0.1 & 0.110333 & 1.10333 & 1.21 \\
0.01 & 0.0101003 & 1.01003 & 1.0201 \\
0.001 & 0.0010010003 & 1.0010003 & 1.002001 \\
\end{bktable}
\bktabcaption{$h$ を小さくすると、1あたりのふえ方は1に近づく。1は、$x=1$ でのグラフの高さ $1^2$。}
\end{bkcenter}

1あたりのふえ方は、1.10333、1.01003、1.0010003と、1に近づきます。1は、$x=1$ でのグラフの高さ $1^2=1$ です。「面積がふえる速さ」が「その場所の高さ」になりました。

\noindent{\sffamily\bfseries\color{bkTheme}図で確かめる}\par
グラフ $y=f(t)$ の、$a$ から $x$ までの面積を $S(x)$ とします。右はしを $x+h$ にすると、面積は、はば $h$ の細い帯の分だけふえます。

\noindent\begin{minipage}[c]{0.5\linewidth}
\begin{tikzpicture}
\begin{axis}[bkplot, width=62mm, height=46mm, xmin=0, xmax=3.4, ymin=0, ymax=3.4, xtick={0.5,2,2.4}, xticklabels={$a$,$x$,$x\!+\!h$}, ytick=\empty, xlabel={$t$}, ylabel={$y$}]
  \addplot[fill=bkBlue!18, draw=none, domain=0.5:2] {0.6+0.35*x^2} \closedcycle;
  \addplot[fill=bkOrange!45, draw=none, domain=2:2.4] {0.6+0.35*x^2} \closedcycle;
  \addplot[bkcurve, domain=0:2.75] {0.6+0.35*x^2};
  \node[font=\sffamily\small] at (axis cs:1.3,0.5) {$S(x)$};
  \draw[bkGray, line width=0.5pt, dashed] (axis cs:2,0) -- (axis cs:2,2.0);
  \node[font=\sffamily\footnotesize, anchor=west] at (axis cs:2.5,1.0) {高さ $f(x)$};
\end{axis}
\end{tikzpicture}
\end{minipage}\hfill
\begin{minipage}[c]{0.46\linewidth}
帯（オレンジ）は、はば $h$、高さほぼ $f(x)$ の長方形なので、
\[ S(x+h)-S(x)\approx f(x)\times h \]
両辺を $h$ でわると
\[ \frac{S(x+h)-S(x)}{h}\approx f(x) \]
\end{minipage}

帯を大きくして見ましょう。帯は、2つの長方形にはさまれています（5-10の「はさみうち」と同じ考え方）。

\begin{bkfigure}
\begin{tikzpicture}[x=1mm, y=1mm]
  % 帯の拡大：左の高さ f(x)=14、右の高さ f(x+h)=24、はば h=24mm
  \fill[bkOrange!40] (0,0) -- (0,14) .. controls (9,17) and (16,20) .. (24,24) -- (24,0) -- cycle;
  \draw[bkBlue!70!black, line width=0.6pt, fill=bkBlue!18, fill opacity=0.7] (0,0) rectangle (24,14);
  \draw[bkRed!80!black, line width=0.8pt, dash pattern=on 2pt off 1.2pt] (0,0) rectangle (24,24);
  \draw[bkBlue!80!black, line width=1.4pt] (-6,10.5) .. controls (-3,12.4) .. (0,14) .. controls (9,17) and (16,20) .. (24,24) .. controls (27,25.6) .. (30,27.5);
  \draw[bkInk, line width=0.6pt] (-8,0) -- (34,0);
  \node[font=\footnotesize, anchor=north, inner sep=1pt] at (0,-1) {$x$};
  \node[font=\footnotesize, anchor=north, inner sep=1pt] at (24,-1) {$x+h$};
  \draw[bkGray, line width=0.5pt, <->] (-3,0) -- (-3,14);
  \node[font=\footnotesize, anchor=east, inner sep=1pt] at (-3.6,7) {$f(x)$};
  \draw[bkGray, line width=0.5pt, <->] (27,0) -- (27,24);
  \node[font=\footnotesize, anchor=west, inner sep=1pt] at (27.6,12) {$f(x+h)$};
  \node[font=\sffamily\footnotesize, anchor=west, align=left, inner sep=0pt] (t1) at (44,22) {{\color{bkRed!80!black}点線の長方形}：$f(x+h)\times h$（大きめ）};
  \node[font=\sffamily\footnotesize, anchor=west, align=left, inner sep=0pt] (t2) at (44,14) {{\color{bkOrange!80!black}帯}：$S(x+h)-S(x)$（ほんとう）};
  \node[font=\sffamily\footnotesize, anchor=west, align=left, inner sep=0pt] (t3) at (44,6) {{\color{bkBlue!70!black}青の長方形}：$f(x)\times h$（小さめ）};
\end{tikzpicture}
\bkfigcaption{グラフが右上がりのとき、帯は、青の長方形より大きく、点線の長方形より小さい。}
\end{bkfigure}

だから、
\[ f(x)\times h<S(x+h)-S(x)<f(x+h)\times h \]
です。全体を $h$ でわると（$h$ はプラスの数なので、大小は変わりません）、
\[ f(x)<\frac{S(x+h)-S(x)}{h}<f(x+h) \]
上の表の $h=0.1$ なら、$1<1.10333<1.21$ です。$h$ を0に近づけると、右はしの $f(x+h)$ は $f(x)$ に近づきます。まん中の数は、$f(x)$ と、$f(x)$ に近づく数にはさまれているので、$f(x)$ に近づきます。左がわは $S(x)$ の微分 $S'(x)$ の定義（5-3）そのものなので、
\[ S'(x)=f(x) \]
\textbf{面積 $S(x)$ を微分すると、グラフの高さ $f(x)$ になる}。これを、\textbf{微分積分学の基本定理}といいます。面積（積分）と傾き（微分）が、逆のはたらきでつながっていることを表しています。\bkidx{びぶんせきぶんがくのきほんていり}{微分積分学の基本定理}

\noindent{\sffamily\bfseries\color{bkTheme}たとえ}\par
へいのペンキぬりを考えます。へいの高さが場所によってちがうとき、横に\qty{1}{cm}ぬり進むごとにふえるぬった面積は、その場所のへいの高さ×\qty{1}{cm}です。高い所では面積が速くふえ、低い所ではゆっくりふえます。

\noindent{\sffamily\bfseries\color{bkTheme}物理でいうと}\par
$v$-$t$ グラフの面積は、進んだ距離でした（5-9）。0秒から $t$ 秒までの面積は、$t$ 秒での位置（出発点から）です。位置のふえる速さは、その時刻の速さです。「面積を微分すると高さ」は、「位置を微分すると速度」（5-5）と同じことを言っています。

\begin{seikaku}
上の「$f(x)$ と $f(x+h)$ にはさまれる」は、グラフが右上がりのときの話です。右下がりのときは大小が入れかわるだけで、同じように $f(x)$ に近づきます。$h$ を0に近づけると $f(x+h)$ が $f(x)$ に近づくには、$f$ がつながったグラフ（とちゅうで切れていない）であることが必要です。この巻であつかう関数は、どれもこの条件をみたしています。
\end{seikaku}

\bksubsection{3}{微分すると $f(x)$ になる関数：原始関数と不定積分}

\label{bk:5-11:genshi}②から、面積 $S(x)$ は「微分すると $f(x)$ になる関数」です。だから、面積を求めるには、微分すると $f(x)$ になる関数をさがせばよいのです。

\noindent{\sffamily\bfseries\color{bkTheme}さがしてみる}\par
微分すると $x^2$ になる関数をさがします。微分すると指数が1へるので（5-4）、指数が1大きい $x^3$ をためします。

\begin{bkcenter}
\begin{bktable}{colspec={Q[c,m,wd=30mm] Q[c,m,wd=30mm] Q[l,m,wd=52mm]}, rowsep=2pt, row{1}={font=\sffamily\bfseries\small}}
ためす関数 & 微分すると & どうする？ \\
$x^3$ & $3x^2$ & $x^2$ の3倍。多すぎるので3でわる \\
$\dfrac{x^3}{3}$ & $\dfrac{3x^2}{3}=x^2$ & ぴったり！ \\
\end{bktable}
\end{bkcenter}

このように、微分すると $f(x)$ になる関数 $F(x)$ を、$f(x)$ の\textbf{原始関数}（げんしかんすう）といいます。「原始」は「もとの」という意味で、微分する前のもとの関数ということです。$x^2$ の原始関数は $\dfrac{x^3}{3}$ です。\bkidx{げんしかんすう}{原始関数}

\noindent{\sffamily\bfseries\color{bkTheme}原始関数は1つではない}\par
$\dfrac{x^3}{3}+1$ や $\dfrac{x^3}{3}-1$ も、微分すると $x^2$ になります。定数は微分すると0だからです（5-4）。グラフで見ると、上下にずらしただけなので、同じ $x$ での傾きは変わりません。

\begin{bkfigure}
\begin{tikzpicture}
\begin{axis}[bkplot, width=84mm, height=58mm, xmin=-0.2, xmax=2.2, ymin=-1.6, ymax=3.9, xtick={1,2}, ytick={1,2,3}, xlabel={$x$}, ylabel={$y$}]
  \addplot[bkcurve, domain=-0.2:1.95] {x^3/3+1};
  \addplot[bkcurve2, domain=-0.2:1.95] {x^3/3};
  \addplot[bkcurve3, domain=-0.2:1.95] {x^3/3-1};
  \foreach \c in {1,0,-1} {\edef\tmp{\noexpand\addplot[bkInk, line width=0.9pt, domain=0.55:1.45] {1/3+\c+1*(x-1)};}\tmp}
  \addplot[only marks, mark=*, mark size=1.5pt, bkInk] coordinates {(1,4/3) (1,1/3) (1,-2/3)};
  \node[font=\footnotesize, text=bkBlue!80!black, anchor=west] at (axis cs:1.97,3.55) {$\frac{x^3}{3}+1$};
  \node[font=\footnotesize, text=bkOrange!80!black, anchor=west] at (axis cs:1.97,2.5) {$\frac{x^3}{3}$};
  \node[font=\footnotesize, text=bkGreen!60!black, anchor=west] at (axis cs:1.97,1.5) {$\frac{x^3}{3}-1$};
\end{axis}
\end{tikzpicture}
\bkfigcaption{上下にずらした3つのグラフ。$x=1$ での接線（黒）は、どれも傾き $1^2=1$ で平行。}
\end{bkfigure}

そこで、原始関数をまとめて、どれでもよい定数を $C$ として
\[ \int f(x)\,\dd x=F(x)+C\qquad（C\text{ は定数}） \]
と書き、\textbf{不定積分}（ふていせきぶん）、$C$ を\textbf{積分定数}といいます。定積分とちがって、$\int$ に下端と上端がついていません。$C$ の分だけ決まらない（定まらない）ので「不定」積分といいます。$C$ の意味は、5-13でくわしく学びます。\bkidx{ふていせきぶん}{不定積分}\bkidx{せきぶんていすう}{積分定数 $C$}

\noindent{\sffamily\bfseries\color{bkTheme}$x^n$ の不定積分}\par
同じように、$1$、$x$、$x^3$ の原始関数もさがして、ならべてみます（$1=x^0$ です。2-7）。

\begin{bkcenter}
\begin{bktable}{colspec={Q[c,m,wd=16mm] Q[c,m,wd=24mm] Q[c,m,wd=27mm] Q[c,m,wd=30mm]}, rowsep=2pt, row{1}={font=\sffamily\bfseries\small}}
$f(x)$ & 原始関数 & 微分して確かめ & 指数 \\
$1=x^0$ & $x=\dfrac{x^1}{1}$ & $(x)'=1$ & $0\to1$、1でわる \\
$x=x^1$ & $\dfrac{x^2}{2}$ & $\dfrac{2x}{2}=x$ & $1\to2$、2でわる \\
$x^2$ & $\dfrac{x^3}{3}$ & $\dfrac{3x^2}{3}=x^2$ & $2\to3$、3でわる \\
$x^3$ & $\dfrac{x^4}{4}$ & $\dfrac{4x^3}{4}=x^3$ & $3\to4$、4でわる \\
\end{bktable}
\bktabcaption{どれも「指数を1ふやして、その数でわる」になっている。}
\end{bkcenter}

きまりが見えました。どの $n$ でも、微分して確かめられます。微分のきまり $(x^n)'=nx^{n-1}$ の $n$ を $n+1$ にかえて使うと、
\[ \left(\frac{x^{n+1}}{n+1}\right)'=\frac{(n+1)x^n}{n+1}=x^n\quad\text{なので}\quad\int x^n\,\dd x=\frac{x^{n+1}}{n+1}+C\quad（n\neq-1） \]
「指数を1ふやして、その数でわる」です。微分（指数を1へらし、もとの指数をかける）のちょうど逆です。$n=-1$（$\frac{1}{x}=x^{-1}$）のときは、指数が $-1+1=0$ になり、0でわることになるので使えません（1-14）。$\frac{1}{x}$ の積分は5-14で学びます。

\noindent{\sffamily\bfseries\color{bkTheme}定数倍と和}\par
定数倍と和は、微分と同じように1つずつ積分できます。理由は2つあります。
\begin{enumerate}[label={\bkcircled[bkTheme]{\arabic*}}, leftmargin=2.4\zw, itemsep=1pt, topsep=2pt]
\item 面積で考える：高さが3倍の長方形をたせば、和も3倍です（5-10の例題6の図）。高さ $f+g$ の長方形は、$f$ と $g$ の長方形を積み重ねたものです。
\item 微分で確かめる：$(3F)'=3F'$、$(F+G)'=F'+G'$（5-4）なので、逆向きにも成り立ちます。
\end{enumerate}

\begin{bkcenter}
\begin{bktable}{colspec={Q[c,m,wd=24mm] Q[c,m,wd=46mm] Q[c,m,wd=34mm]}, rowsep=1.5pt, row{1}={font=\sffamily\bfseries\small}}
$f(x)$ & 不定積分 $\int f(x)\,\dd x$ & 確かめ（微分する） \\
$5$ & $5x+C$ & $(5x)'=5$ \\
$6x$ & $6\times\dfrac{x^2}{2}+C=3x^2+C$ & $3\times2x=6x$ \\
$3x^2+2x$ & $x^3+x^2+C$ & $3x^2+2x$ \\
\end{bktable}
\bktabcaption{不定積分は、微分して、もとにもどるかで確かめられる。}
\end{bkcenter}

\bksubsection{4}{定積分は $F(b)-F(a)$}

\label{bk:5-11:teisekibun}原始関数を使って、定積分を計算する方法を作ります。

\noindent{\sffamily\bfseries\color{bkTheme}数で考える}\par
$y=t^2$ の、1から $x$ までの面積を $S(x)$ とします。5-10の「0からでないときは、大きい面積から小さい面積をひく」を使うと、
\[ S(x)=\frac{x^3}{3}-\frac{1^3}{3}=\frac{x^3}{3}-\frac{1}{3} \]
です。$x=1$ を入れると $\dfrac{1}{3}-\dfrac{1}{3}=0$。1から1までは、はばが0なので面積も0で、合っています。$x=2$ なら $\dfrac{8}{3}-\dfrac{1}{3}=\dfrac{7}{3}$ です。

この $S(x)$ は、原始関数 $\dfrac{x^3}{3}$ に、定数 $-\dfrac{1}{3}$ をたした形です。$-\dfrac{1}{3}$ は「左はしで面積が0になるように」決まった数です。

\noindent{\sffamily\bfseries\color{bkTheme}一般にする}\par
$a$ から $x$ までの面積 $S(x)$ も、微分すると $f(x)$ になる（②）ので、$f(x)$ の原始関数の1つです。だから、どれか1つの原始関数 $F(x)$ を使って、次のように書けます。
\begin{enumerate}[label={\bkcircled[bkTheme]{\arabic*}}, leftmargin=2.4\zw, itemsep=1pt, topsep=2pt]
\item $S(x)=F(x)+C$\quad（原始関数は、$C$ だけちがう）
\item $x=a$ では、はばが0なので面積も0。だから $0=F(a)+C$、つまり $C=-F(a)$
\item ①に入れて、$S(x)=F(x)-F(a)$
\item $x=b$ とすると、$a$ から $b$ までの面積は $S(b)=F(b)-F(a)$
\end{enumerate}
\[ \int_a^b f(x)\,\dd x=F(b)-F(a)=\Bigl[\,F(x)\,\Bigr]_a^b \]
$\Bigl[\,F(x)\,\Bigr]_a^b$ は「$F(b)-F(a)$ を計算する」という書き方です。\textbf{上端を入れた値から、下端を入れた値をひきます}（「あと${}-{}$まえ」、2-3）。

\noindent{\sffamily\bfseries\color{bkTheme}物理でいうと}\par
$v$-$t$ グラフの面積は、位置の変化（変位）でした（5-9）。位置 $x(t)$ は、微分すると速度 $v$ になる関数、つまり $v$ の原始関数です。だから
\[ \int_a^b v\,\dd t=x(b)-x(a)=\text{（あとの位置）}-\text{（まえの位置）} \]
「速度の面積は、あとの位置からまえの位置をひいたもの」。上の式と同じことです。

\noindent{\sffamily\bfseries\color{bkTheme}どの原始関数を使っても、答えは同じ}\par
$x^2$ の原始関数として $\dfrac{x^3}{3}+5$ を使っても、
\[ \Bigl(\frac{8}{3}+5\Bigr)-\Bigl(\frac{1}{3}+5\Bigr)=\frac{8}{3}-\frac{1}{3}=\frac{7}{3} \]
です。$+5$ はひき算で消えます。だから、定積分では $C$ を書かなくてよいのです。

\noindent\textbf{例}\quad $\displaystyle\int_0^1x^2\,\dd x=\Bigl[\,\frac{x^3}{3}\,\Bigr]_0^1=\frac{1^3}{3}-\frac{0^3}{3}=\frac{1}{3}-0=\frac{1}{3}$。5-10で、細い長方形を1000本もたして見つけた $\dfrac{1}{3}$ が、1行で求められました。

\noindent{\sffamily\bfseries\color{bkTheme}下端と上端を入れかえると}\par
$\displaystyle\int_b^a f(x)\,\dd x=F(a)-F(b)=-\bigl(F(b)-F(a)\bigr)$ です。入れかえると符号が反対になります。5-10で「はばがマイナスになるから」と考えたことと、同じ結果です。

\bksubsection{5}{横軸より下の部分}

\label{bk:5-11:shita}定積分では、横軸より下の部分はマイナスとして計算されます（5-10）。$y=x^2-2x$ で調べます。$x^2-2x=x(x-2)$ なので、$x=0$ と $x=2$ でグラフは横軸と交わります（3-10）。

\noindent\begin{minipage}[c]{0.48\linewidth}
\begin{tikzpicture}
\begin{axis}[bkplot, width=62mm, height=62mm, xmin=-0.3, xmax=3.8, ymin=-1.6, ymax=3.6, xtick={2,3}, ytick={1,2,3}, xlabel={$x$}, ylabel={$y$}]
  \addplot[fill=bkRed!25, draw=none, domain=0:2] {x^2-2*x} \closedcycle;
  \addplot[fill=bkBlue!25, draw=none, domain=2:3] {x^2-2*x} \closedcycle;
  \addplot[bkcurve, domain=-0.25:3.25] {x^2-2*x};
  \node[font=\sffamily\small] at (axis cs:1,-0.5) {$-\frac{4}{3}$};
  \node[font=\sffamily\small, anchor=west] at (axis cs:3.2,1.1) {$+\frac{4}{3}$};
  \draw[bkGray, line width=0.4pt] (axis cs:3.22,1.0) -- (axis cs:2.85,0.55);
\end{axis}
\end{tikzpicture}
\end{minipage}\hfill
\begin{minipage}[c]{0.48\linewidth}
赤い部分は横軸より下、青い部分は横軸より上です。原始関数は $\dfrac{x^3}{3}-x^2$ なので、それぞれの定積分は次のようになります。
\end{minipage}
\begin{align*}
\int_0^2(x^2-2x)\,\dd x&=\Bigl(\frac{8}{3}-4\Bigr)-0=-\frac{4}{3}&&\text{（赤）}\\
\int_2^3(x^2-2x)\,\dd x&=(9-9)-\Bigl(-\frac{4}{3}\Bigr)=\frac{4}{3}&&\text{（青）}
\end{align*}

だから、0から3まで通して計算すると、
\[ \int_0^3(x^2-2x)\,\dd x=\Bigl[\,\frac{x^3}{3}-x^2\,\Bigr]_0^3=(9-9)-0=0 \]
となります。でも、面積が0なのではありません。下の部分（$-\frac{4}{3}$）と、上の部分（$+\frac{4}{3}$）が打ち消しあったのです。物理では、これが「もどった分を差し引いた変位」にあたります（5-9）。囲まれた部分の面積を知りたいときは、横軸と交わる所で分けて、それぞれの大きさをたします（$\frac{4}{3}+\frac{4}{3}=\frac{8}{3}$）。

\bksubsection{6}{2つのグラフをならべて見る}

\begin{bkfigure}
\begin{tikzpicture}
\begin{axis}[bkplot, name=u1, width=60mm, height=37mm, xmin=0, xmax=2.4, ymin=0, ymax=4.6, xtick={1,1.5,2}, ytick={1,2.25,4}, xlabel={$x$}, ylabel={$y$}, title={\sffamily\small 高さ\enspace$f(x)=x^2$}]
  \addplot[fill=bkBlue!18, draw=none, domain=0:1.5] {x^2} \closedcycle;
  \addplot[bkcurve, domain=0:2.1] {x^2};
  \draw[bkOrange, line width=1.2pt] (axis cs:1.5,0) -- (axis cs:1.5,2.25);
  \node[font=\sffamily\footnotesize, text=bkOrange!80!black] at (axis cs:1.8,0.45) {高さ};
  \node[font=\sffamily\footnotesize] at (axis cs:1.15,0.3) {面積};
\end{axis}
\begin{axis}[bkplot, at={(u1.east)}, anchor=west, xshift=14mm, width=60mm, height=37mm, xmin=0, xmax=2.4, ymin=0, ymax=3, xtick={1,1.5,2}, ytick={1,2}, xlabel={$x$}, ylabel={$y$}, title={\sffamily\small 面積\enspace$S(x)=\dfrac{x^3}{3}$}]
  \addplot[bkcurve2, domain=0:2.05] {x^3/3};
  \addplot[bkcurve, line width=1pt, domain=1.05:1.95] {1.125+2.25*(x-1.5)};
  \addplot[only marks, mark=*, mark size=1.6pt, bkInk] coordinates {(1.5,1.125)};
  \node[font=\sffamily\footnotesize, align=center, text=bkBlue] at (axis cs:1.75,0.5) {傾き\\2.25};
\end{axis}
\end{tikzpicture}
\bkfigcaption{左の面積（青）を右の高さにかいた。右の傾き＝左の高さ。}
\end{bkfigure}

$x=1.5$ では、右のグラフの接線の傾き $S'(1.5)=1.5^2=2.25$ が、左のグラフの高さ $f(1.5)=2.25$ と一致しています。左のグラフが高い所ほど、右のグラフは急に上がります。左のグラフが0の所（$x=0$）では、右のグラフは平ら（傾き0）です。

\begin{taisetsu}
微分積分学の基本定理\quad 面積 $S(x)$ を微分すると高さ $f(x)$\\
不定積分\quad $\displaystyle\int x^n\,\dd x=\frac{x^{n+1}}{n+1}+C$（指数を1ふやして、その数でわる）\\
定積分\quad $\displaystyle\int_a^b f(x)\,\dd x=\Bigl[\,F(x)\,\Bigr]_a^b=F(b)-F(a)$（上端の値－下端の値）
\end{taisetsu}

\begin{kotoba}
\small
\bkword{原始関数}{微分すると $f(x)$ になる関数。}
\bkword{不定積分・積分定数}{原始関数をまとめて $F(x)+C$ と書いたもの。$C$ が積分定数。}
\bkword{微分積分学の基本定理}{面積を微分すると、もとの高さになる。積分は微分の逆。}
\end{kotoba}

%% ---------------- 例題 ----------------
\bksection{reidai}{例題でたしかめよう}

\begin{reidai}[$x^n$ と定数の不定積分]
次の不定積分を求めましょう。\quad（1）$\displaystyle\int x^3\,\dd x$\quad（2）$\displaystyle\int6x^2\,\dd x$\quad（3）$\displaystyle\int5\,\dd x$\par
\kangaekata ③の「指数を1ふやして、その数でわる」を使います。数の倍（6）はそのまま残します。5は $5x^0$ と考えます。\par
\shiki（1）\nobreak 指数3を4にして、4でわる：$\dfrac{x^4}{4}+C$\par
\hspace*{2.6\zw}（2）\nobreak $6\times\dfrac{x^3}{3}+C=2x^3+C$\par
\hspace*{2.6\zw}（3）\nobreak $5\times\dfrac{x^1}{1}+C=5x+C$\par
\bkstep[bkGreen]{\R{確}{たし}かめ}微分すると、（1）$\dfrac{4x^3}{4}=x^3$、（2）$2\times3x^2=6x^2$、（3）$5$。もとにもどります。\par
\kotae（1）\nobreak $\dfrac{x^4}{4}+C$\quad（2）\nobreak $2x^3+C$\quad（3）\nobreak $5x+C$
\end{reidai}

\begin{reidai}[項がいくつかある不定積分]
$\displaystyle\int(3x^2+2x-4)\,\dd x$ を求めましょう。\par
\kangaekata 和は1つずつ積分できます（③）。項ごとに\bkreidairef{1}と同じことをして、$+C$ は最後に1つだけつけます（$C$ をいくつたしても定数）。\par
\shiki $3x^2\to3\times\dfrac{x^3}{3}=x^3$、\quad$2x\to2\times\dfrac{x^2}{2}=x^2$、\quad$-4\to-4x$\par
\hspace*{2.6\zw}合わせて $x^3+x^2-4x+C$\par
\bkstep[bkGreen]{\R{確}{たし}かめ}$(x^3+x^2-4x+C)'=3x^2+2x-4$\par
\kotae $x^3+x^2-4x+C$
\end{reidai}

\noindent\begin{minipage}[t]{0.485\linewidth}\vspace{0pt}\raggedright
\begin{jibun}
次の不定積分を求めましょう。\quad（1）$\displaystyle\int x^4\,\dd x$\quad（2）$\displaystyle\int4x\,\dd x$\quad（3）$\displaystyle\int3\,\dd x$
\bkjans{（1）\nobreak $\dfrac{x^5}{5}+C$\quad（2）\nobreak $2x^2+C$\quad（3）\nobreak $3x+C$}
\end{jibun}
\end{minipage}\hfill
\begin{minipage}[t]{0.485\linewidth}\vspace{0pt}\raggedright
\begin{jibun}
$\displaystyle\int(x^2-4x+1)\,\dd x$ を求めましょう。
\bkjans{$\dfrac{x^3}{3}-2x^2+x+C$}
\end{jibun}
\end{minipage}

\begin{reidai}[定積分]
$\displaystyle\int_1^3(2x+1)\,\dd x$ を求めましょう。\par
\kangaekata まず原始関数を1つ見つけます。$2x\to x^2$、$1\to x$ なので $F(x)=x^2+x$（$C$ は消えるので書かない）。次に、④の $F(3)-F(1)$ を計算します。\par
\shiki $F(3)=3^2+3=12$、\quad$F(1)=1^2+1=2$\par
\hspace*{2.6\zw}$\Bigl[\,x^2+x\,\Bigr]_1^3=12-2=10$\par
\bkstep[bkGreen]{\R{確}{たし}かめ}グラフ $y=2x+1$ はまっすぐなので、1から3までは、左の高さ3・右の高さ7・はば2の台形。$\dfrac{(3+7)\times2}{2}=10$。合っています。\par
\kotae 10
\end{reidai}

\begin{reidai}[速度から進んだ距離]
速度が $v=3t^2+2$（m/s）の物体が、1秒から2秒までに進んだ距離を求めましょう。\par
\kangaekata 距離は $v$-$t$ グラフの面積（5-9）。④の定積分で求めます。速度がずっとプラスなので、変位と道のりは同じです。\par
\shiki 原始関数は $t^3+2t$\par
\hspace*{2.6\zw}$\displaystyle\int_1^2(3t^2+2)\,\dd t=\Bigl[\,t^3+2t\,\Bigr]_1^2=(8+4)-(1+2)=12-3=9$\par
\kotae \qty{9}{m}
\end{reidai}

\noindent\begin{minipage}[t]{0.485\linewidth}\vspace{0pt}\raggedright
\begin{jibun}
$\displaystyle\int_0^2(3x^2+1)\,\dd x$ を求めましょう。
\bkjans{$\Bigl[\,x^3+x\,\Bigr]_0^2=10-0=10$}
\end{jibun}
\end{minipage}\hfill
\begin{minipage}[t]{0.485\linewidth}\vspace{0pt}\raggedright
\begin{jibun}
速度が $v=10t$（m/s）の物体が、0秒から3秒までに進んだ距離を求めましょう（5-2の $x=5t^2$ とくらべよう）。
\bkjans{$\Bigl[\,5t^2\,\Bigr]_0^3=\qty{45}{m}$}
\end{jibun}
\end{minipage}

\begin{column}{積分の記号の由来}
$\int$ は、ライプニッツが17世紀に使い始めた記号です。ラテン語で「和」を意味する summa の頭文字 S を縦にのばした形で、「細い長方形 $f(x)\,\dd x$ をたし合わせる」という意味がこめられています。微分の $\dd y/\dd x$ も彼の記号です。2つの記号が「逆のはたらき」をうまく表していたことが、その後の数学と物理の発展を助けました。
\end{column}

\begin{reidai}[放物線の下の面積]
$y=4-x^2$ のグラフと横軸で囲まれた部分の面積を求めましょう。\par
\kangaekata まず、グラフが横軸と交わる所（$y=0$ になる所）を求めます。$4-x^2=0$ より $x^2=4$。2乗して4になる数は $x=2$ と $x=-2$（2-8）。そのあいだでは、たとえば $x=0$ で $y=4$ なので、横軸より上です。だから、$-2$ から2までの定積分がそのまま面積です。\par
\begin{center}
\begin{tikzpicture}
\begin{axis}[bkplot, width=52mm, height=34mm, xmin=-2.7, xmax=2.9, ymin=-0.8, ymax=4.8, xtick={-2,2}, ytick={4}, xlabel={$x$}, ylabel={$y$}]
  \addplot[fill=bkBlue!25, draw=none, domain=-2:2] {4-x^2} \closedcycle;
  \addplot[bkcurve, line width=1.2pt, domain=-2.12:2.12] {4-x^2};
  \node[font=\footnotesize, text=bkBlue!80!black, anchor=west] at (axis cs:1.15,3.4) {$y=4-x^2$};
\end{axis}
\end{tikzpicture}
\end{center}
\shiki 原始関数は $4x-\dfrac{x^3}{3}$\par
\hspace*{2.6\zw}上端：$4\times2-\dfrac{2^3}{3}=8-\dfrac{8}{3}=\dfrac{16}{3}$\par
\hspace*{2.6\zw}下端：$4\times(-2)-\dfrac{(-2)^3}{3}=-8+\dfrac{8}{3}=-\dfrac{16}{3}$\par
\hspace*{2.6\zw}$\dfrac{16}{3}-\Bigl(-\dfrac{16}{3}\Bigr)=\dfrac{32}{3}$\par
\bkstep[bkGreen]{\R{確}{たし}かめ}囲まれた部分は、はば4・高さ4の長方形（面積16）の中に入っています。$\dfrac{32}{3}=10.66\cdots$ は16より小さく、ちょうど $16\times\dfrac{2}{3}$ です（放物線の山の面積は、外の長方形の $\dfrac{2}{3}$。2000年以上前にアルキメデスが見つけました）。\par
\kotae $\dfrac{32}{3}$
\end{reidai}

\begin{jibun}
$y=3x-x^2$ のグラフと横軸で囲まれた部分の面積を求めましょう。
\bkjans{$3x-x^2=x(3-x)=0$ より $x=0,\,3$。$\Bigl[\,\dfrac{3x^2}{2}-\dfrac{x^3}{3}\,\Bigr]_0^3=\dfrac{27}{2}-9=\dfrac{9}{2}$}
\end{jibun}

\begin{reidai}[面積の式を微分する]
$S(x)=\displaystyle\int_0^x(t^2+1)\,\dd t$ のとき、$S'(2)$ を求めましょう。\par
\kangaekata ②の基本定理から、$S(x)$ を微分すると、右はし $x$ でのグラフの高さ $x^2+1$ になります。積分の計算はいりません。\par
\shiki $S'(x)=x^2+1$（$t$ を $x$ にかえる）なので、$S'(2)=2^2+1=5$\par
\bkstep[bkGreen]{\R{確}{たし}かめ}$S(x)=\Bigl[\,\dfrac{t^3}{3}+t\,\Bigr]_0^x=\dfrac{x^3}{3}+x$ を微分すると $x^2+1$。同じです。\par
\kotae 5
\end{reidai}

\begin{reidai}[横軸より下の部分：変位と道のり]
速度が $v=2t-4$（m/s）の物体について、0秒から3秒までの変位（位置の変化）と、実際に動いた道のりを求めましょう。\par
\kangaekata $2t-4=0$ になるのは $t=2$。$0<t<2$ では $v<0$（もどる向き）なので、その部分の定積分はマイナスになります（⑤）。変位はそのまま計算します。道のりは、符号の変わる $t=2$ で分けて、大きさをたします（5-9）。原始関数は $t^2-4t$ です。\par
\shiki 変位\quad $\Bigl[\,t^2-4t\,\Bigr]_0^3=(9-12)-0=-3$\par
\hspace*{2.6\zw}0〜2秒\quad $\Bigl[\,t^2-4t\,\Bigr]_0^2=(4-8)-0=-4$\par
\hspace*{2.6\zw}2〜3秒\quad $\Bigl[\,t^2-4t\,\Bigr]_2^3=-3-(-4)=1$\par
\hspace*{2.6\zw}道のりは、大きさをたして $4+1=5$\par
\kotae 変位 \qty{-3}{m}、道のり \qty{5}{m}
\end{reidai}

\noindent\begin{minipage}[t]{0.485\linewidth}\vspace{0pt}\raggedright
\begin{jibun}
$S(x)=\displaystyle\int_1^x(3t^2-t)\,\dd t$ のとき、$S'(3)$ を求めましょう。
\bkjans{$3\times9-3=24$}
\end{jibun}
\end{minipage}\hfill
\begin{minipage}[t]{0.485\linewidth}\vspace{0pt}\raggedright
\begin{jibun}
速度が $v=3-t$（m/s）の物体について、0秒から4秒までの変位と道のりを求めましょう。
\bkjans{変位 $\Bigl[\,3t-\dfrac{t^2}{2}\,\Bigr]_0^4=4$。\mbox{0〜3秒}で4.5、\mbox{3〜4秒}で $-0.5$。道のり\qty{5}{m}}
\end{jibun}
\end{minipage}

\begin{reidai}[2つのグラフにはさまれた面積]
$y=x$ と $y=x^2$ のグラフで囲まれた部分の面積を求めましょう。\par
\noindent\begin{minipage}[t]{0.64\linewidth}\raggedright
\kangaekata 2つのグラフが交わる所は、$x=x^2$ より $x^2-x=0$、$x(x-1)=0$ で、$x=0,\,1$（3-10）。そのあいだでは、たとえば $x=0.5$ で $0.5>0.25$ なので、$y=x$ のほうが上です。右の図のように、はさまれた部分を細い長方形に分けると、高さは「上－下」＝ $x-x^2$ です。だから、$x-x^2$ を0から1まで積分します。
\end{minipage}\hfill
\begin{minipage}[t]{0.33\linewidth}\centering\vspace{0pt}
\begin{tikzpicture}
\begin{axis}[bkplot, width=42mm, height=40mm, xmin=0, xmax=1.25, ymin=0, ymax=1.25, xtick={1}, ytick={1}, xlabel={$x$}, ylabel={$y$}]
  \addplot[fill=bkBlue!20, draw=none, domain=0:1, samples=40] {x} \closedcycle;
  \addplot[fill=white, draw=none, domain=0:1, samples=40] {x^2} \closedcycle;
  \fill[bkOrange!60] (axis cs:0.55,0.3025) rectangle (axis cs:0.62,0.62);
  \addplot[bkcurve, line width=1.2pt, domain=0:1.15] {x};
  \addplot[bkcurve2, line width=1.2pt, domain=0:1.1] {x^2};
  \node[font=\footnotesize, text=bkBlue!80!black, anchor=south east, inner sep=1pt] at (axis cs:0.5,0.62) {$y=x$};
  \node[font=\footnotesize, text=bkOrange!80!black, anchor=west, inner sep=1pt] at (axis cs:0.82,0.42) {$y=x^2$};
\end{axis}
\end{tikzpicture}
\end{minipage}\par
\shiki $\displaystyle\int_0^1(x-x^2)\,\dd x=\Bigl[\,\frac{x^2}{2}-\frac{x^3}{3}\,\Bigr]_0^1=\frac{1}{2}-\frac{1}{3}$\par
\hspace*{2.6\zw}$\displaystyle=\frac{3}{6}-\frac{2}{6}=\frac{1}{6}$\par
\bkstep[bkGreen]{\R{確}{たし}かめ}$y=x$ の下の三角形 $\dfrac{1}{2}$ から、$y=x^2$ の下の $\dfrac{1}{3}$ をひいても $\dfrac{1}{6}$。\par
\kotae $\dfrac{1}{6}$
\end{reidai}

\begin{jibun}
$y=2x$ と $y=x^2$ のグラフで囲まれた部分の面積を求めましょう。
\bkjans{$x=0,\,2$。$\displaystyle\int_0^2(2x-x^2)\,\dd x=4-\dfrac{8}{3}=\dfrac{4}{3}$}
\end{jibun}

%% ---------------- まちがえやすいところ・物理工学 ----------------
\bksection{tsumazuki}{まちがえやすいところをチェック}

\begin{tsumazuki}{不定積分の $+C$ を書き忘れる}
\bkng{$\int2x\,\dd x=x^2$}\bkok{$\int2x\,\dd x=x^2+C$}
$x^2+1$ や $x^2-7$ も、微分すると $2x$ になる正しい答えです。不定積分には必ず $+C$ をつけます。定積分では $C$ が消えるので、いりません。
\end{tsumazuki}

\begin{tsumazuki}{指数をふやしたのに、わるのを忘れる}
\bkng{$\int x^2\,\dd x=x^3+C$}\bkok{$\int x^2\,\dd x=\dfrac{x^3}{3}+C$}
微分して確かめると、$(x^3)'=3x^2$ で、$x^2$ の3倍になってしまいます。不定積分は、いつも微分してもとにもどるかで確かめられます。
\end{tsumazuki}

\begin{tsumazuki}{$F(b)-F(a)$ の順番を逆にする}
\bkng{$\Bigl[\,x^2\,\Bigr]_1^3=1-9=-8$}\bkok{$\Bigl[\,x^2\,\Bigr]_1^3=9-1=8$}
上端（上の数）を入れた値から、下端（下の数）を入れた値をひきます（「あと－まえ」、2-3）。
\end{tsumazuki}

\begin{tsumazuki}{下端の値が0だと思いこむ}
\bkng{$\Bigl[\,x^2+x\,\Bigr]_1^3=9+3=12$}\bkok{$\Bigl[\,x^2+x\,\Bigr]_1^3=12-2=10$}
下端が0でないときは、下端を入れた値も計算して、ひきます。下端が0でも、$F(0)$ が0とは限りません（$F(x)=x+5$ なら $F(0)=5$）。
\end{tsumazuki}

\begin{butsuri}
\begin{itemize}[leftmargin=1.2\zw, itemsep=2pt, topsep=0pt, label={\color{bkPurple}●}]
\item \textbf{ニュートンとライプニッツの発見}\par
毎回細かくたす大変な計算だった積分が、微分の逆として求められるとわかり、物理が一気に進みました。
\item \textbf{位置・速度・加速度}\par
速度を積分すると位置、加速度を積分すると速度。微分と積分を行き来して運動を調べます（5-13）。
\item \textbf{積算計}\par
ガスや水道のメーターは、流れる速さ（流量）を時間で積分して、使った量を表示しています（5-12）。
\end{itemize}
\end{butsuri}

%% ---------------- 章末問題 ----------------
\bkendproblems
\bkproblemgroup{基本}{3}

\bkproblem $\displaystyle\int(4x^3-6x+2)\,\dd x$ を求めましょう。
\bkpans{$x^4-3x^2+2x+C$}

\bkproblem $\displaystyle\int_1^2(3x^2-2x)\,\dd x$ を求めましょう。
\bkpans{$\Bigl[\,x^3-x^2\,\Bigr]_1^2=4-0=4$}

\bkproblem 速度が $v=6t-2$（m/s）の物体が、0秒から3秒までに進んだ距離を求めましょう。
\bkpans{$\Bigl[\,3t^2-2t\,\Bigr]_0^3=\qty{21}{m}$}

\bkproblem $y=4x-x^2$ のグラフと横軸で囲まれた部分の面積を求めましょう。
\bkpans{$\Bigl[\,2x^2-\dfrac{x^3}{3}\,\Bigr]_0^4=\dfrac{32}{3}$}

\bkproblem 速度が $v=2t-6$（m/s）の物体について、0秒から4秒までの変位と道のりを求めましょう。
\bkpans{変位 $\Bigl[\,t^2-6t\,\Bigr]_0^4=\qty{-8}{m}$。\mbox{0〜3秒}で $-9$、\mbox{3〜4秒}で1。道のり\qty{10}{m}}

\bkproblem $y=3x$ と $y=x^2$ のグラフで囲まれた部分の面積を求めましょう。
\bkpans{$\displaystyle\int_0^3(3x-x^2)\,\dd x=\dfrac{27}{2}-9=\dfrac{9}{2}$}

\bkproblem $\displaystyle\int(x^3+6x^2-5)\,\dd x$ を求め、微分して、もとにもどることを確かめましょう。
\bkpans{$\dfrac{x^4}{4}+2x^3-5x+C$（微分すると $x^3+6x^2-5$）}

\bkproblemgroup{チャレンジ}{3}

\bkproblem $S(x)=\displaystyle\int_0^x(t^3-2t)\,\dd t$ のとき、$S'(2)$ を求めましょう。また、$S(x)$ を計算して、微分した結果と同じになることを確かめましょう。
\bkpans{$S'(2)=8-4=4$。$S(x)=\dfrac{x^4}{4}-x^2$ を微分すると $x^3-2x$}

\bkproblem $\displaystyle\int_0^3(x^2-2x)\,\dd x=0$ ですが、グラフと横軸で囲まれた部分の面積は0ではありません。$0\leqq x\leqq2$ と $2\leqq x\leqq3$ に分けて、それぞれの定積分を求めましょう。
\bkpans{$\Bigl[\,\dfrac{x^3}{3}-x^2\,\Bigr]_0^2=-\dfrac{4}{3}$、$\Bigl[\,\dfrac{x^3}{3}-x^2\,\Bigr]_2^3=0-\left(-\dfrac{4}{3}\right)=\dfrac{4}{3}$}

\bkproblem $S(x)=\displaystyle\int_1^x(2t+3)\,\dd t$ を、$x$ の式で表しましょう。また、$S(1)=0$ と、$S'(x)=2x+3$ になることを確かめましょう。
\bkpans{$S(x)=\Bigl[\,t^2+3t\,\Bigr]_1^x=x^2+3x-4$。$S(1)=1+3-4=0$、$S'(x)=2x+3$}

\bkproblem $y=x^2-4$ のグラフと横軸で囲まれた部分の面積を求めましょう。
\bkpans{$\displaystyle\int_{-2}^{2}(x^2-4)\,\dd x=-\dfrac{32}{3}$（横軸より下なのでマイナス）。面積は $\dfrac{32}{3}$}

\begin{furikaeri}
\begin{checklist}
\item 面積のふえる速さが高さであることを説明できる。\hfill{\footnotesize\sffamily\color{bkTheme}→\pageref{bk:5-11:kihon}ページ}
\item 不定積分を $+C$ をつけて書ける。\hfill{\footnotesize\sffamily\color{bkTheme}→\pageref{bk:5-11:genshi}ページ}
\item 定積分を $F(b)-F(a)$ で計算できる。\hfill{\footnotesize\sffamily\color{bkTheme}→\pageref{bk:5-11:teisekibun}ページ}
\end{checklist}
\end{furikaeri}

\bkfukidashi{sora}{積分がかんたんになった！　次は、物理で積分を使ってみよう。}

%% 5-12 物理での積分：仕事と流量（第11回：ページ制限なし。積分の3段「分ける・かける・たす」を、量ごとに図と数で見せる）
\bkchapter{5}{12}{物理での積分：仕事と流量}{\bklevel{3}}{物理での積分：仕事と流量}
\bkmichishirube{2-15,4-12,5-11}{5-13,6-15}

\begin{goal}
\begin{itemize}[label={\color{bkTheme}●}, leftmargin=1.6\zw, itemsep=0pt, topsep=1pt]
\item \bklevel{3}\enspace 力が変わるときの仕事を、$F$-$x$ グラフの面積（積分）で求められる。
\item \bklevel{3}\enspace 流量や電力を時間で積分して、たまった量を求められる。
\item \bklevel{3}\enspace 積分を使って、平均の値を求められる。
\end{itemize}
\end{goal}

\begin{hatena}
\noindent\begin{minipage}[c]{0.58\linewidth}\raggedright
ばねをのばすとき、はじめは軽い力でのびますが、のばすほど強い力が必要になります。ばねの力は、のびに比例するからです（2-15）。

ばねを\qty{10}{cm}のばすのに必要な「仕事」は、いくらでしょう。仕事は「力×距離」で求めます。でも、力が変わり続けるので、どの力を使えばよいかわかりません。
\end{minipage}\hfill
\begin{minipage}[c]{0.38\linewidth}
\begin{tikzpicture}[x=1mm, y=1mm]
  \fill[bkGray!40] (-2,-6) rectangle (0,6);
  \draw[bkInk, line width=0.8pt, decorate, decoration={coil, aspect=0.5, segment length=2.4mm, amplitude=2.2mm}] (0,0) -- (26,0);
  \fill[bkTheme!40] (26,-4) rectangle (32,4);
  \draw[bkInk, line width=0.6pt] (26,-4) rectangle (32,4);
  \draw[bk arrow, bkOrange] (33,0) -- (44,0);
  \node[font=\sffamily\scriptsize, anchor=south] at (39,1) {$F=kx$};
  \draw[bkGray, <->, line width=0.5pt] (18,-8) -- (26,-8);
  \node[font=\sffamily\scriptsize, anchor=north] at (22,-9) {のび $x$};
\end{tikzpicture}
\end{minipage}
\end{hatena}

5-10と5-11で学んだ積分は、「細かく分けて、かけて、たす」計算でした。速度から距離を求めるだけでなく、物理のいろいろな「たまる量」を計算する道具です。この章では、仕事、水の量、電気のエネルギー、電気の量を、この3つの段で求めます。どの量でも、同じ3つの段をくり返すだけです。

\clearpage
%% ---------------- やってみよう・思い出そう ----------------
\begin{yattemiyou}[ばねばかりで、のばす力をはかろう]
\textbf{用意するもの}\par
ばねばかり（\qty{20}{N}まで）、ばね、ものさし
\begin{enumerate}[label={\bkcircled[bkGreen]{\arabic*}}, leftmargin=2.4\zw, itemsep=1pt, topsep=3pt]
\item ばねの一方を固定し、もう一方をばねばかりで引く。のびを\qty{2}{cm}ずつふやしながら、力を読もう。
\item 横軸をのび、たて軸を力にしてグラフをかこう（$F$-$x$ グラフ）。
\end{enumerate}
\begin{bkcenter}
\begin{bktable}{colspec={Q[l,m,wd=30mm] *{6}{Q[c,m,wd=10mm]}}, rowsep=1.5pt}
のび $x$（cm） & 0 & 2 & 4 & 6 & 8 & 10 \\
力 $F$（N） & 0 & 4 & 8 & 12 & 16 & 20 \\
\end{bktable}
\bktabcaption{ある班の結果。力はのびに比例している。}
\end{bkcenter}
\end{yattemiyou}

\noindent \qty{2}{cm}のびるごとに、力は\qty{4}{N}ずつふえています。のびを m（メートル）で表すと、\qty{2}{cm}は
\[ \qty{2}{cm}\div\qty{100}{cm/m}=\qty{0.02}{m} \]
です（\qty{100}{cm/m} は「\qty{1}{m}あたり\qty{100}{cm}」。2-10）。だから、のび\qty{1}{m}あたりの力（ばね定数 $k$）は
\[ k=\qty{4}{N}\div\qty{0.02}{m}=\qty{200}{N/m} \]
です。このばねの力は $F=200x$（$x$ は m、$F$ は N）と書けます。たとえば $x=0.1$ なら $F=200\times0.1=20$ で、表の\qty{20}{N}と合います。\bkidx{ばねていすう}{ばね定数}

\begin{omoidasou}
\bkrecall{4-12}{物を力の向きに動かしたとき、仕事 $W$＝力 $F$×動いた距離 $s$。単位は J（ジュール）。\qty{1}{N}の力で\qty{1}{m}動かしたときの仕事が\qty{1}{J}（1 J＝1 N·m）。}
\bkrecall{2-15}{ばねの力は、のびに比例する：$F=kx$。$k$ はばね定数（N/m）。}
\bkrecall{5-10}{定積分 $\displaystyle\int_a^b f(x)\,\dd x$ は、高さ $f(x)$、はば $\Delta x$ の細い長方形を $a$ から $b$ までたし、はばを0に近づけた先。グラフの下の面積。}
\bkrecall{5-11}{$\displaystyle\int_a^b f(x)\,\dd x=\Bigl[\,F(x)\,\Bigr]_a^b=F(b)-F(a)$（$F$ は微分すると $f$ になる関数）。$\displaystyle\int x\,\dd x=\frac{x^2}{2}+C$}
\bkrecall{2-10}{換算は「1をかける」：\qty{60}{s/min}、\qty{3600}{s/h}、\qty{1000}{W/kW}。単位をそろえてから計算する。}
\end{omoidasou}

%% ---------------- 数学の道具 ----------------
\bksection{doguu}{数学の道具}

\bksubsection{1}{力が一定なら、仕事は長方形の面積}

\label{bk:5-12:shigoto}物を動かすたいへんさは、力が大きいほど、動かす距離が長いほど大きくなります。たとえば、重い荷物を2つ運ぶのは1つのときの2倍たいへんで、2倍の距離を運ぶのも2倍たいへんです。そこで物理では、
\[ \text{仕事}=\text{力}\times\text{力の向きに動いた距離},\qquad W=Fx \]
と決めています（4-12）。力が\qty{1}{N}、距離が\qty{1}{m}なら\qty{1}{J}です。これを\textbf{仕事}といいます。\bkidx{しごと}{仕事}

\noindent\begin{minipage}[c]{0.46\linewidth}
\begin{tikzpicture}
\begin{axis}[bkplot, width=56mm, height=40mm, xmin=0, xmax=2, ymin=0, ymax=26, xtick={1.5}, ytick={19.6}, xlabel={$x$（m）}, ylabel={$F$（N）}]
  \fill[bkOrange!25] (axis cs:0,0) rectangle (axis cs:1.5,19.6);
  \addplot[bkcurve, line width=1.4pt] coordinates {(0,19.6) (1.5,19.6)};
  \node[font=\sffamily\small, fill=bkOrange!25, inner sep=1pt] at (axis cs:0.75,9) {\qty{29.4}{J}};
\end{axis}
\end{tikzpicture}
\end{minipage}\hfill
\begin{minipage}[c]{0.52\linewidth}\raggedright
横軸を位置 $x$、たて軸を力 $F$ にした $F$-$x$ グラフでは、力が一定なら、仕事はたて $F$、横 $x$ の\textbf{長方形の面積}です。$v$-$t$ グラフの長方形の面積が距離になったのと同じ形です（5-9）。単位も N×m＝J になります。
\end{minipage}

\noindent\textbf{例}\quad 質量\qty{2}{kg}の荷物を、ゆっくり一定の速さで\qty{1.5}{m}持ち上げます。地球は、質量\qty{1}{kg}あたり約\qty{9.8}{N}の力（重力）で、物を下に引いています。この9.8は、5-5で調べた落ちる物の加速度 $g=\qty{9.8}{m/s^2}$ と同じ数です。だから荷物にはたらく重力は
\[ F=mg=\qty{2}{kg}\times\qty{9.8}{N/kg}=\qty{19.6}{N} \]
です。ゆっくり一定の速さで持ち上げるときは、この重力と同じ大きさの力で上に引けばよいので、力は\qty{19.6}{N}で一定です。
\[ W=Fx=19.6\times1.5=\qty{29.4}{J} \]
上の図の長方形の面積です。

\bksubsection{2}{力が変わるときは、細かく分けて、かけて、たす}

\label{bk:5-12:sekibun}\bkyattemiyouref のばね（$F=200x$）を、0から\qty{0.1}{m}までのばす仕事を考えます。力が変わるので、「力×距離」をそのまま使えません。そこで、5-10と同じ3つの段で考えます。

\noindent\textbf{段1　細かく分ける}\quad 0から\qty{0.1}{m}までを、5つに等しく分けます。1つのはばは $0.1\div5=\qty{0.02}{m}$ です（長さ0.1を5こに等しく分けるから）。

\noindent\textbf{段2　かける}\quad 細い区間の中では、力はほとんど変わりません。そこで、区間の右はしの力を、その区間の力とみなします。たとえば2番めの区間（0.02から\qty{0.04}{m}）の右はしは $x=0.04$ なので、力は $200\times0.04=\qty{8}{N}$。この区間の小さな仕事は、力×はばで $8\times0.02=\qty{0.16}{J}$ です。

\noindent\textbf{段3　たす}\quad 5つの小さな仕事をたします。

\begin{bkcenter}
\begin{bktable}{colspec={Q[l,m,wd=32mm] *{5}{Q[c,m,wd=12mm]} Q[c,m,wd=12mm]}, rowsep=1.5pt, colsep=4pt, row{1}={bg=bkTheme!10}}
区間の番号 & 1 & 2 & 3 & 4 & 5 & 合計 \\
右はしの位置（m） & 0.02 & 0.04 & 0.06 & 0.08 & 0.1 & \\
右はしの力（N） & 4 & 8 & 12 & 16 & 20 & \\
小さな仕事（J） & 0.08 & 0.16 & 0.24 & 0.32 & 0.40 & 1.2 \\
\end{bktable}
\bktabcaption{右はしの力を使った和。小さな仕事＝右はしの力×0.02。}
\end{bkcenter}

\begin{bkfigure}
\begin{tikzpicture}
\begin{axis}[bkplot, name=s1, width=60mm, height=44mm, xmin=0, xmax=0.115, ymin=0, ymax=24, xtick={0.02,0.04,0.06,0.08,0.1}, xticklabels={,0.04,,0.08,}, ytick={10,20}, xlabel={$x$（m）}, ylabel={$F$（N）}, title={\sffamily\small 右はしの力（大きめ）}, scaled x ticks=false]
  \foreach \k in {1,...,5} {\edef\tmp{\noexpand\fill[bkOrange!30, draw=bkOrange!70, line width=0.5pt] (axis cs:{(\k-1)*0.02},0) rectangle (axis cs:{\k*0.02},{4*\k});}\tmp}
  \addplot[bkcurve, domain=0:0.11] {200*x};
\end{axis}
\begin{axis}[bkplot, at={(s1.east)}, anchor=west, xshift=14mm, width=60mm, height=44mm, xmin=0, xmax=0.115, ymin=0, ymax=24, xtick={0.02,0.04,0.06,0.08,0.1}, xticklabels={,0.04,,0.08,}, ytick={10,20}, xlabel={$x$（m）}, ylabel={$F$（N）}, title={\sffamily\small 左はしの力（小さめ）}, scaled x ticks=false]
  \foreach \k in {1,...,5} {\edef\tmp{\noexpand\fill[bkBlue!25, draw=bkBlue!70, line width=0.5pt] (axis cs:{(\k-1)*0.02},0) rectangle (axis cs:{\k*0.02},{4*(\k-1)});}\tmp}
  \addplot[bkcurve, domain=0:0.11] {200*x};
\end{axis}
\end{tikzpicture}
\bkfigcaption{長方形の面積が小さな仕事。和は、右はしの力で1.2 J、左はしで0.8 J。}
\end{bkfigure}

右はしの力は、その区間でいちばん大きい力なので、1.2 J は少し大きすぎます。左はしの力（いちばん小さい力：0、4、8、12、16 N）を使うと、$0+0.08+0.16+0.24+0.32=\qty{0.8}{J}$ で、少し小さすぎます。ほんとうの仕事は、そのあいだです。分け方を細かくしていくと、次のようになります。

\begin{bkcenter}
\begin{bktable}{colspec={Q[l,m,wd=34mm] *{4}{Q[c,m,wd=17mm]}}, rowsep=1.5pt}
分ける数 & 5 & 10 & 100 & 1000 \\
はば（m） & 0.02 & 0.01 & 0.001 & 0.0001 \\
右はしの和（J） & 1.2 & 1.1 & 1.01 & 1.001 \\
左はしの和（J） & 0.8 & 0.9 & 0.99 & 0.999 \\
\end{bktable}
\bktabcaption{細かく分けるほど、どちらの和も1に近づく。}
\end{bkcenter}

2つの和は、上と下から1にはさまれて近づいていきます。だから、このばねを\qty{0.1}{m}のばす仕事は、ちょうど\qty{1}{J}です。速度から距離を求めたとき（5-10）と、まったく同じ考え方です。

\bksubsection{3}{記号で書くと $\displaystyle\int F\,\dd x$}

\label{bk:5-12:kigou}②の「細かく分けて、かけて、たして、はばを0に近づけた先」を、5-10の定積分の記号で書きます。

\[ W=\int_a^b F(x)\,\dd x \]
$F$-$x$ グラフの、$a$ から $b$ までの面積が仕事です。記号の1つ1つの意味は、下の表のとおりです。

\noindent\textbf{なぜ $\dd x$ を書くの？}\quad 仕事は「力×距離」でした。$F(x)$ が力で、$\dd x$ が（とても小さな）距離です。$\dd x$ を書かないと、かけ算の相手の「距離」がなくなってしまいます。単位も、$F(x)$ の N に $\dd x$ の m がかかって、N×m＝J になります。$\dd x$ は「何をはばにして細かく分けたか」も表しています。ここでは位置 $x$ で分けたので $\dd x$ です。

\begin{bkcenter}
\begin{bktable}{colspec={Q[c,m,wd=26mm] Q[l,m,wd=84mm]}, rowsep=2pt, row{1}={bg=bkTheme!10}}
記号 & 意味 \\
$\displaystyle\int$ & たし合わせる（「和」の sum の S を、たてにのばした形） \\
下の $a$、上の $b$ & はじめの位置（下端）と、おわりの位置（上端） \\
$F(x)$ & 位置 $x$ での力（細い長方形の高さ） \\
$\dd x$ & 小さな距離（細い長方形のはば） \\
$F(x)\,\dd x$ & 小さな仕事（力×小さな距離＝細い長方形の面積） \\
\end{bktable}
\bktabcaption{$\int$ は「インテグラル」と読む（5-10）。}
\end{bkcenter}

\bksubsection{4}{ばねの仕事は $\dfrac{1}{2}kx^2$}

\label{bk:5-12:bane}ばね定数 $k$ のばねを、0から $x$ までのばす仕事を、3つの方法で求めます。どれも同じ答えになります。

\noindent\begin{minipage}[c]{0.44\linewidth}
\begin{tikzpicture}
\begin{axis}[bkplot, width=54mm, height=42mm, xmin=0, xmax=0.125, ymin=0, ymax=25, xtick={0.1}, xticklabels={$x$}, ytick={10,20}, yticklabels={$\frac{kx}{2}$,$kx$}, xlabel={のび}, ylabel={$F$}, scaled x ticks=false]
  \addplot[fill=bkOrange!25, draw=none, domain=0:0.1] {200*x} \closedcycle;
  \addplot[bkcurve, domain=0:0.115] {200*x};
  \draw[bkGreen!70!black, line width=1pt, dashed] (axis cs:0,10) -- (axis cs:0.1,10);
  \node[font=\sffamily\small] at (axis cs:0.072,4.5) {$W$};
  \draw[bkGray, dashed, line width=0.5pt] (axis cs:0.1,0) -- (axis cs:0.1,20);
\end{axis}
\end{tikzpicture}
\end{minipage}\hfill
\begin{minipage}[c]{0.54\linewidth}\raggedright
\textbf{方法1　三角形の面積}\par
$F=kx$ のグラフはまっすぐなので、面積は底辺 $x$、高さ $kx$ の三角形です。
\[ W=\frac{1}{2}\times x\times kx=\frac{1}{2}kx^2 \]
$\dfrac{1}{2}$ は、三角形の面積の「÷2」から来ています。
\end{minipage}

\noindent\textbf{方法2　積分}\quad 力 $ks$（$s$ はとちゅうののび）を、0から $x$ まで積分します。$s$ は、0から $x$ まで動くとちゅうののびです。上端の $x$ と区別するため、別の文字にしています。
\[ W=\int_0^x ks\,\dd s=\Bigl[\,k\times\frac{s^2}{2}\,\Bigr]_0^x=k\times\frac{x^2}{2}-0=\frac{1}{2}kx^2 \]
$\dfrac{1}{2}$ は、こんどは「$s$ を積分すると $\dfrac{s^2}{2}$」（5-11：指数を1ふやして、その数でわる）から出てきました。

\noindent\textbf{方法3　平均の力}\quad 力は0から $kx$ まで、のびに比例して一様にふえます。だから、平均の力は $(0+kx)\div2=\dfrac{kx}{2}$（グラフの緑の点線）です。平均の力で $x$ だけ動かしたと考えると、$\dfrac{kx}{2}\times x=\dfrac{1}{2}kx^2$。三角形を、同じ面積の長方形にならしたことになります。

\noindent\textbf{数で確かめる}\quad \bkyattemiyouref のばね（$k=200$）を $x=0.1$ だけのばすと、
\[ W=\frac{1}{2}\times200\times0.1^2=\frac{1}{2}\times200\times0.01=\frac{1}{2}\times2=\qty{1}{J} \]
②の表で近づいていった1と、ぴったり合います。これが\bkhatenaref の答えです。

\noindent\textbf{まちがえた考え方}\quad 最後の力\qty{20}{N}を使って $20\times0.1=\qty{2}{J}$ とすると、ほんとうの2倍になります。はじめの力0を使えば0 J です。力が変わるときに「どこか1つの力×距離」とすると、まちがえるのです。

のばしたばねは、この仕事の分のエネルギーをたくわえます。手をはなすと、このエネルギーで物をおし出します。これを\textbf{弾性エネルギー}といいます。\bkidx{だんせいえねるぎー}{弾性エネルギー $\frac{1}{2}kx^2$}\bkidx{えねるぎー}{エネルギー}

\bksubsection{5}{流量を積分すると、たまった水の量}

\label{bk:5-12:ryuuryou}じゃ口から1分あたりに出る水の量を\textbf{流量}といい、$Q$（L/min）で表します。L/min は「1分あたり何 L」です。流量が一定なら、たまる水の量は「流量×時間」です。たとえば\qty{10}{L/min}で3分なら、$\qty{10}{L/min}\times\qty{3}{min}=\qty{30}{L}$（min どうしが約分されて L が残る）。速さ×時間＝距離と同じ形です。\bkidx{りゅうりょう}{流量}

\noindent\textbf{流量が変わるとき}\quad ふろのじゃ口を少しずつ開けていくと、流量は $Q=2t$ でふえていきました。$Q$ の単位は L/min、$t$ は開け始めてからの時間（分）です。6分後の流量は、$2\times6=12$ より\qty{12}{L/min}です。6分間で、何 L たまったでしょう。ここでも3つの段を使います。

\begin{enumerate}[label={\bkcircled[bkTheme]{\arabic*}}, leftmargin=2.4\zw, itemsep=1pt, topsep=2pt]
\item \textbf{分ける}：0分から6分までを、1分ずつ6つに分ける。
\item \textbf{かける}：1分のあいだは流量がほぼ一定とみて、右はしの流量×1分。1番めは $2\times1\times1=\qty{2}{L}$、2番めは $2\times2\times1=\qty{4}{L}$、…、6番めは\qty{12}{L}。
\item \textbf{たす}：$2+4+6+8+10+12=\qty{42}{L}$（右はし。大きめ）。左はしの流量では $0+2+4+6+8+10=\qty{30}{L}$（小さめ）。
\end{enumerate}

\noindent\begin{minipage}[c]{0.46\linewidth}
\begin{tikzpicture}
\begin{axis}[bkplot, width=50mm, height=42mm, xmin=0, xmax=7, ymin=0, ymax=14.5, xtick={1,2,...,6}, ytick={4,8,12}, xlabel={$t$（min）}, ylabel={$Q$（L/min）}]
  \foreach \k in {1,...,6} {\edef\tmp{\noexpand\fill[bkBlue!18, draw=bkBlue!60, line width=0.4pt] (axis cs:\k-1,0) rectangle (axis cs:\k,{2*\k});}\tmp}
  \addplot[fill=bkBlue!35, draw=none, domain=0:6] {2*x} \closedcycle;
  \addplot[bkcurve, domain=0:6.6] {2*x};
  \node[font=\sffamily\small, fill=bkBlue!35, inner sep=1pt] at (axis cs:4.3,3) {\qty{36}{L}};
\end{axis}
\end{tikzpicture}
\end{minipage}\hfill
\begin{minipage}[c]{0.52\linewidth}\raggedright
ほんとうの水の量は、42 L と 30 L のあいだです。はばを0に近づけた先は、$Q$-$t$ グラフの下の面積（こい青の三角形）で、
\[ V=\int_0^6 2t\,\dd t=\Bigl[\,t^2\,\Bigr]_0^6=\qty{36}{L} \]
（$t^2$ を微分すると $2t$ なので、$2t$ の原始関数は $t^2$。5-11）三角形の面積でも $\dfrac{1}{2}\times6\times12=36$ で、同じです。
\end{minipage}

\noindent 一般に、時刻 $t_1$ から $t_2$ までにたまる水の量 $V$ は
\[ V=\int_{t_1}^{t_2} Q(t)\,\dd t \]
です。こんどは時間で細かく分けたので、$\dd x$ ではなく $\dd t$（小さな時間）を書きます。$Q(t)\,\dd t$ は「流量×小さな時間＝小さな水の量」です。

\bksubsection{6}{電力を積分すると、電力量}

\label{bk:5-12:denryoku}電気器具が1秒あたりに使うエネルギーを\textbf{電力}といい、単位は W（ワット）です。1秒に\qty{1}{J}使うのが\qty{1}{W}なので、W＝J/s です。たとえば\qty{10}{W}のLED電球は、1秒に\qty{10}{J}のエネルギーを使います。1時間（\qty{3600}{s}）つけておくと、
\[ \qty{10}{J/s}\times\qty{3600}{s}=\qty{36000}{J} \]
です。電力×時間で求めた、使ったエネルギーを\textbf{電力量}といいます。電力が変わるときは、流量のときと同じく、時間で細かく分けてたします。\bkidx{でんりょくりょう}{電力量}
\[ E=\int_{t_1}^{t_2}P(t)\,\dd t \]

\noindent\textbf{kWh（キロワット時）}\quad J では数が大きくなりすぎるので、家庭では電力量の単位に kWh を使います。\qty{1}{kW}で1時間使ったときの電力量が\qty{1}{kWh}です。J に直すには、kW を W に、h を s に換算します。
\[ \qty{1}{kWh}=\qty{1}{kW}\times\qty{1000}{W/kW}\times\qty{1}{h}\times\qty{3600}{s/h}=\qty{3600000}{J}=\qty{3.6e6}{J} \]
（kW と h がそれぞれ約分されて、W×s＝J が残ります。）電力に kW、時間に h（時間）を使って積分すれば、答えは kWh で出てきます。電力に W を使えば、答えは Wh（ワット時。\qty{1}{Wh}は\qty{1}{W}で1時間使った電力量で、\qty{1}{kWh}の $\dfrac{1}{1000}$）です。

\noindent\textbf{例}\quad 電気自動車の充電器の電力は、充電が進むとへっていき、$P=6-t$（kW、$t$ は時間）でした。はじめの2時間の電力量を求めます。$6$ を積分すると $6t$、$t$ を積分すると $\dfrac{t^2}{2}$ なので（5-11）、

\noindent\begin{minipage}[c]{0.44\linewidth}
\begin{tikzpicture}
\begin{axis}[bkplot, width=52mm, height=40mm, xmin=0, xmax=3, ymin=0, ymax=7.5, xtick={1,2}, ytick={4,6}, xlabel={$t$（h）}, ylabel={$P$（kW）}]
  \addplot[fill=bkYellow!45, draw=none, domain=0:2] {6-x} \closedcycle;
  \addplot[bkcurve, domain=0:2.8] {6-x};
  \draw[bkGray, dashed, line width=0.5pt] (axis cs:2,0) -- (axis cs:2,4);
  \node[font=\sffamily\small, fill=bkYellow!45, inner sep=1pt] at (axis cs:1,2) {\qty{10}{kWh}};
\end{axis}
\end{tikzpicture}
\end{minipage}\hfill
\begin{minipage}[c]{0.54\linewidth}\raggedright
\[ \int_0^2(6-t)\,\dd t=\Bigl[\,6t-\frac{t^2}{2}\,\Bigr]_0^2 \]
\[ =\Bigl(12-\frac{4}{2}\Bigr)-0=12-2=\qty{10}{kWh} \]
グラフは $t=0$ で6、$t=2$ で4の直線なので、黄色の台形の面積 $(6+4)\times2\div2=10$ でも確かめられます。
\end{minipage}

\bksubsection{7}{電流を積分すると、流れた電気の量}

\label{bk:5-12:denryuu}電流は、導線のある所を1秒あたりに通りぬける電気の量です。電気の量の単位を C（クーロン）といい、\qty{1}{A}（アンペア）の電流が1秒流れると、\qty{1}{C}の電気が通りぬけます。つまり A＝C/s で、電流は「1秒あたりの電気の量」です。だから、電流 $I$ を時間で積分すると、流れた電気の量 $q$ になります。
\[ q=\int_{t_1}^{t_2}I(t)\,\dd t \]
\noindent\begin{minipage}[c]{0.44\linewidth}
\begin{tikzpicture}
\begin{axis}[bkplot, width=52mm, height=38mm, xmin=0, xmax=13, ymin=0, ymax=0.8, xtick={5,10}, ytick={0.5}, yticklabels={0.5}, xlabel={$t$（s）}, ylabel={$I$（A）}]
  \fill[bkGreen!25] (axis cs:0,0) rectangle (axis cs:10,0.5);
  \addplot[bkcurve, line width=1.4pt] coordinates {(0,0.5) (12,0.5)};
  \draw[bkGray, dashed, line width=0.5pt] (axis cs:10,0) -- (axis cs:10,0.5);
  \node[font=\sffamily\small, fill=bkGreen!25, inner sep=1pt] at (axis cs:5,0.25) {\qty{5}{C}};
\end{axis}
\end{tikzpicture}
\end{minipage}\hfill
\begin{minipage}[c]{0.54\linewidth}\raggedright
\noindent\textbf{例}\quad 電流が\qty{0.5}{A}で一定なら、10秒で
\[ \qty{0.5}{C/s}\times\qty{10}{s}=\qty{5}{C} \]
です。$I$-$t$ グラフでは、たて0.5、横10の長方形（緑）の面積です。電流が変わるときも、$I$-$t$ グラフの下の面積が、流れた電気の量です（\bkreidairef{7}）。
\end{minipage}

\bksubsection{8}{平均の値}

\label{bk:5-12:heikin}自転車が、0秒から2秒まで $v=3t^2$（m/s）で速くなりました。この2秒の平均の速さはいくつでしょう。

\noindent\textbf{方法1　たくさんの値の平均}\quad 2-6の平均は「合計÷個数」でした。2秒を4つに分け、それぞれのまん中の時刻（0.25、0.75、1.25、1.75秒）の速さをはかると、0.1875、1.6875、4.6875、9.1875です。平均は
\[ (0.1875+1.6875+4.6875+9.1875)\div4=15.75\div4=3.9375 \]
8つに分けると3.984…、もっと細かくすると4に近づきます。

\noindent\textbf{方法2　面積を長方形にならす}\quad 進んだ距離は $\displaystyle\int_0^2 3t^2\,\dd t=\Bigl[\,t^3\,\Bigr]_0^2=8$（m）です。2秒で\qty{8}{m}進んだので、平均の速さは $8\div2=\qty{4}{m/s}$。5-1の「全部の距離÷全部の時間」です。

\noindent\begin{minipage}[c]{0.53\linewidth}\centering
\begin{tikzpicture}
\begin{axis}[bkplot, width=60mm, height=38mm, xmin=0, xmax=2.4, ymin=0, ymax=13, xtick={1,2}, ytick={4,12}, xlabel={$t$（s）}, ylabel={$v$（m/s）}]
  \fill[bkOrange!25] (axis cs:0,0) rectangle (axis cs:2,4);
  \addplot[fill=bkBlue!20, fill opacity=0.7, draw=none, domain=0:2] {3*x^2} \closedcycle;
  \addplot[bkcurve, domain=0:2.08] {3*x^2};
  \draw[bkOrange, line width=1.2pt] (axis cs:0,4) -- (axis cs:2,4);
  \node[font=\sffamily\scriptsize, anchor=west, text=bkOrange!80!black] at (axis cs:2.08,4) {平均 4};
\end{axis}
\end{tikzpicture}
\bkfigcaption{青とオレンジの面積は、どちらも8。}
\end{minipage}\hfill
\begin{minipage}[c]{0.44\linewidth}\raggedright
グラフでは、平均の値は、もとのグラフの下の面積と同じ面積をもつ長方形の高さです。でこぼこの砂をならして、平らにした高さと同じです。
\end{minipage}

\noindent\textbf{2つの方法がつながるわけ}\quad 方法1で、はば $\Delta t$ の区間に $n$ こ分けたとき、値の合計÷個数の分子と分母に $\Delta t$ をかけると（分数の大きさは変わらない）、
\[ \frac{\text{値の合計}}{n}=\frac{\text{値の合計}\times\Delta t}{n\times\Delta t}=\frac{\text{細い長方形の面積の和}}{\text{全体のはば}} \]
になります。はばを0に近づけると、分子は積分になります。だから
\[ \text{平均の値}=\frac{1}{b-a}\int_a^b f(t)\,\dd t \]
です。$b-a$ でわるのは、「個数でわる」のかわりに「全体のはば（時間）でわる」からです。

\begin{bkcenter}
\begin{bktable}{colspec={Q[l,m,wd=36mm] Q[l,m,wd=36mm] Q[c,m,wd=36mm]}, rowsep=1pt, row{1}={bg=bkTheme!10}}
1あたりの量 & 積分すると（たまる量） & 単位の計算 \\
速度 $v$（m/s） & 位置の変化（m） & (m/s)×s＝m \\
力 $F$（N） & 仕事 $W$（J） & N×m＝J \\
流量 $Q$（L/min） & 水の量 $V$（L） & (L/min)×min＝L \\
電力 $P$（W） & 電力量 $E$（J） & (J/s)×s＝J \\
電流 $I$（A） & 電気の量 $q$（C） & (C/s)×s＝C \\
\end{bktable}
\bktabcaption{積分は「1あたりの量」から「たまる量」を求める道具。}
\end{bkcenter}

単位の計算では、かけ算で「あたり」（/s や /min）が消えます。力だけは「1あたり」に見えませんが、仕事を位置で細かく分けると「\qty{1}{m}動かすあたりの仕事」が力になっていて、同じ形です。

\begin{column}{スマートフォンの電池の「mAh」}
電池の箱の「\qty{3000}{mAh}（ミリアンペア時）」は、電流（mA）を時間（h）で積分した、とり出せる電気の量です。\qty{1}{mA}の電流が1時間流れた電気の量が\qty{1}{mAh}です。\qty{300}{mA}の電流で使い続けると、
\[ \qty{3000}{mAh}\div\qty{300}{mA}=\qty{10}{h} \]
で、10時間使えます。電流が変わっても、積分した値が\qty{3000}{mAh}になるまで使えます。\qty{1}{mAh}を C に直すと、$\qty{0.001}{A}\times\qty{3600}{s}=\qty{3.6}{C}$ です。
\end{column}

\noindent\textbf{3つの段のまとめ}\quad この章で求めた量は、どれも同じ3つの段で求めています。ちがうのは、何で細かく分けるか（位置か時間か）と、何をかけるかだけです。
\begin{bkcenter}
\begin{bktable}{colspec={Q[l,m,wd=22mm] Q[l,m,wd=30mm] Q[c,m,wd=30mm] Q[c,m,wd=26mm]}, rowsep=2pt, row{1}={bg=bkTheme!10}}
求める量 & 段1　分ける & 段2　かける & 段3　たす \\
仕事 & 位置を細かく（$\dd x$） & 力×小さな距離 $F\,\dd x$ & $\displaystyle\int F\,\dd x$ \\
水の量 & 時間を細かく（$\dd t$） & 流量×小さな時間 $Q\,\dd t$ & $\displaystyle\int Q\,\dd t$ \\
電力量 & 時間を細かく（$\dd t$） & 電力×小さな時間 $P\,\dd t$ & $\displaystyle\int P\,\dd t$ \\
電気の量 & 時間を細かく（$\dd t$） & 電流×小さな時間 $I\,\dd t$ & $\displaystyle\int I\,\dd t$ \\
距離（5-10） & 時間を細かく（$\dd t$） & 速度×小さな時間 $v\,\dd t$ & $\displaystyle\int v\,\dd t$ \\
\end{bktable}
\bktabcaption{どの量も「分ける・かける・たす」。たした先が定積分。}
\end{bkcenter}

\begin{taisetsu}
変わる量の積分 ＝ たまる量（細かく分けて、かけて、たす）\\
\quad 仕事 $\displaystyle\int F\,\dd x$、水の量 $\displaystyle\int Q\,\dd t$、電力量 $\displaystyle\int P\,\dd t$、電気の量 $\displaystyle\int I\,\dd t$\\
ばねの仕事 $\dfrac{1}{2}kx^2$\qquad 平均の値 $=\dfrac{1}{b-a}\displaystyle\int_a^b f\,\dd t$
\end{taisetsu}

\begin{kotoba}
\small
\bkword{仕事}{力×力の向きに動いた距離。力が変わるときは $\int F\,\dd x$。単位は J。}
\bkword{流量・電力・電流}{1分（1秒）あたりの水の量・エネルギー・電気の量。積分すると、たまった量・使った量。}
\bkword{電力量}{電力を時間で積分したもの。kWh や J で表す。}
\end{kotoba}

%% ---------------- 例題 ----------------
\bksection{reidai}{例題でたしかめよう}

\begin{reidai}[一定の力の仕事]
質量\qty{5}{kg}の荷物を、ゆっくり一定の速さで\qty{2}{m}の高さまで持ち上げる仕事を求めましょう（重力は1 kg あたり\qty{9.8}{N}）。\par
\kangaekata 持ち上げる力は荷物の重力と同じで、一定です。だから①の「力×距離」を使います。\par
\shiki 力\quad $F=\qty{5}{kg}\times\qty{9.8}{N/kg}=\qty{49}{N}$\par
\hspace*{2.6\zw}仕事\quad $W=Fx=49\times2=98$\par
\kotae \qty{98}{J}
\end{reidai}

\begin{jibun}
\qty{30}{N}の力で、台車を力の向きに\qty{4}{m}おしました。仕事を求めましょう。
\bkjans{$30\times4=\qty{120}{J}$}
\end{jibun}

\noindent\begin{minipage}[c]{0.46\linewidth}
\begin{tikzpicture}
\begin{axis}[bkplot, width=52mm, height=40mm, xmin=0, xmax=2.6, ymin=0, ymax=62, xtick={1,2}, ytick={49}, xlabel={$x$（m）}, ylabel={$F$（N）}]
  \fill[bkOrange!25] (axis cs:0,0) rectangle (axis cs:2,49);
  \addplot[bkcurve, line width=1.4pt] coordinates {(0,49) (2,49)};
  \node[font=\sffamily\small, fill=bkOrange!25, inner sep=1pt] at (axis cs:1,22) {\qty{98}{J}};
\end{axis}
\end{tikzpicture}
\end{minipage}\hfill
\begin{minipage}[c]{0.52\linewidth}\raggedright
\bkreidairef{1}を $F$-$x$ グラフで見ると、たて\qty{49}{N}、横\qty{2}{m}の長方形です。面積 $49\times2=98$ が仕事です。力が一定のときは、細かく分けなくても、長方形1つで求められます。
\end{minipage}

\begin{reidai}[ばねを、とちゅうからのばす仕事]
ばね定数 $k=\qty{200}{N/m}$ のばねを、のび\qty{0.1}{m}から\qty{0.2}{m}までのばす仕事を求めましょう。\par
\kangaekata 力 $F=200x$ はのびとともに変わるので、③の $\int F\,\dd x$ を使います。はじめの位置0.1が下端、おわりの位置0.2が上端です。\par
\shiki $\displaystyle\int_{0.1}^{0.2}200x\,\dd x=\Bigl[\,200\times\frac{x^2}{2}\,\Bigr]_{0.1}^{0.2}=\Bigl[\,100x^2\,\Bigr]_{0.1}^{0.2}$\par
\hspace*{2.6\zw}$=100\times0.2^2-100\times0.1^2=100\times0.04-100\times0.01=4-1=3$\par
\hspace*{2.6\zw}確かめ：台形の面積（力は20 Nから40 N）\par
\hspace*{2.6\zw}$(20+40)\times0.1\div2=3$\par
\kotae \qty{3}{J}（0から\qty{0.1}{m}のときの3倍。のびているほど力が大きいため）
\end{reidai}

\begin{jibun}
ばね定数 $k=\qty{500}{N/m}$ のばねを、のび0から\qty{0.2}{m}までのばす仕事を求めましょう。
\bkjans{$\dfrac{1}{2}\times500\times0.2^2=\dfrac{1}{2}\times500\times0.04=\qty{10}{J}$}
\end{jibun}

\begin{reidai}[グラフが曲がっている力の仕事]
位置 $x$（m）で力が $F=3x^2$（N）になる物体を、$x=0$ から2まで動かす仕事を求めましょう。\par
\kangaekata $F$-$x$ グラフは曲がっているので、三角形や台形では求められません。③の $\int F\,\dd x$（「力×小さな距離」の和）を、5-11の方法で計算します。\par
\shiki $3x^2$ を積分すると $3\times\dfrac{x^3}{3}=x^3$（指数を1ふやして、その数でわる）\par
\hspace*{2.6\zw}$\displaystyle\int_0^2 3x^2\,\dd x=\Bigl[\,x^3\,\Bigr]_0^2=2^3-0^3=8-0=8$\par
\kotae \qty{8}{J}
\end{reidai}

\begin{jibun}
力が $F=4x+2$（N）のとき、$x=0$ から3まで動かす仕事を求めましょう。
\bkjans{$\Bigl[\,2x^2+2x\,\Bigr]_0^3=18+6=\qty{24}{J}$}
\end{jibun}

\begin{reidai}[流量から水の量]
ポンプでタンクに水を入れます。流量が $Q=3t^2$（L/min、$t$ は分）でふえていきます。0分から4分までに入った水の量を求めましょう。\par
\kangaekata 流量（L/min）を時間（min）で積分すると、水の量。流量が「1分あたり」で、時間も分なので、単位はそろっています。\par
\shiki $\displaystyle\int_0^4 3t^2\,\dd t=\Bigl[\,3\times\frac{t^3}{3}\,\Bigr]_0^4=\Bigl[\,t^3\,\Bigr]_0^4=4^3-0=64$\par
\hspace*{2.6\zw}単位\quad (L/min)×min＝L\par
\kotae \qty{64}{L}
\end{reidai}

\begin{jibun}
大きなポンプの流量が $Q=10-2t$（L/s、$t$ は秒）でへっていき、5秒で止まりました。出た水の量を求めましょう。
\bkjans{$\Bigl[\,10t-t^2\,\Bigr]_0^5=50-25=\qty{25}{L}$}
\end{jibun}

\noindent\begin{minipage}[c]{0.46\linewidth}
\begin{tikzpicture}
\begin{axis}[bkplot, width=54mm, height=38mm, xmin=0, xmax=4.6, ymin=0, ymax=54, xtick={1,2,3,4}, ytick={12,27,48}, xlabel={$t$（min）}, ylabel={$Q$（L/min）}]
  \foreach \k in {1,...,4} {\edef\tmp{\noexpand\fill[bkBlue!15, draw=bkBlue!60, line width=0.4pt] (axis cs:\k-1,0) rectangle (axis cs:\k,{3*\k*\k});}\tmp}
  \addplot[fill=bkBlue!35, draw=none, domain=0:4] {3*x^2} \closedcycle;
  \addplot[bkcurve, domain=0:4.2] {3*x^2};
\end{axis}
\end{tikzpicture}
\end{minipage}\hfill
\begin{minipage}[c]{0.52\linewidth}\raggedright
\bkreidairef{4}を1分ずつ4つに分け、右はしの流量×1分をたすと $3+12+27+48=90$（大きめ）、左はしでは $0+3+12+27=42$（小さめ）。ほんとうの64 L（こい青）は、そのあいだです。
\end{minipage}

\begin{reidai}[電力が一定のときの電力量]
\qty{1.2}{kW}のドライヤーを15分使いました。電力量は何 kWh で、何 J ですか。\par
\kangaekata 電力が一定なので、電力量＝電力×時間（⑥）。kWh で答えるので、時間を h（時間）に直してからかけます。J には⑥の $\qty{1}{kWh}=\qty{3.6e6}{J}$ で直します。\par
\shiki 時間\quad $\qty{15}{min}\div\qty{60}{min/h}=\qty{0.25}{h}$\par
\hspace*{2.6\zw}電力量\quad $\qty{1.2}{kW}\times\qty{0.25}{h}=\qty{0.3}{kWh}$\par
\hspace*{2.6\zw}J に直す\quad $0.3\times\qty{3.6e6}{J}=\qty{1.08e6}{J}$\par
\kotae \qty{0.3}{kWh}、\qty{1.08e6}{J}（108万 J）
\end{reidai}

\begin{jibun}
\qty{800}{W}の電子レンジを3分使いました。電力量は何 kWh で、何 J ですか。
\bkjans{$\qty{0.8}{kW}\times\qty{0.05}{h}=\qty{0.04}{kWh}$、$800\times180=\qty{144000}{J}$}
\end{jibun}

\begin{reidai}[電力量と平均の電力]
電力が $P=6-t$（kW）の充電器で、0時間から4時間まで充電しました。電力量と、平均の電力を求めましょう。\par
\kangaekata 電力を時間で積分すると電力量（⑥）。平均は、⑧の「積分÷時間」です。kW×h なので、単位は kWh。\par
\shiki $\displaystyle\int_0^4(6-t)\,\dd t=\Bigl[\,6t-\frac{t^2}{2}\,\Bigr]_0^4=\Bigl(24-\frac{16}{2}\Bigr)-0=24-8=16$\par
\hspace*{2.6\zw}平均\quad $16\div4=4$\par
\kotae \qty{16}{kWh}、平均\qty{4}{kW}（はじめの6 kWとおわりの2 kWのまん中）
\end{reidai}

\begin{jibun}
電力が $P=2t$（kW）でふえるヒーターを、0時間から3時間まで使いました。電力量と平均の電力を求めましょう。
\bkjans{$\Bigl[\,t^2\,\Bigr]_0^3=\qty{9}{kWh}$、平均 $9\div3=\qty{3}{kW}$}
\end{jibun}

\begin{reidai}[電流から電気の量]
コンデンサー（電気をためる部品）に流れこむ電流が $I=2-0.5t$（A）でへっていき、4秒で0になりました。たまった電気の量 $q$ を求めましょう。\par
\kangaekata 電流（A＝C/s）は1秒あたりの電気の量（⑦）。時間で積分すると電気の量（C）。\par
\shiki $\displaystyle\int_0^4(2-0.5t)\,\dd t=\Bigl[\,2t-0.25t^2\,\Bigr]_0^4$\par
\hspace*{2.6\zw}$=(2\times4-0.25\times16)-0=8-4=4$\par
\hspace*{2.6\zw}確かめ：三角形の面積 $\dfrac{1}{2}\times4\times2=4$\par
\kotae \qty{4}{C}
\end{reidai}

\begin{jibun}
電流が $I=0.3t$（A）でふえていきます。0秒から10秒までに流れた電気の量を求めましょう。
\bkjans{$\Bigl[\,0.15t^2\,\Bigr]_0^{10}=\qty{15}{C}$}
\end{jibun}

\begin{reidai}[変わる量の平均]
信号で止まっていた自転車が走り出し、次の信号で止まりました。速さは $v=6t-t^2$（m/s、0秒から6秒まで）でした。平均の速さを求めましょう。\par
\kangaekata 平均＝積分÷時間（⑧）。はじめと最後の速さはどちらも0なので、「(はじめ＋おわり)÷2」では求められません。\par
\shiki $\displaystyle\int_0^6(6t-t^2)\,\dd t=\Bigl[\,3t^2-\frac{t^3}{3}\,\Bigr]_0^6=\Bigl(3\times36-\frac{216}{3}\Bigr)-0$\par
\hspace*{2.6\zw}$=108-72=36$\par
\hspace*{2.6\zw}平均\quad $36\div6=6$\par
\kotae \qty{6}{m/s}（6秒で\qty{36}{m}進んだ）
\end{reidai}

\begin{jibun}
速度が $v=3t^2$（m/s）の物体の、0秒から3秒までの平均の速さを求めましょう。
\bkjans{$\Bigl[\,t^3\,\Bigr]_0^3=27$。$27\div3=\qty{9}{m/s}$}
\end{jibun}

\noindent\begin{minipage}[c]{0.48\linewidth}
\begin{tikzpicture}
\begin{axis}[bkplot, width=56mm, height=42mm, xmin=0, xmax=7, ymin=0, ymax=11, xtick={1,2,...,6}, ytick={6,9}, xlabel={$t$（s）}, ylabel={$v$（m/s）}]
  \fill[bkOrange!25] (axis cs:0,0) rectangle (axis cs:6,6);
  \addplot[fill=bkBlue!20, fill opacity=0.7, draw=none, domain=0:6] {6*x-x^2} \closedcycle;
  \addplot[bkcurve, domain=0:6] {6*x-x^2};
  \draw[bkOrange, line width=1.2pt] (axis cs:0,6) -- (axis cs:6,6);
\end{axis}
\end{tikzpicture}
\end{minipage}\hfill
\begin{minipage}[c]{0.5\linewidth}\raggedright
\bkreidairef{8}の自転車の $v$-$t$ グラフです。速さは0から\qty{9}{m/s}まで上がって、また0にもどります。青い部分の面積（\qty{36}{m}）を、はば6秒の長方形にならすと、高さは6です（オレンジの線）。これが平均の速さ\qty{6}{m/s}です。
\end{minipage}

\begin{reidai}[電池の容量 mAh]
電流が $I=100+50t$（mA、$t$ は時間）で流れたとき、2時間で使った電気の量は、何 mAh（ミリアンペア時）でしょう。\par
\kangaekata 電流（mA）を時間（h）で積分すると、電気の量（mAh）。\qty{1}{mA}が1時間流れた量が\qty{1}{mAh}です。\par
\shiki $\displaystyle\int_0^2(100+50t)\,\dd t=\Bigl[\,100t+50\times\frac{t^2}{2}\,\Bigr]_0^2=\Bigl[\,100t+25t^2\,\Bigr]_0^2$\par
\hspace*{2.6\zw}$=(100\times2+25\times4)-0=200+100=300$\par
\hspace*{2.6\zw}確かめ：台形の面積（電流は100 mAから200 mA）\par
\hspace*{2.6\zw}$(100+200)\times2\div2=300$\par
\kotae \qty{300}{mAh}
\end{reidai}

\begin{jibun}
電流が $I=400-100t$（mA）でへっていき、4時間で0になりました。使った電気の量を求めましょう。
\bkjans{$\Bigl[\,400t-50t^2\,\Bigr]_0^4=1600-800=\qty{800}{mAh}$}
\end{jibun}

\noindent\begin{minipage}[c]{0.46\linewidth}
\begin{tikzpicture}
\begin{axis}[bkplot, width=54mm, height=36mm, xmin=0, xmax=2.7, ymin=0, ymax=260, xtick={1,2}, ytick={100,200}, xlabel={$t$（h）}, ylabel={$I$（mA）}]
  \addplot[fill=bkGreen!25, draw=none, domain=0:2] {100+50*x} \closedcycle;
  \addplot[bkcurve, domain=0:2.5] {100+50*x};
  \draw[bkGray, dashed, line width=0.5pt] (axis cs:2,0) -- (axis cs:2,200);
  \node[font=\sffamily\small, fill=bkGreen!25, inner sep=1pt] at (axis cs:1,70) {\qty{300}{mAh}};
\end{axis}
\end{tikzpicture}
\end{minipage}\hfill
\begin{minipage}[c]{0.5\linewidth}\raggedright
\bkreidairef{9}の $I$-$t$ グラフです。緑の台形の面積が、使った電気の量\qty{300}{mAh}です。たての単位 mA と横の単位 h をかけるので、面積の単位が mAh になります。
\end{minipage}

\begin{reidai}[はかった値から、たまった量を求める]
ポンプの流量を、1分ごとにはかると、下の表のようでした。0分から4分までに出た水の量を、およそ求めましょう。
\begin{bkcenter}
\begin{bktable}{colspec={Q[l,m,wd=30mm] *{5}{Q[c,m,wd=10mm]}}, rowsep=1pt}
$t$（min） & 0 & 1 & 2 & 3 & 4 \\
$Q$（L/min） & 2 & 6 & 8 & 9 & 9 \\
\end{bktable}
\end{bkcenter}
\kangaekata 式がないので、積分の公式は使えません。そこで、5-10のように、細い図形の面積をたします。1分ごとの区間の中では、流量はほぼまっすぐ変わると考え、はじめとおわりの流量を上底・下底とする\textbf{台形}の面積を、その1分の水の量とします。台形の面積は、左はしの長方形（小さめ）と右はしの長方形（大きめ）の、ちょうどまん中です。\par
\shiki 0〜1分\quad $(2+6)\times1\div2=4$\par
\hspace*{2.6\zw}1〜2分\quad $(6+8)\times1\div2=7$\par
\hspace*{2.6\zw}2〜3分\quad $(8+9)\times1\div2=8.5$\par
\hspace*{2.6\zw}3〜4分\quad $(9+9)\times1\div2=9$\par
\hspace*{2.6\zw}合計\quad $4+7+8.5+9=28.5$\par
\hspace*{2.6\zw}確かめ：右はしの和 $6+8+9+9=32$\par
\hspace*{2.6\zw}\phantom{確かめ：}左はしの和 $2+6+8+9=25$。28.5はそのあいだ\par
\kotae 約\qty{28.5}{L}
\end{reidai}

\begin{jibun}
太陽光パネルの電力を1時間ごとにはかると、0、2、3、3、\qty{2}{kW}（0時間から4時間まで）でした。4時間の電力量を、台形の面積の和でおよそ求めましょう。
\bkjans{$1+2.5+3+2.5=\qty{9}{kWh}$}
\end{jibun}

\noindent\begin{minipage}[c]{0.46\linewidth}
\begin{tikzpicture}
\begin{axis}[bkplot, width=54mm, height=38mm, xmin=0, xmax=4.6, ymin=0, ymax=11, xtick={1,2,3,4}, ytick={2,6,9}, xlabel={$t$（min）}, ylabel={$Q$（L/min）}]
  \addplot[fill=bkGreen!25, draw=bkGreen!60!black, line width=0.4pt] coordinates {(0,0) (0,2) (1,6) (2,8) (3,9) (4,9) (4,0)} -- cycle;
  \draw[bkGreen!60!black, line width=0.4pt] (axis cs:1,0) -- (axis cs:1,6) (axis cs:2,0) -- (axis cs:2,8) (axis cs:3,0) -- (axis cs:3,9);
  \addplot[only marks, mark=*, mark size=1.6pt, bkInk] coordinates {(0,2) (1,6) (2,8) (3,9) (4,9)};
\end{axis}
\end{tikzpicture}
\end{minipage}\hfill
\begin{minipage}[c]{0.5\linewidth}\raggedright
\bkreidairef{10}の表の点を直線でつなぐと、4つの台形（緑）ができます。その面積の和が28.5 L です。工場の積算流量計や家の電力量計（スマートメーター）も、はかった値を、ごく短い時間ごとにこのようにたし合わせています。
\end{minipage}

\begin{reidai}[階段をのぼる仕事]
質量\qty{40}{kg}の人が、高さ\qty{3}{m}（建物の1階分）の階段を、ゆっくり一定の速さでのぼります。重力にさからってする仕事は何 J ですか。また、それは何 kWh ですか（重力は1 kg あたり\qty{9.8}{N}）。\par
\kangaekata 階段はななめに進みますが、体を持ち上げる力は上向きです。力の向き（上向き）に動いた距離は、のぼった高さの\qty{3}{m}です。横に進んだ分は、上向きの力の向きに動いていないので、この仕事に入りません。kWh には、⑥の $\qty{1}{kWh}=\qty{3.6e6}{J}$ でわって直します。\par
\shiki 力\quad $F=\qty{40}{kg}\times\qty{9.8}{N/kg}=\qty{392}{N}$\par
\hspace*{2.6\zw}仕事\quad $W=392\times3=\qty{1176}{J}$\par
\hspace*{2.6\zw}kWh に直す\quad $\qty{1176}{J}\div(\qty{3.6e6}{J/kWh})\approx0.00033$\,kWh\par
\kotae \qty{1176}{J}（約\qty{1.2}{kJ}）、約0.00033\,kWh
\end{reidai}

\noindent\begin{minipage}[c]{0.44\linewidth}\centering
\begin{tikzpicture}[x=1mm, y=1mm]
  \fill[bkGray!25] (0,0) -- (0,6) -- (6,6) -- (6,12) -- (12,12) -- (12,18) -- (18,18) -- (18,24) -- (24,24) -- (24,30) -- (32,30) -- (32,0) -- cycle;
  \draw[bkInk, line width=0.6pt] (-2,0) -- (0,0) -- (0,6) -- (6,6) -- (6,12) -- (12,12) -- (12,18) -- (18,18) -- (18,24) -- (24,24) -- (24,30) -- (32,30);
  \bkperson{(2,6.5)}{13}{bkBlue}
  \draw[bk arrow, bkOrange, line width=0.9pt] (-5,4) -- (-5,30);
  \node[font=\sffamily\scriptsize, anchor=south, text=bkOrange!80!black] at (-5,31) {力（上向き）};
  \draw[bkGray, <->, line width=0.5pt] (35,0) -- (35,30);
  \node[font=\sffamily\scriptsize, anchor=west] at (36.5,15) {\qty{3}{m}};
\end{tikzpicture}
\end{minipage}\hfill
\begin{minipage}[c]{0.53\linewidth}\raggedright
\bkreidairef{11}の仕事1176 J は、電気代にすると約0.01円（1 kWh を31円として、$0.00033\times31\approx0.01$）。1 kWh は、この階段をおよそ3000回のぼる仕事です（$\qty{3.6e6}{J}\div\qty{1176}{J}\approx3061$）。kWh がとても大きな単位だとわかります。
\end{minipage}

\begin{jibun}
質量\qty{50}{kg}の人が、高さ\qty{12}{m}（4階分）の階段をのぼります。重力にさからってする仕事は何 J で、何 kWh ですか（重力は1 kg あたり\qty{9.8}{N}）。
\bkjans{$50\times9.8\times12=\qty{5880}{J}$、$5880\div(3.6\times10^6)\approx\qty{0.0016}{kWh}$}
\end{jibun}

%% ---------------- まちがえやすいところ・物理工学 ----------------
\bksection{tsumazuki}{まちがえやすいところをチェック}

\begin{tsumazuki}{ばねの仕事を $kx\times x$ としてしまう}
\bkng{$W=kx\times x=kx^2$}\bkok{$W=\dfrac{1}{2}kx^2$}
$kx$ は、のびが $x$ になったときの最後の力で、はじめの力は0です。$F$-$x$ グラフは三角形なので、$\dfrac{1}{2}$ がつきます。
\end{tsumazuki}

\begin{tsumazuki}{時間の単位と流量の単位がそろっていない}
\bkng{$Q=\qty{12}{L/min}$ で10秒なので $12\times10=\qty{120}{L}$}\bkok{$\qty{12}{L/min}\times\qty{10}{s}\div\qty{60}{s/min}=\qty{2}{L}$}
流量が「1分あたり」なのに秒をかけると、60倍まちがえます。単位を書いてかければ、min と s が約分できないことに気づきます（2-10）。
\end{tsumazuki}

\begin{tsumazuki}{kW と kWh を同じものと考える}
\bkng{\qty{1.2}{kW}のドライヤーを15分使ったので、電力量は\qty{1.2}{kWh}}\bkok{$\qty{1.2}{kW}\times\qty{0.25}{h}=\qty{0.3}{kWh}$}
kW は電力（1秒あたりに使うエネルギー）で、速さのような「あたり」の量です。kWh は電力×時間で、進んだ距離のような「たまった量」です。時間をかけないと kWh にはなりません。
\end{tsumazuki}

\begin{tsumazuki}{平均を「(はじめ＋おわり)÷2」としてしまう}
\bkng{$v=3t^2$ の0〜2秒の平均は $(0+12)\div2=\qty{6}{m/s}$}\bkok{$8\div2=\qty{4}{m/s}$（積分÷時間）}
「(はじめ＋おわり)÷2」が使えるのは、グラフがまっすぐなとき（台形）だけです。曲がっているときは、積分÷時間で求めます。
\end{tsumazuki}

\begin{column}{電気代を計算しよう}
電力会社は、使った電力量（kWh）に応じて電気代を決めます。1 kWh あたりの料金を約31円とします。\qty{1.2}{kW}のドライヤーを毎日15分（\qty{0.25}{h}）、30日使うと、
\[ \qty{1.2}{kW}\times\qty{0.25}{h}\times30=\qty{9}{kWh},\qquad 9\times31=279 \]
で、1か月の電気代は約280円です。電力（kW）が大きい器具ほど、また長い時間使うほど、電力量がたまって電気代が高くなります。
\end{column}

\begin{butsuri}
\begin{itemize}[leftmargin=1.2\zw, itemsep=1pt, topsep=0pt, label={\color{bkPurple}●}]
\item \textbf{電気のメーター}\\家の電力量計（スマートメーター）は、電力を細かい時間ごとにはかってたし合わせ（積分し）、kWh で表示しています。
\item \textbf{ばねのエネルギー}\\おもちゃの車のぜんまいや、時計のばね、車のサスペンションは、$\frac{1}{2}kx^2$ のエネルギーをたくわえたり、とり出したりしています。
\item \textbf{工場の積算流量計}\\化学工場では、パイプを流れる液体やガスの流量を時間で積分して、原料や製品の量を管理しています。
\end{itemize}
\end{butsuri}

%% ---------------- 章末問題 ----------------
\bkendproblems
\bkproblemgroup{基本}{3}

\bkproblem \qty{20}{N}の力で、物を力の向きに\qty{3}{m}動かしました。仕事を求めましょう。
\bkpans{\qty{60}{J}}

\bkproblem ばね定数 $k=\qty{400}{N/m}$ のばねについて、（1）0から\qty{0.05}{m}、（2）\qty{0.05}{m}から\qty{0.1}{m} のばす仕事を求めましょう。
\bkpans{（1）\nobreak \qty{0.5}{J}\quad（2）\nobreak $2-0.5=\qty{1.5}{J}$}

\bkproblem ポンプの流量が $Q=0.5t$（L/s）でふえていきます。0秒から20秒までに出た水の量を求めましょう。
\bkpans{$\Bigl[\,0.25t^2\,\Bigr]_0^{20}=\qty{100}{L}$}

\bkproblem \qty{1.2}{kW}のドライヤーを10分使った電力量は、何 kWh、何 J ですか。
\bkpans{$1.2\times\dfrac{10}{60}=\qty{0.2}{kWh}$、$0.2\times3.6\times10^6=\qty{7.2e5}{J}$}

\bkproblem 位置 $x$（m）で力が $F=6x$（N）になる物体を、$x=0$ から2まで動かす仕事を求めましょう。
\bkpans{$\Bigl[\,3x^2\,\Bigr]_0^2=\qty{12}{J}$}

\bkproblem 電流が $I=0.2t$（A）でふえていきます。0秒から5秒までに流れた電気の量を求めましょう。
\bkpans{$\Bigl[\,0.1t^2\,\Bigr]_0^5=\qty{2.5}{C}$}

\bkproblemgroup{チャレンジ}{3}

\bkproblem 速度が $v=t^2$（m/s）の物体の、0秒から3秒までの平均の速さを求めましょう。
\bkpans{距離 $\Bigl[\,\dfrac{t^3}{3}\,\Bigr]_0^3=9$。平均 $9\div3=\qty{3}{m/s}$}

\bkproblem 小さな太陽光パネルの電力が、朝6時から夕方6時まで $P=12t-t^2$（W、$t$ は6時からの時間）で変わりました。1日の電力量（Wh）と、平均の電力を求めましょう。
\bkpans{$\Bigl[\,6t^2-\dfrac{t^3}{3}\,\Bigr]_0^{12}=864-576=\qty{288}{Wh}$、平均 $288\div12=\qty{24}{W}$}

\bkproblem スマートフォンから、電流 $I=300+100t$（mA、$t$ は時間）が流れ続けたとします。4時間で使った電気の量は何 mAh ですか。\qty{3000}{mAh}の電池は、4時間もちますか。
\bkpans{$\Bigl[\,300t+50t^2\,\Bigr]_0^4=1200+800=\qty{2000}{mAh}$。3000より少ないので、\mbox{4時間}もつ}

\bkproblem 暖房の強さを変えた部屋の温度が $T=15+6t-t^2$（℃、$t$ は時間）で変わりました。$t=0$ から6までの平均の温度を求めましょう。
\bkpans{$\Bigl[\,15t+3t^2-\dfrac{t^3}{3}\,\Bigr]_0^6=90+108-72=126$。$126\div6=\mbox{$21$\,℃}$}

\begin{furikaeri}
\begin{checklist}
\item 変わる力の仕事を、積分で求められる。\hfill{\footnotesize\sffamily\color{bkTheme}→\pageref{bk:5-12:sekibun}ページ}
\item 流量や電力から、たまった量を求められる。\hfill{\footnotesize\sffamily\color{bkTheme}→\pageref{bk:5-12:ryuuryou}ページ}
\item 積分を使って、平均の値を求められる。\hfill{\footnotesize\sffamily\color{bkTheme}→\pageref{bk:5-12:heikin}ページ}
\end{checklist}
\end{furikaeri}

\bkfukidashi{mina}{たまった量は積分でわかるけど、「はじめにどれだけあったか」がわからないと、いまの量は決まらないよね。}

%% 5-13 積分定数は「はじめの状態」（第11回：ページ制限なし。C がなぜ「はじめの位置」なのかを、表・グラフ・式の3つで見せる）
\bkchapter{5}{13}{積分定数は「はじめの状態」}{\bklevel{3}}{積分定数は「はじめの状態」}
\bkmichishirube{5-5,5-9,5-11}{6-1,6-3}

\begin{goal}
\begin{itemize}[label={\color{bkTheme}●}, leftmargin=1.6\zw, itemsep=0pt, topsep=1pt]
\item \bklevel{3}\enspace 積分定数 $C$ を、はじめの状態（初期条件）から決められる。
\item \bklevel{3}\enspace 加速度から速度、速度から位置を、積分で求められる。
\item \bklevel{3}\enspace $x=x_0+v_0t+\dfrac{1}{2}at^2$ を、積分で導ける。
\end{itemize}
\end{goal}

\begin{hatena}
\noindent\begin{minipage}[c]{0.58\linewidth}\raggedright
AさんとBさんは、同じ時刻に、同じ速さで、同じ向きに歩き出しました。1時間のあいだ、2人の速さはいつも同じでした。

では、1時間後、2人は同じ場所にいるでしょうか。
\end{minipage}\hfill
\begin{minipage}[c]{0.38\linewidth}
\begin{tikzpicture}[x=1mm, y=1mm]
  \draw[bkGray, line width=1pt] (0,0) -- (40,0);
  \draw[bkGray, line width=1pt] (0,12) -- (40,12);
  \bkperson{(2,0.5)}{10}{bkRoute2}
  \bkperson{(14,12.5)}{10}{bkRoute4}
  \draw[bk arrow, bkGray, line width=0.9pt] (8,6) -- (18,6);
  \draw[bk arrow, bkGray, line width=0.9pt] (20,18) -- (30,18);
  \node[font=\sffamily\scriptsize, anchor=east] at (-1.5,5) {A};
  \node[font=\sffamily\scriptsize, anchor=east] at (10.5,17) {B};
\end{tikzpicture}
\end{minipage}
\end{hatena}

速さがいつも同じなら、1時間に進む距離は同じです。でも、Bさんがはじめから\qty{1}{km}先にいたなら、1時間後も\qty{1}{km}先にいます。速さだけでは、位置は決まりません。「はじめにどこにいたか」が必要です。この章では、5-11で出てきた積分の $+C$ が、この「はじめの状態」を表していることを確かめます。

\clearpage
%% ---------------- やってみよう・思い出そう ----------------
\begin{yattemiyou}[同じ速度、ちがうスタート]
速度が $v=2t$（m/s）の物体が3つあります。スタートの位置は、Aが\qty{0}{m}、Bが\qty{3}{m}、Cが\qty{-2}{m}です（Cはスタートの線より\qty{2}{m}後ろ）。
\begin{enumerate}[label={\bkcircled[bkGreen]{\arabic*}}, leftmargin=2.4\zw, itemsep=1pt, topsep=3pt]
\item 0秒から $t$ 秒までに進む距離を、$v$-$t$ グラフの面積で求めよう。
\item それぞれの位置（スタートの位置＋進んだ距離）を、$t=0$〜3 で表にしよう。
\item 横軸を $t$、たて軸を位置にして、3つのグラフをかこう（$x$-$t$ グラフ）。
\end{enumerate}
\begin{bkcenter}
\begin{bktable}{colspec={Q[l,m,wd=34mm] *{4}{Q[c,m,wd=12mm]}}, rowsep=1.5pt}
$t$（s） & 0 & 1 & 2 & 3 \\
進んだ距離 $t^2$（m） & 0 & 1 & 4 & 9 \\
Aの位置（m） & 0 & 1 & 4 & 9 \\
Bの位置（m） & 3 & 4 & 7 & 12 \\
Cの位置（m） & $-2$ & $-1$ & 2 & 7 \\
\end{bktable}
\end{bkcenter}
\end{yattemiyou}

\noindent ①の面積は、底辺 $t$、高さ $2t$（$t$ 秒の速度）の三角形なので、$\dfrac{1}{2}\times t\times2t=t^2$ です（5-9）。5-11の積分で書くと、
\[ \int_0^t 2s\,\dd s=\Bigl[\,s^2\,\Bigr]_0^t=t^2-0=t^2 \]
です（$s$ は0から $t$ まで動く時刻。上端の $t$ と区別するため、別の文字にしています）。どちらの方法でも、進んだ距離は $t^2$ です。

3つのグラフは、同じ形の放物線を、上下にずらしたものになります。どの時刻でも、3つのグラフの傾き（速度）は同じです。それでも位置がちがうのは、スタートの位置がちがうからです。

\begin{omoidasou}
\bkrecall{5-11}{微分すると $f(x)$ になる関数が原始関数。$\displaystyle\int f(x)\,\dd x=F(x)+C$。定数は微分すると0なので、原始関数には $+C$ の自由がある。}
\bkrecall{5-5}{速度は位置を時間で微分したもの、加速度は速度を時間で微分したもの。落ちる物の加速度は $g=\qty{9.8}{m/s^2}$。}
\bkrecall{5-9}{$v$-$t$ グラフの面積は、位置の変化（変位）。}
\bkrecall{2-2}{マイナスは「反対の向き」を表す。$+$ の向きは、はじめに決めて変えない。}
\bkrecall{3-10}{$t^2+2t-15=0$ は、かけて $-15$、たして2になる2つの数（5と $-3$）を見つけて、$(t+5)(t-3)=0$ と因数分解する。かけて0になるのは、どちらかが0のときなので、$t=-5$ か $t=3$。}
\end{omoidasou}

%% ---------------- 数学の道具 ----------------
\bksection{doguu}{数学の道具}

\bksubsection{1}{速度を積分すると、位置が「$+C$」つきでわかる}

\label{bk:5-13:C}位置 $x(t)$ を微分すると速度 $v(t)$ です（5-5）。逆に言うと、位置 $x(t)$ は「微分すると $v(t)$ になる関数」、つまり $v(t)$ の原始関数です（5-11）。だから
\[ x(t)=\int v(t)\,\dd t=F(t)+C \]
と書けます。$v=2t$ なら、$t^2$ を微分すると $2t$ なので、$x=t^2+C$ です。

\noindent\textbf{なぜ $+C$ がつくの？}\quad $t^2+3$ も $t^2-2$ も、微分すると $2t$ です。定数（3や $-2$）は時間がたっても変わらないので、変わる速さは0だからです。表で確かめます。

\begin{bkcenter}
\begin{bktable}{colspec={Q[l,m,wd=40mm] *{4}{Q[c,m,wd=9mm]} Q[c,m,wd=26mm]}, rowsep=1.5pt, row{1}={bg=bkTheme!10}}
$t$（s） & 0 & 1 & 2 & 3 & 1秒ごとのふえ方 \\
$x=t^2$（$C=0$） & 0 & 1 & 4 & 9 & 1，3，5 \\
$x=t^2+3$（$C=3$） & 3 & 4 & 7 & 12 & 1，3，5 \\
$x=t^2-2$（$C=-2$） & $-2$ & $-1$ & 2 & 7 & 1，3，5 \\
\end{bktable}
\bktabcaption{$C$ がちがっても、1秒ごとのふえ方（動き方）はまったく同じ。}
\end{bkcenter}

$C$ の値によって、グラフは上下にずれるだけです。どの $C$ でも、速度（傾き）は同じ $2t$ です。だから、速度だけでは $C$ は決まりません。

\noindent\begin{minipage}[c]{0.5\linewidth}
\begin{tikzpicture}
\begin{axis}[bkplot, width=62mm, height=48mm, xmin=0, xmax=3.4, ymin=-2.6, ymax=13, xtick={1,2,3}, xticklabels={}, ytick={-2,3,9}, xlabel={$t$（s）}, ylabel={$x$（m）}]
  \addplot[bkcurve, domain=0:2.75] {x^2};
  \addplot[bkcurve2, domain=0:2.75] {x^2+3};
  \addplot[bkcurve3, domain=0:2.75] {x^2-2};
  % 目もりの数字は、緑の線（t=1 で −1）と重ならないよう、−2 の線より下に置く
  \node[font=\footnotesize, text=bkInk, anchor=north] at (axis cs:1,-3.2) {$1$};
  \node[font=\footnotesize, text=bkInk, anchor=north] at (axis cs:2,-3.2) {$2$};
  \node[font=\footnotesize, text=bkInk, anchor=north] at (axis cs:3,-3.2) {$3$};
  \node[font=\sffamily\scriptsize, anchor=west, text=bkOrange!80!black, fill=white, inner sep=1pt] at (axis cs:2.85,10.6) {$C=3$（B）};
  \node[font=\sffamily\scriptsize, anchor=west, text=bkBlue, fill=white, inner sep=1pt] at (axis cs:2.85,7.6) {$C=0$（A）};
  \node[font=\sffamily\scriptsize, anchor=west, text=bkGreen!70!black, fill=white, inner sep=1pt] at (axis cs:2.85,5.6) {$C=-2$（C）};
\end{axis}
\end{tikzpicture}
\end{minipage}\hfill
\begin{minipage}[c]{0.46\linewidth}
$x=t^2+C$ に $t=0$ を入れると、
\[ x(0)=0^2+C=C \]
です。つまり、\textbf{積分定数 $C$ は、はじめ（$t=0$）の位置}を表しています。

はじめの位置 $x_0$ がわかれば、$C=x_0$ と決まり、どの時刻の位置も1つに決まります。
\end{minipage}

\noindent\textbf{\bkhatenaref の答え}\quad AさんとBさんは、速さがいつも同じなので、同じ時間に同じ距離だけ進みます。だから、1時間後の2人の差は、はじめの差のままです。Bさんがはじめ\qty{1}{km}先にいたなら1時間後も\qty{1}{km}先、同じ場所から出発したなら同じ場所にいます。位置の式で言えば、2人の式は $C$（はじめの位置）だけがちがうのです。

\bksubsection{2}{はじめの状態：初期条件}

\label{bk:5-13:shoki}$C$ を決めるのに使う「ある時刻の値」を、\textbf{初期条件}（しょきじょうけん）といいます。たいていは $t=0$ の値ですが、ほかの時刻の値でも $C$ は決められます。\bkidx{しょきじょうけん}{初期条件}

\noindent\begin{minipage}[c]{0.5\linewidth}
\begin{tikzpicture}
\begin{axis}[bkplot, width=62mm, height=48mm, xmin=0, xmax=3.4, ymin=-1, ymax=16, xtick={1,2}, ytick={4,8,12}, xlabel={$t$（s）}, ylabel={$x$（m）}]
  \foreach \c in {0,6} {\edef\tmp{\noexpand\addplot[bkthin, domain=0:2.75] {x^2+\c};}\tmp}
  \addplot[bkcurve2, domain=0:2.75] {x^2+3};
  \addplot[only marks, mark=*, mark size=2pt, bkInk] coordinates {(1,4)};
  \node[font=\sffamily\scriptsize, anchor=west, text=bkOrange!80!black, fill=white, inner sep=1pt] at (axis cs:2.82,10.6) {$C=3$};
\end{axis}
\end{tikzpicture}
\end{minipage}\hfill
\begin{minipage}[c]{0.46\linewidth}\raggedright
たとえば、速度 $v=2t$ の物体が、「1秒のとき\qty{4}{m}の所にいた」とわかったとします。$x=t^2+C$ に $t=1$、$x=4$ を入れると
\[ 4=1^2+C,\qquad C=4-1=3 \]
グラフでは、上下にずれた放物線のなかから、点 $(1,\,4)$（黒い点）を通る1本（オレンジ）を選んだことになります。
\end{minipage}

\noindent\textbf{$C$ を決める4つの手順}\quad どの問題でも、次の順に進めば $C$ が決まります。右の列は、いまの例（$v=2t$、1秒のとき\qty{4}{m}）です。
\begin{bkcenter}
\begin{bktable}{colspec={Q[c,m,wd=8mm] Q[l,m,wd=60mm] Q[l,m,wd=50mm]}, rowsep=2pt, row{1}={bg=bkTheme!10}}
 & すること & いまの例 \\
① & 積分して、$+C$ をつける & $x=t^2+C$ \\
② & 初期条件の時刻と値を入れる & $t=1$、$x=4$ を入れて $4=1^2+C$ \\
③ & $C$ について解く & 両辺から1をひいて $C=3$ \\
④ & 確かめる（微分してもとの式、初期条件を入れて合うか） & $(t^2+3)'=2t$、$1^2+3=4$ \\
\end{bktable}
\bktabcaption{①で $+C$ を忘れると、②③ができなくなる。}
\end{bkcenter}

このあと、例題1は $t=0$ の値で、例題2は $t=1$ の値で $C$ を決めます。

\bksubsection{3}{加速度から位置まで：2回積分する}

\label{bk:5-13:toukasoku}加速度 $a$ が一定のとき、2回積分して位置を求めます。5-5で予告した式 $x=x_0+v_0t+\dfrac{1}{2}at^2$ を、ここで導きましょう。

\noindent\textbf{1回め：加速度から速度}\quad 速度を微分すると加速度なので、速度は加速度の原始関数です。$a$ は定数で、$at$ を微分すると $a$ になる（5-4）ので、
\[ v=\int a\,\dd t=at+C_1 \]
$t=0$ を入れると $v(0)=a\times0+C_1=C_1$ なので、$C_1$ は\textbf{はじめの速度}です。これを $v_0$（ブイ・ゼロ）と書くと、$C_1=v_0$ で
\[ v=v_0+at \]
「はじめの速度に、1秒ごとに $a$ ずつたしていく」という意味です。

\noindent\textbf{2回め：速度から位置}\quad 位置は、速度の原始関数です。$v_0$ も $a$ も定数なので、項ごとに積分します（5-11）。$v_0$ を積分すると $v_0t$、$at$ を積分すると $a\times\dfrac{t^2}{2}$（指数を1ふやして、その数でわる）です。
\[ x=\int(v_0+at)\,\dd t=v_0t+\frac{1}{2}at^2+C_2 \]
$t=0$ を入れると $x(0)=C_2$ なので、$C_2$ は\textbf{はじめの位置} $x_0$ です。だから
\[ x=x_0+v_0t+\frac{1}{2}at^2 \]

\noindent\textbf{数と図で確かめる}\quad $a=\qty{2}{m/s^2}$、$v_0=\qty{3}{m/s}$、$x_0=\qty{1}{m}$ の場合を表にします。式は $v=3+2t$、$x=1+3t+t^2$ です。

\begin{bkcenter}
\begin{bktable}{colspec={Q[l,m,wd=50mm] *{4}{Q[c,m,wd=11mm]}}, rowsep=1.5pt}
$t$（s） & 0 & 1 & 2 & 3 \\
速度 $v=3+2t$（m/s） & 3 & 5 & 7 & 9 \\
1秒ごとの台形の面積（m） & & 4 & 6 & 8 \\
位置 $x=1+3t+t^2$（m） & 1 & 5 & 11 & 19 \\
\end{bktable}
\bktabcaption{位置は、はじめの1 m に、$v$-$t$ グラフの台形の面積をつみ上げたもの。}
\end{bkcenter}

\noindent\begin{minipage}[c]{0.46\linewidth}
\begin{tikzpicture}
\begin{axis}[bkplot, width=56mm, height=44mm, xmin=0, xmax=3.6, ymin=0, ymax=11, xtick={1,2,3}, ytick={3,5,7,9}, xlabel={$t$（s）}, ylabel={$v$（m/s）}]
  \addplot[fill=bkBlue!18, draw=bkBlue!60, line width=0.4pt] coordinates {(0,0) (0,3) (1,5) (1,0)} -- cycle;
  \addplot[fill=bkGreen!22, draw=bkGreen!60!black, line width=0.4pt] coordinates {(1,0) (1,5) (2,7) (2,0)} -- cycle;
  \addplot[fill=bkOrange!25, draw=bkOrange!70, line width=0.4pt] coordinates {(2,0) (2,7) (3,9) (3,0)} -- cycle;
  \addplot[bkcurve, domain=0:3.4] {3+2*x};
  \node[font=\sffamily\small] at (axis cs:0.5,1.5) {4};
  \node[font=\sffamily\small] at (axis cs:1.5,2.2) {6};
  \node[font=\sffamily\small] at (axis cs:2.5,3) {8};
\end{axis}
\end{tikzpicture}
\end{minipage}\hfill
\begin{minipage}[c]{0.52\linewidth}\raggedright
速度は1秒ごとに2ずつふえます（加速度2）。0〜1秒の $v$-$t$ グラフは上底3、下底5の台形（青）で、面積は $(3+5)\times1\div2=4$。はじめの位置1に4をたすと5で、式の $x(1)=1+3+1=5$ と合います。次の1秒は緑の台形6をたして11、その次はオレンジの8をたして19です。
\end{minipage}

\noindent $\dfrac{1}{2}at^2$ の $\dfrac{1}{2}$ は、5-9の三角形の面積の「÷2」と同じものです。台形を「長方形 $v_0\times t$」と「三角形 $\dfrac{1}{2}\times t\times at$」に分けると、$v_0t+\dfrac{1}{2}at^2$ になるからです。

2回積分したので、積分定数が2つ出てきました。それを決めるのに、はじめの位置 $x_0$ とはじめの速度 $v_0$ の、2つの初期条件が必要です。たとえば同じ $a=2$ でも、$v_0$ が3と0では、1秒後の速度が5と2でちがいます。$v_0$ だけわかっても、$x_0$ がわからなければ、どこにいるかは決まりません。

\bksubsection{4}{落ちる物・投げ上げた物}

\label{bk:5-13:rakka}\noindent\begin{minipage}[c]{0.6\linewidth}\raggedright
上向きを＋と決めます（2-2）。位置 $x$ は、地面からの高さです。重力は、物を1秒ごとに下向きへ\qty{9.8}{m/s}ずつ速くします。下向きは「＋の反対の向き」なので、加速度は $a=-9.8$ です。

高さ $x_0$ から、上向きの速さ $v_0$ で投げると、③の式に $a=-9.8$ を入れて
\[ v=v_0-9.8t,\qquad x=x_0+v_0t-4.9t^2 \]
4.9は、$\dfrac{1}{2}at^2$ の $\dfrac{1}{2}\times9.8=4.9$ から来ています。
\end{minipage}\hfill
\begin{minipage}[c]{0.36\linewidth}\centering
\begin{tikzpicture}[x=1mm, y=1mm]
  \fill[bkGray!35] (-2,-3) rectangle (40,0);
  \draw[bkInk, line width=0.7pt] (-2,0) -- (40,0);
  \node[font=\sffamily\scriptsize, anchor=north] at (24,-4) {地面（$x=0$）};
  \draw[bk arrow, bkInk, line width=0.8pt] (4,0) -- (4,40);
  \node[font=\sffamily\scriptsize, anchor=south] at (4,40.5) {＋（上）};
  \draw[bkInk, line width=0.5pt] (2.5,28) -- (5.5,28);
  \node[font=\sffamily\scriptsize, anchor=east] at (2,28) {$x_0$};
  \fill[bkOrange!70] (14,28) circle (2.6);
  \draw[bkOrange!80!black, line width=0.5pt] (14,28) circle (2.6);
  \draw[bk arrow, bkBlue, line width=0.9pt] (14,31.5) -- (14,39);
  \node[font=\sffamily\scriptsize, text=bkBlue, anchor=west] at (15.5,36) {$v_0$（上向き）};
  \draw[bk arrow, bkRed, line width=0.9pt] (14,24.5) -- (14,13);
  \node[font=\sffamily\scriptsize, text=bkRed, anchor=west, align=left] at (15.5,18) {加速度\\$a=-9.8$};
\end{tikzpicture}
\end{minipage}

\par\medskip
\noindent\textbf{例}\quad 高さ\qty{20}{m}から、ボールを静かに落とします。「静かに落とす」は、はじめの速度が0ということです。$x_0=20$、$v_0=0$ なので、位置は $x=20-4.9t^2$ です。地面は高さ0なので、$x=0$ になる時刻を求めます。
\begin{align*}
20-4.9t^2&=0\\
4.9t^2&=20\qquad\text{（両辺に $4.9t^2$ をたす）}\\
t^2&=20\div4.9\approx4.08\qquad\text{（両辺を4.9でわる）}
\end{align*}
2回かけて4.08になる正の数は、電卓の√キーで $\sqrt{4.08}\approx2.02$ です（2-8）。$t$ は時間なので、正のほうだけを使います。地面に着くのは約\qty{2.0}{s}後です。

\bksubsection{5}{定積分で書くと「はじめ＋たまった変化」}

\label{bk:5-13:teisekibun}位置の変化は、$v$-$t$ グラフの面積でした（5-9）。はじめの位置にそれをたせば、いまの位置です。
\[ x(t)=x_0+\int_0^t v(s)\,\dd s \]
（$s$ は0から $t$ まで動く時刻です。）おこづかいでたとえると、「いまの貯金＝はじめの貯金＋その後にたまった分」です。たまった分がわかっても、はじめの貯金がわからなければ、いまの貯金はわかりません。$C$ や $x_0$ が「はじめの貯金」にあたります。

「いまの量＝はじめの量＋たまった変化」は、水そうの水の量、温度、電池の電気の量など、変化の速さがわかる量すべてに使えます（5-12で求めた「たまった量」に、はじめの量をたせばよい）。

\begin{bkfigure}
\begin{tikzpicture}[x=1mm, y=1mm, m/.style={draw=bkTheme, line width=1pt, fill=white, rounded corners=2mm, minimum width=28mm, minimum height=11mm, align=center, font=\sffamily\small}]
  \node[m] (a) at (0,0) {加速度 $a$};
  \node[m] (v) at (48,0) {速度 $v$};
  \node[m] (x) at (96,0) {位置 $x$};
  \draw[bk arrow, bkBlue] (a) -- node[above, font=\sffamily\scriptsize] {積分} node[below, font=\sffamily\scriptsize] {$+v_0$} (v);
  \draw[bk arrow, bkBlue] (v) -- node[above, font=\sffamily\scriptsize] {積分} node[below, font=\sffamily\scriptsize] {$+x_0$} (x);
\end{tikzpicture}
\bkfigcaption{積分するたびに、はじめの状態を1つ使う。}
\end{bkfigure}

\bksubsection{6}{表計算で、はじめの状態から先を計算する}

\label{bk:5-13:hyou}式で積分できなくても、5-8の方法で、はじめの状態から先を計算できます。「短い時間 $\Delta t$ のあいだは、速度も加速度もほとんど変わらない」と考えて、次のようにくり返します。
\[ \text{次の速度}=\text{今の速度}+a\times\Delta t,\qquad \text{次の位置}=\text{今の位置}+\text{今の速度}\times\Delta t \]
投げ上げたボール（$a=-10$、$v_0=20$、$x_0=0$）を $\Delta t=0.5$ で計算します。たとえば3行めは、速度 $20+(-10)\times0.5=15$、位置 $0+20\times0.5=10$ です。
\begin{bkcenter}
\bksheet{5}{20}{%
1 & $t$ & 速度 $v$ & 位置 $x$（予想） & $20t-5t^2$ & ずれ \\
2 & 0 & 20 & 0 & 0 & 0 \\
3 & 0.5 & 15 & 10 & 8.75 & 1.25 \\
4 & 1.0 & 10 & 17.5 & 15 & 2.5 \\
5 & 1.5 & 5 & 22.5 & 18.75 & 3.75 \\
6 & 2.0 & 0 & 25 & 20 & 5 \\}
\bktabcaption{はじめの2つの値（$v_0=20$、$x_0=0$）がないと、計算を始められない。}
\end{bkcenter}
表の2行めの「はじめの速度」と「はじめの位置」が、積分定数 $C_1$、$C_2$ にあたります。予想の位置が少し大きいのは、0.5秒のあいだにへっていく速度を、はじめの（大きい）速度のままとして計算したからです。ずれは $\Delta t$ を小さくするとへります（5-8）。

\begin{column}{ゲームの中で、ボールはどう動いている？}
ボールを投げるゲームの中では、コンピュータが⑥と同じ計算をしています。画面は1秒に60回かきかえられるので、$\Delta t$ は $\dfrac{1}{60}$ 秒（約0.017秒）です。1回かきかえるたびに「次の速度＝今の速度＋加速度×$\dfrac{1}{60}$」「次の位置＝今の位置＋今の速度×$\dfrac{1}{60}$」を計算します。ボールを投げた瞬間に決まる「投げた場所」と「投げた速度」が、はじめの状態（$x_0$ と $v_0$）です。同じ場所・同じ速度で投げれば、何回やっても同じ所に落ちます。
\end{column}

\begin{column}{未来は、はじめの状態で決まる？}
200年ほど前、フランスのラプラスは「すべての物の位置と速度がわかれば、運動の法則で未来はすべて計算できる」と考えました。この章の考えを、宇宙全体に広げたものです。でも、はじめのわずかなずれが大きく育つ現象（天気など）もあり、未来の予想には限界があります。
\end{column}

\begin{taisetsu}
積分定数 $C$ は、はじめの状態（初期条件）で決まる。\\
加速度が一定なら\quad $v=v_0+at$、\quad $x=x_0+v_0t+\dfrac{1}{2}at^2$\\
いまの量 ＝ はじめの量 ＋ 変化の速さの積分
\end{taisetsu}

\begin{kotoba}
\small
\bkword{初期条件}{積分定数を決めるのに使う、ある時刻（ふつうは $t=0$）の値。}
\bkword{$x_0$、$v_0$}{はじめの位置、はじめの速度。「エックス・ゼロ」「ブイ・ゼロ」と読む。}
\end{kotoba}

%% ---------------- 例題 ----------------
\bksection{reidai}{例題でたしかめよう}

\begin{reidai}[初期条件で $C$ を決める]
速度が $v=6t+2$（m/s）で、$t=0$ の位置が\qty{5}{m}の物体の、位置 $x$ を $t$ の式で表しましょう。\par
\kangaekata 位置は速度の原始関数（①）。項ごとに積分して $+C$ をつけ、$t=0$ の位置で $C$ を決めます。\par
\shiki $6t$ を積分すると $6\times\dfrac{t^2}{2}=3t^2$、2を積分すると $2t$\par
\hspace*{2.6\zw}$\displaystyle x=\int(6t+2)\,\dd t=3t^2+2t+C$\par
\hspace*{2.6\zw}$t=0$ のとき $x=5$ なので、$5=3\times0^2+2\times0+C$ より $C=5$\par
\hspace*{2.6\zw}確かめ：$x=3t^2+2t+5$ を微分すると $6t+2$\par
\kotae $x=3t^2+2t+5$
\end{reidai}

\begin{jibun}
速度 $v=4t-3$（m/s）、$t=0$ で位置\qty{2}{m}の物体の位置を求めましょう。
\bkjans{$x=2t^2-3t+C$、$C=2$。$x=2t^2-3t+2$}
\end{jibun}

\noindent\begin{minipage}[c]{0.48\linewidth}
\begin{tikzpicture}
\begin{axis}[bkplot, width=58mm, height=44mm, xmin=0, xmax=2.4, ymin=0, ymax=25, xtick={1,2}, ytick={5,10,20}, xlabel={$t$（s）}, ylabel={$x$（m）}]
  \addplot[bkcurve, domain=0:2.15] {3*x^2+2*x+5};
  \addplot[only marks, mark=*, mark size=2pt, bkInk] coordinates {(0,5)};
  \node[font=\sffamily\scriptsize, anchor=west, fill=white, inner sep=1pt] at (axis cs:0.06,15) {出発点 $(0,\,5)$};
  \draw[bkGray, line width=0.5pt] (axis cs:0.12,13.6) -- (axis cs:0.03,6);
\end{axis}
\end{tikzpicture}
\end{minipage}\hfill
\begin{minipage}[c]{0.5\linewidth}\raggedright
\bkreidairef{1}の答え $x=3t^2+2t+5$ のグラフです。$t=0$ のとき $x=5$ で、グラフは高さ5から始まります。$C=5$ は、グラフの出発点の高さです。$C$ を変えると、同じ形のまま、出発点だけが上下に動きます。
\end{minipage}

\begin{reidai}[$t=0$ でない初期条件]
速度 $v=3t^2-4$（m/s）、$t=1$ で位置0の物体の位置を求めましょう。\par
\kangaekata 積分して $+C$ をつけ、$t=1$、$x=0$ を入れて $C$ を決めます（②）。\par
\shiki $\displaystyle x=\int(3t^2-4)\,\dd t=3\times\frac{t^3}{3}-4t+C=t^3-4t+C$\par
\hspace*{2.6\zw}$t=1$ のとき $x=0$ なので、$0=1^3-4\times1+C=-3+C$\par
\hspace*{2.6\zw}両辺に3をたして $C=3$\par
\kotae $x=t^3-4t+3$（はじめの位置は $x(0)=\qty{3}{m}$ だった）
\end{reidai}

\begin{jibun}
速度が $v=t^2+1$（m/s）で、$t=3$ のとき位置が\qty{10}{m}の物体の位置を求めましょう。
\bkjans{$x=\dfrac{t^3}{3}+t+C$、$9+3+C=10$ より $C=-2$。$x=\dfrac{t^3}{3}+t-2$}
\end{jibun}

\begin{reidai}[加速度から位置]
加速度 \qty{-10}{m/s^2}（上向きが＋）で、はじめの位置0、はじめの速度\qty{20}{m/s}の物体の、速度と位置を求めましょう。\par
\kangaekata ③と同じく2回積分し、$C_1$ をはじめの速度で、$C_2$ をはじめの位置で決めます。\par
\shiki $\displaystyle v=\int(-10)\,\dd t=-10t+C_1$。$t=0$ で $v=20$ より $C_1=20$\par
\hspace*{2.6\zw}$\displaystyle x=\int(20-10t)\,\dd t=20t-10\times\frac{t^2}{2}+C_2=20t-5t^2+C_2$\par
\hspace*{2.6\zw}$t=0$ で $x=0$ より $C_2=0$\par
\kotae $v=20-10t$、$x=20t-5t^2$（5-5の投げ上げと同じ式）
\end{reidai}

\begin{jibun}
加速度 \qty{2}{m/s^2}、はじめの速度\qty{3}{m/s}、はじめの位置\qty{1}{m}の物体の、速度と位置を求めましょう。
\bkjans{$v=3+2t$、$x=1+3t+t^2$}
\end{jibun}

\begin{reidai}[いちばん高くなる所]
\bkreidairef{3}のボール（$v=20-10t$、$x=20t-5t^2$）が、いちばん高くなる時刻と、そのときの高さを求めましょう。\par
\kangaekata 上がっているあいだは速度が＋、落ちてくるあいだは $-$ です。＋から $-$ に変わるとき、速度は0になります。その瞬間が、いちばん高い所です（5-7）。\par
\shiki $v=0$ になる時刻\quad $20-10t=0$ より $10t=20$、$t=2$\par
\hspace*{2.6\zw}そのときの高さ\quad $x=20\times2-5\times2^2=40-20=20$\par
\kotae 2秒後、高さ\qty{20}{m}
\end{reidai}

\begin{jibun}
手の高さ\qty{1.5}{m}から、上向きに\qty{10}{m/s}でボールを投げます（$a=-10$）。速度と位置を $t$ の式で表し、いちばん高くなる時刻と高さを求めましょう。
\bkjans{$v=10-10t$、$x=1.5+10t-5t^2$。\mbox{1秒後}、\qty{6.5}{m}}
\end{jibun}

\begin{bkfigure}
\begin{tikzpicture}
\begin{axis}[bkplot, name=g1, width=58mm, height=42mm, xmin=0, xmax=4.5, ymin=0, ymax=25, xtick={1,2,3,4}, ytick={10,20}, xlabel={$t$（s）}, ylabel={$x$（m）}, title={\sffamily\small 位置（高さ）}]
  \addplot[bkcurve, domain=0:4] {20*x-5*x^2};
  \draw[bkGray, dashed, line width=0.5pt] (axis cs:2,0) -- (axis cs:2,20);
  \addplot[only marks, mark=*, mark size=1.8pt, bkInk] coordinates {(2,20)};
\end{axis}
\begin{axis}[bkplot, at={(g1.east)}, anchor=west, xshift=20mm, width=58mm, height=42mm, xmin=0, xmax=4.5, ymin=-24, ymax=24, xtick={1,2,3,4}, ytick={20}, xlabel={$t$（s）}, ylabel={$v$（m/s）}, title={\sffamily\small 速度}]
  \addplot[bkcurve2, domain=0:4] {20-10*x};
  \draw[bkGray, dashed, line width=0.5pt] (axis cs:2,-24) -- (axis cs:2,0);
  \addplot[only marks, mark=*, mark size=1.8pt, bkInk] coordinates {(2,0)};
\end{axis}
\end{tikzpicture}
\bkfigcaption{\bkreidairef{4}のボール。速度が＋から $-$ に変わる2秒（速度0）で、高さがいちばん大きい。}
\end{bkfigure}

\begin{reidai}[ビルから落とす]
高さ\qty{44.1}{m}から、ボールを静かに落とします。地面に着くまでの時間と、そのときの速さを求めましょう（$g=9.8$）。\par
\kangaekata 上向きを＋にすると、④の式で $x_0=44.1$、$v_0=0$ です。$x=44.1-4.9t^2$、$v=-9.8t$。地面は $x=0$。\par
\shiki $44.1-4.9t^2=0$ より $4.9t^2=44.1$、$t^2=44.1\div4.9=9$\par
\hspace*{2.6\zw}2回かけて9になる正の数は3なので $t=3$\par
\hspace*{2.6\zw}$v=-9.8t=-9.8\times3=-29.4$（上向きが＋なので、$-$ は下向き）\par
\kotae 3秒後、下向きに\qty{29.4}{m/s}
\end{reidai}

\begin{jibun}
高さ\qty{19.6}{m}から静かに落としたボールが、地面に着くまでの時間と速さを求めましょう。
\bkjans{$4.9t^2=19.6$ より $t^2=4$、\mbox{2秒後}。$v=-9.8\times2$ で、下向きに\qty{19.6}{m/s}}
\end{jibun}

\noindent\begin{minipage}[c]{0.48\linewidth}
\begin{tikzpicture}
\begin{axis}[bkplot, width=58mm, height=46mm, xmin=0, xmax=3.5, ymin=0, ymax=50, xtick={1,2,3}, ytick={24.5,44.1}, xlabel={$t$（s）}, ylabel={$x$（m）}]
  \addplot[bkcurve, domain=0:3] {44.1-4.9*x^2};
  \addplot[only marks, mark=*, mark size=1.8pt, bkInk] coordinates {(0,44.1) (1,39.2) (2,24.5) (3,0)};
\end{axis}
\end{tikzpicture}
\end{minipage}\hfill
\begin{minipage}[c]{0.5\linewidth}\raggedright
\bkreidairef{5}のボールの高さです。1秒ごとの高さは44.1、39.2、24.5、\qty{0}{m}。1秒ごとに落ちる距離は4.9、14.7、\qty{24.5}{m}で、だんだん大きくなります。速度が1秒ごとに\qty{9.8}{m/s}ずつ速くなるからです。
\end{minipage}

\begin{reidai}[2つの物体が出会う時刻]
物体Aは位置0から、速度 $v=2t$ で＋の向きに動きます。物体Bは位置\qty{15}{m}から、速度\qty{-2}{m/s}で動きます。2つが出会う時刻を求めましょう。\par
\kangaekata 積分と初期条件で、それぞれの位置の式をつくります。出会うのは、位置が等しくなるときです。\par
\shiki A：$x=t^2+C$、$t=0$ で0なので $x=t^2$\par
\hspace*{2.6\zw}B：$x=-2t+C$、$t=0$ で15なので $x=15-2t$\par
\hspace*{2.6\zw}$t^2=15-2t$ より $t^2+2t-15=0$、$(t+5)(t-3)=0$（3-10）\par
\hspace*{2.6\zw}$t=-5$ または $t=3$。時刻は0より後なので $t=3$。位置は $3^2=9$\par
\kotae 3秒後（位置\qty{9}{m}）
\end{reidai}

\begin{jibun}
Aは位置0から速度 $v=4t$ で、Bは位置\qty{30}{m}から速度\qty{-4}{m/s}で動きます。出会う時刻を求めましょう。
\bkjans{A：$x=2t^2$、B：$x=30-4t$。$2t^2=30-4t$、$t^2+2t-15=0$ より\mbox{3秒後}（位置\qty{18}{m}）}
\end{jibun}

\noindent\begin{minipage}[c]{0.48\linewidth}
\begin{tikzpicture}
\begin{axis}[bkplot, width=58mm, height=44mm, xmin=0, xmax=4.2, ymin=0, ymax=18, xtick={1,2,3,4}, ytick={9,15}, xlabel={$t$（s）}, ylabel={$x$（m）}]
  \addplot[bkcurve, domain=0:3.9] {x^2};
  \addplot[bkcurve2, domain=0:3.9] {15-2*x};
  \draw[bkGray, dashed, line width=0.5pt] (axis cs:3,0) -- (axis cs:3,9);
  \addplot[only marks, mark=*, mark size=1.8pt, bkInk] coordinates {(3,9)};
  \node[font=\sffamily\scriptsize, text=bkBlue, anchor=west] at (axis cs:3.92,15.2) {A};
  \node[font=\sffamily\scriptsize, text=bkOrange!80!black, anchor=west] at (axis cs:3.92,7.2) {B};
\end{axis}
\end{tikzpicture}
\end{minipage}\hfill
\begin{minipage}[c]{0.5\linewidth}\raggedright
\bkreidairef{6}の2つの物体の $x$-$t$ グラフです。2本のグラフが交わる点が、同じ時刻に同じ位置にいる「出会い」です。交わる点 $(3,\,9)$ が、計算の答え（3秒後、位置\qty{9}{m}）と合っています。
\end{minipage}

\begin{reidai}[整数をためして解く]
物体Aは位置0から速度 $v=3t^2$ で、物体Bは位置\qty{30}{m}から速度\qty{-1}{m/s}で動きます。2つが出会う時刻を求めましょう（出会う時刻は整数秒です）。\par
\kangaekata 位置の式をつくると、$t^3$ の入った式になり、因数分解の公式が使えません。「整数秒」とわかっているので、$t=1,\,2,\,3,\,\ldots$ を順に入れて、位置が等しくなる時刻をさがします。\par
\shiki A：$x=t^3$（$t^3$ を微分すると $3t^2$）、B：$x=30-t$\par
\hspace*{2.6\zw}$t=1$：Aは1、Bは29\quad $t=2$：Aは8、Bは28\par
\hspace*{2.6\zw}$t=3$：Aは $3^3=27$、Bは $30-3=27$\quad 等しい\par
\kotae 3秒後（位置\qty{27}{m}）
\end{reidai}

\begin{jibun}
Aは位置0から速度 $v=3t^2$ で、Bは位置\qty{12}{m}から速度\qty{-2}{m/s}で動きます。出会う時刻を求めましょう（整数秒です）。
\bkjans{A：$x=t^3$、B：$x=12-2t$。$t=2$ で $8=8$。\mbox{2秒後}（位置\qty{8}{m}）}
\end{jibun}

\noindent\begin{minipage}[c]{0.48\linewidth}
\begin{tikzpicture}
\begin{axis}[bkplot, width=58mm, height=44mm, xmin=0, xmax=4.2, ymin=0, ymax=40, xtick={1,2,3,4}, ytick={10,27}, xlabel={$t$（s）}, ylabel={$x$（m）}]
  \addplot[bkcurve, domain=0:3.35] {x^3};
  \addplot[bkcurve2, domain=0:3.9] {30-x};
  \draw[bkGray, dashed, line width=0.5pt] (axis cs:3,0) -- (axis cs:3,27);
  \addplot[only marks, mark=*, mark size=1.8pt, bkInk] coordinates {(3,27)};
  \node[font=\sffamily\scriptsize, text=bkBlue, anchor=west] at (axis cs:3.38,37.5) {A};
  \node[font=\sffamily\scriptsize, text=bkOrange!80!black, anchor=south] at (axis cs:3.85,28.6) {B};
\end{axis}
\end{tikzpicture}
\end{minipage}\hfill
\begin{minipage}[c]{0.5\linewidth}\raggedright
\bkreidairef{7}のグラフです。Aははじめゆっくりで、だんだん速くなります（速度 $3t^2$）。Bは一定の速さでもどってきます。$t=1,\,2$ ではBのほうが上（先）にいて、$t=3$ で追いつきます。表やグラフで「ためす」のも、りっぱな解き方です。
\end{minipage}

\begin{reidai}[加速度が変わるとき]
加速度が $a=6t$（\unit{m/s^2}）で、はじめの速度\qty{2}{m/s}、はじめの位置0の物体の、速度と位置を求めましょう。2秒後の値も答えましょう。\par
\kangaekata $a$ が変わるので、④の $x=x_0+v_0t+\frac{1}{2}at^2$ は使えません（あれは $a$ が一定のときの式）。③と同じく2回積分します。\par
\shiki $\displaystyle v=\int6t\,\dd t=3t^2+C_1$。$t=0$ で $v=2$ より $C_1=2$\par
\hspace*{2.6\zw}$\displaystyle x=\int(3t^2+2)\,\dd t=t^3+2t+C_2$。$t=0$ で $x=0$ より $C_2=0$\par
\hspace*{2.6\zw}$t=2$ では $v=3\times2^2+2=14$、$x=2^3+2\times2=12$\par
\kotae $v=3t^2+2$、$x=t^3+2t$。2秒後は\qty{14}{m/s}、\qty{12}{m}
\end{reidai}

\begin{jibun}
加速度が $a=12t$（\unit{m/s^2}）で、はじめの速度0、はじめの位置\qty{1}{m}の物体の、速度と位置を求めましょう。1秒後の値も答えましょう。
\bkjans{$v=6t^2$、$x=2t^3+1$。\mbox{1秒後}は\qty{6}{m/s}、\qty{3}{m}}
\end{jibun}

\begin{reidai}[水そうの水位]
水そうの水位が、$\dfrac{\dd h}{\dd t}=0.5-0.1t$（cm/min）の速さで変わります。はじめの水位が\qty{10}{cm}のとき、10分後の水位を求めましょう。\par
\kangaekata いまの水位＝はじめの水位＋変わる速さの積分（⑤）。\par
\shiki $\displaystyle h=10+\int_0^{10}(0.5-0.1t)\,\dd t=10+\Bigl[\,0.5t-0.05t^2\,\Bigr]_0^{10}$\par
\hspace*{2.6\zw}$=10+(0.5\times10-0.05\times10^2)=10+(5-5)=10$\par
\kotae \qty{10}{cm}（はじめの5分は上がり、5分後からは下がって、もとにもどった）
\end{reidai}

\begin{jibun}
お湯の温度が、$\dfrac{\dd T}{\dd t}=3-0.1t$（℃/min）で変わります。はじめが\qty{20}{\degreeCelsius}のとき、10分後の温度を求めましょう。
\bkjans{$20+\Bigl[\,3t-0.05t^2\,\Bigr]_0^{10}=20+25=\qty{45}{\degreeCelsius}$}
\end{jibun}

\begin{reidai}[たまる量がいちばん多くなるとき]
タンクの水の量が $\dfrac{\dd V}{\dd t}=30-6t$（L/min）で変わり、はじめは\qty{50}{L}です。水の量がいちばん多くなるのは何分後で、そのとき何Lですか。\par
\kangaekata 変わる速さが＋のあいだはふえ、$-$ になるとへります。速さが0になる所で、ふえるからへるに変わります（5-7、\bkreidairef{4}と同じ考え）。\par
\shiki $\displaystyle V=\int(30-6t)\,\dd t=30t-3t^2+C$。$t=0$ で $V=50$ より $C=50$\par
\hspace*{2.6\zw}変わる速さ $30-6t=0$ より $t=5$\par
\hspace*{2.6\zw}$V=50+30\times5-3\times5^2=50+150-75=125$\par
\kotae 5分後、\qty{125}{L}
\end{reidai}

\begin{jibun}
水位の変わる速さが $4-2t$（cm/min）で、はじめの水位は\qty{10}{cm}です。水位がいちばん高くなる時刻と水位を求めましょう。
\bkjans{$h=10+4t-t^2$。$4-2t=0$ より\mbox{2分後}、\qty{14}{cm}}
\end{jibun}

\begin{reidai}[あとの時刻の様子から、さかのぼる]
ボールを真上に投げました（$a=-10$、上向きが＋）。投げてから2秒後に、いちばん高い\qty{25}{m}の所に来ました。はじめの速度 $v_0$ と、投げた高さ $x_0$ を求めましょう。\par
\kangaekata 初期条件が「2秒後」の値で2つ（いちばん高い所なので速度0、高さ25）わかっています。②のように、わかっている時刻の値を入れて $C_1$、$C_2$ を決めます。$t=0$ を入れれば、$v_0$ と $x_0$ がわかります（$C_1=v_0$、$C_2=x_0$）。\par
\shiki $\displaystyle v=\int(-10)\,\dd t=-10t+C_1$\par
\hspace*{2.6\zw}$t=2$ で $v=0$ なので $0=-10\times2+C_1$、$C_1=20$\par
\hspace*{2.6\zw}$\displaystyle x=\int(-10t+20)\,\dd t=-5t^2+20t+C_2$\par
\hspace*{2.6\zw}$t=2$ で $x=25$ なので $25=-5\times4+20\times2+C_2$\par
\hspace*{2.6\zw}$25=20+C_2$、両辺から20をひいて $C_2=5$\par
\hspace*{2.6\zw}確かめ：$x=5+20t-5t^2$ に $t=2$ を入れると $5+40-20=25$\par
\kotae はじめの速度\qty{20}{m/s}（上向き）、投げた高さ\qty{5}{m}
\end{reidai}

\begin{jibun}
ボールを真上に投げると（$a=-10$）、1秒後に、いちばん高い\qty{6.5}{m}の所に来ました。はじめの速度と、投げた高さを求めましょう。
\bkjans{$v=-10t+C_1$、$0=-10+C_1$ より $C_1=10$。$x=-5t^2+10t+C_2$、$6.5=-5+10+C_2$ より $C_2=1.5$。はじめの速度\qty{10}{m/s}、高さ\qty{1.5}{m}}
\end{jibun}

%% ---------------- まちがえやすいところ・物理工学 ----------------
\bksection{tsumazuki}{まちがえやすいところをチェック}

\begin{tsumazuki}{$C$ を、いつも0にしてしまう}
\bkng{$v=2t$、$t=0$ で位置\qty{3}{m} → $x=t^2$}\bkok{$x=t^2+C$、$C=3$ → $x=t^2+3$}
$x=t^2$ とすると、はじめの位置を0と決めつけたことになります。初期条件がわかればそれで $C$ を決め、なければ $+C$ を残します。
\end{tsumazuki}

\begin{tsumazuki}{2回積分するのに、$C$ を1つしか使わない}
加速度から位置までは、$C_1$ と $C_2$ の2つが出ます。$C_1$ ははじめの速度、$C_2$ ははじめの位置で決めます。1つの $C$ にまとめてはいけません。
\end{tsumazuki}

\begin{tsumazuki}{速度が変わるのに「位置＝速度×時間」とする}
\bkng{$v=2t$ の3秒後は $v=6$ なので $x=6\times3=\qty{18}{m}$}\bkok{$x=\displaystyle\int_0^3 2t\,\dd t=\qty{9}{m}$}
「速度×時間」は、速度が一定のときだけ使えます。3秒後の速さ6でずっと走っていたわけではないからです。
\end{tsumazuki}

\begin{butsuri}
\begin{itemize}[leftmargin=1.2\zw, itemsep=1pt, topsep=0pt, label={\color{bkPurple}●}]
\item \textbf{慣性航法}\\飛行機やロケットは、加速度センサーの値を2回積分して位置を計算します。はじめの位置と速度を正確に入れておきます。
\item \textbf{物理の法則と初期条件}\\物理の法則は「変化の速さ」（運動なら加速度）を決めます。法則とはじめの状態がそろって、はじめて先のようすが1つに決まります。これが第6巻の微分方程式の考え方です（6-1）。
\item \textbf{ダム}\\入る水と出る水の差を積分して、はじめの水位にたし、これからの水位を予想します。
\end{itemize}
\end{butsuri}

%% ---------------- 章末問題 ----------------
\bkendproblems
\bkproblemgroup{基本}{3}

\bkproblem 速度が $v=6t$（m/s）で、$t=0$ の位置が\qty{4}{m}の物体の位置を求めましょう。
\bkpans{$x=3t^2+4$}

\bkproblem 速度が $v=4t+1$（m/s）で、$t=0$ の位置が\qty{-3}{m}の物体の位置を求めましょう。
\bkpans{$x=2t^2+t-3$}

\bkproblem 速度が $v=3t^2-2t$（m/s）で、$t=2$ のとき位置が\qty{5}{m}の物体の位置を求めましょう。
\bkpans{$x=t^3-t^2+C$、$8-4+C=5$ より $C=1$。$x=t^3-t^2+1$}

\bkproblem 高さ\qty{2}{m}から、上向きに\qty{15}{m/s}でボールを投げます（$a=-10$）。速度と位置を $t$ の式で表し、いちばん高くなる時刻と高さを求めましょう。
\bkpans{$v=15-10t$、$x=2+15t-5t^2$。\mbox{1.5秒後}、\qty{13.25}{m}}

\bkproblem 高さ\qty{78.4}{m}から静かに落としたボールが、地面に着くまでの時間を求めましょう（$g=9.8$）。
\bkpans{$4.9t^2=78.4$ より $t^2=16$、\mbox{4秒後}}

\bkproblem 高さ\qty{4.9}{m}から静かに落としたボールが、地面に着くまでの時間と、そのときの速さを求めましょう（$g=9.8$）。
\bkpans{$4.9t^2=4.9$ より $t^2=1$、\mbox{1秒後}。下向きに\qty{9.8}{m/s}}

\bkproblem 加速度\qty{3}{m/s^2}、はじめの速度\qty{4}{m/s}、はじめの位置0の物体の、2秒後の速度と位置を求めましょう。
\bkpans{$v=4+3t$、$x=4t+1.5t^2$。\qty{10}{m/s}、\qty{14}{m}}

\bkproblem お湯の温度が $\dfrac{\dd T}{\dd t}=2-0.2t$（℃/min）で変わります。はじめが25\,℃のとき、10分後の温度を求めましょう。
\bkpans{$25+\Bigl[\,2t-0.1t^2\,\Bigr]_0^{10}=25+10=\mbox{$35$\,℃}$}

\bkproblemgroup{チャレンジ}{3}

\bkproblem 加速度が $a=6t$（\unit{m/s^2}）で、はじめの速度\qty{-3}{m/s}、はじめの位置0の物体の、速度と位置を求めましょう。2秒後の値も答えましょう。
\bkpans{$v=3t^2-3$、$x=t^3-3t$。2秒後は\qty{9}{m/s}、\qty{2}{m}}

\bkproblem ロケットの加速度が $a=2t$（m/s$^2$）でふえていきます。はじめの速度と位置は0です。6秒後の速度と位置を求めましょう。
\bkpans{$v=t^2$、$x=\dfrac{t^3}{3}$。\qty{36}{m/s}、\qty{72}{m}}

\bkproblem タンクの水の量が、$\dfrac{\dd V}{\dd t}=20-4t$（L/min）で変わり、はじめは\qty{100}{L}です。水の量がいちばん多くなるのは何分後で、そのとき何 L ですか（5-7）。
\bkpans{$V=100+20t-2t^2$。$20-4t=0$ より\mbox{5分後}、\qty{150}{L}}

\bkproblem 水位が $\dfrac{\dd h}{\dd t}=6-2t$（cm/min）で変わり、はじめの水位は\qty{5}{cm}です。5分後の水位と、水位がいちばん高くなる時刻と水位を求めましょう。
\bkpans{$h=5+6t-t^2$。5分後は\qty{10}{cm}。$6-2t=0$ より\mbox{3分後}、\qty{14}{cm}}

\bkproblem 物体Aは位置0から速度 $v=2t$ で、物体Bは位置\qty{8}{m}から速度\qty{2}{m/s}で、同じ向きに動きます。AがBに追いつく時刻と位置を求めましょう。
\bkpans{A：$x=t^2$、B：$x=8+2t$。$t^2-2t-8=0$、$(t-4)(t+2)=0$ より\mbox{4秒後}（位置\qty{16}{m}）}

\bkproblem 物体Aは位置0から速度 $v=3t^2$ で、物体Bは位置\qty{16}{m}から速度\qty{-4}{m/s}で動きます。2つが出会う時刻を求めましょう（出会う時刻は整数秒です）。
\bkpans{A：$x=t^3$、B：$x=16-4t$。$t=2$ でAは8、Bは $16-8=8$。\mbox{2秒後}（位置\qty{8}{m}）}

\bkproblem ボールを真上に投げると（$a=-10$、上向きが＋）、3秒後に、いちばん高い\qty{50}{m}の所に来ました。はじめの速度と、投げた高さを求めましょう。
\bkpans{$v=-10t+C_1$、$0=-30+C_1$ より $C_1=30$。$x=-5t^2+30t+C_2$、$50=-45+90+C_2$ より $C_2=5$。はじめの速度\qty{30}{m/s}、高さ\qty{5}{m}}

\bkproblem 速度が $v=2t+1$（m/s）の物体が、$t=2$ のとき位置\qty{10}{m}にいました。$t=0$ の位置と、$t=4$ の位置を求めましょう。
\bkpans{$x=t^2+t+C$、$4+2+C=10$ より $C=4$。$t=0$ で\qty{4}{m}、$t=4$ で $16+4+4=\qty{24}{m}$}

\begin{furikaeri}
\begin{checklist}
\item 初期条件から、積分定数を決められる。\hfill{\footnotesize\sffamily\color{bkTheme}→\pageref{bk:5-13:shoki}ページ}
\item 加速度から、速度と位置を求められる。\hfill{\footnotesize\sffamily\color{bkTheme}→\pageref{bk:5-13:toukasoku}ページ}
\item 「いまの量＝はじめの量＋変化の積分」を使える。\hfill{\footnotesize\sffamily\color{bkTheme}→\pageref{bk:5-13:teisekibun}ページ}
\end{checklist}
\end{furikaeri}

\bkfukidashi{delta}{第2部はここまで。次の部では、$e$ や $\sin$・$\cos$ という特別な関数を、微分・積分してみるよ。}


%% ======================== 第3部 特別な関数 ========================
\bkPartOpener{3}{第5巻\hspace{0.8em}第3部}{特別な関数}%
  {172/{微分しても形が変わらない関数はある？},
   160/{波を微分すると、どんな形になる？},
   148/{むずかしい関数を、かんたんな式で近似できる？}}%
  {% 指数関数と波（白い線）
   \begin{scope}[white, line width=1.3pt, line cap=round]
     \draw plot[domain=0:3.2, samples=40, smooth] ({118+\x*15},{112+3*exp(\x*0.75)});
     \draw[line width=1pt] plot[domain=0:50, samples=60, smooth] ({118+\x},{128+7*sin(\x*14.4)});
   \end{scope}}%
  {この部では、e・sin・cos という特別な関数を微分・積分するよ。どれも、自然の中でよく出てくる関数なんだ。最後に、この巻の道具を1枚の表にまとめるよ。}
\bkPartContents{3}{%% 自動生成（gen.py）
\bktocchap{5-14}{ふしぎな数 e}{微分しても形が変わらない、特別な関数}{3}{3}{4}{334}{C}
\bktocchap{5-15}{sin・cosを微分すると？}{微分すると、波が1/4周期ずれる}{4}{4}{4}{362}{C}
\bktocchap{5-16}{積分の道具箱：置換積分と部分積分}{むずかしい積分を、知っている形に変える}{2}{4}{4}{392}{C}
\bktocchap{5-17}{近くでは多項式で近似できる：テイラー展開}{小さくゆれる振り子は、$\sin\theta\approx\theta$で計算できる}{1}{4}{4}{420}{C}
\bktocchap{5-18}{まとめ：微分したときのふるまい}{微分方程式の「答えの候補」がそろった}{1}{3}{3}{444}{C}
}%
{ふえる・へる・ゆれるという、自然のいちばん基本的な変化を、微分と積分で書けるようになります。第6巻では、この部の表が「微分方程式の答えの候補」になります。}
\bkpartbegin{3}{特別な関数}
%% 5-14 ふしぎな数 e（第11回：ページ制限なし。X17：複利の「半年で50％」を、1年あたりの割合から1段ずつ説明する）
\bkchapter{5}{14}{ふしぎな数 $e$}{\bklevelrange{3}{4}}{ふしぎな数 e}
\bkmichishirube{3-11,3-12,5-6}{5-17,6-4,9-5}

\begin{goal}
\begin{itemize}[label={\color{bkTheme}●}, leftmargin=1.6\zw, itemsep=0pt, topsep=1pt]
\item \bklevel{3}\enspace $e$ が $\Bigl(1+\dfrac{1}{n}\Bigr)^n$ の近づく先であることを、表計算で確かめられる。
\item \bklevel{3}\enspace $(e^{ax})'=ae^{ax}$ を使って、ふえる量・へる量の変化の速さを求められる。
\item \bklevel{4}\enspace 自然対数 $\ln x$ を微分・積分できる。\bkidx{たいすう}{対数}
\end{itemize}
\end{goal}

\begin{hatena}
\noindent\begin{minipage}[c]{0.68\linewidth}\raggedright
利率が「1年あたり100\,\%」という、ゆめのような銀行があります。1年たつと、預けたお金と同じだけの利息（ふえる分）がつきます。1万円を預けると、1年後に2万円です。

この銀行が、利息を半年ごとに計算してくれることになりました。1年あたりの利率は100\,\%のままです。1年後のお金は、2万円より多くなるでしょうか。毎月、毎日……と、もっと細かく計算してもらえば、いくらでも大きくなるのでしょうか。
\end{minipage}\hfill
\begin{minipage}[c]{0.3\linewidth}\centering
\begin{tikzpicture}[x=1mm, y=1mm]
  \draw[bkInk, line width=0.6pt] (-1,0) -- (36,0);
  \fill[bkYellow!60] (0,0) rectangle (8,12);
  \draw[bkYellow!70!black, line width=0.6pt] (0,0) rectangle (8,12);
  \node[font=\sffamily\scriptsize, anchor=south] at (4,13) {1万円};
  \foreach \px in {13,26} {\draw[bkGray, line width=0.6pt, dash pattern=on 2pt off 1.5pt] (\px,0) rectangle (\px+8,22); \node[font=\sffamily\bfseries\small, text=bkGray] at (\px+4,11) {?};}
  \node[font=\sffamily\scriptsize, anchor=north] at (4,-1) {はじめ};
  \node[font=\sffamily\scriptsize, anchor=north] at (17,-1) {半年};
  \node[font=\sffamily\scriptsize, anchor=north] at (30,-1) {1年};
\end{tikzpicture}
\end{minipage}
\end{hatena}

この章では、半年ごとの計算を1段ずつ追いかけ、その先にあらわれる数 $e$（ネイピア数）を調べます。$e$ は、放射線のへり方、お湯の冷め方、電気回路など、自然のあらゆる所に顔を出します。

%% ---------------- やってみよう・思い出そう ----------------
\begin{yattemiyou}[半年ごとに利息をつけてもらおう]
「1年あたり100\,\%」の利率で、1万円を預けます。銀行は、半年ごとに利息を計算して、預けているお金にたしてくれます。電卓を使ってよいです。
\begin{enumerate}[label={\bkcircled[bkGreen]{\arabic*}}, leftmargin=2.4\zw, itemsep=1pt, topsep=3pt]
\item 半年は、1年の何分の1かな？　半年分の利率は、100\,\%の何分の1になるかな？
\item はじめの半年：1万円に、半年分の利息をつけよう。半年後にはいくら？
\item 次の半年：半年後のお金に、また半年分の利息をつけよう。1年後にはいくら？
\item 1年に1回だけ計算したとき（2万円）とくらべよう。差は、どこから出たのかな？
\end{enumerate}
\begin{bkcenter}
\begin{bktable}{colspec={Q[l,m,wd=40mm] *{3}{Q[c,m,wd=20mm]}}, rowsep=1.5pt}
 & はじめ & 半年後 & 1年後 \\
1年に1回計算する（万円） & 1 & 1 & 2 \\
半年ごとに計算する（万円） & 1 & ？ & ？ \\
\end{bktable}
\end{bkcenter}
\end{yattemiyou}

\noindent ①半年は1年の $\dfrac{1}{2}$ なので、半年分の利率は $100\,\%\times\dfrac{1}{2}=50\,\%$ です。②1万円の50\,\%は0.5万円なので、半年後は $1+0.5=1.5$ 万円。③次の半年も50\,\%です。こんどは1.5万円に50\,\%がつくので、利息は $1.5\times0.5=0.75$ 万円、1年後は $1.5+0.75=2.25$ 万円です。④1年に1回のときより、0.25万円多くなりました。

\begin{bkfigure}
\begin{tikzpicture}[x=1mm, y=1mm]
  \def\u{16}% 1万円の高さ（mm）
  \draw[bkInk, line width=0.7pt] (-4,0) -- (98,0);
  % はじめ
  \fill[bkYellow!60] (0,0) rectangle (14,\u); \draw[bkYellow!70!black, line width=0.6pt] (0,0) rectangle (14,\u);
  % 半年
  \fill[bkYellow!60] (40,0) rectangle (54,\u); \fill[bkOrange!55] (40,\u) rectangle (54,1.5*\u);
  \draw[bkYellow!70!black, line width=0.6pt] (40,0) rectangle (54,\u); \draw[bkOrange!80!black, line width=0.6pt] (40,\u) rectangle (54,1.5*\u);
  % 1年
  \fill[bkYellow!60] (80,0) rectangle (94,\u); \fill[bkOrange!55] (80,\u) rectangle (94,1.5*\u);
  \fill[bkBlue!35] (80,1.5*\u) rectangle (94,2*\u); \fill[bkRed!55] (80,2*\u) rectangle (94,2.25*\u);
  \draw[bkYellow!70!black, line width=0.6pt] (80,0) rectangle (94,\u); \draw[bkOrange!80!black, line width=0.6pt] (80,\u) rectangle (94,1.5*\u);
  \draw[bkBlue!80!black, line width=0.6pt] (80,1.5*\u) rectangle (94,2*\u); \draw[bkRed!80!black, line width=0.6pt] (80,2*\u) rectangle (94,2.25*\u);
  % 合計
  \node[font=\sffamily\small, anchor=south] at (7,\u+1.5) {1万円};
  \node[font=\sffamily\small, anchor=south] at (47,1.5*\u+1.5) {1.5万円};
  \node[font=\sffamily\bfseries\small, anchor=south] at (87,2.25*\u+1.5) {2.25万円};
  % 50％ふえる（×1.5）
  \foreach \a/\b in {16/38, 56/78} {
    \node[font=\sffamily\scriptsize, align=center, anchor=south, text=bkInk] at ({(\a+\b)/2},8.5) {50\,\%ふえる\\（×1.5）};
    \draw[bk arrow, bkGray, line width=0.8pt] (\a,6) -- (\b,6);}
  % 時間の目もり
  \node[font=\sffamily\small, anchor=north] at (7,-1.5) {はじめ};
  \node[font=\sffamily\small, anchor=north] at (47,-1.5) {半年};
  \node[font=\sffamily\small, anchor=north] at (87,-1.5) {1年};
  % 凡例
  \begin{scope}[shift={(0,-14)}]
    \foreach \c/\t/\px/\py in {bkYellow!60/はじめの1万円/0/0, bkOrange!55/1回めの利息0.5万円/52/0, bkBlue!35/2回めの利息のうち、1万円の分0.5万円/0/-6, bkRed!55/2回めの利息のうち、利息の分0.25万円/0/-12} {
      \fill[\c] (\px,\py-1.6) rectangle ++(3.2,3.2);
      \node[font=\sffamily\scriptsize, anchor=west, inner sep=0pt] at (\px+4.8,\py) {\t};}
  \end{scope}
\end{tikzpicture}
\bkfigcaption{半年ごとに計算すると、2回めの利息は、1万円だけでなく、1回めの利息0.5万円にもつく（赤）。}
\end{bkfigure}

\noindent 1年に1回のときとの差の0.25万円（赤）は、1回めの利息0.5万円に、2回めの半年でついた50\,\%の利息（$0.5\times0.5=0.25$）です。利息にも利息がつくのです。このような利息のつけ方を\textbf{複利}（ふくり）といいます。\bkidx{ふくり}{複利}

\begin{omoidasou}
\bkrecall{2-5}{50\,\%ふえる＝もとの1倍＋0.5倍＝1.5倍。20\,\%引き＝0.8倍と同じ考え。}
\bkrecall{2-7}{同じ数を何回もかけるのが累乗。\\ $1.5\times1.5=1.5^2$、\quad $1.25\times1.25\times1.25\times1.25=1.25^4$}
\bkrecall{3-11}{指数関数 $y=a^x$。$a^{x+h}=a^x\times a^h$、$(a^m)^n=a^{mn}$、$a^0=1$。たとえば $2^{3+1}=2^3\times2^1=8\times2=16$}
\bkrecall{3-12}{$\log_a b$ は「$a$ を何乗すると $b$ になるか」で、$a^{\log_a b}=b$。\\ $\log_a(bc)=\log_a b+\log_a c$、\quad $\log_a b^r=r\log_a b$}
\bkrecall{5-6}{合成関数の微分：外側を微分して、内側の微分をかける。}
\bkrecall{5-8}{$h$ が小さいとき、$f(a+h)\approx f(a)+f'(a)h$（接線で近似）}
\end{omoidasou}

\clearpage
%% ---------------- 数学の道具 ----------------
\bksection{doguu}{数学の道具}

\bksubsection{1}{「1年あたり100\,\%」を、半年に分ける}

\label{bk:5-14:bunkatsu}「1年あたり100\,\%」の「あたり」は、時速の「あたり」と同じです。時速\qty{60}{km}は「1時間あたり\qty{60}{km}」でした。30分（0.5時間）走ると、
\[ \qty{60}{km/h}\times\qty{0.5}{h}=\qty{30}{km} \]
進みます。時間が半分なら、進む距離も半分です。利率も同じで、「1年あたり100\,\%」の利率で半年（0.5年）預けると、つく利息は
\[ 100\,\%\times0.5=50\,\% \]
です。\textbf{時間が半分なので、利息の割合も半分の50\,\%}になります。

\noindent\textbf{なぜ「1.5倍」なのか}\quad 50\,\%の利息がつくと、お金は「もとの1倍」に「0.5倍の利息」をたした $1+0.5=1.5$ 倍になります（2-5）。だから、半年ごとに1.5倍が2回で
\[ 1\ \text{万円}\times1.5\times1.5=1\ \text{万円}\times1.5^2=2.25\ \text{万円} \]
です。

\noindent\textbf{まちがえやすい考え方1}\quad 「はじめの半年は50\,\%、次の半年は100\,\%」と考えると、$1.5\times2=3$ 万円になってしまいます。100\,\%は「1年あたり」の割合でした。2回めも半年（0.5年）なので、2回めの利息も50\,\%です。半年ごとの計算では、毎回50\,\%です。

\noindent\textbf{まちがえやすい考え方2}\quad 「50\,\%と50\,\%で、合わせて100\,\%なので2万円」とすると、2回めの利息が1万円だけにつくと考えたことになります。でも、2回めの半年が始まるときには、預けているお金はもう1.5万円です。2回めの利息は、1.5万円の50\,\%（0.75万円）です。

\bksubsection{2}{もっと細かく分けると：$\Bigl(1+\dfrac{1}{n}\Bigr)^n$}

\label{bk:5-14:n}\textbf{3か月ごと（1年に4回）}\quad 3か月は1年の $\dfrac{1}{4}$ なので、1回分の利率は $100\,\%\times\dfrac{1}{4}=25\,\%$、1回ごとに $1+0.25=1.25$ 倍です。

\begin{bkcenter}
\begin{bktable}{colspec={Q[l,m,wd=22mm] Q[c,m,wd=38mm] Q[c,m,wd=30mm]}, rowsep=1.5pt, row{1}={bg=bkTheme!10}}
時期 & 計算 & お金（万円） \\
はじめ & & 1 \\
3か月後 & $1\times1.25$ & 1.25 \\
6か月後 & $1.25\times1.25$ & 1.5625 \\
9か月後 & $1.5625\times1.25$ & 1.953125 \\
1年後 & $1.953125\times1.25$ & 2.44140625 \\
\end{bktable}
\bktabcaption{3か月ごとに1.25倍。1年後は $1.25^4=2.4414\cdots$ 倍。}
\end{bkcenter}

\noindent\textbf{1年を $n$ 回に分ける}\quad 同じように考えると、
\begin{enumerate}[label={\bkcircled[bkTheme]{\arabic*}}, leftmargin=2.4\zw, itemsep=1pt, topsep=2pt]
\item 1回分の時間は1年の $\dfrac{1}{n}$ なので、1回分の利率は $100\,\%\times\dfrac{1}{n}$、つまり $\dfrac{1}{n}$ です（100\,\%＝1）。
\item 1回ごとに、お金は $\Bigl(1+\dfrac{1}{n}\Bigr)$ 倍になります。
\item それを $n$ 回くり返すので、1年後には $\Bigl(1+\dfrac{1}{n}\Bigr)^n$ 倍です。指数の $n$ は「$n$ 回かける」の $n$ です（2-7）。
\end{enumerate}
表計算で、$n$ を大きくしていきます。
\begin{bkcenter}
\bksheet{3}{28}{%
1 & ふやし方 & $n$ & $\left(1+\frac{1}{n}\right)^n$ \\
2 & 1年に1回 & 1 & 2 \\
3 & 半年ごと & 2 & 2.25 \\
4 & 3か月ごと & 4 & 2.4414… \\
5 & 毎月 & 12 & 2.6130… \\
6 & 毎日 & 365 & 2.7145… \\
7 & （細かく） & 10000 & 2.71814… \\
8 & （もっと細かく） & 1000000 & 2.718280… \\}
\bktabcaption{$n$ を大きくすると、2.71828… に近づいていく。}
\end{bkcenter}

\begin{bkfigure}
\begin{tikzpicture}
\begin{axis}[bkplot, width=100mm, height=46mm, xmin=0, xmax=52, ymin=1.8, ymax=3, xtick={1,10,20,30,40,50}, ytick={2,2.25,2.5,2.718}, yticklabels={2,2.25,2.5,$e$}, xlabel={分ける回数 $n$}, ylabel={1年後（万円）}]
  \addplot[bkthin, dashed, domain=0:52] {2.71828};
  \addplot[only marks, mark=*, mark size=1.3pt, bkBlue, samples at={1,2,...,50}] {(1+1/x)^x};
\end{axis}
\end{tikzpicture}
\bkfigcaption{$n$ をふやすと $\bigl(1+\frac{1}{n}\bigr)^n$ はふえていくが、点線の $e=2.718\cdots$ をこえない。}
\end{bkfigure}

細かく分けるほどお金はふえますが、ふえ方はだんだん小さくなります。1回から2回で0.25、2回から4回で約0.19、4回から12回で約0.17ふえますが、365回から1000000回にしても0.004ほどしかふえません。1回分の利息がとても小さくなるので、「利息の利息」も、ほんの少ししかふえないのです。$n$ をどんなに大きくしても3にもとどかず、2.71828… に近づいていきます。この近づく先を $e$ と書きます。
\[ e=\lim_{n\to\infty}\Bigl(1+\frac{1}{n}\Bigr)^n=2.718281828\cdots \]
$e$ は、$\pi$ と同じように、小数がどこまでも続き、くり返さない数です。これが\bkhatenaref の答えです。いくら細かく計算してもらっても、1年後のお金は約2.718万円より大きくはなりません。\bkidx{いー}{$e$（ネイピア数）}

\noindent\textbf{なぜ細かく分けて考えるの？}\quad お金のふえ方は銀行の決まりしだいですが、自然の変化は「1年に1回」まとめて起きるわけではありません。細菌は、いまいる数に比例して、いつも少しずつふえています。お湯は、いつも少しずつ冷めています。分け方をどんどん細かくしていくと、「いつも少しずつ、いまある量に比例してふえる」ようすに近づきます。その行き着く先にあらわれるのが $e$ です。このことを、③からの微分で確かめます。

\bksubsection{3}{指数関数を微分すると、自分に比例する}

\label{bk:5-14:shisuu}$f(x)=a^x$ を微分してみましょう。5-3の決め方どおり、$\dfrac{f(x+h)-f(x)}{h}$ を計算して、$h$ を0に近づけます。指数のきまり（3-11）
\[ a^{x+h}=a^x\times a^h\qquad（たとえば\ 2^{3+1}=2^3\times2^1=8\times2=16） \]
を使うと、
\[ \frac{a^{x+h}-a^x}{h}=\frac{a^x\times a^h-a^x}{h}=a^x\times\frac{a^h-1}{h} \]
最後の式は、分子の2つの項に共通な $a^x$ をくくり出しました。右の $\dfrac{a^h-1}{h}$ には $x$ が入っていません。だから、$a$ が決まれば、どの $x$ でも同じ数です。$h$ を0に近づけたときの近づく先を $k_a$ と書くと、
\[ (a^x)'=k_a\,a^x \]
\textbf{指数関数を微分すると、自分自身の $k_a$ 倍になります}。「いまある量に比例してふえる」（3-11）ことの、微分による表し方です。また、$x=0$ を入れると $a^0=1$ なので $f'(0)=k_a$。$k_a$ は $x=0$ での接線の傾きです。

\noindent\textbf{数で確かめる}\quad $a=2$ で、$h=0.0001$ とします。電卓で $2^{0.0001}=1.0000693\cdots$ なので、
\[ \frac{2^{0.0001}-1}{0.0001}=\frac{0.0000693\cdots}{0.0001}=0.693\cdots \]
$y=2^x$ の各点の傾きと高さをくらべると、どこでも「傾き÷高さ」が0.693です。
\begin{bkcenter}
\begin{bktable}{colspec={Q[l,m,wd=34mm] *{4}{Q[c,m,wd=14mm]}}, rowsep=1.5pt}
$x$ & 0 & 1 & 2 & 3 \\
高さ $2^x$ & 1 & 2 & 4 & 8 \\
傾き & 0.693 & 1.386 & 2.773 & 5.545 \\
傾き÷高さ & 0.693 & 0.693 & 0.693 & 0.693 \\
\end{bktable}
\bktabcaption{$y=2^x$ の傾きは、いつも高さの0.693倍。}
\end{bkcenter}

ほかの底についても、$k_a$ を表計算で $h=0.0001$ として調べました。
\begin{bkcenter}
\begin{bktable}{colspec={Q[c,m,wd=24mm] *{5}{Q[c,m,wd=14mm]}}, rowsep=1.5pt}
底 $a$ & 2 & 2.5 & 2.7 & 2.8 & 3 \\
$k_a\approx\dfrac{a^h-1}{h}$ & 0.693 & 0.916 & 0.993 & 1.030 & 1.099 \\
\end{bktable}
\bktabcaption{底が2.7と2.8のあいだで、$k_a$ がちょうど1になりそう。}
\end{bkcenter}

\bksubsection{4}{$e$ は「微分しても変わらない」底}

\label{bk:5-14:e}$k_a=1$ になる底を見つければ、その指数関数は、微分しても自分自身のままです。その底が、ちょうど②の $e=2.71828\cdots$ です。$a=e$ で計算すると、$\dfrac{e^{0.0001}-1}{0.0001}=1.00005$ で、たしかに1に近くなります。
\[ (e^x)'=e^x \]

\noindent\textbf{なぜ利息の $e$ と同じなのか}\quad $k_e=1$ は、$y=e^x$ の $x=0$ での接線が、点 $(0,\,1)$ を通る傾き1の直線 $y=1+x$ だということです。接点の近くでは、グラフと接線はほとんど重なるので、小さい $h$ では（5-8）
\[ e^h\approx1+h \]
たとえば $e^{0.1}=1.10517\cdots$ と $1.1$、$e^{0.01}=1.01005\cdots$ と $1.01$ です。ここで、$n$ を大きな数として、$h$ に $\dfrac{1}{n}$ を入れると、
\[ e^{\frac{1}{n}}\approx1+\frac{1}{n} \]
です。両辺を $n$ 乗します。左辺は、$(a^m)^n=a^{mn}$（3-11）を使うと
\[ \Bigl(e^{\frac{1}{n}}\Bigr)^n=e^{\frac{1}{n}\times n}=e^1=e \]
なので、
\[ e\approx\Bigl(1+\frac{1}{n}\Bigr)^n \]
$n$ が大きいほど近似は正確になるので、②の表の数が $e$ に近づいたのです。

\begin{bkfigure}
\begin{tikzpicture}
\begin{axis}[bkplot, width=86mm, height=52mm, xmin=-2.2, xmax=2.3, ymin=0, ymax=7.6, xtick={-2,-1,1,2}, ytick={2,4,6}, xlabel={$x$}, ylabel={$y$}, yticklabel style={fill=white, inner sep=1pt, outer sep=2.3pt}]
  \addplot[bkthin, domain=-2.1:2.1] {2^x};
  \addplot[bkthin, dashed, domain=-2.1:1.75] {3^x};
  \addplot[bkcurve, domain=-2.1:2] {exp(x)};
  \addplot[bkcurve2, line width=1pt, domain=-1:1.6] {1+x};
  \node[font=\sffamily\scriptsize, anchor=west] at (axis cs:2.2,4.29) {$2^x$};
  \node[font=\sffamily\scriptsize, anchor=east] at (axis cs:1.6,6.9) {$3^x$};
  \node[font=\sffamily\scriptsize, anchor=west, text=bkBlue] at (axis cs:2.08,7.3) {$e^x$};
  % 格子の線が「＋」の上を通らないように、空いている左上（x<0、y=4〜6）に置く
  \node[font=\sffamily\scriptsize, align=center, text=bkOrange!80!black, fill=white, inner sep=1pt] at (axis cs:-1.5,5) {$y=1+x$\\（傾き1）};
\end{axis}
\end{tikzpicture}
\bkfigcaption{$y=e^x$ は、$(0,\,1)$ での接線の傾きがちょうど1になる指数関数。$2^x$ はそれよりゆるく、$3^x$ は急。}
\end{bkfigure}

\begin{seikaku}
\textbf{$e$ の決め方}\par
この本では、$e$ を「$k_a=1$ になる底」、つまり $\displaystyle\lim_{h\to0}\frac{e^h-1}{h}=1$ となる数として決めました。大学の教科書では、$e=\displaystyle\lim_{n\to\infty}\Bigl(1+\frac{1}{n}\Bigr)^n$ で $e$ を決め、そこから $(e^x)'=e^x$ を導くことが多いです。どちらから始めても、同じ数 $2.71828\cdots$ になり、同じ性質がすべて導けます。また、$e$ は分数で表せない数（2-9）であることが証明されています。
\end{seikaku}

\bksubsection{5}{$e^{ax}$ の微分}

\label{bk:5-14:eax}合成関数の微分（5-6）を使います。内側を $u=ax$、外側を $e^u$ とします。外側 $e^u$ を $u$ で微分すると $e^u$（④）、内側 $ax$ を $x$ で微分すると $a$ です。2つをかけて、
\[ (e^{ax})'=e^{ax}\times a=a\,e^{ax} \]
たとえば $a=2$ なら $(e^{2x})'=2e^{2x}$ で、$x=0$ での傾きは2です。表計算で $\dfrac{e^{0.0002}-1}{0.0001}=2.0002$ と確かめられます。

$a>0$ なら、$e^{ax}$ はふえていき、ふえる速さはいつも自分の $a$ 倍です。$a<0$ ならへっていき、へる速さが自分に比例します。放射線を出す物質の量 $N$ が $N=N_0e^{-\lambda t}$（$\lambda$ はラムダ）でへるのは、「へる速さが、いまある量に比例する」からです。\bkidx{しすうかんすう}{指数関数の微分}

\noindent\textbf{例}\quad $(e^{3x})'=3e^{3x}$、\quad $(5e^{-2t})'=5\times(-2)e^{-2t}=-10e^{-2t}$\par
\noindent\hspace*{2.5\zw}$(e^{x^2})'=e^{x^2}\times(x^2)'=e^{x^2}\times2x$（内側を $u=x^2$ とする）

\bksubsection{6}{自然対数 $\ln$}

\label{bk:5-14:ln}$e$ を底とする対数 $\log_e x$ を\textbf{自然対数}といい、$\ln x$ と書きます（エル・エヌ・エックス。natural logarithm の略）。「$e$ を何乗すると $x$ になるか」です（3-12）。\bkidx{しぜんたいすう}{自然対数 $\ln$}
\[ e^{\ln x}=x,\qquad \ln e^{y}=y,\qquad \ln1=0,\qquad \ln e=1 \]
$\ln1=0$ なのは、$e^0=1$（0乗すると1）だからです。$\ln e=1$ なのは、$e^1=e$ だからです。

たとえば $\ln2=0.693\cdots$ は、「$e$ を約0.693乗すると2になる」ということです。電卓で確かめると、
\[ 2.71828^{0.693}=1.9997\cdots\approx2 \]
です。電卓や表計算では、\texttt{LN(2)} で求められます。

\noindent\textbf{どんな指数関数も $e$ で書ける}\par
$a=e^{\ln a}$ なので、3-11の $(a^m)^n=a^{mn}$ を使うと
\[ a^x=(e^{\ln a})^x=e^{x\ln a} \]
です。たとえば $2=e^{0.693}$ なので、$2^x=e^{0.693x}$ です。これを、⑤の $(e^{ax})'=ae^{ax}$ の $a$ のかわりに $\ln a$ として微分すると、
\[ (a^x)'=(e^{x\ln a})'=\ln a\times e^{x\ln a}=a^x\ln a \]
③の表の $k_a$ は、$\ln a$ だったのです。たしかに
\[ k_2=0.693=\ln2,\qquad k_3=1.099=\ln3 \]
です。

\bksubsection{7}{$\ln x$ の微分：グラフを折り返して考える}

\label{bk:5-14:lnbibun}$y=\ln x$ は、「$x=e^y$」と同じことです。$x$ と $y$ の役割を入れかえた関係なので、$y=\ln x$ のグラフは、$y=e^x$ のグラフを、直線 $y=x$ で折り返したものです（3-12）。

\noindent\begin{minipage}[c]{0.41\linewidth}
\centering
\begin{tikzpicture}
\begin{axis}[bkplot, width=56mm, height=56mm, xmin=-1.3, xmax=4.4, ymin=-1.3, ymax=4.4, xtick={1,2,3,4}, ytick={1,2,3,4}, xticklabels={,2,3,4}, yticklabels={,2,3,4}, xlabel={$x$}, ylabel={$y$}, tick label style={fill=white, inner sep=1pt, outer sep=2.3pt}]
  \addplot[bkthin, dashed, domain=-1.3:4.2] {x};
  \addplot[bkcurve2, domain=-1.3:1.45] {exp(x)};
  \addplot[bkcurve, domain=0.28:4.3] {ln(x)};
  \addplot[bkcurve2, line width=0.8pt, domain=0.15:1.25] {2+2*(x-0.693)};
  \addplot[bkcurve, line width=0.8pt, domain=1:3.2] {0.693+0.5*(x-2)};
  \addplot[only marks, mark=*, mark size=1.4pt, bkInk] coordinates {(0.693,2) (2,0.693)};
  \node[font=\sffamily\scriptsize, text=bkOrange!80!black, anchor=west, fill=white, inner sep=1pt] at (axis cs:1.6,4.0) {$e^x$};
  \node[font=\sffamily\scriptsize, text=bkBlue, anchor=south, fill=white, inner sep=1pt] at (axis cs:3.9,1.75) {$\ln x$};
\end{axis}
\end{tikzpicture}\par\vspace{1pt}
{\sffamily\scriptsize\color{bkOrange!80!black}点 $(\ln2,\,2)$ での傾き2}\par
{\sffamily\scriptsize\color{bkBlue}点 $(2,\,\ln2)$ での傾き $\frac{1}{2}$}
\end{minipage}\hfill
\begin{minipage}[c]{0.56\linewidth}
折り返すと、横の長さとたての長さが入れかわります。だから、接線の傾き（たて÷横）は、逆数になります。$y=e^x$ の点 $(\ln2,\,2)$ での傾きは $e^{\ln2}=2$ なので、折り返した $y=\ln x$ の点 $(2,\,\ln2)$ での傾きは $\dfrac{1}{2}$ です。

どの点でも同じです。$y=e^x$ の点 $(p,\ e^p)$ での傾きは $e^p$。折り返すと、この点は $y=\ln x$ の点 $(e^p,\ p)$ になり、傾きは $\dfrac{1}{e^p}$。この点の $x$ 座標は $e^p$ なので、傾きは $\dfrac{1}{x}$ です。
\end{minipage}
\[ (\ln x)'=\frac{1}{x}\qquad（x>0） \]
\noindent\textbf{数で確かめる}\quad $x=2$ で、$\dfrac{\ln2.001-\ln2}{0.001}=0.49988\cdots$。たしかに $\dfrac{1}{2}$ に近い値です。

\bksubsection{8}{どんな指数でも $(x^r)'=rx^{r-1}$}

\label{bk:5-14:xr}5-6で、$(\sqrt{x})'$ を求めるのに、$(x^n)'=nx^{n-1}$ が分数の $n$ でも使えるとしました。その理由がここでわかります。$x>0$ で、$r$ をどんな数としても、⑥と同じように $x=e^{\ln x}$ なので $x^r=e^{r\ln x}$ です。内側を $u=r\ln x$（微分すると $r\times\dfrac{1}{x}$）とする合成関数の微分で、
\[ (x^r)'=e^{r\ln x}\times\Bigl(r\times\frac{1}{x}\Bigr)=x^r\times\frac{r}{x}=rx^{r-1} \]
最後は、$x^r\div x=x^{r-1}$（2-12の指数のきまり）を使いました。

\bksubsection{9}{$e^{ax}$ と $\frac{1}{x}$ の積分}

\label{bk:5-14:sekibun}微分の表を逆に読むと（5-11）、
\[ \int e^{ax}\,\dd x=\frac{1}{a}e^{ax}+C,\qquad \int\frac{1}{x}\,\dd x=\ln x+C\quad（x>0） \]
\noindent\textbf{なぜ $\dfrac{1}{a}$ がつくの？}\quad $e^{ax}$ を微分すると $a$ 倍になってしまうので、はじめに $\dfrac{1}{a}$ をかけておけば、微分したときにちょうどもとにもどります。微分して確かめます。
\[ \Bigl(\frac{1}{a}e^{ax}\Bigr)'=\frac{1}{a}\times a\,e^{ax}=e^{ax} \]
たとえば
\begin{align*}
a=2\ \text{なら}&\quad \int e^{2x}\,\dd x=\frac{1}{2}e^{2x}+C\\
a=-1\ \text{なら}&\quad \int e^{-x}\,\dd x=\frac{1}{-1}e^{-x}+C=-e^{-x}+C
\end{align*}
です。

\noindent\textbf{$\frac{1}{x}$ の積分}\par
5-11で「$n=-1$ のときは使えない」とした $\displaystyle\int x^{-1}\,\dd x$ の答えが、自然対数でした。⑦で $(\ln x)'=\dfrac{1}{x}$ だったからです。

\noindent\begin{minipage}[c]{0.44\linewidth}
\begin{tikzpicture}
\begin{axis}[bkplot, width=54mm, height=42mm, xmin=0, xmax=2.6, ymin=0, ymax=1.3, xtick={1,1.5,2}, ytick={0.5,1}, xlabel={$x$}, ylabel={$y$}]
  \foreach \k in {0,...,3} {\edef\tmp{\noexpand\fill[bkBlue!20, draw=bkBlue!60, line width=0.4pt] (axis cs:{1+0.25*\k},0) rectangle (axis cs:{1.25+0.25*\k},{1/(1.125+0.25*\k)});}\tmp}
  \addplot[bkcurve, domain=0.8:2.5] {1/x};
\end{axis}
\end{tikzpicture}
\end{minipage}\hfill
\begin{minipage}[c]{0.52\linewidth}\raggedright
面積で確かめます。$y=\dfrac{1}{x}$ の1から2までの下の面積を、はば0.25の長方形4つで近似します。高さは、区間のまん中（1.125、1.375、1.625、1.875）での $y$ の値です。
\end{minipage}

\noindent 4つの長方形の面積をたすと、
\[ \Bigl(\frac{1}{1.125}+\frac{1}{1.375}+\frac{1}{1.625}+\frac{1}{1.875}\Bigr)\times0.25=0.691\cdots \]
積分で求めると $\Bigl[\,\ln x\,\Bigr]_1^2=\ln2-\ln1=0.693\cdots-0=0.693\cdots$ で、ほとんど同じです。

\bksubsection{10}{「自分に比例してふえる」を表計算でつくる}

\label{bk:5-14:euler}$(e^x)'=e^x$ は、「ふえる速さが、いまの値と同じ」ということです。この性質だけを使って、5-8の方法（次の値＝今の値＋変化の速さ×$h$）で、$x=0$ から1まで計算してみましょう。はじめの値は1、変化の速さは「今の値」です。
\[ \text{次の値}=\text{今の値}+\text{今の値}\times h=\text{今の値}\times(1+h) \]
1回ごとに $(1+h)$ 倍になるので、$h=\dfrac{1}{n}$ で $n$ 回くり返すと、$x=1$ での値は $\Bigl(1+\dfrac{1}{n}\Bigr)^n$ です。

\begin{bkcenter}
\begin{bktable}{colspec={Q[c,m,wd=20mm] Q[c,m,wd=20mm] Q[c,m,wd=40mm]}, rowsep=1.5pt}
$h$ & 回数 $n$ & $x=1$ での値 $(1+h)^n$ \\
0.1 & 10 & 2.5937… \\
0.01 & 100 & 2.7048… \\
0.001 & 1000 & 2.7169… \\
\end{bktable}
\bktabcaption{$h$ を小さくすると、$e=2.71828\cdots$ に近づく。}
\end{bkcenter}
\bkhatenaref の銀行の計算（②）と、まったく同じ式になりました。「1年を $n$ 回に分けて、そのたびに今のお金に比例してふやす」ことは、「ふえる速さが自分に比例する」を少しずつ計算することだったのです。$n$ を大きくした極限で、ぴったり $e$ になります。

\noindent\begin{minipage}[c]{0.46\linewidth}
\begin{tikzpicture}
\begin{axis}[bkplot, width=54mm, height=46mm, xmin=0, xmax=1.2, ymin=0.8, ymax=3, xtick={0.25,0.5,0.75,1}, xticklabels={,0.5,,1}, ytick={1,2.44,2.72}, yticklabels={1,2.44,$e$}, xlabel={$x$}, ylabel={値}]
  \addplot[bkcurve, domain=0:1.1] {exp(x)};
  \addplot[bkcurve2, line width=1pt, mark=*, mark size=1.4pt] coordinates {(0,1) (0.25,1.25) (0.5,1.5625) (0.75,1.953125) (1,2.44140625)};
\end{axis}
\end{tikzpicture}
\end{minipage}\hfill
\begin{minipage}[c]{0.5\linewidth}\raggedright
オレンジは $h=0.25$（4回）で、1回ごとに1.25倍した折れ線です。$x=1$ で2.44になり、②の「3か月ごと」と同じ数です。青い $y=e^x$ より少し下にあります。$h$ を小さくすると、折れ線は青い曲線に近づきます。
\end{minipage}

\Needspace{14\baselineskip}
\bksubsection{11}{へっていく $e^{-t/\tau}$ と時定数}

\label{bk:5-14:tau}\noindent\begin{minipage}[c]{0.52\linewidth}
\begin{tikzpicture}
\begin{axis}[bkplot, width=64mm, height=46mm, xmin=0, xmax=5.4, ymin=0, ymax=1.15, xtick={1,2,3,4,5}, xticklabels={$\tau$,$2\tau$,$3\tau$,$4\tau$,$5\tau$}, ytick={0.368,1}, yticklabels={0.37,1}, xlabel={$t$}, ylabel={$\frac{N}{N_0}$}]
  \addplot[bkcurve, domain=0:5.2] {exp(-x)};
  \addplot[bkcurve2, line width=1pt, domain=0:1] {1-x};
  \draw[bkGray, dashed, line width=0.5pt] (axis cs:1,0) -- (axis cs:1,0.368) -- (axis cs:0,0.368);
\end{axis}
\end{tikzpicture}
\end{minipage}\hfill
\begin{minipage}[c]{0.44\linewidth}
へっていく量は、$N=N_0e^{-t/\tau}$ と書くことがよくあります。$\tau$（タウ）は時間の単位をもつ数で、\textbf{時定数}といいます。\bkidx{じていすう}{時定数 $\tau$}
\begin{itemize}[leftmargin=1.2\zw, itemsep=1pt, topsep=2pt, label={\color{bkTheme}●}]
\item $t=\tau$ で、$e^{-1}=0.368$ 倍（約37\,\%）になる。
\item $t=0$ での接線（オレンジ）は、ちょうど $t=\tau$ で0になる。
\item $5\tau$ たつと $e^{-5}=0.0067$ 倍で、ほぼ0。
\end{itemize}
\end{minipage}\par
\noindent\textbf{なぜ接線が $\tau$ で0になるの？}\quad ⑤で $a=-\dfrac{1}{\tau}$ とすると、$N'=N_0\times\Bigl(-\dfrac{1}{\tau}\Bigr)e^{-t/\tau}$ です。$t=0$ では $e^0=1$ なので、へる速さは $\dfrac{N_0}{\tau}$。はじめの量 $N_0$ を、この速さのままへらし続けると、$N_0\div\dfrac{N_0}{\tau}=\tau$ の時間で0になります。

「はじめの速さのままへり続けたら、$\tau$ で0になる」というのが、時定数の意味です。じっさいには、へるにつれてへる速さもおそくなるので、$\tau$ たっても37\,\%残ります。電気回路や計器の反応の速さは、時定数で表します（6-4）。

\begin{bkfigure}
\begin{tikzpicture}
\begin{axis}[bkplot, name=p1, at={(0mm,0mm)}, width=38mm, height=25mm, xmin=0, xmax=5.5, ymin=0, ymax=1.15, xtick=\empty, ytick=\empty, title={\sffamily\scriptsize 放射性物質の量}, title style={yshift=-1.5mm}]
  \addplot[bkcurve, line width=1.2pt, domain=0:5.3] {exp(-0.6*x)};
\end{axis}
\begin{axis}[bkplot, at={(45mm,0mm)}, name=p2, width=38mm, height=25mm, xmin=0, xmax=5.5, ymin=0, ymax=1.15, xtick=\empty, ytick=\empty, title={\sffamily\scriptsize お湯の温度}, title style={yshift=-1.5mm}]
  \addplot[bkcurve2, line width=1.2pt, domain=0:5.3] {0.25+0.75*exp(-0.7*x)};
  \addplot[bkthin, dashed, domain=0:5.3] {0.25};
\end{axis}
\begin{axis}[bkplot, at={(90mm,0mm)}, width=38mm, height=25mm, xmin=0, xmax=5.5, ymin=0, ymax=1.15, xtick=\empty, ytick=\empty, title={\sffamily\scriptsize コンデンサーの電圧}, title style={yshift=-1.5mm}]
  \addplot[bkcurve3, line width=1.2pt, domain=0:5.3] {1-exp(-0.9*x)};
  \addplot[bkthin, dashed, domain=0:5.3] {1};
\end{axis}
\end{tikzpicture}
\bkfigcaption{身のまわりの $e^{-t/\tau}$。横軸は時間。へる・冷める（差がへる）・たまる（足りない分がへる）が、同じ数学で書ける。}
\end{bkfigure}

\begin{column}{ネイピアと対数表}
$e$ の名前のもとになったスコットランドのネイピアは、400年ほど前、かけ算をたし算に変えて計算を楽にする「対数」を考え出し、20年かけて対数の表を作りました。天文学者のケプラーは、この表のおかげで惑星の計算を何年分も短くできたといわれます。$e$ そのものを使い始めたのは、その100年ほど後のスイスのオイラーで、$e$ という文字もオイラーが使い始めました。
\end{column}

\begin{senmonsho}[大学の教科書ふうの書き方]
\begin{bkhonbun}
指数関数 $\exp x\equiv e^x$ は、すべての実数 $x$ に対して $\dfrac{\dd}{\dd x}\exp x=\exp x$、$\exp0=1$ をみたす。その逆関数を $\ln x$ と書く。
\end{bkhonbun}
\bkyaku{$\exp x$（エクスポネンシャル・エックス）は $e^x$ の別の書き方。肩の式が長いときに使う：$e^{-t/\tau}=\exp(-t/\tau)$。「逆関数」は、入力と出力を入れかえた関数（$y=e^x$ と $x=e^y$）。}
\end{senmonsho}

\bkkigocards{\bkkigocard{$e$}{イー（ネイピア数）}{$\bigl(1+\frac{1}{n}\bigr)^n$ の近づく先。$2.718\cdots$。$(e^x)'=e^x$}{Napier's constant}{5-14}}
{\bkkigocard{$\ln x$}{エル・エヌ・エックス}{自然対数。$e$ を何乗すると $x$ になるか}{natural log}{5-14}}

\begin{taisetsu}
「1年あたり100\,\%」を $n$ 回に分けると、1回 $\dfrac{1}{n}$ ずつ、1年で $\Bigl(1+\dfrac{1}{n}\Bigr)^n$ 倍 → $e$ に近づく\\
$(e^x)'=e^x$\qquad $(e^{ax})'=a\,e^{ax}$\qquad $(\ln x)'=\dfrac{1}{x}$\qquad $(a^x)'=a^x\ln a$\\
$\displaystyle\int e^{ax}\,\dd x=\frac{1}{a}e^{ax}+C$\qquad $\displaystyle\int\frac{1}{x}\,\dd x=\ln x+C$
\end{taisetsu}

\begin{kotoba}
\small
\bkword{複利}{利息にも利息がつく、ふえ方。}
\bkword{$e$（ネイピア数）}{$(1+\frac{1}{n})^n$ の近づく先。2.71828…。$e^x$ は微分しても変わらない。}
\bkword{自然対数 $\ln x$}{$e$ を底とする対数。$\log_e x$。}
\bkword{$\exp x$}{$e^x$ の別の書き方。}
\end{kotoba}

%% ---------------- 例題 ----------------
\bksection{reidai}{例題でたしかめよう}

\begin{reidai}[利息を何回かに分ける]
「1年あたり100\,\%」の利率で1万円を預け、利息を3か月ごと（1年に4回）に計算してもらいます。1年後のお金を求めましょう。\par
\kangaekata 3か月は1年の $\dfrac{1}{4}$ なので、1回分の利率は $100\,\%\times\dfrac{1}{4}=25\,\%$ です（①②）。毎回、そのときのお金の25\,\%がつくので、1回ごとに1.25倍になります。\par
\shiki 1回分の利率\quad $100\,\%\times\dfrac{1}{4}=25\,\%$、1回ごとに $1+0.25=1.25$ 倍\par
\hspace*{2.6\zw}4回くり返すので\quad $1\times1.25^4=1.25\times1.25\times1.25\times1.25$\par
\hspace*{2.6\zw}$=2.44140625$\par
\kotae 約2.44万円（1年に1回の2万円、半年ごとの2.25万円より多い）
\end{reidai}

\begin{jibun}
「1年あたり100\,\%」の利率で1万円を預け、利息を4か月ごと（1年に3回）に計算してもらいます。1年後のお金を求めましょう。分数のまま計算してから、小数にしましょう。
\bkjans{1回分 $\dfrac{1}{3}$、$\Bigl(1+\dfrac{1}{3}\Bigr)^3=\Bigl(\dfrac{4}{3}\Bigr)^3=\dfrac{64}{27}\approx2.37$。約2.37万円}
\end{jibun}

\begin{reidai}[$e$ をふくむ関数の微分]
次の関数を微分しましょう。\quad（1）$e^{4x}$\quad（2）$3e^{-x}$\quad（3）$e^{x^2}$\par
\kangaekata（1）（2）は⑤の $(e^{ax})'=ae^{ax}$（（2）は $a=-1$）。（3）は内側 $u=x^2$ の合成関数（5-6）。\par
\shiki（1）$4e^{4x}$\quad（2）$3\times(-1)e^{-x}=-3e^{-x}$\par
\hspace*{2.6\zw}（3）外側 $e^u$ の微分 $e^u$ に、内側 $x^2$ の微分 $2x$ をかけて $e^{x^2}\times2x$\par
\kotae （1）$4e^{4x}$\quad（2）$-3e^{-x}$\quad（3）$2xe^{x^2}$
\end{reidai}

\begin{jibun}
次の関数を微分しましょう。\quad（1）$e^{5x}$\quad（2）$2e^{-3x}$\quad（3）$e^{2x+1}$
\bkjans{（1）\nobreak $5e^{5x}$\quad（2）\nobreak $-6e^{-3x}$\quad（3）\nobreak $2e^{2x+1}$}
\end{jibun}

\begin{reidai}[へる速さ]
ある物質の量が $N=1000e^{-0.1t}$（g、$t$ は時間）でへっていきます。はじめ（$t=0$）のへる速さを求めましょう。\par
\kangaekata へる速さは、$N$ を $t$ で微分した $\dfrac{\dd N}{\dd t}$。⑤の $(e^{ax})'=ae^{ax}$ で、$a=-0.1$ です。\par
\shiki $\dfrac{\dd N}{\dd t}=1000\times(-0.1)e^{-0.1t}=-100e^{-0.1t}$\par
\hspace*{2.6\zw}$t=0$ では $e^0=1$ なので $-100\times1=-100$\par
\kotae 1時間あたり\qty{100}{g}の速さでへる（$\frac{\dd N}{\dd t}=-0.1N$：へる速さは、いまの量の0.1倍）
\end{reidai}

\begin{jibun}
$N=800e^{-0.05t}$ について、$\dfrac{\dd N}{\dd t}$ を求め、$t=0$ での値を答えましょう。
\bkjans{$800\times(-0.05)e^{-0.05t}=-40e^{-0.05t}$、$t=0$ で $-40$}
\end{jibun}

\begin{reidai}[$\ln$ と $a^x$ の微分]
次の関数を微分しましょう。\quad（1）$\ln(2x+1)$\quad（2）$2^x$\par
\kangaekata（1）内側を $u=2x+1$ とする合成関数。外側 $\ln u$ の微分は $\dfrac{1}{u}$（⑦）。（2）$2^x=e^{x\ln2}$ に直す（⑥）。\par
\shiki（1）$\dfrac{1}{2x+1}\times(2x+1)'=\dfrac{1}{2x+1}\times2$\par
\hspace*{2.6\zw}（2）$(2^x)'=(e^{x\ln2})'=e^{x\ln2}\times\ln2=2^x\ln2$\par
\kotae （1）$\dfrac{2}{2x+1}$\quad（2）$2^x\ln2$（$\approx0.693\times2^x$）
\end{reidai}

\begin{jibun}
次の関数を微分しましょう。\quad（1）$\ln(x^2+1)$\quad（2）$10^x$
\bkjans{（1）\nobreak $\dfrac{2x}{x^2+1}$\quad（2）\nobreak $10^x\ln10$}
\end{jibun}

\begin{reidai}[$e^{ax}$ と $\frac{1}{x}$ の積分]
（1）$\displaystyle\int_0^1e^x\,\dd x$\quad（2）$\displaystyle\int_0^1e^{-x}\,\dd x$\quad（3）$\displaystyle\int_1^e\frac{1}{x}\,\dd x$ を求めましょう。\par
\kangaekata ⑨の積分の式で原始関数を求め、上端の値から下端の値をひきます（5-11）。（2）は $a=-1$ なので、$\dfrac{1}{a}=-1$ がつきます。\par
\shiki（1）$\Bigl[\,e^x\,\Bigr]_0^1=e^1-e^0=e-1$\par
\hspace*{2.6\zw}（2）$\Bigl[\,-e^{-x}\,\Bigr]_0^1=(-e^{-1})-(-e^0)=1-e^{-1}$\par
\hspace*{2.6\zw}（3）$\Bigl[\,\ln x\,\Bigr]_1^e=\ln e-\ln1=1-0$\par
\kotae （1）$e-1\approx1.718$\quad（2）$1-e^{-1}\approx0.632$\quad（3）1
\end{reidai}

\begin{jibun}
（1）$\displaystyle\int_0^2e^{x}\,\dd x$\quad（2）$\displaystyle\int_1^3\frac{1}{x}\,\dd x$\quad（3）$\displaystyle\int_0^1e^{2x}\,\dd x$ を求めましょう。
\bkjans{（1）\nobreak $e^2-1\approx6.39$\quad（2）\nobreak $\ln3\approx1.10$\quad（3）\nobreak $\dfrac{e^2-1}{2}\approx3.19$}
\end{jibun}

\begin{reidai}[接線で近似する]
$y=e^x$ のグラフの、$x=0$ の点での接線の式を求めましょう。これを使って $e^{0.2}$ の近似値を求め、電卓の値 1.2214… とくらべましょう。\par
\kangaekata 接線の式は $y=f(a)+f'(a)(x-a)$（5-8）。$f(x)=e^x$、$a=0$ です。\par
\shiki $f(0)=e^0=1$、$f'(x)=e^x$ より $f'(0)=1$\par
\hspace*{2.6\zw}接線は $y=1+1\times(x-0)=x+1$\par
\hspace*{2.6\zw}$x=0.2$ を入れて $e^{0.2}\approx0.2+1=1.2$\par
\kotae 接線 $y=x+1$、$e^{0.2}\approx1.2$（ずれは約0.02）
\end{reidai}

\begin{jibun}
例題6の接線を使って、$e^{-0.1}$ の近似値を求め、電卓の値 0.9048… とくらべましょう。
\bkjans{$e^{-0.1}\approx-0.1+1=0.9$（ずれは約0.005）}
\end{jibun}

\begin{reidai}[半減期]
$N=N_0e^{-0.1t}$（$t$ は時間）でへる物質が、はじめの半分になる時間（半減期）を求めましょう。\par\bkidx{はんげんき}{半減期}
\kangaekata 半分になるのは、$N_0e^{-0.1T}=\dfrac{1}{2}N_0$、つまり $e^{-0.1T}=\dfrac{1}{2}$ のとき。$T$ が指数にあるので、両辺の自然対数をとって $T$ を下ろします（⑥の $\ln e^y=y$）。\par
\shiki 両辺の自然対数をとって\quad $\ln e^{-0.1T}=\ln\dfrac{1}{2}$\par
\hspace*{2.6\zw}左辺は $-0.1T$。右辺は $\ln\dfrac{1}{2}=\ln2^{-1}=-\ln2$（3-12）\par
\hspace*{2.6\zw}だから $-0.1T=-\ln2$\par
\hspace*{2.6\zw}$T=\dfrac{\ln2}{0.1}=\dfrac{0.693}{0.1}=6.93$\par
\kotae 約6.9時間
\end{reidai}

\begin{jibun}
$N=N_0e^{-0.05t}$（$t$ は年）でへる物質の半減期を求めましょう。
\bkjans{\mbox{$\dfrac{\ln2}{0.05}\approx13.9$\,年}}
\end{jibun}

\begin{reidai}[時定数]
コンデンサーにたまった電気の量が、$Q=Q_0e^{-t/2}$（$t$ は秒）でへっていきます。時定数と、半分になるまでの時間を求めましょう。\par
\kangaekata $e^{-t/\tau}$ とくらべて $\tau$ を読む（⑪）。半分の時間は、例題7と同じく両辺の自然対数をとる。\par
\shiki $e^{-t/2}$ なので時定数 $\tau=2$ 秒\par
\hspace*{2.6\zw}半分になる時間 $T$ は $e^{-T/2}=\dfrac{1}{2}$\par
\hspace*{2.6\zw}$-\dfrac{T}{2}=\ln\dfrac{1}{2}=-\ln2$ より $T=2\ln2\approx2\times0.693\approx1.39$\par
\kotae 時定数2秒、半分になるのは約1.4秒後（半減期＝時定数×$\ln2$）
\end{reidai}

\begin{jibun}
温度計の針が、$\theta=\theta_0e^{-t/5}$（$t$ は秒）で目もりに近づいていきます。時定数と、ずれが半分になるまでの時間を求めましょう。
\bkjans{時定数5秒、\mbox{$5\ln2\approx3.5$\,秒}}
\end{jibun}

\begin{reidai}[たえずふえる量]
いまある量に比例して、たえずふえていく量は $y=y_0e^{rt}$ と書けます（$r$ はふえる割合。$(y_0e^{rt})'=ry_0e^{rt}=ry$ なので、ふえる速さはいまの量の $r$ 倍）。100万円が $r=0.02$（1年あたり）でたえずふえるとき、10年後の金額を求めましょう。\par
\kangaekata $y=y_0e^{rt}$ に、$y_0=100$（万円）、$r=0.02$、$t=10$ を入れます。\par
\shiki $100\times e^{0.02\times10}=100\times e^{0.2}=100\times1.2214=122.14$\par
\kotae 約122万円（電卓で $e^{0.2}$ を計算する。表計算では \texttt{EXP(0.2)}）
\end{reidai}

\begin{jibun}
人口が $r=0.05$（1年あたり）でたえずふえる町があります。いま1万人のとき、10年後の人口を求めましょう（$e^{0.5}=1.6487$）。
\bkjans{$10000\times e^{0.5}=16487$。約16487人}
\end{jibun}

\begin{reidai}[コンデンサーにたまる電圧]
電池につないだコンデンサーの電圧が、$V=10\,(1-e^{-t/2})$（V、$t$ は秒）でふえていきます。2秒後の電圧と、$t=0$ での電圧のふえる速さを求めましょう（$e^{-1}=0.368$）。\par
\kangaekata $1-e^{-t/2}$ は、1から「へっていく量」をひいた形。時定数 $\tau=2$ 秒たつと、最後の値の $1-0.368=0.632$（約63\,\%）までたまります。\par
\shiki $t=2$：$10\times(1-e^{-1})=10\times(1-0.368)=6.32$\par
\hspace*{2.6\zw}$\dfrac{\dd V}{\dd t}=10\times\Bigl(-\Bigl(-\dfrac{1}{2}\Bigr)e^{-t/2}\Bigr)=5e^{-t/2}$、$t=0$ で5\par
\kotae 約\qty{6.3}{V}、毎秒\qty{5}{V}（だんだんゆっくりになり、\qty{10}{V}に近づく）
\end{reidai}

\begin{jibun}
コンデンサーの電圧が $V=5\,(1-e^{-t/3})$（V）でふえていきます。3秒後の電圧と、$t=0$ でのふえる速さを求めましょう（$e^{-1}=0.368$）。
\bkjans{$5\times0.632\approx\qty{3.16}{V}$。$\dfrac{\dd V}{\dd t}=\dfrac{5}{3}e^{-t/3}$、$t=0$ で毎秒約\qty{1.67}{V}}
\end{jibun}

\begin{reidai}[2倍になるまでの時間]
細菌の数が $N=50e^{0.35t}$（$t$ は時間）でふえます。数が2倍になるまでの時間（倍加時間）と、$t=0$ でのふえる速さを求めましょう（$\ln2=0.693$）。\par
\kangaekata 半減期（例題7）と同じように、$e^{0.35T}=2$ の両辺の自然対数をとります。\par
\shiki $0.35T=\ln2$ より $T=\dfrac{0.693}{0.35}\approx1.98$\par
\hspace*{2.6\zw}$\dfrac{\dd N}{\dd t}=50\times0.35e^{0.35t}=17.5e^{0.35t}$、$t=0$ で17.5\par
\kotae 約2時間。1時間あたり17.5（いまの数の0.35倍の速さでふえる）
\end{reidai}

\begin{jibun}
$N=200e^{0.1t}$（$t$ は時間）でふえる量の、倍加時間と、$t=0$ でのふえる速さを求めましょう。
\bkjans{\mbox{$T=\dfrac{0.693}{0.1}\approx6.9$\,時間}。$\dfrac{\dd N}{\dd t}=20e^{0.1t}$、$t=0$ で1時間あたり20}
\end{jibun}

\begin{reidai}[$\frac{1}{x}$ の積分と、対数のきまり]
$\displaystyle\int_2^6\frac{1}{x}\,\dd x$ を求め、$\ln$ 1つで表しましょう。\par
\kangaekata ⑨で原始関数は $\ln x$。上端の値から下端の値をひくと $\ln6-\ln2$ になります。対数のきまり $\ln(bc)=\ln b+\ln c$（3-12）を使うと、$\ln6=\ln(3\times2)=\ln3+\ln2$ なので、ひき算がかんたんになります。\par
\shiki $\displaystyle\int_2^6\frac{1}{x}\,\dd x=\Bigl[\,\ln x\,\Bigr]_2^6=\ln6-\ln2$\par
\hspace*{2.6\zw}$\ln6=\ln(3\times2)=\ln3+\ln2$ なので\par
\hspace*{2.6\zw}$\ln6-\ln2=(\ln3+\ln2)-\ln2=\ln3$\par
\kotae $\ln3$（$\approx1.10$。$6\div2=3$ の $\ln$ になる）
\end{reidai}

\begin{jibun}
$\displaystyle\int_5^{10}\frac{1}{x}\,\dd x$ を求め、$\ln$ 1つで表しましょう。
\bkjans{$\ln10-\ln5=(\ln2+\ln5)-\ln5=\ln2\approx0.693$}
\end{jibun}

\begin{reidai}[ふつうの利率を、細かく分ける]
1年あたり6\,\%の利率で、100万円を預けます。（1）1年に1回、（2）毎月（1年に12回）利息を計算するとき、1年後の金額を求めましょう。（3）たえず（いつも少しずつ）ふえるときの金額 $100e^{0.06}$ とくらべましょう（$e^{0.06}=1.06184$）。\par
\kangaekata 1回分の利率は、①②と同じく「1年あたりの利率×1回分の時間（年）」です。毎月なら1か月は $\dfrac{1}{12}$ 年なので、$6\,\%\times\dfrac{1}{12}=0.5\,\%$。1回ごとに $1+0.005=1.005$ 倍になり、それを12回くり返します。\par
\shiki（1）$100\times1.06=106$\par
\hspace*{2.6\zw}（2）$100\times1.005^{12}=100\times1.06168=106.168$\par
\hspace*{2.6\zw}（3）$100\times1.06184=106.184$\par
\kotae （1）106万円\quad（2）約106.17万円\quad（3）約106.18万円（細かく分けるほど、$y_0e^{rt}$ の値に近づく）
\end{reidai}

\noindent 1年あたりの利率を $r$ として $n$ 回に分けると、1年後は $\Bigl(1+\dfrac{r}{n}\Bigr)^n$ 倍です。$n$ を大きくすると、これは $e^r$ に近づきます。\bkreidairef{9}の「たえずふえる量」$y_0e^{rt}$ は、この行き着く先でした（$r=1$ のときが②の $e$）。

\begin{jibun}
1年あたり12\,\%の利率で100万円を預けます。毎月（1年に12回）利息を計算するとき、1回分の利率と、1年後の金額を求めましょう（$1.01^{12}=1.12683$）。
\bkjans{1回分 $12\,\%\times\dfrac{1}{12}=1\,\%$。$100\times1.01^{12}\approx112.68$。約112.68万円}
\end{jibun}

%% ---------------- まちがえやすいところ・物理工学 ----------------
\bksection{tsumazuki}{まちがえやすいところをチェック}

\begin{tsumazuki}{半年ごとの利率を、2回めだけ100\,\%にしてしまう}
\bkng{はじめの半年は50\,\%、次の半年は100\,\%で、$1.5\times2=3$ 万円}\bkok{半年は毎回50\,\%で、$1.5\times1.5=2.25$ 万円}
100\,\%は「1年あたり」の割合です。半年はいつでも1年の半分なので、何回めでも50\,\%です（①）。
\end{tsumazuki}

\begin{tsumazuki}{$e^x$ に $(x^n)'=nx^{n-1}$ を使ってしまう}
\bkng{$(e^x)'=xe^{x-1}$}\bkok{$(e^x)'=e^x$}
$x^n$ は「$x$ を $n$ 回かける」で、変わるのは底の $x$。$e^x$ は「$e$ を $x$ 回かける」で、変わるのは指数の $x$。形は似ていても、まったくちがう関数です。
\end{tsumazuki}

\begin{tsumazuki}{$\ln(a+b)=\ln a+\ln b$ としてしまう}
対数で和になるのは、かけ算のときです：$\ln(ab)=\ln a+\ln b$（3-12）。$\ln(a+b)$ は、それ以上かんたんにはなりません。$\ln(1+1)=\ln2=0.693$ で、$\ln1+\ln1=0$ ではありません。
\end{tsumazuki}

\begin{tsumazuki}{時定数と半減期をまちがえる}
時定数 $\tau$ は「約37\,\%（$e^{-1}$）になる時間」、半減期は「50\,\%になる時間（$\tau\ln2\approx0.69\tau$）」で、ちがう量です。
\end{tsumazuki}

\begin{butsuri}
\begin{itemize}[leftmargin=1.2\zw, itemsep=1pt, topsep=0pt, label={\color{bkPurple}●}]
\item \textbf{放射線と年代測定}\\炭素14は、半減期約5730年で、$e^{-\lambda t}$ の形でへります。木や骨に残っている量から、何年前のものかを調べられます。
\item \textbf{冷めるお湯・コンデンサー}\\お湯と部屋の温度差や、コンデンサーにたまった電気は、$e^{-t/\tau}$ の形でへります。時定数 $\tau$ が、計器や制御の速さを表します（6-4）。
\item \textbf{ラプラス変換}\\制御工学で使うラプラス変換は、$e^{-st}$ をかけて積分する計算です。$e$ の微分・積分が、その土台になります（9-5）。
\end{itemize}
\end{butsuri}


%% ---------------- 章末問題 ----------------
\bkendproblems
\bkproblemgroup{基本}{3}

\bkproblem 「1年あたり100\,\%」の利率で1万円を預け、利息を2か月ごと（1年に6回）に計算してもらいます。1回分の利率は何\,\%ですか。1年後のお金は約何万円ですか（小数第2位まで）。
\bkpans{$\dfrac{100}{6}\approx16.7$\,\%。$\Bigl(1+\dfrac{1}{6}\Bigr)^6\approx2.52$ 万円}

\bkproblem 次の関数を微分しましょう。\quad（1）$e^{6x}$\quad（2）$4e^{-2x}$\quad（3）$e^{x^3}$
\bkpans{（1）\nobreak $6e^{6x}$\quad（2）\nobreak $-8e^{-2x}$\quad（3）\nobreak $3x^2e^{x^3}$}

\bkproblem 次の関数を微分しましょう。\quad（1）$\ln(3x)$\quad（2）$\ln(4x-1)$\quad（3）$3^x$
\bkpans{（1）\nobreak $\dfrac{3}{3x}=\dfrac{1}{x}$\quad（2）\nobreak $\dfrac{4}{4x-1}$\quad（3）\nobreak $3^x\ln3$}

\bkproblem 次の定積分を求めましょう。\quad（1）$\displaystyle\int_0^2e^{-x}\,\dd x$\quad（2）$\displaystyle\int_2^4\frac{1}{x}\,\dd x$
\bkpans{（1）\nobreak $1-e^{-2}\approx0.865$\quad（2）\nobreak $\ln4-\ln2=\ln2\approx0.693$}

\bkproblem $N=500e^{-0.2t}$ について、$\dfrac{\dd N}{\dd t}$ と半減期を求めましょう。
\bkpans{$-100e^{-0.2t}$。半減期 $\dfrac{\ln2}{0.2}\approx3.5$}

\bkproblem コンデンサーの電圧が $V=12\,(1-e^{-t/4})$（V）でふえていきます。4秒後の電圧と、$t=0$ での電圧のふえる速さを求めましょう（$e^{-1}=0.368$）。
\bkpans{$12\times0.632\approx\qty{7.6}{V}$。$\dfrac{\dd V}{\dd t}=3e^{-t/4}$、$t=0$ で毎秒\qty{3}{V}}

\bkproblem 電気の量が $Q=Q_0e^{-t/4}$（$t$ は秒）でへっていきます。時定数と、半分になるまでの時間を求めましょう。
\bkpans{時定数4秒、\mbox{$4\ln2\approx2.8$\,秒}}

\bkproblem 「1年あたり100\,\%」の利率で1万円を預け、利息を毎月（1年に12回）計算してもらいます。1回分の利率は1年の何分の1ですか。1年後のお金は約何万円ですか（電卓で小数第2位まで）。
\bkpans{$\dfrac{1}{12}$。$\Bigl(1+\dfrac{1}{12}\Bigr)^{12}\approx2.61$ 万円}

\bkproblem 次の関数を微分しましょう。\quad（1）$5e^{2x}$\quad（2）$e^{-x^2}$\quad（3）$e^{3x+2}$
\bkpans{（1）\nobreak $10e^{2x}$\quad（2）\nobreak $-2xe^{-x^2}$\quad（3）\nobreak $3e^{3x+2}$}

\bkproblem 100万円が $r=0.03$（1年あたり）でたえずふえるとき、10年後の金額を求めましょう（$e^{0.3}=1.3499$）。
\bkpans{$100\times e^{0.3}=134.99$。約135万円}

\bkproblemgroup{チャレンジ}{4}

\bkproblem $y=e^x$ のグラフの、$x=0$ の点での接線の式を求めましょう。これを使って、$e^{0.1}$ の近似値を求め、電卓の値 1.10517… とくらべましょう。
\bkpans{$y=x+1$。$e^{0.1}\approx1.1$（ずれ約0.005）}

\bkproblem 細菌の数が、$N=100e^{0.7t}$（$t$ は時間）でふえていきます。数が2倍になるまでの時間（倍加時間）と、$t=0$ でのふえる速さを求めましょう。
\bkpans{$e^{0.7T}=2$ より \mbox{$T=\dfrac{\ln2}{0.7}\approx1.0$\,時間}。$\dfrac{\dd N}{\dd t}=70e^{0.7t}$、$t=0$ で1時間あたり70}

\bkproblem 人口が $r=0.07$（1年あたり）でたえずふえる町があります。人口が2倍になるまでの年数を求めましょう（$\ln2=0.693$）。
\bkpans{$e^{0.07T}=2$ より \mbox{$T=\dfrac{0.693}{0.07}=9.9$\,年}}

\bkproblem $\displaystyle\int_0^{\ln2}e^x\,\dd x$ を求めましょう（$e^{\ln2}=2$ を使います）。
\bkpans{$\Bigl[\,e^x\,\Bigr]_0^{\ln2}=e^{\ln2}-e^0=2-1=1$}

\begin{furikaeri}
\begin{checklist}
\item $e$ が $(1+\frac{1}{n})^n$ の近づく先であることを確かめられる。\hfill{\footnotesize\sffamily\color{bkTheme}→\pageref{bk:5-14:n}ページ}
\item $(e^{ax})'=ae^{ax}$ を使える。\hfill{\footnotesize\sffamily\color{bkTheme}→\pageref{bk:5-14:eax}ページ}
\item $\ln x$ を微分・積分できる。\hfill{\footnotesize\sffamily\color{bkTheme}→\pageref{bk:5-14:lnbibun}ページ}
\end{checklist}
\end{furikaeri}

\bkfukidashi{sora}{ふえる・へるは $e$ で書けた。じゃあ、ゆれる・回るは？　次は $\sin$ と $\cos$ だ！}

%% 5-15 sin・cosを微分すると？（第11回：ページ制限なし。説明・図・例題を書き足した）
\newcommand\bkvomidashi[1]{\par\Needspace{6\baselineskip}\noindent\textbf{#1}\par}
\bkchapter{5}{15}{$\sin$・$\cos$を微分すると？}{\bklevel{4}}{sin・cosを微分すると？}
\bkmichishirube{4-2,4-8,5-6}{5-17,6-6,6-10}

\begin{goal}
\begin{itemize}[label={\color{bkTheme}●}, leftmargin=1.6\zw, itemsep=0pt, topsep=1pt]
\item \bklevel{3}\enspace $(\sin t)'=\cos t$、$(\cos t)'=-\sin t$ を、円運動で説明できる。
\item \bklevel{4}\enspace $\sin\omega t$、$\cos\omega t$ を微分・積分できる。
\item \bklevel{4}\enspace 単振動の速度・加速度を求め、$\ddot{x}=-\omega^2x$ を確かめられる。
\item \bklevel{4}\enspace $\sin^2$ の平均から、交流の実効値を求められる。\bkidx{こうりゅう}{交流}
\end{itemize}
\end{goal}

\begin{hatena}
\noindent\begin{minipage}[c]{0.58\linewidth}
観覧車のゴンドラは、一定の速さで円をえがいて回ります。でも、地面から見た高さの変わり方は一定ではありません。いちばん下や上ではほとんど変わらず、横を通るときに、いちばん速く上がったり下がったりします。

高さは $\sin$ や $\cos$ で表せました（4-8）。その「変わる速さ」は、どんな式になるのでしょう。
\end{minipage}\hfill
\begin{minipage}[c]{0.39\linewidth}
\centering
\begin{tikzpicture}[x=0.78mm, y=0.78mm]
  \draw[bkGray, line width=1pt] (-20,-22) -- (20,-22);
  \draw[bkInk!60, line width=0.8pt] (0,0) circle (17);
  \draw[bkInk!60, line width=0.8pt] (0,0) -- (-8,-22) (0,0) -- (8,-22);
  \foreach \a in {0,45,...,315} {\draw[bkInk!40, line width=0.5pt] (0,0) -- (\a:17); \fill[bkRoute2] (\a:17) ++(-1.8,-2.6) rectangle ++(3.6,3);}
  \fill[bkOrange] (0:17) ++(-1.8,-2.6) rectangle ++(3.6,3);
  \draw[bk arrow, bkOrange, line width=1pt] (21,-3) -- (21,7);
  \node[font=\sffamily\scriptsize, anchor=west, align=left] at (22,2) {ここで\\いちばん\\速く上がる};
\end{tikzpicture}
\end{minipage}
\end{hatena}

$\sin$ と $\cos$ は、回るもの・ゆれるもの・波を表す関数です（4-8）。この章では、これらを微分すると、また $\sin$ と $\cos$ が出てくることを、円の上を回る点の動きから確かめます。

\clearpage
%% ---------------- やってみよう・思い出そう ----------------
\begin{yattemiyou}[表計算で $\sin x$ の傾きを調べよう]
$h=0.001$ として、$\dfrac{\sin(x+h)-\sin x}{h}$ を計算します（5-3の「平均の変化率で $h$ を小さくする」と同じ）。$x$ はラジアン（4-2）です。表計算の \texttt{SIN} は、ラジアンで計算します。\bkidx{らじあん}{ラジアン}
\begin{bkcenter}
\bksheet{4}{22}{%
1 & $x$ & $\sin x$ & 傾き（$h=0.001$） & $\cos x$ \\
2 & 0 & 0 & 1.0000 & 1 \\
3 & $\frac{\pi}{6}=0.5236$ & 0.5 & 0.8658 & 0.8660 \\
4 & $\frac{\pi}{4}=0.7854$ & 0.7071 & 0.7068 & 0.7071 \\
5 & $\frac{\pi}{3}=1.0472$ & 0.8660 & 0.4996 & 0.5 \\
6 & $\frac{\pi}{2}=1.5708$ & 1 & $-0.0005$ & 0 \\
7 & $\pi=3.1416$ & 0 & $-1.0000$ & $-1$ \\}
\bktabcaption{傾きの列と、$\cos x$ の列がほとんど同じになる。}
\end{bkcenter}
\end{yattemiyou}

\noindent $\sin x$ の傾きは、$\cos x$ にとても近い値になりました。小さなずれ（0.8658 と 0.8660 など）は、$h$ が0でない分です。$\sin x$ を微分すると $\cos x$ になりそうです。この章では、そのわけを円運動で考えます。

\begin{omoidasou}
\bkrecall{4-2}{弧度法：角 $\theta$（rad）＝弧の長さ÷半径。半径1の円なら、角と弧の長さは同じ数。1周は $2\pi$\,rad（360°）。円の接線は、半径に垂直。}
\bkrecall{4-8}{半径1の円の上で、$x$ 軸から角 $\theta$ の所にある点の座標が $(\cos\theta,\ \sin\theta)$。\\半径 $r$ の円の上を、1秒に $\omega$ ラジアン（角速度 $\omega$）で回る点は\\$(r\cos\omega t,\ r\sin\omega t)$。周期 $T=\dfrac{2\pi}{\omega}$}
\bkrecall{4-11}{速度は、向きと大きさをもつ量（ベクトル）。$x$ 方向の成分と $y$ 方向の成分に分けて考えられる。}
\bkrecall{5-5}{速度は位置を時間で微分したもの、加速度は速度を微分したもの。力＝質量×加速度。}
\bkrecall{5-6}{合成関数の微分：外側を微分して、内側の微分をかける。}
\end{omoidasou}

\clearpage
%% ---------------- 数学の道具 ----------------
\bksection{doguu}{数学の道具}

\bksubsection{1}{円の上を回る点の速さと向き}

\label{bk:5-15:en}半径1の円の上を、点Pが「1秒に1ラジアン」の割合で回ります。0秒のとき、Pは $x$ 軸の上（角0）にいます。$t$ 秒たつと、Pは角 $t$ の所にいるので、Pの座標は $(\cos t,\ \sin t)$ です（4-8）。

\noindent\begin{minipage}[c]{0.46\linewidth}
\centering
\begin{tikzpicture}[x=1mm, y=1mm]
  \draw[bkGray, line width=0.6pt, -{Stealth[length=4pt]}] (-26,0) -- (28,0) node[anchor=north, font=\sffamily\scriptsize] {$x$};
  \draw[bkGray, line width=0.6pt, -{Stealth[length=4pt]}] (0,-26) -- (0,34) node[anchor=east, font=\sffamily\scriptsize] {$y$};
  \draw[bkBlue, line width=1.2pt] (0,0) circle (21);
  \def\ang{35}
  \draw[bkInk, line width=0.8pt] (0,0) -- (\ang:21);
  \fill[bkInk] (\ang:21) circle (1.1);
  \draw[bk arrow, bkOrange, line width=1.5pt] (\ang:21) -- ++(\ang+90:15);
  \draw[bkInk, line width=0.5pt] (0,0) ++(0:6) arc (0:\ang:6);
  \node[font=\sffamily\scriptsize] at ({\ang/2}:8.5) {$t$};
  \node[font=\sffamily\scriptsize, anchor=west] at ($(\ang:21)+(1,-1.5)$) {P};
  \node[font=\sffamily\scriptsize, text=bkOrange!80!black, anchor=south west, fill=white, inner sep=1pt] at ($(\ang:21)+(\ang+90:15)+(1,-1)$) {速度（長さ1）};
  \draw[bkInk, line width=0.4pt] (\ang:18) -- ++(\ang+90:3) -- ++(\ang:3);
\end{tikzpicture}
\end{minipage}\hfill
\begin{minipage}[c]{0.5\linewidth}
\begin{itemize}[leftmargin=1.2\zw, itemsep=3pt, topsep=0pt, label={\color{bkTheme}●}]
\item \textbf{速さ}\quad 弧の長さ＝半径×角（4-2）です。1秒で角は1ふえるので、1秒に進む弧の長さは $1\times1=1$。だから\textbf{速さは1}です。
\item \textbf{向き}\quad Pは円にそって動くので、向きは円の接線の向きです。接線は半径に垂直です（4-2）。だから速度の矢印は、\textbf{半径の向きを左回りに90°回した向き}です（Pは左回りに回っているから）。
\end{itemize}
\end{minipage}\par\vspace{4pt}
この「長さ1で、半径に垂直な矢印」を、$x$ 方向（横）と $y$ 方向（たて）の2つに分けてみましょう（4-11）。

\bksubsection{2}{速度の矢印を、たてと横に分ける}

\label{bk:5-15:bunkai}下の図の左は、半径（長さ1）を、横 $\cos t$ とたて $\sin t$ の直角三角形の斜めの辺と見たものです。速度の矢印は、この半径を90°回したものです。だから、\textbf{この直角三角形をまるごと90°回した三角形}が、速度の矢印の下にできます。

\begin{bkfigure}
\begin{tikzpicture}[x=1mm, y=1mm]
  \def\R{30}
  \draw[bkGray, line width=0.6pt, -{Stealth[length=4pt]}] (-6,0) -- (40,0) node[anchor=west, font=\sffamily\scriptsize] {$x$};
  \draw[bkGray, line width=0.6pt, -{Stealth[length=4pt]}] (0,-4) -- (0,46) node[anchor=east, font=\sffamily\scriptsize] {$y$};
  \draw[bkBlue!50, line width=0.8pt] (-4:\R) arc (-4:48:\R);
  % 半径の三角形（横 cos t は青、たて sin t はオレンジ）
  \fill[bkBlue!8] (0,0) -- (25.98,0) -- (25.98,15) -- cycle;
  \draw[bkBlue, line width=1.6pt] (0,0) -- (25.98,0);
  \draw[bkOrange, line width=1.6pt] (25.98,0) -- (25.98,15);
  \draw[bkInk, line width=0.9pt] (0,0) -- (25.98,15);
  \fill[bkInk] (25.98,15) circle (1);
  \node[font=\sffamily\scriptsize, anchor=north] at (13,-1.5) {$\cos t$};
  \node[font=\sffamily\scriptsize, anchor=east] at (24.4,4.6) {$\sin t$};
  \draw[bkInk, line width=0.5pt] (7,0) arc (0:30:7);
  \node[font=\sffamily\scriptsize] at (15:9.8) {$t$};
  \node[font=\sffamily\scriptsize, anchor=north west] at (26.8,15) {P};
  % 速度の三角形（90°回した。横 sin t はオレンジ、たて cos t は青）
  \fill[bkOrange!8] (25.98,15) -- (10.98,15) -- (10.98,40.98) -- cycle;
  \draw[bkOrange, line width=1.6pt] (25.98,15) -- (10.98,15);
  \draw[bkBlue, line width=1.6pt] (10.98,15) -- (10.98,40.98);
  \draw[bk arrow, bkInk, line width=1.1pt] (25.98,15) -- (10.98,40.98);
  \node[font=\sffamily\scriptsize, anchor=south] at (17.5,15.6) {$\sin t$};
  \node[font=\sffamily\scriptsize, anchor=east] at (9.4,28) {$\cos t$};
  \node[font=\sffamily\scriptsize, anchor=west, fill=white, inner sep=1pt] at (20.5,30) {速度（長さ1）};
  % 右の説明
  \node[font=\sffamily\small, anchor=north west, align=left, text width=66mm] at (52,44) {%
  \textcolor{bkBlue!80!black}{青い辺}の長さは、どちらも $\cos t$。\\
  \textcolor{bkOrange!80!black}{オレンジの辺}の長さは、どちらも $\sin t$。\\[4pt]
  速度の矢印は、\\
  \quad 横に $\sin t$ だけ\textbf{左}（$x$ がへる向き）\\
  \quad たてに $\cos t$ だけ\textbf{上}（$y$ がふえる向き）\\[4pt]
  だから、速度の成分は\\
  \quad $(x\text{ 方向},\ y\text{ 方向})=(-\sin t,\ \cos t)$};
\end{tikzpicture}
\bkfigcaption{半径の三角形（青い面）を90°回すと、速度の三角形（オレンジの面）になる。辺の長さはそのまま、たてと横が入れかわる。}
\end{bkfigure}

\bkvomidashi{数で確かめる}
$t=\dfrac{\pi}{6}$（30°）のとき、$\cos\dfrac{\pi}{6}=0.866$、$\sin\dfrac{\pi}{6}=0.5$ なので、Pは $(0.866,\ 0.5)$ にいます。速度の成分は $(-0.5,\ 0.866)$ のはずです。

ほんとうにそうか、0.001秒だけ時間を進めて調べます。
\begin{bkcenter}
\begin{bktable}{colspec={Q[l,m,wd=34mm] Q[c,m,wd=30mm] Q[c,m,wd=30mm]}, rowsep=1.5pt}
 & $x=\cos t$ & $y=\sin t$ \\
$t=\frac{\pi}{6}$ & 0.866025 & 0.500000 \\
$t=\frac{\pi}{6}+0.001$ & 0.865525 & 0.500866 \\
変化÷0.001 & $-0.500$ & $0.866$ \\
\end{bktable}
\bktabcaption{0.001秒のあいだの変化を0.001でわると、1秒あたりの変化（速度）になる。}
\end{bkcenter}
$x$ は1秒あたり0.5の割合でへり、$y$ は1秒あたり0.866の割合でふえています。図から考えた $(-\sin t,\ \cos t)=(-0.5,\ 0.866)$ とぴったり合います。

\bksubsection{3}{座標を微分すると、速度の成分}

\label{bk:5-15:bibun}5-5で、まっすぐ動く物の速度は「位置を時間で微分したもの」でした。平面の上を動く点でも、横の動きとたての動きを別々に見れば、同じことがいえます。
\begin{itemize}[leftmargin=1.2\zw, itemsep=2pt, topsep=2pt, label={\color{bkTheme}●}]
\item Pの $x$ 座標（横の位置）がふえる速さ＝速度の $x$ 成分
\item Pの $y$ 座標（たての位置、高さ）がふえる速さ＝速度の $y$ 成分
\end{itemize}
Pの $x$ 座標は $\cos t$、$y$ 座標は $\sin t$ でした。そして、速度の成分は $(-\sin t,\ \cos t)$ でした。2つを組にして読むと、
\begin{bkcenter}
\begin{bktable}{colspec={Q[l,m,wd=30mm] Q[c,m,wd=30mm] Q[c,m,wd=50mm]}, rowsep=2pt}
 & 位置（座標） & 位置を微分したもの（速度） \\
$x$ 方向（横） & $\cos t$ & $-\sin t$ \\
$y$ 方向（たて） & $\sin t$ & $\cos t$ \\
\end{bktable}
\end{bkcenter}
つまり、
\[ (\sin t)'=\cos t,\qquad (\cos t)'=-\sin t \]
$\cos t$ の微分に $-$ がつくのは、Pが左回りに上がっていくとき（$0<t<\frac{\pi}{2}$）、横の位置 $x$ がへっていくからです。\bkidx{さいん・こさいんのびぶん}{$\sin$・$\cos$ の微分}

\bkvomidashi{\bkhatenaref の答え}
観覧車のゴンドラの高さが $\sin t$ で表されるとき、高さの変わる速さは $\cos t$ です。$\cos t$ の大きさがいちばん大きい（1 か $-1$）のは、$t=0$ や $t=\pi$、つまりゴンドラが中心と同じ高さ（真横）を通るときです。このとき、速度の矢印はまっすぐ上か下を向いていて、速さ1がそのまま高さの変化になります。いちばん上（$t=\frac{\pi}{2}$）やいちばん下では $\cos t=0$ で、速度の矢印は真横を向いているので、高さは一瞬変わりません。

\noindent\begin{minipage}[c]{0.57\linewidth}
\textbf{いちばん下からはかった高さ}\par
観覧車のゴンドラの高さを、いちばん下からはかってみます。半径\qty{50}{m}の観覧車の中心は、いちばん下から\qty{50}{m}の高さです。いちばん下から角 $\theta$ だけ回ったゴンドラは、中心より $50\cos\theta$ だけ下にいます（右の図の直角三角形）。だから高さは
\[ y=50-50\cos\theta \]
$\theta=0$（いちばん下）で $y=0$、$\theta=\pi$（いちばん上）で $y=50+50=100$。合っています。
\end{minipage}\hfill
\begin{minipage}[c]{0.4\linewidth}
\centering
\begin{tikzpicture}[x=1mm, y=1mm]
  \draw[bkGray, line width=0.8pt] (-26,0) -- (24,0);
  \draw[bkInk!50, line width=0.8pt] (0,20) circle (20);
  \draw[bkGray, dashed, line width=0.5pt] (0,20) -- (0,0);
  \def\th{50}
  \draw[bkInk, line width=0.8pt] (0,20) -- ({20*sin(\th)},{20-20*cos(\th)});
  \draw[bkOrange, line width=1.4pt] (0,20) -- (0,{20-20*cos(\th)});
  \draw[bkGray, line width=0.5pt] (0,{20-20*cos(\th)}) -- ({20*sin(\th)},{20-20*cos(\th)});
  \fill[bkRoute2] ({20*sin(\th)},{20-20*cos(\th)}) ++(-1.8,-2.6) rectangle ++(3.6,3);
  \draw[bkInk, line width=0.5pt] (0,14) arc (-90:{-90+\th}:6);
  \node[font=\sffamily\scriptsize] at ($(0,20)+({-90+\th/2}:8.8)$) {$\theta$};
  \node[font=\sffamily\scriptsize, text=bkOrange!80!black, anchor=east] at (-1.2,13.5) {$50\cos\theta$};
  \draw[bkGray, line width=0.5pt, <->] (-22,0) -- (-22,20);
  \node[font=\sffamily\scriptsize, anchor=east] at (-22.8,10) {50};
  \fill[bkInk] (0,20) circle (0.8);
\end{tikzpicture}
\end{minipage}
例題6では、この式を使って、上がる速さを計算します。


\bksubsection{4}{微分すると、波が $\frac{1}{4}$ 周期ずれる}

\label{bk:5-15:zure}$\sin t$ と $\cos t$ のグラフを重ねてみます。

\begin{bkfigure}
\begin{tikzpicture}[tl/.style={font=\sffamily\footnotesize, text=bkInk, anchor=base, yshift=-6mm}]
\begin{axis}[bkplot, width=120mm, height=46mm, xmin=0, xmax=7, ymin=-1.15, ymax=1.35, xtick={1.5708,3.1416,4.7124,6.2832}, xticklabels={}, ytick={-1,1}, xlabel={$t$}, ylabel={$y$}]
  \addplot[bkcurve, domain=0:6.6, samples=120] {sin(deg(x))};
  \addplot[bkcurve2, domain=0:6.6, samples=120] {cos(deg(x))};
  % 目もりの数字は、波と重ならないように −1 の線より下に置く。分数も π・2π と同じ字の大きさにし、ベースラインをそろえる
  \node[tl] at (axis cs:1.5708,-1.15) {$\dfrac{\pi}{2}$};
  \node[tl] at (axis cs:3.1416,-1.15) {$\pi$};
  \node[tl] at (axis cs:4.7124,-1.15) {$\dfrac{3\pi}{2}$};
  \node[tl] at (axis cs:6.2832,-1.15) {$2\pi$};
  \draw[bk arrow, bkGray, line width=0.8pt] (axis cs:0,1.15) -- (axis cs:1.5708,1.15);
  \node[font=\sffamily\scriptsize, anchor=south] at (axis cs:0.785,1.15) {$\frac{1}{4}$ 周期};
  \node[font=\sffamily\scriptsize, text=bkBlue, anchor=west] at (axis cs:6.75,0.31) {$\sin t$};
  \node[font=\sffamily\scriptsize, text=bkOrange!80!black, anchor=west] at (axis cs:6.75,0.95) {$\cos t$};
\end{axis}
\end{tikzpicture}
\bkfigcaption{$\sin t$（青）の傾きのグラフが $\cos t$（オレンジ）。$\cos t$ の山は $t=0$、$\sin t$ の山は $t=\frac{\pi}{2}$。$\cos t$ のほうが $\frac{\pi}{2}$ だけ早く山が来る。}
\end{bkfigure}
\begin{itemize}[leftmargin=1.2\zw, itemsep=2pt, topsep=2pt, label={\color{bkTheme}●}]
\item $\sin t$ が0を上向きに通る所（$t=0$）では、青い線がいちばん急に上がっていて、傾きは最大の1。そこで $\cos t$ はちょうど山（1）です。
\item $\sin t$ が山（$t=\frac{\pi}{2}$）では、青い線は一瞬平らで、傾きは0。そこで $\cos t=0$ です。
\item $\sin t$ が0を下向きに通る所（$t=\pi$）では、傾きは最も急な下り（$-1$）。そこで $\cos t$ は谷（$-1$）です。
\end{itemize}
$\cos t$ は、$\sin t$ と同じ形の波を、$\frac{\pi}{2}$ だけ左にずらしたものです。式で書くと $\cos t=\sin\bigl(t+\frac{\pi}{2}\bigr)$。数で確かめましょう。
\begin{bkcenter}
\begin{bktable}{colspec={Q[l,m,wd=30mm] *{4}{Q[c,m,wd=18mm]}}, rowsep=1.5pt}
$t$ & 0 & $\dfrac{\pi}{6}$ & $\dfrac{\pi}{3}$ & $\dfrac{\pi}{2}$ \\
$\cos t$ & 1 & 0.866 & 0.5 & 0 \\
$t+\dfrac{\pi}{2}$ & $\dfrac{\pi}{2}$ & $\dfrac{2\pi}{3}$ & $\dfrac{5\pi}{6}$ & $\pi$ \\
$\sin\Bigl(t+\dfrac{\pi}{2}\Bigr)$ & 1 & 0.866 & 0.5 & 0 \\
\end{bktable}
\end{bkcenter}
1周期は $2\pi$ なので、$\frac{\pi}{2}$ は $2\pi$ の $\frac{1}{4}$ です。\textbf{$\sin t$ を微分すると、同じ形の波が $\frac{1}{4}$ 周期だけ先に進みます}。

微分するたびに $\frac{1}{4}$ 周期ずつ進むので、4回微分すると1周期進んで、もとの $\sin t$ にもどります。

\begin{bkfigure}
\begin{tikzpicture}[x=1mm, y=1mm, nd/.style={draw=bkTheme, line width=0.9pt, rounded corners=3pt, fill=bkTheme!8, minimum width=20mm, minimum height=8mm, font=\small}]
  \node[nd] (a) at (0,13) {$\sin t$};
  \node[nd] (b) at (34,0) {$\cos t$};
  \node[nd] (c) at (0,-13) {$-\sin t$};
  \node[nd] (d) at (-34,0) {$-\cos t$};
  \draw[bk arrow, bkOrange, line width=1.1pt] (a.east) to[bend left=20] node[pos=0.5, anchor=south west, font=\sffamily\scriptsize] {微分} (b.north);
  \draw[bk arrow, bkOrange, line width=1.1pt] (b.south) to[bend left=20] node[pos=0.5, anchor=north west, font=\sffamily\scriptsize] {微分} (c.east);
  \draw[bk arrow, bkOrange, line width=1.1pt] (c.west) to[bend left=20] node[pos=0.5, anchor=north east, font=\sffamily\scriptsize] {微分} (d.south);
  \draw[bk arrow, bkOrange, line width=1.1pt] (d.north) to[bend left=20] node[pos=0.5, anchor=south east, font=\sffamily\scriptsize] {微分} (a.west);
  \node[font=\sffamily\scriptsize, text=bkGray, align=center] at (0,0) {右回り：微分\\左回り：積分};
\end{tikzpicture}
\bkfigcaption{微分すると右回りに1つ進む。4回でもとにもどる。積分（$+C$ は省略）は左回り。}
\end{bkfigure}
図の3つめと4つめは、1つめと2つめに $-$ をつけたものです。定数倍（ここでは $-1$ 倍）はそのまま前に出る（5-4）ので、
\begin{align*}
(-\sin t)'&=-(\sin t)'=-\cos t\\
(-\cos t)'&=-(\cos t)'=-(-\sin t)=\sin t
\end{align*}

\bksubsection{5}{角速度 $\omega$ をふくむとき}

\label{bk:5-15:omega}\bkidx{かくそくど}{角速度}1秒に2ラジアン回る点を考えます。$t$ 秒後の角は $2t$ なので、高さは $\sin2t$ です。2倍の速さで回るので、高さも2倍の速さで変わるはずです。

合成関数の微分（5-6）で確かめます。内側を $u=2t$、外側を $\sin u$ とすると、
\[ (\sin2t)'=\underbrace{\cos u}_{\text{外側の微分}}\times\underbrace{(2t)'}_{\text{内側の微分}}=\cos2t\times2=2\cos2t \]
前に出た2は、「1秒に2ラジアン回る」の2です。表計算でも確かめましょう。$t=0.3$ では、
\begin{align*}
\text{傾き（$h=0.001$）}&\quad \frac{\sin(2\times0.301)-\sin(2\times0.3)}{0.001}=1.6495\\
\text{微分の式}&\quad 2\cos(2\times0.3)=2\cos0.6=1.6507
\end{align*}
で、ほとんど同じです。

1秒に $\omega$ ラジアン回るときは、角が $\omega t$ です。同じように内側を $u=\omega t$（微分すると $\omega$）として、
\[ (\sin\omega t)'=\omega\cos\omega t,\qquad (\cos\omega t)'=-\omega\sin\omega t \]
$\omega$（オメガ）は\textbf{角速度}といい、単位は rad/s です。

\bkvomidashi{半径が $r$ のとき}
半径 $r$ の円では、座標はどちらも半径1のときの $r$ 倍です。$r$ は定数なので、微分しても前にそのまま残ります（5-4）。
\[ \text{座標}\ (r\cos\omega t,\ r\sin\omega t)\quad\xrightarrow{\ \text{微分}\ }\quad\text{速度}\ (-r\omega\sin\omega t,\ r\omega\cos\omega t) \]
速さ（速度の矢印の長さ）は $r\omega$ です。1秒に回る角が $\omega$ なので、1秒に進む弧の長さは「半径×角」で $r\omega$（4-2）。微分した結果と合っています。速く回る（$\omega$ が大きい）ほど、大きい円（$r$ が大きい）ほど、速く動きます。

\bksubsection{6}{2回微分すると、もとの $-\omega^2$ 倍}

\label{bk:5-15:nikai}$\sin\omega t$ を2回微分してみます。1回めは⑤のとおり $\omega\cos\omega t$。2回めは、$\omega$ を前に置いたまま $\cos\omega t$ を微分します。
\[ (\sin\omega t)''=(\omega\cos\omega t)'=\omega\times(-\omega\sin\omega t)=-\omega^2\sin\omega t \]
$\omega$ が2回かかって $\omega^2$、$\cos$ の微分で $-$ がついて、$-\omega^2$ になりました。$\cos\omega t$ でも同じで、\textbf{2回微分すると、もとの関数の $-\omega^2$ 倍}になります。ゆれの位置を $x=A\sin\omega t$（$A$ はゆれの大きさ）とすると、$\ddot{x}$（$x$ の2回微分、5-5）は
\[ \ddot{x}=-\omega^2x \]
\textbf{加速度が、位置に比例して、いつも位置と反対向き}です。$x$ が正（中心より右）なら加速度は負（左向き）、$x$ が負なら加速度は正（右向き）。いつも中心のほうへ引きもどされます。

\bkvomidashi{ばねのおもり}
ばねは、のびに比例する力（2-15の $F=kx$）で、のびと反対向きにおもりを引きもどします。向きまでこめて書くと $F=-kx$。力＝質量×加速度（5-5）なので、
\[ m\ddot{x}=-kx\qquad\text{両辺を }m\text{ でわると}\qquad \ddot{x}=-\frac{k}{m}\,x \]
これは $\omega^2=\dfrac{k}{m}$ とした $\ddot{x}=-\omega^2x$ と同じ形です。だから、ばねのおもりは $\sin$ の形でゆれます。このゆれを\textbf{単振動}といいます。たとえば質量\qty{0.5}{kg}のおもりと、ばね定数\qty{8}{N/m}のばねなら、$\omega^2=8\div0.5=16$ で $\omega=4$、周期は $\dfrac{2\pi}{4}\approx\qty{1.57}{s}$ です。くわしくは6-6で調べます。\bkidx{たんしんどう}{単振動}

\begin{bkfigure}
\begin{tikzpicture}[tl/.style={font=\sffamily\footnotesize, text=bkInk, anchor=base, yshift=-6mm}]
\begin{axis}[bkplot, width=120mm, height=46mm, xmin=0, xmax=6.6, ymin=-1.15, ymax=1.5, xtick={1.5708,3.1416,4.7124,6.2832}, xticklabels={}, ytick={-1,1}, xlabel={$t$}]
  \addplot[bkcurve, domain=0:6.4, samples=120] {sin(deg(x))};
  \addplot[bkcurve2, domain=0:6.4, samples=120] {cos(deg(x))};
  \addplot[bkcurve3, domain=0:6.4, samples=120] {-sin(deg(x))};
  \node[font=\sffamily\scriptsize, text=bkBlue, anchor=south west, inner sep=1pt] at (axis cs:1.7,1.12) {位置 $x$};
  \node[font=\sffamily\scriptsize, text=bkOrange!80!black, anchor=south west, inner sep=1pt] at (axis cs:0.12,1.12) {速度 $v$};
  \node[font=\sffamily\scriptsize, text=bkGreen!70!black, anchor=south west, inner sep=1pt] at (axis cs:4.84,1.12) {加速度 $a$};
  \draw[bkGray, dashed, line width=0.5pt] (axis cs:1.5708,-1.1) -- (axis cs:1.5708,1.0);
  \draw[bkGray, dashed, line width=0.5pt] (axis cs:3.1416,-1.1) -- (axis cs:3.1416,1.0);
  % 目もりの数字と「はし」「中心」は、波と重ならないように −1 の線より下に置く
  \node[tl] at (axis cs:1.5708,-1.15) {$\dfrac{\pi}{2}$};
  \node[tl] at (axis cs:3.1416,-1.15) {$\pi$};
  \node[tl] at (axis cs:4.7124,-1.15) {$\dfrac{3\pi}{2}$};
  \node[tl] at (axis cs:6.2832,-1.15) {$2\pi$};
  \node[tl, font=\sffamily\scriptsize, yshift=-6.5mm] at (axis cs:1.5708,-1.15) {はし};
  \node[tl, font=\sffamily\scriptsize, yshift=-6.5mm] at (axis cs:3.1416,-1.15) {中心};
\end{axis}
\end{tikzpicture}
\bkfigcaption{$x=\sin t$ のとき、$v=\cos t$、$a=-\sin t$。はし（$x$ が最大）では速度0で加速度が最大、中心では速度が最大で加速度0。}
\end{bkfigure}
図の緑の線（加速度）は、青い線（位置）を上下にひっくり返した形です。これが「$-1$ 倍」（ここでは $\omega=1$ なので $-\omega^2=-1$）の意味です。

\bksubsection{7}{$\sin$・$\cos$ の積分}

\label{bk:5-15:sekibun}積分は微分の逆でした（5-11）。微分の式を、右から左へ逆に読みます。
\begin{bkcenter}
\begin{bktable}{colspec={Q[c,m,wd=40mm] Q[c,m,wd=58mm]}, rowsep=2pt}
微分の式 & 逆に読んだ積分の式 \\
$(-\cos t)'=\sin t$ & $\displaystyle\int\sin t\,\dd t=-\cos t+C$ \\
$(\sin t)'=\cos t$ & $\displaystyle\int\cos t\,\dd t=\sin t+C$ \\
\end{bktable}
\end{bkcenter}
$\sin t$ の積分に $-$ がつくのは、$\cos t$ を微分すると $-\sin t$ になるからです。$-\cos t$ を微分すれば、$-$ が2つで $+\sin t$ になります。

\bkvomidashi{角速度 $\omega$ があるとき}
$\cos\omega t$ を微分すると $-\omega\sin\omega t$ で、ほしい $\sin\omega t$ の $-\omega$ 倍になってしまいます。そこで、はじめに $-\dfrac{1}{\omega}$ をかけておくと、微分したとき $\omega$ と打ち消しあいます。
\[ \Bigl(-\frac{1}{\omega}\cos\omega t\Bigr)'=-\frac{1}{\omega}\times(-\omega\sin\omega t)=\sin\omega t \]
だから、
\[ \int\sin\omega t\,\dd t=-\frac{1}{\omega}\cos\omega t+C,\qquad \int\cos\omega t\,\dd t=\frac{1}{\omega}\sin\omega t+C \]
\textbf{微分では $\omega$ をかけ、積分では $\omega$ でわる}。微分と積分は逆のはたらきなので、$\omega$ のあつかいも逆になります。

\bkvomidashi{例}
$\sin t$ の山1つ分（0から $\pi$ まで）の面積を求めます。$-\cos t$ の、上端 $\pi$ での値から下端0での値をひきます（5-11）。$\cos\pi=-1$、$\cos0=1$ なので、
\begin{align*}
\int_0^\pi\sin t\,\dd t&=\Bigl[\,-\cos t\,\Bigr]_0^\pi=(-\cos\pi)-(-\cos0)\\
&=\{-(-1)\}-(-1)=1+1=2
\end{align*}
5-10のように、細い長方形をたしても確かめられます。0から $\pi$ を $n$ 等分し、まん中の高さで長方形をつくってたすと、$n=6$ で2.023、$n=12$ で2.006、$n=100$ で2.0001。2に近づきます。

\begin{bkfigure}
\begin{tikzpicture}
\begin{axis}[bkplot, name=L, width=64mm, height=36mm, xmin=0, xmax=3.6, ymin=0, ymax=1.3, xtick={1.5708,3.1416}, xticklabels={$\dfrac{\pi}{2}$,$\pi$}, ytick={1}, xlabel={$t$}, title={\small $\sin t$ の山1つ分}, title style={yshift=-1mm}]
  \addplot[fill=bkBlue!15, draw=none, domain=0:3.1416, samples=80] {sin(deg(x))} \closedcycle;
  \addplot[bkcurve, domain=0:3.1416, samples=80] {sin(deg(x))};
  \node[font=\sffamily\scriptsize, text=bkBlue!80!black] at (axis cs:1.5708,0.4) {面積2};
\end{axis}
\begin{axis}[bkplot, at={(L.east)}, anchor=west, xshift=10mm, width=64mm, height=36mm, xmin=0, xmax=3.6, ymin=0, ymax=1.3, xtick={1.5708,3.1416}, xticklabels={$\dfrac{\pi}{2}$,$\pi$}, ytick={1}, xlabel={$t$}, title={\small $\sin2t$ の山1つ分}, title style={yshift=-1mm}]
  \addplot[fill=bkOrange!25, draw=none, domain=0:1.5708, samples=60] {sin(deg(2*x))} \closedcycle;
  \addplot[bkcurve2, domain=0:1.5708, samples=60] {sin(deg(2*x))};
  \node[font=\sffamily\scriptsize, text=bkOrange!80!black] at (axis cs:0.785,0.4) {面積1};
\end{axis}
\end{tikzpicture}
\bkfigcaption{$\sin2t$（オレンジ）の山は、$\sin t$（青）の山と高さは同じで、はばが半分。だから面積も半分の1。積分で $\frac{1}{\omega}=\frac{1}{2}$ がかかるのは、このため。}
\end{bkfigure}
$\sin2t$ の山1つ分は0から $\frac{\pi}{2}$ までです（$2t$ が0から $\pi$ まで動くから）。計算すると、
\[ \int_0^{\frac{\pi}{2}}\sin2t\,\dd t=\Bigl[\,-\frac{1}{2}\cos2t\,\Bigr]_0^{\frac{\pi}{2}}=\Bigl(-\frac{1}{2}\cos\pi\Bigr)-\Bigl(-\frac{1}{2}\cos0\Bigr)=\frac{1}{2}+\frac{1}{2}=1 \]
図の「はばが半分なので面積も半分」と合います。

\bksubsection{8}{なぜラジアンを使うのか}

\label{bk:5-15:rad}角を度（°）ではかると、どうなるでしょう。30°から1°ふやしたときの $\sin$ のふえ方は
\[ \sin31^\circ-\sin30^\circ=0.51504-0.50000=0.01504 \]
で、$\cos30^\circ=0.866$ とは合いません。1°は $\dfrac{\pi}{180}$ ラジアン（$\pi$ ラジアンが180°だから）なので、1°あたりのふえ方は $0.866\times\dfrac{\pi}{180}=0.0151$ になります。度で微分すると、いつも $\dfrac{\pi}{180}$ という余分な数がかかるのです。

ラジアンで角をはかると、半径1の円で「角＝弧の長さ」になり、①で見たように速さがちょうど1になるので、$(\sin t)'=\cos t$ ときれいな形になります。微分や積分では、角はいつもラジアンで表します。電卓や表計算で $\sin$ を計算するときも、角がラジアンになっているか確かめましょう。

\bksubsection{9}{$\sin^2$ の平均：交流の電力}

\label{bk:5-15:jikko}\bkvomidashi{交流の電力は、行ったり来たりする}
家のコンセントの電圧は、$V=V_0\sin\omega t$ のように $+$ と $-$ を行ったり来たりする交流です（$V_0$ は山の高さ）。抵抗 $R$ をつなぐと、電流は $I=\dfrac{V}{R}$（2-15）。電力は電圧×電流なので、
\[ P=VI=V\times\frac{V}{R}=\frac{V^2}{R}=\frac{V_0^2}{R}\sin^2\omega t \]
$\sin^2\omega t$ は $(\sin\omega t)^2$ のことです。2乗なので、電圧が $-$ のときも電力は $+$ で、0から $\dfrac{V_0^2}{R}$ のあいだを行ったり来たりします。では、平均するといくらでしょう。

\bkvomidashi{数で見当をつける}
$\sin^2t$ の値を、1周期（0から $2\pi$）のあいだ、$\frac{\pi}{6}$（30°）ごとに12こ調べます。
\begin{bkcenter}
\begin{bktable}{colspec={Q[l,m,wd=16mm] *{6}{Q[c,m,wd=13mm]}}, rowsep=1.5pt}
角 & 0° & 30° & 60° & 90° & 120° & 150° \\
$\sin t$ & 0 & 0.5 & 0.87 & 1 & 0.87 & 0.5 \\
$\sin^2t$ & 0 & 0.25 & 0.75 & 1 & 0.75 & 0.25 \\
\end{bktable}\par\vspace{4pt}
\begin{bktable}{colspec={Q[l,m,wd=16mm] *{6}{Q[c,m,wd=13mm]}}, rowsep=1.5pt}
角 & 180° & 210° & 240° & 270° & 300° & 330° \\
$\sin t$ & 0 & $-0.5$ & $-0.87$ & $-1$ & $-0.87$ & $-0.5$ \\
$\sin^2t$ & 0 & 0.25 & 0.75 & 1 & 0.75 & 0.25 \\
\end{bktable}
\end{bkcenter}
後半の6こは、$\sin t$ が $-$ になっても、2乗すると前半と同じ値になります。$\sin^2t$ の12この合計は $(0+0.25+0.75+1+0.75+0.25)\times2=6$ なので、平均は $6\div12=0.5$。どうやら平均は $\frac{1}{2}$ になりそうです。0.25と0.75、0と1のように、平均より小さい値と大きい値が、ちょうど組になっているからです。

\noindent\begin{minipage}[c]{0.5\linewidth}
\begin{tikzpicture}
\begin{axis}[bkplot, width=66mm, height=36mm, xmin=0, xmax=6.6, ymin=0, ymax=1.15, xtick={3.1416,6.2832}, xticklabels={$\pi$,$2\pi$}, ytick={0.5,1}, yticklabels={$\dfrac{1}{2}$,1}, xlabel={$t$}]
  \addplot[bkcurve, domain=0:6.2832, samples=120, fill=bkBlue!12] {sin(deg(x))^2} \closedcycle;
  \addplot[bkcurve2, domain=0:6.2832] {0.5};
\end{axis}
\end{tikzpicture}
\end{minipage}\hfill
\begin{minipage}[c]{0.46\linewidth}
グラフで見ると、$\sin^2t$（青）の山の、オレンジの線より上の部分を切りとって、谷にうめると、ちょうど高さ $\frac{1}{2}$ の平らな形になります。平均が $\frac{1}{2}$ というのは、このことです。
\end{minipage}\par\vspace{4pt}
これを積分で確かめます。そのために、$\sin^2$ を、積分のできる形に変えます。

\bkvomidashi{2倍角の式で、$\sin^2$ を $\cos2\theta$ に変える}
4-9の加法定理 $\cos(a+b)=\cos a\cos b-\sin a\sin b$ で、$a$ も $b$ も $\theta$ にすると、
\[ \cos2\theta=\cos\theta\cos\theta-\sin\theta\sin\theta=\cos^2\theta-\sin^2\theta \]
$\sin^2\theta+\cos^2\theta=1$（4-8。半径1の円で三平方の定理）から $\cos^2\theta=1-\sin^2\theta$ なので、これを入れると、
\[ \cos2\theta=(1-\sin^2\theta)-\sin^2\theta=1-2\sin^2\theta \]
これを \textbf{2倍角の式}といいます。$\theta=\frac{\pi}{6}$（30°）で確かめると、左は $\cos60^\circ=0.5$、右は $1-2\times0.5^2=1-0.5=0.5$。合っています。\bkidx{にばいかく}{2倍角}

この式を、$\sin^2\theta$ について解きます（2-14）。$2\sin^2\theta$ を左へ、$\cos2\theta$ を右へ移すと $2\sin^2\theta=1-\cos2\theta$。両辺を2でわって、
\[ \sin^2\theta=\frac{1-\cos2\theta}{2},\qquad\text{同じようにして}\qquad \cos^2\theta=\frac{1+\cos2\theta}{2} \]
（$\cos^2\theta$ のほうは、$\cos2\theta=\cos^2\theta-\sin^2\theta$ の $\sin^2\theta$ に $1-\cos^2\theta$ を入れて作ります。$\cos2\theta=2\cos^2\theta-1$ となるので、$\cos^2\theta$ について解きます。）

\bkvomidashi{1周期で積分する}
$\sin^2t=\dfrac{1-\cos2t}{2}$ を、2つに分けて1つずつ積分します。
\[ \frac{1-\cos2t}{2}=\frac{1}{2}-\frac{1}{2}\cos2t \]
\begin{itemize}[leftmargin=1.2\zw, itemsep=1pt, topsep=2pt, label={\color{bkTheme}●}]
\item $\dfrac{1}{2}$（定数）の積分は $\dfrac{1}{2}t=\dfrac{t}{2}$（5-11）。
\item $\cos2t$ の積分は、⑦で $\omega=2$ として $\dfrac{1}{2}\sin2t$。だから $\dfrac{1}{2}\cos2t$ の積分は、その $\dfrac{1}{2}$ 倍の $\dfrac{1}{4}\sin2t=\dfrac{\sin2t}{4}$。
\end{itemize}
\begin{align*}
\int_0^{2\pi}\sin^2t\,\dd t&=\int_0^{2\pi}\Bigl(\frac{1}{2}-\frac{1}{2}\cos2t\Bigr)\dd t=\Bigl[\,\frac{t}{2}-\frac{\sin2t}{4}\,\Bigr]_0^{2\pi}\\
&=\Bigl(\frac{2\pi}{2}-\frac{\sin4\pi}{4}\Bigr)-\Bigl(0-\frac{\sin0}{4}\Bigr)=(\pi-0)-(0-0)=\pi
\end{align*}
平均は、積分÷区間の長さ（5-12）なので、$\pi\div2\pi=\dfrac{1}{2}$。表で見当をつけた値と合いました。\textbf{$\sin^2$ の1周期の平均は $\frac{1}{2}$}です。$\cos^2$ も同じで、平均は $\frac{1}{2}$ です。

\bkvomidashi{実効値}
だから、交流の平均の電力は $\dfrac{V_0^2}{R}\times\dfrac{1}{2}=\dfrac{V_0^2}{2R}$ です。では、同じ抵抗に直流の電圧 $V_{\mathrm{e}}$ をかけて、同じ平均の電力 $\dfrac{V_{\mathrm{e}}^2}{R}$ を出すには、$V_{\mathrm{e}}$ をいくらにすればよいでしょう。2つの電力を等しいとおくと、
\[ \frac{V_{\mathrm{e}}^2}{R}=\frac{V_0^2}{2R}\qquad\text{両辺に }R\text{ をかけて}\qquad V_{\mathrm{e}}^2=\frac{V_0^2}{2}\qquad\text{つまり}\qquad V_{\mathrm{e}}=\frac{V_0}{\sqrt2} \]
この $V_{\mathrm{e}}$ を、交流の\textbf{実効値}といいます。「はたらき（電力）で見ると、直流の何Vと同じか」を表す値です。\bkidx{じっこうち}{実効値}

家のコンセントの「100\,V」は実効値です。電圧の山の高さは、
\[ V_0=\sqrt2\,V_{\mathrm{e}}=1.414\times100\approx\qty{141}{V} \]
もあります。電気製品は、この山の高さにもたえられるように作られています。

\bksubsection{10}{表計算で、ばねを動かしてみよう}

\label{bk:5-15:sim}5-13では、速度から少しずつ位置を進めて、動きを計算しました。ばねのおもりでも同じことができます。$\omega=1$（$\ddot{x}=-x$）、はじめ $x=1$、$v=0$（はしで手をはなす）とし、$\Delta t=0.1$ 秒ごとに、次の順で計算します。
\begin{enumerate}[label={\bkcircled[bkTheme]{\arabic*}}, leftmargin=2.4\zw, itemsep=1pt, topsep=2pt]
\item 加速度を求める：$a=-x$（⑥の $\ddot{x}=-\omega^2x$ で $\omega=1$）
\item 新しい速度：$v+a\times\Delta t$（加速度は「1秒あたりの速度の変化」だから）
\item 新しい位置：$x+(\text{新しい速度})\times\Delta t$（速度は「1秒あたりの位置の変化」だから）
\end{enumerate}
たとえば1行めは、$a=-1$、新しい速度 $0+(-1)\times0.1=-0.1$、新しい位置 $1+(-0.1)\times0.1=0.99$ です。

\noindent\begin{minipage}[c]{0.44\linewidth}
\centering
\begin{tblr}{colspec={Q[c,m]Q[c,m]Q[c,m]Q[c,m]}, row{1}={bg=bkTheme!12, font=\sffamily\footnotesize}, cells={font=\footnotesize}, hline{1,2,Z}={0.6pt,bkGray}, hline{8}={0.4pt,bkGray,dashed}, colsep=3pt, rowsep=0.5pt}
$t$ & $x$ & $v$ & $\cos t$ \\
0 & 1 & 0 & 1 \\
0.1 & 0.990 & $-0.100$ & 0.995 \\
0.2 & 0.970 & $-0.199$ & 0.980 \\
0.3 & 0.940 & $-0.296$ & 0.955 \\
0.4 & 0.901 & $-0.390$ & 0.921 \\
0.5 & 0.853 & $-0.480$ & 0.878 \\
3.1 & $-1.001$ & $-0.040$ & $-0.999$ \\
6.3 & 0.999 & $-0.019$ & 1.000 \\
\end{tblr}
\end{minipage}\hfill
\begin{minipage}[c]{0.54\linewidth}
\begin{tikzpicture}
\begin{axis}[bkplot, width=70mm, height=40mm, xmin=0, xmax=6.6, ymin=-1.2, ymax=1.2, xtick={3.1416,6.2832}, xticklabels={}, ytick={-1,1}, xlabel={$t$}]
  \addplot[bkcurve, domain=0:6.4, samples=100] {cos(deg(x))};
  % 目もりの数字は、格子の線が通らないように −1 の線より下に置く
  \node[font=\sffamily\footnotesize, text=bkInk, anchor=base, yshift=-4mm] at (axis cs:3.1416,-1.2) {$\pi$};
  \node[font=\sffamily\footnotesize, text=bkInk, anchor=base, yshift=-4mm] at (axis cs:6.2832,-1.2) {$2\pi$};
  \addplot[only marks, mark=*, mark size=1.3pt, bkOrange] coordinates {(0.0,1.000) (0.3,0.940) (0.6,0.797) (0.9,0.582) (1.2,0.315) (1.5,0.020) (1.8,-0.277) (2.1,-0.549) (2.4,-0.772) (2.7,-0.926) (3.0,-0.997) (3.3,-0.979) (3.6,-0.874) (3.9,-0.690) (4.2,-0.445) (4.5,-0.160) (4.8,0.139) (5.1,0.426) (5.4,0.675) (5.7,0.863) (6.0,0.975) (6.3,0.999)};
\end{axis}
\end{tikzpicture}
\end{minipage}\par\vspace{2pt}
{\small オレンジの点が表計算の $x$、青い線が $\cos t$。}\par
「加速度は位置の $-1$ 倍」というきまりだけから計算したのに、点は $\cos t$ の波にぴったりのり、約6.3秒（$2\pi\approx6.28$）で1周してもどってきました。$\ddot{x}=-\omega^2x$ というきまりが、$\sin$・$\cos$ のゆれを生み出していることが分かります。きまり（微分方程式）から動きを求めることは、第6巻の中心です（6-6）。

\begin{motto}[4]
\textbf{加法定理から導く}\par
4-9の加法定理 $\sin(x+h)=\sin x\cos h+\cos x\sin h$ を使うと、
\[ \frac{\sin(x+h)-\sin x}{h}=\sin x\cdot\frac{\cos h-1}{h}+\cos x\cdot\frac{\sin h}{h} \]
$h=0.01$ のとき、表計算で計算すると、
\[ \frac{\sin h}{h}=0.99998,\qquad \frac{\cos h-1}{h}=-0.005 \]
$h$ をもっと0に近づけると、
\[ \frac{\sin h}{h}\to1,\qquad \frac{\cos h-1}{h}\to0 \]
となるので、$(\sin x)'=\sin x\times0+\cos x\times1=\cos x$ です。半径1の円で、小さい角 $h$ の弧の長さ（$h$）と高さ（$\sin h$）がほぼ等しいことが、$\dfrac{\sin h}{h}\to1$ の意味です。
\end{motto}

\begin{senmonsho}[大学の教科書ふうの書き方]
\begin{bkhonbun}
単振動の変位を $x(t)=A\cos(\omega t+\phi)$ とおくと、速度と加速度は
\[ \dot{x}=-A\omega\sin(\omega t+\phi),\qquad \ddot{x}=-A\omega^2\cos(\omega t+\phi)=-\omega^2x \tag{2.3} \]
となる。$A$ を振幅、$\omega$ を角振動数、$\phi$ を初期位相という。\bkidx{しんぷく}{振幅}\bkidx{かくしんどうすう}{角振動数}\bkidx{しょきいそう}{初期位相}
\end{bkhonbun}
\bkyaku{ゆれの位置を $A\cos(\omega t+\phi)$ とする。速度は、これを微分したもの。加速度は、もう一度微分したもので、位置の $-\omega^2$ 倍になる。$\phi$（ファイ）は「$t=0$ で波がどこから始まるか」を決める角。}
\end{senmonsho}

\begin{column}{サインは「弓のつる」から}
$\sin$ の名前は、長い旅をしてきました。1500年ほど前のインドの天文学では、円の弧を弓に、弦（弧の両はしを結ぶ線分）を弓のつるに見立て、その半分の長さを「つるの半分」とよびました。これが半径1の円での $\sin$ です。アラビア語に訳されたときに音だけが写され、のちにヨーロッパでラテン語に訳すとき、にた形の「ふところ・入り江」という言葉と読みまちがえられ、ラテン語の sinus（入り江）になったといわれます。
\end{column}

\bkkigocards{\bkkigocard{$\phi$}{ファイ}{初期位相。$t=0$ での角}{phase}{5-15}}
{\bkkigocard{$V_{\mathrm{e}}$}{ブイ・イー}{実効値。平均の電力が同じになる直流の電圧。$V_0/\sqrt2$}{effective value}{5-15}}

\begin{taisetsu}
\def\tsl#1{\makebox[3.2\zw][l]{\sffamily\bfseries\small\color{bkTheme}#1}}
\tsl{微分}$(\sin\omega t)'=\omega\cos\omega t$\qquad $(\cos\omega t)'=-\omega\sin\omega t$\\[3pt]
\tsl{}2回微分すると、もとの $-\omega^2$ 倍（$\ddot{x}=-\omega^2x$）\\[3pt]
\tsl{積分}$\displaystyle\int\sin\omega t\,\dd t=-\frac{1}{\omega}\cos\omega t+C$\qquad $\displaystyle\int\cos\omega t\,\dd t=\frac{1}{\omega}\sin\omega t+C$\\[3pt]
\tsl{平均}$\sin^2$ の1周期の平均は $\dfrac{1}{2}$\qquad 実効値 $V_{\mathrm{e}}=\dfrac{V_0}{\sqrt2}$\\[3pt]
\tsl{}角は、いつもラジアンではかる。
\end{taisetsu}

\begin{kotoba}
\small
\bkword{角速度}{1秒に回る角 $\omega$（rad/s）。周期は $\frac{2\pi}{\omega}$。}
\bkword{単振動}{位置が $A\sin\omega t$ や $A\cos\omega t$ で表されるゆれ。$\ddot{x}=-\omega^2x$ をみたす。}
\bkword{振幅・角振動数}{ゆれの大きさ $A$、1秒に進む角 $\omega$。}
\bkword{実効値}{交流と同じ平均の電力になる直流の値。山の高さの $\frac{1}{\sqrt2}$ 倍。}
\end{kotoba}

%% ---------------- 例題 ----------------
\bksection{reidai}{例題でたしかめよう}

\begin{reidai}[$\sin$・$\cos$ をふくむ関数の微分]
次の関数を微分しましょう。\quad（1）$3\sin2t$\quad（2）$-\cos5t$\quad（3）$\sin(3t+1)$\quad（4）$\sin t^2$\par
\kangaekata 内側（角の部分）を $u$ とし、「外側の微分×内側の微分」（5-6）。前の数（定数倍）はそのまま残す（5-4）。\par
\shiki（1）内側 $u=2t$、内側の微分は2。$3\times\cos2t\times2=6\cos2t$\par
\hspace*{2.6\zw}（2）内側 $u=5t$、内側の微分は5。$-1\times(-\sin5t)\times5=5\sin5t$\par
\hspace*{2.6\zw}（3）内側 $u=3t+1$、内側の微分は $(3t+1)'=3$。$\cos(3t+1)\times3$\par
\hspace*{2.6\zw}（4）内側 $u=t^2$、内側の微分は $2t$。$\cos t^2\times2t$\par
\kotae （1）$6\cos2t$\quad（2）$5\sin5t$\quad（3）$3\cos(3t+1)$\quad（4）$2t\cos t^2$
\end{reidai}
（2）では、$\cos$ の微分の $-$ と、はじめからある $-$ で、$-$ が2つになって $+$ になります。（3）の $+1$ は定数なので、微分すると0。内側の微分は3だけです。

\begin{jibun}
次の関数を微分しましょう。\quad（1）$4\cos3t$\quad（2）$\sin\dfrac{t}{2}$\quad（3）$\cos(2t+1)$\quad（4）$-2\sin4t$
\bkjans{（1）\nobreak $-12\sin3t$\quad（2）\nobreak $\dfrac{1}{2}\cos\dfrac{t}{2}$\quad（3）\nobreak $-2\sin(2t+1)$\quad（4）\nobreak $-8\cos4t$}
\end{jibun}

\begin{reidai}[ばねのおもりの速度と加速度]
ばねのおもりの位置が $x=0.1\sin4t$（m）です。速度と加速度を求め、速さが最大になる位置とその速さ、加速度の大きさが最大になる位置とその大きさを答えましょう。\par
\kangaekata 速度は位置の微分、加速度は速度の微分（5-5）。$\sin$ も $\cos$ も、大きさの最大は1。\par
\shiki $v=0.1\times4\cos4t=0.4\cos4t$\par
\hspace*{2.6\zw}$a=0.4\times(-4\sin4t)=-1.6\sin4t$\par
\hspace*{2.6\zw}速さ：$\cos4t=\pm1$ のとき最大の0.4。このとき $\sin4t=0$ で $x=0$\par
\hspace*{2.6\zw}加速度：$\sin4t=\pm1$ のとき最大の1.6。このとき $x=\pm0.1$\par
\kotae 中心（$x=0$）を通るとき速さ最大\qty{0.4}{m/s}。はし（$x=\pm\qty{0.1}{m}$）で加速度の大きさ最大\qty{1.6}{m/s^2}
\end{reidai}
$\cos4t=\pm1$ のとき $\sin4t=0$ になるのは、$\sin^2+\cos^2=1$ だからです。速さが最大の所と、加速度が最大の所は、⑥の図のように $\frac{1}{4}$ 周期ずれています。

\begin{jibun}
$x=0.05\cos10t$（m）の速度・加速度と、速さの最大値・加速度の大きさの最大値を求めましょう。
\bkjans{$v=-0.5\sin10t$、$a=-5\cos10t$。速さ最大\qty{0.5}{m/s}、加速度最大\qty{5}{m/s^2}}
\end{jibun}

\begin{reidai}[円運動の速さと加速度]
半径\qty{0.2}{m}の円の上を毎秒5回転する点の、角速度 $\omega$ と速さを求めましょう。$x$ 座標 $x=0.2\cos\omega t$ の加速度の大きさの最大値も求めましょう。\par
\kangaekata 1回転は $2\pi$\,rad（4-2）。速さは $r\omega$（⑤）。加速度は、$x$ を2回微分して求めます（⑥）。\par
\shiki 1秒に5回転なので、1秒に回る角は $\omega=5\times2\pi=10\pi$\par
\hspace*{2.6\zw}速さは $r\omega=0.2\times10\pi=2\pi\approx6.3$\par
\hspace*{2.6\zw}$\dot{x}=-0.2\omega\sin\omega t$、$\ddot{x}=-0.2\omega^2\cos\omega t$\par
\hspace*{2.6\zw}$\cos\omega t=\pm1$ のとき最大で、$0.2\times(10\pi)^2=0.2\times100\pi^2=20\pi^2$\par
\kotae $\omega=10\pi\approx31$\,rad/s、速さ約\qty{6.3}{m/s}、加速度最大約\qty{197}{m/s^2}
\end{reidai}

\begin{jibun}
半径\qty{0.3}{m}の円の上を、1秒に1回転する点の $\omega$ と速さ、$x=0.3\cos\omega t$ の加速度の大きさの最大値を求めましょう（$\pi^2\approx9.87$）。
\bkjans{$\omega=2\pi$\,rad/s、速さ $0.6\pi\approx\qty{1.9}{m/s}$、加速度最大 $0.3\times4\pi^2\approx\qty{11.8}{m/s^2}$}
\end{jibun}

\begin{reidai}[$\sin$・$\cos$ の定積分]
$\displaystyle\int_0^{\frac{\pi}{2}}\cos t\,\dd t$ と $\displaystyle\int_0^{\pi}\sin2t\,\dd t$ を求めましょう。\par
\kangaekata ⑦の積分の式で原始関数を作り、上端の値から下端の値をひきます（5-11）。$\sin2t$ は $\omega=2$ なので $\frac{1}{2}$ がかかる。\par
\shiki $\displaystyle\int_0^{\frac{\pi}{2}}\cos t\,\dd t=\Bigl[\,\sin t\,\Bigr]_0^{\frac{\pi}{2}}=\sin\frac{\pi}{2}-\sin0=1-0=1$\par\vspace{2pt}
\hspace*{2.6\zw}$\displaystyle\int_0^{\pi}\sin2t\,\dd t=\Bigl[\,-\frac{1}{2}\cos2t\,\Bigr]_0^{\pi}=\Bigl(-\frac{1}{2}\cos2\pi\Bigr)-\Bigl(-\frac{1}{2}\cos0\Bigr)$\par\vspace{2pt}
\hspace*{4.6\zw}$\displaystyle=-\frac{1}{2}+\frac{1}{2}=0$\par
\kotae 1、0
\end{reidai}
2つめが0になるのは、$\sin2t$ が0から $\pi$ でちょうど1周期で、横軸より上の山（$+1$）と下の谷（$-1$）の面積が打ち消しあうからです（5-11）。

\begin{jibun}
$\displaystyle\int_0^{\pi}\sin t\,\dd t$ と $\displaystyle\int_0^{\frac{\pi}{4}}\cos2t\,\dd t$ を求めましょう。
\bkjans{2、$\Bigl[\,\dfrac{1}{2}\sin2t\,\Bigr]_0^{\frac{\pi}{4}}=\dfrac{1}{2}\sin\dfrac{\pi}{2}-0=\dfrac{1}{2}$}
\end{jibun}

\begin{reidai}[$\ddot{x}=-\omega^2x$ を確かめる]
$x=3\cos2t$ が、$\ddot{x}=-4x$ をみたすことを確かめましょう。\par
\kangaekata $x$ を2回微分した $\ddot{x}$ が、$-4x$ と同じ式になるかを見ます（⑥）。\par
\shiki $\dot{x}=3\times(-\sin2t)\times2=-6\sin2t$\par
\hspace*{2.6\zw}$\ddot{x}=-6\times\cos2t\times2=-12\cos2t$\par
\hspace*{2.6\zw}$-12\cos2t=-4\times3\cos2t=-4x$\par
\kotae みたす（$\omega=2$、$\omega^2=4$）
\end{reidai}

\begin{jibun}
$x=5\sin3t$ は $\ddot{x}=-kx$ をみたします。$k$ の値を求めましょう。
\bkjans{$\dot{x}=15\cos3t$、$\ddot{x}=-45\sin3t=-9x$ より $k=9$}
\end{jibun}

\begin{reidai}[観覧車の高さの変わる速さ]
半径\qty{50}{m}、1周\qty{20}{min}の観覧車があります。ゴンドラの高さは、いちばん下を出発してから $t$ 秒後に $y=50-50\cos\omega t$（m）です。上がる速さの最大値と、それが何分後かを求めましょう。最も速く下がるのは何分後ですか。\par
\kangaekata 1周の時間を秒に直して $\omega$ を求める。上がる速さは、高さ $y$ を時間で微分した $\dot{y}$（5-5）。\par
\shiki 1周は $\qty{20}{min}\times\qty{60}{s/min}=\qty{1200}{s}$\par
\hspace*{2.6\zw}1200秒で $2\pi$ 回るので $\omega=\dfrac{2\pi}{1200}$\par
\hspace*{2.6\zw}$\dot{y}=0-50\times(-\omega\sin\omega t)=50\omega\sin\omega t$\par
\hspace*{2.6\zw}最大は $\sin\omega t=1$ のとき\par
\hspace*{4.6\zw}$50\omega=50\times\dfrac{2\pi}{1200}=\dfrac{\pi}{12}\approx0.26$\par
\hspace*{2.6\zw}$\omega t=\dfrac{\pi}{2}$ は $2\pi$ の $\dfrac{1}{4}$ なので、$\qty{20}{min}\times\dfrac{1}{4}=\qty{5}{min}$\par
\hspace*{2.6\zw}最も速く下がるのは $\sin\omega t=-1$ のとき\par
\hspace*{4.6\zw}$\omega t=\dfrac{3\pi}{2}$（$\dfrac{3}{4}$ 周）で、$\qty{20}{min}\times\dfrac{3}{4}=\qty{15}{min}$\par
\kotae 上がる速さの最大は約\qty{0.26}{m/s}で5分後（真横を通るとき）。最も速く下がるのは15分後
\end{reidai}

\begin{jibun}
半径\qty{30}{m}、1周\qty{15}{min}の観覧車で、上がる速さの最大値を求めましょう。また、最も速く下がるのは何分何秒後ですか。
\bkjans{$\qty{15}{min}\times\qty{60}{s/min}=\qty{900}{s}$。$30\times\dfrac{2\pi}{900}=\dfrac{\pi}{15}\approx\qty{0.21}{m/s}$。最も速く下がるのは $15\times\dfrac{3}{4}=11.25$ 分後で、11分15秒後}
\end{jibun}

\begin{reidai}[$\sin^2$・$\cos^2$ の積分]
$\displaystyle\int_0^{\pi}\sin^2t\,\dd t$ を求めましょう。\par
\kangaekata 2倍角の式で、$\sin^2t$ を積分できる $\cos2t$ の形に直す（⑨）。\par
\shiki $\sin^2t=\dfrac{1-\cos2t}{2}=\dfrac{1}{2}-\dfrac{1}{2}\cos2t$\par\vspace{2pt}
\hspace*{2.6\zw}$\displaystyle\int_0^{\pi}\Bigl(\frac{1}{2}-\frac{1}{2}\cos2t\Bigr)\dd t=\Bigl[\,\frac{t}{2}-\frac{\sin2t}{4}\,\Bigr]_0^{\pi}=\Bigl(\frac{\pi}{2}-\frac{\sin2\pi}{4}\Bigr)-0=\frac{\pi}{2}$\par
\kotae $\dfrac{\pi}{2}$（区間の長さ $\pi$ の $\frac{1}{2}$。平均が $\frac{1}{2}$ だから）
\end{reidai}

\begin{jibun}
$\displaystyle\int_0^{\frac{\pi}{2}}\cos^2t\,\dd t$ を求めましょう。
\bkjans{$\displaystyle\int_0^{\frac{\pi}{2}}\frac{1+\cos2t}{2}\,\dd t=\Bigl[\,\frac{t}{2}+\frac{\sin2t}{4}\,\Bigr]_0^{\frac{\pi}{2}}=\frac{\pi}{4}$}
\end{jibun}

\begin{reidai}[交流の実効値と電力]
（1）実効値\qty{100}{V}のコンセントに、\qty{20}{\ohm}の抵抗をつなぎます。電圧の最大値 $V_0$ と、平均の電力を求めましょう。\par
（2）電圧の最大値が\qty{200}{V}の交流の実効値を求めましょう（$\sqrt2=1.41$）。\par
\kangaekata ⑨より、$V_0=\sqrt2\,V_{\mathrm{e}}$、$V_{\mathrm{e}}=\dfrac{V_0}{\sqrt2}$。平均の電力は、実効値と同じ電圧の直流のときと同じ。\par
\shiki（1）$V_0=\sqrt2\times100=1.414\times100\approx141$\par
\hspace*{4.6\zw}平均の電力は $\dfrac{V_{\mathrm{e}}^2}{R}=\dfrac{100^2}{20}=\dfrac{10000}{20}=500$\par
\hspace*{2.6\zw}（2）$V_{\mathrm{e}}=\dfrac{200}{\sqrt2}=\dfrac{200}{1.41}\approx141.8$\par
\kotae （1）約\qty{141}{V}、\qty{500}{W}\quad（2）約\qty{142}{V}
\end{reidai}

\begin{jibun}
最大値\qty{170}{V}の交流の実効値を求めましょう。また、\qty{50}{\ohm}の抵抗での平均の電力を求めましょう（$\sqrt2=1.41$）。
\bkjans{$V_{\mathrm{e}}=\dfrac{170}{1.41}\approx\qty{120}{V}$、平均の電力 $\dfrac{V_0^2}{2R}=\dfrac{170^2}{2\times50}=\qty{289}{W}$}
\end{jibun}

%% ---------------- まちがえやすいところ・物理工学 ----------------
\bksection{tsumazuki}{まちがえやすいところをチェック}

\begin{tsumazuki}{$\cos$ の微分の $-$ をわすれる}
\bkng{$(\cos t)'=\sin t$}\bkok{$(\cos t)'=-\sin t$}
$t=0$ の少しあとでは、$\cos t$ は1からへっていきます（傾きが負）。$\sin t$ は正なので、$-$ がないと合いません。グラフを思いうかべて確かめましょう。
\end{tsumazuki}

\begin{tsumazuki}{$\omega$ を前に出しわすれる・積分でかけてしまう}
\bkng{$\displaystyle\int\cos3t\,\dd t=3\sin3t+C$}\bkok{$\displaystyle\int\cos3t\,\dd t=\frac{1}{3}\sin3t+C$}
微分では $\omega$ をかけ、積分では $\omega$ でわります。答えを微分して確かめましょう。$\Bigl(\dfrac{1}{3}\sin3t\Bigr)'=\dfrac{1}{3}\times3\cos3t=\cos3t$ で、もとにもどります。
\end{tsumazuki}

\begin{tsumazuki}{$\sin^2t$ と $\sin t^2$ を取りちがえる}
$\sin^2t$ は $(\sin t)^2$（$\sin t$ を2乗する）。$\sin t^2$ は $\sin(t^2)$（角が $t^2$）。$t=2$ なら、$\sin^2t=(\sin2)^2=0.83$、$\sin t^2=\sin4=-0.76$ で、まったくちがう数です。
\end{tsumazuki}

\begin{tsumazuki}{電卓や表計算を「度」のまま使う}
$\sin1$ は、度のモードで0.0175、ラジアンのモードで0.8415。式に入れるのはラジアンの値です。電卓の表示（DEG／RAD）を確かめましょう。
\end{tsumazuki}

\begin{butsuri}
\begin{itemize}[leftmargin=1.2\zw, itemsep=1pt, topsep=0pt, label={\color{bkPurple}●}]
\item \textbf{ばねと振り子}\par\noindent 速さは中心で最大、はしで0。加速度ははしで最大、中心で0。$\sin$ と $\cos$ の $\frac{1}{4}$ 周期のずれが、このちがいを生みます。
\item \textbf{交流とコンデンサー}\par\noindent 家のコンセントの電圧は $V_0\sin\omega t$ の交流です。コンデンサーに流れる電流は、電圧の変わる速さ（電圧の微分）に比例するので、$\cos\omega t$ の形になり、電圧より $\frac{1}{4}$ 周期先に進みます（位相のずれ、6-10）。\par
\begin{tikzpicture}[tl/.style={font=\sffamily\footnotesize, text=bkInk, anchor=base, yshift=-6mm}]
\begin{axis}[bkplot, width=118mm, height=31mm, xmin=0, xmax=6.8, ymin=-1.15, ymax=1.15, xtick={1.5708,3.1416,4.7124,6.2832}, xticklabels={}, ytick=\empty, xlabel={$\omega t$}]
  \addplot[bkcurve, domain=0:6.6, samples=120] {sin(deg(x))};
  \addplot[bkcurve2, domain=0:6.6, samples=120] {0.8*cos(deg(x))};
  \node[tl] at (axis cs:1.5708,-1.15) {$\dfrac{\pi}{2}$};
  \node[tl] at (axis cs:3.1416,-1.15) {$\pi$};
  \node[tl] at (axis cs:4.7124,-1.15) {$\dfrac{3\pi}{2}$};
  \node[tl] at (axis cs:6.2832,-1.15) {$2\pi$};
  % 「電圧」「電流」は、格子の線が通らないように、格子の上（ymax の外）に置く
  \node[font=\sffamily\scriptsize, text=bkBlue, anchor=south west, inner sep=1pt, yshift=1.5mm] at (axis cs:1.70,1.15) {電圧};
  \node[font=\sffamily\scriptsize, text=bkOrange!80!black, anchor=south east, inner sep=1pt, yshift=1.5mm] at (axis cs:6.25,1.15) {電流（電圧より $\frac{1}{4}$ 周期先に山が来る）};
\end{axis}
\end{tikzpicture}
\item \textbf{音とスピーカー}\par\noindent スピーカーの振動板は $\sin$ の形でゆれて音を出します。振動板の加速度の大きさは $A\omega^2$ で、振動数の2乗に比例するので、同じ振幅なら、高い音ほど大きな加速度が必要です。
\end{itemize}
\end{butsuri}

%% ---------------- 章末問題 ----------------
\bkendproblems
\bkproblemgroup{基本}{3}

\bkproblem 次の関数を微分しましょう。\quad（1）$2\sin5t$\quad（2）$-\cos4t$\quad（3）$\sin(3t-1)$
\bkpans{（1）\nobreak $10\cos5t$\quad（2）\nobreak $4\sin4t$\quad（3）\nobreak $3\cos(3t-1)$}

\bkproblem $x=0.2\cos5t$（m）のおもりについて、速度と加速度を求めましょう。速さの最大値はいくつですか。
\bkpans{$v=-\sin5t$、$a=-5\cos5t$。速さ最大\qty{1}{m/s}}

\bkproblem 次の定積分を求めましょう。\\
（1）$\displaystyle\int_0^{\pi}\sin t\,\dd t$\qquad（2）$\displaystyle\int_0^{\frac{\pi}{2}}\cos t\,\dd t$\qquad（3）$\displaystyle\int_0^{\frac{\pi}{3}}\sin3t\,\dd t$
\bkpans{（1）\nobreak 2\quad（2）\nobreak 1\quad（3）\nobreak \mbox{$\Bigl[\,-\dfrac{1}{3}\cos3t\,\Bigr]_0^{\frac{\pi}{3}}=\dfrac{2}{3}$}}

\bkproblem $x=2\sin6t$ が $\ddot{x}=-kx$ をみたすときの $k$ を求めましょう。
\bkpans{$k=36$}

\bkproblem $\displaystyle\int_0^{2\pi}\cos^2t\,\dd t$ を求めましょう。また、$\cos^2t$ の1周期の平均はいくつですか。
\bkpans{$\Bigl[\,\dfrac{t}{2}+\dfrac{\sin2t}{4}\,\Bigr]_0^{2\pi}=\pi$、平均 $\dfrac{1}{2}$}

\bkproblemgroup{チャレンジ}{4}

\bkproblem 半径\qty{0.5}{m}の円の上を、1秒に2回転する点があります。$\omega$ と速さを求めましょう。また、点の $x$ 座標 $x=0.5\cos\omega t$ の加速度の大きさの最大値を求めましょう。
\bkpans{$\omega=4\pi$\,rad/s、速さ $0.5\times4\pi=2\pi\approx\qty{6.3}{m/s}$。加速度最大 $0.5\times16\pi^2=8\pi^2\approx\qty{79}{m/s^2}$}

\bkproblem 半径\qty{40}{m}、1周\qty{16}{min}の観覧車があります。いちばん下を出発してから $t$ 秒後のゴンドラの高さは $y=40-40\cos\omega t$（m）です。上がる速さの最大値を求めましょう。また、高さが最も速く下がるのは何分後ですか。
\bkpans{\mbox{$\qty{16}{min}\times\qty{60}{s/min}=\qty{960}{s}$}、\mbox{$\omega=\dfrac{2\pi}{960}$}。最大 $40\omega=\dfrac{\pi}{12}\approx\qty{0.26}{m/s}$。最も速く下がるのは $\frac{3}{4}$ 周の12分後}

\bkproblem 実効値\qty{200}{V}の交流の、電圧の最大値を求めましょう。また、\qty{40}{\ohm}の抵抗での平均の電力を求めましょう。
\bkpans{$200\sqrt2\approx\qty{283}{V}$、$\dfrac{200^2}{40}=\qty{1000}{W}$}

\begin{furikaeri}
\begin{checklist}
\item 円運動の速度で、$\sin$・$\cos$ の微分を説明できる。\hfill{\footnotesize\sffamily\color{bkTheme}→\pageref{bk:5-15:en}ページ}
\item $\sin\omega t$、$\cos\omega t$ を微分・積分できる。\hfill{\footnotesize\sffamily\color{bkTheme}→\pageref{bk:5-15:omega}ページ}
\item 単振動の $\ddot{x}=-\omega^2x$ を確かめられる。\hfill{\footnotesize\sffamily\color{bkTheme}→\pageref{bk:5-15:nikai}ページ}
\item $\sin^2$ を積分して、交流の実効値を求められる。\hfill{\footnotesize\sffamily\color{bkTheme}→\pageref{bk:5-15:jikko}ページ}
\end{checklist}
\end{furikaeri}

\bkfukidashi{mina}{$e$ も $\sin$ も、微分の計算ができるようになった。次は、積分のむずかしい形にちょうせんだよ。}

%% 5-16 積分の道具箱：置換積分と部分積分（第11回：ページ制限なし。積分の道具を1段ずつ書き足した）
\newcommand\bkvpmidashi[1]{\par\Needspace{6\baselineskip}\noindent\textbf{#1}\par}
\bkchapter{5}{16}{積分の道具箱：\\置換積分と部分積分}{\bklevel{4}}{積分の道具箱：置換積分と部分積分}
\bkmichishirube{5-6,5-11,5-14}{6-4,9-1,9-5}

\begin{goal}
\begin{itemize}[label={\color{bkTheme}●}, leftmargin=1.6\zw, itemsep=0pt, topsep=1pt]
\item \bklevel{4}\enspace 合成関数の微分を逆に使って、置換積分ができる。
\item \bklevel{4}\enspace 積の微分を逆に使って、部分積分ができる。
\item \bklevel{4}\enspace 流れた電気の量や円の面積など、表にない形の積分を求められる。
\end{itemize}
\end{goal}

\begin{hatena}
\noindent\begin{minipage}[c]{0.535\linewidth}
ロボットのモーターのスイッチを入れると、電流が $i=2t\,e^{-t}$（A、$t$ は秒）のように、ふえてからへっていきました。はじめの3秒間に流れた電気の量（グラフの面積）は、何C（クーロン）でしょう。

流れた電気の量は、電流の積分です（5-12）。でも、$t$ と $e^{-t}$ のかけ算の積分は、まだ表にありません。
\end{minipage}\hfill
\begin{minipage}[c]{0.425\linewidth}
\begin{tikzpicture}
\begin{axis}[bkplot, width=56mm, height=38mm, xmin=0, xmax=4.2, ymin=0, ymax=0.95, xtick={1,2,3,4}, ytick={0.5}, xlabel={$t$}, ylabel={$i$}]
  \addplot[bkcurve, domain=0:3, samples=60, fill=bkYellow!45, draw=none] {2*x*exp(-x)} \closedcycle;
  \addplot[bkcurve, domain=0:4.1, samples=80] {2*x*exp(-x)};
  \node[font=\sffamily\bfseries\small] at (axis cs:1.5,0.35) {？};
\end{axis}
\end{tikzpicture}
\end{minipage}
\end{hatena}

積分は「微分すると、その式になる関数」をさがす計算でした（5-11）。この章では、微分の2つの道具（5-6）を「逆向き」に使って、さがし方の道具を2つ作ります。

\clearpage
%% ---------------- やってみよう・思い出そう ----------------
\begin{yattemiyou}[表計算で、面積をおおよそ求めよう]
5-10と同じように、0から3を、はば0.5の細い長方形6本に分け、長方形の面積をたします。高さは、それぞれの区間のまん中の $t$ での $i$ の値にします。
\begin{bkcenter}
\bksheet{4}{22}{%
1 & 区間のまん中 $t$ & $i=2t\,e^{-t}$ & はば & 長方形の面積 \\
2 & 0.25 & 0.3894 & 0.5 & 0.1947 \\
3 & 0.75 & 0.7085 & 0.5 & 0.3543 \\
4 & 1.25 & 0.7163 & 0.5 & 0.3581 \\
5 & 1.75 & 0.6082 & 0.5 & 0.3041 \\
6 & 2.25 & 0.4743 & 0.5 & 0.2371 \\
7 & 2.75 & 0.3516 & 0.5 & 0.1758 \\
8 & & & 合計 & 1.6242 \\}
\bktabcaption{はばを0.1にすると1.6026、0.01にすると1.6017。1.60…に近づいていく。}
\end{bkcenter}
\end{yattemiyou}

\noindent 表計算なら、およそ1.60Cと分かります。でも、「1.60…」のもとになる、正確な式は何でしょう。この章の道具を使うと、答えは $2-8e^{-3}$ と分かります（⑩と例題8）。

\begin{omoidasou}
\bkrecall{5-4}{定数倍は前に出せる：$(cf)'=cf'$。積分でも $\int cf\,\dd x=c\int f\,\dd x$}
\bkrecall{5-6}{積の微分 $(fg)'=f'g+fg'$。合成関数の微分 $\dfrac{\dd y}{\dd x}=\dfrac{\dd y}{\dd u}\times\dfrac{\dd u}{\dd x}$（外側の微分×内側の微分）}
\bkrecall{5-8}{$x$ が少し（$\Delta x$）変わると、$y$ は約 $f'(x)\,\Delta x$ 変わる。$\Delta y\approx f'(x)\,\Delta x$}
\bkrecall{5-11}{微分して $f(x)$ になる関数 $F(x)$ を見つければ、$\displaystyle\int_a^bf(x)\,\dd x=F(b)-F(a)$。$\displaystyle\int x^n\,\dd x=\frac{x^{n+1}}{n+1}+C$}
\bkrecall{5-14}{$\displaystyle\int e^{ax}\,\dd x=\frac{1}{a}e^{ax}+C$、$\displaystyle\int\frac{1}{x}\,\dd x=\ln x+C$（$x>0$）}
\bkrecall{5-15}{$\displaystyle\int\sin t\,\dd t=-\cos t+C$、$\displaystyle\int\cos t\,\dd t=\sin t+C$、$\cos^2\theta=\dfrac{1+\cos2\theta}{2}$}
\end{omoidasou}

\Needspace{14\baselineskip}
%% ---------------- 数学の道具 ----------------
\bksection{doguu}{数学の道具}

\bksubsection{1}{積分は「微分の答え」から見つける}

\label{bk:5-16:hajime}5-11では、微分の表を逆に読んで、積分の表を作りました。
\begin{bkcenter}
\begin{bktable}{colspec={Q[c,m,wd=44mm] Q[c,m,wd=50mm]}, rowsep=2pt}
微分の式（左を微分すると右） & 積分の式（右を積分すると左） \\
$(x^3)'=3x^2$ & $\displaystyle\int3x^2\,\dd x=x^3+C$ \\
$(e^{2x})'=2e^{2x}$ & $\displaystyle\int2e^{2x}\,\dd x=e^{2x}+C$ \\
$(\sin x)'=\cos x$ & $\displaystyle\int\cos x\,\dd x=\sin x+C$ \\
\end{bktable}
\end{bkcenter}
積分したい式が、表の右の列にそのまま出ていれば、答えはすぐ分かります。では、$2x(x^2+1)^3$ や $x\,e^x$ はどうでしょう。こんな形は、表にはありません。

でも、5-6で、積の微分と合成関数の微分を使って、たくさんの式を微分しました。\textbf{その微分の答えを集めておけば、それを逆に読んで積分できる}はずです。この章の2つの道具は、この考え方から生まれます。
\begin{bkcenter}
\begin{bktable}{colspec={Q[l,m,wd=46mm] Q[l,m,wd=46mm]}, rowsep=2pt}
微分の道具（5-6） & 逆に使った積分の道具 \\
合成関数の微分 & 置換積分（②〜⑥） \\
積の微分 & 部分積分（⑦〜⑩） \\
\end{bktable}
\end{bkcenter}

\bksubsection{2}{合成関数の微分を、逆に読む}

\label{bk:5-16:gyaku}合成関数の微分で、$(x^2+1)^4$ を微分してみます。かっこの中を内側 $u=x^2+1$、4乗を外側 $u^4$ と見ます（5-6）。
\begin{align*}
\bigl\{(x^2+1)^4\bigr\}'&=\underbrace{4u^3}_{\text{外側の微分}}\times\underbrace{(x^2+1)'}_{\text{内側の微分}}=4(x^2+1)^3\times2x\\
&=8x(x^2+1)^3
\end{align*}
この式を逆に読むと、積分の式が1つ手に入ります。
\[ \int8x(x^2+1)^3\,\dd x=(x^2+1)^4+C \]
積分の中の $8x(x^2+1)^3$ を、色で分けて見てみましょう。

\begin{bkfigure}
\begin{tikzpicture}[x=1mm, y=1mm]
  \node[font=\large, inner sep=1pt, anchor=base east] (p1) at (0,0) {$\displaystyle\int\ 4\,$};
  \node[font=\large, inner sep=1pt, anchor=base west, text=bkBlue] (p2) at (p1.base east) {$(x^2+1)^3$};
  \node[font=\large, inner sep=1pt, anchor=base west] (p3) at (p2.base east) {$\,\times\,$};
  \node[font=\large, inner sep=1pt, anchor=base west, text=bkOrange!80!black] (p4) at (p3.base east) {$2x$};
  \node[font=\large, inner sep=1pt, anchor=base west] (p5) at (p4.base east) {$\,\dd x$};
  \node[bk node, draw=bkBlue, fill=bkBlue!6, font=\sffamily\footnotesize, align=center, anchor=north] (k1) at (-22,-14) {内側 $u=x^2+1$ が\\入った式};
  \node[bk node, draw=bkOrange, fill=bkOrange!6, font=\sffamily\footnotesize, align=center, anchor=north] (k2) at (40,-14) {内側の微分\\$(x^2+1)'=2x$};
  \draw[bk link, bkBlue] (k1.north) -- (p2.south);
  \draw[bk link, bkOrange] (k2.north) -- (p4.south);
\end{tikzpicture}
\bkfigcaption{積分の中に、「内側の式」と「内側の微分」が両方入っている。これが、合成関数の微分を逆に読める目じるし。}
\end{bkfigure}

図では、$8x$ を $4\times2x$ と分けて、「内側の微分 $2x$」が見えるようにしました。では、前の数が4でなく1のとき、$\displaystyle\int2x(x^2+1)^3\,\dd x$ はどうなるでしょう。$2x(x^2+1)^3$ は $8x(x^2+1)^3$ のちょうど $\frac{1}{4}$ です。だから答えも $\frac{1}{4}$ にすればよいはずです（定数倍は前に出せる、5-4）。
\[ \int2x(x^2+1)^3\,\dd x=\frac{1}{4}\int8x(x^2+1)^3\,\dd x=\frac{(x^2+1)^4}{4}+C \]
微分して確かめましょう。$\Bigl\{\dfrac{(x^2+1)^4}{4}\Bigr\}'=\dfrac{1}{4}\times8x(x^2+1)^3=2x(x^2+1)^3$。積分の中身にもどりました。

\bksubsection{3}{置換積分：$u$ におきかえる}

\label{bk:5-16:chikan}②では、答えを知っている式から逆にたどりました。こんどは、答えを知らないところから、手順だけで $\displaystyle\int2x(x^2+1)^3\,\dd x$ を求めます。
\begin{bkcenter}
% 第6回：3つの枠の大きさ（はば40mm）と上端をそろえ、中の字は上からそろえて組む。矢印は枠の名前どうしでつなぐ
\begin{tikzpicture}[x=1mm, y=1mm, st/.style={bk node, draw=bkTheme, fill=bkTheme!6, align=left, anchor=north west, font=\small}]
  \node[st] (a) at (0,0) {\parbox[t][23.5mm][t]{33mm}{\textbf{①おく}\par\vspace{2pt}内側を\par $u=x^2+1$ とおく。}};
  \node[st] (b) at (49.5,0) {\parbox[t][23.5mm][t]{33mm}{\textbf{②書きかえる}\par\vspace{2pt}$\dd u=2x\,\dd x$\par\vspace{2pt}$\displaystyle\int u^3\,\dd u$}};
  \node[st] (c) at (99,0) {\parbox[t][23.5mm][t]{33mm}{\textbf{③積分してもどす}\par\vspace{2pt}$\dfrac{u^4}{4}+C$\par\vspace{2pt}$=\dfrac{(x^2+1)^4}{4}+C$}};
  \draw[bk link, bkTheme] (a) -- (b);
  \draw[bk link, bkTheme] (b) -- (c);
\end{tikzpicture}
\end{bkcenter}
いちばん大事なのは、②の $\dd u=2x\,\dd x$ です。2通りの見方で説明します。

\bkvpmidashi{見方1：微分の記号を分数のように使う}
$u=x^2+1$ を $x$ で微分すると $\dfrac{\dd u}{\dd x}=2x$。両辺に $\dd x$ をかけると、$\dd u=2x\,\dd x$ になります。$\dfrac{\dd u}{\dd x}$ はほんとうの分数ではありませんが、もとは小さな変化の比 $\dfrac{\Delta u}{\Delta x}$ なので、このように計算してよいのです（5-6のコラム）。分数のように書きかえられるのは、ライプニッツ（5-3）が、計算しやすいように記号を工夫したおかげです。

\bkvpmidashi{見方2：数で確かめる}
$x$ が1から1.01まで、0.01だけ変わるとき、$u=x^2+1$ はどれだけ変わるでしょう。
\[ u(1.01)-u(1)=(1.01^2+1)-(1^2+1)=2.0201-2=0.0201 \]
一方、$2x\,\dd x$ に $x=1$、$\dd x=0.01$ を入れると、$2\times1\times0.01=0.02$。ほとんど同じです。「$x$ が少し（$\dd x$）変わると、$u$ はその $2x$ 倍だけ変わる」という意味です（5-8の $\Delta y\approx f'(x)\,\Delta x$ と同じ）。

\bkvpmidashi{書きかえ}
積分の中を、$u$ の式とそれ以外に分けて、1つずつ書きかえます。
\[ \int\underbrace{(x^2+1)^3}_{u^3}\ \underbrace{2x\,\dd x}_{\dd u}=\int u^3\,\dd u \]
$u$ だけの式になったので、5-11の「指数を1ふやして、その数でわる」で積分できます。$\dfrac{u^4}{4}+C$。最後に、$u$ を $x^2+1$ にもどして、$\dfrac{(x^2+1)^4}{4}+C$。②で求めた答えと同じになりました。

このように、積分の中の式を $u$ におきかえて、知っている形にする方法を\textbf{置換積分}といいます。「置換」は「おきかえる」という意味です。\bkidx{ちかんせきぶん}{置換積分}まとめると、
\[ \int(\text{$u$ の式})\times u'\,\dd x=\int(\text{$u$ の式})\,\dd u \]
です。$u'\,\dd x$（内側の微分×$\dd x$）が、まとめて $\dd u$ になります。

\bksubsection{4}{係数が合わないとき}

\label{bk:5-16:keisu}$\displaystyle\int x\,e^{x^2}\,\dd x$ を求めます。$e$ の肩の $x^2$ を内側と見て、$u=x^2$ とおきます。
\begin{enumerate}[label={\bkcircled[bkTheme]{\arabic*}}, leftmargin=2.4\zw, itemsep=1pt, topsep=2pt]
\item $u=x^2$ を微分すると $\dfrac{\dd u}{\dd x}=2x$ なので、$\dd u=2x\,\dd x$。
\item でも、積分の中には $x\,\dd x$ しかありません。$2x\,\dd x$ の半分です。そこで、$\dd u=2x\,\dd x$ の両辺を2でわって、$x\,\dd x=\dfrac{1}{2}\dd u$。
\item 書きかえると、$\displaystyle\int e^{u}\cdot\frac{1}{2}\,\dd u=\frac{1}{2}\int e^u\,\dd u=\frac{1}{2}e^u+C$（$\frac{1}{2}$ は定数なので前に出せる）。
\item $u$ をもどして、$\dfrac{1}{2}e^{x^2}+C$。
\end{enumerate}
確かめ：$\Bigl(\dfrac{1}{2}e^{x^2}\Bigr)'=\dfrac{1}{2}\times e^{x^2}\times2x=x\,e^{x^2}$。合っています。

\bkvpmidashi{使えないとき}
2や $\frac{1}{2}$ のような\textbf{数}のずれは、このように調整できます。でも、$\displaystyle\int e^{x^2}\,\dd x$ のように $x$ そのものが足りないときは、置換積分は使えません。$\dd x=\dfrac{\dd u}{2x}$ と書きかえても、分母に $x$ が残って、「$u$ だけの式」にならないからです（$x$ は定数ではないので、前に出せません）。

\bkvpmidashi{よく出る形：分母の微分が分子にある}
$\displaystyle\int\frac{2x}{x^2+1}\,\dd x$ では、分母 $x^2+1$ を微分すると分子の $2x$ になっています。そこで $u=x^2+1$（いつも正）とおくと、$\dd u=2x\,\dd x$ で、
\[ \int\frac{2x}{x^2+1}\,\dd x=\int\frac{1}{x^2+1}\cdot2x\,\dd x=\int\frac{1}{u}\,\dd u=\ln u+C=\ln(x^2+1)+C \]
（$\displaystyle\int\frac{1}{u}\,\dd u=\ln u+C$ は5-14。）確かめ：$\{\ln(x^2+1)\}'=\dfrac{1}{x^2+1}\times2x$（5-14と合成関数の微分）。

\bksubsection{5}{定積分の置換は、はんいも書きかえる}

\label{bk:5-16:hanni}定積分 $\displaystyle\int_0^12x(x^2+1)^3\,\dd x$ を求めます。

\bkvpmidashi{やり方1：$u$ のはんいで計算する}
$u=x^2+1$ とおくと、積分の中は $u$ だけの式になります。このとき、「$x$ が0から1まで」も、「$u$ がいくつからいくつまで」に書きかえます。積分の横軸が、$x$ から $u$ に変わったからです。
\begin{bkcenter}
\begin{bktable}{colspec={Q[l,m,wd=30mm] Q[c,m,wd=26mm] Q[c,m,wd=26mm]}, rowsep=2pt}
 & 下端 & 上端 \\
$x$ & 0 & 1 \\
$u=x^2+1$ & $0^2+1=1$ & $1^2+1=2$ \\
\end{bktable}
\end{bkcenter}
\[ \int_0^12x(x^2+1)^3\,\dd x=\int_1^2u^3\,\dd u=\Bigl[\,\frac{u^4}{4}\,\Bigr]_1^2=\frac{2^4}{4}-\frac{1^4}{4}=\frac{16}{4}-\frac{1}{4}=\frac{15}{4} \]
$x$ にもどさなくても、$u$ のはんいで計算すれば答えが出ます。

\bkvpmidashi{やり方2：$x$ にもどしてから計算する}
③の不定積分 $\dfrac{(x^2+1)^4}{4}$ を使って、$x$ のはんいで計算しても、
\[ \Bigl[\,\frac{(x^2+1)^4}{4}\,\Bigr]_0^1=\frac{2^4}{4}-\frac{1^4}{4}=\frac{15}{4} \]
で、同じ答えです。やり方1の $u=2$、$u=1$ は、$x=1$、$x=0$ のときの $x^2+1$ の値だったので、同じ計算になるのです。

\begin{bkfigure}
\begin{tikzpicture}
\begin{axis}[bkplot, name=L, width=60mm, height=38mm, xmin=0, xmax=1.15, ymin=0, ymax=19, xtick={0,1}, ytick={8,16}, xlabel={$x$}, title={\small $2x(x^2+1)^3$ を $x$：0〜1で}, title style={yshift=-1mm}]
  \addplot[bkcurve, domain=0:1, samples=50, fill=bkBlue!15] {2*x*(x^2+1)^3} \closedcycle;
  \addplot[bkcurve, domain=0:1.04, samples=50] {2*x*(x^2+1)^3};
\end{axis}
\begin{axis}[bkplot, at={(L.east)}, anchor=west, xshift=16mm, width=60mm, height=38mm, xmin=0.8, xmax=2.2, ymin=0, ymax=9.5, xtick={1,2}, ytick={4,8}, xlabel={$u$}, title={\small $u^3$ を $u$：1〜2で}, title style={yshift=-1mm}]
  \addplot[bkcurve2, domain=1:2, samples=50, fill=bkOrange!18] {x^3} \closedcycle;
  \addplot[bkcurve2, domain=0.8:2.1, samples=50] {x^3};
\end{axis}
\node[font=\sffamily\small] at ($(L.east)+(8mm,0)$) {$=$};
\end{tikzpicture}
\bkfigcaption{形も横のはんいもちがうのに、色の部分の面積はどちらも $\frac{15}{4}=3.75$（細い長方形1000本でたしても3.7500）。$x$ の目もりを $u=x^2+1$ の目もりにつけかえると、同じ面積を別の見方で計算できる。}
\end{bkfigure}
左の図で、$x$ のあいだ $\dd x$ の細い長方形（高さ $2x(x^2+1)^3$）の面積は、$(x^2+1)^3\times2x\,\dd x=u^3\,\dd u$ で、右の図の、はば $\dd u$ の細い長方形の面積と同じです。細い長方形が1本ずつ同じ面積なので、全部たした面積も同じになります。

\bksubsection{6}{$x$ を角でおきかえる：円の面積}

\label{bk:5-16:en}置換積分は、「$x$ を別の文字の式でおきかえる」向きにも使えます。半径 $r$ の円の面積を、積分で求めてみましょう。

\bkvpmidashi{円を細い長方形に分ける}
\noindent\begin{minipage}[c]{0.4\linewidth}
\centering
\begin{tikzpicture}[x=1mm, y=1mm]
  \fill[bkBlue!8] (0,0) circle (18);
  \draw[bkBlue, line width=1pt] (0,0) circle (18);
  \draw[bkGray, line width=0.5pt, -{Stealth[length=3pt]}] (-21,0) -- (22,0) node[anchor=west, font=\sffamily\scriptsize] {$x$};
  \draw[bkGray, line width=0.5pt, -{Stealth[length=3pt]}] (0,-21) -- (0,22) node[anchor=south, font=\sffamily\scriptsize] {$y$};
  \fill[bkTheme!35] (9.52,-14.74) rectangle (11.12,14.74);
  \draw[bkTheme!80!black, line width=0.5pt] (9.52,-14.74) rectangle (11.12,14.74);
  \draw[bkInk, line width=0.8pt] (0,0) -- (10.32,14.74);
  \draw[bkBlue!80!black, line width=1.8pt] (0,-0.1) -- (10.32,-0.1);
  \draw[bkOrange, line width=1.8pt] (12.3,0) -- (12.3,14.74);
  \draw[bkInk, line width=0.5pt] (0,6) arc (90:55:6);
  \node[font=\sffamily\scriptsize] at (72.5:8.6) {$\theta$};
  \node[font=\sffamily\scriptsize] at (6.9,6.0) {$r$};
  \fill[bkInk] (10.32,14.74) circle (0.8);
\end{tikzpicture}
\end{minipage}\hfill
\begin{minipage}[c]{0.56\linewidth}
円を、はば $\dd x$ のたての細長い長方形に分けます。位置 $x$ での長方形の上はしは円の上の点なので、中心からの距離は $r$。三平方の定理（4-4）から、上はしの高さ $y$ は
\[ x^2+y^2=r^2\quad\text{より}\quad y=\sqrt{r^2-x^2} \]
長方形は上下に同じだけあるので、高さは $2\sqrt{r^2-x^2}$。たとえば $r=5$、$x=3$ なら $y=\sqrt{25-9}=4$ で、高さ8です。
\end{minipage}\par\vspace{3pt}
この長方形を $x=-r$ から $r$ まで集めると、円の面積は
\[ S=\int_{-r}^{r}2\sqrt{r^2-x^2}\,\dd x \]
$\sqrt{\phantom{x}}$ の中がひき算で、このままでは積分できません。

\bkvpmidashi{図を見て、角 $\theta$ でおきかえる}
図の半径（長さ $r$）が、$y$ 軸となす角を $\theta$ とします。半径を斜めの辺とする直角三角形を見ると、
\begin{itemize}[leftmargin=1.2\zw, itemsep=1pt, topsep=2pt, label={\color{bkTheme}●}]
\item \textcolor{bkBlue!80!black}{横の長さ（青）} $x=r\sin\theta$
\item \textcolor{bkOrange!80!black}{たての長さ（オレンジ）} $\sqrt{r^2-x^2}=r\cos\theta$
\end{itemize}
（$\theta$ を $y$ 軸からはかったので、横が $\sin$、たてが $\cos$ になります。）だから、$x=r\sin\theta$ とおけば、$\sqrt{\phantom{x}}$ がはずれて $r\cos\theta$ になるはずです。式でも確かめます。$\sin^2\theta+\cos^2\theta=1$（4-8）から
\[ r^2-x^2=r^2-r^2\sin^2\theta=r^2(1-\sin^2\theta)=r^2\cos^2\theta,\qquad \sqrt{r^2\cos^2\theta}=r\cos\theta \]
（$-\frac{\pi}{2}\leqq\theta\leqq\frac{\pi}{2}$ では $\cos\theta\geqq0$ なので、$\sqrt{\phantom{x}}$ をはずしても $+$ のまま。）

\bkvpmidashi{$\dd x$ とはんいを書きかえる}
$x=r\sin\theta$ を $\theta$ で微分すると $\dfrac{\dd x}{\dd\theta}=r\cos\theta$（5-15）。だから $\dd x=r\cos\theta\,\dd\theta$ です。はんいは、図で半径を左はしから右はしまで動かして考えます。
\begin{bkcenter}
\begin{bktable}{colspec={Q[l,m,wd=30mm] *{3}{Q[c,m,wd=20mm]}}, rowsep=2pt}
 & 左はし & まん中 & 右はし \\
$x$ & $-r$ & 0 & $r$ \\
$\theta$ & $-\frac{\pi}{2}$ & 0 & $\frac{\pi}{2}$ \\
\end{bktable}
\end{bkcenter}
（$r\sin\bigl(-\frac{\pi}{2}\bigr)=-r$、$r\sin\frac{\pi}{2}=r$ で確かめられます。）

\bkvpmidashi{計算する}
\begin{align*}
S&=\int_{-\frac{\pi}{2}}^{\frac{\pi}{2}}\underbrace{2r\cos\theta}_{2\sqrt{r^2-x^2}}\cdot\underbrace{r\cos\theta\,\dd\theta}_{\dd x}=2r^2\int_{-\frac{\pi}{2}}^{\frac{\pi}{2}}\cos^2\theta\,\dd\theta=2r^2\int_{-\frac{\pi}{2}}^{\frac{\pi}{2}}\frac{1+\cos2\theta}{2}\,\dd\theta\\
&=2r^2\Bigl[\,\frac{\theta}{2}+\frac{\sin2\theta}{4}\,\Bigr]_{-\frac{\pi}{2}}^{\frac{\pi}{2}}=2r^2\Bigl\{\Bigl(\frac{\pi}{4}+\frac{\sin\pi}{4}\Bigr)-\Bigl(-\frac{\pi}{4}+\frac{\sin(-\pi)}{4}\Bigr)\Bigr\}\\
&=2r^2\Bigl(\frac{\pi}{4}+\frac{\pi}{4}\Bigr)=2r^2\times\frac{\pi}{2}=\pi r^2
\end{align*}
（$\cos^2\theta=\dfrac{1+\cos2\theta}{2}$ は5-15。$\sin\pi=\sin(-\pi)=0$。）小学校から使ってきた「円の面積＝半径×半径×円周率」が、積分で確かめられました。

\bksubsection{7}{積の微分を、逆に読む：部分積分}

\label{bk:5-16:bubun}こんどは、積の微分を逆に読みます。まず、数の式で見てみましょう。$x\,e^x$ を積の微分（5-6）で微分すると、
\[ (x\,e^x)'=\underbrace{(x)'e^x}_{1\cdot e^x}+\underbrace{x\,(e^x)'}_{x\,e^x}=e^x+x\,e^x \]
これを逆に読むと、「$e^x+x\,e^x$ を積分すると $x\,e^x$」。2つのたし算を1つずつ積分すると、
\[ \int e^x\,\dd x+\int x\,e^x\,\dd x=x\,e^x \]
（$+C$ は最後にまとめてつけます。）$\int e^x\,\dd x$ を右へ移すと（$\int e^x\,\dd x=e^x$ は5-14）、
\[ \int x\,e^x\,\dd x=x\,e^x-\int e^x\,\dd x=x\,e^x-e^x+C \]
表になかった $\int x\,e^x\,\dd x$ が求められました。微分して確かめると、
\[ (xe^x-e^x)'=(e^x+xe^x)-e^x=xe^x \]
で、合っています。

\bkvpmidashi{文字で書くと}
同じことを、$f$ と $g$ で書きます。積の微分 $(fg)'=f'g+fg'$ の両辺を積分すると、左は「微分したものを積分する」のでもとの $fg$ にもどります（5-11）。
\[ f\,g=\int f'g\,\dd x+\int f\,g'\,\dd x \]
$\int f'g\,\dd x$ を左へ移して、左右を入れかえると、
\[ \int f\,g'\,\dd x=f\,g-\int f'\,g\,\dd x \]
これを\textbf{部分積分}といいます。左の積分が求めにくくても、右の $\int f'g\,\dd x$ が簡単になれば、答えが出ます。上の例は、$f=x$、$g'=e^x$（$g=e^x$）としたものです。\bkidx{ぶぶんせきぶん}{部分積分}

\bkvpmidashi{使い方の手順}
\begin{enumerate}[label={\bkcircled[bkTheme]{\arabic*}}, leftmargin=2.4\zw, itemsep=1pt, topsep=2pt]
\item かけ算の2つのうち、どちらを $f$（微分するほう）、どちらを $g'$（積分するほう）にするか決める（⑧）。
\item $f'$（$f$ を微分したもの）と、$g$（$g'$ を積分したもの。$+C$ はつけない）を求める。
\item $fg-\int f'g\,\dd x$ に入れて、残った積分を計算する。最後に $+C$ をつける。
\end{enumerate}

\begin{bkfigure}
\begin{tikzpicture}[x=1mm, y=1mm]
  \fill[bkBlue!14] (0,0) -- (0,24) -- plot[domain=0:1, samples=30] ({50*\x}, {24+12*\x*\x}) -- (50,0) -- cycle;
  \fill[bkOrange!20] (0,24) -- plot[domain=0:1, samples=30] ({50*\x}, {24+12*\x*\x}) -- (0,36) -- cycle;
  \draw[bkInk, line width=1.1pt] plot[domain=0:1, samples=30] ({50*\x}, {24+12*\x*\x});
  \draw[bkInk, line width=0.8pt] (0,0) rectangle (50,36);
  \node[font=\small, text=bkBlue!70!black] at (25,12) {$\int f\,\dd g$};
  \node[font=\small, text=bkOrange!80!black] at (12,31) {$\int g\,\dd f$};
  \node[font=\sffamily\scriptsize, anchor=north] at (25,-1) {横：$g$};
  \node[font=\sffamily\scriptsize, anchor=east] at (-1,18) {たて：$f$};
  \node[font=\sffamily\small, anchor=west, align=left] at (56,18) {全体の長方形＝$f\,g$\\青＋オレンジ＝$f\,g$\\だから\enspace 青＝$f\,g-{}$オレンジ};
\end{tikzpicture}
\bkfigcaption{図で見る部分積分。横に $g$、たてに $f$ をとって、$f$ と $g$ が変わっていくようすを曲線でかく。青い部分は、はば $\dd g$・高さ $f$ のたての細い長方形を集めたもの。オレンジの部分は、はば $\dd f$・長さ $g$ の横の細い長方形を集めたもの。$g'\,\dd x=\dd g$、$f'\,\dd x=\dd f$ なので、部分積分は「青＝全体${}-{}$オレンジ」（はじめ $g=0$ の場合）。}
\end{bkfigure}

\bksubsection{8}{どちらを $f$ にするか}

\label{bk:5-16:erabu}部分積分では、\textbf{微分すると簡単になるもの}を $f$、\textbf{積分しやすいもの}を $g'$ にします。$\displaystyle\int x\,e^x\,\dd x$ で、逆に $f=e^x$、$g'=x$（$g=\frac{x^2}{2}$）とすると、
\[ \int x\,e^x\,\dd x=\frac{x^2}{2}e^x-\int\frac{x^2}{2}e^x\,\dd x \]
となり、残った積分に $x^2$ が出てきて、はじめよりむずかしくなります。$x$ は微分すると1になって消えるので、$f$ にするのがよいのです。よく出る形の選び方は、次の表のとおりです。
\begin{bkcenter}
\begin{tblr}{colspec={Q[l,m]Q[l,m]Q[l,m]}, row{1}={bg=bkTheme!12, font=\sffamily\small}, hline{1,2,Z}={0.6pt,bkGray}, hline{3-5}={0.3pt,bkGray}, colsep=6pt}
積分の形 & $f$（微分する） & $g'$（積分する） \\
$x\,e^{ax}$ & $x$ → 1 & $e^{ax}$ → $\frac{1}{a}e^{ax}$ \\
$x\sin x$、$x\cos x$ & $x$ → 1 & $\sin x$ → $-\cos x$、$\cos x$ → $\sin x$ \\
$x^2e^x$ & $x^2$ → $2x$（2回使う） & $e^x$ → $e^x$ \\
$\ln x$ & $\ln x$ → $\frac{1}{x}$ & 1 → $x$ \\
\end{tblr}
\end{bkcenter}

\bkvpmidashi{$\ln x$ の積分}
$\ln x$ は1つの関数だけで、かけ算に見えません。でも、$\ln x=\ln x\times1$ と見ると、部分積分が使えます。$\ln x$ は積分のしかたが分かりませんが、微分すると $\frac{1}{x}$ と簡単になるので、$f=\ln x$、$g'=1$（$g=x$）とします。
\[ \int\ln x\,\dd x=\underbrace{\ln x}_{f}\cdot\underbrace{x}_{g}-\int\underbrace{\frac{1}{x}}_{f'}\cdot\underbrace{x}_{g}\,\dd x=x\ln x-\int1\,\dd x=x\ln x-x+C \]
確かめ：$(x\ln x-x)'=\Bigl(1\cdot\ln x+x\cdot\dfrac{1}{x}\Bigr)-1=\ln x+1-1=\ln x$。合っています。

\bksubsection{9}{定積分での部分積分}

\label{bk:5-16:teisekibun}定積分では、$fg$ の部分にも、はんいをつけます。
\[ \int_a^bf\,g'\,\dd x=\Bigl[\,f\,g\,\Bigr]_a^b-\int_a^bf'\,g\,\dd x \]
$\Bigl[\,f\,g\,\Bigr]_a^b$ は、$f\,g$ の $b$ での値から $a$ での値をひいたものです（5-11）。

\bksubsection{10}{はてな？の答え}

\label{bk:5-16:kotae}\bkhatenaref の電気の量 $\displaystyle\int_0^32t\,e^{-t}\,\dd t$ を求めましょう。
\begin{enumerate}[label={\bkcircled[bkTheme]{\arabic*}}, leftmargin=2.4\zw, itemsep=2pt, topsep=2pt]
\item $2t$ は微分すると2と簡単になるので $f=2t$。$e^{-t}$ を $g'$ にします。
\item $f'=2$。$g$ は $e^{-t}$ を積分したもので、5-14の $\int e^{at}\,\dd t=\frac{1}{a}e^{at}$ で $a=-1$ とすると、$g=-e^{-t}$。
\item 部分積分に入れます。$f'g=2\times(-e^{-t})=-2e^{-t}$ なので、ひき算の $-$ と合わせて $+$ になります。
\begin{align*}
\int_0^32t\,e^{-t}\,\dd t&=\Bigl[\,2t\times(-e^{-t})\,\Bigr]_0^3-\int_0^3(-2e^{-t})\,\dd t\\
&=\Bigl[\,-2t\,e^{-t}\,\Bigr]_0^3+\int_0^32e^{-t}\,\dd t
\end{align*}
\item 前の部分：$t=3$ で $-6e^{-3}$、$t=0$ で $0$。ひいて $-6e^{-3}$。
\item 後ろの部分：$2e^{-t}$ の原始関数は $-2e^{-t}$ なので、
\[ \Bigl[\,-2e^{-t}\,\Bigr]_0^3=-2e^{-3}-(-2e^0)=-2e^{-3}+2 \]
\item 合わせて $-6e^{-3}-2e^{-3}+2=2-8e^{-3}$。
\end{enumerate}
$e^{-3}=0.0498$ なので、$2-8\times0.0498=2-0.398=1.602$。答えは約 \qty{1.60}{C} です。表計算の1.6017と合いました。

\bksubsection{11}{積分では、どの道具を使う？}

\label{bk:5-16:erabikata}積分の式を見たら、次の順に考えます。微分の章（5-6）の「どの道具を使う？」と同じように、上から順に問いかけます。
\begin{bkfigure}
\begin{tikzpicture}[x=1mm, y=1mm, q/.style={bk node, minimum width=50mm, minimum height=13mm, font=\small},
  r/.style={bk node, draw=bkOrange, fill=bkOrange!8, minimum width=46mm, minimum height=13mm, font=\small}]
  \node[q] (q1) at (25,0) {積分の表（5-11・5-14・5-15）に\\そのままある？};
  \node[r] (a1) at (97,0) {表を逆に読んで積分};
  \node[q] (q2) at (25,-22) {「内側の式」と\\「内側の微分」が入っている？};
  \node[r] (a2) at (97,-22) {置換積分（②〜⑥）};
  \node[q] (q3) at (25,-44) {$x$ や $\ln x$ と、$e^x$・$\sin$・$\cos$ の\\かけ算になっている？};
  \node[r] (a3) at (97,-44) {部分積分（⑦〜⑩）};
  \node[r, minimum width=50mm] (a4) at (25,-66) {式を変形してから考えなおす\\（2倍角の式など、5-15）};
  \draw[bk link] (q1) -- node[bk link label, above] {はい} (a1);
  \draw[bk link] (q1) -- node[bk link label, right] {いいえ} (q2);
  \draw[bk link] (q2) -- node[bk link label, above] {はい} (a2);
  \draw[bk link] (q2) -- node[bk link label, right] {いいえ} (q3);
  \draw[bk link] (q3) -- node[bk link label, above] {はい} (a3);
  \draw[bk link] (q3) -- node[bk link label, right] {いいえ} (a4);
\end{tikzpicture}
\bkfigcaption{積分の道具の選び方。どの道具を使っても、さいごは答えを微分して、もとの式にもどるか確かめる。}
\end{bkfigure}
たとえば $\displaystyle\int\sin^2t\,\dd t$ は、表にもなく、内側の微分もなく、かけ算の形でもないので、2倍角の式 $\sin^2t=\dfrac{1-\cos2t}{2}$ で形を変えてから積分しました（5-15）。

\begin{motto}[4]
\textbf{2回使うと、もとにもどる}\par
$I=\displaystyle\int e^x\sin x\,\dd x$ に部分積分を2回使います。1回めは $f=\sin x$、$g'=e^x$。2回めは $f=\cos x$、$g'=e^x$ です。
\begin{align*}
I&=e^x\sin x-\int e^x\cos x\,\dd x\\
&=e^x\sin x-\Bigl(e^x\cos x+\int e^x\sin x\,\dd x\Bigr)=e^x(\sin x-\cos x)-I
\end{align*}
（2行めのかっこの中は、$\int e^x\cos x\,\dd x=e^x\cos x-\int(-\sin x)e^x\,\dd x$ です。）右にも $I$ が出てきました。$I$ を左に移すと $2I=e^x(\sin x-\cos x)$、つまり $I=\dfrac{1}{2}e^x(\sin x-\cos x)+C$。ゆれながら止まる動き（減衰振動、6-7・6-12）の計算で使います。
\end{motto}

\begin{senmonsho}[大学の教科書ふうの書き方]
\begin{bkhonbun}
部分積分の公式 $\int u\,\dd v=uv-\int v\,\dd u$ を用いると、次式を得る。
\[ \int_0^\infty e^{-st}\frac{\dd f}{\dd t}\,\dd t=\Bigl[\,e^{-st}f(t)\,\Bigr]_0^\infty+s\int_0^\infty e^{-st}f(t)\,\dd t \]
\end{bkhonbun}
\bkyaku{$u$ と $v$ は、この章の $f$ と $g$ のこと。$\dd v=g'\,\dd x$。$\int_0^\infty$ は「0からどこまでも」積分すること。$e^{-st}f'(t)$ の積分が、$f(t)$ の積分の $s$ 倍に変わる。これが9-5のラプラス変換で「微分がかけ算に変わる」しくみ。}
\end{senmonsho}

\begin{taisetsu}
\def\tsl#1{\par\addvspace{3pt}\noindent\hangindent=5.2\zw\hangafter=1\makebox[5.2\zw][l]{\sffamily\bfseries\small\color{bkTheme}#1}}
\tsl{置換積分}中身 $u$ とその微分を見つけ、$\dd u=u'\,\dd x$ で書きかえる。定積分は、はんいも $u$ に。
\tsl{部分積分}$\displaystyle\int f\,g'\,\dd x=f\,g-\int f'\,g\,\dd x$\par
\noindent\hspace*{5.2\zw}$f$：微分で簡単になるもの、$g'$：積分しやすいもの
\tsl{確かめ}答えを微分して、積分の中身にもどるか確かめる。
\end{taisetsu}

\begin{kotoba}
\small
\bkword{置換積分}{積分の中の式を $u$ におきかえて、簡単な形にする方法。合成関数の微分の逆。}
\bkword{部分積分}{積の微分の逆。$\int fg'\,\dd x=fg-\int f'g\,\dd x$。}
\end{kotoba}

%% ---------------- 例題 ----------------
\bksection{reidai}{例題でたしかめよう}

\begin{reidai}[置換積分]
$\displaystyle\int3x^2(x^3+2)^4\,\dd x$ を求めましょう。\par
\kangaekata かっこの中 $x^3+2$ を微分すると $3x^2$ で、それが前にある。だから $u=x^3+2$ とおく（③）。\par
\shiki $u=x^3+2$ とおくと $\dfrac{\dd u}{\dd x}=3x^2$ より $\dd u=3x^2\,\dd x$\par
\hspace*{2.6\zw}$\displaystyle\int\underbrace{(x^3+2)^4}_{u^4}\,\underbrace{3x^2\,\dd x}_{\dd u}=\int u^4\,\dd u=\frac{u^5}{5}+C$\par
\hspace*{2.6\zw}$u$ を $x^3+2$ にもどす。確かめ：\par
\hspace*{4.6\zw}$\Bigl\{\dfrac{(x^3+2)^5}{5}\Bigr\}'=\dfrac{5(x^3+2)^4\times3x^2}{5}=3x^2(x^3+2)^4$\par
\kotae $\dfrac{(x^3+2)^5}{5}+C$
\end{reidai}
置換積分では、さいごに $u$ を $x$ の式にもどすのをわすれないようにしましょう。問題は $x$ の式なので、答えも $x$ の式で書きます。もどしたら、微分して積分の中身にもどるかを確かめると安心です。

\begin{jibun}
$\displaystyle\int2x(x^2-1)^5\,\dd x$ を求めましょう。
\bkjans{$u=x^2-1$、$\dd u=2x\,\dd x$。$\dfrac{(x^2-1)^6}{6}+C$}
\end{jibun}

\begin{reidai}[係数を調整する置換積分]
（1）$\displaystyle\int t\,e^{-t^2}\,\dd t$\quad（2）$\displaystyle\int\cos x\,\sin^3x\,\dd x$\quad（3）$\displaystyle\int\sin x\,\cos^3x\,\dd x$ を求めましょう。\par
\kangaekata（1）$u=-t^2$ とおき、数のずれを④のように調整。（2）$\cos x$ は $\sin x$ の微分なので $u=\sin x$。（3）$\cos x$ の微分は $-\sin x$ なので $u=\cos x$ とし、$-$ を調整。\par
\shiki（1）$\dd u=-2t\,\dd t$ より $t\,\dd t=-\dfrac{1}{2}\dd u$。$\displaystyle\int e^u\Bigl(-\frac{1}{2}\Bigr)\dd u=-\frac{1}{2}e^u+C$\par
\hspace*{2.6\zw}（2）$\dd u=\cos x\,\dd x$。$\displaystyle\int u^3\,\dd u=\frac{u^4}{4}+C$\par
\hspace*{2.6\zw}（3）$\dd u=-\sin x\,\dd x$ より $\sin x\,\dd x=-\dd u$。\par
\hspace*{2.6\zw}\phantom{（3）}$\displaystyle\int u^3(-1)\,\dd u=-\frac{u^4}{4}+C$\par
\kotae （1）$-\dfrac{1}{2}e^{-t^2}+C$\quad（2）$\dfrac{\sin^4x}{4}+C$\quad（3）$-\dfrac{\cos^4x}{4}+C$
\end{reidai}
（3）を微分して確かめると、
\[ \Bigl(-\frac{\cos^4x}{4}\Bigr)'=-\frac{4\cos^3x\times(-\sin x)}{4}=\sin x\cos^3x \]
$-$ が2つで $+$ になって、もとにもどります。

\begin{jibun}
（1）$\displaystyle\int x\cos(x^2)\,\dd x$\quad（2）$\displaystyle\int\sin x\,\cos^2x\,\dd x$ を求めましょう。
\bkjans{（1）\nobreak $u=x^2$、$x\,\dd x=\frac{1}{2}\dd u$。$\dfrac{1}{2}\sin(x^2)+C$\quad（2）\nobreak $u=\cos x$、$\sin x\,\dd x=-\dd u$。$-\dfrac{\cos^3x}{3}+C$}
\end{jibun}

\begin{reidai}[分母の微分が分子にある形]
（1）$\displaystyle\int\frac{e^x}{e^x+1}\,\dd x$\quad（2）$\displaystyle\int\frac{x}{x^2+4}\,\dd x$ を求めましょう。\par
\kangaekata 分子が分母の微分（またはその数倍）になっている。④の「よく出る形」で $u=$（分母）。\par
\shiki（1）$u=e^x+1$（いつも正）とおくと $\dd u=e^x\,\dd x$。$\displaystyle\int\frac{1}{u}\,\dd u=\ln u+C$\par
\hspace*{2.6\zw}（2）$u=x^2+4$（いつも正）とおくと $\dd u=2x\,\dd x$ より $x\,\dd x=\dfrac{1}{2}\dd u$\par
\hspace*{4.6\zw}$\displaystyle\int\frac{1}{u}\cdot\frac{1}{2}\,\dd u=\frac{1}{2}\ln u+C$\par
\kotae （1）$\ln(e^x+1)+C$\quad（2）$\dfrac{1}{2}\ln(x^2+4)+C$
\end{reidai}
「いつも正」をたしかめるのは、$\ln u$ は $u>0$ のときだけ考えるからです（5-14）。$e^x$ も $x^2$ も0以上なので、$e^x+1$ や $x^2+4$ はいつも正です。（2）のように、分子が分母の微分のちょうど何倍か（ここでは $\frac{1}{2}$ 倍）になっているときは、その数を前に出してから積分します。

\begin{jibun}
$\displaystyle\int\frac{4x}{x^2+3}\,\dd x$ を求めましょう。
\bkjans{$u=x^2+3$、$\dd u=2x\,\dd x$ より $4x\,\dd x=2\,\dd u$。$2\ln(x^2+3)+C$}
\end{jibun}

\begin{reidai}[定積分の置換]
（1）$\displaystyle\int_0^1x(x^2+1)^2\,\dd x$\quad（2）$\displaystyle\int_0^{\frac{\pi}{2}}\sin^2t\,\cos t\,\dd t$ を求めましょう。\par
\kangaekata $u$ におきかえ、はんいも $u$ の値に書きかえる（⑤）。\par
\shiki（1）$u=x^2+1$、$x\,\dd x=\dfrac{1}{2}\dd u$。$x:0\to1$ のとき $u:1\to2$\par\vspace{2pt}
\hspace*{4.6\zw}$\displaystyle\frac{1}{2}\int_1^2u^2\,\dd u=\frac{1}{2}\Bigl[\,\frac{u^3}{3}\,\Bigr]_1^2=\frac{1}{2}\times\Bigl(\frac{8}{3}-\frac{1}{3}\Bigr)=\frac{1}{2}\times\frac{7}{3}=\frac{7}{6}$\par
\hspace*{2.6\zw}（2）$u=\sin t$、$\dd u=\cos t\,\dd t$。$t:0\to\frac{\pi}{2}$ のとき $u:\sin0=0\to\sin\frac{\pi}{2}=1$\par\vspace{2pt}
\hspace*{4.6\zw}$\displaystyle\int_0^1u^2\,\dd u=\Bigl[\,\frac{u^3}{3}\,\Bigr]_0^1=\frac{1}{3}-0=\frac{1}{3}$\par
\kotae （1）$\dfrac{7}{6}$\quad（2）$\dfrac{1}{3}$
\end{reidai}
（2）では、$t$ が0から $\frac{\pi}{2}$ まで動くと、$u=\sin t$ は0から1まで動きます。はんいの書きかえは、下端と上端のそれぞれに、おいた式（ここでは $\sin t$）を計算するだけです。（1）は、$x$ の式にもどしてから計算しても同じです。
\[ \Bigl[\,\frac{(x^2+1)^3}{6}\,\Bigr]_0^1=\frac{8}{6}-\frac{1}{6}=\frac{7}{6} \]

\begin{jibun}
$\displaystyle\int_0^12x(x^2+1)^3\,\dd x$ と $\displaystyle\int_0^{\frac{\pi}{2}}\sin t\,\cos t\,\dd t$ を求めましょう。
\bkjans{$u=x^2+1$（$1\to2$）：$\dfrac{15}{4}$。$u=\sin t$（$0\to1$）：$\Bigl[\,\dfrac{u^2}{2}\,\Bigr]_0^1=\dfrac{1}{2}$}
\end{jibun}

\begin{reidai}[$x=r\sin\theta$ とおく置換]
$\displaystyle\int_0^2\sqrt{4-x^2}\,\dd x$ を求めましょう。また、それは何の面積ですか。\par
\kangaekata ⑥の円の面積と同じく、$r=2$ として $x=2\sin\theta$ とおき、$\sqrt{\phantom{x}}$ をはずす。\par
\shiki $x=2\sin\theta$ とおくと\par
\hspace*{4.6\zw}$\sqrt{4-x^2}=\sqrt{4-4\sin^2\theta}=\sqrt{4\cos^2\theta}=2\cos\theta$\par
\hspace*{2.6\zw}$\dd x=2\cos\theta\,\dd\theta$。$x:0\to2$ のとき $\theta:0\to\frac{\pi}{2}$（$2\sin\frac{\pi}{2}=2$）\par\vspace{3pt}
\hspace*{2.6\zw}$\displaystyle\int_0^{\frac{\pi}{2}}2\cos\theta\cdot2\cos\theta\,\dd\theta=4\int_0^{\frac{\pi}{2}}\frac{1+\cos2\theta}{2}\,\dd\theta$\par\vspace{3pt}
\hspace*{4.6\zw}$\displaystyle=4\Bigl[\,\frac{\theta}{2}+\frac{\sin2\theta}{4}\,\Bigr]_0^{\frac{\pi}{2}}=4\times\Bigl(\frac{\pi}{4}+0\Bigr)-0=\pi$\par
\kotae $\pi$（半径2の円の $\frac{1}{4}$ の面積。$\pi\times2^2\div4=\pi$ と合う）
\end{reidai}
\noindent\begin{minipage}[c]{0.3\linewidth}
\centering
\begin{tikzpicture}[x=1mm, y=1mm]
  \fill[bkBlue!15] (0,0) -- (20,0) arc (0:90:20) -- cycle;
  \draw[bkBlue, line width=1pt] (20,0) arc (0:90:20);
  \draw[bkGray, line width=0.5pt, -{Stealth[length=3pt]}] (-2,0) -- (25,0) node[anchor=west, font=\sffamily\scriptsize] {$x$};
  \draw[bkGray, line width=0.5pt, -{Stealth[length=3pt]}] (0,-2) -- (0,25) node[anchor=south, font=\sffamily\scriptsize] {$y$};
  \node[font=\sffamily\scriptsize, anchor=north] at (20,-1) {2};
  \node[font=\sffamily\scriptsize, anchor=east] at (-1,20) {2};
\end{tikzpicture}
\end{minipage}\hfill
\begin{minipage}[c]{0.66\linewidth}
$y=\sqrt{4-x^2}$ は、両辺を2乗すると $x^2+y^2=4$ で、半径2の円の上半分です。0から2までなので、左の図の青い部分（円の $\frac{1}{4}$）の面積を求めたことになります。
\end{minipage}

\begin{jibun}
$\displaystyle\int_0^3\sqrt{9-x^2}\,\dd x$ を求めましょう。
\bkjans{$x=3\sin\theta$、$\dd x=3\cos\theta\,\dd\theta$。$\displaystyle\int_0^{\frac{\pi}{2}}9\cos^2\theta\,\dd\theta=\dfrac{9\pi}{4}$（半径3の円の $\frac{1}{4}$）}
\end{jibun}

\begin{reidai}[部分積分]
（1）$\displaystyle\int x\cos x\,\dd x$\quad（2）$\displaystyle\int x\,e^{2x}\,\dd x$ を求めましょう。\par
\kangaekata $x$ とほかの関数の積なので、⑦の部分積分。$f$ は微分すると1になる $x$（⑧）。\par
\shiki（1）$f=x$、$g'=\cos x$ とする。$f'=1$、$g=\sin x$\par
\hspace*{4.6\zw}$\displaystyle\int x\cos x\,\dd x=x\sin x-\int1\cdot\sin x\,\dd x$\par
\hspace*{4.6\zw}$\displaystyle\phantom{\int x\cos x\,\dd x}=x\sin x-(-\cos x)+C$\par
\hspace*{2.6\zw}（2）$f=x$、$g'=e^{2x}$ とする。$f'=1$、$g=\dfrac{1}{2}e^{2x}$（5-14）\par
\hspace*{4.6\zw}$\displaystyle\int x\,e^{2x}\,\dd x=\frac{1}{2}x\,e^{2x}-\int\frac{1}{2}e^{2x}\,\dd x=\frac{1}{2}x\,e^{2x}-\frac{1}{4}e^{2x}+C$\par
\kotae （1）$x\sin x+\cos x+C$\quad（2）$\dfrac{1}{2}xe^{2x}-\dfrac{1}{4}e^{2x}+C$
\end{reidai}
（1）を微分して確かめます。$x\sin x$ は積の微分（5-6）で微分します。
\[ (x\sin x+\cos x)'=(1\cdot\sin x+x\cos x)-\sin x=x\cos x \]
もとの式にもどりました。

\begin{jibun}
（1）$\displaystyle\int x\,e^{-x}\,\dd x$\quad（2）$\displaystyle\int x\sin x\,\dd x$ を求めましょう。
\bkjans{（1）\nobreak $g=-e^{-x}$。$-xe^{-x}-e^{-x}+C$\quad（2）\nobreak $g=-\cos x$。$-x\cos x+\sin x+C$}
\end{jibun}

\begin{reidai}[$\ln x$ の積分]
$\displaystyle\int_1^e\ln x\,\dd x$ を求めましょう。\par
\kangaekata ⑧の $\ln x$ の積分 $x\ln x-x$ を使い、上端の値から下端の値をひく。\par
\shiki $\ln e=1$（$e^1=e$）、$\ln1=0$（$e^0=1$）より\par
\hspace*{2.6\zw}$x=e$ で $e\cdot1-e=0$、$x=1$ で $1\cdot0-1=-1$\par
\hspace*{2.6\zw}$\Bigl[\,x\ln x-x\,\Bigr]_1^e=0-(-1)=1$\par
\kotae 1
\end{reidai}

\begin{jibun}
$\displaystyle\int_1^2\ln x\,\dd x$ を求めましょう（$\ln2=0.693$）。
\bkjans{$\Bigl[\,x\ln x-x\,\Bigr]_1^2=(2\ln2-2)-(0-1)=2\ln2-1\approx0.386$}
\end{jibun}

\begin{reidai}[定積分での部分積分：流れた電気の量]
電流 $i=2t\,e^{-t}$（A）が流れました。はじめの3秒間に流れた電気の量を求めましょう（$e^{-3}=0.0498$）。\par
\kangaekata 電気の量は電流の積分（5-12）。積なので⑨の部分積分（⑩と同じ手順）。\par
\shiki $f=2t$、$g'=e^{-t}$。$f'=2$、$g=-e^{-t}$\par
\hspace*{2.6\zw}$\begin{aligned}[t]\int_0^32t\,e^{-t}\,\dd t&=\Bigl[\,-2te^{-t}\,\Bigr]_0^3+\int_0^32e^{-t}\,\dd t=-6e^{-3}+\Bigl[\,-2e^{-t}\,\Bigr]_0^3\\
&=-6e^{-3}-2e^{-3}+2=2-8e^{-3}=2-8\times0.0498\end{aligned}$\par
\kotae $2-8e^{-3}\approx\qty{1.60}{C}$
\end{reidai}

\begin{jibun}
$\displaystyle\int_0^{\pi}t\sin t\,\dd t$ を求めましょう。
\bkjans{$f=t$、$g=-\cos t$。$\Bigl[\,-t\cos t\,\Bigr]_0^{\pi}+\displaystyle\int_0^{\pi}\cos t\,\dd t=\pi+0=\pi$}
\end{jibun}

\begin{reidai}[部分積分を2回使う]
$\displaystyle\int x^2e^x\,\dd x$ を求めましょう。\par
\kangaekata ⑧の表の $x^2e^x$ の行。1回めで $x^2$ が $2x$ になり、2回めで定数になる。\par
\shiki 1回め：$f=x^2$、$g'=e^x$。$f'=2x$、$g=e^x$\par
\hspace*{4.6\zw}$\displaystyle\int x^2e^x\,\dd x=x^2e^x-\int2x\,e^x\,\dd x$\par
\hspace*{2.6\zw}2回め：⑦で求めた $\displaystyle\int x\,e^x\,\dd x=xe^x-e^x$ を使う\par
\hspace*{4.6\zw}$=x^2e^x-2(xe^x-e^x)+C=x^2e^x-2xe^x+2e^x+C$\par
\kotae $e^x(x^2-2x+2)+C$
\end{reidai}

\begin{jibun}
$\displaystyle\int_0^1x^2e^x\,\dd x$ を求めましょう（$e=2.718$）。
\bkjans{$\Bigl[\,e^x(x^2-2x+2)\,\Bigr]_0^1=e\times1-1\times2=e-2\approx0.718$}
\end{jibun}

%% ---------------- まちがえやすいところ・物理工学 ----------------
\bksection{tsumazuki}{まちがえやすいところをチェック}

\begin{tsumazuki}{定積分の置換で、はんいを書きかえない}
\bkng{$\displaystyle\int_0^12x(x^2+1)^3\,\dd x=\int_0^1u^3\,\dd u$\quad（はんいが $x$ のまま）}
\bkok{$\displaystyle\int_0^12x(x^2+1)^3\,\dd x=\int_1^2u^3\,\dd u$}
はんいは $x$ の値のままではなく、$u=x^2+1$ の値（1と2）に書きかえます。×のまま計算すると $\frac{1}{4}$ になり、正しい $\frac{15}{4}$ とずれます。
\end{tsumazuki}

\begin{tsumazuki}{$\dd x$ を $\dd u$ に書きかえわすれる}
$\displaystyle\int x\,e^{x^2}\,\dd x$ で、$e^{x^2}$ を $e^u$ にしただけで $\int e^u\,\dd x$ としてはいけません。$x\,\dd x=\frac{1}{2}\dd u$ まで書きかえて、積分の中を $u$ だけにします。
\end{tsumazuki}

\begin{tsumazuki}{部分積分の $-$ をわすれる}
\bkng{$\displaystyle\int x\cos x\,\dd x=x\sin x+\int\sin x\,\dd x$}
\bkok{$\displaystyle\int x\cos x\,\dd x=x\sin x-\int\sin x\,\dd x$}
部分積分は「$fg$ ひく $\int f'g$」です。答えを微分して、もとの式にもどるか確かめましょう。計算が進むほど式が長くなったら、$f$ と $g'$ の選び方も見直します（⑧）。
\end{tsumazuki}

\begin{butsuri}
\begin{itemize}[leftmargin=1.2\zw, itemsep=1pt, topsep=0pt, label={\color{bkPurple}●}]
\item \begin{minipage}[t]{0.53\linewidth}\raggedright
\par\addvspace{3pt}\textbf{流れた電気の量}\par\noindent スイッチを入れたあとの電流は、$t\,e^{-at}$ の形のことが多くあります。$i=2te^{-t}$ を、0から $t$ 秒まで部分積分すると、$t$ 秒までに流れた電気の量は\par\vspace{3pt}
{\centering $Q=2-2(t+1)e^{-t}$\par}\vspace{3pt}
となり、2Cに近づきます。
\end{minipage}\hfill
\begin{minipage}[t]{0.45\linewidth}
\begin{tikzpicture}[baseline={([yshift=-3.5mm]current bounding box.north)}]
\begin{axis}[bkplot, scale only axis, width=36mm, height=27mm, xmin=0, xmax=6, ymin=0, ymax=2.4, xtick={1,2,...,5}, ytick={1,2}, xlabel={$t$（秒）}]
  \addplot[bkcurve, domain=0:5.6, samples=80] {2*x*exp(-x)};
  \addplot[bkcurve2, domain=0:5.4, samples=80] {2-2*(x+1)*exp(-x)};
  \node[font=\sffamily\footnotesize, text=bkBlue, anchor=south west, fill=bkPurpleBg, inner sep=1pt] at (axis cs:3.0,0.45) {電流 $i$（A）};
  \node[font=\sffamily\footnotesize, text=bkOrange!80!black, anchor=north west, align=left, inner sep=1pt] at (axis cs:6.1,2.0) {流れた量\\[-2pt]$Q$（C）};
\end{axis}
\end{tikzpicture}
\end{minipage}\par\vspace{9pt}
\item \textbf{ラプラス変換（9-5）}\par\noindent 制御工学の中心の道具。$e^{-st}$ をかけて積分すると、微分が $s$ のかけ算に変わります。そのしくみが、部分積分です。
\item \textbf{フーリエ級数（9-1）}\par\noindent $\int t\sin nt\,\dd t$ のような積分を部分積分で計算して、音や振動を $\sin$ と $\cos$ に分けます。
\end{itemize}
\end{butsuri}

%% ---------------- 章末問題 ----------------
\bkendproblems
\bkproblemgroup{基本}{3}

\bkproblem 次の不定積分を求めましょう。\quad（1）$\displaystyle\int4x(x^2+5)^3\,\dd x$\quad（2）$\displaystyle\int x\,e^{3x^2}\,\dd x$\quad（3）$\displaystyle\int\cos x\,\sin^2x\,\dd x$
\bkpans{（1）\nobreak $\dfrac{(x^2+5)^4}{2}+C$\quad（2）\nobreak $\dfrac{1}{6}e^{3x^2}+C$\quad（3）\nobreak $\dfrac{\sin^3x}{3}+C$}

\bkproblem $\displaystyle\int\frac{2x+1}{x^2+x+1}\,\dd x$ を求めましょう（$x^2+x+1$ はいつも正）。
\bkpans{$u=x^2+x+1$。$\ln(x^2+x+1)+C$}

\bkproblem 次の定積分を求めましょう。\quad（1）$\displaystyle\int_0^2x(x^2+1)^2\,\dd x$\quad（2）$\displaystyle\int_0^1t\,e^{t^2}\,\dd t$
\bkpans{（1）\nobreak $\dfrac{1}{2}\Bigl[\,\dfrac{u^3}{3}\,\Bigr]_1^5=\dfrac{62}{3}$\quad（2）\nobreak $\dfrac{1}{2}\Bigl[\,e^u\,\Bigr]_0^1=\dfrac{e-1}{2}\approx0.859$}

\bkproblem 次の不定積分を求めましょう。\quad（1）$\displaystyle\int x\,e^{2x}\,\dd x$\quad（2）$\displaystyle\int3x\cos x\,\dd x$
\bkpans{（1）\nobreak $\dfrac{1}{2}xe^{2x}-\dfrac{1}{4}e^{2x}+C$\quad（2）\nobreak $3x\sin x+3\cos x+C$}

\bkproblem $\displaystyle\int_0^{\frac{\pi}{2}}x\cos x\,\dd x$ と $\displaystyle\int_1^e2\ln x\,\dd x$ を求めましょう。
\bkpans{$\Bigl[\,x\sin x+\cos x\,\Bigr]_0^{\frac{\pi}{2}}=\dfrac{\pi}{2}-1\approx0.571$。$2\times1=2$}

\bkproblem 次を求めましょう。\quad（1）$\displaystyle\int3x^2(x^3+1)^2\,\dd x$\quad（2）$\displaystyle\int x\sin2x\,\dd x$
\bkpans{（1）\nobreak $u=x^3+1$。$\dfrac{(x^3+1)^3}{3}+C$\quad（2）\nobreak $f=x$、$g=-\dfrac{1}{2}\cos2x$。$-\dfrac{1}{2}x\cos2x+\dfrac{1}{4}\sin2x+C$}

\bkproblem $\displaystyle\int_0^1x\,e^x\,\dd x$ を求めましょう。
\bkpans{$\Bigl[\,xe^x-e^x\,\Bigr]_0^1=(e-e)-(0-1)=1$}

\bkproblemgroup{チャレンジ}{4}

\bkproblem 電流 $i=4t\,e^{-2t}$（A）が流れました。はじめの1秒間に流れた電気の量を求めましょう（$e^{-2}=0.135$）。
\bkpans{$\Bigl[\,-2te^{-2t}\,\Bigr]_0^1+\displaystyle\int_0^12e^{-2t}\,\dd t=-2e^{-2}+(1-e^{-2})=1-3e^{-2}\approx\qty{0.59}{C}$}

\bkproblem $\displaystyle\int_0^1x^2e^{-x}\,\dd x$ を求めましょう（$e^{-1}=0.368$）。
\bkpans{$\displaystyle\int x^2e^{-x}\,\dd x=-e^{-x}(x^2+2x+2)+C$。$\Bigl[\,\cdots\,\Bigr]_0^1=2-5e^{-1}\approx0.160$}

\bkproblem $\displaystyle\int_0^{\pi}\sin^2t\,\cos t\,\dd t$ と $\displaystyle\int_0^{\frac{\pi}{2}}\sin^2t\,\cos t\,\dd t$ を求め、ちがいのわけを考えましょう。
\bkpans{$u=\sin t$。0から $\pi$ では $u:0\to0$ で0。0から $\frac{\pi}{2}$ では $u:0\to1$ で $\dfrac{1}{3}$。$\frac{\pi}{2}$ から $\pi$ では $\cos t<0$ で、面積が打ち消しあう。}

\bkproblem $\displaystyle\int_0^1\frac{x}{x^2+1}\,\dd x$ と $\displaystyle\int_0^{\frac{\pi}{2}}x\sin x\,\dd x$ を求めましょう（$\ln2=0.693$）。
\bkpans{$u=x^2+1$（$1\to2$）で $\dfrac{1}{2}\Bigl[\,\ln u\,\Bigr]_1^2=\dfrac{1}{2}\ln2\approx0.347$。$\Bigl[\,-x\cos x+\sin x\,\Bigr]_0^{\frac{\pi}{2}}=1$}

\bkproblem $\displaystyle\int_{-2}^{2}\sqrt{4-x^2}\,\dd x$ を、$x=2\sin\theta$ とおいて求めましょう。答えが、半径2の半円の面積と同じになることを確かめましょう。
\bkpans{$\displaystyle\int_{-\frac{\pi}{2}}^{\frac{\pi}{2}}4\cos^2\theta\,\dd\theta=4\times\dfrac{\pi}{2}=2\pi$。半円 $\dfrac{1}{2}\times\pi\times2^2=2\pi$}

\begin{furikaeri}
\begin{checklist}
\item 置換積分で、$\dd u$ への書きかえができる。\hfill{\footnotesize\sffamily\color{bkTheme}→\pageref{bk:5-16:chikan}ページ}
\item 定積分の置換で、はんいを書きかえられる。\hfill{\footnotesize\sffamily\color{bkTheme}→\pageref{bk:5-16:hanni}ページ}
\item 部分積分で、$f$ と $g'$ を選んで計算できる。\hfill{\footnotesize\sffamily\color{bkTheme}→\pageref{bk:5-16:bubun}ページ}
\end{checklist}
\end{furikaeri}

\bkfukidashi{delta}{積分の道具がそろった。次は、テイラー展開だ。}

%% 5-17 近くでは多項式で近似できる：テイラー展開（第11回：ページ制限なし。なぜその数になるかを書き足した）
\newcommand\bkvqmidashi[1]{\par\Needspace{6\baselineskip}\noindent\textbf{#1}\par}
\bkchapter{5}{17}{近くでは多項式で近似できる：\\テイラー展開}{\bklevel{4}}{近くでは多項式で近似できる：テイラー展開}
\bkmichishirube{5-8,5-14,5-15}{6-6,6-10,9-1}

\begin{goal}
\begin{itemize}[label={\color{bkTheme}●}, leftmargin=1.6\zw, itemsep=0pt, topsep=1pt]
\item \bklevel{3}\enspace 1次近似（接線）に2次の項をたすと、近似がよくなることを説明できる。
\item \bklevel{4}\enspace $e^x$、$\sin x$、$\cos x$ のテイラー展開を作り、近似値を計算できる。
\item \bklevel{4}\enspace $\sin\theta\approx\theta$ を使って、振り子の周期を求められる。
\end{itemize}
\end{goal}

\begin{hatena}
\noindent\begin{minipage}[c]{0.6\linewidth}
電卓で $\sin0.5$ を計算すると、0.4794255… が一瞬で出てきます。

でも、電卓やコンピュータの中の回路がもともとできるのは、たし算・ひき算・かけ算・わり算のような計算と、そのくり返しだけです。円もグラフも書けない計算機は、どうやって $\sin$ の値を出しているのでしょう。
\end{minipage}\hfill
\begin{minipage}[c]{0.36\linewidth}
\centering
\begin{tikzpicture}[x=1mm, y=1mm]
  \draw[bkInk!70, line width=1pt, rounded corners=3pt, fill=bkPaper] (0,0) rectangle (30,44);
  \draw[bkInk!50, line width=0.6pt, fill=white] (3,33) rectangle (27,41);
  \node[font=\ttfamily\small, anchor=east] at (26,37) {0.4794255};
  \foreach \i in {0,1,2,3} \foreach \j in {0,1,2} {\draw[bkInk!40, fill=bkTheme!10, rounded corners=1pt] (3+\i*6.5,3+\j*7) rectangle ++(5,5);}
  \foreach \i in {2,3} {\draw[bkInk!40, fill=bkTheme!10, rounded corners=1pt] (3+\i*6.5,24) rectangle ++(5,5);}
  \draw[bkInk!40, fill=bkTheme!10, rounded corners=1pt] (3,24) rectangle (14.5,29);
  \node[font=\sffamily\footnotesize] at (8.75,26.5) {sin};
  \node[font=\sffamily\bfseries\small] at (38,22) {？};
\end{tikzpicture}
\end{minipage}
\end{hatena}

答えのかぎは、「なめらかな関数は、ある点の近くでは、$x$ と $x^2$ と $x^3$ と……の、たし算の式（多項式）で表せる」という考えです。たし算とかけ算の式なら、計算機でも計算できます。

\clearpage
%% ---------------- やってみよう・思い出そう ----------------
\begin{yattemiyou}[表計算で、$e^{0.5}$ に近づけよう]
$e^{0.5}$ の正しい値は 1.6487212…。次の式に $x=0.5$ を入れて、項を1つずつたしていきます。$3!$ は $1\times2\times3=6$、$4!$ は $1\times2\times3\times4=24$、$5!$ は $1\times2\times3\times4\times5=120$ のことです（くわしくは③）。
\[ 1+x+\frac{x^2}{2}+\frac{x^3}{3!}+\frac{x^4}{4!}+\frac{x^5}{5!}+\cdots \]
\begin{bkcenter}
\bksheet{3}{28}{%
1 & たした項 & 項の値 & ここまでの合計 \\
2 & 1 & 1 & 1 \\
3 & $x$ & 0.5 & 1.5 \\
4 & $\frac{x^2}{2}$ & 0.125 & 1.625 \\
5 & $\frac{x^3}{6}$ & 0.020833 & 1.645833 \\
6 & $\frac{x^4}{24}$ & 0.002604 & 1.648438 \\
7 & $\frac{x^5}{120}$ & 0.000260 & 1.648698 \\}
\bktabcaption{項をたすたびに、正しい値 1.648721… に近づく。}
\end{bkcenter}
\end{yattemiyou}

\noindent たす項は、どんどん小さくなり、合計は正しい値にどんどん近づきます。かけ算とたし算だけで、$e^{0.5}$ が計算できました。この式はどこから来たのでしょう。なぜ2、6、24、120でわるのでしょう。

\begin{omoidasou}
\bkrecall{5-4}{$(x^n)'=nx^{n-1}$。定数の微分は0。定数倍は前に出せる。}
\bkrecall{5-7}{2回微分 $f''(x)$ は、グラフの曲がり方を表す。$f''>0$ なら下に凸（上に曲がる）。}
\bkrecall{5-8}{1次近似：$x$ が $a$ に近いとき $f(x)\approx f(a)+f'(a)(x-a)$（接線で近似）}
\bkrecall{5-14}{$(e^x)'=e^x$、$e^0=1$。$(x^r)'=rx^{r-1}$（$r$ が分数でもよい）。$(\ln x)'=\dfrac{1}{x}$}
\bkrecall{5-15}{$\sin$ と $\cos$ は、微分するたびに $\sin\to\cos\to-\sin\to-\cos\to\sin$ と回る。}
\end{omoidasou}

\clearpage
%% ---------------- 数学の道具 ----------------
\bksection{doguu}{数学の道具}

\bksubsection{1}{接線で近似する（1次近似）}

\label{bk:5-17:ichiji}5-8では、グラフを接線で近似しました。$x=0$ の近くでは、接線は「点 $(0,\ f(0))$ を通り、傾き $f'(0)$ の直線」なので、
\[ f(x)\approx f(0)+f'(0)\,x \]
$e^x$ では、$f(0)=e^0=1$。$(e^x)'=e^x$（5-14）なので $f'(0)=e^0=1$。だから $e^x\approx1+x$ です。どのくらい合うか、表で見てみましょう。
\begin{bkcenter}
\begin{bktable}{colspec={Q[l,m,wd=40mm] *{3}{Q[c,m,wd=20mm]}}, rowsep=1.5pt}
$x$ & 0.1 & 0.5 & 1 \\
正しい値 $e^x$ & 1.10517 & 1.64872 & 2.71828 \\
1次近似 $1+x$ & 1.1 & 1.5 & 2 \\
ずれ & 0.005 & 0.149 & 0.718 \\
\end{bktable}
\end{bkcenter}
$x=0.1$ ではよく合いますが、$x$ が0からはなれるほど、ずれが大きくなります。接線はまっすぐなのに、$e^x$ のグラフは上に曲がっている（下に凸、5-7）ので、だんだんはなれていってしまうのです。

\bksubsection{2}{曲がり方も合わせる（2次近似）}

\label{bk:5-17:niji}\bkidx{にじきんじ}{2次近似}そこで、まっすぐな直線のかわりに、曲がった放物線を使います。放物線 $p(x)=a_0+a_1x+a_2x^2$ には、自由に決められる数が3つ（$a_0$、$a_1$、$a_2$）あるので、$x=0$ で次の3つを $f(x)$ に合わせられます。
\begin{bkcenter}
\begin{bktable}{colspec={Q[l,m,wd=56mm] Q[l,m,wd=56mm]}, rowsep=1.5pt}
合わせるもの & 表すもの \\
値 $p(0)=f(0)$ & 同じ点を通る \\
1回微分 $p'(0)=f'(0)$ & 傾きが同じ（接線が同じ） \\
2回微分 $p''(0)=f''(0)$ & 曲がり方が同じ \\
\end{bktable}
\end{bkcenter}
$p(x)$ を微分していきます（5-4）。
\[ p(x)=a_0+a_1x+a_2x^2,\qquad p'(x)=a_1+2a_2x,\qquad p''(x)=2a_2 \]
それぞれに $x=0$ を入れると、$x$ のつく項は0になって、
\[ p(0)=a_0,\qquad p'(0)=a_1,\qquad p''(0)=2a_2 \]
これを $f(0)$、$f'(0)$、$f''(0)$ と等しくすると、$a_0=f(0)$、$a_1=f'(0)$、そして $2a_2=f''(0)$ から $a_2=\dfrac{f''(0)}{2}$。
\[ f(x)\approx f(0)+f'(0)\,x+\frac{f''(0)}{2}\,x^2 \]
$x^2$ の項で2でわるのは、$x^2$ を2回微分すると2が出てくるので、それを打ち消すためです。

$e^x$ は何回微分しても $e^x$ なので $f''(0)=e^0=1$。だから $e^x\approx1+x+\dfrac{x^2}{2}$ です。
\begin{bkcenter}
\begin{bktable}{colspec={Q[l,m,wd=40mm] *{3}{Q[c,m,wd=20mm]}}, rowsep=1.5pt}
$x$ & 0.1 & 0.5 & 1 \\
正しい値 $e^x$ & 1.10517 & 1.64872 & 2.71828 \\
2次近似 $1+x+\frac{x^2}{2}$ & 1.105 & 1.625 & 2.5 \\
ずれ & 0.0002 & 0.024 & 0.218 \\
\end{bktable}
\end{bkcenter}
1次近似より、ずれがずっと小さくなりました。たとえば $x=0.5$ では、ずれが0.149から0.024にへりました。

\begin{bkfigure}
\begin{tikzpicture}
\begin{axis}[bkplot, scale only axis, width=84mm, height=40mm, xmin=-2, xmax=2, ymin=-0.6, ymax=6.2, xtick={-2,-1,1,2}, ytick={2,4,6}, ylabel={$y$},
    axis x line=bottom, x axis line style={draw=none}, xticklabel style={yshift=-0.5mm},
    yticklabel style={fill=white, inner sep=1pt, outer sep=2.3pt}]
  \draw[bkInk, line width=0.8pt, -{Stealth[length=5pt]}] (axis cs:-2,0) -- (axis cs:2.12,0) node[anchor=west, font=\sffamily\small, text=bkInk] {$x$};
  \addplot[bkcurve, domain=-2:1.82, samples=80, line width=2.2pt] {exp(x)};
  \addplot[bkthin, domain=-1.6:2] {1+x};
  \addplot[bkcurve2, domain=-2:2, samples=60] {1+x+x^2/2};
  \addplot[bkcurve3, domain=-2:1.95, samples=60] {1+x+x^2/2+x^3/6};
  \node[font=\sffamily\footnotesize, text=bkBlue, anchor=east] at (axis cs:1.6,6.35) {$e^x$};
  \node[font=\sffamily\footnotesize, text=bkGreen!70!black, anchor=west, fill=white, inner sep=1pt] at (axis cs:2.05,6.1) {$1+x+\frac{x^2}{2}+\frac{x^3}{6}$};
  \node[font=\sffamily\footnotesize, text=bkOrange!80!black, anchor=west, fill=white, inner sep=1pt] at (axis cs:2.05,5) {$1+x+\frac{x^2}{2}$};
  \node[font=\sffamily\footnotesize, text=bkGray, anchor=west, fill=white, inner sep=1pt] at (axis cs:2.05,3) {$1+x$};
\end{axis}
\end{tikzpicture}
\bkfigcaption{$e^x$（太い青）と、その近似。項をふやすほど、$x=0$ から遠くまでよく合う。}
\end{bkfigure}

\bksubsection{3}{何回でも合わせる：階乗とテイラー展開}

\label{bk:5-17:kaijo}同じように $x^3$、$x^4$、… の項をたして、3回微分、4回微分、… まで合わせていきます。まず、$x^3$ を3回微分してみましょう。
\[ x^3\ \xrightarrow{\text{微分}}\ 3x^2\ \xrightarrow{\text{微分}}\ 3\times2x\ \xrightarrow{\text{微分}}\ 3\times2\times1=6 \]
$x^4$ を4回微分すると、$4\times3\times2\times1=24$。指数が前に出て、1ずつへりながら、かけ算がつながっていきます。一般に、$x^n$ を $n$ 回微分すると、1から $n$ までの整数をすべてかけた数になります。これを $n!$ と書き、「$n$ の\textbf{階乗}」と読みます。$0!=1$ と決めておきます。\bkidx{かいじょう}{階乗}
\begin{bkcenter}
\begin{tblr}{colspec={Q[c,m]*{7}{Q[c,m]}}, row{1}={bg=bkTheme!12, font=\sffamily\small}, hline{1,2,Z}={0.6pt,bkGray}, colsep=7pt}
$n$ & 0 & 1 & 2 & 3 & 4 & 5 & 6 \\
$n!$ & 1 & 1 & 2 & 6 & 24 & 120 & 720 \\
\end{tblr}
\end{bkcenter}
表の右へ1つ進むと、$n$ をかけることになります。たとえば $5!=4!\times5$ なので、$24\times5=120$ です。

\bkvqmidashi{3次の多項式で確かめる}
$p(x)=a_0+a_1x+a_2x^2+a_3x^3$ を3回微分し、$x=0$ を入れます。
\begin{bkcenter}
\begin{bktable}{colspec={Q[l,m,wd=12mm] Q[l,m,wd=62mm] Q[c,m,wd=26mm]}, rowsep=1.5pt}
 & 式 & $x=0$ を入れると \\
$p$ & $a_0+a_1x+a_2x^2+a_3x^3$ & $a_0$ \\
$p'$ & $a_1+2a_2x+3a_3x^2$ & $a_1$ \\
$p''$ & $2a_2+3\times2\,a_3x$ & $2a_2$ \\
$p'''$ & $3\times2\times1\,a_3$ & $6a_3=3!\,a_3$ \\
\end{bktable}
\end{bkcenter}
3回微分すると、$x^2$ までの項は消えて（定数の微分は0）、$x^3$ の項だけが $3!\,a_3$ として残ります。だから3回微分を $f'''(0)$ に合わせるには $a_3=\dfrac{f'''(0)}{3!}$ にします。もっと次数の高い項があっても、$x=0$ を入れたとき $x$ が残って0になるので、じゃまをしません。

同じように、$a_nx^n$ の項を $n$ 回微分すると $a_n\times n!$ なので、$n$ 回微分を合わせるには $a_n=\dfrac{f^{(n)}(0)}{n!}$ にします（$f^{(n)}$ は $f$ を $n$ 回微分したもの）。
\[ f(x)=f(0)+f'(0)\,x+\frac{f''(0)}{2!}\,x^2+\frac{f'''(0)}{3!}\,x^3+\cdots+\frac{f^{(n)}(0)}{n!}\,x^n+\cdots \]
これを $f(x)$ の\textbf{テイラー展開}（$x=0$ のまわりの展開）といいます。項をふやすほど、$x=0$ から遠くまで、よく合うようになります。\bkidx{ていらーてんかい}{テイラー展開}

\bksubsection{4}{$e^x$ のテイラー展開}

\label{bk:5-17:e}$e^x$ は何回微分しても $e^x$ で、$x=0$ ではいつも $e^0=1$。だから、③の式の $f(0)$、$f'(0)$、$f''(0)$、… が全部1になって、
\[ e^x=1+x+\frac{x^2}{2!}+\frac{x^3}{3!}+\frac{x^4}{4!}+\cdots \]
やってみようの式は、これでした。$x=1$ を入れると、$e^1=e$ が計算できます。
\begin{bkcenter}
\begin{bktable}{colspec={Q[l,m,wd=21mm] *{8}{Q[c,m,wd=12mm]}}, rowsep=1.5pt, colsep=1pt, cells={font=\footnotesize}}
たす項 & 1 & 1 & $\frac{1}{2}$ & $\frac{1}{6}$ & $\frac{1}{24}$ & $\frac{1}{120}$ & $\frac{1}{720}$ & $\frac{1}{5040}$ \\
ここまでの合計 & 1 & 2 & 2.5 & 2.6667 & 2.7083 & 2.7167 & 2.71806 & 2.71825 \\
\end{bktable}
\bktabcaption{$e=2.718281\cdots$ に近づいていく。}
\end{bkcenter}
5-14の $\Bigl(1+\dfrac{1}{n}\Bigr)^n$ では、$n=10000$ にしても2.71814…で、小数3けたまでしか合いませんでした。こちらは、わずか8項のたし算で、小数4けたまで合います。

\bksubsection{5}{$\sin x$ と $\cos x$ のテイラー展開}

\label{bk:5-17:sin}$\sin x$ を微分していくと $\cos x$、$-\sin x$、$-\cos x$、$\sin x$、… と回ります（5-15）。$\sin0=0$、$\cos0=1$ を使って、$x=0$ での値を表にします。
\begin{bkcenter}
\begin{tblr}{colspec={Q[l,m]*{6}{Q[c,m]}}, row{1}={bg=bkTheme!12, font=\sffamily\small}, hline{1,2,Z}={0.6pt,bkGray}, colsep=7pt}
 & $f$ & $f'$ & $f''$ & $f'''$ & $f^{(4)}$ & $f^{(5)}$ \\
$\sin x$ & 0 & 1 & 0 & $-1$ & 0 & 1 \\
$\cos x$ & 1 & 0 & $-1$ & 0 & 1 & 0 \\
\end{tblr}
\end{bkcenter}
たとえば $\sin x$ の3回微分は $-\cos x$ なので、$x=0$ で $-\cos0=-1$ です。③の式に入れると、0の所の項は消えるので、
\[ \sin x=x-\frac{x^3}{3!}+\frac{x^5}{5!}-\frac{x^7}{7!}+\cdots,\qquad \cos x=1-\frac{x^2}{2!}+\frac{x^4}{4!}-\frac{x^6}{6!}+\cdots \]
$\sin x$ は奇数の次数だけ、$\cos x$ は偶数の次数だけで、$+$ と $-$ が交代します。

\bkvqmidashi{確かめ}
$\sin x$ の展開を、項ごとに微分してみます。たとえば $\dfrac{x^3}{3!}$ を微分すると、
\[ \frac{3x^2}{3!}=\frac{3x^2}{1\times2\times3}=\frac{x^2}{1\times2}=\frac{x^2}{2!} \]
のように、階乗が1つへります。
\[ \Bigl(x-\frac{x^3}{3!}+\frac{x^5}{5!}-\cdots\Bigr)'=1-\frac{3x^2}{3!}+\frac{5x^4}{5!}-\cdots=1-\frac{x^2}{2!}+\frac{x^4}{4!}-\cdots \]
ちゃんと $\cos x$ の展開になります（$(\sin x)'=\cos x$ と合う）。\bkidx{じすう}{次数}

\bkvqmidashi{\bkhatenaref の答え}
$\sin x$ の展開のはじめの3項に $x=0.5$ を入れます。$0.5^3=0.125$、$0.5^5=0.03125$ です。
\begin{align*}
\sin0.5&\approx0.5-\frac{0.5^3}{3!}+\frac{0.5^5}{5!}=0.5-\frac{0.125}{6}+\frac{0.03125}{120}\\
&=0.5-0.0208333+0.0002604=0.4794271
\end{align*}
3項で、電卓の0.4794255と小数5けたまで合います。たし算・ひき算・かけ算・わり算だけで、$\sin$ の値が出せました。コンピュータの中の $\sin$ や $\exp$ も、このような多項式の考えで計算しています（実際には、計算を速く正確にするための工夫をさらに加えています。電卓には、たし算と「けたをずらす」計算だけで $\sin$ を求める別の方法を使うものもあります）。

\begin{bkfigure}
\begin{tikzpicture}
\begin{axis}[bkplot, scale only axis, width=84mm, height=38mm, xmin=-3.6, xmax=3.6, ymin=-2, ymax=2, xtick={-3.1416,-1.5708,1.5708,3.1416}, xticklabels={$-\pi$,$-\dfrac{\pi}{2}$,$\dfrac{\pi}{2}$,$\pi$}, ytick={-1,1},
    axis x line=bottom, x axis line style={draw=none}, xticklabel style={anchor=base, yshift=-5.5mm},
    yticklabel style={fill=white, inner sep=1pt, outer sep=2.3pt}]
  \draw[bkInk, line width=0.8pt, -{Stealth[length=5pt]}] (axis cs:-3.6,0) -- (axis cs:3.75,0) node[anchor=west, font=\sffamily\small, text=bkInk] {$x$};
  \addplot[bkcurve, domain=-3.6:3.6, samples=120, line width=2.2pt] {sin(deg(x))};
  \addplot[bkthin, domain=-2:2] {x};
  \addplot[bkcurve2, domain=-3.0:3.0, samples=80] {x-x^3/6};
  \addplot[bkcurve3, domain=-3.6:3.6, samples=80] {x-x^3/6+x^5/120};
  \node[font=\sffamily\footnotesize, text=bkBlue, anchor=east] at (axis cs:-3.65,0.44) {$\sin x$};
  \node[font=\sffamily\footnotesize, text=bkGray, anchor=west, fill=white, inner sep=1pt] at (axis cs:2.05,2) {$x$};
  \node[font=\sffamily\footnotesize, text=bkOrange!80!black, anchor=west, fill=white, inner sep=1pt] at (axis cs:3.07,-1.5) {$x-\frac{x^3}{6}$};
  \node[font=\sffamily\footnotesize, text=bkGreen!70!black, anchor=west, fill=white, inner sep=1pt] at (axis cs:3.67,0.86) {$x-\frac{x^3}{6}+\frac{x^5}{120}$};
\end{axis}
\end{tikzpicture}
\bkfigcaption{$\sin x$（太い青）と、そのテイラー展開を1項、2項、3項で止めたもの。3項なら、$-\frac{\pi}{2}$ から $\frac{\pi}{2}$ くらいまで、ほとんど重なる。}
\end{bkfigure}

\bksubsection{6}{小さい $x$ では、はじめの項だけでよい}

\label{bk:5-17:chiisai}$x$ が小さいと、$x^2$、$x^3$ はもっと小さくなります。たとえば $x=0.1$ なら $x^2=0.01$、$x^3=0.001$。そこで、はじめの1つか2つの項だけを使う近似がよく使われます。
\begin{bkcenter}
\begin{tblr}{colspec={Q[l,m]Q[l,m]Q[l,m]}, row{1}={bg=bkTheme!12, font=\sffamily\small}, hline{1,2,Z}={0.6pt,bkGray}, hline{3,4,5}={0.3pt,bkGray}, colsep=5.5pt}
近似の式 & 次の項（ずれの目安） & 使われる所 \\
$\sin\theta\approx\theta$ & $-\dfrac{\theta^3}{6}$ & 振り子、光の屈折、小さな傾き \\
$\cos\theta\approx1-\dfrac{\theta^2}{2}$ & $+\dfrac{\theta^4}{24}$ & 振り子の高さ、アンテナの向き \\
$e^x\approx1+x$ & $+\dfrac{x^2}{2}$ & 小さな利率、温度による変化 \\
$(1+x)^r\approx1+rx$ & $+\dfrac{r(r-1)}{2}x^2$ & 熱による伸び、小さなずれの計算 \\
\end{tblr}
\end{bkcenter}

\bkvqmidashi{$(1+x)^r$ の近似の作り方}
$f(x)=(1+x)^r$ を微分します。内側 $1+x$ の微分は1なので、合成関数の微分（5-6）と $(x^r)'=rx^{r-1}$（5-14）から、
\[ f'(x)=r(1+x)^{r-1},\qquad f''(x)=r(r-1)(1+x)^{r-2} \]
$x=0$ を入れると、$1+0=1$ は何乗しても1なので、$f(0)=1$、$f'(0)=r$、$f''(0)=r(r-1)$。②の式に入れると、
\[ (1+x)^r\approx1+rx+\frac{r(r-1)}{2}x^2 \]
はじめの2項が $1+rx$、3つめの項が「次の項」です。

\bkvqmidashi{$\sin\theta\approx\theta$ はどのくらい合う？}
表で確かめましょう（$\theta$ はラジアン。度に $\frac{\pi}{180}$ をかけて求める）。
\begin{bkcenter}
\begin{tblr}{colspec={Q[c,m]*{5}{Q[c,m]}}, row{1}={bg=bkTheme!12, font=\sffamily\small}, hline{1,2,Z}={0.6pt,bkGray}, colsep=7pt}
角（度） & 5° & 10° & 20° & 30° & 45° \\
$\theta$（rad） & 0.0873 & 0.1745 & 0.3491 & 0.5236 & 0.7854 \\
$\sin\theta$ & 0.0872 & 0.1736 & 0.3420 & 0.5 & 0.7071 \\
ずれ（\%） & 0.13 & 0.51 & 2.1 & 4.7 & 11 \\
\end{tblr}
\end{bkcenter}
ずれの割合は、次の項の大きさ $\dfrac{\theta^3}{6}$ を、近似の値 $\theta$ でわった $\dfrac{\theta^3}{6}\div\theta=\dfrac{\theta^2}{6}$ くらいです。10°（0.1745 rad）なら $\dfrac{0.1745^2}{6}=0.0051$ で0.5\,\%、20°で2\,\%。表と合います。「どこまで近似してよいか」は、次の項の大きさで見積もります。

\bksubsection{7}{$x=a$ のまわりの展開}

\label{bk:5-17:a}$\sqrt{4.1}$ のように、0から遠い所の値を知りたいときは、近くの計算しやすい点 $a$（ここでは4）のまわりで展開します。$a$ からのずれ $x-a$ を新しい小さな数と見て、$x-a$ の多項式で近似するのです。考え方は③と同じで、$x=a$ で値・傾き・曲がり方…を合わせると、
\[ f(x)=f(a)+f'(a)(x-a)+\frac{f''(a)}{2!}(x-a)^2+\frac{f'''(a)}{3!}(x-a)^3+\cdots \]
$a=0$ にすると、③の式にもどります。

\bkvqmidashi{例：$\sqrt{4.1}$ を求める}
$f(x)=\sqrt{x}=x^{\frac{1}{2}}$、$a=4$ とします（$\sqrt4=2$ とすぐ分かるから）。$x-a=4.1-4=0.1$ です。
\begin{align*}
f'(x)&=\frac{1}{2}x^{-\frac{1}{2}}=\frac{1}{2\sqrt{x}}, & f'(4)&=\frac{1}{2\times2}=\frac{1}{4}\\
f''(x)&=\frac{1}{2}\times\Bigl(-\frac{1}{2}\Bigr)x^{-\frac{3}{2}}=-\frac{1}{4\,x\sqrt{x}}, & f''(4)&=-\frac{1}{4\times4\times2}=-\frac{1}{32}
\end{align*}
\[ \sqrt{4.1}\approx2+\frac{1}{4}\times0.1-\frac{1}{32}\times\frac{0.1^2}{2}=2+0.025-0.00015625=2.02484 \]
電卓の値は2.0248457。2次までで、小数5けた合いました。

\Needspace{12\baselineskip}
\bksubsection{8}{振り子は、小さくゆれれば単振動}

\label{bk:5-17:furiko}長さ $L$ の糸につるした、質量 $m$ のおもりの振り子です。

\noindent\begin{minipage}[c]{0.4\linewidth}
\centering
\begin{tikzpicture}[x=1mm, y=1mm]
  \fill[bkGray!40] (-14,0) rectangle (14,2);
  \draw[bkGray, dashed, line width=0.5pt] (0,0) -- (0,-36.5);
  \def\th{25}
  \draw[bkInk, line width=0.9pt] (0,0) -- (-90+\th:36);
  \draw[bkInk, line width=0.5pt] (0,-9) arc (-90:-90+\th:9);
  \node[font=\sffamily\scriptsize] at (-90+\th/2:12) {$\theta$};
  \node[font=\sffamily\scriptsize, anchor=west, xshift=1mm] at (-90+\th:18) {$L$};
  \draw[bkBlue!50, dashed, line width=0.6pt] (-90-\th:36) arc (-90-\th:-90+\th:36);
  \fill[bkOrange] (-90+\th:36) circle (2.6);
  \draw[bk arrow, bkInk, line width=1pt] (-90+\th:36) -- ++(0,-14) node[anchor=north, font=\sffamily\scriptsize] {$mg$};
  \draw[bk arrow, bkRed, line width=1.2pt] (-90+\th:36) -- ++(180+\th:8);
  \node[anchor=north east, font=\sffamily\scriptsize, text=bkRed, inner sep=0pt] at (9,-38.2) {$mg\sin\theta$};
\end{tikzpicture}
\end{minipage}\hfill
\begin{minipage}[c]{0.56\linewidth}
おもりにはたらく重力は $mg$（質量×重力加速度。$g=\qty{9.8}{m/s^2}$）で、真下を向いています。この力を、糸の向きと、円の弧にそった向きの2つに分けます（4-6の「斜面の重力を2つに分ける」と同じ）。

糸と真下の向きのあいだの角が $\theta$ なので、弧にそった向きの成分は $mg\sin\theta$（赤い矢印）です。糸の向きの成分は、糸が引く力といっしょに円の中心の向きにはたらくだけで、おもりを弧にそって進めたりもどしたりはしません。弧にそってふれを引きもどすのは、この $mg\sin\theta$ だけです。
\end{minipage}\par\vspace{3pt}
弧にそってはかった位置は $s=L\theta$（弧の長さ＝半径×角、4-2）です。$L$ は定数なので、2回微分すると、弧にそった加速度は $\ddot{s}=L\ddot{\theta}$。力＝質量×加速度（5-5）から、
\[ mL\ddot{\theta}=-mg\sin\theta \]
（$-$ は、ふれと反対向きに引きもどすから。）両辺を $mL$ でわると、
\[ \ddot{\theta}=-\frac{g}{L}\sin\theta \]
質量 $m$ は両辺から消えました。このままでは、5-15の単振動の形 $\ddot{x}=-\omega^2x$（右が $x$ そのものの何倍か）ではありません。でも、ふれが小さければ $\sin\theta\approx\theta$（⑥）なので、
\[ \ddot{\theta}\approx-\frac{g}{L}\,\theta \]
これは $\omega^2=\dfrac{g}{L}$、つまり $\omega=\sqrt{\dfrac{g}{L}}$ の単振動です。周期は5-15の $T=\dfrac{2\pi}{\omega}$ から、
\[ T=\frac{2\pi}{\omega}=2\pi\div\sqrt{\frac{g}{L}}=2\pi\sqrt{\frac{L}{g}} \]
（$\sqrt{\frac{g}{L}}$ でわることは、$\sqrt{\frac{L}{g}}$ をかけることと同じ。）この式には、ふれはば $\theta$ も質量 $m$ も入っていません。\textbf{小さいふれの周期は、ふれはばにもおもりの質量にもよらず、糸の長さだけで決まります}。2-8で「振り子の周期は糸の長さの $\sqrt{\phantom{x}}$ に比例する」と表で調べたことの理由が、ここで分かりました。たとえば $L=\qty{1}{m}$ なら $T=2\times3.14\times\sqrt{\dfrac{1}{9.8}}\approx\qty{2.0}{s}$ です。\bkidx{ふりこのしゅうき}{振り子の周期}

ふれが大きくなると $\sin\theta\approx\theta$ がずれてきて、周期は少し長くなります。
\begin{bkfigure}
\begin{tikzpicture}
\begin{axis}[bkplot, width=110mm, height=34mm, xmin=0, xmax=95, ymin=0.98, ymax=1.22, xtick={0,10,20,30,45,60,90}, xticklabels={0°,10°,20°,30°,45°,60°,90°}, ytick={1,1.1,1.2}, yticklabels={1倍,1.1倍,1.2倍}, xlabel={ふれはば}]
  \addplot[bkcurve, mark=*, mark size=1.4pt] coordinates {(0,1) (10,1.0019) (20,1.0077) (30,1.0174) (40,1.0313) (50,1.0498) (60,1.0732) (70,1.1021) (80,1.1375) (90,1.1803)};
  % ラベルは、横の格子の線（1.1倍・1.2倍）のあいだに置く
  \node[font=\sffamily\footnotesize, anchor=west, fill=white, inner sep=1pt] at (axis cs:3,1.15) {10°では0.2\,\%長いだけ};
  \node[font=\sffamily\footnotesize, anchor=east, fill=white, inner sep=1pt] at (axis cs:76,1.172) {90°では18\,\%長い};
\end{axis}
\end{tikzpicture}
\bkfigcaption{ふれはばと周期（小さいふれの周期を1倍とし、$\sin\theta$ のまま計算した値）。}
\end{bkfigure}

\Needspace{12\baselineskip}
\begin{motto}[5]
\textbf{オイラーの公式への入口}\par
$e^x$ の展開の $x$ に、「2乗すると $-1$ になる数」$j$（6-8）をかけた $j\theta$ を入れます。$j^2=-1$、$j^3=j^2\times j=-j$、$j^4=(j^2)^2=1$ を使って、$j$ のつかない項とつく項に分けると、
\[ e^{j\theta}=\Bigl(1-\frac{\theta^2}{2!}+\frac{\theta^4}{4!}-\cdots\Bigr)+j\Bigl(\theta-\frac{\theta^3}{3!}+\frac{\theta^5}{5!}-\cdots\Bigr)=\cos\theta+j\sin\theta \]
かっこの中は、⑤の $\cos\theta$ と $\sin\theta$ の展開そのものです。指数関数と三角関数が、1つの式でつながります。これが6-10のオイラーの公式で、交流回路の計算の土台になります。
\end{motto}

\begin{senmonsho}[大学の教科書ふうの書き方]
\begin{bkhonbun}
$f(x)$ の $x=a$ のまわりのテイラー展開は $\displaystyle f(x)=\sum_{n=0}^{\infty}\frac{f^{(n)}(a)}{n!}(x-a)^n$ である。特に $|\theta|\ll1$ のとき $\sin\theta=\theta+O(\theta^3)$ であり、振り子の運動方程式は線形化される。\bkidx{せんけいか}{線形化}
\end{bkhonbun}
\bkyaku{$\sum$ は5-10のたし算の記号。$\infty$ は「どこまでも」。$|\theta|\ll1$ は「$\theta$ が1よりずっと小さい」。$O(\theta^3)$（オーダー）は「$\theta^3$ くらいの小さな量」。「線形化」は、$\sin\theta$ を $\theta$ のような1次の式におきかえて、解きやすくすること。}
\end{senmonsho}

\bkkigocards{\bkkigocard{$n!$}{エヌのかいじょう}{1から $n$ までの整数をすべてかけた数。$0!=1$}{factorial}{5-17}}
{\bkkigocard{$f^{(n)}(x)$}{エフ・エヌかいびぶん}{$f(x)$ を $n$ 回微分したもの}{$n$th derivative}{5-17}}

\begin{taisetsu}
$\displaystyle f(x)=f(0)+f'(0)\,x+\frac{f''(0)}{2!}x^2+\frac{f'''(0)}{3!}x^3+\cdots$\\
$\displaystyle e^x=1+x+\frac{x^2}{2!}+\cdots$\quad $\displaystyle\sin x=x-\frac{x^3}{3!}+\cdots$\quad $\displaystyle\cos x=1-\frac{x^2}{2!}+\cdots$\\
小さい $\theta$ では $\sin\theta\approx\theta$。振り子の周期 $T=2\pi\sqrt{L/g}$
\end{taisetsu}

\begin{kotoba}
\small
\bkword{階乗}{1から $n$ までの整数をすべてかけた数 $n!$。$x^n$ を $n$ 回微分すると $n!$。}
\bkword{テイラー展開}{関数を、ある点の近くで多項式（たし算とかけ算の式）で表すこと。}
\bkword{線形化}{$\sin\theta\approx\theta$ のように、1次の式で近似して、式を解きやすくすること。}
\end{kotoba}

%% ---------------- 例題 ----------------
\bksection{reidai}{例題でたしかめよう}

\begin{reidai}[階乗の計算]
次の値を求めましょう。\quad（1）$5!$\quad（2）$\dfrac{6!}{4!}$\quad（3）$\dfrac{1}{4!}$ を小数で\par
\kangaekata $n!$ は1から $n$ までの整数をかけたもの（③）。わり算は、同じかけ算を約分する。\par
\shiki（1）$1\times2\times3\times4\times5=120$\par
\hspace*{2.6\zw}（2）$\dfrac{6!}{4!}=\dfrac{1\times2\times3\times4\times5\times6}{1\times2\times3\times4}=5\times6=30$\par
\hspace*{2.6\zw}（3）$4!=24$ なので、$\dfrac{1}{24}=0.041666\cdots$\par
\kotae （1）120\quad（2）30\quad（3）0.041666…
\end{reidai}

\begin{jibun}
次の値を求めましょう。\quad（1）$6!$\quad（2）$\dfrac{7!}{5!}$\quad（3）$\dfrac{1}{6!}$ を小数で
\bkjans{（1）\nobreak 720\quad（2）\nobreak $6\times7=42$\quad（3）\nobreak 0.0013888…}
\end{jibun}

\begin{reidai}[2次近似を作る]
次の関数の、$x=0$ のまわりの2次近似を求めましょう。\quad（1）$e^{2x}$\quad（2）$(1+x)^{10}$（これを使って $1.01^{10}$ も）\quad（3）$\sqrt{1+x}$\par
\kangaekata $x=0$ での $f$、$f'$、$f''$ の値を求め、②の式 $f(0)+f'(0)x+\frac{f''(0)}{2}x^2$ に入れる。\par
\shiki（1）$f'=2e^{2x}$、$f''=4e^{2x}$ より $f(0)=1$、$f'(0)=2$、$f''(0)=4$。$1+2x+\dfrac{4}{2}x^2$\par
\hspace*{2.6\zw}（2）$f'=10(1+x)^9$、$f''=90(1+x)^8$ より $f(0)=1$、$f'(0)=10$、$f''(0)=90$\par
\hspace*{4.6\zw}$1+10x+\dfrac{90}{2}x^2$。$x=0.01$ で $1+0.1+0.0045$（電卓1.10462）\par
\hspace*{2.6\zw}（3）$(1+x)^{\frac{1}{2}}$ と見る。$f'=\dfrac{1}{2}(1+x)^{-\frac{1}{2}}$、$f''=\dfrac{1}{2}\times\Bigl(-\dfrac{1}{2}\Bigr)(1+x)^{-\frac{3}{2}}$\par
\hspace*{4.6\zw}$f(0)=1$、$f'(0)=\dfrac{1}{2}$、$f''(0)=-\dfrac{1}{4}$。$1+\dfrac{1}{2}x+\dfrac{-\frac{1}{4}}{2}x^2$\par
\kotae （1）$1+2x+2x^2$\quad（2）$1+10x+45x^2$、$1.01^{10}\approx1.1045$\quad（3）$1+\dfrac{x}{2}-\dfrac{x^2}{8}$
\end{reidai}
（3）で $x=0.1$ とすると、$\sqrt{1.1}\approx1+0.05-0.00125=1.04875$。電卓の値1.04881とよく合います。

\begin{jibun}
$e^{-x}$ と $(1+x)^{\frac{1}{3}}$ の、$x=0$ のまわりの2次近似を求めましょう。
\bkjans{$e^{-x}\approx1-x+\dfrac{x^2}{2}$。$(1+x)^{\frac{1}{3}}$：$f'(0)=\frac{1}{3}$、$f''(0)=\frac{1}{3}\times\bigl(-\frac{2}{3}\bigr)=-\frac{2}{9}$ より \mbox{$1+\dfrac{x}{3}-\dfrac{x^2}{9}$}}
\end{jibun}

\begin{reidai}[$\sin$・$\cos$ の値を計算する]
$\sin0.2$ をテイラー展開の2項で、$\cos0.2$ を3項で計算しましょう。\par
\kangaekata ⑤の $\sin x$、$\cos x$ の展開の、はじめの項に $x=0.2$ を入れる。$0.2^2=0.04$、$0.2^3=0.008$、$0.2^4=0.0016$。\par
\shiki $\sin0.2\approx0.2-\dfrac{0.2^3}{3!}=0.2-\dfrac{0.008}{6}$\par
\hspace*{4.6\zw}$=0.2-0.001333=0.198667$\par
\hspace*{2.6\zw}$\cos0.2\approx1-\dfrac{0.2^2}{2!}+\dfrac{0.2^4}{4!}=1-\dfrac{0.04}{2}+\dfrac{0.0016}{24}$\par
\hspace*{4.6\zw}$=1-0.02+0.0000667=0.980067$\par
\kotae $\sin0.2\approx0.198667$（電卓0.198669）、$\cos0.2\approx0.980067$（電卓0.980067）
\end{reidai}

\begin{jibun}
$\sin0.3$ を2項で、$\cos0.3$ を3項（$1-\frac{x^2}{2!}+\frac{x^4}{4!}$）で計算しましょう。
\bkjans{$\sin0.3\approx0.3-0.0045=0.2955$（電卓0.29552）。$\cos0.3\approx1-0.045+0.0003375=0.955338$（電卓0.955336）}
\end{jibun}

\begin{reidai}[$e^x$ の値を計算する]
$e^{0.1}$ をテイラー展開の3項で、$e^{-0.1}$ を4項で計算しましょう。\par
\kangaekata ④の $e^x$ の展開に $x=0.1$、$x=-0.1$ を入れる。$x$ がマイナスのときは、奇数乗の項がマイナスになる（$(-0.1)^3=-0.001$）。\par
\shiki $e^{0.1}\approx1+0.1+\dfrac{0.1^2}{2}=1+0.1+0.005=1.105$\par
\hspace*{2.6\zw}$e^{-0.1}\approx1+(-0.1)+\dfrac{(-0.1)^2}{2}+\dfrac{(-0.1)^3}{6}$\par
\hspace*{4.6\zw}$=1-0.1+0.005-0.000167=0.904833$\par
\kotae $e^{0.1}\approx1.105$（電卓1.10517）、$e^{-0.1}\approx0.904833$（電卓0.904837）
\end{reidai}

\begin{jibun}
$e^{-0.2}$ を、テイラー展開の4項で計算しましょう。
\bkjans{$1-0.2+0.02-0.001333=0.818667$（電卓0.818731）}
\end{jibun}

\begin{reidai}[振り子の周期]
長さ\qty{1}{m}の振り子の、小さなゆれの周期を求めましょう（$g=\qty{9.8}{m/s^2}$、$\pi=3.14$）。\par
\kangaekata ⑧の周期の式 $T=2\pi\sqrt{\frac{L}{g}}$ に、$L=1$、$g=9.8$ を入れる。\par
\shiki $\dfrac{L}{g}=\dfrac{1}{9.8}=0.102$、$\sqrt{0.102}=0.319$\par
\hspace*{2.6\zw}$T=2\times3.14\times0.319=2.00$\par
\kotae 約2.0秒
\end{reidai}

\begin{jibun}
周期1秒の振り子の、糸の長さは何cmでしょう（$g=\qty{9.8}{m/s^2}$、$\pi^2=9.87$）。
\bkjans{$T=2\pi\sqrt{\dfrac{L}{g}}$ の両辺を2乗して $L=g\Bigl(\dfrac{T}{2\pi}\Bigr)^2=\dfrac{9.8\times1^2}{4\times9.87}\approx\qty{0.248}{m}$。約25\,cm}
\end{jibun}

\begin{reidai}[近似のずれを見積もる]
15°（0.2618 rad）での $\sin\theta\approx\theta$ のずれの割合を、次の項から見積もりましょう。\par
\kangaekata ずれの割合は、次の項 $\frac{\theta^3}{6}$ を $\theta$ でわった $\frac{\theta^2}{6}$（⑥）。\par
\shiki ずれの割合 $\approx\dfrac{\theta^2}{6}=\dfrac{0.2618^2}{6}=\dfrac{0.0685}{6}=0.0114$\par
\kotae 約1.1\,\%（電卓で計算した $\sin15^\circ=0.2588$ とくらべても1.2\,\%）
\end{reidai}

\begin{jibun}
$\cos\theta\approx1$ と近似したときのずれ（次の項 $\frac{\theta^2}{2}$）が0.01以下になるのは、$\theta$ が何radまでですか。
\bkjans{$\dfrac{\theta^2}{2}\leqq0.01$ より $\theta^2\leqq0.02$、$\theta\leqq\sqrt{0.02}\approx0.14$\,rad（約8°）}
\end{jibun}

\begin{reidai}[$x=a$ のまわりの展開]
$\sqrt{9.2}$ を、$a=9$ のまわりの2次近似で求めましょう。\par
\kangaekata ⑦の展開で、$f(x)=\sqrt{x}$、$a=9$、$x-a=0.2$ とする。$\sqrt9=3$。\par
\shiki $f(9)=3$、$f'(9)=\dfrac{1}{2\sqrt9}=\dfrac{1}{6}$、$f''(9)=-\dfrac{1}{4\times9\times\sqrt9}=-\dfrac{1}{108}$\par
\hspace*{2.6\zw}$3+\dfrac{1}{6}\times0.2-\dfrac{1}{108}\times\dfrac{0.2^2}{2}=3+0.033333-0.000185$\par
\kotae 3.033148（電卓3.033150）
\end{reidai}

\begin{jibun}
$\sqrt{16.4}$ を、$a=16$ のまわりの1次近似と2次近似で求めましょう。
\bkjans{$f'(16)=\frac{1}{8}$、$f''(16)=-\frac{1}{256}$。1次：$4+0.05=4.05$。2次：$4.05-\dfrac{0.16}{512}=4.049688$（電卓4.049691）}
\end{jibun}

\begin{reidai}[展開に式を入れる]
（1）$\cos2x$\quad（2）$e^{-x^2}$ のテイラー展開を、0でない項を3つずつ書きましょう。\par
\kangaekata ⑤の $\cos u$、④の $e^u$ の展開の $u$ に、$2x$ や $-x^2$ をかっこごと入れる（2-13）。\par
\shiki（1）$\cos u=1-\dfrac{u^2}{2!}+\dfrac{u^4}{4!}-\cdots$ に $u=2x$。$\dfrac{(2x)^2}{2!}=\dfrac{4x^2}{2}=2x^2$、$\dfrac{(2x)^4}{4!}=\dfrac{16x^4}{24}=\dfrac{2}{3}x^4$\par
\hspace*{2.6\zw}（2）$e^u=1+u+\dfrac{u^2}{2!}+\cdots$ に $u=-x^2$。$(-x^2)^2=x^4$ なので $\dfrac{(-x^2)^2}{2}=\dfrac{x^4}{2}$\par
\kotae （1）$1-2x^2+\dfrac{2}{3}x^4-\cdots$\quad（2）$1-x^2+\dfrac{x^4}{2}-\cdots$
\end{reidai}

\begin{jibun}
$\sin3x$ と $e^{2x^2}$ のテイラー展開を、0でない項を3つずつ書きましょう。
\bkjans{$\sin3x=3x-\dfrac{9}{2}x^3+\dfrac{81}{40}x^5-\cdots$。$e^{2x^2}=1+2x^2+2x^4+\cdots$}
\end{jibun}

\begin{reidai}[$\ln(1+x)$ の展開]
$\ln(1+x)$ を3次の項まで展開し、$\ln1.1$ を計算しましょう。\par
\kangaekata 3回微分して $x=0$ を入れ、③の式に入れる（5-14の $(\ln x)'=\frac{1}{x}$ と、合成関数の微分）。\par
\shiki $f'=\dfrac{1}{1+x}=(1+x)^{-1}$、$f''=-(1+x)^{-2}$、$f'''=2(1+x)^{-3}$\par
\hspace*{2.6\zw}$x=0$ で $f=\ln1=0$、$f'=1$、$f''=-1$、$f'''=2$ なので\par
\hspace*{2.6\zw}$\ln(1+x)\approx0+x+\dfrac{-1}{2!}x^2+\dfrac{2}{3!}x^3=x-\dfrac{x^2}{2}+\dfrac{x^3}{3}$\par
\hspace*{2.6\zw}$x=0.1$ で $0.1-0.005+0.000333$\par
\kotae $x-\dfrac{x^2}{2}+\dfrac{x^3}{3}$、$\ln1.1\approx0.09533$（電卓0.09531）
\end{reidai}

\begin{jibun}
$\ln1.2$ を、$\ln(1+x)$ の展開の3項で計算しましょう。
\bkjans{$0.2-0.02+0.002667=0.18267$（電卓0.18232）}
\end{jibun}

\begin{reidai}[次の項までたして、近似をよくする]
$\sqrt{1.02}$ を、$(1+x)^r\approx1+rx$ と、次の項 $\dfrac{r(r-1)}{2}x^2$ までたした式で計算しましょう。\par
\kangaekata $\sqrt{1.02}=(1+0.02)^{\frac{1}{2}}$ と見て、⑥の表の式を使う。\par
\shiki $r=\dfrac{1}{2}$、$x=0.02$。1次は $1+\dfrac{1}{2}\times0.02=1.01$\par
\hspace*{2.6\zw}次の項は $\dfrac{\frac{1}{2}\times(-\frac{1}{2})}{2}\times0.02^2=-\dfrac{1}{8}\times0.0004=-0.00005$\par
\kotae 1次で1.01、次の項までで1.00995（電卓1.0099505）
\end{reidai}

\begin{jibun}
$1.01^{20}$ を、$1+rx$ と、次の項までたした式で計算しましょう。
\bkjans{$1+20\times0.01=1.2$。次の項 $\dfrac{20\times19}{2}\times0.0001=0.019$ で1.219（電卓1.22019）}
\end{jibun}

\begin{reidai}[近似を式に入れる：振り子の高さ]
長さ\qty{2}{m}の振り子が、$\theta=0.2$\,rad ふれました。おもりは、いちばん下から何cm上がっていますか。$\cos\theta\approx1-\dfrac{\theta^2}{2}$ を使って求めましょう。\par
\kangaekata 支点からおもりまでの高さの差は $L\cos\theta$（直角三角形）。だから、いちばん下からの高さは $h=L-L\cos\theta=L(1-\cos\theta)$。ここに近似の式を入れる。\par
\shiki $1-\cos\theta\approx1-\Bigl(1-\dfrac{\theta^2}{2}\Bigr)=\dfrac{\theta^2}{2}$ なので $h\approx\dfrac{L\theta^2}{2}$\par
\hspace*{2.6\zw}$h\approx\dfrac{2\times0.2^2}{2}=0.04$（m）\par
\kotae 約\qty{4}{cm}（電卓で $2(1-\cos0.2)=0.0399$）
\end{reidai}

\begin{jibun}
長さ\qty{0.8}{m}の振り子が0.3\,rad ふれたときの、いちばん下からの高さを、同じ近似で求めましょう。
\bkjans{$h\approx\dfrac{0.8\times0.3^2}{2}=\qty{0.036}{m}$。約\qty{3.6}{cm}（電卓0.0357）}
\end{jibun}

%% ---------------- まちがえやすいところ・物理工学 ----------------
\bksection{tsumazuki}{まちがえやすいところをチェック}

\begin{tsumazuki}{$\sin\theta\approx\theta$ に度の値を入れる}
\bkng{$\sin10^\circ\approx10$}\bkok{$\sin10^\circ\approx0.1745$（$10^\circ=0.1745$\,rad）}
テイラー展開は、5-15の $(\sin x)'=\cos x$ から作ったので、ラジアンでしか成り立ちません。
\end{tsumazuki}

\begin{tsumazuki}{階乗でわるのをわすれる}
\bkng{$\dfrac{f'''(0)}{3}x^3$}\bkok{$\dfrac{f'''(0)}{3!}x^3=\dfrac{f'''(0)}{6}x^3$}
$x^3$ を3回微分すると $3\times2\times1=6$ が出てくるので、6でわります。3でわるのではありません。$\dfrac{f''(0)}{2}x^2$ の2も、ほんとうは $2!$ です。
\end{tsumazuki}

\begin{tsumazuki}{$x$ が大きいのに、少ない項で近似する}
テイラー展開は「$x=a$ の近く」で合う式です。$\sin3\approx3$ は大まちがい（正しくは0.14）。遠い所では、項をたくさんとるか、近い点 $a$ のまわりで展開しなおします。
\end{tsumazuki}

\begin{column}{振り子時計とガリレオ}
ガリレオは、教会のランプのゆれを脈ではかり、ゆれが小さくなっても周期が変わらないことに気づいたといわれます。のちにホイヘンスが振り子時計を作り、時計はとても正確になりました。ただし周期が変わらないのは、ふれが小さいとき（$\sin\theta\approx\theta$）だけです（⑧の図）。
\end{column}

\begin{butsuri}
\begin{itemize}[leftmargin=1.2\zw, itemsep=1pt, topsep=0pt, label={\color{bkPurple}●}]
\item \textbf{振り子とばね}\par\noindent 振り子は、ふれが小さいときだけ単振動。ばねも、のびが小さいときだけ $F=kx$ です。「小さいときは1次で近似」が、物理の式の多くを支えています。
\item \textbf{トランジスタと「小信号」}\par\noindent ダイオードやトランジスタでは、電流が電圧の指数関数（$e^{V/V_T}$ の形。$V_T$ は室温で約\qty{0.026}{V}）で変わります。電圧がほんの少しだけ変わるときは、$e^x\approx1+x$ のように1次で近似して、増幅の計算をします。
\item \textbf{ロボットと制御（9-10）}\par\noindent 手のひらに立てたほうき（倒立振子）の式にも $\sin\theta$ が入ります。立った姿勢のまわりで $\sin\theta\approx\theta$ と1次の式にし（線形化）、第9巻の方法で制御します。
\end{itemize}
\end{butsuri}

%% ---------------- 章末問題 ----------------
\bkendproblems
\bkproblemgroup{基本}{3}

\bkproblem 次の値を求めましょう。\quad（1）$4!$\quad（2）$\dfrac{6!}{4!}$\quad（3）$\dfrac{1}{5!}$ を小数で
\bkpans{（1）\nobreak 24\quad（2）\nobreak 30\quad（3）\nobreak 0.008333…}

\bkproblem $e^{3x}$ と $\ln(1+x)$ の、$x=0$ のまわりの2次近似を求めましょう。
\bkpans{$e^{3x}\approx1+3x+\dfrac{9}{2}x^2$。$\ln(1+x)$：$f'=\frac{1}{1+x}$、$f''=-\frac{1}{(1+x)^2}$ より \mbox{$x-\dfrac{x^2}{2}$}}

\bkproblem テイラー展開を使って、次の値を小数4けたまで求めましょう。\newline（1）$\sin0.1$（2項）\quad（2）$e^{0.2}$（4項）
\bkpans{（1）\nobreak $0.1-0.000167=0.0998$\quad（2）\nobreak $1+0.2+0.02+0.001333=1.2213$}

\bkproblem 長さ\qty{2.45}{m}の振り子の周期を求めましょう（$g=\qty{9.8}{m/s^2}$）。
\bkpans{\mbox{$T=2\pi\sqrt{0.25}=\pi\approx3.1$\,秒}}

\bkproblem $\sqrt{25.5}$ を、$a=25$ のまわりの1次近似で求めましょう。
\bkpans{$5+\dfrac{1}{10}\times0.5=5.05$（電卓5.04975）}

\bkproblem $\ln1.05$ を、$\ln(1+x)$ の展開の2項で計算しましょう。
\bkpans{$0.05-0.00125=0.04875$（電卓0.04879）}

\bkproblem 長さ\qty{0.5}{m}の振り子が小さくゆれるときの周期を求めましょう。ただし、$g=\qty{9.8}{m/s^2}$、$\sqrt{0.051}=0.226$ とします。
\bkpans{\mbox{$2\pi\sqrt{\dfrac{0.5}{9.8}}=2\times3.14\times0.226\approx1.4$\,秒}}

\bkproblemgroup{チャレンジ}{4}

\bkproblem $\cos\theta\approx1-\dfrac{\theta^2}{2}$ を使って、長さ $L$ の振り子が角 $\theta$ だけふれたときの、いちばん下からの高さ $h=L(1-\cos\theta)$ を近似しましょう。$L=\qty{1}{m}$、$\theta=0.1$\,rad のとき、何cmですか。
\bkpans{$h\approx\dfrac{L\theta^2}{2}$。$\dfrac{1\times0.01}{2}=\qty{0.005}{m}=\qty{0.5}{cm}$}

\bkproblem $(1+x)^r\approx1+rx$ を使って、$1.02^{10}$ と $\sqrt{1.06}$ の近似値を求めましょう。また、次の項 $\dfrac{r(r-1)}{2}x^2$ までたすと、$1.02^{10}$ はいくつになりますか（電卓では1.21899）。
\bkpans{$1.02^{10}\approx1.2$、$\sqrt{1.06}\approx1.03$。次の項 $45\times0.0004=0.018$ をたすと1.218}

\bkproblem $\cos x$ のテイラー展開の各項を微分すると、$-\sin x$ のテイラー展開になることを確かめましょう。また、$e^x$ の展開を項ごとに微分すると、どうなりますか。
\bkpans{$\Bigl(1-\dfrac{x^2}{2!}+\dfrac{x^4}{4!}-\cdots\Bigr)'=-x+\dfrac{x^3}{3!}-\cdots=-\Bigl(x-\dfrac{x^3}{3!}+\cdots\Bigr)$（$\frac{n}{n!}=\frac{1}{(n-1)!}$）。$e^x$ は、微分してももとの展開と同じ（$(e^x)'=e^x$）}

\bkproblem $\sin1$ を、テイラー展開の2項と3項で計算し、電卓の値0.841471とくらべましょう。ずれが0.001より小さくなるのは、何項からですか。
\bkpans{2項：$1-\dfrac{1}{6}=0.8333$（ずれ0.008）。3項：$+\dfrac{1}{120}$ で0.84167（ずれ0.0002）。3項から}

\begin{furikaeri}
\begin{checklist}
\item 2次近似で、曲がり方まで合わせる意味を説明できる。\hfill{\footnotesize\sffamily\color{bkTheme}→\pageref{bk:5-17:niji}ページ}
\item $e^x$・$\sin x$・$\cos x$ の展開を書き、値を計算できる。\hfill{\footnotesize\sffamily\color{bkTheme}→\pageref{bk:5-17:sin}ページ}
\item $\sin\theta\approx\theta$ から、振り子の周期を求められる。\hfill{\footnotesize\sffamily\color{bkTheme}→\pageref{bk:5-17:furiko}ページ}
\end{checklist}
\end{furikaeri}

\bkfukidashi{sora}{第5巻の道具が、これで全部そろった。最後に、表にまとめておこう！}

%% 5-18 まとめ：微分したときのふるまい（第11回：ページ制限なし。表の1行ずつに「なぜ」と図を書き足した）
\newcommand\bkvromidashi[1]{\par\Needspace{6\baselineskip}\noindent\textbf{#1}\par}
\bkchapter{5}{18}{まとめ：微分したときのふるまい}{\bklevel{3}}{まとめ：微分したときのふるまい}
\bkmichishirube{5-14,5-15,5-17}{6-1,6-4,6-6}

\begin{goal}
\begin{itemize}[label={\color{bkTheme}●}, leftmargin=1.6\zw, itemsep=0pt, topsep=1pt]
\item \bklevel{3}\enspace 定数・$t^n$・$e^{at}$・$\sin\omega t$・$\cos\omega t$・$\ln t$ の微分と積分を、表にまとめて使える。
\item \bklevel{3}\enspace 「微分すると自分の $a$ 倍」「2回微分すると自分の $-\omega^2$ 倍」になる関数を答えられる。
\item \bklevel{3}\enspace はじめの状態から、式の中の数を決められる。
\end{itemize}
\end{goal}

\begin{hatena}
\noindent\begin{minipage}[c]{0.58\linewidth}
第6巻では、「変わる速さが、いまの状態で決まる」という式から、これからの動きを求めます。たとえば、
\begin{itemize}[leftmargin=1.2\zw, itemsep=0pt, topsep=2pt, label={\color{bkYellow!70!black}●}]
\item 微分すると、自分の $-2$ 倍になる関数は？
\item 2回微分すると、自分の $-9$ 倍になる関数は？
\end{itemize}
こんななぞなぞに答えるには、どんな関数を知っていればよいでしょう。
\end{minipage}\hfill
\begin{minipage}[c]{0.38\linewidth}
\centering
\begin{tikzpicture}[x=1mm, y=1mm]
  \foreach \i/\c in {0/bkBlue,1/bkOrange,2/bkPurple,3/bkTheme} {\draw[\c, line width=1pt, rounded corners=2pt, fill=\c!10] (\i*11,0) rectangle ++(9,13); \node[font=\sffamily\bfseries\small, text=\c] at (\i*11+4.5,6.5) {？};}
  \node[font=\sffamily\scriptsize, anchor=north] at (20,-2) {答えの候補のカード};
\end{tikzpicture}
\end{minipage}
\end{hatena}

\clearpage
第5巻で、微分と積分の道具と、$e^{at}$・$\sin$・$\cos$ という大切な関数がそろいました。この章では、それらを「微分したときのふるまい」で整理し、第6巻で使う「答えの候補リスト」を作ります。

%% ---------------- やってみよう・思い出そう ----------------
\begin{yattemiyou}[関数あてクイズ]
左のせつめいに合う関数を、右から選びましょう（ノートに線で結ぶ）。答えは微分して確かめます。\par\vspace{7pt}
\begin{bkcenter}
\begin{tblr}{colspec={Q[l,m]Q[c,m]Q[l,m]}, hline{1,Z}={0.6pt,bkGray}, rowsep=3pt, colsep=6pt}
（1）微分すると0 & & ア\enspace $e^{3t}$ \\
（2）微分すると、いつも2 & & イ\enspace $\sin2t$ \\
（3）微分すると、自分の3倍 & & ウ\enspace 5（定数） \\
（4）2回微分すると、自分の $-4$ 倍 & & エ\enspace $\ln t$ \\
（5）微分すると $\dfrac{1}{t}$ & & オ\enspace $2t+1$ \\
\end{tblr}
\end{bkcenter}
\end{yattemiyou}

\noindent 答えは（1）ウ、（2）オ、（3）ア、（4）イ、（5）エです。1つずつ微分して確かめます。
\begin{bkcenter}
\begin{bktable}{colspec={Q[l,m,wd=14mm] Q[l,m,wd=62mm] Q[l,m,wd=36mm]}, rowsep=1.5pt}
 & 微分して確かめる & 使ったきまり \\
（1）ウ & $(5)'=0$ & 定数の微分は0（5-4） \\
（2）オ & $(2t+1)'=2+0=2$ & $(t)'=1$（5-4） \\
（3）ア & $(e^{3t})'=3e^{3t}$。自分（$e^{3t}$）の3倍 & $(e^{at})'=ae^{at}$（5-14） \\
（4）イ & $(\sin2t)'=2\cos2t$、$(2\cos2t)'=-4\sin2t$ & 5-15 \\
（5）エ & $(\ln t)'=\dfrac{1}{t}$ & 5-14 \\
\end{bktable}
\end{bkcenter}
（4）は $\cos2t$ でも正解です（$(\cos2t)''=-4\cos2t$）。このように「せつめい（微分したときのふるまい）」から関数を当てることが、第6巻の微分方程式を解くことそのものです。

\begin{omoidasou}
\bkrecall{5-4}{$(t^n)'=nt^{n-1}$。定数の微分は0。定数倍は前に出せる。}
\bkrecall{5-13}{積分すると $+C$ がつく。$C$ は、はじめの状態（初期条件）で決まる。}
\bkrecall{5-14}{$(e^{at})'=ae^{at}$、$e^0=1$、$(\ln t)'=\dfrac{1}{t}$。$e^x=y$ なら $x=\ln y$}
\bkrecall{5-15}{$(\sin\omega t)'=\omega\cos\omega t$、$(\cos\omega t)'=-\omega\sin\omega t$。2回微分で $-\omega^2$ 倍}
\bkrecall{5-17}{テイラー展開。$e^t=1+t+\frac{t^2}{2!}+\frac{t^3}{3!}+\cdots$}
\end{omoidasou}

新しい計算はほとんどありません。これまでの章をふり返りながら、1つずつ確かめていきましょう。

%% ---------------- 数学の道具 ----------------
\bksection{doguu}{数学の道具}

\bksubsection{1}{微分と積分のふるまい表}

\label{bk:5-18:hyo}第5巻で出てきた関数を、1つの表にまとめます。$C$ は積分定数です（5-11）。
\begin{bkcenter}
\small
\begin{tblr}{colspec={Q[c,m]Q[c,m]Q[c,m]Q[c,m]Q[l,m]}, row{1}={bg=bkTheme!12, font=\sffamily\small}, hline{1,2,Z}={0.6pt,bkGray}, hline{3-7}={0.3pt,bkGray}, rowsep=2pt, colsep=5pt}
関数 $f(t)$ & 微分 $f'(t)$ & 2回微分 $f''(t)$ & 積分 $\int f\,\dd t$ & ふるまい \\
$c$（定数） & 0 & 0 & $ct+C$ & 変わらない \\
$t^n$ & $nt^{n-1}$ & $n(n-1)t^{n-2}$ & $\dfrac{t^{n+1}}{n+1}+C$ & 次数が1つ下がる \\
$e^{at}$ & $ae^{at}$ & $a^2e^{at}$ & $\dfrac{1}{a}e^{at}+C$ & 自分の $a$ 倍 \\
$\sin\omega t$ & $\omega\cos\omega t$ & $-\omega^2\sin\omega t$ & $-\dfrac{1}{\omega}\cos\omega t+C$ & $\frac{1}{4}$ 周期進む \\
$\cos\omega t$ & $-\omega\sin\omega t$ & $-\omega^2\cos\omega t$ & $\dfrac{1}{\omega}\sin\omega t+C$ & $\frac{1}{4}$ 周期進む \\
$\ln t$（$t>0$） & $\dfrac{1}{t}$ & $-\dfrac{1}{t^2}$ & $t\ln t-t+C$ & ゆっくりふえる \\
\end{tblr}
\end{bkcenter}
表の1行ずつ、どこで作ったかと、なぜそうなるかを短くふり返ります。
\begin{itemize}[leftmargin=1.2\zw, itemsep=2pt, topsep=2pt, label={\color{bkTheme}●}]
\item \textbf{定数}\quad 変わらない量は、変わる速さが0（5-4）。積分は、高さ $c$ の長方形の面積 $ct$（5-9）。
\item \textbf{$t^n$}\quad 微分は5-4の表から見つけた「指数を前に出して1へらす」。2回微分は、それをもう1回使ったもの（$nt^{n-1}$ の指数 $n-1$ を前に出す）。積分は逆に「指数を1ふやして、その数でわる」（5-11）。
\item \textbf{$e^{at}$}\quad 微分するたびに $a$ が1つかかる（5-14）ので、2回微分では $a\times a=a^2$。積分では逆に $a$ でわる。
\item \textbf{$\sin\omega t$・$\cos\omega t$}\quad 円運動の速度から作った（5-15）。微分するたびに $\omega$ がかかり、$\frac{1}{4}$ 周期ずつ進む。2回微分で $\omega^2$ がかかり、$-$ がつく。
\item \textbf{$\ln t$}\quad $e^t$ のグラフを折り返して微分した（5-14）。積分は部分積分で求めた（5-16）。2回微分は、$\dfrac{1}{t}=t^{-1}$ を $t^n$ の行で $n=-1$ として、$-1\times t^{-2}=-\dfrac{1}{t^2}$。
\end{itemize}
積や合成関数は、積の微分・合成関数の微分（5-6）、置換積分・部分積分（5-16）で、この表の形にもちこみます。

\begin{bkfigure}
\begin{tikzpicture}
\begin{scope}[x=1mm, y=1mm]
  \foreach \i/\ttl/\ax in {0/{$e^{0.4t}$：ふえる}/0, 1/{$e^{-t}$：へる}/0, 2/{$\sin t$：ゆれる}/10, 3/{$\ln t$：ゆっくりふえる}/6} {
    \begin{scope}[xshift=\i*33mm]
      \draw[bkGray, line width=0.6pt, -{Stealth[length=3pt]}] (0,0) -- (0,20);
      \draw[bkGray, line width=0.6pt, -{Stealth[length=3pt]}] (0,\ax) -- (27,\ax) node[anchor=west, font=\sffamily\scriptsize] {$t$};
      \node[font=\sffamily\footnotesize, anchor=north] at (13,-2) {\ttl};
    \end{scope}}
  \draw[bkBlue, line width=1.4pt] plot[domain=0:3.6, samples=30] ({7*\x}, {4*exp(0.4*\x)});
  \draw[bkBlue, line width=1.4pt, xshift=33mm] plot[domain=0:3.6, samples=30] ({7*\x}, {16*exp(-\x)});
  \draw[bkBlue, line width=1.4pt, xshift=66mm] plot[domain=0:6.283, samples=50] ({3.8*\x}, {10+7*sin(\x r)});
  \draw[bkBlue, line width=1.4pt, xshift=99mm] plot[domain=0.35:10, samples=50] ({2.4*\x}, {6+4.6*ln(\x)});
\end{scope}
\end{tikzpicture}
\bkfigcaption{4つの「ふるまい」。ふえる・へる・ゆれる・ゆっくりふえる。自然や機械の変化の多くは、この形か、その組み合わせになる。}
\end{bkfigure}

\bksubsection{2}{なぞなぞの答え：自分にもどる関数}

\label{bk:5-18:modoru}表の中で、微分しても「もとの形」にもどってくる関数が、とくに大切です。
\begin{bkcenter}
\small
\begin{tblr}{width=\linewidth, colspec={X[l,m]Q[c,m]Q[c,m]Q[l,m]}, row{1}={bg=bkTheme!12, font=\sffamily\small}, hline{1,2,Z}={0.6pt,bkGray}, hline{3-5}={0.3pt,bkGray}, rowsep=3pt, colsep=5pt}
なぞなぞ & 式 & 答えの関数 & 例 \\
微分すると自分の $a$ 倍 & $f'=af$ & $Ce^{at}$ & ふえる・へる（6-4） \\
2回微分で自分の $-\omega^2$ 倍 & $f''=-\omega^2f$ & $A\sin\omega t+B\cos\omega t$ & ゆれる（6-6） \\
2回微分で自分の $a^2$ 倍 & $f''=a^2f$ & $Ae^{at}+Be^{-at}$ & たおれる \\
微分すると定数 $v$ & $f'=v$ & $vt+C$ & 等速（6-3） \\
\end{tblr}
\end{bkcenter}
$C$、$A$、$B$ はどんな数でもかまいません。定数倍は、微分しても前にそのまま残るからです（5-4）。表の答えが正しいことを、1行ずつ微分して確かめます。
\begin{itemize}[leftmargin=1.2\zw, itemsep=2pt, topsep=2pt, label={\color{bkTheme}●}]
\item $f=Ce^{at}$ なら $f'=C\times ae^{at}=a\times(Ce^{at})=af$。
\item $f=Ae^{at}+Be^{-at}$ なら、
\[ f'=aAe^{at}-aBe^{-at},\qquad f''=a^2Ae^{at}+a^2Be^{-at}=a^2f \]
$e^{-at}$ は1回微分するごとに $-a$ がかかり、2回で $(-a)\times(-a)=a^2$ になるからです。
\item $f=vt+C$ なら $f'=v$。
\end{itemize}
（$A\sin\omega t+B\cos\omega t$ は、このあとの「2つをたしても答えになる」で確かめます。）「たおれる」は、手のひらに立てたほうきのように、少しかたむくと、かたむきがどんどん大きくなる動きです。$e^{at}$ の部分が、時間とともにどんどん大きくなります。

\bkvromidashi{\bkhatenaref の答え}
\begin{itemize}[leftmargin=1.2\zw, itemsep=2pt, topsep=2pt, label={\color{bkTheme}●}]
\item 「微分すると、自分の $-2$ 倍」は $f'=-2f$。表の1行めで $a=-2$ として、$f=Ce^{-2t}$。
\item 「2回微分すると、自分の $-9$ 倍」は $f''=-9f$。表の2行めとくらべると $\omega^2=9$。2乗して9になる正の数は3なので $\omega=3$。答えは $f=A\sin3t+B\cos3t$。
\end{itemize}
$\omega=9$ ではなく $\omega=3$ になることに注意しましょう。$\sin3t$ を2回微分すると、3が2回かかって $3\times3=9$ になるからです。

\bkvromidashi{2つをたしても答えになる}
$\sin3t$ と $\cos3t$ は、どちらも $f''=-9f$ をみたします。では、たしたもの $f=A\sin3t+B\cos3t$ はどうでしょう。2回微分すると、
\[ f''=A(-9\sin3t)+B(-9\cos3t)=-9(A\sin3t+B\cos3t)=-9f \]
で、たしたものも答えになります。$-9$ を前にくくり出せたのは、どちらの項にも同じ $-9$ がかかったからです。このように「答えどうしをたしても答え」になり、「答えを何倍しても答え」になる式を\textbf{線形}な式といいます。第6巻・第9巻の中心になる考えです。\bkidx{せんけい}{線形}

\noindent 線形でない式では、こうはなりません。たとえば $f'=f^2$（変わる速さが自分の2乗）では、$f$ が答えでも $2f$ は答えになりません（$2f$ を微分すると $2f'=2f^2$ ですが、$(2f)^2=4f^2$ なので合わない）。

\bksubsection{3}{数を決めるのは、はじめの状態}

\label{bk:5-18:shoki}$Ce^{at}$ の $C$ や、$A\sin\omega t+B\cos\omega t$ の $A$、$B$ は、なぞなぞだけでは決まりません。5-13と同じように、はじめの状態（初期条件）で決めます。

\begin{bkfigure}
\begin{tikzpicture}
\begin{axis}[bkplot, width=80mm, height=50mm, xmin=0, xmax=12.8, ymin=0, ymax=135, xtick={0,5,10}, ytick={40,80,120}, xlabel={$t$}, ylabel={$f$}]
  \addplot[bkthin, domain=0:12.5, samples=40] {120*exp(-0.1*x)};
  \addplot[bkcurve, domain=0:12.5, samples=40] {80*exp(-0.1*x)};
  \addplot[bkthin, domain=0:12.5, samples=40] {40*exp(-0.1*x)};
  \addplot[only marks, mark=*, mark size=1.8pt, bkOrange] coordinates {(0,80)};
  \draw[bkGray, line width=0.4pt] (axis cs:12.6,34.4) -- (axis cs:13.5,54);
  \draw[bkGray, line width=0.4pt] (axis cs:12.6,22.9) -- (axis cs:13.5,36);
  \draw[bkGray, line width=0.4pt] (axis cs:12.6,11.5) -- (axis cs:13.5,18);
  \node[font=\sffamily\scriptsize, anchor=west, text=bkGray] at (axis cs:13.6,54) {$C=120$};
  \node[font=\sffamily\scriptsize, anchor=west, text=bkBlue] at (axis cs:13.6,36) {$C=80$};
  \node[font=\sffamily\scriptsize, anchor=west, text=bkGray] at (axis cs:13.6,18) {$C=40$};
\end{axis}
\end{tikzpicture}
\bkfigcaption{3本の曲線は、どれも $f'=-0.1f$ をみたす $f=Ce^{-0.1t}$。$C$ がちがうと、スタートの高さがちがう。「$t=0$ で80」と決めると、オレンジの点を通る1本（青）だけが選ばれる。}
\end{bkfigure}

\bkvromidashi{例1}
$f'=-0.1f$、$f(0)=80$ のとき。答えの形は $f=Ce^{-0.1t}$ です。$t=0$ を入れると、$e^0=1$ なので、
\[ f(0)=Ce^{-0.1\times0}=Ce^0=C\times1=C\qquad\text{これが80なので}\quad C=80 \]
よって $f=80e^{-0.1t}$。$C$ は、ちょうど「はじめの値」になります。

\bkvromidashi{例2}
$\ddot{x}=-4x$、$x(0)=0.1$、$\dot{x}(0)=0$ のとき（ばねを0.1だけのばして、静かに手をはなした）。$-\omega^2=-4$ から $\omega=2$ なので、答えの形は $x=A\sin2t+B\cos2t$ です。
\begin{enumerate}[label={\bkcircled[bkTheme]{\arabic*}}, leftmargin=2.4\zw, itemsep=1pt, topsep=2pt]
\item はじめの位置：$t=0$ を入れる。$\sin0=0$、$\cos0=1$ なので、
\[ x(0)=A\times0+B\times1=B \]
これが0.1なので $B=0.1$。
\item はじめの速度：まず微分して $\dot{x}=2A\cos2t-2B\sin2t$。$t=0$ を入れると、
\[ \dot{x}(0)=2A\times1-2B\times0=2A \]
これが0なので $A=0$。
\end{enumerate}
だから $x=0.1\cos2t$ です。はしから手をはなすと、$\cos$ の形（山から始まる波）でゆれます。

\noindent 1回微分の式では数が1つ（$C$）、2回微分の式では数が2つ（$A$、$B$）。1回積分するごとに $C$ が1つ出てくる（5-13）ので、2回微分の式には、はじめの「位置」と「速度」の2つが必要です。この2つで、ばねの動きはただ1つに決まります。

\bksubsection{4}{ずらしても使える}

\label{bk:5-18:zurasu}お湯が冷めるとき、冷める速さは「部屋の温度との差」に比例します（6-4）。部屋が20\,℃なら、比例定数を $k$ として $T'=-k(T-20)$。このままでは表にない形です。

そこで、部屋の温度との差 $u=T-20$ を主役にします。$T=80$ なら $u=60$、$T=50$ なら $u=30$ です。
\begin{bkfigure}
\begin{tikzpicture}
\begin{axis}[bkplot, width=86mm, height=46mm, xmin=0, xmax=65, ymin=0, ymax=90, xtick={0,20,40,60}, ytick={20,50,80}, xlabel={$t$（分）}, ylabel={$T$}]
  \addplot[bkthin, dashed, domain=0:62] {20};
  \addplot[bkcurve, domain=0:62, samples=50] {20+60*exp(-0.05*x)};
  \draw[bkOrange, line width=1pt, <->] (axis cs:20,20.8) -- (axis cs:20,41.3);
  \node[font=\sffamily\scriptsize, anchor=east, text=bkOrange!80!black] at (axis cs:18.3,30) {差 $u$};
  \node[font=\sffamily\scriptsize, anchor=north east, text=bkGray] at (axis cs:62,16) {部屋 20\,℃};
\end{axis}
\end{tikzpicture}
\bkfigcaption{お湯の温度 $T$（青）と、部屋の温度との差 $u$（オレンジの矢印）。$T$ と $u$ はいつも20ちがうだけ。}
\end{bkfigure}
$T$ と $u$ は、いつも20ちがうだけなので、$u$ のグラフは $T$ のグラフを20だけ下にずらしたものです。ずらしても傾きは変わりません。式では、20は定数で微分すると0なので、$u'=T'-0=T'$ です。
$T'=-k(T-20)$ の左の $T'$ を $u'$ に、右の $T-20$ を $u$ にかえると、
\[ u'=-ku \]
これは表の「微分すると自分の $a$ 倍」（$a=-k$）なので、答えは $u=Ce^{-kt}$。$u=T-20$ をもどすと、
\[ T-20=Ce^{-kt},\qquad\text{つまり}\qquad T=20+Ce^{-kt} \]
「差」が $e^{-kt}$ でへり、温度は20\,℃に近づきます（20\,℃より下にはなりません）。5-16の置換と同じ「おきかえ」の考えです。

\bksubsection{5}{ふるまいをくらべる：どれがいちばん速くふえる？}

\label{bk:5-18:kurabe}$t$ が大きくなったときのふえ方をくらべます。
\begin{bkcenter}
\begin{tblr}{colspec={Q[c,m]*{7}{Q[r,m]}}, row{1}={bg=bkTheme!12, font=\sffamily\small}, hline{1,2,Z}={0.6pt,bkGray}, colsep=5pt, cells={font=\small}}
$t$ & 1 & 2 & 3 & 4 & 5 & 10 & 20 \\
$\ln t$ & 0 & 0.69 & 1.10 & 1.39 & 1.61 & 2.30 & 3.00 \\
$t^3$ & 1 & 8 & {\bfseries 27} & {\bfseries 64} & 125 & 1000 & 8000 \\
$e^t$ & 2.72 & 7.39 & 20.1 & 54.6 & {\bfseries 148} & {\bfseries 22026} & {\bfseries 485165195} \\
\end{tblr}
\end{bkcenter}
表をよく見ましょう。$t=3$、4では $t^3$ のほうが大きいのですが、$t=5$ で $e^t$ が追いこし、そのあとは、どんどん引きはなします。$t=20$ では、$e^t$ は $t^3$ の6万倍以上です。

なぜ、最後は $e^t$ が勝つのでしょう。5-17のテイラー展開を思い出します。
\[ e^t=1+t+\frac{t^2}{2!}+\frac{t^3}{3!}+\frac{t^4}{4!}+\frac{t^5}{5!}+\cdots \]
$e^t$ の中には、$t^3$ より次数の高い $\dfrac{t^4}{24}$ がふくまれています。$\dfrac{t^4}{24}$ は $t^3$ の $\dfrac{t}{24}$ 倍なので、$t$ が24より大きくなれば、この1つの項だけで $t^3$ をこえます。そのうえ $t^5$、$t^6$、… の項もあるので、いずれ必ず $e^t$ が勝つのです。逆に $\ln t$ はとてもゆっくりふえます。$t$ が10倍になっても、$\ln t$ は2.3ふえるだけです（$\ln10=2.30$）。$t$ が大きいとき、
\[ \ln t\ \ll\ t^n\ \ll\ e^t\qquad\text{（$\ll$ は「ずっと小さい」）} \]
同じ理由で、$t^3e^{-t}=\dfrac{t^3}{e^t}$ は、$t$ が大きくなると0に近づきます（$t=10$ で0.045、$t=20$ で0.00002）。「$e^{-t}$ がかかると、$t^n$ がかかっていても、やがて0に近づく」ことは、9-5のラプラス変換でも使います。

\begin{bkfigure}
\begin{tikzpicture}[x=1mm, y=1mm, nd/.style={draw=bkTheme, line width=0.8pt, rounded corners=3pt, fill=bkTheme!6, font=\small, align=center, inner sep=3pt, minimum height=13mm, text width=26mm}]
  \node[nd] (a) at (0,14) {平均の速さ\\（5-1）};
  \node[nd] (b) at (36,14) {微分\\（5-3〜5-8）};
  \node[nd] (c) at (72,14) {積分（微分の逆）\\（5-9〜5-13）};
  \node[nd, fill=bkOrange!12, draw=bkOrange] (d) at (108,14) {$e$・$\sin$・$\cos$\\（5-14〜5-17）};
  \node[nd, fill=bkPurple!10, draw=bkPurple, text width=108mm] (e) at (54,-8) {第6巻：微分方程式\enspace「変わり方のきまり」から未来を求める};
  \draw[bk arrow, bkTheme, line width=1pt, shorten <=1.5mm, shorten >=1.5mm] (a) -- (b);
  \draw[bk arrow, bkTheme, line width=1pt, shorten <=1.5mm, shorten >=1.5mm] (b) -- (c);
  \draw[bk arrow, bkTheme, line width=1pt, shorten <=1.5mm, shorten >=1.5mm] (c) -- (d);
  \draw[bk arrow, bkPurple, line width=1.2pt, shorten <=1.5mm, shorten >=1.5mm] (d.south) -- (e.north -| d.south);
  \draw[bk arrow, bkPurple, line width=1.2pt, shorten <=1.5mm, shorten >=1.5mm] (b.south) -- (e.north -| b.south);
\end{tikzpicture}%
\bkfigcaption{第5巻の地図。微分と積分の道具と、答えの候補の関数をもって、第6巻へ。}
\end{bkfigure}

\Needspace{12\baselineskip}
\bksubsection{6}{第5巻の計算のきまり 一覧}

\label{bk:5-18:kimari}第6巻からは、ここにあるきまりを、ことわりなしに使います。わからなくなったら、右の章にもどりましょう。
\begin{bkcenter}
\small
\begin{tblr}{width=\linewidth, colspec={Q[l,m]X[l,m]Q[c,m]}, row{1}={bg=bkTheme!12, font=\sffamily\small}, hline{1,2,Z}={0.6pt,bkGray}, hline{3-12}={0.3pt,bkGray}, rowsep=2.5pt, colsep=5pt, column{1}={font=\sffamily\small}}
きまり & 式 & 章 \\
微分の定義 & $f'(x)=\displaystyle\lim_{h\to0}\frac{f(x+h)-f(x)}{h}$（接線の傾き・瞬間の変わる速さ） & 5-3 \\
たし算・定数倍 & $(f+g)'=f'+g'$、$(cf)'=cf'$ & 5-4 \\
積・合成関数 & $(fg)'=f'g+fg'$、$\dfrac{\dd y}{\dd x}=\dfrac{\dd y}{\dd u}\cdot\dfrac{\dd u}{\dd x}$ & 5-6 \\
増減と最大・最小 & $f'>0$ でふえる、$f'<0$ でへる。$f'=0$ の所が山や谷の候補 & 5-7 \\
1次近似 & $f(a+h)\approx f(a)+f'(a)\,h$ & 5-8 \\
定積分 & 細かく分けた長方形の面積の和 $\displaystyle\int_a^bf\,\dd x=\lim_{\Delta x\to0}\sum f(x_k)\,\Delta x$ & 5-10 \\
微積分の基本定理 & $F'=f$ なら $\displaystyle\int_a^bf\,\dd x=F(b)-F(a)$ & 5-11 \\
初期条件 & 積分の $+C$ は、はじめの状態で決まる & 5-13 \\
置換・部分積分 & $\dd u=u'\,\dd x$ で書きかえ。$\displaystyle\int fg'\,\dd x=fg-\int f'g\,\dd x$ & 5-16 \\
テイラー展開 & $f(x)=f(0)+f'(0)x+\dfrac{f''(0)}{2!}x^2+\cdots$ & 5-17 \\
\end{tblr}
\end{bkcenter}

\begin{motto}[4]
\textbf{$f'=af$ の答えは $Ce^{at}$ だけ}\par
ほかに答えはないのでしょうか。$f'=af$ をみたす $f$ について、$g=f\,e^{-at}$ を積の微分（5-6）で微分すると、
\[ g'=f'e^{-at}+f\cdot(-a)e^{-at}=af\,e^{-at}-af\,e^{-at}=0 \]
（$f'=af$ を入れました。）微分して0になるのは定数なので、$g=C$。$g=fe^{-at}$ の両辺に $e^{at}$ をかけると、$e^{-at}\times e^{at}=e^0=1$ なので $f=Ce^{at}$。答えはこれだけだと分かります。
\end{motto}

\begin{senmonsho}[大学の教科書ふうの書き方]
\begin{bkhonbun}
微分方程式 $\dfrac{\dd y}{\dd t}=ay$ の一般解は $y=Ce^{at}$（$C$ は任意定数）であり、初期条件 $y(0)=y_0$ のもとで特殊解 $y=y_0e^{at}$ が一意に定まる。\bkidx{いっぱんかい}{一般解}\bkidx{とくしゅかい}{特殊解}
\end{bkhonbun}
\bkyaku{「一般解」は $C$ がついたままの答え、「任意定数」はどんな数でもよい定数 $C$、「特殊解」は $C$ を決めた答え。「一意に定まる」は「ただ1つに決まる」。}
\end{senmonsho}

\begin{taisetsu}
\def\tsl#1{\makebox[5.2\zw][l]{\sffamily\bfseries\small\color{bkTheme}#1}}
\tsl{答えの候補}$f'=af$ → $Ce^{at}$\qquad $f''=-\omega^2f$ → $A\sin\omega t+B\cos\omega t$\\[3pt]
\tsl{}$f'=v$ → $vt+C$\qquad $f''=a^2f$ → $Ae^{at}+Be^{-at}$\\[3pt]
\tsl{数の決め方}$C$、$A$、$B$ は、はじめの状態（$t=0$ の値）で決める。\\[3pt]
\tsl{ふえ方}$\ln t\ \ll\ t^n\ \ll\ e^t$
\end{taisetsu}

\begin{kotoba}
\small
\bkword{線形}{答えどうしをたしても、答えを何倍しても答えになる式の性質。}
\bkword{初期条件}{はじめ（$t=0$）の位置や速度。式の中の数を決める。}
\end{kotoba}

%% ---------------- 例題 ----------------
\bksection{reidai}{例題でたしかめよう}

\begin{reidai}[自分に比例してへる]
$f'=-3f$、$f(0)=5$ をみたす関数を求めましょう。また、$t=1$ での値を求めましょう（$e^{-3}=0.0498$）。\par
\kangaekata 「微分すると自分の $-3$ 倍」なので、②の表から答えの形は $Ce^{-3t}$。$C$ は $t=0$ の値で決める（③）。\par
\shiki 答えの形は $f=Ce^{-3t}$。$t=0$ で $f(0)=Ce^0=C=5$\par
\hspace*{2.6\zw}確かめ：$f'=5\times(-3)e^{-3t}=-3\times5e^{-3t}=-3f$\par
\hspace*{2.6\zw}$t=1$ で $5e^{-3}=5\times0.0498=0.249$\par
\kotae $f=5e^{-3t}$、$t=1$ で約0.25
\end{reidai}

\begin{jibun}
$f'=0.5f$、$f(0)=2$ をみたす関数を求めましょう。$t=2$ での値も求めましょう（$e=2.718$）。
\bkjans{$f=2e^{0.5t}$。$t=2$ で $2e^{1}=2e\approx5.44$}
\end{jibun}

\begin{reidai}[ゆれる動きを決める]
$\ddot{x}=-9x$、$x(0)=0.2$、$\dot{x}(0)=0$ をみたす $x$ を求めましょう。また、周期を求めましょう。\par
\kangaekata 「2回微分で自分の $-9$ 倍」なので $\omega^2=9$、$\omega=3$（②）。$A$、$B$ ははじめの位置と速度で決める（③）。周期は $\frac{2\pi}{\omega}$（5-15）。\par
\shiki 答えの形は $x=A\sin3t+B\cos3t$。\par
\hspace*{2.6\zw}$x(0)=A\times0+B\times1=B=0.2$\par
\hspace*{2.6\zw}$\dot{x}=3A\cos3t-3B\sin3t$ より $\dot{x}(0)=3A=0$、$A=0$\par
\hspace*{2.6\zw}周期 $\dfrac{2\pi}{3}=\dfrac{6.28}{3}\approx2.09$\par
\kotae $x=0.2\cos3t$、周期約2.1
\end{reidai}

\begin{jibun}
$\ddot{x}=-16x$、$x(0)=0$、$\dot{x}(0)=2$ をみたす $x$ を求めましょう。周期も求めましょう。
\bkjans{$x=A\sin4t+B\cos4t$。$B=0$、$4A=2$ より $x=\dfrac{1}{2}\sin4t$。周期 $\dfrac{2\pi}{4}\approx1.57$}
\end{jibun}

\begin{reidai}[位置も速度も0でないとき]
$\ddot{x}=-4x$、$x(0)=0.1$、$\dot{x}(0)=0.6$ をみたす $x$ を求めましょう。\par
\kangaekata $\omega^2=4$ より $\omega=2$。はじめの位置で $B$ を、はじめの速度で $A$ を決める（③の例2と同じ手順）。\par
\shiki 答えの形は $x=A\sin2t+B\cos2t$\par
\hspace*{2.6\zw}位置：$x(0)=A\times0+B\times1=B$。$B=0.1$\par
\hspace*{2.6\zw}速度：$\dot{x}=2A\cos2t-2B\sin2t$ より\par
\hspace*{2.6\zw}$\dot{x}(0)=2A\times1-2B\times0=2A$\par
\hspace*{4.6\zw}$2A=0.6$ より $A=0.3$\par
\kotae $x=0.3\sin2t+0.1\cos2t$
\end{reidai}
$\sin$ の項は速度で、$\cos$ の項は位置で決まります。$t=0$ で $\sin0=0$ なので、$\sin$ の項は位置に効かず、微分すると $\cos$ になるので速度に効くからです。

\begin{jibun}
$\ddot{x}=-9x$、$x(0)=0.2$、$\dot{x}(0)=0.3$ をみたす $x$ を求めましょう。
\bkjans{$x=A\sin3t+B\cos3t$。$B=0.2$、$3A=0.3$ より $A=0.1$。$x=0.1\sin3t+0.2\cos3t$}
\end{jibun}

\begin{reidai}[式をみたすか確かめる]
$y=3e^{2t}+e^{-2t}$ が $y''=4y$ をみたすことを確かめましょう。\par
\kangaekata ①の表で2回微分し、$y''$ が $y$ の4倍になるかを見る。\par
\shiki $y'=3\times2e^{2t}+(-2)e^{-2t}=6e^{2t}-2e^{-2t}$\par
\hspace*{2.6\zw}$y''=6\times2e^{2t}-2\times(-2)e^{-2t}=12e^{2t}+4e^{-2t}$\par
\hspace*{2.6\zw}$12e^{2t}+4e^{-2t}=4(3e^{2t}+e^{-2t})=4y$\par
\kotae $y''=4y$ をみたす
\end{reidai}

\begin{jibun}
$y=2\sin3t-\cos3t$ が $y''=-9y$ をみたすことを確かめましょう。
\bkjans{$y'=6\cos3t+3\sin3t$、$y''=-18\sin3t+9\cos3t=-9(2\sin3t-\cos3t)=-9y$}
\end{jibun}

\begin{reidai}[冷めるお湯]
部屋の温度は20\,℃です。お湯の温度 $T$ が、$T'=-0.05(T-20)$ にしたがって冷めます（$t$ は分）。はじめ80\,℃のとき、$T$ の式と20分後の温度を求めましょう。また、50\,℃になるのは何分後ですか（$e^{-1}=0.368$、$\ln2=0.693$）。\par
\kangaekata ④のように、部屋との温度差 $T-20$ が $e^{-0.05t}$ でへる。時刻は、両辺の自然対数をとって求める（5-14）。\par
\shiki $T=20+Ce^{-0.05t}$。$t=0$ で $20+C\times1=80$ より $C=60$\par
\hspace*{2.6\zw}20分後：$T=20+60e^{-0.05\times20}=20+60e^{-1}$\par
\hspace*{4.6\zw}$=20+60\times0.368=20+22.1$\par
\hspace*{2.6\zw}50\,℃：$20+60e^{-0.05t}=50$ より $60e^{-0.05t}=30$、$e^{-0.05t}=\dfrac{1}{2}$\par
\hspace*{2.6\zw}両辺の $\ln$ をとる：$-0.05t=\ln\dfrac{1}{2}=-\ln2$\par
\hspace*{4.6\zw}$t=\dfrac{\ln2}{0.05}=\dfrac{0.693}{0.05}=13.9$\par
\kotae $T=20+60e^{-0.05t}$、約42\,℃、約14分後
\end{reidai}
$\ln\dfrac{1}{2}=-\ln2$ は、$e^{-0.693}=\dfrac{1}{e^{0.693}}=\dfrac{1}{2}$ だからです（3-12の $\log_ab^r=r\log_ab$ で $r=-1$）。50\,℃は、はじめの差60\,℃が半分の30\,℃になった時刻です。

\begin{jibun}
部屋の温度25\,℃で、$T'=-0.1(T-25)$（$t$ は分）、はじめ90\,℃のお湯の、10分後の温度を求めましょう（$e^{-1}=0.368$）。
\bkjans{$T=25+65e^{-0.1t}$。$t=10$ で $25+65\times0.368\approx48.9\,℃$}
\end{jibun}

\begin{reidai}[表を使って微分・積分する]
$f(t)=3t^2+2e^{-t}+\sin2t$ を微分しましょう。また、$\displaystyle\int_0^{\pi}f(t)\,\dd t$ を求めましょう（$\pi^3=31.0$、$e^{-\pi}=0.043$）。\par
\kangaekata 項ごとに①の表を使う（定数倍は前に出す、5-4）。定積分は上端と下端の値の差（5-11）。\par
\shiki $f'=3\times2t+2\times(-e^{-t})+2\cos2t=6t-2e^{-t}+2\cos2t$\par
\hspace*{2.6\zw}原始関数：$3\times\dfrac{t^3}{3}+2\times(-e^{-t})-\dfrac{1}{2}\cos2t=t^3-2e^{-t}-\dfrac{1}{2}\cos2t$\par
\hspace*{2.6\zw}$\begin{aligned}[t]\int_0^{\pi}f(t)\,\dd t&=\Bigl(\pi^3-2e^{-\pi}-\frac{1}{2}\cos2\pi\Bigr)-\Bigl(0-2e^0-\frac{1}{2}\cos0\Bigr)\\
&=\Bigl(\pi^3-2e^{-\pi}-\frac{1}{2}\Bigr)-\Bigl(-2-\frac{1}{2}\Bigr)=\pi^3+2-2e^{-\pi}\end{aligned}$\par
\kotae $f'=6t-2e^{-t}+2\cos2t$、積分は約 $31.0+2-0.09=32.9$
\end{reidai}
$\sin2t$ の部分は、
\[ \Bigl(-\frac{1}{2}\cos2\pi\Bigr)-\Bigl(-\frac{1}{2}\cos0\Bigr)=-\frac{1}{2}+\frac{1}{2}=0 \]
で、答えに残りませんでした。0から $\pi$ までで $\sin2t$ はちょうど1周期なので、山と谷の面積が打ち消しあうからです（5-15）。このように、項ごとに分けて計算すると、どの項がどれだけ答えにきいているかも分かります。

\begin{jibun}
$g(t)=4t^3-e^{2t}+3\cos t$ を微分しましょう。また、$\displaystyle\int_0^1g(t)\,\dd t$ を求めましょう（$e^2=7.389$、$\sin1=0.841$）。
\bkjans{$g'=12t^2-2e^{2t}-3\sin t$。$\Bigl[\,t^4-\dfrac{1}{2}e^{2t}+3\sin t\,\Bigr]_0^1=1-\dfrac{e^2}{2}+3\sin1+\dfrac{1}{2}\approx0.33$}
\end{jibun}

\begin{reidai}[たおれる動き]
少しかたむいた棒のかたむき $\theta$（rad）が、$\ddot{\theta}=4\theta$ にしたがって大きくなります。$\theta(0)=0.01$、$\dot{\theta}(0)=0$ のとき、$\theta$ の式と、2秒後のかたむきを求めましょう（$e^4=54.6$、$e^{-4}=0.018$）。\par
\kangaekata 「2回微分で自分の4倍」なので、②の表で $a^2=4$、$a=2$。答えの形は $Ae^{2t}+Be^{-2t}$。$A$、$B$ ははじめの状態で決める（③）。\par
\shiki $\theta=Ae^{2t}+Be^{-2t}$。$t=0$ で $e^0=1$ なので $\theta(0)=A+B=0.01$\par
\hspace*{2.6\zw}$\dot{\theta}=2Ae^{2t}-2Be^{-2t}$ より $\dot{\theta}(0)=2A-2B=0$、つまり $A=B$\par
\hspace*{2.6\zw}$A+B=0.01$、$A=B$ より $A=B=0.005$\par
\hspace*{2.6\zw}2秒後：$0.005\times(e^4+e^{-4})=0.005\times(54.6+0.018)\approx0.273$\par
\kotae $\theta=0.005(e^{2t}+e^{-2t})$、約0.27\,rad（はじめの27倍）
\end{reidai}
$e^{-2t}$ の部分はすぐに小さくなり、$e^{2t}$ の部分がどんどん大きくなります。はじめのかたむきがほんの少しでも、たおれる動きは止まらないことが分かります。

\begin{jibun}
$\ddot{x}=9x$、$x(0)=0$、$\dot{x}(0)=0.6$ をみたす $x$ を求め、$t=1$ での値を求めましょう（$e^3=20.09$、$e^{-3}=0.0498$）。
\bkjans{$x=Ae^{3t}+Be^{-3t}$。$A+B=0$、$3A-3B=0.6$ より $A=0.1$、$B=-0.1$。$x=0.1(e^{3t}-e^{-3t})$、$t=1$ で約2.00}
\end{jibun}

%% ---------------- まちがえやすいところ・物理工学 ----------------
\bksection{tsumazuki}{まちがえやすいところをチェック}

\begin{tsumazuki}{$\omega$ と $\omega^2$ を取りちがえる}
\bkng{$\ddot{x}=-9x$ の答えは $\sin9t$}
\bkok{$\ddot{x}=-9x$ の答えは $\sin3t$}
2回微分で前に出るのは $\omega^2$ なので、$\omega^2=9$ から $\omega=3$ です。$\sin9t$ を2回微分すると $-81\sin9t$ で、$-9$ 倍になりません。
\end{tsumazuki}

\begin{tsumazuki}{はじめの状態を、とちゅうで使う}
$C$ を決めるのは $t=0$ の値です。$T=20+Ce^{-0.05t}$ で $T(0)=80$ なら、$e^0=1$ なので $20+C=80$。$t$ に0を入れてから $C$ を求めます。
\end{tsumazuki}

\begin{tsumazuki}{2回微分の式で、数を1つしか決めない}
$\ddot{x}=-4x$ の答えの形は $A\sin2t+B\cos2t$ で、数が2つあります。はじめの位置だけでは $B$ しか決まりません。はじめの速度も使って $A$ を決めます。
\end{tsumazuki}

\begin{butsuri}
\begin{itemize}[leftmargin=1.2\zw, itemsep=1pt, topsep=0pt, label={\color{bkPurple}●}]
\item \textbf{ふえる・へる（$e^{at}$）}\par\noindent 冷めるお湯、コンデンサーの放電、薬の濃度、人口。答えはいつも $Ce^{at}$ の形（6-4）。
\item \textbf{ゆれる（$\sin$・$\cos$）}\par\noindent ばね、振り子、コイルとコンデンサーをつないだ回路（LC回路）の電気の振動、建物のゆれ（6-6）。
\item \textbf{ゆれながら止まる（$e^{at}$ と $\sin$ のかけ算）}\par\noindent 計器の針、車のサスペンション（6-7・6-12）。2つの候補をかけ算した形が答えになります。
\item \textbf{制御工学（第9巻）}\par\noindent ラプラス変換の表（9-6）にならぶのも、この章の表と同じ関数です。
\end{itemize}
\end{butsuri}

%% ---------------- 章末問題 ----------------
\bkendproblems
\bkproblemgroup{基本}{3}

\bkproblem 次のなぞなぞに合う関数を、1つずつ答えましょう。\newline
（1）微分すると、自分の $-4$ 倍\quad（2）2回微分すると、自分の $-25$ 倍\newline
（3）微分すると、いつも $-3$
\bkpans{（1）\nobreak $Ce^{-4t}$\quad（2）\nobreak $A\sin5t+B\cos5t$\quad（3）\nobreak $-3t+C$}

\bkproblem $f'=-0.2f$、$f(0)=50$ をみたす関数を求め、$t=5$ での値を求めましょう（$e^{-1}=0.368$）。
\bkpans{$f=50e^{-0.2t}$。$t=5$ で $50e^{-1}\approx18.4$}

\bkproblem $\ddot{x}=-25x$、$x(0)=0.05$、$\dot{x}(0)=0$ をみたす $x$ を求めましょう。周期も求めましょう。
\bkpans{$x=0.05\cos5t$。周期 $\dfrac{2\pi}{5}\approx1.26$}

\bkproblem $y=e^{3t}-2e^{-3t}$ が $y''=9y$ をみたすことを確かめましょう。
\bkpans{$y''=9e^{3t}-18e^{-3t}=9(e^{3t}-2e^{-3t})=9y$}

\bkproblem $h(t)=2t^3-3e^{2t}+4\sin t$ を微分しましょう。また、$\displaystyle\int_0^{\pi}4\sin t\,\dd t$ を求めましょう。
\bkpans{$h'=6t^2-6e^{2t}+4\cos t$。$\Bigl[\,-4\cos t\,\Bigr]_0^{\pi}=4+4=8$}

\bkproblemgroup{チャレンジ}{3}

\bkproblem 部屋の温度22\,℃で、$T'=-0.04(T-22)$（$t$ は分）にしたがって冷めるコーヒーがあります。はじめ82\,℃です。$T$ の式を求め、25分後の温度と、温度が50\,℃になる時刻を求めましょう（$e^{-1}=0.368$、$\ln\frac{15}{7}=0.762$）。
\bkpans{$T=22+60e^{-0.04t}$。25分後 $22+60\times0.368\approx44\,℃$。$60e^{-0.04t}=28$ より \mbox{$t=\dfrac{\ln\frac{60}{28}}{0.04}=\dfrac{0.762}{0.04}\approx19$\,分後}}

\bkproblem $\ddot{x}=-4x$、$x(0)=0.3$、$\dot{x}(0)=0.8$ をみたす $x$ を求めましょう。
\bkpans{$x=A\sin2t+B\cos2t$。$B=0.3$、$2A=0.8$ より $A=0.4$。$x=0.4\sin2t+0.3\cos2t$}

\bkproblem 部屋の温度20\,℃の部屋に、5\,℃の冷たい飲み物をおきました。温度 $T$ は $T'=-0.1(T-20)$（$t$ は分）にしたがってあたたまります。$T$ の式と、10分後の温度を求めましょう（$e^{-1}=0.368$）。
\bkpans{$T=20+Ce^{-0.1t}$。$t=0$ で $20+C=5$ より $C=-15$。$T=20-15e^{-0.1t}$。10分後 $20-15\times0.368\approx14.5\,℃$}

\bkproblem $\ddot{x}=x$、$x(0)=2$、$\dot{x}(0)=0$ をみたす $x$ を求めましょう。
\bkpans{$x=Ae^{t}+Be^{-t}$。$A+B=2$、$A-B=0$ より $A=B=1$。$x=e^t+e^{-t}$}

\begin{furikaeri}
\begin{checklist}
\item ふるまい表を使って、微分・積分できる。\hfill{\footnotesize\sffamily\color{bkTheme}→\pageref{bk:5-18:hyo}ページ}
\item 「自分にもどる」関数を答えられる。\hfill{\footnotesize\sffamily\color{bkTheme}→\pageref{bk:5-18:modoru}ページ}
\item はじめの状態から、式の中の数を決められる。\hfill{\footnotesize\sffamily\color{bkTheme}→\pageref{bk:5-18:shoki}ページ}
\end{checklist}
\end{furikaeri}

\bkfukidashi{delta}{第5巻、おつかれさま！　答えの候補のカードはそろった。第6巻では、これを使って、ものの動きを式で求めよう。}


%% ======================== 巻末 ========================
%% 第5巻 巻末（こたえ・記号カード・定数と単位・さくいん・奥付）
\makeatletter
%% ---------------- こたえ ----------------
\bkAnswerPages

%% ---------------- 記号カード ----------------
\bkKigoPages{%
\bkkigorow{$\Delta$}{デルタ}{変化量。「あと$-$まえ」}{2-3}%
\bkkigorow{$f(x)$}{エフ・エックス}{関数 $f$ に $x$ を入れた値}{3-3}%
\bkkigorow{$\theta$}{シータ}{角の大きさ（ラジアン）}{4-2}%
\bkkigorow{$\omega$}{オメガ}{角速度。1秒に回る角（rad/s）}{4-8}%
}

%% 記号カードの2ページ目：ギリシャ文字（第6回：一覧のつづきのすぐあとに置く）
\def\bkbkhead#1{\par\addvspace{4mm}\noindent{\sffamily\bfseries\large\color{bkTheme}#1}\par\nobreak\vspace{2mm}}
\bkbkhead{ギリシャ文字の読み方}
{\small 工学の本では、ギリシャ文字がたくさん使われます。大文字・小文字・読み方と、よく使われる意味です。}\par\vspace{2mm}
\begin{bkcenter}
\footnotesize
\begin{tblr}{colspec={Q[c,m,wd=7mm] Q[c,m,wd=7mm] Q[l,m,wd=17mm] Q[l,m,wd=30mm] Q[c,m,wd=7mm] Q[c,m,wd=7mm] Q[l,m,wd=17mm] Q[l,m,wd=30mm]}, row{1}={bg=bkVol!14, font=\sffamily\bfseries\footnotesize}, hlines={0.4pt,bkLine}, vline{5}={0.8pt,bkVol}, rowsep=0.8pt, colsep=2.5pt}
大 & 小 & 読み方 & よく使う意味 & 大 & 小 & 読み方 & よく使う意味 \\
$\mathrm{A}$ & $\alpha$ & アルファ & 角、角加速度 & $\mathrm{N}$ & $\nu$ & ニュー & 振動数 \\
$\mathrm{B}$ & $\beta$ & ベータ & 角 & $\Xi$ & $\xi$ & クシー & （変数） \\
$\Gamma$ & $\gamma$ & ガンマ & （係数） & $\mathrm{O}$ & $o$ & オミクロン & ― \\
$\Delta$ & $\delta$ & デルタ & 変化量、小さな量 & $\Pi$ & $\pi$ & パイ & 円周率 \\
$\mathrm{E}$ & $\varepsilon$ & イプシロン & 小さな数、誘電率 & $\mathrm{P}$ & $\rho$ & ロー & 密度 \\
$\mathrm{Z}$ & $\zeta$ & ゼータ & 減衰の割合 & $\Sigma$ & $\sigma$ & シグマ & 和、ばらつき \\
$\mathrm{H}$ & $\eta$ & イータ & 効率 & $\mathrm{T}$ & $\tau$ & タウ & 時定数、トルク \\
$\Theta$ & $\theta$ & シータ & 角 & $\Upsilon$ & $\upsilon$ & ウプシロン & ― \\
$\mathrm{I}$ & $\iota$ & イオタ & ― & $\Phi$ & $\phi$ & ファイ & 位相、磁束 \\
$\mathrm{K}$ & $\kappa$ & カッパ & （係数） & $\mathrm{X}$ & $\chi$ & カイ & ― \\
$\Lambda$ & $\lambda$ & ラムダ & 波長 & $\Psi$ & $\psi$ & プサイ & （角） \\
$\mathrm{M}$ & $\mu$ & ミュー & 摩擦係数、マイクロ & $\Omega$ & $\omega$ & オメガ & 角速度（$\Omega$ は抵抗の単位） \\
\end{tblr}
\end{bkcenter}
{\footnotesize 大文字が英語のアルファベットと同じ形の文字（A、B、E など）は、記号としてはほとんど使いません。$\varepsilon$ は $\epsilon$ と書くこともあります。}

\bkbkhead{大きな数・小さな数を表す記号（接頭語）}
\begin{bkcenter}
\small
\begin{tblr}{colspec={Q[l,m,wd=14mm] *{6}{Q[c,m,wd=18mm]}}, column{1}={bg=bkVol!14, font=\sffamily\bfseries\footnotesize}, hlines={0.4pt,bkLine}, rowsep=1.2pt, colsep=2pt}
記号 & G & M & k & m & $\upmu$ & n \\
読み方 & ギガ & メガ & キロ & ミリ & マイクロ & ナノ \\
倍 & $10^9$ & $10^6$ & $10^3$ & $10^{-3}$ & $10^{-6}$ & $10^{-9}$ \\
例 & GHz & MW & km、kWh & mm、mA & $\upmu$m、$\upmu$F & nm \\
\end{tblr}
\end{bkcenter}

%% 第6回：記号カードの3ページ目。ノートに書くとき、ほかの字とまぎれやすい記号
\clearpage
\bkbkhead{記号を書くときのちゅうい}
{\small この本では、問題をノートに解きます。手で書くと、ほかの字とまぎれやすい記号があります。書くときのコツです。}\par\vspace{2mm}
\begin{bkcenter}
\begin{tblr}{colspec={Q[c,m,wd=22mm] Q[l,m,wd=34mm] Q[l,m,wd=68mm]}, row{1}={bg=bkVol!14, font=\sffamily\bfseries\footnotesize}, hlines={0.4pt,bkLine}, rowsep=3pt, cells={font=\small}}
記号 & まぎれやすいもの & 書くときのコツ \\
$\displaystyle\int_a^b$ & S、$f$ & 上下に長い、細い S の形。$a$（下端）は右下、$b$（上端）は右上に小さく書く \\
$\displaystyle\sum_{k=1}^{n}$ & E & 上と下の横線を同じ長さにし、まん中をとがらせる。$k=1$ は下、$n$ は上 \\
$\Delta x$ & 三角形の記号 △、$\Delta\times x$ & $\Delta x$ で1つの記号（「$x$ の変化量」）。$\Delta$ と $x$ をかけるのではない \\
$\dd x$、$\dfrac{\dd y}{\dd x}$ & 文字の $d$ をかけたもの & $\dd x$ もひとまとまり。$\dfrac{\dd y}{\dd x}$ は分数の形をした1つの記号（$d$ で約分しない） \\
$\theta$ & 0、O & たて長の丸の、まん中に横線を1本 \\
$\omega$ & $w$ & 下を丸く2つ。$w$ のようにとがらせない \\
$f'(x)$、$f''(x)$ & 1、$f^1$ & 右かたに小さく短い線。2回微分は2本 \\
$e$ & 文字の $e$ & 2.718… の決まった数。$e^x$ の $e$ は数なので、ほかの文字に変えない \\
$\infty$ & 8、$\propto$ & 8 を横にたおした、閉じた形 \\
$\displaystyle\lim_{n\to\infty}$ & ― & $n\to\infty$ は lim の真下に小さく書く \\
\end{tblr}
\end{bkcenter}


%% ---------------- 定数・単位・公式 ----------------
\bkbacktitle{よく使う定数・単位・公式}
\def\bkbkhead#1{\par\addvspace{4mm}\noindent{\sffamily\bfseries\large\color{bkTheme}#1}\par\nobreak\vspace{2mm}}
\bksubsection{1}{この巻で使う定数}
\begin{bkcenter}
\small
\begin{tblr}{colspec={Q[c,m,wd=14mm] Q[l,m,wd=34mm] Q[l,m,wd=34mm] Q[c,m,wd=14mm]}, row{1}={bg=bkVol!14, font=\sffamily\bfseries\footnotesize}, hlines={0.4pt,bkLine}, rowsep=1.6pt}
記号 & 名前 & 値 & 章 \\
$g$ & 重力加速度 & \qty{9.8}{m/s^2} & 5-5 \\
$\pi$ & 円周率 & 3.14159… & 4-2 \\
$e$ & ネイピア数 & 2.71828… & 5-14 \\
$\ln 2$ & 2の自然対数 & 0.693… & 5-14 \\
$\ln 10$ & 10の自然対数 & 2.303… & 5-14 \\
$\sqrt2$ & 2の平方根 & 1.414… & 2-8 \\
\end{tblr}
\end{bkcenter}

\bksubsection{2}{この巻で使う単位}
\begin{bkcenter}
\small
\begin{tblr}{colspec={Q[l,m,wd=27mm] Q[l,m,wd=27mm] Q[l,m,wd=56.5mm] Q[c,m,wd=12mm]}, row{1}={bg=bkVol!14, font=\sffamily\bfseries\footnotesize}, hlines={0.4pt,bkLine}, rowsep=1.6pt}
量 & 単位 & 意味・ほかの単位との関係 & 章 \\
速度 & m/s & 1秒に進む距離。$\qty{1}{m/s}=\qty{3.6}{km/h}$ & 5-1 \\
加速度 & \unit{m/s^2} & 1秒に変わる速度 & 5-5 \\
力 & N（ニュートン） & $\qty{1}{N}=\qty{1}{kg.m/s^2}$ & 5-12 \\
仕事・エネルギー & J（ジュール） & $\qty{1}{J}=\qty{1}{N.m}$ & 5-12 \\
仕事率・電力 & W（ワット） & $\qty{1}{W}=\qty{1}{J/s}$ & 5-12 \\
電力量 & kWh & $\qty{1}{kWh}=\qty{3.6e6}{J}$ & 5-12 \\
電気の量 & C（クーロン） & \qty{1}{A}の電流が1秒流れたときの量 & 5-12 \\
流量 & L/s & 1秒に流れる水の量 & 5-12 \\
角 & rad（ラジアン） & 弧の長さ÷半径。$180^\circ=\pi$\,rad & 4-2 \\
角速度 & rad/s & 1秒に回る角 & 5-15 \\
\end{tblr}
\end{bkcenter}


\clearpage
\bksubsection{3}{微分と積分の表}
\begin{bkcenter}
\small
\begin{tblr}{colspec={Q[c,m,wd=26mm] Q[c,m,wd=32mm] Q[c,m,wd=44mm] Q[c,m,wd=18mm]}, row{1}={bg=bkVol!14, font=\sffamily\bfseries\footnotesize}, hlines={0.4pt,bkLine}, rowsep=2pt}
関数 & 微分すると & 積分すると（$+C$ を省略） & 章 \\
$c$（定数） & $0$ & $cx$ & 5-4 \\
$x^n$ & $nx^{n-1}$ & $\dfrac{x^{n+1}}{n+1}$（$n\neq-1$） & 5-4・5-11 \\
$e^{ax}$ & $ae^{ax}$ & $\dfrac{1}{a}e^{ax}$ & 5-14 \\
$\dfrac{1}{x}$ & $-\dfrac{1}{x^2}$ & $\ln x$（$x>0$） & 5-14 \\
$\ln x$ & $\dfrac{1}{x}$ & $x\ln x-x$ & 5-14・5-16 \\
$\sin\omega x$ & $\omega\cos\omega x$ & $-\dfrac{1}{\omega}\cos\omega x$ & 5-15 \\
$\cos\omega x$ & $-\omega\sin\omega x$ & $\dfrac{1}{\omega}\sin\omega x$ & 5-15 \\
\end{tblr}
\end{bkcenter}

\bksubsection{4}{計算のきまり}
\begin{bkcenter}
\small
\begin{tblr}{colspec={Q[l,m,wd=28mm] Q[l,m,wd=82.5mm] Q[c,m,wd=16mm]}, row{1}={bg=bkVol!14, font=\sffamily\bfseries\footnotesize}, hlines={0.4pt,bkLine}, rowsep=2pt, column{1}={font=\sffamily}}
きまり & 式 & 章 \\
積の微分 & $(fg)'=f'g+fg'$ & 5-6 \\
合成関数の微分 & $\diff{y}{x}=\diff{y}{u}\cdot\diff{u}{x}$ & 5-6 \\
1次近似 & $f(a+h)\approx f(a)+f'(a)\,h$ & 5-8 \\
微積分の基本定理 & $F'=f$ なら $\displaystyle\int_a^bf\,\dd x=F(b)-F(a)$ & 5-11 \\
置換積分 & $u=g(x)$ とおき、$\dd u=g'(x)\,\dd x$ で書きかえる。定積分は、はんいも $u$ に & 5-16 \\
部分積分 & $\displaystyle\int fg'\,\dd x=fg-\int f'g\,\dd x$ & 5-16 \\
2倍角から & $\sin^2\theta=\dfrac{1-\cos2\theta}{2}$、\quad $\cos^2\theta=\dfrac{1+\cos2\theta}{2}$ & 4-9・5-15 \\
テイラー展開 & $f(x)=f(0)+f'(0)\,x+\dfrac{f''(0)}{2!}x^2+\cdots$ & 5-17 \\
 & $e^x=1+x+\dfrac{x^2}{2!}+\cdots$、\ $\sin x=x-\dfrac{x^3}{3!}+\cdots$、\ $\cos x=1-\dfrac{x^2}{2!}+\cdots$ & 5-17 \\
\end{tblr}
\end{bkcenter}

\clearpage
\bksubsection{5}{よく出る関数のグラフ}
{\small この巻で微分・積分した関数のグラフの形です。グラフの形を思いうかべると、微分の符号（ふえる・へる）や、積分の値のおよそがわかります。}\par\vspace{2mm}
\begin{bkcenter}
\def\bkgr#1#2#3{\begin{minipage}[t]{0.32\linewidth}\centering
\begin{tikzpicture}
\begin{axis}[bkplot, width=46mm, height=38mm, #1]
#2
\end{axis}
\end{tikzpicture}\par\vspace{1mm}
{\footnotesize #3}
\end{minipage}}
\bkgr{xmin=-2.2, xmax=2.2, ymin=-1, ymax=4.5, xtick={-2,0,2}, ytick={2,4}}{\addplot[bkcurve, domain=-2.1:2.1, samples=60] {x^2};}{$y=x^2$\quad 微分 $2x$\\$x=0$ で傾き0（谷）}\hfill
\bkgr{xmin=-1.8, xmax=1.8, ymin=-4, ymax=4, xtick={-1,1}, ytick={-2,2}}{\addplot[bkcurve, domain=-1.58:1.58, samples=60] {x^3};}{$y=x^3$\quad 微分 $3x^2$\\いつも傾きは0以上}\hfill
\bkgr{xmin=0, xmax=4.2, ymin=0, ymax=4.5, xtick={1,2,3,4}, ytick={1,2,3,4}}{\addplot[bkcurve, domain=0.23:4.1, samples=60] {1/x};}{$y=\dfrac{1}{x}$（$x>0$）\quad 微分 $-\dfrac{1}{x^2}$\\積分すると $\ln x$}

\vspace{4mm}
\bkgr{xmin=-2.2, xmax=2.2, ymin=0, ymax=7.5, xtick={-2,0,2}, ytick={1,4}}{\addplot[bkcurve, domain=-2.1:2, samples=60] {exp(x)}; \addplot[bkcurve2, domain=-2:2.1, samples=60] {exp(-x)};}{$y=e^x$（青）・$e^{-x}$\mbox{（オレンジ）}\\微分しても形が変わらない}\hfill
\bkgr{xmin=0, xmax=6.2, ymin=-2.2, ymax=2.2, xtick={1,2.718}, xticklabels={1,$e$}, ytick={-1,1}}{\addplot[bkcurve, domain=0.12:6, samples=80] {ln(x)};}{$y=\ln x$\quad 微分 $\dfrac{1}{x}$\\$(1,\,0)$ を通り、ゆっくりふえる}\hfill
\bkgr{xmin=0, xmax=6.6, ymin=-1.3, ymax=1.3, xtick={1.5708,3.1416,6.2832}, xticklabels={$\frac{\pi}{2}$,$\pi$,$2\pi$}, ytick={-1,1}}{\addplot[bkcurve, domain=0:6.4, samples=80] {sin(deg(x))}; \addplot[bkcurve2, domain=0:6.4, samples=80] {cos(deg(x))};}{$y=\sin x$（青）・$\cos x$\mbox{（オレンジ）}\\微分すると $\frac{1}{4}$ 周期進む}
\end{bkcenter}

\vspace{4mm}
\begin{bkcenter}
\small
\begin{tblr}{colspec={Q[l,m,wd=40mm] Q[l,m,wd=88mm]}, row{1}={bg=bkVol!14, font=\sffamily\bfseries\footnotesize}, hlines={0.4pt,bkLine}, rowsep=1.6pt}
グラフのようす & 微分・積分では \\
右上がり（ふえる） & 微分した値（傾き）が正 \\
右下がり（へる） & 微分した値が負 \\
山・谷（いちばん高い・低い所） & 微分した値が0（前後で符号が変わる） \\
下に凸・上に凸 & 2回微分した値が正・負 \\
横軸より上・下 & 定積分（面積）が正・負 \\
\end{tblr}
\end{bkcenter}

\clearpage
\bksubsection{6}{表計算の関数}
{\small この巻の「やってみよう」で使った、表計算ソフトの関数です。角はラジアンで入れます。}\par\vspace{2mm}
\begin{bkcenter}
\small
\begin{tblr}{colspec={Q[l,m,wd=40mm] Q[l,m,wd=53.5mm] Q[l,m,wd=32mm]}, row{1}={bg=bkVol!14, font=\sffamily\bfseries\footnotesize}, hlines={0.4pt,bkLine}, rowsep=2pt}
書き方 & 意味 & 例 \\
\texttt{=A2\textasciicircum 3} & A2の3乗 & 2なら8 \\
\texttt{=SQRT(A2)} & A2の平方根 & 9なら3 \\
\texttt{=EXP(A2)} & $e$ のA2乗 & 1なら2.71828… \\
\texttt{=LN(A2)} & A2の自然対数 $\ln$ & 1なら0 \\
\texttt{=SIN(A2)}\quad\texttt{=COS(A2)} & $\sin$、$\cos$（A2はラジアン） & 0なら0と1 \\
\texttt{=PI()} & 円周率 $\pi$ & 3.14159… \\
\texttt{=RADIANS(A2)} & 度をラジアンに直す & 180なら3.14159… \\
\texttt{=SUM(B2:B7)} & B2からB7までの合計（$\Sigma$） & 面積のたし算（5-10） \\
\texttt{=B2+(C2*0.1)} & となりのセルを使って計算 & オイラー法（5-8・5-13） \\
\end{tblr}
\end{bkcenter}

\bksubsection{7}{表計算で計算するときの注意}
\begin{itemize}[leftmargin=1.4\zw, itemsep=2pt, topsep=2pt, label={\color{bkTheme}●}]
\item 角は、度ではなくラジアンで入れます。度の値は \texttt{RADIANS} でなおします（5-15）。
\item $h$ を小さくしすぎると、計算機のけた数の限界で、かえって答えがずれることがあります。$h=0.001$ くらいが目安です（5-3）。
\item 式をとなりのセルにコピーすると、セルの番地がいっしょにずれます。ずらしたくない番地は \texttt{\$B\$1} のように \texttt{\$} をつけます。
\item 表計算の答えは、いつも「近似値」です。式で求めた答えとくらべて、どれくらい合っているかを確かめるのが、この本のやり方です。
\end{itemize}

%% ---------------- チェック表 ----------------
\bkbacktitle{できるようになったことチェック表}
{\small 各章の「ふりかえり」を集めた表です。章を学び終えたら1回め、次の巻に進む前に2回めのチェックをしよう。できなかった所は、ページにもどって読みなおします。}\par\vspace{2mm}
%% 自動生成（tools/mkchecklist.py）。各章の「ふりかえり」の項目を集めたもの
\begingroup\footnotesize
\def\bkcb{\tikz[baseline=-0.5ex]\draw[bkTheme, line width=0.7pt, rounded corners=0.5pt] (0,-0.3em) rectangle (0.8em,0.5em);}
\begin{longtblr}{colspec={Q[c,m,wd=9mm] X[l,m] Q[c,m,wd=10mm] Q[c,m,wd=9mm] Q[c,m,wd=9mm]}, width=\linewidth, rowhead=1, colsep=3pt,
  row{1}={bg=bkVol!14, font=\sffamily\bfseries\scriptsize}, hlines={0.4pt,bkLine}, rowsep=0.7pt}
章 & できるようになったこと & ページ & 1回め & 2回め \\
{\sffamily\bfseries\color{bkTheme}5-1} & 平均の速さを $\tfrac{\Delta x}{\Delta t}$ で求められる。 & \pageref{bk:5-1:heikin} & \bkcb & \bkcb \\
 & 平均の速さが、2点を結ぶ直線の傾きだと説明できる。 & \pageref{bk:5-1:katamuki} & \bkcb & \bkcb \\
 & あいだを短くして、瞬間の速さを見積もれる。 & \pageref{bk:5-1:shunkan} & \bkcb & \bkcb \\
{\sffamily\bfseries\color{bkTheme}5-2} & $\Delta t$ を小さくして、平均の速さを表で計算できる。 & \pageref{bk:5-2:hyou} & \bkcb & \bkcb \\
 & 平均の速さを $h$ の式で表せる。 & \pageref{bk:5-2:h} & \bkcb & \bkcb \\
 & $x=at^2$ の瞬間の速さを求められる。 & \pageref{bk:5-2:a} & \bkcb & \bkcb \\
{\sffamily\bfseries\color{bkTheme}5-3} & $\lim$ の記号を読み、極限の値を求められる。 & \pageref{bk:5-3:lim} & \bkcb & \bkcb \\
 & 微分係数を、接線の傾きとして求められる。 & \pageref{bk:5-3:bibunkeisu} & \bkcb & \bkcb \\
 & 導関数を求め、接線の式を書ける。 & \pageref{bk:5-3:shiki} & \bkcb & \bkcb \\
{\sffamily\bfseries\color{bkTheme}5-4} & 表の結果から、微分のきまりを見つけられる。 & \pageref{bk:5-4:kimari0} & \bkcb & \bkcb \\
 & $(x^n)'=nx^{n-1}$ を使って微分できる。 & \pageref{bk:5-4:kimari} & \bkcb & \bkcb \\
 & 定数倍と和のきまりで、多項式を微分できる。 & \pageref{bk:5-4:wa} & \bkcb & \bkcb \\
{\sffamily\bfseries\color{bkTheme}5-5} & 速度と速さのちがいを説明できる。 & \pageref{bk:5-5:sokudo} & \bkcb & \bkcb \\
 & 位置を2回微分して、加速度を求められる。 & \pageref{bk:5-5:nikai} & \bkcb & \bkcb \\
 & $\dot{x}$、$\ddot{x}$、$\tfrac{\dd^2x}{\dd t^2}$ を読める。 & \pageref{bk:5-5:nikai} & \bkcb & \bkcb \\
{\sffamily\bfseries\color{bkTheme}5-6} & 積の微分を、長方形の面積で説明できる。 & \pageref{bk:5-6:seki} & \bkcb & \bkcb \\
 & 合成関数の内側と外側を見分けられる。 & \pageref{bk:5-6:gousei} & \bkcb & \bkcb \\
 & 合成関数の微分を使って計算できる。 & \pageref{bk:5-6:chain} & \bkcb & \bkcb \\
{\sffamily\bfseries\color{bkTheme}5-7} & 導関数の符号から、上がる・下がるを判断できる。 & \pageref{bk:5-7:fugou} & \bkcb & \bkcb \\
 & 増減表をかいて、極大・極小を求められる。 & \pageref{bk:5-7:zougen} & \bkcb & \bkcb \\
 & 範囲のはしもくらべて、最大・最小を求められる。 & \pageref{bk:5-7:saidai} & \bkcb & \bkcb \\
{\sffamily\bfseries\color{bkTheme}5-8} & 次の位置を、今の位置＋速度×$\Delta t$ で予想できる。 & \pageref{bk:5-8:tsugi} & \bkcb & \bkcb \\
 & 接線を使って、1次近似ができる。 & \pageref{bk:5-8:sessen} & \bkcb & \bkcb \\
 & 測定のずれが、計算にひびく大きさを見積もれる。 & \pageref{bk:5-8:gosa} & \bkcb & \bkcb \\
{\sffamily\bfseries\color{bkTheme}5-9} & 速さが変わるときも、細かく分けて、かけて、たして距離を求められる。 & \pageref{bk:5-9:wakeru} & \bkcb & \bkcb \\
 & 長方形・三角形・台形の面積で、進んだ距離がわかる。 & \pageref{bk:5-9:sankaku} & \bkcb & \bkcb \\
 & マイナスの面積の意味（変位と道のり）を説明できる。 & \pageref{bk:5-9:minus} & \bkcb & \bkcb \\
{\sffamily\bfseries\color{bkTheme}5-10} & 細い長方形の和で、面積を近似できる。 & \pageref{bk:5-10:hasami} & \bkcb & \bkcb \\
 & $\Sigma$ の記号を読み書きできる。 & \pageref{bk:5-10:sigma} & \bkcb & \bkcb \\
 & 定積分の意味を説明し、$x^2$ の定積分を求められる。 & \pageref{bk:5-10:teisekibun} & \bkcb & \bkcb \\
{\sffamily\bfseries\color{bkTheme}5-11} & 面積のふえる速さが高さであることを説明できる。 & \pageref{bk:5-11:kihon} & \bkcb & \bkcb \\
 & 不定積分を $+C$ をつけて書ける。 & \pageref{bk:5-11:genshi} & \bkcb & \bkcb \\
 & 定積分を $F(b)-F(a)$ で計算できる。 & \pageref{bk:5-11:teisekibun} & \bkcb & \bkcb \\
{\sffamily\bfseries\color{bkTheme}5-12} & 変わる力の仕事を、積分で求められる。 & \pageref{bk:5-12:sekibun} & \bkcb & \bkcb \\
 & 流量や電力から、たまった量を求められる。 & \pageref{bk:5-12:ryuuryou} & \bkcb & \bkcb \\
 & 積分を使って、平均の値を求められる。 & \pageref{bk:5-12:heikin} & \bkcb & \bkcb \\
{\sffamily\bfseries\color{bkTheme}5-13} & 初期条件から、積分定数を決められる。 & \pageref{bk:5-13:shoki} & \bkcb & \bkcb \\
 & 加速度から、速度と位置を求められる。 & \pageref{bk:5-13:toukasoku} & \bkcb & \bkcb \\
 & 「いまの量＝はじめの量＋変化の積分」を使える。 & \pageref{bk:5-13:teisekibun} & \bkcb & \bkcb \\
{\sffamily\bfseries\color{bkTheme}5-14} & $e$ が $(1+\frac{1}{n})^n$ の近づく先であることを確かめられる。 & \pageref{bk:5-14:n} & \bkcb & \bkcb \\
 & $(e^{ax})'=ae^{ax}$ を使える。 & \pageref{bk:5-14:eax} & \bkcb & \bkcb \\
 & $\ln x$ を微分・積分できる。 & \pageref{bk:5-14:lnbibun} & \bkcb & \bkcb \\
{\sffamily\bfseries\color{bkTheme}5-15} & 円運動の速度で、$\sin$・$\cos$ の微分を説明できる。 & \pageref{bk:5-15:en} & \bkcb & \bkcb \\
 & $\sin\omega t$、$\cos\omega t$ を微分・積分できる。 & \pageref{bk:5-15:omega} & \bkcb & \bkcb \\
 & 単振動の $\ddot{x}=-\omega^2x$ を確かめられる。 & \pageref{bk:5-15:nikai} & \bkcb & \bkcb \\
 & $\sin^2$ を積分して、交流の実効値を求められる。 & \pageref{bk:5-15:jikko} & \bkcb & \bkcb \\
{\sffamily\bfseries\color{bkTheme}5-16} & 置換積分で、$\dd u$ への書きかえができる。 & \pageref{bk:5-16:chikan} & \bkcb & \bkcb \\
 & 定積分の置換で、はんいを書きかえられる。 & \pageref{bk:5-16:hanni} & \bkcb & \bkcb \\
 & 部分積分で、$f$ と $g'$ を選んで計算できる。 & \pageref{bk:5-16:bubun} & \bkcb & \bkcb \\
{\sffamily\bfseries\color{bkTheme}5-17} & 2次近似で、曲がり方まで合わせる意味を説明できる。 & \pageref{bk:5-17:niji} & \bkcb & \bkcb \\
 & $e^x$・$\sin x$・$\cos x$ の展開を書き、値を計算できる。 & \pageref{bk:5-17:sin} & \bkcb & \bkcb \\
 & $\sin\theta\approx\theta$ から、振り子の周期を求められる。 & \pageref{bk:5-17:furiko} & \bkcb & \bkcb \\
{\sffamily\bfseries\color{bkTheme}5-18} & ふるまい表を使って、微分・積分できる。 & \pageref{bk:5-18:hyo} & \bkcb & \bkcb \\
 & 「自分にもどる」関数を答えられる。 & \pageref{bk:5-18:modoru} & \bkcb & \bkcb \\
 & はじめの状態から、式の中の数を決められる。 & \pageref{bk:5-18:shoki} & \bkcb & \bkcb \\
\end{longtblr}
\endgroup



%% ---------------- つぎの巻へ（第6回で追加：第6巻の予告と、この巻の道具がこの先どこで使われるか） ----------------
\bkbacktitle{つぎの巻へ}
\newcommand\bknextpill[1]{\bk@pill{#1}}% 第11回：章の番号の札は、その章の巻の色（5-○は緑、6-○は青緑）
\def\bknextrow{1}
\newcommand\bknextcard[4]{%
  \begin{tcolorbox}[enhanced, nobeforeafter, box align=top, width=0.485\linewidth, colframe=bkVol6, colback=white, boxrule=0.9pt, arc=2.5mm,
    left=3mm, right=3mm, top=2mm, bottom=2.5mm, equal height group=bknext\bknextrow]
  \bknextpill{#1}\enspace{\sffamily\bfseries\small #2}\par\vspace{1.5mm}
  {\sffamily\bfseries\small\color{bkVol6!80!black}#3}\par\vspace{2mm}
  {\footnotesize\setlength{\parindent}{0pt}{\sffamily\bfseries\color{bkInk}使う道具}\enspace #4}
  \end{tcolorbox}}
微分と積分の道具がそろいました。第6巻「未来を計算する」では、この2つを組み合わせて、「次の瞬間を決めるルール」（微分方程式）から未来を予想します。第6巻の章をのぞいてみましょう。\par\vspace{2.5mm}
\noindent\bknextcard{6-2}{表計算で解こう：オイラー法}{80 ℃のお湯は、10分後に何℃？}{短い時間なら、かけ算で近似できる（5-8）。それを表計算で何回もくり返す。}\hfill
\bknextcard{6-3}{積分するだけで解ける形}{落ちるボールの位置を、式で求める}{積分は微分の逆（5-11）と、積分定数は「はじめの状態」（5-13）。}\par\vspace{1.8mm}
\def\bknextrow{2}\noindent\bknextcard{6-4}{自分に比例してへる：冷めるお湯}{熱いほど速く冷める。答えは指数関数}{微分しても形が変わらない数 $e$（5-14）。答えが指数関数になる。}\hfill
\bknextcard{6-6}{ばねと振り子}{小さなゆれなら、周期はゆれの大きさによらない}{$\sin\cdot\cos$ の微分（5-15）と、$\sin\theta\approx\theta$ の近似（5-17）。}\par\vspace{1.8mm}
\def\bknextrow{3}\noindent\bknextcard{6-10}{オイラーの公式}{指数関数と三角関数が、1本の式でつながる}{$e^x$、$\sin x$、$\cos x$ のテイラー展開（5-17）をならべてくらべる。}\hfill
\bknextcard{6-13}{ゆらし続けると：強制振動と共振}{小さな力でも、タイミングが合うと大きくゆれる}{いちばん大きくなる所を、微分が0の点で見つける（5-7）。}\par
\bkfukidashi{delta}{第5巻の道具は、この先もずっと使うよ。チェック表で自信のない章を見直してから、第6巻へ進もう！}

\clearpage
\bkbkhead{この巻の道具は、この先でこう使う}
{\small 第5巻の章と、その道具を使う第6巻・第9巻の章です（章のはじめの「道しるべ」の「この章の先に」から集めました）。}\par\vspace{3mm}
\begin{bkcenter}
\begin{tblr}{colspec={Q[l,m,wd=52mm] Q[l,m,wd=80mm]}, row{1}={bg=bkVol!14, font=\sffamily\bfseries\footnotesize}, hlines={0.4pt,bkLine}, rowsep=1.5pt, cells={font=\footnotesize}}
この巻の章 & この道具を使う章 \\
{\bknextpill{5-5}\ 位置→速度→加速度} & {\bknextpill{6-1}\ 微分方程式ってなに？\quad \bknextpill{6-15}\ まとめ：物理は「微分方程式＋はじめの状態」で書ける} \\
{\bknextpill{5-7}\ いちばん高いのはどこ？：最大と最小} & {\bknextpill{6-13}\ ゆらし続けると：強制振動と共振} \\
{\bknextpill{5-8}\ 短い時間なら、かけ算で近似できる} & {\bknextpill{6-2}\ 表計算で解こう：オイラー法} \\
{\bknextpill{5-11}\ 積分は微分の逆} & {\bknextpill{6-3}\ 積分するだけで解ける形} \\
{\bknextpill{5-12}\ 物理での積分：仕事と流量} & {\bknextpill{6-15}\ まとめ：物理は「微分方程式＋はじめの状態」で書ける} \\
{\bknextpill{5-13}\ 積分定数は「はじめの状態」} & {\bknextpill{6-1}\ 微分方程式ってなに？\quad \bknextpill{6-3}\ 積分するだけで解ける形} \\
{\bknextpill{5-14}\ ふしぎな数 e} & {\bknextpill{6-4}\ 自分に比例してへる：冷めるお湯\quad \bknextpill{9-5}\ 微分をかけ算に変える翻訳機：ラプラス変換} \\
{\bknextpill{5-15}\ sin・cosを微分すると？} & {\bknextpill{6-6}\ ばねと振り子\quad \bknextpill{6-10}\ オイラーの公式} \\
{\bknextpill{5-16}\ 積分の道具箱：置換積分と部分積分} & {\bknextpill{6-4}\ 自分に比例してへる：冷めるお湯\quad \bknextpill{9-1}\ 波をたし合わせて形を作る：フーリエ級数\quad \bknextpill{9-5}\ 微分をかけ算に変える翻訳機：ラプラス変換} \\
{\bknextpill{5-17}\ 近くでは多項式で近似できる：テイラー展開} & {\bknextpill{6-6}\ ばねと振り子\quad \bknextpill{6-10}\ オイラーの公式\quad \bknextpill{9-1}\ 波をたし合わせて形を作る：フーリエ級数} \\
{\bknextpill{5-18}\ まとめ：微分したときのふるまい} & {\bknextpill{6-1}\ 微分方程式ってなに？\quad \bknextpill{6-4}\ 自分に比例してへる：冷めるお湯\quad \bknextpill{6-6}\ ばねと振り子} \\
\end{tblr}
\end{bkcenter}

%% ---------------- さくいん ----------------
\bkbacktitle{さくいん}
{\small 数字はページです。その語がはじめて出てくるページや、くわしく説明しているページをのせています。}\par\vspace{2mm}
\begin{multicols}{2}\small\setlength{\parskip}{0pt}
\directlua{bk.printidx()}
\end{multicols}

%% ---------------- 奥付 ----------------
\bkpadfour
\global\bk@decofalse\bkfolio{false}
\null\vfill
\noindent\begin{minipage}{\linewidth}\small
{\sffamily\bfseries\Large\color{bkDeep}物理を理解するための数学　第5巻}\par\vspace{1mm}
{\sffamily\bfseries\large 変化をつかむ　微分と積分}\par\vspace{3mm}
{\color{bkVol}\hrule height 1pt}\vspace{3mm}
\begin{tabular}{@{}l@{\hspace{6mm}}l@{}}
編者 & 「物理を理解するための数学」編集委員会 \\
 & 委員長　小学校長（まとめ役） \\
 & 委員　数学・電気工学・機械工学・教育学の大学教授、幼稚園教諭、\\
 & \phantom{委員　}特別支援学校教諭、心理士、科学解説者、小学校教諭 \\
デザイン & ユニバーサルデザイン担当委員 \\
挿絵 & イラスト担当委員 \\
版 & 見本版（第6回編集委員会で修正）　2026年10月 \\
\end{tabular}\par\vspace{3mm}
{\color{bkVol}\hrule height 1pt}\vspace{2mm}
{\footnotesize この本の問題の答えは、巻末の「こたえ」にあります。くわしい解き方は、別冊の問題集と指導書にのせます。}
\end{minipage}
\makeatother

